% !TEX program = lualatex
\documentclass[lualatex,ja=standard,enablejfam=true,base=11pt]{bxjsbook}

% ===============================
% ===============================
% ---------- SETTINGS ----------
% ===============================
% ===============================

% ===================
% Typesetting setup
% ===================

% Load the typesetting
% typesetting.tex
% =====================
% Engine: LuaLaTeX + Japanese
% =====================

\usepackage{luatexja} % Japanese support

% treat all Japanese chars uniformly}
\ltjsetparameter{jacharrange={-2}} 

% --- 禁則処理 ---
\ltjsetparameter{prebreakpenalty={`ぁ,10000}}
\ltjsetparameter{prebreakpenalty={`ぃ,10000}}
\ltjsetparameter{prebreakpenalty={`ぅ,10000}}
\ltjsetparameter{prebreakpenalty={`ぇ,10000}}
\ltjsetparameter{prebreakpenalty={`ぉ,10000}}
\ltjsetparameter{prebreakpenalty={`ゃ,10000}}
\ltjsetparameter{prebreakpenalty={`ゅ,10000}}
\ltjsetparameter{prebreakpenalty={`ょ,10000}}
\ltjsetparameter{prebreakpenalty={`っ,10000}}
\ltjsetparameter{prebreakpenalty={`ァ,10000}}
\ltjsetparameter{prebreakpenalty={`ィ,10000}}
\ltjsetparameter{prebreakpenalty={`ゥ,10000}}
\ltjsetparameter{prebreakpenalty={`ェ,10000}}
\ltjsetparameter{prebreakpenalty={`ォ,10000}}
\ltjsetparameter{prebreakpenalty={`ャ,10000}}
\ltjsetparameter{prebreakpenalty={`ュ,10000}}
\ltjsetparameter{prebreakpenalty={`ョ,10000}}
\ltjsetparameter{prebreakpenalty={`ッ,10000}}

% --- punctuation ---
\ltjsetparameter{prebreakpenalty={`、,10000}} % U+3001
\ltjsetparameter{prebreakpenalty={`。,10000}} % U+3002
\ltjsetparameter{prebreakpenalty={`，,10000}} % U+FF0C
\ltjsetparameter{prebreakpenalty={`．,10000}} % U+FF0E

% --- closing brackets ---
\ltjsetparameter{prebreakpenalty={`）,10000}} % U+FF09
\ltjsetparameter{prebreakpenalty={`〕,10000}} % U+3015
\ltjsetparameter{prebreakpenalty={`］,10000}} % U+FF3D
\ltjsetparameter{prebreakpenalty={`｝,10000}} % U+FF5D
\ltjsetparameter{prebreakpenalty={`〉,10000}} % U+3009
\ltjsetparameter{prebreakpenalty={`》,10000}} % U+300B
\ltjsetparameter{prebreakpenalty={`」,10000}} % U+300D
\ltjsetparameter{prebreakpenalty={`』,10000}} % U+300F
\ltjsetparameter{prebreakpenalty={`】,10000}} % U+3011

\ltjsetparameter{
  kanjiskip  = 0pt plus 0.05em minus 0.03em,
  xkanjiskip = 0.12em plus 0.06em minus 0.03em
}
\emergencystretch=2em


\usepackage{luatexja-fontspec} % fontspec style interface
% !TEX program = lualatex

% ===================
% Latin / Roman font
% ===================
\setmainfont[
	Ligatures        = TeX,
	Numbers          = OldStyle,
	Scale            = MatchLowercase,
	Kerning          = On
]{TeX Gyre Pagella}

% ==============
% Japanese fonts
% ==============
\setmainjfont[
	Renderer = HarfBuzz,
	BoldFont = {Noto Serif CJK JP Bold},
	YokoFeatures = {JFM=prop},
	TateFeatures = {JFM=prop}
]{Noto Serif CJK JP Regular}

\newjfontfamily\EmbJpFont[
	Renderer = HarfBuzz,
	BoldFont = {Noto Sans CJK JP Bold},
	YokoFeatures = {JFM=prop},
	TateFeatures = {JFM=prop}
]{Noto Sans CJK JP Regular}

% =========
% IPA font
% =========
\newfontfamily\ipafont[
	Scale            = MatchLowercase,
	Ligatures        = TeX,
	Kerning          = On
]{Charis SIL}

% Define font for pitch variation diacritics
\newfontfamily\pitchvar{DejaVu Serif}

% Mathtools
\usepackage{unicode-math}
\setmathfont{TeX Gyre Pagella Math}


% =================
% ---- LAYOUT ----
% =================

\geometry{
	paper=a4paper,
	top=22mm,
	bottom=22mm,
	left=22mm,
	right=22mm,
    headsep=6mm,
	footskip=10mm,
	includehead,
	includefoot
}
\setlength{\headheight}{19pt}
\setlength{\columnsep}{6mm}
\setlength{\columnseprule}{0pt}

% ======================
% ---- Bibliography ----
% ======================

% Load the typesetting
% bib-config.tex

% ===============================
% ---- BiBiography Settings ----
% ===============================

% Biblatex + biber (LuaLaTeX-friendly)
\usepackage[
	backend=biber,
	style=authoryear,
	natbib=true,         % keeps \citep, \citet, etc.
	sorting=nyt,         % name-year-title
	sortlocale={ja_JP}     % Japanese collation in biber (ICU)
	giveninits=false,   % do NOT abbreviate given names
	dashed=false,       % do NOT replace repeated authors with dash
	maxbibnames=99,     % print all author names in bibliography
	maxcitenames=2      % (optional) keep citations short; set 99 if you want full in-text too
]{biblatex}

\addbibresource{main.bib}

% If you want Author:Year (your natbib had a ":"-ish intent)
\DeclareDelimFormat{nameyeardelim}{\addcolon\space}

% If you want (A, 2020, 15) for postnotes:
\DeclareDelimFormat{postnotedelim}{\addcomma\space}

\DeclareNameAlias{default}{family-given}

\DeclareDelimFormat{revsdnamepunct}{\addspace}

% Article-like titles: Japanese -> 「...」, otherwise -> “...”
\DeclareFieldFormat[article,incollection,inproceedings]{title}{%
	\iffieldequalstr{langid}{japanese}
		{「#1」}
		{\mkbibquote{#1}}%
}



% Page style
\AtBeginBibliography{\thispagestyle{plain}}

% ===================
% ---- PACKAGES ---- 
% ===================

% 色々
\usepackage{etoolbox}

% Header/footer customization
\usepackage{fancyhdr} 

% Title formatting and customization
\usepackage{titlesec}

% Package for ruby
\usepackage{pxrubrica}

% Draw Tree package
\usepackage{tikz,tikz-qtree}
\usetikzlibrary{trees} % this is to allow the fork right path
\usetikzlibrary{fit,shapes}

% Graphics, Figures and Tables related packages
\usepackage{graphicx,wrapfig,caption}
\usepackage{booktabs,hhline}
\usepackage{makecell,multirow,multicol,arydshln} % Several lines in cell, multirow/col, dashed lines
\usepackage{placeins}

% Internal navigation
\usepackage[hidelinks,unicode=true]{hyperref}
\usepackage{bookmark}

% Package for enumeration
\usepackage[shortlabels]{enumitem}

% Linguistics related Packages
\usepackage{linguex} % Examples
\renewcommand{\firstrefdash}{} % Suppress first '-' when referring to subexample

% For conditional formatting
\usepackage{xifthen}

% Use of color
\usepackage{color}

% Check layout if needed
%\usepackage{layout}

% ===========================
% ---- TITLE FORMATTING ---- 
% ===========================

% Format Chapter titles
\titleformat{\chapter}[block]{\large\bfseries\filcenter}{章}{0pt}{}
\titlespacing*{\chapter}{0pt}{*0}{*6}

% Removing the chapter for section
\renewcommand{\thesection}{\arabic{section}.}

% Format section titles
\titleformat{\section}[block]{\normalfont\normalsize\bfseries}{\normalfont\thesection}{4pt}{}
\titlespacing{\section}{0pt}{10pt plus 3pt minus 1pt}{*0}

% Format subsection titles
\titleformat{\subsection}[block]{\normalfont\normalsize\bfseries}{\normalfont\thesubsection}{4pt}{}

% Table of contents print level
\setcounter{tocdepth}{1}

% Level of numbering
\setcounter{secnumdepth}{4}

% Caption formatting
\DeclareCaptionFormat{custom}{#1 #3}
\captionsetup{within=none,format=custom}

% Indent first paragraph
\usepackage{indentfirst}

% ===================================
% ---- HEADER/FOOTER FORMATTING ---- 
% ===================================

% Compute textwidth for header and footer settings
\usepackage{calc}
\newlength{\DictBlockWidth}

\newcommand{\SetDictBlockWidth}{%
	\setlength{\DictBlockWidth}{\textwidth}
}

% Redefine the plain style
\fancypagestyle{plain}{

    % Adjust size of header and footer to the size of \textwidth
    \fancyheadoffset[LE,RO]{\DictBlockWidth}
    \fancyfootoffset[LE,RO]{\DictBlockWidth}

    \fancyhf{} % Clear header and footer
    \fancyfoot[LE,RO]{\textbf{\small{\textsf{\thepage}}}} 
    \renewcommand{\headrulewidth}{0pt} % no rule under the header
}

% Definition of style for dictionary contents
\fancypagestyle{dictionary}{

    % Adjust size of header and footer to the size of \textwidth
    \fancyheadoffset[LE,RO]{\DictBlockWidth}
    \fancyfootoffset[LE,RO]{\DictBlockWidth}

    \fancyhf{} % Flush header and footer
    \fancyhead[LE]{\large \textsf{\rightmark}} % Even page header
    \fancyhead[RO]{\large \textsf{\leftmark}} % Odd page header
    \renewcommand{\headrulewidth}{1pt} % Rule under header
    
    \fancyfoot[LE,RO]{\textbf{\small{\textsf{\thepage}}}} % Page printed in Footer
    \renewcommand{\footrulewidth}{0pt} % Rule under footer
}

% ==================================================
% ---- DICTIONARY/ORTHOGRAPHY RELATED COMMANDS ----
% ==================================================

% Load the dictionary commands
%\input{settings/d_commands_old}
% d_commands.tex

% =====================================
% ---- Orthographic Commands & Others ----
% =====================================

% Commands for IPA font and IPA embedded in Japanese text
\newcommand\ip[1]{{\ipafont #1}}
\newcommand\ipe[1]{{\space\ipafont #1\space}} % % Embedded with japanese text at both side
\newcommand\ipr[1]{{\ipafont #1\space}} % Embedded with japanese text at right side
\newcommand\ipl[1]{{\space\ipafont #1}} % Embedded with japanese text at left side

% Command for embedded Ryukyuan in Japanese
%\newcommand\emb[1]{{\EmbJpFont\addfontfeatures{LetterSpace=1.0} #1}}
\newcommand\emb[1]{{\EmbJpFont #1}}

% Command for outputting roman number
\makeatletter
\newcommand*{\rom}[1]{\expandafter\@slowromancap\romannumeral #1@類}
\makeatother

% Definition of space for wakachigaki 
% NOTE: never use 全角スペース as such
%\newcommand{\ws}{\nobreak\hskip0.75em plus 0.1em minus 0.08em\relax}
\newcommand{\ws}{\space\kern0.35em}

% Kerneling of 「゜」
%\newcommand{\m}{\nobreak\hspace{-2pt}゜\hspace{-5pt plus 0.5pt minus 0.5pt}}
\newcommand{\m}{\nobreak\kern-2pt゜\kern-5pt}

\newcommand{\sms}{\kern0.05em\raisebox{-0.3ex}{\scalebox{0.7}{ス}}\kern0.05em}
\newcommand{\smz}{\kern0.05em\raisebox{-0.3ex}{\scalebox{0.7}{ズ}}\kern0.05em}

% Definition for the correct display of 「～」
\newcommand{\ntilde}{\nobreak\hspace{1pt plus 0.2pt minus 0.2pt}～\hspace{0.6pt plus 0.1pt minus 0.1pt}}

% Commands for pitch variation diacritics
\newcommand{\josho}{\pitchvar{\hspace{0pt}⸢\hspace{0pt}}}
\newcommand{\kako}{\pitchvar{\hspace{0pt}⸣\hspace{0pt}}}
\newcommand{\heiban}{\ip{\hspace{0pt}\textbf{\nolinebreak ¯}\hspace{0pt}}}

\DeclareRobustCommand{\tl}[1]{\nobreak\kern.02em\hbox{\ip{#1}}}

% Command for ①②③... 
% NOTE: used in MEANING(S) numbering
\newcommand{\numbering}[1]{%
  \ifcase#1\or ①\or ②\or ③\or ④\or ⑤\or ⑥\or ⑦\or ⑧\or ⑨\or ⑩\or ⑪\or ⑫\or ⑬\or ⑭\or ⑮\or ⑯\or ⑰\or ⑱\or ⑲\or ⑳\or ㉑\or ㉒\or ㉓\or ㉔\or ㉕\or ㉖\or ㉗\or ㉘\or ㉙\or ㉚\fi
  \hspace{0.3em}%
}

% Command for letter heading
% Principle: unbreakable vertical box with fixed height

\makeatletter
\newcommand{\LetterHead}[1]{%
  \par\noindent

  % Build an unbreakable box with *exact* spacing (tune these two):
  \setbox0=\vbox{%
    \vskip 1.0\baselineskip
    \hbox to \linewidth{\hfil \Large\bfseries —\,#1\,—\hfil}%
    \vskip 1.0\baselineskip
  }%

  % Decide remaining space (multicols column vs normal page)
  \dimen0=\ht0 \advance\dimen0 by \dp0 % needed height

  \ifdefined\@colroom
    % Inside multicols: remaining space in current column is \@colroom
    \ifdim\@colroom<\dimen0
      \columnbreak
    \fi
  \else
    % Normal single-column: remaining space on page
    \dimen2=\dimexpr\pagegoal-\pagetotal\relax
    \ifdim\dimen2<\dimen0
      \pagebreak
    \fi
  \fi

  % Place the box (unbreakable, predictable)
  \nointerlineskip\box0\par
}
\makeatother

% =====================================
% ---- DICTIONARY ENTRY STRUCTURE ----
% =====================================

\usepackage{keyval}

% New entry command
% 1st param: Entry head
% 2st param: Meanings block
% 3st param: Entry tail
\newcommand{\entry}[3]{%
	\par\begingroup
    \leavevmode
	\hangindent=0.3cm
	#1% 
	%\hspace{0.2em plus 0.1em minus 0.05em}%
	\kern0.08em%
    \ignorespaces#2#3%
    \par\endgroup
}

% Keyed Entry head
\makeatletter

% Slots for Entry head
\newcommand{\ryu@head@kana}{} % Kana orthograpy
\newcommand{\ryu@head@ipa}{} % IPA/Romaji orthograpghy
\newcommand{\ryu@head@pos}{} % Part of Speech
\newcommand{\ryu@head@cat}{} % Semantic Category 
\newcommand{\ryu@head@acc}{} % Accent Class/information
\newcommand{\ryu@head@rui}{} % Conjugation Class
\newcommand{\ryu@head@para}{} % Conjugation Paradigm

% Keys
\define@key{ryuhead}{kana}{\def\ryu@head@kana{#1}}
\define@key{ryuhead}{ipa}{\def\ryu@head@ipa{#1}}
\define@key{ryuhead}{pos} {\def\ryu@head@pos{#1}}
\define@key{ryuhead}{cat}{\def\ryu@head@cat{#1}}
\define@key{ryuhead}{acc}{\def\ryu@head@acc{#1}}
\define@key{ryuhead}{rui} {\def\ryu@head@rui{#1}}
\define@key{ryuhead}{para} {\def\ryu@head@para{#1}}

% Keyed Entry head command
\newcommand{\headPrint}[1][]{%
  % reset every time
  \let\ryu@head@kana\@empty%
  \let\ryu@head@ipa\@empty%
  \let\ryu@head@pos\@empty%
  \let\ryu@head@cat\@empty%
  \let\ryu@head@acc\@empty%
  \let\ryu@head@rui\@empty%
  \let\ryu@head@para\@empty%
  % parse keys if any
  \setkeys{ryuhead}{#1}%
  %
  % Kana orthography + header markings
  \ifx\ryu@head@kana\@empty\else%
    \markboth{\ryu@head@kana}{\ryu@head@kana}%
    %\textbf{\ryu@head@kana}%
    {\EmbJpFont\bfseries \ryu@head@kana}%
  \fi%
  % IPA/Romaji orthography
  \ifx\ryu@head@ipa\@empty\else%
    \enspace{\ip{[\ryu@head@ipa]}}%
  \fi%
  % POS block (print as one unit)
  \ifx\ryu@head@pos\@empty\else%
    \kern-0.1em［\ryu@head@pos%
    \ifx\ryu@head@rui\@empty\else\ip{\ryu@head@rui}\fi%
    ］\kern-0.1em%
  \fi%
  % Accent class
  \ifx\ryu@head@acc\@empty\else%
    \kern-0.1em［\ryu@head@acc］\kern-0.1em%
  \fi%
  % Semantic category
  \ifx\ryu@head@cat\@empty\else%
    \kern-0.1em［\ryu@head@cat］\kern-0.1em%
  \fi%
  % Paradigm
  \ifx\ryu@head@para\@empty\else%
    \kern-0.1em［{\EmbJpFont \ryu@head@para}］\kern-0.1em%
  \fi%  
  }%
\makeatother

% Meaning(s) block
\makeatletter
\newcount\ryu@meaningcount
\newcount\ryu@meaningindex

% 2 passes required for the confitioning of numbering 
% i.e. no numbering if only one meaning in meaning block
\newcommand{\MeaningsBlock}[1]{%
  \begingroup%
    % pass 1: count
    \ryu@meaningcount=0\relax%
    \long\def\ryumeaning##1##2##3{%
      \advance\ryu@meaningcount by 1\relax%
    }%
    #1%
    %
    % pass 2: print
    \ryu@meaningindex=0\relax%
    \long\def\ryumeaning##1##2##3{%
      \ifnum\ryu@meaningcount>1\relax%
        \advance\ryu@meaningindex by 1\relax%
        \numbering{\the\ryu@meaningindex}%
      \fi%
      \MeaningPrint{##1}{##2}{##3}%
    }%
    \ignorespaces#1%
  \endgroup%
}
\makeatother

\makeatletter
% Command for printing the meaning
\newcommand{\MeaningPrint}[3]{%
	\ifblank{#1}{}{#1。}% gloss (相当語・訳語)
	\ifblank{#2}{}{#2。}% Explanation ()
	\ifblank{#3}{}{#3}% Example Sentence
}
\makeatother

% Entry Tail
\makeatletter

\newcommand{\ryu@tail@note}{}
\newcommand{\ryu@tail@allo}{}
\newcommand{\ryu@tail@syn}{}

\define@key{ryutail}{note}{\def\ryu@tail@note{#1}}
\define@key{ryutail}{allo}{\def\ryu@tail@allo{#1}}
\define@key{ryutail}{syn} {\def\ryu@tail@syn{#1}}

\newcommand{\tailPrint}[1][]{%
  % reset every time
  \let\ryu@tail@note\@empty%
  \let\ryu@tail@allo\@empty%
  \let\ryu@tail@syn\@empty%
  %
  % parse keys if any
  \setkeys{ryutail}{#1}%
  %
  % output only if nonempty
  \ifx\ryu@tail@note\@empty\else［備］\ryu@tail@note。\fi%
  \ifx\ryu@tail@allo\@empty\else［同］\ryu@tail@allo。\fi%
  \ifx\ryu@tail@syn \@empty\else［類］\ryu@tail@syn 。\fi%
}

\makeatother

% Command for inline sentence example
\newcommand{\inlineExample}[1]{#1。}

% Command for sentence example
% 1st param: Source language orthography
% 2nd param: Target language (Japanese) translation 
\newcommand{\reibun}[2]{\emb{#1}（#2）。}

% Definition of reverse lookup structure
% 1st para: Reading (読み)
% 2nd para: Kanji Orthography (漢字表記)
% 3rd para: Word-form (語形)

% Structure command for reverse lookup
\newcommand{\rentry}[3]{%
  \par\begingroup
  \leavevmode
  \markboth{\unexpanded{#1}}{\unexpanded{#1}}%
  \hangindent=1em
  \noindent
  \ifblank{#2}%
    {#1\hspace{0.4em}}%
    {#1【#2】\hspace{0.4em}}%
  #3%
  \par\endgroup
}


% =========================
% ---- MISCELLANEOUS ----
% =========================

% Shortcut for superscript text
\let\sp\textsuperscript

% Footnote numeration system from default *1, *2, *3 to 1, 2 ,3
\renewcommand{\thefootnote}{\arabic{footnote}}

% Float-only page layout (概説 tables)
\makeatletter
\setlength{\@fptop}{0pt}
\setlength{\@fpsep}{8pt}
\setlength{\@fpbot}{0pt}
\makeatother

% ===========================
% ===========================
% -------- CONTENTS --------
% ===========================
% ===========================

% ===========================
% ---- TITLE AND AUTHOR ----
% ===========================

\title{RyuTeX：琉和辞典用のLaTeXテンプレート\\\small{（2026年03月22日）}}
\author{セリック・ケナン}
\date{}


% =============================
% ------ FRONT MATTER ---------
% =============================

\begin{document}
\maketitle      
\frontmatter

\pagestyle{plain}

% Table of contents
\begingroup
  \renewcommand{\baselinestretch}{0.95}\selectfont
  \tableofcontents\thispagestyle{empty}
\endgroup

% =============================
% ------ MAIN MATTER ----------
% =============================

\pagestyle{empty} % Suppress header and footer
\mainmatter
\addcontentsline{toc}{part}{解説}\pagestyle{empty}\thispagestyle{empty}
\part*{解説}

% Revert to plain style
\pagestyle{plain}

% Include the explanation and the legend
% !TEX root = ../main.tex
\chapter*{概説}\label{gaisetu}\addcontentsline{toc}{chapter}{概説}\thispagestyle{plain}\setcounter{chapter}{1}

\section{水納島方言}
南琉球宮古語\ruby[g]{水納}{みんな}\ruby{島}{じま}方言（自称〈みんなふつ\heiban{}〉）は沖縄県宮古郡多良間村水納島で伝統的に話されていたことばである。1962年より集団移住の政策が取られたため、島民のほとんどが宮古島の高野集落に移った。現在も水納島方言の多くの話し手たちが高野や宮古島の他の集落に住んでいる。しかし、島民たちが集落を離れて60年間経っている今は水納島の伝統的なことばが全く継承されておらず、その上に伝統的な方言の最後の継承者が少なくなっている。

%水納島は多良間島の北に浮かぶ小島で、宮古島と石垣島の間に位置している。宮古島から約67キロメートル、石垣島から約35キロメートル離れている。水納島


水納島方言は系統的に宮古語に属しており、その中で特に多良間島方言（\ruby[g]{仲筋}{なかすじ}方言、\ruby[g]{塩川}{しおかわ}方言）に近いことが分かっている（図1）\citep{lawrence03:taramakeito,pellard09:phdthesis,celik20:phdthesis,celik20:minna}。

\begin{figure}[ht!]\centering
\begin{tikzpicture}
\tikzset{level distance=0.8cm}
\tikzset{sibling distance=3cm}
%\coordinate
\node{日琉諸語}
  child { node {琉球諸語} 
    child { node{南琉球} 
      child { node{八重山語} }
      child { node{宮古語}
       child { node{多良間諸方言} 
        child { node{\underline{水納島方言}} }
        child { node{多良間島方言} }
      }
       child { node{共通宮古語} }
      } 
    }
    child { node{北琉球} } 
  }
  child { node{日本語}
  }
;
\end{tikzpicture}
\caption{水納島方言の系統}\label{fg:keito}
\end{figure}


多良間諸方言を構成する水納島方言と多良間島方言は音韻論が大きく異なるものの、それ以外の点においてよく似ており、相互理解が完全に成り立っている。明治と大正生まれの水納島方言の報告\citep{zenshuraku:minna}を見ると、明治期と大正期との間に大きな音変化があったと推測される。つまり、〈イ゜〉の母音が〈さ〉〈ざ〉〈つぁ〉の行を除いて〈い〉の母音と完全に合流している。明治生まれの話者が〈ぽーキ゜〉「箒」と発音していたところ、現在は〈ぽーき〉と発音する。一方、興味深いことに水納島方言のアクセント体系の方が多良間島方言に比べ保守的な側面が多い。


% 話者情報

\section{発音}

水納島方言の子音体系を表\ref{tb:shiin}に示す。注意点として、〈さ〉行、〈ざ〉行と〈つぁ〉行の発音の揺れがある。特に〈さ〉行と〈ざ〉行は非口蓋化の発音（さ、すぅ、せ、そ／ざ、ずぅ、ぜ、ぞ）と口蓋化した発音（しゃ、しゅ、しぇ、しょ／じゃ、じゅ、じぇ、じょ）のとで激しく揺れている。〈つぁ〉行も揺れる場合があるが、〈ちゃー〉「茶」のように揺れない語もある。

\begin{table}[ht]\caption{子音体系}\label{tb:shiin}
%\resizebox{\textwidth}{!}{    
\centering
\def\arraystretch{0.95}
     \begin{tabular}{llllll}
     \toprule
        & 両唇・唇歯音 & 歯茎音 & 軟口蓋音 & 声門音 \\
        \midrule
        破裂音 & ぱ \ip{[pa]}  & た \ip{[ta]}  & か \ip{[ka]} &   \\
        & ば \ip{[ba]} & だ\ip{[da]} & が \ip{[ga]} \\[2pt]
        鼻音 & ま \ip{[ma]} & な \ip{[na]} & &  \\[2pt]
        流音 & & ら \ip{[ra\ntilde{}ɾa]} & & \\[2pt]
        摩擦音 & ふぁ \ip{[fa]} & さ \ip{[sa\ntilde{}ɕa]}& & （は \ip{[ha]}） \\
        & ヴぁ \ip{[va\ntilde{}ʋa]} & & \\[2pt]
        破擦音 & & つぁ \ip{[tsa\ntilde{}tɕa]} & & \\
        & & ざ \ip{[dza\ntilde{}dʑa]} & & \\[2pt]
        接近音 & & や \ip{[ja]} & & \\
        %\midrule
        %二重子音 & & & & & & \\
     \bottomrule
 \end{tabular}
% }
\end{table}

もう1つの注意点として、〈ん〉と〈ム〉の区別が曖昧になっていることが挙げられる。つまり、〈ム〉がよく〈ん〉と発音される。例えば、過去接辞の〈たム〉が〈たん〉ともよく発音されるわけである。

水納島の母音体系を表\ref{tb:boin}に示す。注意点として〈イ\m{}〉の母音がある。この母音は〈さ〉行、〈ざ〉行、〈つぁ〉行の子音としか共起しない。つまり、〈す〉〈ず〉〈つ〉以外の音節がない。また、〈い〉の発音と揺れている。〈うす〉「牛」は「うし」、〈みず〉「水」は〈みじ〉、〈みつ〉「道」は〈みち〉とも発音できるわけである。ただし、元の〈い〉の母音は基本的に揺れず、〈ちび〉「尻」は〈ちび〉とだけ発音される。

\begin{table}[ht]\caption{母音体系}\label{tb:boin}
%\resizebox{\textwidth}{!}{    
\centering
\def\arraystretch{0.95}
     \begin{tabular}{lllllll}
     \toprule
        狭 & い \ip{[i]} & いー \ip{[iː]} & イ\m{} \ip{[ʉ]} & イ\m{}ー \ip{[ʉː]}  & う \ip{[u]} & うー \ip{[uː]}  \\[2pt]
        半狭 & え \ip{[e]} & えー \ip{[eː]} & & & & おー \ip{[oː]} \\[2pt]
        広 & & & & & あ \ip{[a]} & あー \ip{[aː]} \\[2pt]
        \bottomrule
 \end{tabular}
% }
\end{table}


\FloatBarrier

\section{アクセント}

水納島方言は主に3つのアクセント型（a型、b型、c型）を区別している\citep{celik20:minna}。これらのアクセント型はピッチ（声の高さ）の局所的な変動（下降または上昇）の有無と位置によって対立している。a型はピッチの変動を伴わないのに対して、b型は遅めのピッチの変動、c型は早めのピッチの変動を伴う。それに加えて、数語にしか観察されない特殊型も存在する。この特殊型は2つのピッチの変動（下降と上昇）を伴う。それぞれのアクセント型の例を\Next[]に示す（ピッチ変動のないことを「\heiban{}」、ピッチの下降を「\kako{}」、ピッチの上昇を「\josho{}」の記号で表す）。

\ex. 
    \makebox[4em][l]{a型：}\makebox[11em][l]{ぶとぅまい\heiban{}　ぶらーん}「夫もいない」\\
    \makebox[4em][l]{b型：}\makebox[11em][l]{みどぅムま\kako{}い　ぶらーん}「女性もいない」\\
    \makebox[4em][l]{c型：}\makebox[11em][l]{う\kako{}やまい　ぶらーん}「お父さんもいない」\\
    \makebox[4em][l]{特殊型：}\makebox[11em][l]{ヴぇー\kako{}い\josho{}まい　あらん}「これでもない」


しかし、多良間方言と同様に、各アクセント型の実現が環境によって異なる。上の\Last[]の例は発話のはじめという環境に該当しているが、発話の頭にc型の〈ん\kako{}め〉「もう」を入れてみると、\Next[]のように各アクセント型の実現が変わってくる（特殊型は省略する）。つまり、a型が低く平らに発音されるのに対して、b型とc型はそれぞれ遅めと早めのピッチの上昇で実現している。

\ex. 
    a型：\makebox[15em][l]{ん\kako{}め　ぶとぅまい\heiban{}　ぶらーん}「もう夫もいない」\\
    b型：\makebox[15em][l]{ん\kako{}め　みどぅムまい　\josho{}ぶらーん}「もう女性もいない」\\
    c型：\makebox[15em][l]{ん\kako{}め　うや\josho{}まい　ぶらーん}「もうお父さんもいない」


\LLast[]におけるアクセント型の実現を「下降調」、\Last[]におけるものを「上昇調」と呼ぶ。各アクセント型の実現がこのように環境によって逆転するという特徴は日琉諸語の中でもたいへん珍しく、他に報告がないようである。

\section{動詞の活用クラス}

動詞は2種類の規則的な活用類、すなわち\rom{1}類と\rom{2}類に分かれる。\rom{1}類に属する動詞は「書く」「読む」のように日本語の五段動詞に対応しており、命令形が〈い〉に終わり、否定形において〈あ〉が現れる。\rom{2}類に属する動詞は「見る」「下りる」のように一段動詞に対応しており、命令形が〈る〉に終わり、否定形において〈あ〉が現れない。それに加えて、混合型の動詞〈すに〉「死ぬ」と不規則な活用を示す変則型の動詞〈しー〉「する」と〈きー〉「来る」の動詞がある。表\ref{tb:katuyorui}にそれぞれの活用類の代表例を示す。

\begin{table}[ht]\caption{活用類の代表例}\label{tb:katuyorui}
\resizebox{\textwidth}{!}{    
\centering
\def\arraystretch{0.9}
     \begin{tabular}{lllll}
     \toprule
        活用類 & 動詞 & 基本形  & 命令形 & 否定形　\\
        \midrule
        \rom{1} &「書く」& かき  & かき & かかん \\
         & 「漕ぐ」&くぎ& くぎ & くがん \\
        & 「飛ぶ」 &とぅび& とぅび& とぅばん\\
        & 「読む」 &ゆム& ゆみ & ゆまん\\
        & 「待つ」 &まつ／まてぃ& まてぃ & またん\\
        & 「成す」 &なす & なし & なしゃん\\
        & 「成る」 &ない & なり & ならん \\
        & 「実る」 &みーり & みーり & みーらん \\
        & 「降る」 &っふぃ & っふぃ& っふぁん\\
        & 「売る」&っヴぃ & っヴぃ& っヴぁん\\
        & 「移る」 &うっつ & うっち & うっちゃん\\
        & 「寄る」& ゆっず &ゆっじ & ゆっじゃん\\
        & 「切る」 &きー & きし & きしゃん\\
        & 「言う」 &いー & いじ & いじゃん \\
        & 「買う」 &こー & けー & かーん \\
        & 「結う」 &ゆー & いぇー & やーん \\
        %\hhline{~----}
        \rom{1}（変則） &「ある」& あい &－& ねーん \\
        & 「\ntilde{}だ」 &あい &－& あらん \\
        & 「居る」 &ぶい & ぶり & ぶら（ー）ん\\
        \hline
        \rom{2} &「見る」& みー & みーる & みーん\\
         & 「貰う」& いぇー & いぇーる & いぇーん \\
         & 「下る」& うり & うりる & うりん\\
        %\hhline{~---}
        \rom{2}（変則）&「寝る」&にに &ににる／にんる&ににん\\
        \midrule
        混合型 & 「死ぬ」&すに & すにる／すんる & すなん \\
        \midrule
        変則型 & 「する」&しー & しる & しゅ（ー）ん \\
         & 「来る」 &きー& くー  & く（ー）ん \\
     \bottomrule
 \end{tabular}
}
\end{table}


% !TEX root = ../main.tex
%\clearpage
\nocite{shimojik17:taramajiten,tomihama13:irabujiten,95:okinawakogodaijiten}

\nocite{jarosz15:phdthesis,takahashi93:taramagoi,takahashi94:taramagoi,takahashi95:taramagoi,hirayama83:miyakokisogoi,pellard09:phdthesis,pellard15:linguisticarcheology}

\nocite{kokuritu63:okinawagojiten,yonaguni19:yonagunijiten,nevskiy05:notes,ryudai68:miyakochosa,taramason73:sonshi,tokuyamacelik20:taramajiten}

\nocite{taramasonshi00,taramasonshi86,taramasonshi05,taramasonshi93,taramasonshi89,momohara19:taramadokuhon}

\nocite{maeara11:taketomijiten,miyagi03:ishigakijiten,95:okinawakogodaijiten,kokuritu63:okinawagojiten,shimojik11:minnashiryo,shimojik12:minnameishi,shimojik17:minnaadjectives,kiku05:yoronjiten}

\begingroup
\pagestyle{plain}
\printbibliography[
	title={参考文献},
	heading=bibintoc
]
\endgroup
% !TEX root = ../main.tex
\chapter*{凡例}\label{hanrei}\addcontentsline{toc}{chapter}{凡例}\setcounter{section}{0}\thispagestyle{plain}

\section{収録語}

本語彙集は第二著者、南琉球宮古語水納島方言の母語話者である大浦辰夫（昭和14年生）のことばを収録したものである。全部で5768項目（発音の揺れも含む）を収録した。

\section{見出し語}

見出し語はアクセント付き仮名表記、\ip{IPA}（国際音声記号）表記、品詞、意味範疇、活用、意味記述、用例、備考、発音の揺れ、類似語から構成されている。

\section{配列}

項目は仮名表記を基に五十音順に並べた（配列の際は分かち書きの空白を無視している）。従って〈ム〉から始まる語は〈む〉から始まる語と同じところにある。仮名表記が同じである場合、品詞を基準とし、接頭辞、接尾辞、助詞、終助詞、助数詞、数詞語根、連体詞、語素、名詞、動詞、形容詞、ナ形容詞、副詞、重複形、擬態語・擬音語、感動詞、接続詞、連語、諺の順で項目を並べた。

\section{仮名表記}

仮名表記は『伊良部方言辞典』の表記を踏襲した『南琉球宮古語多良間方言辞典』の表記に従った。ただし、第二著者の語感に従って有声唇歯摩擦音の二重子音\ip{[vv]}を「っヴ」に改めた。

概説で説明した通り、〈ム〉の発音が必ずしも安定していない。また、子音が続く場合、短音の〈ム〉と〈ん〉が後続の子音の調音点に同化し、完全に中和すると思われる（逆に言うとほとんどの環境において〈ム〉と〈ん〉が対立しない）。従って、本語彙集では環境で予測できない〈ム〉だけを〈ム〉と表記した。例えば、「女性」は〈みどぅム〉と発音されたため、見出し語として〈みどぅム〉を立てたが、「女子生徒」は子音が続くため、〈みどぅんしーとぅ〉のように表記した。また、項目の表記を決める際に、実際の発音によったので、同じ語根を含む複合語などの間には表記の統一を行っていない。例えば、単独の「女性」は〈みどぅム〉だが、「美人」は〈あぱらぎみどぅん〉となっている。

また、水納島方言では〈わ〉と〈ヴぁ〉の区別がないと思われるが、十分な観察ができなかったこともあり、単子音の場合は慣習的に〈わ〉を使った。

\section{アクセント}

アクセント型の認定（主に名詞、動詞、形容詞、擬態語・擬音語）は第一著者が行った。アクセント型が分かるように仮名表記にアクセント記号を付与した。平板型のa型は「\heiban{}」の記号で示した。ピッチの変動のある型は「\kako{}」と「\josho{}」の記号を用いてピッチの下降および上昇を示した。提示した下降と上昇の位置は発話のはじめに2拍の接語が続いた発音による。ただし、b型はピッチの変動が後続の接語で実現するため、語の後に「\kako{}」を付与して示した。アクセント型がピッチの変動がある発音とピッチの変動がない発音とで揺れる場合はピッチの記号を（　）で囲んだ。

各項目のアクセント型の認定は必ず複数のアクセント資料の観察に基づいて行った。名詞と擬態語・擬音語は下降調と上昇調の両方の発音を収録し、認定した。動詞は下降調における進行形と過去否定形の発音を収録し、認定した。形容詞は〈～さどぅ　あい〉、〈～さまい　ねーん〉と〈～むぬ〉の発音を収録し、認定した。十分な確信が得られなかった場合は再調査を実施してアクセント型を認定したが、アクセント型の判断が最後まで下しきれかった項目はアクセントを省略した（未調査の項目も若干数ある）。

\section{IPA表記}

音韻的な解釈に基づき、\ip{IPA}（国際音声記号）の適切な記号を選んだ。母音の無声化は音声的な特徴であると考え、表記していない。また、〈しみんやー〉\ip{[ɕiminjaː]}「セメントの家」のように、\ip{IPA}表記において音節境界が曖昧になる場合、音節境界を「\ip{.}」で示した（\ip{[ɕimin.jaː]}）。

\section{品詞}

見出し語に対し次のような品詞を設定し、括弧の中にある略号を用いて示した：助詞（助）、終助詞（終）、助数詞（助数）、数詞語根（数）、連体詞（連体）、名詞（名）、動詞（動）、形容詞（形）、ナ形容詞（ナ形）、副詞（副）、重複形（重複）、擬態語・擬音語（擬）、感動詞（感）、接続詞（接続）。それに加えて、接頭辞（接頭）、接尾辞（接尾）、語素（語素）、連語（連語）、諺（諺）の範疇も設けた。

動詞は品詞の後に活用クラス（「\rom{1}」、「\rom{2}」、「変則型」、「混合型」）も示した。活用クラスの認定は多良間方言との対応語の活用クラスを考慮しながら、過去否定形に基づいて行った。動詞によって否定形が揺れるため、活用クラスの認定を控えた項目もある。

\section{意味範疇}

一部の語に対して次の意味範疇を設けて示した：干支、貝、蟹、疑問詞（疑）、魚、指示詞（指）、植物（植）、人名、代名詞（代）、地名（地）、鳥、虫、屋号。

\section{活用}

動詞は接続形（〜して）と否定形（〜しない）の活用形を示した。活用形は見出し語に合うように、ある程度規範的に語形を整えた。

\section{意味記述}

意味は相当する共通語を示した上、必要な場合に解説を加えた。多義語は「\numbering{1}\numbering{2}\numbering{3}...」など番号を付けて意味を分けて示した。

\section{用例}

収録語の意味の理解を助けるために657点の用例を加えた。可能な限りピッチの上り・下がりを記そうとしたが、確信が得られなかった用例はピッチ記号を省略した。用例は発音通りに掲載しているため、見出し語の語形と用例における語形が一致しないことがある。

\section{備考}

必要に応じて、見出し語に関する備考を加えた。

\section{発音の揺れ}

日本語の「おととい」・「おとつい」のように同じ語に対して複数の発音が許容される場合、それぞれの発音を別の項目に立て、それぞれの項目の中に異なる発音を［同］で示し、掲載した。しかし、次の2つの場合においては異なる方針を取った。第一に、水納島方言では〈す〉〈つ〉〈ず〉が〈し〉〈ち〉〈じ〉の自由変異を持つため、原則より古い発音である前者の発音のみを立項した。例えば〈みず〉\ws{}～\ws{}〈みじ〉「水」に対して〈みず〉のみを立てた。第二に、〈さ〉行、〈ざ〉行、部分的に〈つぁ〉行が非口蓋化の発音（さ、せ、すぅ、そ）と口蓋化した発音（しゃ、しぇ、しゅ、しょ）で激しく揺れているが、一々発音の揺れを立項せず、調査の際に得られた発音をそのまま見出し語の表記として使用した。従って、本語彙集を引くに当たって、両方の文字（さ・しゃなど）で単語を探す必要がある（または共通語引きを使う）。%

\section{類似語}

見出し語と同様の意味を持つ語は［類］で示し、掲載した。

\section{日本語引き}

便宜のため、本語彙集の後に日本語引きを用意した。

\begin{table}[p]\caption*{仮名・発音記号一覧}\label{tb:kanaichiran}
\resizebox{\textwidth}{!}{    
\def\arraystretch{1}
     \begin{tabular}{|cc|cc|cc|cc|cc|cc|}
     %\hline
     \Xhline{2\arrayrulewidth}
     %\toprule
あ&\small{\ip{[a]}}&い&\small{\ip{[i]}}&& &う&\small{\ip{[u]}}&え&\small{\ip{[e]}}&お&\small{\ip{[o]}}\\
か&\small{\ip{[ka]}}&き&\small{\ip{[ki]}}&&&く&\small{\ip{[ku]}}&け&\small{\ip{[ke]}}&こ&\small{\ip{[ko]}}\\
が&\small{\ip{[ga]}}&ぎ&\small{\ip{[gi]}}&&&ぐ&\small{\ip{[gu]}}&げ&\small{\ip{[ge]}}&ご&\small{\ip{[go]}}\\
きゃ&\small{\ip{[kja]}}&&&&&きゅ&\small{\ip{[kju]}}&&&きょ&\small{\ip{[kjo]}}\\
ぎゃ&\small{\ip{[gja]}}&&&&&ぎゅ&\small{\ip{[gju]}}&&&ぎょ&\small{\ip{[gjo]}}\\
さ&\small{\ip{[sa]}}&&&す&\small{\ip{[sʉ]}}&すぅ&\small{\ip{[su]}}&せ&\small{\ip{[se]}}&そ&\small{\ip{[so]}}\\
ざ&\small{\ip{[dza]}}&&&ず&\small{\ip{[dzʉ]}}&ずぅ&\small{\ip{[dzu]}}&ぜ&\small{\ip{[dze]}}&ぞ&\small{\ip{[dzo]}}\\
しゃ&\small{\ip{[ɕa]}}&し&\small{\ip{[ɕi]}}&&&しゅ&\small{\ip{[ɕu]}}&しぇ&\small{\ip{[ɕe]}}&しょ&\small{\ip{[ɕo]}}\\
じゃ&\small{\ip{[dʑa]}}&じ&\small{\ip{[dʑi]}}&&&じゅ&\small{\ip{[dʑu]}}&じぇ&\small{\ip{[dʑe]}}&じょ&\small{\ip{[dʑo]}}\\
た&\small{\ip{[ta]}}&てぃ&\small{\ip{[ti]}}&&&とぅ&\small{\ip{[tu]}}&て&\small{\ip{[te]}}&と&\small{\ip{[to]}}\\
だ&\small{\ip{[da]}}&でぃ&\small{\ip{[di]}}&&&どぅ&\small{\ip{[du]}}&で&\small{\ip{[de]}}&ど&\small{\ip{[do]}}\\
ちゃ&\small{\ip{[tɕa]}}&ち&\small{\ip{[tɕi]}}&つ&\small{\ip{[tsʉ]}}&ちゅ&\small{\ip{[tɕu]}}&ちぇ&\small{\ip{[tɕe]}}&ちょ&\small{\ip{[tɕo]}}\\
つぁ&\small{\ip{[tsa]}}&&&&&つぅ&\small{\ip{[tsu]}}&つぇ&\small{\ip{[tse]}}&つぉ&\small{\ip{[tso]}}\\
な&\small{\ip{[na]}}&に&\small{\ip{[ni]}}&&&ぬ&\small{\ip{[nu]}}&ね&\small{\ip{[ne]}}&の&\small{\ip{[no]}}\\
にゃ&\small{\ip{[nja]}}&&&&&にゅ&\small{\ip{[nju]}}&&&にょ&\small{\ip{[njo]}}\\
は&\small{\ip{[ha]}}&ひ&\small{\ip{[hi]}}&&&ほぅ&\small{\ip{[hu]}}&へ&\small{\ip{[he]}}&ほ&\small{\ip{[ho]}}\\
ば&\small{\ip{[ba]}}&び&\small{\ip{[bi]}}&&&ぶ&\small{\ip{[bu]}}&べ&\small{\ip{[be]}}&ぼ&\small{\ip{[bo]}}\\
ぱ&\small{\ip{[pa]}}&ぴ&\small{\ip{[pi]}}&&&ぷ&\small{\ip{[pu]}}&ぺ&\small{\ip{[pe]}}&ぽ&\small{\ip{[po]}}\\
ひゃ&\small{\ip{[hja]}}&&&&&ひゅ&\small{\ip{[hju]}}&&&ひょ&\small{\ip{[hjo]}}\\
びゃ&\small{\ip{[bja]}}&&&&&びゅ&\small{\ip{[bju]}}&&&びょ&\small{\ip{[bjo]}}\\
ぴゃ&\small{\ip{[pja]}}&&&&&ぴゅ&\small{\ip{[pju]}}&&&ぴょ&\small{\ip{[pjo]}}\\
ふぁ&\small{\ip{[fa]}}&ふぃ&\small{\ip{[fi]}}&&&ふ&\small{\ip{[fu]}}&ふぇ&\small{\ip{[fe]}}&ふぉ&\small{\ip{[fo]}}\\
ヴぁ&\small{\ip{[va]}}&ヴぃ&\small{\ip{[vi]}}&&&ヴぅ&\small{\ip{[vu]}}&ヴぇ&\small{\ip{[ve]}}&ヴぉ&\small{\ip{[vo]}}\\
ま&\small{\ip{[ma]}}&み&\small{\ip{[mi]}}&&&む&\small{\ip{[mu]}}&め&\small{\ip{[me]}}&も&\small{\ip{[mo]}}\\
みゃ&\small{\ip{[mja]}}&&&&&みゅ&\small{\ip{[mju]}}&&&みょ&\small{\ip{[mjo]}}\\
や&\small{\ip{[ja]}}&&&&&ゆ&\small{\ip{[ju]}}&いぇ&\small{\ip{[je]}}&よ&\small{\ip{[jo]}}\\
ら&\small{\ip{[ra]}}&り&\small{\ip{[ri]}}&&&る&\small{\ip{[ru]}}&れ&\small{\ip{[re]}}&ろ&\small{\ip{[ro]}}\\
っら&\small{\ip{[ɭɭa]}}&&&&&&&&&っろ&\small{\ip{[ɭɭo]}}\\
りゃ&\small{\ip{[rja]}}&&&&&りゅ&\small{\ip{[rju]}}&&&りょ&\small{\ip{[rjo]}}\\
わ&\small{\ip{[ʋa]}}&&&&&&&&&を&\small{\ip{[ʋo]}}\\
ん&\small{\ip{[n]}}&ム&\small{\ip{[m]}}&ヴ&\small{\ip{[v]}}&&&&&&\\
ー&\small{\ip{[ː]}}&\multicolumn{3}{c}{っ\hspace{2.5pt}（子音を重ねる）}&&&&&&&\\
     %\hline
     %
     \Xhline{2\arrayrulewidth}
     %\bottomrule
 \end{tabular}
 }
\end{table}



%\clearpage

% =====================
% ----- DICTIONARY ----
% =====================

% Add title of dictionary part (if necessary)
\addcontentsline{toc}{part}{本文}\pagestyle{empty}\thispagestyle{empty}
\part*{本文}
 
\begingroup
  % Include the dictionary's contents
  \cleardoublepage % or \clearpage in oneside 
  \pagestyle{dictionary}
  \phantomsection
  \addcontentsline{toc}{chapter}{琉球・日本語}
  \SetDictBlockWidth

  % Suppress paragraph indent for the entries
  \setlength{\parindent}{0pt}
  
  % Set the size of paragraph skip between entries
  \setlength{\parskip}{6pt}

  % Adjustment of between-line space
  %\renewcommand{\baselineskip}{0.9}

  % Switch to multi-column mode
  % starred multicols for not balancing the columns 
  \begin{multicols*}{2}
    \LetterHead{あ}
\entry
	{\headPrint[kana={あー},ipa={aː},pos={感}]}%
	{%
		\MeaningsBlock{%
	\ryumeaning%
		{\ruby[g]{嗚呼}{ああ}}{}{}%
	}%
	}%
	{\tailPrint[]}%
\entry
	{\headPrint[kana={あーい},ipa={aːi},pos={感}]}%
	{%
		\MeaningsBlock{%
	\ryumeaning%
		{いいえ。\ruby[g]{嫌}{いや}}{}{}%
	}%
	}%
	{\tailPrint[]}%
\entry
	{\headPrint[kana={あーいし\tl{˥}},ipa={aː.iʃi},pos={名},acc={H}]}%
	{%
		\MeaningsBlock{%
	\ryumeaning%
		{砂岩}{「粟石」の義}{}%
	}%
	}%
	{\tailPrint[]}%
\entry
	{\headPrint[kana={あーさ\tl{˥}},ipa={aːsa},pos={名},acc={H}]}%
	{%
		\MeaningsBlock{%
	\ryumeaning%
		{ヒトエグサ}{海藻名。食用}{}%
	}%
	}%
	{\tailPrint[]}%
\entry
	{\headPrint[kana={あーしみるん\tl{˥}},ipa={aːʃi̥miruɴ},pos={動},rui={3},para={あーすぅむぬ},acc={H}]}%
	{%
		\MeaningsBlock{%
	\ryumeaning%
		{会わせる}{会うようにさせる}{}%
	}%
	}%
	{\tailPrint[]}%
\entry
	{\headPrint[kana={あーすん\tl{˥}},ipa={aːsuɴ},pos={動},rui={2},para={あーはぬ},acc={H}]}%
	{%
		\MeaningsBlock{%
	\ryumeaning%
		{合わせる}{}{}%
	\ryumeaning%
		{戦わせる}{}{}%
	}%
	}%
	{\tailPrint[]}%
\entry
	{\headPrint[kana={あーらすん},ipa={aːrasuɴ},pos={動}]}%
	{%
		\MeaningsBlock{%
	\ryumeaning%
		{研ぐ}{臼の歯を削る。のこぎりの歯を研ぐ}{}%
	}%
	}%
	{\tailPrint[]}%
\entry
	{\headPrint[kana={あーらすん\tl{˥˩}},ipa={aːrasuɴ},pos={動},rui={2},para={あーらはぬ},acc={F}]}%
	{%
		\MeaningsBlock{%
	\ryumeaning%
		{騒がす。騒がしくする}{}{}%
	\ryumeaning%
		{急き立てる。急がせる}{}{}%
	}%
	}%
	{\tailPrint[]}%
\entry
	{\headPrint[kana={あーらすん\tl{˥˩}},ipa={aːrasuɴ},pos={動},rui={2},para={あーらはぬ},acc={F}]}%
	{%
		\MeaningsBlock{%
	\ryumeaning%
		{蒸す}{餅や菓子類などを蒸す}{}%
	}%
	}%
	{\tailPrint[]}%
\entry
	{\headPrint[kana={あーり\tl{˥˩}},ipa={aːri},pos={名},acc={F}]}%
	{%
		\MeaningsBlock{%
	\ryumeaning%
		{蟻}{}{}%
	}%
	}%
	{\tailPrint[]}%
\entry
	{\headPrint[kana={あーり\tl{˥˩}},ipa={aːri},pos={名},acc={F}]}%
	{%
		\MeaningsBlock{%
	\ryumeaning%
		{東。東方}{}{}%
	}%
	}%
	{\tailPrint[]}%
\entry
	{\headPrint[kana={あーるん\tl{˥˩}},ipa={aːruɴ},pos={動},rui={1},para={あーらぬ},acc={F}]}%
	{%
		\MeaningsBlock{%
	\ryumeaning%
		{騒ぐ}{}{}%
	}%
	}%
	{\tailPrint[]}%
\entry
	{\headPrint[kana={あい},ipa={ai},pos={感}]}%
	{%
		\MeaningsBlock{%
	\ryumeaning%
		{おや。はて}{}{}%
	}%
	}%
	{\tailPrint[]}%
\entry
	{\headPrint[kana={あい\tl{˥˩}},ipa={ai},pos={名},acc={F}]}%
	{%
		\MeaningsBlock{%
	\ryumeaning%
		{藍。染料}{}{}%
	}%
	}%
	{\tailPrint[]}%
\entry
	{\headPrint[kana={あいさち\tl{˥˩}},ipa={aisatʃi},pos={名},acc={F}]}%
	{%
		\MeaningsBlock{%
	\ryumeaning%
		{挨拶}{}{}%
	}%
	}%
	{\tailPrint[]}%
\entry
	{\headPrint[kana={あいじ\tl{˥}},ipa={aidʒi},pos={名},acc={H}]}%
	{%
		\MeaningsBlock{%
	\ryumeaning%
		{蜻蛉}{}{}%
	}%
	}%
	{\tailPrint[]}%
\entry
	{\headPrint[kana={あいず\tl{˥}},ipa={aidʒi},pos={名},acc={H}]}%
	{%
		\MeaningsBlock{%
	\ryumeaning%
		{合図}{}{}%
	}%
	}%
	{\tailPrint[]}%
\entry
	{\headPrint[kana={あいずみ},ipa={aidzumi},pos={名}]}%
	{%
		\MeaningsBlock{%
	\ryumeaning%
		{藍染}{}{}%
	}%
	}%
	{\tailPrint[]}%
\entry
	{\headPrint[kana={あいち\tl{˥}},ipa={aitʃi},pos={名},acc={H}]}%
	{%
		\MeaningsBlock{%
	\ryumeaning%
		{\ruby[g]{木}{き}\ruby[g]{槌}{づち}}{}{}%
	}%
	}%
	{\tailPrint[]}%
\entry
	{\headPrint[kana={あいてぃ\tl{˥}},ipa={aiti},pos={名},acc={H}]}%
	{%
		\MeaningsBlock{%
	\ryumeaning%
		{相手}{}{}%
	}%
	}%
	{\tailPrint[]}%
\entry
	{\headPrint[kana={あいなー\tl{˥}},ipa={ainaː},pos={名},acc={H}]}%
	{%
		\MeaningsBlock{%
	\ryumeaning%
		{花嫁}{}{}%
	}%
	}%
	{\tailPrint[]}%
\entry
	{\headPrint[kana={あいなー\tl{˥}\ws{}すん\tl{˥˩}},ipa={ainaː suɴ},pos={句},acc={H F}]}%
	{%
		\MeaningsBlock{%
	\ryumeaning%
		{\ruby[g]{娶}{めと}る}{嫁をもらう。嫁として縁組する}{}%
	}%
	}%
	{\tailPrint[]}%
\entry
	{\headPrint[kana={\josho{}あいなー\kako{}よい},ipa={ainaːjoi},pos={名},acc={H+L}]}%
	{%
		\MeaningsBlock{%
	\ryumeaning%
		{結婚式。結婚祝い}{}{}%
	}%
	}%
	{\tailPrint[]}%
\entry
	{\headPrint[kana={あいま\tl{˥}},ipa={aima},pos={名},acc={H}]}%
	{%
		\MeaningsBlock{%
	\ryumeaning%
		{あいま。隙間}{}{}%
	}%
	}%
	{\tailPrint[]}%
\entry
	{\headPrint[kana={あいるん\tl{˥}},ipa={airuɴ},pos={動},rui={3},para={あるぬ},acc={H}]}%
	{%
		\MeaningsBlock{%
	\ryumeaning%
		{会える。逢える}{}{}%
	}%
	}%
	{\tailPrint[]}%
\entry
	{\headPrint[kana={あいるん\tl{˥}},ipa={airuɴ},pos={動},acc={H}]}%
	{%
		\MeaningsBlock{%
	\ryumeaning%
		{和える。混ぜる。混ぜ合わせる}{}{}%
	}%
	}%
	{\tailPrint[]}%
\entry
	{\headPrint[kana={あか\tl{˥}},ipa={aka},pos={名},acc={H}]}%
	{%
		\MeaningsBlock{%
	\ryumeaning%
		{\ruby[g]{淦}{あか}}{船底に溜まった海水。船舶用語}{}%
	}%
	}%
	{\tailPrint[]}%
\entry
	{\headPrint[kana={あがー},ipa={agaː},pos={感}]}%
	{%
		\MeaningsBlock{%
	\ryumeaning%
		{いたっ。痛い}{}{}%
	}%
	}%
	{\tailPrint[]}%
\entry
	{\headPrint[kana={あかーいる},ipa={akaːiru},pos={名}]}%
	{%
		\MeaningsBlock{%
	\ryumeaning%
		{赤色}{}{}%
	}%
	}%
	{\tailPrint[]}%
\entry
	{\headPrint[kana={あ\josho{}かー\kako{}し},ipa={akaːʃi},pos={副},acc={LHL}]}%
	{%
		\MeaningsBlock{%
	\ryumeaning%
		{赤く}{}{}%
	}%
	}%
	{\tailPrint[]}%
\entry
	{\headPrint[kana={あ\josho{}かー\kako{}し\ws{}なるん\tl{˩˥}},ipa={akaːʃi naruɴ},pos={句},acc={HL R}]}%
	{%
		\MeaningsBlock{%
	\ryumeaning%
		{赤くなる。赤らむ}{}{}%
	}%
	}%
	{\tailPrint[]}%
\entry
	{\headPrint[kana={あかー\ws{}すん},ipa={akaː suɴ},pos={句}]}%
	{%
		\MeaningsBlock{%
	\ryumeaning%
		{赤くなる。赤らむ}{人についてだけ言う}{}%
	}%
	}%
	{\tailPrint[]}%
\entry
	{\headPrint[kana={あがしぃ\tl{˥˩}},ipa={agaʃi},pos={名},acc={F}]}%
	{%
		\MeaningsBlock{%
	\ryumeaning%
		{松の根の芯}{昔は燈明に使われた}{}%
	}%
	}%
	{\tailPrint[]}%
\entry
	{\headPrint[kana={あかしきぶし},ipa={akaʃi̥kibuʃi},pos={名}]}%
	{%
		\MeaningsBlock{%
	\ryumeaning%
		{明けの明星}{}{}%
	}%
	}%
	{\tailPrint[]}%
\entry
	{\headPrint[kana={あかしきん\tl{˥}},ipa={akaʃi̥kin},pos={名},acc={H}]}%
	{%
		\MeaningsBlock{%
	\ryumeaning%
		{\ruby[g]{暁}{あかつき}}{}{}%
	}%
	}%
	{\tailPrint[]}%
\entry
	{\headPrint[kana={あがすん\tl{˥˩}},ipa={agasuɴ},pos={動},rui={2},para={あがはぬ},acc={F}]}%
	{%
		\MeaningsBlock{%
	\ryumeaning%
		{仲裁する}{}{}%
	}%
	}%
	{\tailPrint[]}%
\entry
	{\headPrint[kana={あがたま\tl{˥}},ipa={agatama},pos={名},acc={H}]}%
	{%
		\MeaningsBlock{%
	\ryumeaning%
		{赤ん坊。赤子}{}{}%
	}%
	}%
	{\tailPrint[note={前部要素に下降型の\emb{あが}「赤い」を含んでいると考えられるが、アクセントが平進型となっている。これは\emb{たま}「指小辞」の影響によると思われる}]}%
\entry
	{\headPrint[kana={あがだん\tl{˥˩}},ipa={agadaɴ},pos={名},acc={F}]}%
	{%
		\MeaningsBlock{%
	\ryumeaning%
		{赤ダニ}{牛などに付くダニ}{}%
	}%
	}%
	{\tailPrint[]}%
\entry
	{\headPrint[kana={あがでーぐに\tl{˥˩}},ipa={agadeːguni},pos={名},acc={F}]}%
	{%
		\MeaningsBlock{%
	\ryumeaning%
		{ニンジン}{「赤い大根」の義}{}%
	}%
	}%
	{\tailPrint[]}%
\entry
	{\headPrint[kana={あがぱーち\tl{˥˩}},ipa={agapaːtʃi},pos={名},acc={F}]}%
	{%
		\MeaningsBlock{%
	\ryumeaning%
		{スズメバチ}{}{}%
	}%
	}%
	{\tailPrint[]}%
\entry
	{\headPrint[kana={あが\tl{˥˩}はん},ipa={agahaɴ},pos={形},acc={F}]}%
	{%
		\MeaningsBlock{%
	\ryumeaning%
		{赤い}{}{}%
	}%
	}%
	{\tailPrint[]}%
\entry
	{\headPrint[kana={あがまーみ\tl{˥˩}},ipa={agamaːmi},pos={名},acc={F}]}%
	{%
		\MeaningsBlock{%
	\ryumeaning%
		{\ruby[g]{小豆}{あずき}}{}{}%
	}%
	}%
	{\tailPrint[]}%
\entry
	{\headPrint[kana={あがますん\tl{˥˩}},ipa={agamasuɴ},pos={動},acc={F}]}%
	{%
		\MeaningsBlock{%
	\ryumeaning%
		{赤くする}{}{}%
	}%
	}%
	{\tailPrint[]}%
\entry
	{\headPrint[kana={あがみるん\tl{˥}},ipa={agamiruɴ},pos={動},acc={H}]}%
	{%
		\MeaningsBlock{%
	\ryumeaning%
		{崇める。敬う}{}{}%
	}%
	}%
	{\tailPrint[]}%
\entry
	{\headPrint[kana={あがむん\tl{˥˩}},ipa={agamuɴ},pos={動},rui={1},para={あがまぬ},acc={F}]}%
	{%
		\MeaningsBlock{%
	\ryumeaning%
		{赤らむ。赤くなる}{}{}%
	}%
	}%
	{\tailPrint[]}%
\entry
	{\headPrint[kana={あがやー},ipa={agajaː},pos={感}]}%
	{%
		\MeaningsBlock{%
	\ryumeaning%
		{残念。惜しい}{}{}%
	}%
	}%
	{\tailPrint[]}%
\entry
	{\headPrint[kana={あがよー},ipa={agajoː},pos={感}]}%
	{%
		\MeaningsBlock{%
	\ryumeaning%
		{とても痛い。痛み嘆くさま。嘆き悲しむさま}{\emb{あがよーあがよー}はその強調}{}%
	}%
	}%
	{\tailPrint[]}%
\entry
	{\headPrint[kana={あがらすん\tl{˥˩}},ipa={agarasuɴ},pos={動},rui={2},para={あがらはぬ},acc={F}]}%
	{%
		\MeaningsBlock{%
	\ryumeaning%
		{明るくする。照らす}{}{}%
	\ryumeaning%
		{（夜を）明かす}{}{}%
	}%
	}%
	{\tailPrint[]}%
\entry
	{\headPrint[kana={あがり\tl{˥˩}},ipa={agari},pos={名},acc={F}]}%
	{%
		\MeaningsBlock{%
	\ryumeaning%
		{灯り}{}{}%
	}%
	}%
	{\tailPrint[]}%
\entry
	{\headPrint[kana={あがりゃん},ipa={agarjaɴ},pos={動（継）}]}%
	{%
		\MeaningsBlock{%
	\ryumeaning%
		{明るい}{}{}%
	}%
	}%
	{\tailPrint[]}%
\entry
	{\headPrint[kana={あがるん\tl{˥˩}},ipa={agaruɴ},pos={動},rui={1},para={あがらぬ},acc={F}]}%
	{%
		\MeaningsBlock{%
	\ryumeaning%
		{明るくなる}{}{}%
	\ryumeaning%
		{明ける}{}{}%
	}%
	}%
	{\tailPrint[]}%
\entry
	{\headPrint[kana={あがるん\tl{˥˩}},ipa={agaruɴ},pos={動},rui={1},para={あがらぬ},acc={F}]}%
	{%
		\MeaningsBlock{%
	\ryumeaning%
		{上がる。昇る}{\emb{しな\ws{}あがるん}「太陽が昇る」}{}%
	}%
	}%
	{\tailPrint[]}%
\entry
	{\headPrint[kana={あがん\tl{˩˥}},ipa={agaɴ},pos={名},acc={R}]}%
	{%
		\MeaningsBlock{%
	\ryumeaning%
		{さつま芋。甘藷}{}{}%
	}%
	}%
	{\tailPrint[]}%
\entry
	{\headPrint[kana={あが\tl{˥˩}んた\tl{˩˥}},ipa={aganta},pos={名},acc={F+R}]}%
	{%
		\MeaningsBlock{%
	\ryumeaning%
		{赤土。赤い粘土}{}{}%
	}%
	}%
	{\tailPrint[]}%
\entry
	{\headPrint[kana={あがんたま\tl{˥}},ipa={agantama},pos={名},acc={H}]}%
	{%
		\MeaningsBlock{%
	\ryumeaning%
		{赤ん坊。赤子}{}{}%
	}%
	}%
	{\tailPrint[note={前部要素に下降型の\emb{あが}「赤い」を含んでいると考えられるが、アクセントが平進型となっている。これは\emb{たま}「指小辞」の影響によると思われる}]}%
\entry
	{\headPrint[kana={あぎごっこー\tl{˥˩}},ipa={agigokkoː},pos={名},acc={F}]}%
	{%
		\MeaningsBlock{%
	\ryumeaning%
		{三十三年忌}{最後の法事。\emb{あぎしょっこー}とも}{}%
	}%
	}%
	{\tailPrint[]}%
\entry
	{\headPrint[kana={あきさみよー},ipa={aki̥samijoː},pos={感}]}%
	{%
		\MeaningsBlock{%
	\ryumeaning%
		{なんたることか}{}{}%
	}%
	}%
	{\tailPrint[]}%
\entry
	{\headPrint[kana={あぎさり},ipa={agisari},pos={名}]}%
	{%
		\MeaningsBlock{%
	\ryumeaning%
		{明け方。夜明け前}{}{}%
	}%
	}%
	{\tailPrint[]}%
\entry
	{\headPrint[kana={あぎひ\tl{˥˩}},ipa={agihi},pos={名},acc={F}]}%
	{%
		\MeaningsBlock{%
	\ryumeaning%
		{空家}{}{}%
	}%
	}%
	{\tailPrint[]}%
\entry
	{\headPrint[kana={あぎまーすん\tl{˥˩}},ipa={agimaːsuɴ},pos={動},rui={2},para={あぎまーはぬ},acc={F}]}%
	{%
		\MeaningsBlock{%
	\ryumeaning%
		{騙す。欺く}{}{}%
	}%
	}%
	{\tailPrint[]}%
\entry
	{\headPrint[kana={あぎやしき\tl{˥˩}},ipa={agijaʃi̥ki},pos={名},acc={F}]}%
	{%
		\MeaningsBlock{%
	\ryumeaning%
		{空き屋敷}{}{}%
	}%
	}%
	{\tailPrint[]}%
\entry
	{\headPrint[kana={あきりるん\tl{˥˩}},ipa={akiriruɴ},pos={動},rui={3},acc={F}]}%
	{%
		\MeaningsBlock{%
	\ryumeaning%
		{飽きる}{}{}%
	}%
	}%
	{\tailPrint[]}%
\entry
	{\headPrint[kana={あきりん\tl{˥˩}},ipa={aki̥riɴ},pos={動},rui={3},acc={F}]}%
	{%
		\MeaningsBlock{%
	\ryumeaning%
		{飽きる}{}{}%
	}%
	}%
	{\tailPrint[]}%
\entry
	{\headPrint[kana={あきるん\tl{˥˩}},ipa={akiruɴ},pos={動},rui={3},acc={F}]}%
	{%
		\MeaningsBlock{%
	\ryumeaning%
		{呆れ果てる。呆れる}{}{}%
	}%
	}%
	{\tailPrint[]}%
\entry
	{\headPrint[kana={あぎるん\tl{˥˩}},ipa={agiruɴ},pos={動},rui={3},acc={F}]}%
	{%
		\MeaningsBlock{%
	\ryumeaning%
		{（戸を）開ける}{}{}%
	}%
	}%
	{\tailPrint[]}%
\entry
	{\headPrint[kana={あぎるん\tl{˥˩}},ipa={agiruɴ},pos={動},rui={3},acc={F}]}%
	{%
		\MeaningsBlock{%
	\ryumeaning%
		{（油で）揚げる}{}{}%
	}%
	}%
	{\tailPrint[]}%
\entry
	{\headPrint[kana={あぎん\tl{˥˩}},ipa={agiɴ},pos={動},rui={3},para={あぐぬ},acc={F}]}%
	{%
		\MeaningsBlock{%
	\ryumeaning%
		{開ける}{}{}%
	}%
	}%
	{\tailPrint[]}%
\entry
	{\headPrint[kana={あく\tl{˥}},ipa={aku},pos={名},acc={H}]}%
	{%
		\MeaningsBlock{%
	\ryumeaning%
		{悪いこと。悪口}{}{}%
	}%
	}%
	{\tailPrint[]}%
\entry
	{\headPrint[kana={あぐん\tl{˥˩}},ipa={aguɴ},pos={動},rui={1},para={あがぬ},acc={F}]}%
	{%
		\MeaningsBlock{%
	\ryumeaning%
		{開く}{}{}%
	}%
	}%
	{\tailPrint[]}%
\entry
	{\headPrint[kana={あごん\tl{˥˩}},ipa={agoɴ},pos={名},acc={F}]}%
	{%
		\MeaningsBlock{%
	\ryumeaning%
		{アコウ}{樹木名}{}%
	}%
	}%
	{\tailPrint[]}%
\entry
	{\headPrint[kana={あざーぎるん\tl{˥}},ipa={adzaːgiruɴ},pos={動},acc={H}]}%
	{%
		\MeaningsBlock{%
	\ryumeaning%
		{片付けてきれいにする}{}{}%
	}%
	}%
	{\tailPrint[]}%
\entry
	{\headPrint[kana={あざぎしゃ\tl{˥}はん},ipa={adzagiʃahaɴ},pos={形},para={あざきしゃへぬ},acc={H}]}%
	{%
		\MeaningsBlock{%
	\ryumeaning%
		{きれい好きだ。清潔だ}{}{}%
	}%
	}%
	{\tailPrint[]}%
\entry
	{\headPrint[kana={あざぎしゃ\tl{˥}はん},ipa={adzagiʃahaɴ},pos={形},acc={H}]}%
	{%
		\MeaningsBlock{%
	\ryumeaning%
		{こざっぱりする}{}{}%
	}%
	}%
	{\tailPrint[]}%
\entry
	{\headPrint[kana={あさどぅり\tl{˥}},ipa={asaduri},pos={名},acc={H}]}%
	{%
		\MeaningsBlock{%
	\ryumeaning%
		{\ruby[g]{朝凪}{あさなぎ}}{}{}%
	}%
	}%
	{\tailPrint[]}%
\entry
	{\headPrint[kana={あさ\tl{˥˩}はん},ipa={asahaɴ},pos={形},para={あさへぬ},acc={F}]}%
	{%
		\MeaningsBlock{%
	\ryumeaning%
		{浅い}{海や川の水深をさす}{}%
	}%
	}%
	{\tailPrint[]}%
\entry
	{\headPrint[kana={あざまぐ\tl{˥}},ipa={adzamagu},pos={名},acc={H}]}%
	{%
		\MeaningsBlock{%
	\ryumeaning%
		{按司}{島の古語。古代の豪族の名につく}{}%
	}%
	}%
	{\tailPrint[]}%
\entry
	{\headPrint[kana={あさむぬ\tl{˥}},ipa={asamunu},pos={名},acc={H}]}%
	{%
		\MeaningsBlock{%
	\ryumeaning%
		{朝食}{}{}%
	}%
	}%
	{\tailPrint[]}%
\entry
	{\headPrint[kana={あさらごー\tl{˥}},ipa={asaragoː},pos={名},acc={H}]}%
	{%
		\MeaningsBlock{%
	\ryumeaning%
		{潮干狩り}{}{}%
	}%
	}%
	{\tailPrint[]}%
\entry
	{\headPrint[kana={あざらはーん},ipa={adzarahaːɴ},pos={形},acc={？}]}%
	{%
		\MeaningsBlock{%
	\ryumeaning%
		{（棘木やつるが多くて）通りにくい。（荒れて）通れない}{}{}%
	}%
	}%
	{\tailPrint[]}%
\entry
	{\headPrint[kana={あさるん\tl{˥}},ipa={asaruɴ},pos={動},rui={1},acc={H}]}%
	{%
		\MeaningsBlock{%
	\ryumeaning%
		{漁る。探す}{}{}%
	}%
	}%
	{\tailPrint[]}%
\entry
	{\headPrint[kana={あざん\tl{˥˩}},ipa={adzaɴ},pos={名},acc={F}]}%
	{%
		\MeaningsBlock{%
	\ryumeaning%
		{薊}{海辺の植物名}{}%
	}%
	}%
	{\tailPrint[]}%
\entry
	{\headPrint[kana={あし\tl{˥}},ipa={aʃi},pos={名},acc={H}]}%
	{%
		\MeaningsBlock{%
	\ryumeaning%
		{汗}{}{}%
	}%
	}%
	{\tailPrint[]}%
\entry
	{\headPrint[kana={あじ\tl{˥˩}},ipa={adʒi},pos={名},acc={F}]}%
	{%
		\MeaningsBlock{%
	\ryumeaning%
		{味}{}{}%
	}%
	}%
	{\tailPrint[]}%
\entry
	{\headPrint[kana={あしきるん\tl{˥}},ipa={aʃi̥kiruɴ},pos={動},rui={3},para={あすくぬ},acc={H}]}%
	{%
		\MeaningsBlock{%
	\ryumeaning%
		{預ける}{}{}%
	}%
	}%
	{\tailPrint[]}%
\entry
	{\headPrint[kana={あしけー\tl{˥}},ipa={aʃi̥keː},pos={名},acc={H}]}%
	{%
		\MeaningsBlock{%
	\ryumeaning%
		{シャコガイ}{貝の種類}{}%
	}%
	}%
	{\tailPrint[]}%
\entry
	{\headPrint[kana={あじ\ws{}すん},ipa={adʒi suɴ},pos={句},rui={2b},para={あじさぬ}]}%
	{%
		\MeaningsBlock{%
	\ryumeaning%
		{味わう}{味をみる}{}%
	}%
	}%
	{\tailPrint[]}%
\entry
	{\headPrint[kana={あした\tl{˥}},ipa={aʃi̥ta},pos={名},acc={H}]}%
	{%
		\MeaningsBlock{%
	\ryumeaning%
		{下駄}{東の村では\emb{あすたん}と言う}{}%
	}%
	}%
	{\tailPrint[]}%
\entry
	{\headPrint[kana={あしぴ\tl{˥˩}},ipa={aʃi̥pi},pos={名},acc={F}]}%
	{%
		\MeaningsBlock{%
	\ryumeaning%
		{遊び}{大人が歌や踊りで楽しむこと。宗教的な行事にも言う}{}%
	}%
	}%
	{\tailPrint[]}%
\entry
	{\headPrint[kana={あし\ws{}ふきん},ipa={aʃi fu̥kiɴ},pos={句},rui={3},para={あしふくぬ}]}%
	{%
		\MeaningsBlock{%
	\ryumeaning%
		{汗ばむ。汗をかく}{}{}%
	}%
	}%
	{\tailPrint[]}%
\entry
	{\headPrint[kana={あし\tl{˥}ふさ\tl{˥}はーん},ipa={aʃifu̥sahaːɴ},pos={形},acc={H+H}]}%
	{%
		\MeaningsBlock{%
	\ryumeaning%
		{汗臭い}{}{}%
	}%
	}%
	{\tailPrint[]}%
\entry
	{\headPrint[kana={あじまぎ\tl{˥}},ipa={adʒimagi},pos={名},acc={H}]}%
	{%
		\MeaningsBlock{%
	\ryumeaning%
		{たすき}{}{}%
	}%
	}%
	{\tailPrint[]}%
\entry
	{\headPrint[kana={あしみじ\tl{˥}},ipa={aʃimidʒi},pos={名},acc={H}]}%
	{%
		\MeaningsBlock{%
	\ryumeaning%
		{汗水}{}{}%
	}%
	}%
	{\tailPrint[]}%
\entry
	{\headPrint[kana={あしみるん\tl{˥}},ipa={aʃi̥miruɴ},pos={動},rui={3},para={あすむぬ},acc={H}]}%
	{%
		\MeaningsBlock{%
	\ryumeaning%
		{集める}{}{}%
	}%
	}%
	{\tailPrint[]}%
\entry
	{\headPrint[kana={あしゃぼー\tl{˥}},ipa={aʃaboː},pos={名},acc={H}]}%
	{%
		\MeaningsBlock{%
	\ryumeaning%
		{汗疹}{}{}%
	}%
	}%
	{\tailPrint[]}%
\entry
	{\headPrint[kana={あすぅとぅ\tl{˥}},ipa={asɨ̥tu},pos={名},acc={H}]}%
	{%
		\MeaningsBlock{%
	\ryumeaning%
		{明後日}{}{}%
	}%
	}%
	{\tailPrint[]}%
\entry
	{\headPrint[kana={あすかるん\tl{˥}},ipa={asu̥karuɴ},pos={動},rui={1},para={あすからぬ},acc={H}]}%
	{%
		\MeaningsBlock{%
	\ryumeaning%
		{預かる}{}{}%
	}%
	}%
	{\tailPrint[]}%
\entry
	{\headPrint[kana={あすぺー\tl{˥}},ipa={asu̥peː},pos={名},acc={H}]}%
	{%
		\MeaningsBlock{%
	\ryumeaning%
		{泡}{}{}%
	}%
	}%
	{\tailPrint[]}%
\entry
	{\headPrint[kana={あずま\tl{˥˩}はん},ipa={adzumahaɴ},pos={形},para={あじまへぬ},acc={F}]}%
	{%
		\MeaningsBlock{%
	\ryumeaning%
		{甘い}{}{}%
	}%
	}%
	{\tailPrint[]}%
\entry
	{\headPrint[kana={あすまるん\tl{˥}},ipa={asu̥maruɴ},pos={動},rui={1},para={あすまらぬ},acc={H}]}%
	{%
		\MeaningsBlock{%
	\ryumeaning%
		{集まる}{}{}%
	}%
	}%
	{\tailPrint[]}%
\entry
	{\headPrint[kana={あた\tl{˥˩}},ipa={ata},pos={副},acc={F}]}%
	{%
		\MeaningsBlock{%
	\ryumeaning%
		{急に。突然}{}{}%
	}%
	}%
	{\tailPrint[]}%
\entry
	{\headPrint[kana={あ\josho{}たー\kako{}すま},ipa={ataːsu̥ma},pos={副},acc={LHL}]}%
	{%
		\MeaningsBlock{%
	\ryumeaning%
		{いっとき。\ruby[g]{暫}{しばら}く}{}{}%
	}%
	}%
	{\tailPrint[]}%
\entry
	{\headPrint[kana={あたあみ},ipa={ata.ami},pos={名}]}%
	{%
		\MeaningsBlock{%
	\ryumeaning%
		{\ruby[g]{俄}{にわか}雨}{}{}%
	}%
	}%
	{\tailPrint[]}%
\entry
	{\headPrint[kana={あたしに},ipa={ataʃini},pos={名}]}%
	{%
		\MeaningsBlock{%
	\ryumeaning%
		{急死}{}{}%
	}%
	}%
	{\tailPrint[]}%
\entry
	{\headPrint[kana={\josho{}あためー\kako{}すん},ipa={atameːsuɴ},pos={動},rui={2},para={あためーさぬ},acc={H L}]}%
	{%
		\MeaningsBlock{%
	\ryumeaning%
		{保管する。大事にしまう}{}{}%
	}%
	}%
	{\tailPrint[]}%
\entry
	{\headPrint[kana={あだら\tl{˥˩}},ipa={adara},pos={-},acc={F}]}%
	{%
		\MeaningsBlock{%
	\ryumeaning%
		{汚い}{体が汚れているときに}{}%
	}%
	}%
	{\tailPrint[]}%
\entry
	{\headPrint[kana={あたら\tl{˥˩}さん},ipa={atarasaɴ},pos={形},acc={F}]}%
	{%
		\MeaningsBlock{%
	\ryumeaning%
		{惜しい。可愛らしい}{}{}%
	}%
	}%
	{\tailPrint[]}%
\entry
	{\headPrint[kana={あたらは\tl{˥˩}},ipa={ataraha},pos={句},acc={F}]}%
	{%
		\MeaningsBlock{%
	\ryumeaning%
		{大切に}{}{}%
	}%
	}%
	{\tailPrint[]}%
\entry
	{\headPrint[kana={あだり\tl{˥˩}},ipa={adari},pos={副},acc={F}]}%
	{%
		\MeaningsBlock{%
	\ryumeaning%
		{いたずらに。無為に}{}{}%
	}%
	}%
	{\tailPrint[]}%
\entry
	{\headPrint[kana={あたりゃん},ipa={atarjaɴ},pos={句}]}%
	{%
		\MeaningsBlock{%
	\ryumeaning%
		{当たり。正しい}{正解の場合に使用する。\emb{あたるん}の完了形}{}%
	}%
	}%
	{\tailPrint[]}%
\entry
	{\headPrint[kana={あたるん\tl{˥˩}},ipa={ataruɴ},pos={動},rui={1},para={あたらぬ},acc={F}]}%
	{%
		\MeaningsBlock{%
	\ryumeaning%
		{当たる。正しい。つり合う}{合格にも言う}{}%
	}%
	}%
	{\tailPrint[]}%
\entry
	{\headPrint[kana={あたるん\tl{˥˩}},ipa={ataruɴ},pos={動},rui={1},para={あたらぬ},acc={F}]}%
	{%
		\MeaningsBlock{%
	\ryumeaning%
		{身に応える}{}{}%
	}%
	}%
	{\tailPrint[]}%
\entry
	{\headPrint[kana={あちぃゆ\tl{˥}},ipa={atʃiju},pos={名},acc={H}]}%
	{%
		\MeaningsBlock{%
	\ryumeaning%
		{湯。熱湯}{}{}%
	}%
	}%
	{\tailPrint[]}%
\entry
	{\headPrint[kana={あちすん\tl{˥˩}},ipa={atʃi̥suɴ},pos={動},rui={2},para={あちさぬ},acc={F}]}%
	{%
		\MeaningsBlock{%
	\ryumeaning%
		{当てにする。期待する}{}{}%
	}%
	}%
	{\tailPrint[]}%
\entry
	{\headPrint[kana={あちらいん\tl{˥}},ipa={atʃiraiɴ},pos={動},rui={2},para={あちらはぬ},acc={H}]}%
	{%
		\MeaningsBlock{%
	\ryumeaning%
		{誂える}{}{}%
	}%
	}%
	{\tailPrint[]}%
\entry
	{\headPrint[kana={あちるん\tl{˥˩}},ipa={atʃiruɴ},pos={動},rui={3},para={あとぅぬ},acc={F}]}%
	{%
		\MeaningsBlock{%
	\ryumeaning%
		{当てる}{}{}%
	}%
	}%
	{\tailPrint[]}%
\entry
	{\headPrint[kana={あつぁーすん\tl{˥}},ipa={atsaːsuɴ},pos={動},rui={2},para={あつぁーはぬ},acc={H}]}%
	{%
		\MeaningsBlock{%
	\ryumeaning%
		{\ruby[g]{温}{あたた}める}{}{}%
	}%
	}%
	{\tailPrint[]}%
\entry
	{\headPrint[kana={あつぁあすとぅ\tl{˥}},ipa={atsa.asu̥tu},pos={名},acc={H}]}%
	{%
		\MeaningsBlock{%
	\ryumeaning%
		{明後日}{}{}%
	\ryumeaning%
		{近いうち}{}{}%
	}%
	}%
	{\tailPrint[]}%
\entry
	{\headPrint[kana={あつぁすとぅむち\tl{˥}},ipa={atsasu̥tumutʃi},pos={名},acc={H}]}%
	{%
		\MeaningsBlock{%
	\ryumeaning%
		{明日の朝。翌朝}{}{}%
	}%
	}%
	{\tailPrint[]}%
\entry
	{\headPrint[kana={あつぁは\tl{˥}\ws{}なるん\tl{˩˥}},ipa={atsaha naruɴ},pos={句},rui={1},para={あつぁは　ならぬ},acc={H R}]}%
	{%
		\MeaningsBlock{%
	\ryumeaning%
		{熱くなる。温まる}{}{}%
	}%
	}%
	{\tailPrint[]}%
\entry
	{\headPrint[kana={あつぁむさ\tl{˥}はん},ipa={atsamusahaɴ},pos={形},acc={H}]}%
	{%
		\MeaningsBlock{%
	\ryumeaning%
		{暑がりやだ}{}{}%
	}%
	}%
	{\tailPrint[]}%
\entry
	{\headPrint[kana={あつぁゆ\tl{˥}},ipa={atsaju},pos={名},acc={H}]}%
	{%
		\MeaningsBlock{%
	\ryumeaning%
		{明日の夜。明晩}{}{}%
	}%
	}%
	{\tailPrint[]}%
\entry
	{\headPrint[kana={あっつぁー\tl{˥}},ipa={attsaː},pos={名},acc={H}]}%
	{%
		\MeaningsBlock{%
	\ryumeaning%
		{あす。明日}{}{}%
	}%
	}%
	{\tailPrint[]}%
\entry
	{\headPrint[kana={あっつぁすん\tl{˥}},ipa={attsasuɴ},pos={動},rui={2},para={あつぁはぬ},acc={H}]}%
	{%
		\MeaningsBlock{%
	\ryumeaning%
		{（食べ物を）\ruby[g]{温}{あたた}める}{}{}%
	}%
	}%
	{\tailPrint[]}%
\entry
	{\headPrint[kana={あっつぁ\tl{˥}はん},ipa={attsahaɴ},pos={形},para={あっあへぬ},acc={H}]}%
	{%
		\MeaningsBlock{%
	\ryumeaning%
		{暑い。熱い}{}{}%
	}%
	}%
	{\tailPrint[]}%
\entry
	{\headPrint[kana={あ\josho{}てー\kako{}し},ipa={ateːʃi},pos={副},acc={LHL}]}%
	{%
		\MeaningsBlock{%
	\ryumeaning%
		{温かく。温く}{}{}%
	}%
	}%
	{\tailPrint[]}%
\entry
	{\headPrint[kana={あとぅ\tl{˥}},ipa={atu},pos={名},acc={H}]}%
	{%
		\MeaningsBlock{%
	\ryumeaning%
		{後}{}{}%
	\ryumeaning%
		{跡}{}{}%
	}%
	}%
	{\tailPrint[]}%
\entry
	{\headPrint[kana={あどぅ\tl{˥}},ipa={adu},pos={名},acc={H}]}%
	{%
		\MeaningsBlock{%
	\ryumeaning%
		{\ruby[g]{踵}{かかと}}{}{}%
	}%
	}%
	{\tailPrint[]}%
\entry
	{\headPrint[kana={あとぅあとぅ\tl{˩˥}},ipa={atu.atu},pos={副},acc={重複}]}%
	{%
		\MeaningsBlock{%
	\ryumeaning%
		{後々。行く末}{}{}%
	}%
	}%
	{\tailPrint[]}%
\entry
	{\headPrint[kana={あとぅ\tl{˥}ちぎ\tl{˥˩}},ipa={atutʃigi},pos={名},acc={H+F}]}%
	{%
		\MeaningsBlock{%
	\ryumeaning%
		{後継ぎ。後継者}{}{}%
	}%
	}%
	{\tailPrint[]}%
\entry
	{\headPrint[kana={あとさき\tl{˥˩}},ipa={atosḁki},pos={名},acc={F}]}%
	{%
		\MeaningsBlock{%
	\ryumeaning%
		{後先。前後}{}{}%
	}%
	}%
	{\tailPrint[note={アクセントは「平板型」と「下降型」の2単位の可能性もある}]}%
\entry
	{\headPrint[kana={あなどるん\tl{˥}},ipa={anadoruɴ},pos={動},rui={1},para={あなどらぬ},acc={H}]}%
	{%
		\MeaningsBlock{%
	\ryumeaning%
		{侮る。嘲る}{}{}%
	}%
	}%
	{\tailPrint[]}%
\entry
	{\headPrint[kana={\josho{}あなばり\kako{}ひー},ipa={anabarihiː},pos={名},acc={H+L}]}%
	{%
		\MeaningsBlock{%
	\ryumeaning%
		{掘っ立て小屋}{}{}%
	}%
	}%
	{\tailPrint[]}%
\entry
	{\headPrint[kana={あなぶ\tl{˥}},ipa={anabu},pos={名},acc={H}]}%
	{%
		\MeaningsBlock{%
	\ryumeaning%
		{沼。溜め池}{牛馬の飲み水や農具の洗い場、農道にそって点在}{}%
	}%
	}%
	{\tailPrint[]}%
\entry
	{\headPrint[kana={あば\tl{˥}},ipa={aba},pos={名},acc={H}]}%
	{%
		\MeaningsBlock{%
	\ryumeaning%
		{油。食油。燃料油}{}{}%
	}%
	}%
	{\tailPrint[]}%
\entry
	{\headPrint[kana={あばずー\tl{˥}さーん},ipa={abadzuːsaːɴ},pos={形},acc={H}]}%
	{%
		\MeaningsBlock{%
	\ryumeaning%
		{脂っこい}{}{}%
	}%
	}%
	{\tailPrint[]}%
\entry
	{\headPrint[kana={あばたしみるん},ipa={abataʃi̥miruɴ},pos={動},rui={3},para={あばたすむぬ}]}%
	{%
		\MeaningsBlock{%
	\ryumeaning%
		{急き立てる}{}{}%
	}%
	}%
	{\tailPrint[]}%
\entry
	{\headPrint[kana={あばだり\tl{˥}},ipa={abadari},pos={名},acc={H}]}%
	{%
		\MeaningsBlock{%
	\ryumeaning%
		{裸。裸体}{}{}%
	}%
	}%
	{\tailPrint[]}%
\entry
	{\headPrint[kana={あばちかんち\tl{˥}\ws{}すん\tl{˥˩}},ipa={abatʃikantʃi suɴ},pos={句},acc={H F}]}%
	{%
		\MeaningsBlock{%
	\ryumeaning%
		{大慌てする。慌てふためく}{}{}%
	}%
	}%
	{\tailPrint[]}%
\entry
	{\headPrint[kana={あばちるん\tl{˥˩}},ipa={abatʃiruɴ},pos={動},rui={3},acc={F}]}%
	{%
		\MeaningsBlock{%
	\ryumeaning%
		{慌てる。急ぐ}{}{}%
	}%
	}%
	{\tailPrint[]}%
\entry
	{\headPrint[kana={あばりしゃ\tl{˥}はん},ipa={abariʃahaɴ},pos={形},para={あばりしゃへぬ},acc={H}]}%
	{%
		\MeaningsBlock{%
	\ryumeaning%
		{（容姿が）美しい。美人だ}{人の容姿の美しさに言う}{}%
	}%
	}%
	{\tailPrint[]}%
\entry
	{\headPrint[kana={あばんぐん\tl{˥˩}},ipa={abaŋguɴ},pos={動},rui={1},para={あばんがぬ},acc={F}]}%
	{%
		\MeaningsBlock{%
	\ryumeaning%
		{仰向く。仰向けになる}{}{}%
	}%
	}%
	{\tailPrint[]}%
\entry
	{\headPrint[kana={あぴら\tl{˥}},ipa={api̥ra},pos={名},acc={H}]}%
	{%
		\MeaningsBlock{%
	\ryumeaning%
		{家鴨}{}{}%
	}%
	}%
	{\tailPrint[]}%
\entry
	{\headPrint[kana={あぶ\tl{˥˩}},ipa={abu},pos={名},acc={F}]}%
	{%
		\MeaningsBlock{%
	\ryumeaning%
		{ドリーネ}{雨水の吸い込み穴}{}%
	}%
	}%
	{\tailPrint[]}%
\entry
	{\headPrint[kana={あふぁ\tl{˥}さん},ipa={afasaɴ},pos={形},para={あふぁさーへぬ},acc={H}]}%
	{%
		\MeaningsBlock{%
	\ryumeaning%
		{薄味だ。塩味不足}{}{}%
	}%
	}%
	{\tailPrint[]}%
\entry
	{\headPrint[kana={あふくん\tl{˥˩}},ipa={afu̥kuɴ},pos={動},rui={1},para={あふかぬ},acc={F}]}%
	{%
		\MeaningsBlock{%
	\ryumeaning%
		{息切れになる。息がはずむ}{}{}%
	}%
	}%
	{\tailPrint[]}%
\entry
	{\headPrint[kana={あぶし\tl{˥}},ipa={abuʃi},pos={名},acc={H}]}%
	{%
		\MeaningsBlock{%
	\ryumeaning%
		{畦。田の畔}{}{}%
	}%
	}%
	{\tailPrint[]}%
\entry
	{\headPrint[kana={あふらん},ipa={afuraɴ},pos={動（継）},rui={3},para={あふるぬ},acc={-}]}%
	{%
		\MeaningsBlock{%
	\ryumeaning%
		{溢れる}{}{}%
	}%
	}%
	{\tailPrint[]}%
\entry
	{\headPrint[kana={あぶるん\tl{˥}},ipa={aburuɴ},pos={動},rui={1},para={あぶらぬ},acc={H}]}%
	{%
		\MeaningsBlock{%
	\ryumeaning%
		{炙る}{}{}%
	\ryumeaning%
		{焼く}{}{}%
	}%
	}%
	{\tailPrint[]}%
\entry
	{\headPrint[kana={あぶゎ\tl{˥}},ipa={abwa},pos={名},acc={H}]}%
	{%
		\MeaningsBlock{%
	\ryumeaning%
		{母。母親}{}{}%
	}%
	}%
	{\tailPrint[]}%
\entry
	{\headPrint[kana={あま\tl{˥}},ipa={ama},pos={名},acc={H}]}%
	{%
		\MeaningsBlock{%
	\ryumeaning%
		{姉}{}{}%
	}%
	}%
	{\tailPrint[]}%
\entry
	{\headPrint[kana={あまいるん\tl{˥˩}},ipa={amairuɴ},pos={動},rui={3},para={あまいるぬ},acc={F}]}%
	{%
		\MeaningsBlock{%
	\ryumeaning%
		{喜ぶ。楽しむ}{}{}%
	}%
	}%
	{\tailPrint[]}%
\entry
	{\headPrint[kana={あまおしき\tl{˥}},ipa={ama.oʃi̥ki},pos={名},acc={H}]}%
	{%
		\MeaningsBlock{%
	\ryumeaning%
		{雨天。曇天}{}{}%
	}%
	}%
	{\tailPrint[]}%
\entry
	{\headPrint[kana={あま\tl{˥}くま\tl{˥˩}},ipa={amaku̥ma},pos={副},acc={H+F}]}%
	{%
		\MeaningsBlock{%
	\ryumeaning%
		{あちらこちら}{}{}%
	}%
	}%
	{\tailPrint[]}%
\entry
	{\headPrint[kana={あまじ\tl{˥}},ipa={amadʒi},pos={名},acc={H}]}%
	{%
		\MeaningsBlock{%
	\ryumeaning%
		{髪。頭髪}{}{}%
	}%
	}%
	{\tailPrint[]}%
\entry
	{\headPrint[kana={あますな\tl{˥}},ipa={amasu̥na},pos={名},acc={H}]}%
	{%
		\MeaningsBlock{%
	\ryumeaning%
		{サトウキビ。甘蔗}{}{}%
	}%
	}%
	{\tailPrint[]}%
\entry
	{\headPrint[kana={あますん},ipa={amasuɴ},pos={動},rui={2},para={あまはぬ}]}%
	{%
		\MeaningsBlock{%
	\ryumeaning%
		{浴びせる}{}{}%
	}%
	}%
	{\tailPrint[]}%
\entry
	{\headPrint[kana={あまだり\tl{˥}},ipa={amadari},pos={名},acc={H}]}%
	{%
		\MeaningsBlock{%
	\ryumeaning%
		{庇}{}{}%
	}%
	}%
	{\tailPrint[]}%
\entry
	{\headPrint[kana={あまだりん\tl{˥}},ipa={amadariɴ},pos={動},acc={H}]}%
	{%
		\MeaningsBlock{%
	\ryumeaning%
		{滴る}{}{}%
	}%
	}%
	{\tailPrint[]}%
\entry
	{\headPrint[kana={あまっすくる\tl{˥}},ipa={amasu̥kuru},pos={名},acc={H}]}%
	{%
		\MeaningsBlock{%
	\ryumeaning%
		{頭。頭部}{}{}%
	}%
	}%
	{\tailPrint[]}%
\entry
	{\headPrint[kana={あまっすくる\tl{˥}やみ\tl{˩˥}},ipa={amasu̥kurujami},pos={句},acc={H+R}]}%
	{%
		\MeaningsBlock{%
	\ryumeaning%
		{頭が痛い。頭痛}{}{}%
	}%
	}%
	{\tailPrint[]}%
\entry
	{\headPrint[kana={あまぱんぎ\tl{˥}},ipa={amapaŋgi},pos={名},acc={H}]}%
	{%
		\MeaningsBlock{%
	\ryumeaning%
		{軒端。軒先}{}{}%
	}%
	}%
	{\tailPrint[]}%
\entry
	{\headPrint[kana={あまふもん\tl{˥}},ipa={amafu̥moɴ},pos={名},acc={H}]}%
	{%
		\MeaningsBlock{%
	\ryumeaning%
		{雨雲}{}{}%
	}%
	}%
	{\tailPrint[]}%
\entry
	{\headPrint[kana={あまみじ\tl{˥˩}},ipa={amamidʒi},pos={名},acc={F}]}%
	{%
		\MeaningsBlock{%
	\ryumeaning%
		{真水。淡水}{}{}%
	}%
	}%
	{\tailPrint[]}%
\entry
	{\headPrint[kana={あまむり\tl{˥}},ipa={amamuri},pos={名},acc={H}]}%
	{%
		\MeaningsBlock{%
	\ryumeaning%
		{雨漏り}{}{}%
	}%
	}%
	{\tailPrint[]}%
\entry
	{\headPrint[kana={あまらすん\tl{˥}},ipa={amarasuɴ},pos={動},rui={2},para={あまらはぬ},acc={H}]}%
	{%
		\MeaningsBlock{%
	\ryumeaning%
		{余す。余らす}{}{}%
	}%
	}%
	{\tailPrint[]}%
\entry
	{\headPrint[kana={あまり\tl{˥}},ipa={amari},pos={名},acc={H}]}%
	{%
		\MeaningsBlock{%
	\ryumeaning%
		{余り。余分}{}{}%
	}%
	}%
	{\tailPrint[]}%
\entry
	{\headPrint[kana={あまるん\tl{˥}},ipa={amaruɴ},pos={動},rui={1},para={あまらぬ},acc={H}]}%
	{%
		\MeaningsBlock{%
	\ryumeaning%
		{余る}{}{}%
	}%
	}%
	{\tailPrint[]}%
\entry
	{\headPrint[kana={あまん\tl{˥}},ipa={amaɴ},pos={名},acc={H}]}%
	{%
		\MeaningsBlock{%
	\ryumeaning%
		{ヤドカリ}{}{}%
	}%
	}%
	{\tailPrint[]}%
\entry
	{\headPrint[kana={あみ\tl{˥}},ipa={ami},pos={名},acc={H}]}%
	{%
		\MeaningsBlock{%
	\ryumeaning%
		{雨}{「梅雨」は\emb{ゆどあみ}}{}%
	}%
	}%
	{\tailPrint[]}%
\entry
	{\headPrint[kana={あみじわー\tl{˥}},ipa={amidʒiwaː},pos={名},acc={H}]}%
	{%
		\MeaningsBlock{%
	\ryumeaning%
		{祈年祭}{次年の豊作を祈る豊年祭}{}%
	}%
	}%
	{\tailPrint[]}%
\entry
	{\headPrint[kana={あみにげー\tl{˥}},ipa={aminigeː},pos={名},acc={H}]}%
	{%
		\MeaningsBlock{%
	\ryumeaning%
		{雨乞い祈願}{\emb{ふつぁまらー}「仮面神」が登場}{}%
	}%
	}%
	{\tailPrint[]}%
\entry
	{\headPrint[kana={あみん\tl{˥˩}},ipa={amiɴ},pos={動},rui={3},para={あむぬ},acc={F}]}%
	{%
		\MeaningsBlock{%
	\ryumeaning%
		{浴びる}{}{}%
	}%
	}%
	{\tailPrint[]}%
\entry
	{\headPrint[kana={あむん\tl{˥}},ipa={amuɴ},pos={動},rui={1},para={あまぬ},acc={H}]}%
	{%
		\MeaningsBlock{%
	\ryumeaning%
		{編む}{紐で編む。竹の時は\emb{ふむん}}{}%
	}%
	}%
	{\tailPrint[]}%
\entry
	{\headPrint[kana={あやかーるん\tl{˥}},ipa={ajakaːruɴ},pos={動},rui={1},acc={H}]}%
	{%
		\MeaningsBlock{%
	\ryumeaning%
		{肖る}{}{}%
	}%
	}%
	{\tailPrint[]}%
\entry
	{\headPrint[kana={あやかるん\tl{˥}},ipa={ajakaruɴ},pos={動},rui={1},para={あやからぬ},acc={H}]}%
	{%
		\MeaningsBlock{%
	\ryumeaning%
		{肖る}{似ることを願う}{}%
	}%
	}%
	{\tailPrint[]}%
\entry
	{\headPrint[kana={あやっ\tl{˥˩}さーん},ipa={ajassaːɴ},pos={形},para={あやっさーへぬ},acc={F}]}%
	{%
		\MeaningsBlock{%
	\ryumeaning%
		{怪しい。疑わしい}{}{}%
	}%
	}%
	{\tailPrint[]}%
\entry
	{\headPrint[kana={あやっふぁ\tl{˥}さーん},ipa={ajaffwasaːɴ},pos={形},acc={H}]}%
	{%
		\MeaningsBlock{%
	\ryumeaning%
		{薄暗い}{}{}%
	}%
	}%
	{\tailPrint[]}%
\entry
	{\headPrint[kana={あやっふゎみ\tl{˥}},ipa={ajaffwami},pos={名},acc={H}]}%
	{%
		\MeaningsBlock{%
	\ryumeaning%
		{たそがれ時}{}{}%
	}%
	}%
	{\tailPrint[]}%
\entry
	{\headPrint[kana={あやみるん\tl{˥}},ipa={ajamiruɴ},pos={動},rui={3},para={あやむぬ},acc={H}]}%
	{%
		\MeaningsBlock{%
	\ryumeaning%
		{傷つける。損なう}{}{}%
	}%
	}%
	{\tailPrint[]}%
\entry
	{\headPrint[kana={あよー\tl{˥}},ipa={ajoː},pos={名},acc={H}]}%
	{%
		\MeaningsBlock{%
	\ryumeaning%
		{古謡の一種}{}{}%
	}%
	}%
	{\tailPrint[note={宮古語の\emb{あやぐ}～\emb{あーぐ}「（一般的な）歌」に対応。\emb{あやぐ}＞\emb{あやう}＞\emb{あよー}のように変化した}]}%
\entry
	{\headPrint[kana={あら\tl{˥˩}},ipa={ara},pos={名},acc={F}]}%
	{%
		\MeaningsBlock{%
	\ryumeaning%
		{粗。（白米中の）籾粒}{}{}%
	}%
	}%
	{\tailPrint[]}%
\entry
	{\headPrint[kana={あらすん\tl{˥˩}},ipa={arasuɴ},pos={動},rui={2},para={あらはぬ},acc={F}]}%
	{%
		\MeaningsBlock{%
	\ryumeaning%
		{荒らす}{}{}%
	}%
	}%
	{\tailPrint[]}%
\entry
	{\headPrint[kana={あらすん\tl{˥˩}},ipa={arasuɴ},pos={動},rui={2},para={あらはぬ},acc={F}]}%
	{%
		\MeaningsBlock{%
	\ryumeaning%
		{洗う。洗濯する}{}{}%
	}%
	}%
	{\tailPrint[]}%
\entry
	{\headPrint[kana={あらた},ipa={arata},pos={名}]}%
	{%
		\MeaningsBlock{%
	\ryumeaning%
		{東方。東側}{}{}%
	}%
	}%
	{\tailPrint[]}%
\entry
	{\headPrint[kana={あらだてぃるん\tl{˥˩}},ipa={aradatiruɴ},pos={動},rui={3},para={あらだつぬ},acc={F}]}%
	{%
		\MeaningsBlock{%
	\ryumeaning%
		{荒立てる}{}{}%
	}%
	}%
	{\tailPrint[]}%
\entry
	{\headPrint[kana={あらたみるん\tl{˥}},ipa={aratḁmiruɴ},pos={動},rui={3},para={あらたむぬ},acc={H}]}%
	{%
		\MeaningsBlock{%
	\ryumeaning%
		{改める。改善する}{}{}%
	}%
	}%
	{\tailPrint[]}%
\entry
	{\headPrint[kana={あらとぅし\tl{˥}},ipa={aratu̥ʃi},pos={名},acc={H}]}%
	{%
		\MeaningsBlock{%
	\ryumeaning%
		{新年}{}{}%
	}%
	}%
	{\tailPrint[]}%
\entry
	{\headPrint[kana={あらなん\tl{˥˩}},ipa={aranaɴ},pos={名},acc={F}]}%
	{%
		\MeaningsBlock{%
	\ryumeaning%
		{荒波。荒海}{大波は\emb{ぶーなん}\emb{むさん}}{}%
	}%
	}%
	{\tailPrint[]}%
\entry
	{\headPrint[kana={あらぬ},ipa={aranu},pos={句}]}%
	{%
		\MeaningsBlock{%
	\ryumeaning%
		{いいえ}{}{}%
	}%
	}%
	{\tailPrint[]}%
\entry
	{\headPrint[kana={あらぬ\tl{˥}},ipa={aranu},pos={動},rui={irr},acc={H}]}%
	{%
		\MeaningsBlock{%
	\ryumeaning%
		{違う。そうではない}{}{}%
	}%
	}%
	{\tailPrint[]}%
\entry
	{\headPrint[kana={あら\tl{˥˩}はん},ipa={arahaɴ},pos={形},para={あらは　ねーぬ},acc={F}]}%
	{%
		\MeaningsBlock{%
	\ryumeaning%
		{（波、動作、粉末などが）荒い。粗い}{}{}%
	}%
	}%
	{\tailPrint[]}%
\entry
	{\headPrint[kana={あら\tl{˥˩}むに\tl{˩˥}},ipa={aramuni},pos={名},acc={F+R}]}%
	{%
		\MeaningsBlock{%
	\ryumeaning%
		{荒い言葉。暴言}{}{}%
	}%
	}%
	{\tailPrint[]}%
\entry
	{\headPrint[kana={あらもーぎ\tl{˥˩}},ipa={aramoːgi},pos={名},acc={F}]}%
	{%
		\MeaningsBlock{%
	\ryumeaning%
		{大儲け}{}{}%
	}%
	}%
	{\tailPrint[]}%
\entry
	{\headPrint[kana={あらわりるん\tl{˥}},ipa={arawariruɴ},pos={動},rui={3},para={あらわるぬ},acc={H}]}%
	{%
		\MeaningsBlock{%
	\ryumeaning%
		{現れる}{}{}%
	}%
	}%
	{\tailPrint[]}%
\entry
	{\headPrint[kana={ありかた\tl{˥˩}},ipa={arikḁta},pos={名},acc={F}]}%
	{%
		\MeaningsBlock{%
	\ryumeaning%
		{東側。東方}{}{}%
	}%
	}%
	{\tailPrint[]}%
\entry
	{\headPrint[kana={ありしあぎるん\tl{˥˩}},ipa={ariʃi.agiruɴ},pos={動},rui={3},para={ありしあぐぬ},acc={F}]}%
	{%
		\MeaningsBlock{%
	\ryumeaning%
		{開墾する}{}{}%
	}%
	}%
	{\tailPrint[]}%
\entry
	{\headPrint[kana={ありしぴてー\tl{˥˩}},ipa={ariʃipi̥teː},pos={名},acc={F}]}%
	{%
		\MeaningsBlock{%
	\ryumeaning%
		{開墾畑地}{}{}%
	}%
	}%
	{\tailPrint[]}%
\entry
	{\headPrint[kana={ありすむち\tl{˥˩}},ipa={arisu̥mutʃi},pos={名},acc={F}]}%
	{%
		\MeaningsBlock{%
	\ryumeaning%
		{米粉の蒸し菓子}{}{}%
	}%
	}%
	{\tailPrint[]}%
\entry
	{\headPrint[kana={ありむら},ipa={arimura},pos={名}]}%
	{%
		\MeaningsBlock{%
	\ryumeaning%
		{東村}{北村と南村をまとめた呼び方}{}%
	}%
	}%
	{\tailPrint[]}%
\entry
	{\headPrint[kana={ありるん\tl{˥˩}},ipa={ariruɴ},pos={動},rui={3},para={あーるぬ},acc={F}]}%
	{%
		\MeaningsBlock{%
	\ryumeaning%
		{荒れる}{}{}%
	\ryumeaning%
		{暴れる}{}{}%
	}%
	}%
	{\tailPrint[]}%
\entry
	{\headPrint[kana={あるぐん\tl{˥}},ipa={aruguɴ},pos={動},rui={1},para={あるがぬ},acc={H}]}%
	{%
		\MeaningsBlock{%
	\ryumeaning%
		{歩く}{}{}%
	}%
	}%
	{\tailPrint[]}%
\entry
	{\headPrint[kana={あるふた\tl{˥˩}},ipa={arufu̥ta},pos={名},acc={F}]}%
	{%
		\MeaningsBlock{%
	\ryumeaning%
		{\ruby[g]{塵}{ちり}。ごみ}{}{}%
	}%
	}%
	{\tailPrint[]}%
\entry
	{\headPrint[kana={あわり\tl{˥}},ipa={awari},pos={名},acc={H}]}%
	{%
		\MeaningsBlock{%
	\ryumeaning%
		{難儀。苦労}{}{}%
	}%
	}%
	{\tailPrint[]}%
\entry
	{\headPrint[kana={あん\tl{˥}},ipa={aɴ},pos={名},acc={H}]}%
	{%
		\MeaningsBlock{%
	\ryumeaning%
		{粟}{}{}%
	}%
	}%
	{\tailPrint[]}%
\entry
	{\headPrint[kana={あん\tl{˥}},ipa={aɴ},pos={名},acc={H}]}%
	{%
		\MeaningsBlock{%
	\ryumeaning%
		{網}{魚網など}{}%
	}%
	}%
	{\tailPrint[]}%
\entry
	{\headPrint[kana={あん\tl{˥}},ipa={aɴ},pos={名},acc={H}]}%
	{%
		\MeaningsBlock{%
	\ryumeaning%
		{\ruby[g]{餡}{あん}。餡子}{}{}%
	}%
	}%
	{\tailPrint[]}%
\entry
	{\headPrint[kana={あん\tl{˥}},ipa={aɴ},pos={動},rui={irr},para={ねーぬ},acc={H}]}%
	{%
		\MeaningsBlock{%
	\ryumeaning%
		{有る}{}{}%
	}%
	}%
	{\tailPrint[]}%
\entry
	{\headPrint[kana={あんぎるん\tl{˥˩}},ipa={aŋgiruɴ},pos={動},rui={3},para={あんぐぬ},acc={F}]}%
	{%
		\MeaningsBlock{%
	\ryumeaning%
		{上げる}{高いところにあげる}{}%
	}%
	}%
	{\tailPrint[]}%
\entry
	{\headPrint[kana={あんざーりるん\tl{˥}},ipa={andzaːriruɴ},pos={動},rui={3},para={あんざーるぬ},acc={H}]}%
	{%
		\MeaningsBlock{%
	\ryumeaning%
		{もつれる}{}{}%
	}%
	}%
	{\tailPrint[]}%
\entry
	{\headPrint[kana={あんざらすん\tl{˥}},ipa={andzarasuɴ},pos={動},rui={2},para={あんざらはぬ},acc={H}]}%
	{%
		\MeaningsBlock{%
	\ryumeaning%
		{交差させる}{}{}%
	}%
	}%
	{\tailPrint[]}%
\entry
	{\headPrint[kana={あんざらだーぐ\tl{˥}},ipa={andzaradaːgu},pos={名},acc={H}]}%
	{%
		\MeaningsBlock{%
	\ryumeaning%
		{釣竿の工夫}{釣針の根かかり防止工夫}{}%
	}%
	}%
	{\tailPrint[]}%
\entry
	{\headPrint[kana={あんざるん\tl{˥}},ipa={andzaruɴ},pos={動},rui={1},para={あんざらぬ},acc={H}]}%
	{%
		\MeaningsBlock{%
	\ryumeaning%
		{絡まる。もつれる。こんがらがる}{}{}%
	}%
	}%
	{\tailPrint[]}%
\entry
	{\headPrint[kana={あんじるん\tl{˥}},ipa={andʒiruɴ},pos={動},rui={2},acc={H}]}%
	{%
		\MeaningsBlock{%
	\ryumeaning%
		{交差させる。あざなう}{}{}%
	}%
	}%
	{\tailPrint[]}%
\entry
	{\headPrint[kana={あんだしぃ\tl{˥}},ipa={andaʃi},pos={名},acc={H}]}%
	{%
		\MeaningsBlock{%
	\ryumeaning%
		{アダンの気根}{縄のよい材料}{}%
	}%
	}%
	{\tailPrint[]}%
\entry
	{\headPrint[kana={あんだに\tl{˥}},ipa={andani},pos={名},acc={H}]}%
	{%
		\MeaningsBlock{%
	\ryumeaning%
		{\ruby[g]{阿旦}{あだん}}{}{}%
	}%
	}%
	{\tailPrint[]}%
\entry
	{\headPrint[kana={あんだみしゅ\tl{˥}},ipa={andamiʃu},pos={名},acc={H}]}%
	{%
		\MeaningsBlock{%
	\ryumeaning%
		{油味噌}{味噌を油でいため豚肉など入れた保存食}{}%
	}%
	}%
	{\tailPrint[]}%
\entry
	{\headPrint[kana={あんむち\tl{˥}},ipa={ammutʃi},pos={名},acc={H}]}%
	{%
		\MeaningsBlock{%
	\ryumeaning%
		{\ruby[g]{餡}{あん}\ruby[g]{餅}{もち}}{}{}%
	}%
	}%
	{\tailPrint[]}%
\LetterHead{い}
\entry
	{\headPrint[kana={いー\tl{˥}},ipa={iː},pos={名},acc={H}]}%
	{%
		\MeaningsBlock{%
	\ryumeaning%
		{飯。ご飯}{}{}%
	}%
	}%
	{\tailPrint[]}%
\entry
	{\headPrint[kana={\josho{}いー\kako{}くとぅ},ipa={iːku̥tu},pos={名},acc={H+L}]}%
	{%
		\MeaningsBlock{%
	\ryumeaning%
		{良いこと。慶事}{}{}%
	}%
	}%
	{\tailPrint[]}%
\entry
	{\headPrint[kana={いーしきるん\tl{˥˩}},ipa={iːʃi̥kiruɴ},pos={動},rui={3},para={いーすくぬ},acc={F}]}%
	{%
		\MeaningsBlock{%
	\ryumeaning%
		{告げる。知らせる}{}{}%
	}%
	}%
	{\tailPrint[]}%
\entry
	{\headPrint[kana={いーのーすん\tl{˩˥}},ipa={iːnoːsuɴ},pos={動},rui={2},para={いーのーはぬ},acc={R}]}%
	{%
		\MeaningsBlock{%
	\ryumeaning%
		{言いなおす}{}{}%
	}%
	}%
	{\tailPrint[]}%
\entry
	{\headPrint[kana={\josho{}いー\kako{}ば},ipa={iːba},pos={副},acc={HL}]}%
	{%
		\MeaningsBlock{%
	\ryumeaning%
		{好都合}{}{}%
	}%
	}%
	{\tailPrint[]}%
\entry
	{\headPrint[kana={いーばぎ\tl{˥˩}},ipa={iːbagi},pos={名},acc={F}]}%
	{%
		\MeaningsBlock{%
	\ryumeaning%
		{言い訳}{}{}%
	}%
	}%
	{\tailPrint[]}%
\entry
	{\headPrint[kana={いーまーり\tl{˥}},ipa={iːmaːri},pos={名},acc={H}]}%
	{%
		\MeaningsBlock{%
	\ryumeaning%
		{飯茶碗}{}{}%
	}%
	}%
	{\tailPrint[]}%
\entry
	{\headPrint[kana={いーまかすん\tl{˥˩}},ipa={iːmakasuɴ},pos={動},rui={2},para={いーまかはぬ},acc={F}]}%
	{%
		\MeaningsBlock{%
	\ryumeaning%
		{言い負かす}{}{}%
	}%
	}%
	{\tailPrint[]}%
\entry
	{\headPrint[kana={いかすく},ipa={ikasu̥ku},pos={副}]}%
	{%
		\MeaningsBlock{%
	\ryumeaning%
		{どれほど}{}{}%
	}%
	}%
	{\tailPrint[]}%
\entry
	{\headPrint[kana={いがすん\tl{˥}},ipa={igasuɴ},pos={動},rui={2},para={いがはぬ},acc={H}]}%
	{%
		\MeaningsBlock{%
	\ryumeaning%
		{生かす}{}{}%
	\ryumeaning%
		{蘇らせる}{}{}%
	}%
	}%
	{\tailPrint[]}%
\entry
	{\headPrint[kana={いがばり\tl{˥˩}},ipa={igabari},pos={名},acc={F}]}%
	{%
		\MeaningsBlock{%
	\ryumeaning%
		{胸焼け}{}{}%
	}%
	}%
	{\tailPrint[]}%
\entry
	{\headPrint[kana={いがばりゃん},ipa={igabarjaɴ},pos={句}]}%
	{%
		\MeaningsBlock{%
	\ryumeaning%
		{（胸のあたりが何となく）変な具合だ。（胸がむかついて）吐き気がする}{}{}%
	}%
	}%
	{\tailPrint[]}%
\entry
	{\headPrint[kana={いがふちるん\tl{˥˩}},ipa={igafu̥tʃiruɴ},pos={動},rui={3},para={いがふとぅぬ},acc={F}]}%
	{%
		\MeaningsBlock{%
	\ryumeaning%
		{（水を）ぶっかける}{}{}%
	}%
	}%
	{\tailPrint[]}%
\entry
	{\headPrint[kana={いぎだる},ipa={igidaru},pos={句},para={いぐぬ}]}%
	{%
		\MeaningsBlock{%
	\ryumeaning%
		{生きている}{}{}%
	}%
	}%
	{\tailPrint[]}%
\entry
	{\headPrint[kana={いきら\tl{˥}さん},ipa={iki̥rasaɴ},pos={形},para={いきらへぬ},acc={H}]}%
	{%
		\MeaningsBlock{%
	\ryumeaning%
		{少ない。わずか}{}{}%
	}%
	}%
	{\tailPrint[]}%
\entry
	{\headPrint[kana={いぎるん\tl{˥}},ipa={igiruɴ},pos={動},rui={3},para={いぐぬ},acc={H}]}%
	{%
		\MeaningsBlock{%
	\ryumeaning%
		{（花などを）生ける}{}{}%
	}%
	}%
	{\tailPrint[]}%
\entry
	{\headPrint[kana={いぎるん\tl{˥}},ipa={igiruɴ},pos={動},rui={1},para={いぎらぬ},acc={H}]}%
	{%
		\MeaningsBlock{%
	\ryumeaning%
		{生きる}{}{}%
	}%
	}%
	{\tailPrint[]}%
\entry
	{\headPrint[kana={いきろー\tl{˥}},ipa={iki̥roː},pos={名},acc={H}]}%
	{%
		\MeaningsBlock{%
	\ryumeaning%
		{生霊。呪い}{}{}%
	}%
	}%
	{\tailPrint[]}%
\entry
	{\headPrint[kana={いぎん\tl{˥}},ipa={igiɴ},pos={動},rui={3},acc={H}]}%
	{%
		\MeaningsBlock{%
	\ryumeaning%
		{生きる}{}{}%
	}%
	}%
	{\tailPrint[]}%
\entry
	{\headPrint[kana={いげーるん\tl{˥˩}},ipa={igeːruɴ},pos={動},rui={1},para={いげーらぬ},acc={F}]}%
	{%
		\MeaningsBlock{%
	\ryumeaning%
		{行き会う}{}{}%
	}%
	}%
	{\tailPrint[]}%
\entry
	{\headPrint[kana={いさすぅま\tl{˥}},ipa={isasu̥ma},pos={名},acc={H}]}%
	{%
		\MeaningsBlock{%
	\ryumeaning%
		{石垣島}{地名}{}%
	}%
	}%
	{\tailPrint[]}%
\entry
	{\headPrint[kana={いざり\tl{˥˩}},ipa={idzari},pos={名},acc={F}]}%
	{%
		\MeaningsBlock{%
	\ryumeaning%
		{漁り}{}{}%
	}%
	}%
	{\tailPrint[]}%
\entry
	{\headPrint[kana={いざんだ\tl{˩˥}},ipa={idzanda},pos={副},acc={R}]}%
	{%
		\MeaningsBlock{%
	\ryumeaning%
		{一生懸命}{}{}%
	}%
	}%
	{\tailPrint[]}%
\entry
	{\headPrint[kana={いし\tl{˥}},ipa={iʃi},pos={名},acc={H}]}%
	{%
		\MeaningsBlock{%
	\ryumeaning%
		{息}{}{}%
	}%
	}%
	{\tailPrint[]}%
\entry
	{\headPrint[kana={いし\tl{˥˩}},ipa={iʃi},pos={名},acc={F}]}%
	{%
		\MeaningsBlock{%
	\ryumeaning%
		{石}{}{}%
	}%
	}%
	{\tailPrint[]}%
\entry
	{\headPrint[kana={いじぃ\tl{˥˩}},ipa={idʒi},pos={名},acc={F}]}%
	{%
		\MeaningsBlock{%
	\ryumeaning%
		{意地。勇気}{}{}%
	}%
	}%
	{\tailPrint[]}%
\entry
	{\headPrint[kana={いじぃぬ\tl{˥˩}\ws{}ねーぬ\tl{˩˥}},ipa={idʒinu neːnu},pos={句},acc={F R}]}%
	{%
		\MeaningsBlock{%
	\ryumeaning%
		{小胆だ。臆病だ}{勇気がない}{}%
	}%
	}%
	{\tailPrint[]}%
\entry
	{\headPrint[kana={いじぃ\ws{}ねーぬ},ipa={idʒi neːnu},pos={句}]}%
	{%
		\MeaningsBlock{%
	\ryumeaning%
		{小胆だ。臆病だ}{勇気がない}{}%
	}%
	}%
	{\tailPrint[]}%
\entry
	{\headPrint[kana={いしうし\tl{˥˩}},ipa={iʃi.uʃi},pos={名},acc={F}]}%
	{%
		\MeaningsBlock{%
	\ryumeaning%
		{石臼}{沖縄では引き臼が主}{}%
	}%
	}%
	{\tailPrint[]}%
\entry
	{\headPrint[kana={いしかぼ\tl{˥}},ipa={iʃi̥kabo},pos={名},acc={H}]}%
	{%
		\MeaningsBlock{%
	\ryumeaning%
		{ハリセンボン}{とげに覆われた魚}{}%
	}%
	}%
	{\tailPrint[]}%
\entry
	{\headPrint[kana={いしなが\tl{˥}},ipa={iʃi̥naga},pos={名},acc={H}]}%
	{%
		\MeaningsBlock{%
	\ryumeaning%
		{背中}{}{}%
	}%
	}%
	{\tailPrint[]}%
\entry
	{\headPrint[kana={いしふく\tl{˥˩}},ipa={iʃifuku},pos={名},acc={F}]}%
	{%
		\MeaningsBlock{%
	\ryumeaning%
		{小石。礫}{}{}%
	}%
	}%
	{\tailPrint[]}%
\entry
	{\headPrint[kana={いしまーし\tl{˥˩}},ipa={iʃimaːʃi},pos={名},acc={F}]}%
	{%
		\MeaningsBlock{%
	\ryumeaning%
		{石垣}{}{}%
	}%
	}%
	{\tailPrint[]}%
\entry
	{\headPrint[kana={いじみるん\tl{˥˩}},ipa={idʒimiruɴ},pos={動},rui={3},para={いじむぬ},acc={F}]}%
	{%
		\MeaningsBlock{%
	\ryumeaning%
		{虐める。虐待する}{}{}%
	}%
	}%
	{\tailPrint[]}%
\entry
	{\headPrint[kana={いしゃがーたま\tl{˥}},ipa={iʃagaːtama},pos={名},acc={H}]}%
	{%
		\MeaningsBlock{%
	\ryumeaning%
		{幼い子。幼児}{}{}%
	}%
	}%
	{\tailPrint[]}%
\entry
	{\headPrint[kana={いしゃが\tl{˥}はん},ipa={iʃagahaɴ},pos={形},para={いしゃがへぬ},acc={H}]}%
	{%
		\MeaningsBlock{%
	\ryumeaning%
		{小さい}{}{}%
	}%
	}%
	{\tailPrint[]}%
\entry
	{\headPrint[kana={いしゃば\tl{˥˩}},ipa={iʃaba},pos={名},acc={F}]}%
	{%
		\MeaningsBlock{%
	\ryumeaning%
		{オニオコゼ}{刺に毒を持つ魚}{}%
	}%
	}%
	{\tailPrint[]}%
\entry
	{\headPrint[kana={いしゃん\tl{˥}},ipa={iʃaɴ},pos={名},acc={H}]}%
	{%
		\MeaningsBlock{%
	\ryumeaning%
		{医者}{}{}%
	}%
	}%
	{\tailPrint[]}%
\entry
	{\headPrint[kana={いしゃんが\tl{˥}\ws{}はかるん\tl{˥}},ipa={iʃaŋga hakaruɴ},pos={句},para={いしゃんがはからぬ},acc={H H}]}%
	{%
		\MeaningsBlock{%
	\ryumeaning%
		{治療する}{医者にかかるの義}{}%
	}%
	}%
	{\tailPrint[]}%
\entry
	{\headPrint[kana={いしょー\tl{˥˩}},ipa={iʃoː},pos={名},acc={F}]}%
	{%
		\MeaningsBlock{%
	\ryumeaning%
		{衣装}{踊りなどの衣装を指す}{}%
	}%
	}%
	{\tailPrint[]}%
\entry
	{\headPrint[kana={いしょー\tl{˥˩}},ipa={iʃoː},pos={名},acc={F}]}%
	{%
		\MeaningsBlock{%
	\ryumeaning%
		{漁労}{}{}%
	}%
	}%
	{\tailPrint[note={日本語の「磯」に対応}]}%
\entry
	{\headPrint[kana={いしょーだぐ\tl{˥˩}},ipa={iʃoːdagu},pos={名},acc={F}]}%
	{%
		\MeaningsBlock{%
	\ryumeaning%
		{漁具}{}{}%
	}%
	}%
	{\tailPrint[]}%
\entry
	{\headPrint[kana={いしょーぶさ\tl{˥˩}},ipa={iʃoːbusa},pos={名},acc={F}]}%
	{%
		\MeaningsBlock{%
	\ryumeaning%
		{神行事の漁労係}{}{}%
	}%
	}%
	{\tailPrint[]}%
\entry
	{\headPrint[kana={いしょん\tl{˥˩}},ipa={iʃoɴ},pos={名},acc={F}]}%
	{%
		\MeaningsBlock{%
	\ryumeaning%
		{砂}{}{}%
	}%
	}%
	{\tailPrint[]}%
\entry
	{\headPrint[kana={いしん\tl{˥˩}たま\tl{˥}},ipa={iʃintama},pos={名},acc={F+H}]}%
	{%
		\MeaningsBlock{%
	\ryumeaning%
		{小石。砂利}{\emb{たま}は小さなものの愛称}{}%
	}%
	}%
	{\tailPrint[]}%
\entry
	{\headPrint[kana={いすがすん\tl{˥}},ipa={isugasuɴ},pos={動},rui={2},para={いすがはぬ},acc={H}]}%
	{%
		\MeaningsBlock{%
	\ryumeaning%
		{急がせる}{}{}%
	}%
	}%
	{\tailPrint[]}%
\entry
	{\headPrint[kana={いすぐん\tl{˥}},ipa={isuguɴ},pos={動},rui={1},para={いすがぬ},acc={H}]}%
	{%
		\MeaningsBlock{%
	\ryumeaning%
		{急ぐ}{}{}%
	}%
	}%
	{\tailPrint[]}%
\entry
	{\headPrint[kana={いすぱん\tl{˥}},ipa={isu̥paɴ},pos={名},acc={H}]}%
	{%
		\MeaningsBlock{%
	\ryumeaning%
		{一番}{一番座。一番狂言など}{}%
	}%
	}%
	{\tailPrint[]}%
\entry
	{\headPrint[kana={いすぱんこんぎ\tl{˥}},ipa={isu̥paŋkoŋgi},pos={名},acc={H}]}%
	{%
		\MeaningsBlock{%
	\ryumeaning%
		{一番狂言}{豊作祈願の狂言}{}%
	}%
	}%
	{\tailPrint[]}%
\entry
	{\headPrint[kana={いすぱんどぅし\tl{˥}},ipa={isu̥panduʃi},pos={名},acc={H}]}%
	{%
		\MeaningsBlock{%
	\ryumeaning%
		{親友}{一番親しい友人の意}{}%
	}%
	}%
	{\tailPrint[]}%
\entry
	{\headPrint[kana={いすむし\tl{˥}},ipa={isu̥muʃi},pos={名},acc={H}]}%
	{%
		\MeaningsBlock{%
	\ryumeaning%
		{生き物。動物}{}{}%
	}%
	}%
	{\tailPrint[]}%
\entry
	{\headPrint[kana={いた\tl{˥}},ipa={ita},pos={名},acc={H}]}%
	{%
		\MeaningsBlock{%
	\ryumeaning%
		{板}{}{}%
	}%
	}%
	{\tailPrint[]}%
\entry
	{\headPrint[kana={いたすきばら\tl{˥}},ipa={itasu̥kibara},pos={名},acc={H}]}%
	{%
		\MeaningsBlock{%
	\ryumeaning%
		{悪霊払いの一つ}{悪霊を払う行事。お盆の翌日に行う}{}%
	}%
	}%
	{\tailPrint[]}%
\entry
	{\headPrint[kana={いたちり\tl{˥˩}},ipa={itatʃiri},pos={-},acc={F}]}%
	{%
		\MeaningsBlock{%
	\ryumeaning%
		{（液体や小粒の物を）勢いよく移すさま}{}{}%
	}%
	}%
	{\tailPrint[]}%
\entry
	{\headPrint[kana={いたちるん\tl{˥˩}},ipa={itatʃiruɴ},pos={動},acc={F}]}%
	{%
		\MeaningsBlock{%
	\ryumeaning%
		{（液体を残りなく）\ruby[g]{零}{こぼ}す}{液体を皆こぼすときに言う}{}%
	\ryumeaning%
		{空にする。空ける}{全部移してしまう}{}%
	}%
	}%
	{\tailPrint[]}%
\entry
	{\headPrint[kana={いたふに\tl{˥}},ipa={itafu̥ni},pos={名},acc={H}]}%
	{%
		\MeaningsBlock{%
	\ryumeaning%
		{くり舟}{}{}%
	}%
	}%
	{\tailPrint[]}%
\entry
	{\headPrint[kana={いたますん\tl{˥}},ipa={itamasuɴ},pos={動},rui={2},para={いたまはぬ},acc={H}]}%
	{%
		\MeaningsBlock{%
	\ryumeaning%
		{損なう。傷つける}{}{}%
	}%
	}%
	{\tailPrint[]}%
\entry
	{\headPrint[kana={いたみるん\tl{˥}},ipa={itamiruɴ},pos={動},acc={H}]}%
	{%
		\MeaningsBlock{%
	\ryumeaning%
		{損なう。傷つける}{}{}%
	}%
	}%
	{\tailPrint[]}%
\entry
	{\headPrint[kana={いためー\tl{˥}},ipa={itameː},pos={名},acc={H}]}%
	{%
		\MeaningsBlock{%
	\ryumeaning%
		{床下}{}{}%
	}%
	}%
	{\tailPrint[]}%
\entry
	{\headPrint[kana={いたんだ\tl{˥}},ipa={itanda},pos={名},acc={H}]}%
	{%
		\MeaningsBlock{%
	\ryumeaning%
		{只のもの。無代}{}{}%
	}%
	}%
	{\tailPrint[]}%
\entry
	{\headPrint[kana={いちぃ\tl{˥˩}},ipa={itʃi},pos={名},acc={F}]}%
	{%
		\MeaningsBlock{%
	\ryumeaning%
		{\ruby[g]{何時}{いつ}}{}{}%
	}%
	}%
	{\tailPrint[]}%
\entry
	{\headPrint[kana={い\josho{}ちぃ\kako{}ん},ipa={itʃiɴ},pos={副},acc={LHL}]}%
	{%
		\MeaningsBlock{%
	\ryumeaning%
		{何時も。常に}{}{}%
	}%
	}%
	{\tailPrint[]}%
\entry
	{\headPrint[kana={いちふ\tl{˥}},ipa={itʃifu},pos={名},acc={H}]}%
	{%
		\MeaningsBlock{%
	\ryumeaning%
		{従兄弟}{}{}%
	}%
	}%
	{\tailPrint[]}%
\entry
	{\headPrint[kana={\josho{}いちふ\kako{}ぶい},ipa={itʃi̥fubui},pos={名},acc={H+L}]}%
	{%
		\MeaningsBlock{%
	\ryumeaning%
		{叔父叔母の孫}{}{}%
	}%
	}%
	{\tailPrint[]}%
\entry
	{\headPrint[kana={いちむん\tl{˥˩}},ipa={itʃimuɴ},pos={名},acc={F}]}%
	{%
		\MeaningsBlock{%
	\ryumeaning%
		{一門。一族}{}{}%
	}%
	}%
	{\tailPrint[]}%
\entry
	{\headPrint[kana={いちゅ\tl{˥}},ipa={itʃu},pos={名},acc={H}]}%
	{%
		\MeaningsBlock{%
	\ryumeaning%
		{絹}{}{}%
	}%
	}%
	{\tailPrint[]}%
\entry
	{\headPrint[kana={いちゅすぬ},ipa={itʃusu̥nu},pos={名}]}%
	{%
		\MeaningsBlock{%
	\ryumeaning%
		{絹の着物}{}{}%
	}%
	}%
	{\tailPrint[]}%
\entry
	{\headPrint[kana={いっし\tl{˥}},ipa={iʃʃi},pos={名},acc={H}]}%
	{%
		\MeaningsBlock{%
	\ryumeaning%
		{五つ}{}{}%
	}%
	}%
	{\tailPrint[]}%
\entry
	{\headPrint[kana={いっふぇー\tl{˥}},ipa={iffeː},pos={名},acc={H}]}%
	{%
		\MeaningsBlock{%
	\ryumeaning%
		{ものもらい}{目の腫物}{}%
	}%
	}%
	{\tailPrint[]}%
\entry
	{\headPrint[kana={いとぅ\tl{˥}},ipa={itu},pos={名},acc={H}]}%
	{%
		\MeaningsBlock{%
	\ryumeaning%
		{糸}{}{}%
	}%
	}%
	{\tailPrint[]}%
\entry
	{\headPrint[kana={いなー\tl{˥}},ipa={inaː},pos={名},acc={H}]}%
	{%
		\MeaningsBlock{%
	\ryumeaning%
		{海}{}{}%
	}%
	}%
	{\tailPrint[]}%
\entry
	{\headPrint[kana={いなが\tl{˥}},ipa={inaga},pos={名},acc={H}]}%
	{%
		\MeaningsBlock{%
	\ryumeaning%
		{海}{}{}%
	}%
	}%
	{\tailPrint[]}%
\entry
	{\headPrint[kana={いなさ\tl{˥}ひー\tl{˩˥}},ipa={inasahiː},pos={名},acc={H+R}]}%
	{%
		\MeaningsBlock{%
	\ryumeaning%
		{昔の船待ち小屋}{}{}%
	}%
	}%
	{\tailPrint[]}%
\entry
	{\headPrint[kana={いなしき\tl{˥}},ipa={inaʃi̥ki},pos={名},acc={H}]}%
	{%
		\MeaningsBlock{%
	\ryumeaning%
		{杵}{}{}%
	}%
	}%
	{\tailPrint[]}%
\entry
	{\headPrint[kana={いなむん},ipa={inamuɴ},pos={感}]}%
	{%
		\MeaningsBlock{%
	\ryumeaning%
		{無念。残念}{}{}%
	}%
	}%
	{\tailPrint[]}%
\entry
	{\headPrint[kana={\josho{}いなん\kako{}ぱた},ipa={inampḁta},pos={名},acc={H+L}]}%
	{%
		\MeaningsBlock{%
	\ryumeaning%
		{海辺。海岸付近}{}{}%
	}%
	}%
	{\tailPrint[]}%
\entry
	{\headPrint[kana={いに\tl{˥}},ipa={ini},pos={名},acc={H}]}%
	{%
		\MeaningsBlock{%
	\ryumeaning%
		{稲}{}{}%
	}%
	}%
	{\tailPrint[]}%
\entry
	{\headPrint[kana={いぬ},ipa={inu},pos={名}]}%
	{%
		\MeaningsBlock{%
	\ryumeaning%
		{犬}{}{}%
	}%
	}%
	{\tailPrint[]}%
\entry
	{\headPrint[kana={いぬ},ipa={inu},pos={名}]}%
	{%
		\MeaningsBlock{%
	\ryumeaning%
		{\ruby[g]{戌}{いぬ}}{十二支の戌(いぬ)}{}%
	}%
	}%
	{\tailPrint[]}%
\entry
	{\headPrint[kana={いのー\tl{˥˩}},ipa={inoː},pos={名},acc={F}]}%
	{%
		\MeaningsBlock{%
	\ryumeaning%
		{礁池}{リーフ内側の浅い海}{}%
	}%
	}%
	{\tailPrint[]}%
\entry
	{\headPrint[kana={いのー\tl{˥˩}},ipa={inoː},pos={名},acc={F}]}%
	{%
		\MeaningsBlock{%
	\ryumeaning%
		{竜巻}{}{}%
	}%
	}%
	{\tailPrint[]}%
\entry
	{\headPrint[kana={いばち\tl{˥˩}},ipa={ibatʃi},pos={名},acc={F}]}%
	{%
		\MeaningsBlock{%
	\ryumeaning%
		{種取祭の山盛飯}{}{}%
	}%
	}%
	{\tailPrint[]}%
\entry
	{\headPrint[kana={いばりすくん\tl{˥˩}},ipa={ibarisu̥kuɴ},pos={動},rui={1},acc={F}]}%
	{%
		\MeaningsBlock{%
	\ryumeaning%
		{威張る。高尚ぶる}{}{}%
	}%
	}%
	{\tailPrint[]}%
\entry
	{\headPrint[kana={いばるん\tl{˥˩}},ipa={ibaruɴ},pos={動},rui={1},para={いばらぬ},acc={F}]}%
	{%
		\MeaningsBlock{%
	\ryumeaning%
		{威張る}{}{}%
	}%
	}%
	{\tailPrint[]}%
\entry
	{\headPrint[kana={いび\tl{˥˩}},ipa={ibi},pos={名},acc={F}]}%
	{%
		\MeaningsBlock{%
	\ryumeaning%
		{\ruby[g]{海老}{えび}}{}{}%
	}%
	}%
	{\tailPrint[]}%
\entry
	{\headPrint[kana={いびるん\tl{˥˩}},ipa={ibiruɴ},pos={動},rui={3},para={いぶぬ},acc={F}]}%
	{%
		\MeaningsBlock{%
	\ryumeaning%
		{植える}{}{}%
	}%
	}%
	{\tailPrint[]}%
\entry
	{\headPrint[kana={いふつぁ\tl{˥}},ipa={ifu̥tsa},pos={名},acc={H}]}%
	{%
		\MeaningsBlock{%
	\ryumeaning%
		{戦争。戦い}{}{}%
	}%
	}%
	{\tailPrint[]}%
\entry
	{\headPrint[kana={いふつぁー\tl{˥}\ws{}すん\tl{˥˩}},ipa={ifu̥tsaː suɴ},pos={句},acc={H F}]}%
	{%
		\MeaningsBlock{%
	\ryumeaning%
		{戦争する}{}{}%
	}%
	}%
	{\tailPrint[]}%
\entry
	{\headPrint[kana={いふなー\tl{˥˩}},ipa={ifu̥naː},pos={名},acc={F}]}%
	{%
		\MeaningsBlock{%
	\ryumeaning%
		{変だ。おかしい}{}{}%
	}%
	}%
	{\tailPrint[]}%
\entry
	{\headPrint[kana={いふなー\ws{}やっさー\tl{˥˩}},ipa={ifu̥naː jassaː},pos={句},acc={F}]}%
	{%
		\MeaningsBlock{%
	\ryumeaning%
		{変だ。変わっている}{悪い意味に使う}{}%
	}%
	}%
	{\tailPrint[]}%
\entry
	{\headPrint[kana={いべー\tl{˥˩}},ipa={ibeː},pos={名},acc={F}]}%
	{%
		\MeaningsBlock{%
	\ryumeaning%
		{位牌}{}{}%
	}%
	}%
	{\tailPrint[]}%
\entry
	{\headPrint[kana={いましみるん\tl{˥}},ipa={imaʃi̥miruɴ},pos={動},rui={3},para={いますむぬ},acc={H}]}%
	{%
		\MeaningsBlock{%
	\ryumeaning%
		{戒める。説教する}{}{}%
	}%
	}%
	{\tailPrint[]}%
\entry
	{\headPrint[kana={いましみん\tl{˥}},ipa={imaʃi̥miɴ},pos={動},rui={3},para={いますむぬ},acc={H}]}%
	{%
		\MeaningsBlock{%
	\ryumeaning%
		{戒める。説教する}{}{}%
	}%
	}%
	{\tailPrint[]}%
\entry
	{\headPrint[kana={いみ\tl{˥}},ipa={imi},pos={名},acc={H}]}%
	{%
		\MeaningsBlock{%
	\ryumeaning%
		{夢}{}{}%
	}%
	}%
	{\tailPrint[]}%
\entry
	{\headPrint[kana={いみ\tl{˥}},ipa={imi},pos={名},acc={H}]}%
	{%
		\MeaningsBlock{%
	\ryumeaning%
		{忌。喪}{}{}%
	}%
	}%
	{\tailPrint[]}%
\entry
	{\headPrint[kana={いみうち},ipa={imi.utʃi},pos={名}]}%
	{%
		\MeaningsBlock{%
	\ryumeaning%
		{忌中}{}{}%
	}%
	}%
	{\tailPrint[]}%
\entry
	{\headPrint[kana={いみはかるん\tl{˥}},ipa={imihḁkaruɴ},pos={動},rui={1},para={いみはからぬ},acc={H}]}%
	{%
		\MeaningsBlock{%
	\ryumeaning%
		{喪に服する。忌中だ}{}{}%
	}%
	}%
	{\tailPrint[]}%
\entry
	{\headPrint[kana={いみるん\tl{˥}},ipa={imiruɴ},pos={動},acc={H}]}%
	{%
		\MeaningsBlock{%
	\ryumeaning%
		{催促する。督促する。ねだる}{}{}%
	}%
	}%
	{\tailPrint[]}%
\entry
	{\headPrint[kana={いや\tl{˥}},ipa={ija},pos={名},acc={H}]}%
	{%
		\MeaningsBlock{%
	\ryumeaning%
		{父。父親}{}{}%
	}%
	}%
	{\tailPrint[]}%
\entry
	{\headPrint[kana={いやぐ\tl{˥}},ipa={ijagu},pos={名},acc={H}]}%
	{%
		\MeaningsBlock{%
	\ryumeaning%
		{櫂}{}{}%
	}%
	}%
	{\tailPrint[]}%
\entry
	{\headPrint[kana={いら\tl{˥}},ipa={ira},pos={名},acc={H}]}%
	{%
		\MeaningsBlock{%
	\ryumeaning%
		{（粟刈りの）小鎌}{}{}%
	}%
	}%
	{\tailPrint[]}%
\entry
	{\headPrint[kana={いら\tl{˥}},ipa={ira},pos={名},acc={H}]}%
	{%
		\MeaningsBlock{%
	\ryumeaning%
		{クラゲ}{}{}%
	}%
	}%
	{\tailPrint[]}%
\entry
	{\headPrint[kana={いらた},ipa={irata},pos={名}]}%
	{%
		\MeaningsBlock{%
	\ryumeaning%
		{西方}{}{}%
	}%
	}%
	{\tailPrint[]}%
\entry
	{\headPrint[kana={いらぶん\tl{˥}},ipa={irabuɴ},pos={動},rui={1},para={いらばぬ},acc={H}]}%
	{%
		\MeaningsBlock{%
	\ryumeaning%
		{選ぶ。選択する}{}{}%
	}%
	}%
	{\tailPrint[]}%
\entry
	{\headPrint[kana={いり\tl{˥}},ipa={iri},pos={名},acc={H}]}%
	{%
		\MeaningsBlock{%
	\ryumeaning%
		{錐}{}{}%
	}%
	}%
	{\tailPrint[]}%
\entry
	{\headPrint[kana={いり\tl{˥˩}},ipa={iri},pos={名},acc={F}]}%
	{%
		\MeaningsBlock{%
	\ryumeaning%
		{西。西方}{}{}%
	}%
	}%
	{\tailPrint[]}%
\entry
	{\headPrint[kana={いりかーるん\tl{˥˩}},ipa={irikaːruɴ},pos={動},rui={1},para={いりかーらぬ},acc={F}]}%
	{%
		\MeaningsBlock{%
	\ryumeaning%
		{入れ代わる}{}{}%
	}%
	}%
	{\tailPrint[]}%
\entry
	{\headPrint[kana={いりかい\tl{˥˩}\ws{}すん\tl{˥˩}},ipa={irikai suɴ},pos={句},acc={F F}]}%
	{%
		\MeaningsBlock{%
	\ryumeaning%
		{入れ換えする。入れ換える}{}{}%
	}%
	}%
	{\tailPrint[]}%
\entry
	{\headPrint[kana={いりかいるん\tl{˥˩}},ipa={irikairuɴ},pos={動},acc={F}]}%
	{%
		\MeaningsBlock{%
	\ryumeaning%
		{入れ換える}{}{}%
	}%
	}%
	{\tailPrint[]}%
\entry
	{\headPrint[kana={いりぎ\tl{˥}},ipa={irigi},pos={名},acc={H}]}%
	{%
		\MeaningsBlock{%
	\ryumeaning%
		{鱗}{}{}%
	}%
	}%
	{\tailPrint[]}%
\entry
	{\headPrint[kana={いりぐん\tl{˥}},ipa={iriguɴ},pos={動},rui={1},para={いりがぬ},acc={H}]}%
	{%
		\MeaningsBlock{%
	\ryumeaning%
		{煎る}{}{}%
	}%
	}%
	{\tailPrint[]}%
\entry
	{\headPrint[kana={いりしな\tl{˥˩}},ipa={iriʃi̥na},pos={名},acc={F}]}%
	{%
		\MeaningsBlock{%
	\ryumeaning%
		{入日}{日光が室内に差し込むこと}{}%
	}%
	}%
	{\tailPrint[]}%
\entry
	{\headPrint[kana={いりびたるん\tl{˥˩}},ipa={iribitaruɴ},pos={動},rui={1},para={いりびたらぬ},acc={F}]}%
	{%
		\MeaningsBlock{%
	\ryumeaning%
		{入り浸る}{}{}%
	}%
	}%
	{\tailPrint[]}%
\entry
	{\headPrint[kana={いり\tl{˥˩}むぐ\tl{˩˥}},ipa={irimugu},pos={名},acc={F+R}]}%
	{%
		\MeaningsBlock{%
	\ryumeaning%
		{入り婿}{}{}%
	}%
	}%
	{\tailPrint[]}%
\entry
	{\headPrint[kana={いりむち\tl{˥˩}},ipa={irimutʃi},pos={名},acc={F}]}%
	{%
		\MeaningsBlock{%
	\ryumeaning%
		{西表。西表島}{八重山諸島の島の名}{}%
	}%
	}%
	{\tailPrint[]}%
\entry
	{\headPrint[kana={いりむら},ipa={irimura},pos={名}]}%
	{%
		\MeaningsBlock{%
	\ryumeaning%
		{西村}{前村、名石村、冨嘉村など西の村々を指す}{}%
	}%
	}%
	{\tailPrint[]}%
\entry
	{\headPrint[kana={いりむん\tl{˥˩}},ipa={irimuɴ},pos={名},acc={F}]}%
	{%
		\MeaningsBlock{%
	\ryumeaning%
		{入れ物。容器}{}{}%
	}%
	}%
	{\tailPrint[]}%
\entry
	{\headPrint[kana={いりん\tl{˥˩}},ipa={iriɴ},pos={動},rui={3},para={いるぬ},acc={F}]}%
	{%
		\MeaningsBlock{%
	\ryumeaning%
		{入れる}{}{}%
	}%
	}%
	{\tailPrint[]}%
\entry
	{\headPrint[kana={いる\tl{˥}},ipa={iru},pos={名},acc={H}]}%
	{%
		\MeaningsBlock{%
	\ryumeaning%
		{色。色彩}{}{}%
	}%
	}%
	{\tailPrint[]}%
\entry
	{\headPrint[kana={いるすそー\tl{˥}はん},ipa={irussoːhaɴ},pos={形},acc={H}]}%
	{%
		\MeaningsBlock{%
	\ryumeaning%
		{色が白い}{色白だ}{}%
	}%
	}%
	{\tailPrint[]}%
\entry
	{\headPrint[kana={いるん\tl{˥}},ipa={iruɴ},pos={動},rui={1},para={いらぬ},acc={H}]}%
	{%
		\MeaningsBlock{%
	\ryumeaning%
		{射る。撃つ}{}{}%
	}%
	}%
	{\tailPrint[]}%
\entry
	{\headPrint[kana={いるん\tl{˥˩}},ipa={iruɴ},pos={動},rui={1},para={いらぬ},acc={F}]}%
	{%
		\MeaningsBlock{%
	\ryumeaning%
		{要る。必要だ}{}{}%
	}%
	}%
	{\tailPrint[]}%
\entry
	{\headPrint[kana={いわり\tl{˥˩}},ipa={iwari},pos={名},acc={F}]}%
	{%
		\MeaningsBlock{%
	\ryumeaning%
		{謂れ。由来}{}{}%
	}%
	}%
	{\tailPrint[]}%
\entry
	{\headPrint[kana={いん\tl{˥}},ipa={iɴ},pos={名},acc={H}]}%
	{%
		\MeaningsBlock{%
	\ryumeaning%
		{印。印鑑}{}{}%
	}%
	}%
	{\tailPrint[]}%
\entry
	{\headPrint[kana={いん\tl{˥}},ipa={in},pos={名},acc={H}]}%
	{%
		\MeaningsBlock{%
	\ryumeaning%
		{犬}{}{}%
	}%
	}%
	{\tailPrint[]}%
\entry
	{\headPrint[kana={いん\tl{˥}},ipa={in},pos={名},acc={H}]}%
	{%
		\MeaningsBlock{%
	\ryumeaning%
		{\ruby[g]{戌}{いぬ}}{十二支の戌(いぬ)}{}%
	}%
	}%
	{\tailPrint[]}%
\entry
	{\headPrint[kana={いん\tl{˥˩}},ipa={iɴ},pos={名},acc={F}]}%
	{%
		\MeaningsBlock{%
	\ryumeaning%
		{洞穴。洞窟}{}{}%
	}%
	}%
	{\tailPrint[]}%
\entry
	{\headPrint[kana={\josho{}いん\kako{}どぅり},ipa={induri},pos={名},acc={H+L}]}%
	{%
		\MeaningsBlock{%
	\ryumeaning%
		{海鳥}{}{}%
	}%
	}%
	{\tailPrint[]}%
\entry
	{\headPrint[kana={いんぬ\tl{˥˩}\ws{}まら\tl{˩˥}},ipa={innu mara},pos={句},acc={F R}]}%
	{%
		\MeaningsBlock{%
	\ryumeaning%
		{鍾乳石}{}{}%
	}%
	}%
	{\tailPrint[]}%
\LetterHead{う}
\entry
	{\headPrint[kana={うー},ipa={uː},pos={接頭}]}%
	{%
		\MeaningsBlock{%
	\ryumeaning%
		{幾～。何～}{助数詞や可算名詞と結合し、数の疑問詞を作る。\emb{うーび}「幾つ」、\emb{うーむし}「何回」など}{}%
	}%
	}%
	{\tailPrint[]}%
\entry
	{\headPrint[kana={うー\tl{˥}},ipa={uː},pos={名},acc={H}]}%
	{%
		\MeaningsBlock{%
	\ryumeaning%
		{\ruby[g]{卯}{う}}{十二支の卯}{}%
	}%
	}%
	{\tailPrint[]}%
\entry
	{\headPrint[kana={うーぐとぅ},ipa={uːgutu},pos={名}]}%
	{%
		\MeaningsBlock{%
	\ryumeaning%
		{大事。一大事}{}{}%
	}%
	}%
	{\tailPrint[]}%
\entry
	{\headPrint[kana={うーちぃ\tl{˥˩}},ipa={uːtʃi},pos={名},acc={F}]}%
	{%
		\MeaningsBlock{%
	\ryumeaning%
		{幾つ}{}{}%
	\ryumeaning%
		{何歳}{}{}%
	}%
	}%
	{\tailPrint[]}%
\entry
	{\headPrint[kana={うーとーとぅ},ipa={uːtoːtu},pos={感}]}%
	{%
		\MeaningsBlock{%
	\ryumeaning%
		{ああ尊し}{祈りの冒頭のことば}{}%
	}%
	}%
	{\tailPrint[]}%
\entry
	{\headPrint[kana={うーび\tl{˥˩}},ipa={uːbi},pos={名},acc={F}]}%
	{%
		\MeaningsBlock{%
	\ryumeaning%
		{幾つ。幾ら}{}{}%
	}%
	}%
	{\tailPrint[]}%
\entry
	{\headPrint[kana={うーむし},ipa={uːmuʃi},pos={名}]}%
	{%
		\MeaningsBlock{%
	\ryumeaning%
		{何回}{}{}%
	}%
	}%
	{\tailPrint[]}%
\entry
	{\headPrint[kana={うい\tl{˥}},ipa={ui},pos={名},acc={H}]}%
	{%
		\MeaningsBlock{%
	\ryumeaning%
		{銛}{}{}%
	}%
	}%
	{\tailPrint[]}%
\entry
	{\headPrint[kana={うい\tl{˥}},ipa={ui},pos={動},para={うわぬ},acc={H}]}%
	{%
		\MeaningsBlock{%
	\ryumeaning%
		{泳ぐ}{\emb{ういっしゃん}「泳げる」}{}%
	}%
	}%
	{\tailPrint[]}%
\entry
	{\headPrint[kana={うい\tl{˥}},ipa={ui},pos={名},acc={H}]}%
	{%
		\MeaningsBlock{%
	\ryumeaning%
		{老い。老いること}{}{}%
	}%
	}%
	{\tailPrint[]}%
\entry
	{\headPrint[kana={うい\tl{˥˩}},ipa={ui},pos={名},acc={F}]}%
	{%
		\MeaningsBlock{%
	\ryumeaning%
		{上。上方。目上}{}{}%
	}%
	}%
	{\tailPrint[]}%
\entry
	{\headPrint[kana={ういが\tl{˥˩}},ipa={uiga},pos={句},acc={F}]}%
	{%
		\MeaningsBlock{%
	\ryumeaning%
		{上へ。上方へ}{}{}%
	}%
	}%
	{\tailPrint[]}%
\entry
	{\headPrint[kana={ういくむん\tl{˥}},ipa={uiku̥muɴ},pos={動},rui={1},para={ういくまぬ},acc={H}]}%
	{%
		\MeaningsBlock{%
	\ryumeaning%
		{追い込む}{}{}%
	}%
	}%
	{\tailPrint[]}%
\entry
	{\headPrint[kana={ういし\tl{˥}},ipa={uiʃi},pos={名},acc={H}]}%
	{%
		\MeaningsBlock{%
	\ryumeaning%
		{訓戒}{「御意志」の意味で国王の命令か}{}%
	}%
	}%
	{\tailPrint[]}%
\entry
	{\headPrint[kana={ういし\ws{}うがますん},ipa={uiʃi ugamasuɴ},pos={句},rui={2},para={ういしうがまはぬ}]}%
	{%
		\MeaningsBlock{%
	\ryumeaning%
		{訓戒する}{「御意志を拝ませる」の意味で国王の命令を伝えたことからか}{}%
	}%
	}%
	{\tailPrint[]}%
\entry
	{\headPrint[kana={ういすくん\tl{˥˩}},ipa={uisu̥kuɴ},pos={動},rui={1},para={ういすかぬ},acc={F}]}%
	{%
		\MeaningsBlock{%
	\ryumeaning%
		{追いつく}{}{}%
	}%
	}%
	{\tailPrint[]}%
\entry
	{\headPrint[kana={ういなん\tl{˥}},ipa={uinaɴ},pos={名},acc={H}]}%
	{%
		\MeaningsBlock{%
	\ryumeaning%
		{牝牛}{}{}%
	}%
	}%
	{\tailPrint[]}%
\entry
	{\headPrint[kana={ういぬぐん\tl{˥˩}},ipa={uinuguɴ},pos={動},rui={1},para={ういぬがぬ},acc={F}]}%
	{%
		\MeaningsBlock{%
	\ryumeaning%
		{追い抜く。追い越す}{}{}%
	}%
	}%
	{\tailPrint[]}%
\entry
	{\headPrint[kana={ういぬ\tl{˥˩}\ws{}ぴぃとぅ\tl{˥˩}},ipa={uinu pi̥tu},pos={句},acc={F F}]}%
	{%
		\MeaningsBlock{%
	\ryumeaning%
		{上役。上司}{}{}%
	}%
	}%
	{\tailPrint[]}%
\entry
	{\headPrint[kana={ういぱー\tl{˥}},ipa={uipaː},pos={名},acc={H}]}%
	{%
		\MeaningsBlock{%
	\ryumeaning%
		{曽祖母}{}{}%
	}%
	}%
	{\tailPrint[]}%
\entry
	{\headPrint[kana={ういぱろーん\tl{˥˩}},ipa={uiparoːɴ},pos={動},rui={2},acc={F}]}%
	{%
		\MeaningsBlock{%
	\ryumeaning%
		{追っ払う}{}{}%
	}%
	}%
	{\tailPrint[]}%
\entry
	{\headPrint[kana={ういぶや\tl{˥}},ipa={uibuja},pos={名},acc={H}]}%
	{%
		\MeaningsBlock{%
	\ryumeaning%
		{曽祖父}{}{}%
	}%
	}%
	{\tailPrint[]}%
\entry
	{\headPrint[kana={ういぷら\tl{˥}},ipa={uipura},pos={名},acc={H}]}%
	{%
		\MeaningsBlock{%
	\ryumeaning%
		{\ruby[g]{耄}{もう}\ruby[g]{碌}{ろく}。痴呆}{}{}%
	}%
	}%
	{\tailPrint[]}%
\entry
	{\headPrint[kana={ういぶり\tl{˥}},ipa={uiburi},pos={動},rui={1},para={ういぶらぬ},acc={H}]}%
	{%
		\MeaningsBlock{%
	\ryumeaning%
		{老いぼれる}{}{}%
	}%
	}%
	{\tailPrint[]}%
\entry
	{\headPrint[kana={ういぶん\tl{˥˩}},ipa={uibuɴ},pos={動},acc={F}]}%
	{%
		\MeaningsBlock{%
	\ryumeaning%
		{及ぶ。叶う}{}{}%
	}%
	}%
	{\tailPrint[]}%
\entry
	{\headPrint[kana={ういまーすん\tl{˥˩}},ipa={uimaːsuɴ},pos={動},rui={2b},para={ういまーさぬ},acc={F}]}%
	{%
		\MeaningsBlock{%
	\ryumeaning%
		{追い回す。追い立てる}{}{}%
	}%
	}%
	{\tailPrint[]}%
\entry
	{\headPrint[kana={ういるん\tl{˥}},ipa={uiruɴ},pos={動},rui={3},para={ういぬ},acc={H}]}%
	{%
		\MeaningsBlock{%
	\ryumeaning%
		{老いる。年を取る}{}{}%
	}%
	}%
	{\tailPrint[]}%
\entry
	{\headPrint[kana={ういわっきるん\tl{˥˩}},ipa={uiwakkiruɴ},pos={動},rui={3},para={ういわっくぬ},acc={F}]}%
	{%
		\MeaningsBlock{%
	\ryumeaning%
		{追い散らす}{}{}%
	}%
	}%
	{\tailPrint[]}%
\entry
	{\headPrint[kana={ういんぐん\tl{˥˩}},ipa={uiŋguɴ},pos={動},rui={1},para={ういんがぬ},acc={F}]}%
	{%
		\MeaningsBlock{%
	\ryumeaning%
		{追っかけていく}{}{}%
	}%
	}%
	{\tailPrint[]}%
\entry
	{\headPrint[kana={ういんだすん\tl{˥}},ipa={uindasuɴ},pos={動},rui={2},para={ういんだはぬ},acc={H}]}%
	{%
		\MeaningsBlock{%
	\ryumeaning%
		{追い出す}{}{}%
	}%
	}%
	{\tailPrint[]}%
\entry
	{\headPrint[kana={うか\tl{˥}},ipa={uka},pos={名},acc={H}]}%
	{%
		\MeaningsBlock{%
	\ryumeaning%
		{負債。負い目}{}{}%
	}%
	}%
	{\tailPrint[]}%
\entry
	{\headPrint[kana={うかさ\tl{˥}はーん},ipa={ukasahaːɴ},pos={形},acc={H}]}%
	{%
		\MeaningsBlock{%
	\ryumeaning%
		{醜い。見苦しい}{}{}%
	}%
	}%
	{\tailPrint[]}%
\entry
	{\headPrint[kana={うがじぃ\tl{˥}},ipa={ugadʒi},pos={名},acc={H}]}%
	{%
		\MeaningsBlock{%
	\ryumeaning%
		{神占い}{米粒で占う}{}%
	}%
	}%
	{\tailPrint[]}%
\entry
	{\headPrint[kana={うがすん\tl{˥}},ipa={ugasuɴ},pos={動},rui={2},para={うがはぬ},acc={H}]}%
	{%
		\MeaningsBlock{%
	\ryumeaning%
		{起こす}{}{}%
	}%
	}%
	{\tailPrint[]}%
\entry
	{\headPrint[kana={うがすん\tl{˥}},ipa={ugasuɴ},pos={動},rui={2},para={うがはぬ},acc={H}]}%
	{%
		\MeaningsBlock{%
	\ryumeaning%
		{動かす}{}{}%
	\ryumeaning%
		{移動する}{}{}%
	}%
	}%
	{\tailPrint[]}%
\entry
	{\headPrint[kana={うかっと},ipa={ukatto},pos={副}]}%
	{%
		\MeaningsBlock{%
	\ryumeaning%
		{うっかり}{}{}%
	}%
	}%
	{\tailPrint[]}%
\entry
	{\headPrint[kana={うかっと\ws{}すん},ipa={ukatto suɴ},pos={句}]}%
	{%
		\MeaningsBlock{%
	\ryumeaning%
		{軽率だ}{}{}%
	}%
	}%
	{\tailPrint[]}%
\entry
	{\headPrint[kana={うかっと\ws{}すん\tl{˥˩}},ipa={ukatto suɴ},pos={句},acc={HL F}]}%
	{%
		\MeaningsBlock{%
	\ryumeaning%
		{ぼんやりする。うっかりする}{}{}%
	}%
	}%
	{\tailPrint[]}%
\entry
	{\headPrint[kana={うがみ\tl{˥}},ipa={ugami},pos={名},acc={H}]}%
	{%
		\MeaningsBlock{%
	\ryumeaning%
		{拝み。拝むこと}{}{}%
	\ryumeaning%
		{祈願}{}{}%
	}%
	}%
	{\tailPrint[]}%
\entry
	{\headPrint[kana={うがむん\tl{˥}},ipa={ugamuɴ},pos={動},rui={1},para={うがまぬ},acc={H}]}%
	{%
		\MeaningsBlock{%
	\ryumeaning%
		{拝む}{}{}%
	}%
	}%
	{\tailPrint[]}%
\entry
	{\headPrint[kana={うがり\tl{˥˩}},ipa={ugari},pos={名},acc={F}]}%
	{%
		\MeaningsBlock{%
	\ryumeaning%
		{高地。小高い台地}{}{}%
	}%
	}%
	{\tailPrint[]}%
\entry
	{\headPrint[kana={うがん\tl{˥}},ipa={ugaɴ},pos={名},acc={H}]}%
	{%
		\MeaningsBlock{%
	\ryumeaning%
		{祈願}{}{}%
	}%
	}%
	{\tailPrint[]}%
\entry
	{\headPrint[kana={うがんじゅ},ipa={ugandʒu},pos={名}]}%
	{%
		\MeaningsBlock{%
	\ryumeaning%
		{拝所}{}{}%
	}%
	}%
	{\tailPrint[]}%
\entry
	{\headPrint[kana={うがんちゅむん\tl{˥}},ipa={ugantʃumuɴ},pos={動},rui={1},para={うがんちゅまぬ},acc={H}]}%
	{%
		\MeaningsBlock{%
	\ryumeaning%
		{くたばる。参る}{}{}%
	}%
	}%
	{\tailPrint[]}%
\entry
	{\headPrint[kana={うき\tl{˥}},ipa={uki},pos={名},acc={H}]}%
	{%
		\MeaningsBlock{%
	\ryumeaning%
		{木桶。水桶}{}{}%
	}%
	}%
	{\tailPrint[]}%
\entry
	{\headPrint[kana={うきしゃ\tl{˥}はん},ipa={uki̥ʃahaɴ},pos={形},para={うきしゃへぬ},acc={H}]}%
	{%
		\MeaningsBlock{%
	\ryumeaning%
		{恰好悪い。不美人}{}{}%
	}%
	}%
	{\tailPrint[]}%
\entry
	{\headPrint[kana={うきとぅるん\tl{˥}},ipa={ukitu̥ruɴ},pos={動},rui={1},para={うきとらぬ},acc={H}]}%
	{%
		\MeaningsBlock{%
	\ryumeaning%
		{受け取る}{}{}%
	\ryumeaning%
		{引き受ける}{}{}%
	}%
	}%
	{\tailPrint[]}%
\entry
	{\headPrint[kana={うきむつん\tl{˥}},ipa={ukimutsuɴ},pos={動},para={うきむつぁぬ},acc={H}]}%
	{%
		\MeaningsBlock{%
	\ryumeaning%
		{受け持つ}{}{}%
	}%
	}%
	{\tailPrint[]}%
\entry
	{\headPrint[kana={うきるん\tl{˥}},ipa={ukiruɴ},pos={動},para={うきらぬ},acc={H}]}%
	{%
		\MeaningsBlock{%
	\ryumeaning%
		{受ける。請け負う}{}{}%
	}%
	}%
	{\tailPrint[]}%
\entry
	{\headPrint[kana={うぎるん},ipa={ugiruɴ},pos={動}]}%
	{%
		\MeaningsBlock{%
	\ryumeaning%
		{起きる。目覚める}{}{}%
	}%
	}%
	{\tailPrint[]}%
\entry
	{\headPrint[kana={うぎん\tl{˥}},ipa={ugiɴ},pos={動},rui={3},para={うぐぬ},acc={H}]}%
	{%
		\MeaningsBlock{%
	\ryumeaning%
		{起きる。目覚める}{}{}%
	}%
	}%
	{\tailPrint[]}%
\entry
	{\headPrint[kana={うく\tl{˥˩}},ipa={uku},pos={名},acc={F}]}%
	{%
		\MeaningsBlock{%
	\ryumeaning%
		{奥。奥の方}{}{}%
	}%
	}%
	{\tailPrint[]}%
\entry
	{\headPrint[kana={うくたるん\tl{˥˩}},ipa={ukutaruɴ},pos={動},rui={1},para={うくたらぬ},acc={F}]}%
	{%
		\MeaningsBlock{%
	\ryumeaning%
		{怠る。油断する}{}{}%
	}%
	}%
	{\tailPrint[]}%
\entry
	{\headPrint[kana={うくびょー\tl{˥}さーん},ipa={ukubjoːsaːɴ},pos={形},acc={H}]}%
	{%
		\MeaningsBlock{%
	\ryumeaning%
		{臆病だ}{}{}%
	}%
	}%
	{\tailPrint[]}%
\entry
	{\headPrint[kana={うくらすん\tl{˥˩}},ipa={ukurasuɴ},pos={動},rui={2},para={うくらはぬ},acc={F}]}%
	{%
		\MeaningsBlock{%
	\ryumeaning%
		{遅らせる}{}{}%
	}%
	}%
	{\tailPrint[]}%
\entry
	{\headPrint[kana={うくりん\tl{˥˩}},ipa={ukuriɴ},pos={動},rui={3},acc={F}]}%
	{%
		\MeaningsBlock{%
	\ryumeaning%
		{遅れる}{}{}%
	}%
	}%
	{\tailPrint[]}%
\entry
	{\headPrint[kana={うぐるぴん\tl{˥˩}},ipa={ugurupiɴ},pos={名},acc={F}]}%
	{%
		\MeaningsBlock{%
	\ryumeaning%
		{お盆の最終日}{}{}%
	}%
	}%
	{\tailPrint[]}%
\entry
	{\headPrint[kana={うくるん\tl{˥}},ipa={ukurun},pos={動},rui={1},para={うくらぬ},acc={H}]}%
	{%
		\MeaningsBlock{%
	\ryumeaning%
		{怒る。叱る}{}{}%
	}%
	}%
	{\tailPrint[]}%
\entry
	{\headPrint[kana={うくるん\tl{˥}},ipa={ukurun},pos={動},rui={1},para={うくらぬ},acc={H}]}%
	{%
		\MeaningsBlock{%
	\ryumeaning%
		{起こる}{}{}%
	}%
	}%
	{\tailPrint[]}%
\entry
	{\headPrint[kana={うぐるん\tl{˥˩}},ipa={uguruɴ},pos={動},rui={1},para={うぐらぬ},acc={F}]}%
	{%
		\MeaningsBlock{%
	\ryumeaning%
		{送る。葬式をする}{}{}%
	}%
	}%
	{\tailPrint[]}%
\entry
	{\headPrint[kana={うぐん\tl{˥}},ipa={uguɴ},pos={動},rui={1},para={うがぬ},acc={H}]}%
	{%
		\MeaningsBlock{%
	\ryumeaning%
		{動く}{}{}%
	}%
	}%
	{\tailPrint[]}%
\entry
	{\headPrint[kana={うさぎ\tl{˥˩}むぬ\tl{˩˥}},ipa={usagimunu},pos={名},acc={F+R}]}%
	{%
		\MeaningsBlock{%
	\ryumeaning%
		{献上品}{}{}%
	}%
	}%
	{\tailPrint[]}%
\entry
	{\headPrint[kana={うさぎるん\tl{˥˩}},ipa={usagiruɴ},pos={動},para={うさぎぬ},acc={F}]}%
	{%
		\MeaningsBlock{%
	\ryumeaning%
		{差し上げる}{敬語}{}%
	}%
	}%
	{\tailPrint[]}%
\entry
	{\headPrint[kana={うさまるん\tl{˥}},ipa={usamaruɴ},pos={動},rui={1},para={うさまらぬ},acc={H}]}%
	{%
		\MeaningsBlock{%
	\ryumeaning%
		{治まる}{}{}%
	}%
	}%
	{\tailPrint[]}%
\entry
	{\headPrint[kana={うさみるん\tl{˥}},ipa={usamiruɴ},pos={動},rui={3},para={うさむぬ},acc={H}]}%
	{%
		\MeaningsBlock{%
	\ryumeaning%
		{納める。納入する}{}{}%
	\ryumeaning%
		{治める。統治する}{}{}%
	}%
	}%
	{\tailPrint[]}%
\entry
	{\headPrint[kana={うし\tl{˥}},ipa={uʃi},pos={名},acc={H}]}%
	{%
		\MeaningsBlock{%
	\ryumeaning%
		{臼}{搗き臼と引き臼がある}{}%
	}%
	}%
	{\tailPrint[]}%
\entry
	{\headPrint[kana={うし\tl{˥˩}},ipa={uʃi},pos={名},acc={F}]}%
	{%
		\MeaningsBlock{%
	\ryumeaning%
		{牛}{「雄牛」は\emb{ぐちぇー}、「牝牛」は\emb{ういなん}}{}%
	}%
	}%
	{\tailPrint[]}%
\entry
	{\headPrint[kana={うし\tl{˥˩}},ipa={uʃi},pos={名},acc={F}]}%
	{%
		\MeaningsBlock{%
	\ryumeaning%
		{\ruby[g]{丑}{うし}}{十二支の丑}{}%
	}%
	}%
	{\tailPrint[]}%
\entry
	{\headPrint[kana={うじ\tl{˥}},ipa={udʒi},pos={名},acc={H}]}%
	{%
		\MeaningsBlock{%
	\ryumeaning%
		{腕}{}{}%
	}%
	}%
	{\tailPrint[]}%
\entry
	{\headPrint[kana={うじ\tl{˥}},ipa={udʒi},pos={名},acc={H}]}%
	{%
		\MeaningsBlock{%
	\ryumeaning%
		{蛆}{}{}%
	}%
	}%
	{\tailPrint[]}%
\entry
	{\headPrint[kana={うしあんぎるん\tl{˥˩}},ipa={uʃi.aŋgiruɴ},pos={動},rui={3},para={うしあんぐぬ},acc={F}]}%
	{%
		\MeaningsBlock{%
	\ryumeaning%
		{押し上げる}{}{}%
	}%
	}%
	{\tailPrint[]}%
\entry
	{\headPrint[kana={うしかいすん\tl{˥˩}},ipa={uʃi̥kaisuɴ},pos={動},rui={2b},para={うしかいさぬ},acc={F}]}%
	{%
		\MeaningsBlock{%
	\ryumeaning%
		{押し返す}{}{}%
	}%
	}%
	{\tailPrint[]}%
\entry
	{\headPrint[kana={うしくみるん\tl{˥˩}},ipa={uʃi̥kumiruɴ},pos={動},rui={3},para={うしくむぬ},acc={F}]}%
	{%
		\MeaningsBlock{%
	\ryumeaning%
		{押し込める}{}{}%
	}%
	}%
	{\tailPrint[]}%
\entry
	{\headPrint[kana={うしくむん\tl{˥˩}},ipa={uʃi̥kumuɴ},pos={動},rui={1},para={うしくまぬ},acc={F}]}%
	{%
		\MeaningsBlock{%
	\ryumeaning%
		{押し込む。詰め込む。へし込む}{力を入れて強く押し込む}{}%
	}%
	}%
	{\tailPrint[]}%
\entry
	{\headPrint[kana={うしくるばすん\tl{˥˩}},ipa={uʃi̥ku̥rubasuɴ},pos={動},rui={2},para={うしくるばはぬ},acc={F}]}%
	{%
		\MeaningsBlock{%
	\ryumeaning%
		{押し転ばす。押し転がす}{}{}%
	}%
	}%
	{\tailPrint[]}%
\entry
	{\headPrint[kana={うししきるん\tl{˥˩}},ipa={uʃiʃi̥kiruɴ},pos={動},rui={3},para={うしすくぬ},acc={F}]}%
	{%
		\MeaningsBlock{%
	\ryumeaning%
		{押し付ける。強いる}{}{}%
	}%
	}%
	{\tailPrint[]}%
\entry
	{\headPrint[kana={うししきん\tl{˥˩}},ipa={uʃiʃi̥kiɴ},pos={動},rui={3},para={うしぃすくぬ},acc={F}]}%
	{%
		\MeaningsBlock{%
	\ryumeaning%
		{押さえ込む}{}{}%
	}%
	}%
	{\tailPrint[]}%
\entry
	{\headPrint[kana={うしたま\tl{˥}},ipa={uʃi̥tama},pos={名},acc={H}]}%
	{%
		\MeaningsBlock{%
	\ryumeaning%
		{彼ら。あいつら}{}{}%
	}%
	}%
	{\tailPrint[]}%
\entry
	{\headPrint[kana={うしとーすん},ipa={uʃi̥toːsuɴ},pos={動},rui={2},para={うしとーはぬ}]}%
	{%
		\MeaningsBlock{%
	\ryumeaning%
		{突き飛ばす}{}{}%
	\ryumeaning%
		{押し倒す}{}{}%
	}%
	}%
	{\tailPrint[]}%
\entry
	{\headPrint[kana={うしとぅ\tl{˥}},ipa={uʃi̥tu},pos={名},acc={H}]}%
	{%
		\MeaningsBlock{%
	\ryumeaning%
		{年寄り。老人}{}{}%
	}%
	}%
	{\tailPrint[]}%
\entry
	{\headPrint[kana={うしとぅ\tl{˥}\ws{}なるん\tl{˩˥}},ipa={uʃi̥tu naruɴ},pos={句},rui={1},para={うしとならぬ},acc={H+R}]}%
	{%
		\MeaningsBlock{%
	\ryumeaning%
		{老人になる。歳を取る}{}{}%
	}%
	}%
	{\tailPrint[]}%
\entry
	{\headPrint[kana={うしとぅぱー\tl{˥}},ipa={uʃi̥tupaː},pos={名},acc={H}]}%
	{%
		\MeaningsBlock{%
	\ryumeaning%
		{老婆。おばあさん}{}{}%
	}%
	}%
	{\tailPrint[]}%
\entry
	{\headPrint[kana={うしとぅぶや},ipa={uʃi̥tubuja},pos={名}]}%
	{%
		\MeaningsBlock{%
	\ryumeaning%
		{老爺。おじいさん}{}{}%
	}%
	}%
	{\tailPrint[]}%
\entry
	{\headPrint[kana={うしぴ\tl{˥}},ipa={uʃi̥pi},pos={名},acc={H}]}%
	{%
		\MeaningsBlock{%
	\ryumeaning%
		{風呂敷}{}{}%
	}%
	}%
	{\tailPrint[]}%
\entry
	{\headPrint[kana={うしむどぅすん\tl{˥˩}},ipa={uʃimudusuɴ},pos={動},rui={2b},para={うしむどぅさぬ},acc={F}]}%
	{%
		\MeaningsBlock{%
	\ryumeaning%
		{押し戻す}{}{}%
	}%
	}%
	{\tailPrint[]}%
\entry
	{\headPrint[kana={うじゃうじゃ},ipa={udʒa.udʒa},pos={擬}]}%
	{%
		\MeaningsBlock{%
	\ryumeaning%
		{うようよ}{虫などのたくさんいるさま}{}%
	}%
	}%
	{\tailPrint[]}%
\entry
	{\headPrint[kana={うしゅがなし\tl{˥˩}},ipa={uʃuganaʃi},pos={名},acc={F}]}%
	{%
		\MeaningsBlock{%
	\ryumeaning%
		{琉球国王（御主加那志）。国王}{}{}%
	}%
	}%
	{\tailPrint[]}%
\entry
	{\headPrint[kana={うしゅく},ipa={uʃu̥ku},pos={副}]}%
	{%
		\MeaningsBlock{%
	\ryumeaning%
		{それほど}{}{}%
	}%
	}%
	{\tailPrint[]}%
\entry
	{\headPrint[kana={うしゆしるん},ipa={uʃijuʃiruɴ},pos={動}]}%
	{%
		\MeaningsBlock{%
	\ryumeaning%
		{押し寄せる}{}{}%
	}%
	}%
	{\tailPrint[]}%
\entry
	{\headPrint[kana={うず\tl{˥˩}},ipa={udzu},pos={名},acc={F}]}%
	{%
		\MeaningsBlock{%
	\ryumeaning%
		{布団}{}{}%
	}%
	}%
	{\tailPrint[]}%
\entry
	{\headPrint[kana={うすくみん\tl{˥˩}},ipa={usu̥kumiɴ},pos={動},rui={3},para={うすくむぬ},acc={F}]}%
	{%
		\MeaningsBlock{%
	\ryumeaning%
		{仕舞う。収納する}{}{}%
	}%
	}%
	{\tailPrint[]}%
\entry
	{\headPrint[kana={うすくん\tl{˩˥}},ipa={usu̥kuɴ},pos={動},rui={1},para={うすかぬ},acc={R}]}%
	{%
		\MeaningsBlock{%
	\ryumeaning%
		{置く}{}{}%
	}%
	}%
	{\tailPrint[]}%
\entry
	{\headPrint[kana={うすな\tl{˥}},ipa={usu̥na},pos={名},acc={H}]}%
	{%
		\MeaningsBlock{%
	\ryumeaning%
		{沖縄。沖縄本島}{地名}{}%
	}%
	}%
	{\tailPrint[]}%
\entry
	{\headPrint[kana={うすなすん\tl{˥˩}},ipa={usu̥nasuɴ},pos={動},rui={2},para={うすなはぬ},acc={F}]}%
	{%
		\MeaningsBlock{%
	\ryumeaning%
		{なくす。失う}{}{}%
	}%
	}%
	{\tailPrint[]}%
\entry
	{\headPrint[kana={うすぴらがすん\tl{˥˩}},ipa={usu̥piragasuɴ},pos={動},rui={2},acc={F}]}%
	{%
		\MeaningsBlock{%
	\ryumeaning%
		{踏み潰す}{}{}%
	}%
	}%
	{\tailPrint[]}%
\entry
	{\headPrint[kana={うすふかすん\tl{˥˩}},ipa={usufu̥kasuɴ},pos={動},rui={2},para={うすふかはぬ},acc={F}]}%
	{%
		\MeaningsBlock{%
	\ryumeaning%
		{伏せる。下に向けて置く}{}{}%
	}%
	}%
	{\tailPrint[]}%
\entry
	{\headPrint[kana={うすふくん\tl{˥˩}},ipa={usufu̥kuɴ},pos={動},rui={1},para={うすふかぬ},acc={F}]}%
	{%
		\MeaningsBlock{%
	\ryumeaning%
		{俯く}{}{}%
	}%
	}%
	{\tailPrint[]}%
\entry
	{\headPrint[kana={うすまさん},ipa={usumasaɴ},pos={形},acc={-}]}%
	{%
		\MeaningsBlock{%
	\ryumeaning%
		{ものすごい。恐ろしい}{}{}%
	}%
	}%
	{\tailPrint[]}%
\entry
	{\headPrint[kana={うずまるん\tl{˥˩}},ipa={udzumaruɴ},pos={動},rui={1},para={うずまらぬ},acc={F}]}%
	{%
		\MeaningsBlock{%
	\ryumeaning%
		{埋もれる。埋まる}{}{}%
	}%
	}%
	{\tailPrint[]}%
\entry
	{\headPrint[kana={うずみるん\tl{˥˩}},ipa={udzumiruɴ},pos={動},rui={3},acc={F}]}%
	{%
		\MeaningsBlock{%
	\ryumeaning%
		{埋める}{}{}%
	}%
	}%
	{\tailPrint[]}%
\entry
	{\headPrint[kana={うずら\tl{˥}},ipa={udzura},pos={名},acc={H}]}%
	{%
		\MeaningsBlock{%
	\ryumeaning%
		{ウズラ}{小鳥の一種}{}%
	}%
	}%
	{\tailPrint[]}%
\entry
	{\headPrint[kana={うすん\tl{˥˩}},ipa={usuɴ},pos={動},rui={2},para={うさぬ},acc={F}]}%
	{%
		\MeaningsBlock{%
	\ryumeaning%
		{押す。押さえる}{}{}%
	\ryumeaning%
		{覆う}{}{}%
	}%
	}%
	{\tailPrint[]}%
\entry
	{\headPrint[kana={うせー\tl{˥}},ipa={useː},pos={名},acc={H}]}%
	{%
		\MeaningsBlock{%
	\ryumeaning%
		{\ruby[g]{肴}{さかな}}{}{}%
	}%
	}%
	{\tailPrint[]}%
\entry
	{\headPrint[kana={うせーるん\tl{˥˩}},ipa={useːruɴ},pos={動},acc={F}]}%
	{%
		\MeaningsBlock{%
	\ryumeaning%
		{軽蔑する}{}{}%
	}%
	}%
	{\tailPrint[]}%
\entry
	{\headPrint[kana={うせーん\tl{˥˩}},ipa={useːɴ},pos={動},rui={3?},para={うせーるぬ},acc={F}]}%
	{%
		\MeaningsBlock{%
	\ryumeaning%
		{貶す。馬鹿にする}{}{}%
	}%
	}%
	{\tailPrint[]}%
\entry
	{\headPrint[kana={うた\tl{˥˩}},ipa={uta},pos={名},acc={F}]}%
	{%
		\MeaningsBlock{%
	\ryumeaning%
		{歌。民謡}{}{}%
	}%
	}%
	{\tailPrint[]}%
\entry
	{\headPrint[kana={うたーり\tl{˥˩}},ipa={utaːri},pos={名},acc={F}]}%
	{%
		\MeaningsBlock{%
	\ryumeaning%
		{いくたり。何人}{}{}%
	}%
	}%
	{\tailPrint[]}%
\entry
	{\headPrint[kana={うだぎ},ipa={udagi},pos={副}]}%
	{%
		\MeaningsBlock{%
	\ryumeaning%
		{その程度。それだけ}{}{}%
	}%
	}%
	{\tailPrint[]}%
\entry
	{\headPrint[kana={うたげーるん\tl{˥˩}},ipa={utageːruɴ},pos={動},para={うたげーぬ},acc={F}]}%
	{%
		\MeaningsBlock{%
	\ryumeaning%
		{疑う}{疑問に思う}{}%
	}%
	}%
	{\tailPrint[]}%
\entry
	{\headPrint[kana={うたごーん\tl{˥˩}},ipa={utagoːɴ},pos={動},para={うたがーぬ},acc={F}]}%
	{%
		\MeaningsBlock{%
	\ryumeaning%
		{疑う}{}{}%
	}%
	}%
	{\tailPrint[]}%
\entry
	{\headPrint[kana={うたさんしん\tl{˥˩}},ipa={utasanʃiɴ},pos={名},acc={F}]}%
	{%
		\MeaningsBlock{%
	\ryumeaning%
		{音曲。歌三線}{}{}%
	}%
	}%
	{\tailPrint[]}%
\entry
	{\headPrint[kana={うたすん\tl{˥}},ipa={utasuɴ},pos={動},rui={2},para={うたはぬ},acc={H}]}%
	{%
		\MeaningsBlock{%
	\ryumeaning%
		{落とす}{}{}%
	}%
	}%
	{\tailPrint[]}%
\entry
	{\headPrint[kana={うだち\tl{˥}},ipa={udatʃi},pos={名},acc={H}]}%
	{%
		\MeaningsBlock{%
	\ryumeaning%
		{海への降り口}{}{}%
	}%
	}%
	{\tailPrint[]}%
\entry
	{\headPrint[kana={うたま\tl{˥}},ipa={utama},pos={名},acc={H}]}%
	{%
		\MeaningsBlock{%
	\ryumeaning%
		{子。子供}{「子供たち」は\emb{うたまんじ}}{}%
	}%
	}%
	{\tailPrint[]}%
\entry
	{\headPrint[kana={うたま\tl{˥}\ws{}すかなすん\tl{˥}},ipa={utama su̥kanasuɴ},pos={句},rui={2},acc={H H}]}%
	{%
		\MeaningsBlock{%
	\ryumeaning%
		{子育てする}{}{}%
	}%
	}%
	{\tailPrint[]}%
\entry
	{\headPrint[kana={うたま\tl{˥}\ws{}なすん\tl{˩˥}},ipa={utama nasuɴ},pos={句},rui={2},acc={H R}]}%
	{%
		\MeaningsBlock{%
	\ryumeaning%
		{子供を産む。出産する}{}{}%
	}%
	}%
	{\tailPrint[]}%
\entry
	{\headPrint[kana={\josho{}うたま\kako{}んじ},ipa={utamandʒi},pos={名},acc={H+L}]}%
	{%
		\MeaningsBlock{%
	\ryumeaning%
		{子たち。子供たち}{}{}%
	}%
	}%
	{\tailPrint[]}%
\entry
	{\headPrint[kana={うちあみ\tl{˥}},ipa={utʃi.ami},pos={名},acc={H}]}%
	{%
		\MeaningsBlock{%
	\ryumeaning%
		{（室内に）打ち込む雨}{}{}%
	}%
	}%
	{\tailPrint[]}%
\entry
	{\headPrint[kana={うちあん\tl{˥}},ipa={utʃi.aɴ},pos={名},acc={H}]}%
	{%
		\MeaningsBlock{%
	\ryumeaning%
		{投げ網。投網}{}{}%
	}%
	}%
	{\tailPrint[]}%
\entry
	{\headPrint[kana={うちくむん\tl{˥}},ipa={utʃi̥ku̥muɴ},pos={動},rui={1},para={うちくまぬ},acc={H}]}%
	{%
		\MeaningsBlock{%
	\ryumeaning%
		{投げ込む。放り込む}{}{}%
	}%
	}%
	{\tailPrint[]}%
\entry
	{\headPrint[kana={うちくる\tl{˥˩}},ipa={utʃi̥kuru},pos={名},acc={F}]}%
	{%
		\MeaningsBlock{%
	\ryumeaning%
		{押入れ}{}{}%
	}%
	}%
	{\tailPrint[]}%
\entry
	{\headPrint[kana={うちざ\tl{˥}},ipa={utʃidza},pos={名},acc={H}]}%
	{%
		\MeaningsBlock{%
	\ryumeaning%
		{兄弟姉妹}{「姉妹」は\emb{ぶなり}という}{}%
	}%
	}%
	{\tailPrint[]}%
\entry
	{\headPrint[kana={うちざ\tl{˥}まり\tl{˥˩}},ipa={utʃidzamari},pos={名},acc={H+F}]}%
	{%
		\MeaningsBlock{%
	\ryumeaning%
		{親戚。親類}{}{}%
	}%
	}%
	{\tailPrint[]}%
\entry
	{\headPrint[kana={うちすくん\tl{˥}},ipa={utʃi̥su̥kuɴ},pos={動},rui={1},para={うちすかぬ},acc={H}]}%
	{%
		\MeaningsBlock{%
	\ryumeaning%
		{落ち付く}{}{}%
	}%
	}%
	{\tailPrint[]}%
\entry
	{\headPrint[kana={うちな\tl{˥˩}},ipa={utʃina},pos={名},acc={F}]}%
	{%
		\MeaningsBlock{%
	\ryumeaning%
		{内。内側}{}{}%
	}%
	}%
	{\tailPrint[]}%
\entry
	{\headPrint[kana={うちばっさん},ipa={utʃibassaɴ},pos={動（継）}]}%
	{%
		\MeaningsBlock{%
	\ryumeaning%
		{うち忘れる}{忘れるを強めた語}{}%
	}%
	}%
	{\tailPrint[]}%
\entry
	{\headPrint[kana={うちまかすん},ipa={utʃimakasuɴ},pos={動}]}%
	{%
		\MeaningsBlock{%
	\ryumeaning%
		{打ち負かす}{}{}%
	}%
	}%
	{\tailPrint[]}%
\entry
	{\headPrint[kana={うちょー\tl{˥}},ipa={utʃoː},pos={名},acc={H}]}%
	{%
		\MeaningsBlock{%
	\ryumeaning%
		{}{獲った魚を紐に通す漁具}{}%
	}%
	}%
	{\tailPrint[]}%
\entry
	{\headPrint[kana={うちん\tl{˥}},ipa={utʃiɴ},pos={動},rui={3},para={うとぅぬ},acc={H}]}%
	{%
		\MeaningsBlock{%
	\ryumeaning%
		{落ちる}{}{}%
	}%
	}%
	{\tailPrint[]}%
\entry
	{\headPrint[kana={うっす\tl{˥}},ipa={ussu},pos={名},acc={H}]}%
	{%
		\MeaningsBlock{%
	\ryumeaning%
		{後頭部}{}{}%
	}%
	}%
	{\tailPrint[]}%
\entry
	{\headPrint[kana={うつすん\tl{˥}},ipa={utsusuɴ},pos={動},rui={2b},para={うつさぬ},acc={H}]}%
	{%
		\MeaningsBlock{%
	\ryumeaning%
		{写す}{}{}%
	\ryumeaning%
		{感染させる}{}{}%
	}%
	}%
	{\tailPrint[]}%
\entry
	{\headPrint[kana={うったいん\tl{˥}},ipa={uttaiɴ},pos={動},rui={3},acc={H}]}%
	{%
		\MeaningsBlock{%
	\ryumeaning%
		{訴える}{}{}%
	}%
	}%
	{\tailPrint[]}%
\entry
	{\headPrint[kana={うったち},ipa={uttatʃi},pos={名}]}%
	{%
		\MeaningsBlock{%
	\ryumeaning%
		{打ち身。ぶっつける}{}{}%
	}%
	}%
	{\tailPrint[]}%
\entry
	{\headPrint[kana={うっつぁーすん\tl{˥}},ipa={uttsaːsuɴ},pos={動},rui={2},para={うつぁーさぬ},acc={H}]}%
	{%
		\MeaningsBlock{%
	\ryumeaning%
		{打ち付ける}{}{}%
	}%
	}%
	{\tailPrint[]}%
\entry
	{\headPrint[kana={うつるん\tl{˥}},ipa={utsuruɴ},pos={動},rui={1},para={うつらぬ},acc={H}]}%
	{%
		\MeaningsBlock{%
	\ryumeaning%
		{移る}{}{}%
	\ryumeaning%
		{引っ越す}{}{}%
	\ryumeaning%
		{（病気が）うつる。感染する}{}{}%
	\ryumeaning%
		{映る}{}{}%
	\ryumeaning%
		{似合う}{}{}%
	}%
	}%
	{\tailPrint[]}%
\entry
	{\headPrint[kana={うつん\tl{˥}},ipa={utsuɴ},pos={動},rui={2},para={うたぬ},acc={H}]}%
	{%
		\MeaningsBlock{%
	\ryumeaning%
		{打つ}{}{}%
	}%
	}%
	{\tailPrint[]}%
\entry
	{\headPrint[kana={うつん\tl{˥}},ipa={utsuɴ},pos={動},para={うたぬ},acc={H}]}%
	{%
		\MeaningsBlock{%
	\ryumeaning%
		{射撃する}{}{}%
	}%
	}%
	{\tailPrint[]}%
\entry
	{\headPrint[kana={うとー\tl{˥˩}},ipa={utoː},pos={名},acc={F}]}%
	{%
		\MeaningsBlock{%
	\ryumeaning%
		{音。物音}{}{}%
	}%
	}%
	{\tailPrint[]}%
\entry
	{\headPrint[kana={うどーし},ipa={udoːʃi},pos={動}]}%
	{%
		\MeaningsBlock{%
	\ryumeaning%
		{脅す}{}{}%
	}%
	}%
	{\tailPrint[]}%
\entry
	{\headPrint[kana={うとぅーとぅ\tl{˥}},ipa={utuːtu},pos={名},acc={H}]}%
	{%
		\MeaningsBlock{%
	\ryumeaning%
		{年下の兄弟。弟。妹}{}{}%
	}%
	}%
	{\tailPrint[]}%
\entry
	{\headPrint[kana={うどぅがすん\tl{˥}},ipa={udugasuɴ},pos={動},rui={2},para={うどぅがはぬ},acc={H}]}%
	{%
		\MeaningsBlock{%
	\ryumeaning%
		{驚かす}{}{}%
	\ryumeaning%
		{どやす}{}{}%
	}%
	}%
	{\tailPrint[]}%
\entry
	{\headPrint[kana={うとぅさた\tl{˥˩}},ipa={utu̥sḁta},pos={名},acc={F}]}%
	{%
		\MeaningsBlock{%
	\ryumeaning%
		{音沙汰。音信。消息}{}{}%
	}%
	}%
	{\tailPrint[]}%
\entry
	{\headPrint[kana={うとぅだが\tl{˥˩}はん},ipa={utudagahaɴ},pos={形},para={うとだがへぬ},acc={F}]}%
	{%
		\MeaningsBlock{%
	\ryumeaning%
		{名高い。有名だ}{}{}%
	}%
	}%
	{\tailPrint[]}%
\entry
	{\headPrint[kana={うとぅなさ\tl{˥}はん},ipa={utunasahaɴ},pos={形},para={うとなさへぬ},acc={H}]}%
	{%
		\MeaningsBlock{%
	\ryumeaning%
		{大人しい。穏和だ}{}{}%
	}%
	}%
	{\tailPrint[]}%
\entry
	{\headPrint[kana={うとぅるいるん\tl{˥}},ipa={uturuiruɴ},pos={動},para={うとぅるあぬ},acc={H}]}%
	{%
		\MeaningsBlock{%
	\ryumeaning%
		{衰える}{}{}%
	}%
	}%
	{\tailPrint[]}%
\entry
	{\headPrint[kana={うどぅるぎ\tl{˥}},ipa={udurugi},pos={名},acc={H}]}%
	{%
		\MeaningsBlock{%
	\ryumeaning%
		{驚き}{}{}%
	}%
	}%
	{\tailPrint[]}%
\entry
	{\headPrint[kana={うどぅるぐん\tl{˥}},ipa={uduruguɴ},pos={動},rui={1},para={うどぅるがぬ},acc={H}]}%
	{%
		\MeaningsBlock{%
	\ryumeaning%
		{驚く}{}{}%
	}%
	}%
	{\tailPrint[]}%
\entry
	{\headPrint[kana={うとぅるさ\tl{˥}はん},ipa={uturusahaɴ},pos={形},para={うとるさへぬ},acc={H}]}%
	{%
		\MeaningsBlock{%
	\ryumeaning%
		{怖い。恐ろしい}{}{}%
	}%
	}%
	{\tailPrint[]}%
\entry
	{\headPrint[kana={うとぅるん\tl{˥}},ipa={uturuɴ},pos={動},rui={1},para={うとぅらぬ},acc={H}]}%
	{%
		\MeaningsBlock{%
	\ryumeaning%
		{劣る}{}{}%
	}%
	}%
	{\tailPrint[]}%
\entry
	{\headPrint[kana={\josho{}うなー\kako{}ぐ},ipa={unaːgu},pos={名},acc={HL}]}%
	{%
		\MeaningsBlock{%
	\ryumeaning%
		{安心。安堵}{}{}%
	}%
	}%
	{\tailPrint[]}%
\entry
	{\headPrint[kana={うなーぐ\ws{}しゃーん},ipa={unaːgu ʃaːɴ},pos={句},acc={HL HL}]}%
	{%
		\MeaningsBlock{%
	\ryumeaning%
		{（心配事がなくなり）ほっとする}{}{}%
	}%
	}%
	{\tailPrint[]}%
\entry
	{\headPrint[kana={\josho{}うなー\kako{}ぐ\ws{}なるん\tl{˩˥}},ipa={unaːgu naruɴ},pos={句},rui={1},para={うなーぐならぬ},acc={HL R}]}%
	{%
		\MeaningsBlock{%
	\ryumeaning%
		{安心する。安堵する。ほっとする}{}{}%
	}%
	}%
	{\tailPrint[]}%
\entry
	{\headPrint[kana={うなん\tl{˥}},ipa={unaɴ},pos={名},acc={H}]}%
	{%
		\MeaningsBlock{%
	\ryumeaning%
		{ウナギ}{}{}%
	}%
	}%
	{\tailPrint[]}%
\entry
	{\headPrint[kana={うぬ\tl{˥˩}},ipa={unu},pos={連体},acc={F}]}%
	{%
		\MeaningsBlock{%
	\ryumeaning%
		{あの。その}{\emb{うぬ-ぴとぅ}「あの人」など}{}%
	}%
	}%
	{\tailPrint[]}%
\entry
	{\headPrint[kana={うぬ\ws{}まーま},ipa={unumaːma},pos={句}]}%
	{%
		\MeaningsBlock{%
	\ryumeaning%
		{そのまま}{}{}%
	}%
	}%
	{\tailPrint[]}%
\entry
	{\headPrint[kana={うぶい\tl{˥}},ipa={ubui},pos={名},acc={H}]}%
	{%
		\MeaningsBlock{%
	\ryumeaning%
		{覚え。記憶}{}{}%
	}%
	}%
	{\tailPrint[]}%
\entry
	{\headPrint[kana={うぶいるん\tl{˥}},ipa={ubuiruɴ},pos={動},acc={H}]}%
	{%
		\MeaningsBlock{%
	\ryumeaning%
		{覚える}{}{}%
	}%
	}%
	{\tailPrint[]}%
\entry
	{\headPrint[kana={うぶいん},ipa={ubuiɴ},pos={動}]}%
	{%
		\MeaningsBlock{%
	\ryumeaning%
		{覚える}{}{}%
	}%
	}%
	{\tailPrint[]}%
\entry
	{\headPrint[kana={うぶりるん\tl{˥˩}},ipa={uburiruɴ},pos={動},para={うぶるぬ},acc={F}]}%
	{%
		\MeaningsBlock{%
	\ryumeaning%
		{溺れる}{}{}%
	}%
	}%
	{\tailPrint[]}%
\entry
	{\headPrint[kana={うむいあたるん},ipa={umuiataruɴ}]}%
	{%
		\MeaningsBlock{%
	\ryumeaning%
		{思い当たる}{}{}%
	}%
	}%
	{\tailPrint[]}%
\entry
	{\headPrint[kana={うむいきすん\tl{˥}},ipa={umuiki̥suɴ},pos={動},acc={H}]}%
	{%
		\MeaningsBlock{%
	\ryumeaning%
		{思い切る。決心する}{}{}%
	\ryumeaning%
		{諦める}{}{}%
	\ryumeaning%
		{果たす}{}{}%
	}%
	}%
	{\tailPrint[]}%
\entry
	{\headPrint[kana={うむいくがりん\tl{˥}},ipa={umuikugariɴ},pos={動},para={うむいくがるぬ},acc={H}]}%
	{%
		\MeaningsBlock{%
	\ryumeaning%
		{思い焦がれる。深く思う}{}{}%
	}%
	}%
	{\tailPrint[]}%
\entry
	{\headPrint[kana={うむいくむん\tl{˥}},ipa={umuiku̥muɴ},pos={動},rui={1},para={うむいくまぬ},acc={H}]}%
	{%
		\MeaningsBlock{%
	\ryumeaning%
		{思い込む}{}{}%
	}%
	}%
	{\tailPrint[]}%
\entry
	{\headPrint[kana={うむい\tl{˥}すくん\tl{˥˩}},ipa={umuisu̥kuɴ},pos={動},rui={1},para={うむいすかぬ},acc={H F}]}%
	{%
		\MeaningsBlock{%
	\ryumeaning%
		{思い付く}{}{}%
	}%
	}%
	{\tailPrint[]}%
\entry
	{\headPrint[kana={うむいたつん\tl{˥}},ipa={umuitḁtsuɴ},pos={動},para={うむいたたぬ},acc={H}]}%
	{%
		\MeaningsBlock{%
	\ryumeaning%
		{思い立つ}{}{}%
	}%
	}%
	{\tailPrint[]}%
\entry
	{\headPrint[kana={うむいつみるん\tl{˥}},ipa={umuitsu̥miruɴ},pos={動},rui={3},para={うむいつむぬ},acc={H}]}%
	{%
		\MeaningsBlock{%
	\ryumeaning%
		{思い詰める}{}{}%
	}%
	}%
	{\tailPrint[]}%
\entry
	{\headPrint[kana={うむいつみん\tl{˥}},ipa={umuitsu̥miɴ},pos={動},rui={3},para={うむいつむぬ},acc={H}]}%
	{%
		\MeaningsBlock{%
	\ryumeaning%
		{思い詰める}{}{}%
	}%
	}%
	{\tailPrint[]}%
\entry
	{\headPrint[kana={うむい\tl{˥}ぬぐすん\tl{˩˥}},ipa={umuinugusuɴ},pos={動},rui={2},para={うむい　ぬぐはぬ},acc={H R}]}%
	{%
		\MeaningsBlock{%
	\ryumeaning%
		{思い残す}{}{}%
	}%
	}%
	{\tailPrint[]}%
\entry
	{\headPrint[kana={うむい\tl{˥}のーすん\tl{˩˥}},ipa={umuinoːsuɴ},pos={動},rui={2b},para={うむいのーさぬ},acc={H+R}]}%
	{%
		\MeaningsBlock{%
	\ryumeaning%
		{思い直す}{}{}%
	}%
	}%
	{\tailPrint[]}%
\entry
	{\headPrint[kana={うむいんだすん\tl{˥}},ipa={umuindasuɴ},pos={動},para={うむいんだるぬ},acc={H}]}%
	{%
		\MeaningsBlock{%
	\ryumeaning%
		{思い出す}{}{}%
	}%
	}%
	{\tailPrint[]}%
\entry
	{\headPrint[kana={うむっさーすん},ipa={umussaːsuɴ},pos={動}]}%
	{%
		\MeaningsBlock{%
	\ryumeaning%
		{楽しむ}{}{}%
	}%
	}%
	{\tailPrint[]}%
\entry
	{\headPrint[kana={うや\tl{˥}},ipa={uja},pos={名},acc={H}]}%
	{%
		\MeaningsBlock{%
	\ryumeaning%
		{親。肉親}{王府時代には役人をも指す}{}%
	}%
	}%
	{\tailPrint[]}%
\entry
	{\headPrint[kana={うやかた\tl{˥}},ipa={ujakḁta},pos={名},acc={H}]}%
	{%
		\MeaningsBlock{%
	\ryumeaning%
		{親方}{}{}%
	}%
	}%
	{\tailPrint[]}%
\entry
	{\headPrint[kana={うやぎ\tl{˥}},ipa={ujagi},pos={名},acc={H}]}%
	{%
		\MeaningsBlock{%
	\ryumeaning%
		{富裕。金持ち}{}{}%
	}%
	}%
	{\tailPrint[]}%
\entry
	{\headPrint[kana={\josho{}うやぎ\kako{}ひー},ipa={ujagihiː},pos={名},acc={H+L}]}%
	{%
		\MeaningsBlock{%
	\ryumeaning%
		{富裕者。金持ちの家}{}{}%
	}%
	}%
	{\tailPrint[]}%
\entry
	{\headPrint[kana={うやぐ\tl{˥}},ipa={ujagu},pos={名},acc={H}]}%
	{%
		\MeaningsBlock{%
	\ryumeaning%
		{親族。親戚}{}{}%
	}%
	}%
	{\tailPrint[note={親子の転訛か}]}%
\entry
	{\headPrint[kana={うやだり\tl{˥}},ipa={ujadari},pos={名},acc={H}]}%
	{%
		\MeaningsBlock{%
	\ryumeaning%
		{公務。公事}{王府時代役人の指示した公事、公務}{}%
	}%
	}%
	{\tailPrint[]}%
\entry
	{\headPrint[kana={うやっすり\tl{˥}},ipa={ujassuri},pos={名},acc={H}]}%
	{%
		\MeaningsBlock{%
	\ryumeaning%
		{祈り。祈願}{}{}%
	}%
	}%
	{\tailPrint[]}%
\entry
	{\headPrint[kana={うやぴぃとぅ\tl{˥}},ipa={ujapi̥tu},pos={名},acc={H}]}%
	{%
		\MeaningsBlock{%
	\ryumeaning%
		{先祖の霊。先祖神}{}{}%
	}%
	}%
	{\tailPrint[]}%
\entry
	{\headPrint[kana={\josho{}うや\kako{}ふぁ},ipa={ujafa},pos={名},acc={H+L}]}%
	{%
		\MeaningsBlock{%
	\ryumeaning%
		{親子}{}{}%
	}%
	}%
	{\tailPrint[]}%
\entry
	{\headPrint[kana={うやまいるん\tl{˥˩}},ipa={ujamairuɴ},pos={動},para={うやまーるぬ},acc={F}]}%
	{%
		\MeaningsBlock{%
	\ryumeaning%
		{敬う。尊敬する}{}{}%
	}%
	}%
	{\tailPrint[]}%
\entry
	{\headPrint[kana={うやめーるん\tl{˥˩}},ipa={ujameːruɴ},pos={動},rui={3},para={うやめーぬ},acc={F}]}%
	{%
		\MeaningsBlock{%
	\ryumeaning%
		{敬う。尊敬する}{}{}%
	}%
	}%
	{\tailPrint[]}%
\entry
	{\headPrint[kana={うゆぶん\tl{˥˩}},ipa={ujubuɴ},pos={動},rui={1},para={うゆばぬ},acc={F}]}%
	{%
		\MeaningsBlock{%
	\ryumeaning%
		{及ぶ}{}{}%
	\ryumeaning%
		{相応する}{}{}%
	}%
	}%
	{\tailPrint[]}%
\entry
	{\headPrint[kana={うら\tl{˥}},ipa={ura},pos={名},acc={H}]}%
	{%
		\MeaningsBlock{%
	\ryumeaning%
		{湾}{王府時代は「蔵元」を指した}{}%
	}%
	}%
	{\tailPrint[]}%
\entry
	{\headPrint[kana={うらがいすん\tl{˥}},ipa={uragaisuɴ},pos={動},acc={H}]}%
	{%
		\MeaningsBlock{%
	\ryumeaning%
		{裏返す}{}{}%
	}%
	}%
	{\tailPrint[]}%
\entry
	{\headPrint[kana={うらげーすん\tl{˥}},ipa={urageːsuɴ},pos={動},rui={2b},para={うらげーさぬ},acc={H}]}%
	{%
		\MeaningsBlock{%
	\ryumeaning%
		{裏返す}{}{}%
	}%
	}%
	{\tailPrint[]}%
\entry
	{\headPrint[kana={うらすん\tl{˥}},ipa={urasuɴ},pos={動},rui={2},para={うらはぬ},acc={H}]}%
	{%
		\MeaningsBlock{%
	\ryumeaning%
		{降ろす。下ろす}{}{}%
	}%
	}%
	{\tailPrint[]}%
\entry
	{\headPrint[kana={うらすん\tl{˥}},ipa={urasuɴ},pos={動},rui={2},para={うらはぬ},acc={H}]}%
	{%
		\MeaningsBlock{%
	\ryumeaning%
		{刻む}{刃物で切って細かくする}{}%
	}%
	}%
	{\tailPrint[]}%
\entry
	{\headPrint[kana={うらみしゃはん},ipa={uramiʃahaɴ}]}%
	{%
		\MeaningsBlock{%
	\ryumeaning%
		{羨ましい}{}{}%
	}%
	}%
	{\tailPrint[]}%
\entry
	{\headPrint[kana={うらんだー\tl{˥}},ipa={urandaː},pos={名},acc={H}]}%
	{%
		\MeaningsBlock{%
	\ryumeaning%
		{欧米人。西洋人}{}{}%
	}%
	}%
	{\tailPrint[]}%
\entry
	{\headPrint[kana={うり\tl{˥˩}},ipa={uri},pos={名},acc={F}]}%
	{%
		\MeaningsBlock{%
	\ryumeaning%
		{あれ。それ}{}{}%
	}%
	}%
	{\tailPrint[]}%
\entry
	{\headPrint[kana={うりー\tl{˥}},ipa={uriː},pos={名},acc={H}]}%
	{%
		\MeaningsBlock{%
	\ryumeaning%
		{潤い}{畑が潤うこと}{}%
	}%
	}%
	{\tailPrint[]}%
\entry
	{\headPrint[kana={うりげー\tl{˥}},ipa={urigeː},pos={名},acc={H}]}%
	{%
		\MeaningsBlock{%
	\ryumeaning%
		{降り井戸}{降りて水を汲む式の井戸}{}%
	}%
	}%
	{\tailPrint[]}%
\entry
	{\headPrint[kana={うりじん\tl{˥}},ipa={uridʒiɴ},pos={名},acc={H}]}%
	{%
		\MeaningsBlock{%
	\ryumeaning%
		{春。春季}{「潤いの季節」から}{}%
	}%
	}%
	{\tailPrint[]}%
\entry
	{\headPrint[kana={うりん\tl{˥}},ipa={uriɴ},pos={名},acc={H}]}%
	{%
		\MeaningsBlock{%
	\ryumeaning%
		{胡瓜}{}{}%
	}%
	}%
	{\tailPrint[]}%
\entry
	{\headPrint[kana={うる\tl{˥˩}},ipa={uru},pos={名},acc={F}]}%
	{%
		\MeaningsBlock{%
	\ryumeaning%
		{\ruby[g]{珊瑚}{さんご}。\ruby[g]{珊瑚}{さんご}の石}{}{}%
	}%
	}%
	{\tailPrint[]}%
\entry
	{\headPrint[kana={うるばたぎ\tl{˥}},ipa={urubatagi},pos={名},acc={H}]}%
	{%
		\MeaningsBlock{%
	\ryumeaning%
		{田虫}{皮膚病名}{}%
	}%
	}%
	{\tailPrint[]}%
\entry
	{\headPrint[kana={うるん\tl{˥˩}},ipa={uruɴ},pos={動},rui={1},para={うらぬ},acc={F}]}%
	{%
		\MeaningsBlock{%
	\ryumeaning%
		{織る}{機を織ること}{}%
	}%
	}%
	{\tailPrint[]}%
\entry
	{\headPrint[kana={うるんがに\tl{˩˥}},ipa={uruŋgani},pos={名},acc={R}]}%
	{%
		\MeaningsBlock{%
	\ryumeaning%
		{指輪}{}{}%
	}%
	}%
	{\tailPrint[]}%
\entry
	{\headPrint[kana={うるんぺー\tl{˥˩}},ipa={urumpeː},pos={名},acc={F}]}%
	{%
		\MeaningsBlock{%
	\ryumeaning%
		{石灰}{「サンゴの灰」の義}{}%
	}%
	}%
	{\tailPrint[]}%
\entry
	{\headPrint[kana={うわー\tl{˥}},ipa={uwaː},pos={名},acc={H}]}%
	{%
		\MeaningsBlock{%
	\ryumeaning%
		{豚}{}{}%
	}%
	}%
	{\tailPrint[]}%
\entry
	{\headPrint[kana={うわるん\tl{˥˩}},ipa={uwaruɴ},pos={動},rui={1},para={うわらぬ},acc={F}]}%
	{%
		\MeaningsBlock{%
	\ryumeaning%
		{終わる。仕舞う}{}{}%
	}%
	}%
	{\tailPrint[]}%
\entry
	{\headPrint[kana={うわん\tl{˥}ひー\tl{˩˥}},ipa={uwaɴhiː},pos={名},acc={H+R}]}%
	{%
		\MeaningsBlock{%
	\ryumeaning%
		{豚小屋}{}{}%
	}%
	}%
	{\tailPrint[]}%
\entry
	{\headPrint[kana={うん\tl{˥}},ipa={uɴ},pos={名},acc={H}]}%
	{%
		\MeaningsBlock{%
	\ryumeaning%
		{鬼}{}{}%
	}%
	}%
	{\tailPrint[]}%
\entry
	{\headPrint[kana={うん\tl{˥˩}},ipa={uɴ},pos={名},acc={F}]}%
	{%
		\MeaningsBlock{%
	\ryumeaning%
		{運。幸運}{}{}%
	}%
	}%
	{\tailPrint[]}%
\entry
	{\headPrint[kana={うん\tl{˥˩}},ipa={uɴ},pos={名},acc={F}]}%
	{%
		\MeaningsBlock{%
	\ryumeaning%
		{海栗}{}{}%
	}%
	}%
	{\tailPrint[]}%
\entry
	{\headPrint[kana={うんしゅく\tl{˩˥}},ipa={unʃu̥ku},pos={副},acc={R}]}%
	{%
		\MeaningsBlock{%
	\ryumeaning%
		{それほど}{}{}%
	}%
	}%
	{\tailPrint[]}%
\entry
	{\headPrint[kana={うんつぇー},ipa={untseː},pos={名}]}%
	{%
		\MeaningsBlock{%
	\ryumeaning%
		{エンサイ}{野菜名}{}%
	}%
	}%
	{\tailPrint[]}%
\entry
	{\headPrint[kana={\josho{}うん\kako{}とぅ},ipa={untu},pos={副},acc={HL}]}%
	{%
		\MeaningsBlock{%
	\ryumeaning%
		{しっかり。精一杯}{}{}%
	}%
	}%
	{\tailPrint[]}%
\entry
	{\headPrint[kana={うんどぅん\tl{˥}},ipa={unduɴ},pos={名},acc={H}]}%
	{%
		\MeaningsBlock{%
	\ryumeaning%
		{}{位牌をまとめた大きな位牌}{}%
	}%
	}%
	{\tailPrint[]}%
\entry
	{\headPrint[kana={うんぬ\tl{˥˩}\ws{}ねーぬ\tl{˩˥}},ipa={unnu neːnu},pos={句},acc={F R}]}%
	{%
		\MeaningsBlock{%
	\ryumeaning%
		{運がない。不運だ}{}{}%
	}%
	}%
	{\tailPrint[]}%
\LetterHead{え}
\entry
	{\headPrint[kana={えー},ipa={eː},pos={感}]}%
	{%
		\MeaningsBlock{%
	\ryumeaning%
		{そう}{軽い肯定}{}%
	}%
	}%
	{\tailPrint[]}%
\entry
	{\headPrint[kana={えー\tl{˥˩}},ipa={eː},pos={名},acc={F}]}%
	{%
		\MeaningsBlock{%
	\ryumeaning%
		{絵。図画。絵画}{}{}%
	}%
	}%
	{\tailPrint[]}%
\entry
	{\headPrint[kana={えーさ},ipa={eːsa},pos={感}]}%
	{%
		\MeaningsBlock{%
	\ryumeaning%
		{そうだよ}{}{}%
	}%
	}%
	{\tailPrint[]}%
\entry
	{\headPrint[kana={\josho{}えー\kako{}\ws{}しみるん\tl{˥˩}},ipa={eː ʃi̥miruɴ},pos={句},acc={H L}]}%
	{%
		\MeaningsBlock{%
	\ryumeaning%
		{そうさせる}{}{}%
	}%
	}%
	{\tailPrint[]}%
\entry
	{\headPrint[kana={えー\tl{˥}\ws{}すん\tl{˥˩}},ipa={eː suɴ},pos={句},acc={H F}]}%
	{%
		\MeaningsBlock{%
	\ryumeaning%
		{そうする}{}{}%
	}%
	}%
	{\tailPrint[]}%
\entry
	{\headPrint[kana={えーちゅー},ipa={eːtʃuː},pos={感}]}%
	{%
		\MeaningsBlock{%
	\ryumeaning%
		{～だそうだ}{}{}%
	}%
	}%
	{\tailPrint[]}%
\entry
	{\headPrint[kana={えーなー},ipa={eːnaː},pos={感}]}%
	{%
		\MeaningsBlock{%
	\ryumeaning%
		{そうか。なるほど}{}{}%
	}%
	}%
	{\tailPrint[]}%
\entry
	{\headPrint[kana={\josho{}えー\kako{}ぬ},ipa={eːnu},pos={句},acc={HL}]}%
	{%
		\MeaningsBlock{%
	\ryumeaning%
		{そんな。そんなこと}{}{}%
	}%
	}%
	{\tailPrint[]}%
\entry
	{\headPrint[kana={\josho{}えー\kako{}ばぎる},ipa={eːbagiru},pos={句},acc={HL}]}%
	{%
		\MeaningsBlock{%
	\ryumeaning%
		{そうしか}{\emb{えーばぎる-なる}など}{}%
	}%
	}%
	{\tailPrint[]}%
\entry
	{\headPrint[kana={\josho{}えー\kako{}\ws{}やちゃら},ipa={eː jatʃara},pos={句},acc={H L}]}%
	{%
		\MeaningsBlock{%
	\ryumeaning%
		{そうなら。それなら}{}{}%
	}%
	}%
	{\tailPrint[]}%
\entry
	{\headPrint[kana={\josho{}えー\kako{}\ws{}やばん},ipa={eː jabaɴ},pos={句},acc={H L}]}%
	{%
		\MeaningsBlock{%
	\ryumeaning%
		{そうでも}{}{}%
	}%
	}%
	{\tailPrint[]}%
\entry
	{\headPrint[kana={えーるやろ},ipa={eːrujaro},pos={感}]}%
	{%
		\MeaningsBlock{%
	\ryumeaning%
		{そうだ。その通りだ}{}{}%
	}%
	}%
	{\tailPrint[]}%
\entry
	{\headPrint[kana={\josho{}え\kako{}した},ipa={eʃi̥ta},pos={句},acc={HL}]}%
	{%
		\MeaningsBlock{%
	\ryumeaning%
		{そうして}{}{}%
	}%
	}%
	{\tailPrint[]}%
\entry
	{\headPrint[kana={えすが},ipa={esu̥ka},pos={句}]}%
	{%
		\MeaningsBlock{%
	\ryumeaning%
		{だが。しかし。けれども}{}{}%
	}%
	}%
	{\tailPrint[]}%
\entry
	{\headPrint[kana={\josho{}えす\kako{}がら},ipa={esu̥gara},pos={句},acc={HL}]}%
	{%
		\MeaningsBlock{%
	\ryumeaning%
		{だから。そうだから}{}{}%
	}%
	}%
	{\tailPrint[]}%
\entry
	{\headPrint[kana={えちー\tl{˥˩}},ipa={etʃiː},pos={句},acc={F}]}%
	{%
		\MeaningsBlock{%
	\ryumeaning%
		{そうだから}{}{}%
	}%
	}%
	{\tailPrint[]}%
\entry
	{\headPrint[kana={えちる},ipa={etʃiru},pos={接続},acc={-}]}%
	{%
		\MeaningsBlock{%
	\ryumeaning%
		{そうだから}{}{}%
	}%
	}%
	{\tailPrint[]}%
\entry
	{\headPrint[kana={えにすかすん\tl{˥˩}},ipa={enisu̥kasuɴ},pos={動},rui={2},para={えにすかはぬ},acc={F+F}]}%
	{%
		\MeaningsBlock{%
	\ryumeaning%
		{言って聞かせる。指図する。説き聞かせる。教訓する}{}{}%
	}%
	}%
	{\tailPrint[]}%
\entry
	{\headPrint[kana={えぬぱち},ipa={enupatʃi̥},pos={句}]}%
	{%
		\MeaningsBlock{%
	\ryumeaning%
		{その\ruby[g]{筈}{はず}}{}{}%
	}%
	}%
	{\tailPrint[]}%
\entry
	{\headPrint[kana={えぬふぁ\tl{˥˩}},ipa={enufa},pos={名},acc={F}]}%
	{%
		\MeaningsBlock{%
	\ryumeaning%
		{アイゴ}{魚の種類}{}%
	}%
	}%
	{\tailPrint[]}%
\entry
	{\headPrint[kana={えぬ\ws{}むぬ},ipa={enu munu},pos={句}]}%
	{%
		\MeaningsBlock{%
	\ryumeaning%
		{そんなもの}{}{}%
	}%
	}%
	{\tailPrint[]}%
\entry
	{\headPrint[kana={えぬん\tl{˥˩}},ipa={enuɴ},pos={動},rui={1},para={えなぬ},acc={F}]}%
	{%
		\MeaningsBlock{%
	\ryumeaning%
		{言う}{}{}%
	}%
	}%
	{\tailPrint[]}%
\entry
	{\headPrint[kana={えん\tl{˩˥}},ipa={eɴ},pos={名},acc={R}]}%
	{%
		\MeaningsBlock{%
	\ryumeaning%
		{来年。来る年}{東村の発音}{}%
	}%
	}%
	{\tailPrint[]}%
\entry
	{\headPrint[kana={えんだりるん\tl{˥}},ipa={endariruɴ},pos={動},rui={3},para={えんだるぬ},acc={H}]}%
	{%
		\MeaningsBlock{%
	\ryumeaning%
		{喧嘩する。争う}{}{}%
	}%
	}%
	{\tailPrint[]}%
\entry
	{\headPrint[kana={えんだりん\tl{˥}},ipa={endariɴ},pos={動},rui={3},para={えんだるぬ},acc={H}]}%
	{%
		\MeaningsBlock{%
	\ryumeaning%
		{喧嘩する。争う}{}{}%
	}%
	}%
	{\tailPrint[]}%
\LetterHead{お}
\entry
	{\headPrint[kana={おー},ipa={oː},pos={感}]}%
	{%
		\MeaningsBlock{%
	\ryumeaning%
		{はい}{目上に対しての返事。敬語。目下には\emb{ん}と返事する}{}%
	}%
	}%
	{\tailPrint[]}%
\entry
	{\headPrint[kana={おーがーるん\tl{˥˩}},ipa={oːgaːruɴ},pos={動},rui={1},para={おーがーらぬ},acc={F}]}%
	{%
		\MeaningsBlock{%
	\ryumeaning%
		{浮かぶ。浮かれる}{}{}%
	}%
	}%
	{\tailPrint[]}%
\entry
	{\headPrint[kana={おーがるん\tl{˥˩}},ipa={oːgaruɴ},pos={動},rui={1},para={おーがらぬ},acc={F}]}%
	{%
		\MeaningsBlock{%
	\ryumeaning%
		{浮かぶ。浮遊する}{}{}%
	}%
	}%
	{\tailPrint[]}%
\entry
	{\headPrint[kana={おーぎるん\tl{˥˩}},ipa={oːgiruɴ},pos={動},rui={3},para={おーぐぬ},acc={F}]}%
	{%
		\MeaningsBlock{%
	\ryumeaning%
		{浮ける。浮かべる。浮かせる。漂わす}{}{}%
	}%
	}%
	{\tailPrint[]}%
\entry
	{\headPrint[kana={お\josho{}ー\kako{}し},ipa={oːʃi},pos={副},acc={LHL}]}%
	{%
		\MeaningsBlock{%
	\ryumeaning%
		{青々と。緑色に}{}{}%
	}%
	}%
	{\tailPrint[]}%
\entry
	{\headPrint[kana={おーし\ws{}おるん},ipa={oːʃi oruɴ},pos={句},acc={-}]}%
	{%
		\MeaningsBlock{%
	\ryumeaning%
		{召し上がる。お上がりになる}{「食べる」の尊敬語}{}%
	}%
	}%
	{\tailPrint[]}%
\entry
	{\headPrint[kana={おー\kako{}し\ws{}なるん\tl{˩˥}},ipa={oːʃi naruɴ},pos={句},acc={HL R}]}%
	{%
		\MeaningsBlock{%
	\ryumeaning%
		{青ばむ}{青味を帯びる}{}%
	}%
	}%
	{\tailPrint[]}%
\entry
	{\headPrint[kana={おーしゃ\tl{˥}},ipa={oːʃa},pos={名},acc={H}]}%
	{%
		\MeaningsBlock{%
	\ryumeaning%
		{村番所。村事務所}{王府時代の村番所}{}%
	}%
	}%
	{\tailPrint[]}%
\entry
	{\headPrint[kana={おーじゃなー\tl{˥}},ipa={oːdʒanaː},pos={名},acc={H}]}%
	{%
		\MeaningsBlock{%
	\ryumeaning%
		{青大将}{ヘビの種類。無毒}{}%
	}%
	}%
	{\tailPrint[]}%
\entry
	{\headPrint[kana={おーすむぬ\tl{˥˩}},ipa={osumunu},pos={名},acc={F}]}%
	{%
		\MeaningsBlock{%
	\ryumeaning%
		{献上品}{}{}%
	}%
	}%
	{\tailPrint[]}%
\entry
	{\headPrint[kana={おーすん\tl{˥}},ipa={oːsuɴ},pos={動},rui={2},para={おーはぬ},acc={H}]}%
	{%
		\MeaningsBlock{%
	\ryumeaning%
		{（牛馬に荷物を）背負わす}{}{}%
	}%
	}%
	{\tailPrint[]}%
\entry
	{\headPrint[kana={おーすん\tl{˥˩}},ipa={osuɴ},pos={動},rui={2},para={おさぬ},acc={F}]}%
	{%
		\MeaningsBlock{%
	\ryumeaning%
		{差し上げる}{謙譲語}{}%
	}%
	}%
	{\tailPrint[]}%
\entry
	{\headPrint[kana={おーばー\tl{˥˩}},ipa={oːbaː},pos={名},acc={F}]}%
	{%
		\MeaningsBlock{%
	\ryumeaning%
		{余り。余分}{}{}%
	}%
	}%
	{\tailPrint[]}%
\entry
	{\headPrint[kana={おーぱーと\tl{˥}},ipa={oːpaːto},pos={名},acc={H}]}%
	{%
		\MeaningsBlock{%
	\ryumeaning%
		{青鳩}{}{}%
	}%
	}%
	{\tailPrint[]}%
\entry
	{\headPrint[kana={おーふつぁ\tl{˥˩}はーん},ipa={oːfu̥tsahaːɴ},pos={形},acc={F}]}%
	{%
		\MeaningsBlock{%
	\ryumeaning%
		{青臭い。生臭い}{魚や牛，豚などの肉の臭さなどに言う}{}%
	}%
	}%
	{\tailPrint[]}%
\entry
	{\headPrint[kana={おーむん\tl{˥˩}},ipa={oːmuɴ},pos={動},acc={F}]}%
	{%
		\MeaningsBlock{%
	\ryumeaning%
		{青む。青ばむ。青くなる}{}{}%
	}%
	}%
	{\tailPrint[]}%
\entry
	{\headPrint[kana={おーらが\tl{˥˩}},ipa={oːraga},pos={句},acc={F}]}%
	{%
		\MeaningsBlock{%
	\ryumeaning%
		{風上へ}{風上へ向かうこと}{}%
	}%
	}%
	{\tailPrint[]}%
\entry
	{\headPrint[kana={おーるん\tl{˥}},ipa={oːruɴ},pos={動},rui={1},para={おーらぬ},acc={H}]}%
	{%
		\MeaningsBlock{%
	\ryumeaning%
		{来られる。いらっしゃる}{「来る」「行く」などの尊敬語}{}%
	}%
	}%
	{\tailPrint[]}%
\entry
	{\headPrint[kana={おこらすん\tl{˥}},ipa={okorasuɴ},pos={動},rui={2},para={おこらはぬ},acc={H}]}%
	{%
		\MeaningsBlock{%
	\ryumeaning%
		{（人を）刺激して怒らせる}{}{}%
	}%
	}%
	{\tailPrint[]}%
\entry
	{\headPrint[kana={おしき\tl{˥˩}},ipa={oʃi̥ki},pos={名},acc={F}]}%
	{%
		\MeaningsBlock{%
	\ryumeaning%
		{天気。天候}{}{}%
	}%
	}%
	{\tailPrint[]}%
\entry
	{\headPrint[kana={おった\tl{˥}},ipa={otta},pos={名},acc={H}]}%
	{%
		\MeaningsBlock{%
	\ryumeaning%
		{蛙}{オタマジャクシは\emb{たらぐ}}{}%
	}%
	}%
	{\tailPrint[]}%
\entry
	{\headPrint[kana={おび\tl{˥}},ipa={obi},pos={副},acc={H}]}%
	{%
		\MeaningsBlock{%
	\ryumeaning%
		{それだけ。それまで。終わり}{}{}%
	}%
	}%
	{\tailPrint[]}%
\entry
	{\headPrint[kana={おま\tl{˥}はん},ipa={omahaɴ},pos={形},para={おまへぬ},acc={H}]}%
	{%
		\MeaningsBlock{%
	\ryumeaning%
		{気分が悪い。頭痛がする}{}{}%
	}%
	}%
	{\tailPrint[]}%
\entry
	{\headPrint[kana={おわるん\tl{˥˩}},ipa={owaruɴ},pos={動},rui={1},para={おわらぬ},acc={F}]}%
	{%
		\MeaningsBlock{%
	\ryumeaning%
		{終わる。終了する}{}{}%
	}%
	}%
	{\tailPrint[]}%
\entry
	{\headPrint[kana={おん\tl{˥}},ipa={oɴ},pos={名},acc={H}]}%
	{%
		\MeaningsBlock{%
	\ryumeaning%
		{恩。恩義}{}{}%
	}%
	}%
	{\tailPrint[]}%
\entry
	{\headPrint[kana={おん\tl{˥}},ipa={oɴ},pos={名},acc={H}]}%
	{%
		\MeaningsBlock{%
	\ryumeaning%
		{\ruby[g]{団扇}{うちわ}。\ruby[g]{扇}{おうぎ}}{}{}%
	}%
	}%
	{\tailPrint[]}%
\entry
	{\headPrint[kana={おんぎ\tl{˥}},ipa={oŋgi},pos={名},acc={H}]}%
	{%
		\MeaningsBlock{%
	\ryumeaning%
		{\ruby[g]{団扇}{うちわ}。\ruby[g]{扇}{おうぎ}}{}{}%
	}%
	}%
	{\tailPrint[]}%
\entry
	{\headPrint[kana={おんぐん\tl{˥}},ipa={oŋguɴ},pos={動},rui={1},para={おんがぬ},acc={H}]}%
	{%
		\MeaningsBlock{%
	\ryumeaning%
		{扇ぐ}{}{}%
	}%
	}%
	{\tailPrint[]}%
\entry
	{\headPrint[kana={おんだ\tl{˥}},ipa={onda},pos={名},acc={H}]}%
	{%
		\MeaningsBlock{%
	\ryumeaning%
		{\ruby[g]{畚}{もっこ}}{縄製の運搬具}{}%
	}%
	}%
	{\tailPrint[]}%
\entry
	{\headPrint[kana={おんだ\tl{˥}},ipa={onda},pos={名},acc={H}]}%
	{%
		\MeaningsBlock{%
	\ryumeaning%
		{海水浴}{}{}%
	}%
	}%
	{\tailPrint[]}%
\LetterHead{か}
\entry
	{\headPrint[kana={が},ipa={ga},pos={助}]}%
	{%
		\MeaningsBlock{%
	\ryumeaning%
		{～の方へ}{\emb{ありが}「東の方へ」など}{}%
	}%
	}%
	{\tailPrint[]}%
\entry
	{\headPrint[kana={かー\tl{˥˩}},ipa={kaː},pos={名},acc={F}]}%
	{%
		\MeaningsBlock{%
	\ryumeaning%
		{匂い。香り}{}{}%
	}%
	}%
	{\tailPrint[]}%
\entry
	{\headPrint[kana={がー\tl{˩˥}},ipa={gaː},pos={名},acc={R}]}%
	{%
		\MeaningsBlock{%
	\ryumeaning%
		{\ruby[g]{我}{が}。根性。忍耐}{}{}%
	}%
	}%
	{\tailPrint[]}%
\entry
	{\headPrint[kana={がーがー},ipa={gaːgaː},pos={擬}]}%
	{%
		\MeaningsBlock{%
	\ryumeaning%
		{がやがや}{やかましくしゃべるさま}{}%
	\ryumeaning%
		{かーかー}{カラスが鳴くさま}{}%
	}%
	}%
	{\tailPrint[]}%
\entry
	{\headPrint[kana={かーき\tl{˥}},ipa={kaːki},pos={名},acc={H}]}%
	{%
		\MeaningsBlock{%
	\ryumeaning%
		{約束。賭け}{}{}%
	}%
	}%
	{\tailPrint[]}%
\entry
	{\headPrint[kana={かーぎ\tl{˥}},ipa={kaːgi},pos={名},acc={H}]}%
	{%
		\MeaningsBlock{%
	\ryumeaning%
		{鉤の手}{}{}%
	}%
	}%
	{\tailPrint[]}%
\entry
	{\headPrint[kana={がーし},ipa={gaːʃi̥},pos={助詞}]}%
	{%
		\MeaningsBlock{%
	\ryumeaning%
		{～だけ。～ばかり}{}{}%
	}%
	}%
	{\tailPrint[]}%
\entry
	{\headPrint[kana={がーし\tl{˥˩}},ipa={gaːʃi},pos={名},acc={F}]}%
	{%
		\MeaningsBlock{%
	\ryumeaning%
		{飢饉。餓死}{飢饉は蛾死につながつた}{}%
	}%
	}%
	{\tailPrint[]}%
\entry
	{\headPrint[kana={がーずさはん},ipa={gaːdzusahaɴ},pos={形},acc={-}]}%
	{%
		\MeaningsBlock{%
	\ryumeaning%
		{忍耐力が強い。我慢強い}{}{}%
	\ryumeaning%
		{強情だ}{}{}%
	}%
	}%
	{\tailPrint[]}%
\entry
	{\headPrint[kana={かーすん\tl{˥˩}},ipa={kaːsuɴ},pos={動},rui={2b},para={かーさぬ},acc={F}]}%
	{%
		\MeaningsBlock{%
	\ryumeaning%
		{匂う。匂いがする}{}{}%
	}%
	}%
	{\tailPrint[]}%
\entry
	{\headPrint[kana={かーち\tl{˥}},ipa={kaːtʃi},pos={名},acc={H}]}%
	{%
		\MeaningsBlock{%
	\ryumeaning%
		{夏至}{}{}%
	}%
	}%
	{\tailPrint[]}%
\entry
	{\headPrint[kana={かーぬ\ws{}すさはぬ},ipa={kaːnu ssahanu},pos={句}]}%
	{%
		\MeaningsBlock{%
	\ryumeaning%
		{匂いが強い}{}{}%
	}%
	}%
	{\tailPrint[]}%
\entry
	{\headPrint[kana={かーま\tl{˥˩}むがし\tl{˩˥}},ipa={kaːmamugaʃi},pos={名},acc={F+R}]}%
	{%
		\MeaningsBlock{%
	\ryumeaning%
		{大昔。太古}{}{}%
	}%
	}%
	{\tailPrint[]}%
\entry
	{\headPrint[kana={かーみん\tl{˥}},ipa={kaːmiɴ},pos={名},acc={H}]}%
	{%
		\MeaningsBlock{%
	\ryumeaning%
		{一重瞼}{}{}%
	}%
	}%
	{\tailPrint[]}%
\entry
	{\headPrint[kana={かーら\tl{˥}},ipa={kaːra},pos={名},acc={H}]}%
	{%
		\MeaningsBlock{%
	\ryumeaning%
		{船の竜骨。キール}{}{}%
	}%
	}%
	{\tailPrint[]}%
\entry
	{\headPrint[kana={かーら\tl{˥}},ipa={kaːra},pos={名},acc={H}]}%
	{%
		\MeaningsBlock{%
	\ryumeaning%
		{瓦}{}{}%
	}%
	}%
	{\tailPrint[]}%
\entry
	{\headPrint[kana={かーら\tl{˥˩}},ipa={kaːra},pos={名},acc={F}]}%
	{%
		\MeaningsBlock{%
	\ryumeaning%
		{川。河川}{}{}%
	}%
	}%
	{\tailPrint[]}%
\entry
	{\headPrint[kana={がーら\tl{˩˥}だま\tl{˥˩}},ipa={gaːradama},pos={名},acc={R+F}]}%
	{%
		\MeaningsBlock{%
	\ryumeaning%
		{\ruby[g]{勾}{まが}\ruby[g]{玉}{たま}}{}{}%
	}%
	}%
	{\tailPrint[]}%
\entry
	{\headPrint[kana={かーら\tl{˥}ひー\tl{˩˥}},ipa={kaːrahiː},pos={名},acc={H+R}]}%
	{%
		\MeaningsBlock{%
	\ryumeaning%
		{瓦葺きの家}{}{}%
	}%
	}%
	{\tailPrint[]}%
\entry
	{\headPrint[kana={かーり\tl{˥˩}},ipa={kaːri},pos={名},acc={F}]}%
	{%
		\MeaningsBlock{%
	\ryumeaning%
		{代わり}{}{}%
	}%
	}%
	{\tailPrint[]}%
\entry
	{\headPrint[kana={がーりすきん\tl{˩˥}},ipa={gaːrisu̥kiɴ},pos={動},rui={3},para={がーりすくぬ},acc={R}]}%
	{%
		\MeaningsBlock{%
	\ryumeaning%
		{威張る}{気勢を上げる}{}%
	}%
	}%
	{\tailPrint[]}%
\entry
	{\headPrint[kana={がーん},ipa={gaːɴ},pos={感}]}%
	{%
		\MeaningsBlock{%
	\ryumeaning%
		{まさか。嘘だ}{強い否定}{}%
	}%
	}%
	{\tailPrint[]}%
\entry
	{\headPrint[kana={がい\tl{˥˩}},ipa={gai},pos={名},acc={F}]}%
	{%
		\MeaningsBlock{%
	\ryumeaning%
		{害}{}{}%
	}%
	}%
	{\tailPrint[]}%
\entry
	{\headPrint[kana={がいじぃ\tl{˩˥}},ipa={gaidʒi},pos={名},acc={R}]}%
	{%
		\MeaningsBlock{%
	\ryumeaning%
		{容器名}{茅と竹ひごで編む物入れ}{}%
	}%
	}%
	{\tailPrint[]}%
\entry
	{\headPrint[kana={がい\tl{˥˩}\ws{}すん\tl{˥˩}},ipa={gai suɴ},pos={句},rui={2},para={がい　さぬ}]}%
	{%
		\MeaningsBlock{%
	\ryumeaning%
		{害する}{}{}%
	}%
	}%
	{\tailPrint[]}%
\entry
	{\headPrint[kana={かぎ\tl{˥}},ipa={kagi},pos={名},acc={H}]}%
	{%
		\MeaningsBlock{%
	\ryumeaning%
		{影。陰}{}{}%
	}%
	}%
	{\tailPrint[]}%
\entry
	{\headPrint[kana={かきあうん\tl{˥}},ipa={kaki.auɴ},pos={動},para={かきあわぬ},acc={H}]}%
	{%
		\MeaningsBlock{%
	\ryumeaning%
		{談判する。抗議する}{}{}%
	}%
	}%
	{\tailPrint[]}%
\entry
	{\headPrint[kana={かぎじん\tl{˥}},ipa={kagidʒiɴ},pos={名},acc={H}]}%
	{%
		\MeaningsBlock{%
	\ryumeaning%
		{陰膳}{旅の人の無事を祈るため}{}%
	}%
	}%
	{\tailPrint[]}%
\entry
	{\headPrint[kana={かぐ\tl{˥}},ipa={kagu},pos={名},acc={H}]}%
	{%
		\MeaningsBlock{%
	\ryumeaning%
		{\ruby[g]{籠}{かご}}{}{}%
	}%
	}%
	{\tailPrint[]}%
\entry
	{\headPrint[kana={がく},ipa={gaku},pos={名},acc={-}]}%
	{%
		\MeaningsBlock{%
	\ryumeaning%
		{学校}{「学校」の略化}{}%
	}%
	}%
	{\tailPrint[]}%
\entry
	{\headPrint[kana={がくたま},ipa={gaku̥tama},pos={名}]}%
	{%
		\MeaningsBlock{%
	\ryumeaning%
		{学童。生徒}{「学童」の方言化}{}%
	}%
	}%
	{\tailPrint[]}%
\entry
	{\headPrint[kana={がくむん\tl{˥}},ipa={gaku̥muɴ},pos={名},acc={H}]}%
	{%
		\MeaningsBlock{%
	\ryumeaning%
		{学問}{}{}%
	}%
	}%
	{\tailPrint[]}%
\entry
	{\headPrint[kana={かさにるん\tl{˥˩}},ipa={kḁsaniruɴ},pos={動},acc={F}]}%
	{%
		\MeaningsBlock{%
	\ryumeaning%
		{重ねる}{}{}%
	}%
	}%
	{\tailPrint[]}%
\entry
	{\headPrint[kana={かさばるん\tl{˥˩}},ipa={kḁsabaruɴ},pos={動},rui={1},para={かつぁばらぬ},acc={F}]}%
	{%
		\MeaningsBlock{%
	\ryumeaning%
		{重なる}{}{}%
	}%
	}%
	{\tailPrint[]}%
\entry
	{\headPrint[kana={かさびるん\tl{˥˩}},ipa={kḁsabiruɴ},pos={動},rui={1},para={かつぁびらぬ},acc={F}]}%
	{%
		\MeaningsBlock{%
	\ryumeaning%
		{重ねる}{}{}%
	}%
	}%
	{\tailPrint[]}%
\entry
	{\headPrint[kana={がざまに\tl{˩˥}},ipa={gadzamani},pos={名},acc={R}]}%
	{%
		\MeaningsBlock{%
	\ryumeaning%
		{ガジュマル}{植物名}{}%
	}%
	}%
	{\tailPrint[]}%
\entry
	{\headPrint[kana={かざまやー\tl{˥}},ipa={kadzamajaː},pos={名},acc={H}]}%
	{%
		\MeaningsBlock{%
	\ryumeaning%
		{風車}{}{}%
	}%
	}%
	{\tailPrint[]}%
\entry
	{\headPrint[kana={かじ\tl{˥}},ipa={kadʒi},pos={名},acc={H}]}%
	{%
		\MeaningsBlock{%
	\ryumeaning%
		{数}{}{}%
	}%
	}%
	{\tailPrint[]}%
\entry
	{\headPrint[kana={かしー\tl{˥}},ipa={kaʃiː},pos={名},acc={H}]}%
	{%
		\MeaningsBlock{%
	\ryumeaning%
		{加勢。手伝い}{}{}%
	}%
	}%
	{\tailPrint[]}%
\entry
	{\headPrint[kana={がしぃどぅしぃ},ipa={gaʃiduʃi},pos={名}]}%
	{%
		\MeaningsBlock{%
	\ryumeaning%
		{餓死の年。凶作の年}{収穫が不作の年}{}%
	}%
	}%
	{\tailPrint[]}%
\entry
	{\headPrint[kana={\josho{}かしがー\kako{}ふくる},ipa={kaʃi̥gaːfu̥kuru},pos={名},acc={H+L}]}%
	{%
		\MeaningsBlock{%
	\ryumeaning%
		{南京袋}{}{}%
	}%
	}%
	{\tailPrint[]}%
\entry
	{\headPrint[kana={かしき\tl{˥}},ipa={kaʃi̥ki},pos={名},acc={H}]}%
	{%
		\MeaningsBlock{%
	\ryumeaning%
		{おこわ。こわ飯}{}{}%
	}%
	}%
	{\tailPrint[]}%
\entry
	{\headPrint[kana={かしきるん\tl{˥}},ipa={kaʃi̥kiruɴ},pos={動},rui={3},para={かしくぬん},acc={H}]}%
	{%
		\MeaningsBlock{%
	\ryumeaning%
		{駆ける。全力で走る}{}{}%
	}%
	}%
	{\tailPrint[]}%
\entry
	{\headPrint[kana={かじまやー\tl{˥}},ipa={kadʒimajaː},pos={名},acc={H}]}%
	{%
		\MeaningsBlock{%
	\ryumeaning%
		{九十七歳の長寿祝い}{}{}%
	}%
	}%
	{\tailPrint[]}%
\entry
	{\headPrint[kana={かしみるん\tl{˥˩}},ipa={kaʃi̥miruɴ},pos={動},rui={3},para={かすむぬ},acc={F}]}%
	{%
		\MeaningsBlock{%
	\ryumeaning%
		{売る。販売する}{}{}%
	}%
	}%
	{\tailPrint[]}%
\entry
	{\headPrint[kana={かしら\tl{˥}},ipa={kaʃira},pos={名},acc={H}]}%
	{%
		\MeaningsBlock{%
	\ryumeaning%
		{頭。頭領}{旗頭}{}%
	}%
	}%
	{\tailPrint[]}%
\entry
	{\headPrint[kana={かずいるん\tl{˥}},ipa={kadzuiruɴ},pos={動},acc={H}]}%
	{%
		\MeaningsBlock{%
	\ryumeaning%
		{数える}{}{}%
	}%
	}%
	{\tailPrint[]}%
\entry
	{\headPrint[kana={\josho{}がす\kako{}た},ipa={gasu̥ta},pos={副},acc={HL}]}%
	{%
		\MeaningsBlock{%
	\ryumeaning%
		{\ruby[g]{皆}{みんな}。全部}{}{}%
	}%
	}%
	{\tailPrint[]}%
\entry
	{\headPrint[kana={かた\tl{˥}},ipa={kḁta},pos={名},acc={H}]}%
	{%
		\MeaningsBlock{%
	\ryumeaning%
		{肩}{}{}%
	}%
	}%
	{\tailPrint[]}%
\entry
	{\headPrint[kana={かた\tl{˥˩}},ipa={kḁta},pos={名},acc={F}]}%
	{%
		\MeaningsBlock{%
	\ryumeaning%
		{型}{}{}%
	}%
	}%
	{\tailPrint[]}%
\entry
	{\headPrint[kana={かたうや\tl{˥}},ipa={kḁta.uja},pos={名},acc={H}]}%
	{%
		\MeaningsBlock{%
	\ryumeaning%
		{片親}{}{}%
	}%
	}%
	{\tailPrint[]}%
\entry
	{\headPrint[kana={かたが\tl{˥}},ipa={kḁtaga},pos={名},acc={H}]}%
	{%
		\MeaningsBlock{%
	\ryumeaning%
		{風除け。庇護}{}{}%
	}%
	}%
	{\tailPrint[]}%
\entry
	{\headPrint[kana={がたがた},ipa={gatagata},pos={擬}]}%
	{%
		\MeaningsBlock{%
	\ryumeaning%
		{がたがた}{震えるさま}{}%
	\ryumeaning%
		{揺れ動くさま}{}{}%
	}%
	}%
	{\tailPrint[]}%
\entry
	{\headPrint[kana={かたぐ\tl{˥}},ipa={kḁtagu},pos={名},acc={H}]}%
	{%
		\MeaningsBlock{%
	\ryumeaning%
		{片方。片一方}{}{}%
	}%
	}%
	{\tailPrint[]}%
\entry
	{\headPrint[kana={かたすか\tl{˥}},ipa={kḁtasu̥ka},pos={名},acc={H}]}%
	{%
		\MeaningsBlock{%
	\ryumeaning%
		{片付け。整理}{}{}%
	}%
	}%
	{\tailPrint[]}%
\entry
	{\headPrint[kana={かたずぐん\tl{˥}},ipa={kḁtadzuguɴ},pos={動},rui={1},para={かたずがぬ},acc={H}]}%
	{%
		\MeaningsBlock{%
	\ryumeaning%
		{きちんと整理される}{}{}%
	}%
	}%
	{\tailPrint[]}%
\entry
	{\headPrint[kana={かたすぷる\tl{˥}やみ\tl{˩˥}},ipa={kḁtasu̥purujami},pos={名},acc={H+R}]}%
	{%
		\MeaningsBlock{%
	\ryumeaning%
		{片頭痛}{}{}%
	}%
	}%
	{\tailPrint[]}%
\entry
	{\headPrint[kana={かたすん\tl{˥˩}},ipa={kḁtasuɴ},pos={動},rui={2},para={かたさぬ},acc={F}]}%
	{%
		\MeaningsBlock{%
	\ryumeaning%
		{組する。味方する}{}{}%
	}%
	}%
	{\tailPrint[]}%
\entry
	{\headPrint[kana={かたち\tl{˥˩}},ipa={kḁtatʃi},pos={名},acc={F}]}%
	{%
		\MeaningsBlock{%
	\ryumeaning%
		{形。姿。恰好}{}{}%
	}%
	}%
	{\tailPrint[]}%
\entry
	{\headPrint[kana={かたな\tl{˥}},ipa={kḁtana},pos={名},acc={H}]}%
	{%
		\MeaningsBlock{%
	\ryumeaning%
		{刀。包丁}{}{}%
	}%
	}%
	{\tailPrint[]}%
\entry
	{\headPrint[kana={かた\tl{˥˩}はん},ipa={kḁtahaɴ},pos={形},para={かたへぬ},acc={F}]}%
	{%
		\MeaningsBlock{%
	\ryumeaning%
		{濃い。濃厚だ。密集している}{}{}%
	}%
	}%
	{\tailPrint[]}%
\entry
	{\headPrint[kana={かたぱん\tl{˥}},ipa={kḁtapaɴ},pos={名},acc={H}]}%
	{%
		\MeaningsBlock{%
	\ryumeaning%
		{片足}{}{}%
	}%
	}%
	{\tailPrint[]}%
\entry
	{\headPrint[kana={かたふかすん},ipa={katafu̥kasuɴ},pos={動}]}%
	{%
		\MeaningsBlock{%
	\ryumeaning%
		{傾ける}{}{}%
	}%
	}%
	{\tailPrint[]}%
\entry
	{\headPrint[kana={かたふきるん\tl{˥}},ipa={katafu̥kiruɴ},pos={動},rui={1},para={かたふかぬ},acc={H}]}%
	{%
		\MeaningsBlock{%
	\ryumeaning%
		{傾ける}{}{}%
	}%
	}%
	{\tailPrint[]}%
\entry
	{\headPrint[kana={かたふくん\tl{˥}},ipa={kḁtafu̥kuɴ},pos={動},rui={1},para={かたふかぬ},acc={H}]}%
	{%
		\MeaningsBlock{%
	\ryumeaning%
		{傾く。傾斜する}{}{}%
	}%
	}%
	{\tailPrint[]}%
\entry
	{\headPrint[kana={かたぶり\tl{˥}},ipa={kḁtaburi},pos={名},acc={H}]}%
	{%
		\MeaningsBlock{%
	\ryumeaning%
		{片降り}{局所的に降る雨}{}%
	}%
	}%
	{\tailPrint[]}%
\entry
	{\headPrint[kana={かたまるん\tl{˥˩}},ipa={kḁtamaruɴ},pos={動},rui={1},para={かたまらぬ},acc={F}]}%
	{%
		\MeaningsBlock{%
	\ryumeaning%
		{固まる。硬くなる}{}{}%
	}%
	}%
	{\tailPrint[]}%
\entry
	{\headPrint[kana={かたみるん\tl{˥}},ipa={kḁtamiruɴ},pos={動},rui={3},para={かたむぬ},acc={H}]}%
	{%
		\MeaningsBlock{%
	\ryumeaning%
		{（肩に）担ぐ}{}{}%
	}%
	}%
	{\tailPrint[]}%
\entry
	{\headPrint[kana={かたみるん\tl{˥˩}},ipa={kḁtamiruɴ},pos={動},rui={3},para={かたむぬ},acc={F}]}%
	{%
		\MeaningsBlock{%
	\ryumeaning%
		{固める}{}{}%
	}%
	}%
	{\tailPrint[]}%
\entry
	{\headPrint[kana={かたるん\tl{˥˩}},ipa={kḁtaruɴ},pos={動},rui={1},para={かたらぬ},acc={F}]}%
	{%
		\MeaningsBlock{%
	\ryumeaning%
		{語る。話す。語らう}{}{}%
	}%
	}%
	{\tailPrint[]}%
\entry
	{\headPrint[kana={かたん\tl{˥˩}},ipa={kḁtaɴ},pos={名},acc={F}]}%
	{%
		\MeaningsBlock{%
	\ryumeaning%
		{バッタ}{昆虫名}{}%
	}%
	}%
	{\tailPrint[]}%
\entry
	{\headPrint[kana={かち\tl{˥}},ipa={kḁtʃi},pos={名},acc={H}]}%
	{%
		\MeaningsBlock{%
	\ryumeaning%
		{すじ。繊維}{}{}%
	}%
	}%
	{\tailPrint[]}%
\entry
	{\headPrint[kana={かち\tl{˥}},ipa={kḁtʃi},pos={名},acc={H}]}%
	{%
		\MeaningsBlock{%
	\ryumeaning%
		{勝ち。勝つこと}{}{}%
	}%
	}%
	{\tailPrint[]}%
\entry
	{\headPrint[kana={かち\tl{˥}},ipa={kḁtʃi},pos={名},acc={H}]}%
	{%
		\MeaningsBlock{%
	\ryumeaning%
		{海栗}{}{}%
	}%
	}%
	{\tailPrint[]}%
\entry
	{\headPrint[kana={かち\tl{˥˩}},ipa={kḁtʃi},pos={名},acc={F}]}%
	{%
		\MeaningsBlock{%
	\ryumeaning%
		{舵}{}{}%
	}%
	}%
	{\tailPrint[]}%
\entry
	{\headPrint[kana={かち\tl{˥˩}},ipa={kḁtʃi},pos={名},acc={F}]}%
	{%
		\MeaningsBlock{%
	\ryumeaning%
		{風}{}{}%
	}%
	}%
	{\tailPrint[]}%
\entry
	{\headPrint[kana={かちぃら\tl{˥˩}},ipa={kḁtʃira},pos={名},acc={F}]}%
	{%
		\MeaningsBlock{%
	\ryumeaning%
		{蔓}{}{}%
	}%
	}%
	{\tailPrint[]}%
\entry
	{\headPrint[kana={かちふき\tl{˥˩}},ipa={kḁtʃifu̥ki},pos={名},acc={F}]}%
	{%
		\MeaningsBlock{%
	\ryumeaning%
		{暴風。台風}{「風吹き」の意味}{}%
	}%
	}%
	{\tailPrint[]}%
\entry
	{\headPrint[kana={かちまーり\tl{˥˩}},ipa={kḁtʃimaːri},pos={名},acc={F}]}%
	{%
		\MeaningsBlock{%
	\ryumeaning%
		{風廻り}{風向の急変のこと}{}%
	}%
	}%
	{\tailPrint[]}%
\entry
	{\headPrint[kana={かちみるん\tl{˥}},ipa={kḁtʃimiɴ},pos={動},rui={3},acc={H}]}%
	{%
		\MeaningsBlock{%
	\ryumeaning%
		{保管する。大事にしまう}{}{}%
	}%
	}%
	{\tailPrint[]}%
\entry
	{\headPrint[kana={かちむん\tl{˥}},ipa={kḁtʃimuɴ},pos={名},acc={H}]}%
	{%
		\MeaningsBlock{%
	\ryumeaning%
		{おかず}{}{}%
	}%
	}%
	{\tailPrint[note={直訳「糧物」から}]}%
\entry
	{\headPrint[kana={かちゅ\tl{˥˩}},ipa={kḁtʃu},pos={名},acc={F}]}%
	{%
		\MeaningsBlock{%
	\ryumeaning%
		{\ruby[g]{鰹}{かつお}}{}{}%
	}%
	}%
	{\tailPrint[]}%
\entry
	{\headPrint[kana={かちゅぶし\tl{˥˩}},ipa={kḁtʃubuʃi},pos={名},acc={F}]}%
	{%
		\MeaningsBlock{%
	\ryumeaning%
		{\ruby[g]{鰹}{かつお}節}{}{}%
	}%
	}%
	{\tailPrint[]}%
\entry
	{\headPrint[kana={かちゅぶに\tl{˥˩}},ipa={kḁtʃubuni},pos={名},acc={F}]}%
	{%
		\MeaningsBlock{%
	\ryumeaning%
		{\ruby[g]{鰹}{かつお}漁船}{}{}%
	}%
	}%
	{\tailPrint[]}%
\entry
	{\headPrint[kana={かちょ\tl{˥˩}はん},ipa={kḁtʃohaɴ},pos={形},para={かちょへぬ},acc={F}]}%
	{%
		\MeaningsBlock{%
	\ryumeaning%
		{風が強い}{}{}%
	}%
	}%
	{\tailPrint[]}%
\entry
	{\headPrint[kana={かちりるん\tl{˥}},ipa={kḁtʃiriruɴ},pos={動},rui={3},para={かつるぬ},acc={H}]}%
	{%
		\MeaningsBlock{%
	\ryumeaning%
		{飢える。餓える}{ひもじい思いをする}{}%
	}%
	}%
	{\tailPrint[]}%
\entry
	{\headPrint[kana={かちりん\tl{˥}},ipa={kḁtʃiriɴ},pos={動},rui={3},para={かつるぬ},acc={H}]}%
	{%
		\MeaningsBlock{%
	\ryumeaning%
		{飢える。餓える}{}{}%
	}%
	}%
	{\tailPrint[]}%
\entry
	{\headPrint[kana={かつぁ\tl{˥}},ipa={kḁtsa},pos={名},acc={H}]}%
	{%
		\MeaningsBlock{%
	\ryumeaning%
		{笠}{}{}%
	}%
	}%
	{\tailPrint[]}%
\entry
	{\headPrint[kana={かつぁ\tl{˥˩}},ipa={kḁtsa},pos={名},acc={F}]}%
	{%
		\MeaningsBlock{%
	\ryumeaning%
		{蚊帳}{}{}%
	}%
	}%
	{\tailPrint[]}%
\entry
	{\headPrint[kana={かつぁなるん\tl{˥˩}},ipa={kḁtsanaruɴ},pos={動},rui={1},para={かつぁならぬ},acc={F}]}%
	{%
		\MeaningsBlock{%
	\ryumeaning%
		{重なる}{}{}%
	}%
	}%
	{\tailPrint[]}%
\entry
	{\headPrint[kana={かつぁびるん\tl{˥˩}},ipa={kḁtsabiruɴ},pos={動},acc={F}]}%
	{%
		\MeaningsBlock{%
	\ryumeaning%
		{重ねる}{}{}%
	}%
	}%
	{\tailPrint[]}%
\entry
	{\headPrint[kana={かつぁま\tl{˥˩}はーん},ipa={kḁtsamahaːɴ},pos={形},acc={F}]}%
	{%
		\MeaningsBlock{%
	\ryumeaning%
		{喧しい}{}{}%
	}%
	}%
	{\tailPrint[]}%
\entry
	{\headPrint[kana={かつぁみるん\tl{˥}},ipa={kḁtsamiruɴ},pos={動},rui={3},para={かつぁむぬ},acc={H}]}%
	{%
		\MeaningsBlock{%
	\ryumeaning%
		{仕舞い込む。隠す}{}{}%
	}%
	}%
	{\tailPrint[]}%
\entry
	{\headPrint[kana={かつぁむん\tl{˥˩}},ipa={kḁtsamuɴ},pos={動},acc={F}]}%
	{%
		\MeaningsBlock{%
	\ryumeaning%
		{嵩む}{数，量が増える}{}%
	}%
	}%
	{\tailPrint[]}%
\entry
	{\headPrint[kana={かつぁり\tl{˥˩}},ipa={kḁtsari},pos={名},acc={F}]}%
	{%
		\MeaningsBlock{%
	\ryumeaning%
		{飾り。装飾}{}{}%
	}%
	}%
	{\tailPrint[]}%
\entry
	{\headPrint[kana={かつぇー\tl{˥˩}},ipa={kḁtseː},pos={名},acc={F}]}%
	{%
		\MeaningsBlock{%
	\ryumeaning%
		{鍛冶。鍛冶屋}{}{}%
	}%
	}%
	{\tailPrint[]}%
\entry
	{\headPrint[kana={かつぇーぶな\tl{˥˩}},ipa={kḁtseːbuna},pos={名},acc={F}]}%
	{%
		\MeaningsBlock{%
	\ryumeaning%
		{鍛冶屋祭。ふいご祭}{}{}%
	}%
	}%
	{\tailPrint[]}%
\entry
	{\headPrint[kana={がっきゃ\tl{˩˥}},ipa={gakkja},pos={名},acc={R}]}%
	{%
		\MeaningsBlock{%
	\ryumeaning%
		{鎌}{}{}%
	}%
	}%
	{\tailPrint[]}%
\entry
	{\headPrint[kana={かつん\tl{˥}},ipa={kḁtsuɴ},pos={動},para={かつぁぬ},acc={H}]}%
	{%
		\MeaningsBlock{%
	\ryumeaning%
		{勝つ}{}{}%
	}%
	}%
	{\tailPrint[]}%
\entry
	{\headPrint[kana={かとーし\tl{˥˩}},ipa={kḁtoːʃi},pos={名},acc={F}]}%
	{%
		\MeaningsBlock{%
	\ryumeaning%
		{梳き櫛}{}{}%
	}%
	}%
	{\tailPrint[]}%
\entry
	{\headPrint[kana={かどぅ\tl{˥}},ipa={kadu},pos={名},acc={H}]}%
	{%
		\MeaningsBlock{%
	\ryumeaning%
		{角。隅}{}{}%
	}%
	}%
	{\tailPrint[]}%
\entry
	{\headPrint[kana={かな\tl{˥}},ipa={kḁna},pos={名},acc={H}]}%
	{%
		\MeaningsBlock{%
	\ryumeaning%
		{鉋}{大工道具}{}%
	}%
	}%
	{\tailPrint[]}%
\entry
	{\headPrint[kana={かな\tl{˥}さーん},ipa={kḁnasaːɴ},pos={形},acc={H}]}%
	{%
		\MeaningsBlock{%
	\ryumeaning%
		{可愛い。愛らしい}{}{}%
	}%
	}%
	{\tailPrint[]}%
\entry
	{\headPrint[kana={かなぱ\tl{˥˩}},ipa={kḁnapa},pos={名},acc={F}]}%
	{%
		\MeaningsBlock{%
	\ryumeaning%
		{大きな葉}{バナナの葉など}{}%
	}%
	}%
	{\tailPrint[]}%
\entry
	{\headPrint[kana={かなばりん\tl{˥}},ipa={kḁnabariɴ},pos={名},acc={H}]}%
	{%
		\MeaningsBlock{%
	\ryumeaning%
		{ヒョウタン。夕顔の実}{}{}%
	}%
	}%
	{\tailPrint[]}%
\entry
	{\headPrint[kana={かなぶ\tl{˥˩}},ipa={kḁnabu},pos={名},acc={F}]}%
	{%
		\MeaningsBlock{%
	\ryumeaning%
		{えびずる。野ブドウ}{}{}%
	}%
	}%
	{\tailPrint[]}%
\entry
	{\headPrint[kana={かに\tl{˥˩}},ipa={kḁni},pos={名},acc={F}]}%
	{%
		\MeaningsBlock{%
	\ryumeaning%
		{金属}{}{}%
	}%
	}%
	{\tailPrint[]}%
\entry
	{\headPrint[kana={かに\tl{˥˩}},ipa={kḁni},pos={名},acc={F}]}%
	{%
		\MeaningsBlock{%
	\ryumeaning%
		{鐘。鉦}{}{}%
	}%
	}%
	{\tailPrint[]}%
\entry
	{\headPrint[kana={かにじん\tl{˥˩}},ipa={kḁnidʒiɴ},pos={名},acc={F}]}%
	{%
		\MeaningsBlock{%
	\ryumeaning%
		{硬貨}{}{}%
	}%
	}%
	{\tailPrint[]}%
\entry
	{\headPrint[kana={かにち\tl{˥˩}},ipa={kḁnitʃi},pos={名},acc={F}]}%
	{%
		\MeaningsBlock{%
	\ryumeaning%
		{金槌。玄翁}{}{}%
	}%
	}%
	{\tailPrint[]}%
\entry
	{\headPrint[kana={かにるん\tl{˥}},ipa={kḁniruɴ},pos={動},rui={3},para={かぬぬ},acc={H}]}%
	{%
		\MeaningsBlock{%
	\ryumeaning%
		{おんぶする。担ぐ}{}{}%
	}%
	}%
	{\tailPrint[]}%
\entry
	{\headPrint[kana={かぬめー\tl{˥}},ipa={kḁnumeː},pos={名},acc={H}]}%
	{%
		\MeaningsBlock{%
	\ryumeaning%
		{神行事}{}{}%
	}%
	}%
	{\tailPrint[]}%
\entry
	{\headPrint[kana={かねー},ipa={kḁneː},pos={-},acc={-}]}%
	{%
		\MeaningsBlock{%
	\ryumeaning%
		{すぐれた。有能な}{}{}%
	}%
	}%
	{\tailPrint[]}%
\entry
	{\headPrint[kana={かねーるん\tl{˥}},ipa={kḁneruɴ},pos={動},rui={1},para={かねーらぬ},acc={H}]}%
	{%
		\MeaningsBlock{%
	\ryumeaning%
		{叶える}{}{}%
	}%
	}%
	{\tailPrint[]}%
\entry
	{\headPrint[kana={かのーん\tl{˥}},ipa={kḁnoːɴ},pos={動},rui={1},para={かのーらぬ},acc={H}]}%
	{%
		\MeaningsBlock{%
	\ryumeaning%
		{叶う。実現する}{}{}%
	}%
	}%
	{\tailPrint[]}%
\entry
	{\headPrint[kana={かぱすん\tl{˥˩}},ipa={kḁpasuɴ},pos={動},rui={2},para={かぱはぬ},acc={F}]}%
	{%
		\MeaningsBlock{%
	\ryumeaning%
		{嗅がせる}{}{}%
	}%
	}%
	{\tailPrint[]}%
\entry
	{\headPrint[kana={かぱ\tl{˥˩}はん},ipa={kḁpahaɴ},pos={形},acc={F}]}%
	{%
		\MeaningsBlock{%
	\ryumeaning%
		{芳しい。薫る}{良い香りがする}{}%
	}%
	}%
	{\tailPrint[]}%
\entry
	{\headPrint[kana={かふちるん\tl{˥}},ipa={kafu̥tʃiruɴ},pos={動},rui={3},para={かふつぬ},acc={H}]}%
	{%
		\MeaningsBlock{%
	\ryumeaning%
		{被せる。上からかける。覆う}{}{}%
	}%
	}%
	{\tailPrint[]}%
\entry
	{\headPrint[kana={かぷる\tl{˥}},ipa={kḁpuru},pos={名},acc={H}]}%
	{%
		\MeaningsBlock{%
	\ryumeaning%
		{コウモリ}{}{}%
	}%
	}%
	{\tailPrint[]}%
\entry
	{\headPrint[kana={かぷるさな\tl{˥}},ipa={kḁpurusḁna},pos={名},acc={H}]}%
	{%
		\MeaningsBlock{%
	\ryumeaning%
		{洋傘。蘭傘}{}{}%
	}%
	}%
	{\tailPrint[]}%
\entry
	{\headPrint[kana={かぷん\tl{˥}},ipa={kḁpuɴ},pos={動},rui={1},para={かぱぬ},acc={H}]}%
	{%
		\MeaningsBlock{%
	\ryumeaning%
		{被る}{}{}%
	\ryumeaning%
		{茂る。生い茂る。繁茂する}{}{}%
	}%
	}%
	{\tailPrint[]}%
\entry
	{\headPrint[kana={かぷん\tl{˥˩}},ipa={kḁpuɴ},pos={動},rui={1},para={かぱぬ},acc={F}]}%
	{%
		\MeaningsBlock{%
	\ryumeaning%
		{（匂いを）嗅ぐ}{}{}%
	}%
	}%
	{\tailPrint[]}%
\entry
	{\headPrint[kana={かぺー\tl{˥}},ipa={kḁpeː},pos={名},acc={H}]}%
	{%
		\MeaningsBlock{%
	\ryumeaning%
		{\ruby[g]{黴}{かび}}{}{}%
	}%
	}%
	{\tailPrint[]}%
\entry
	{\headPrint[kana={かま\tl{˥˩}},ipa={kḁma},pos={名},acc={F}]}%
	{%
		\MeaningsBlock{%
	\ryumeaning%
		{\ruby[g]{釜}{かま}}{}{}%
	}%
	}%
	{\tailPrint[]}%
\entry
	{\headPrint[kana={かまいるん\tl{˥}},ipa={kḁmairuɴ},pos={動},rui={1},para={かまわぬ},acc={H}]}%
	{%
		\MeaningsBlock{%
	\ryumeaning%
		{構える。身構える}{}{}%
	}%
	}%
	{\tailPrint[]}%
\entry
	{\headPrint[kana={かまぐ\tl{˥˩}},ipa={kḁmagu},pos={名},acc={F}]}%
	{%
		\MeaningsBlock{%
	\ryumeaning%
		{かまぼこ}{}{}%
	}%
	}%
	{\tailPrint[]}%
\entry
	{\headPrint[kana={かまさ\tl{˥}},ipa={kḁmasa},pos={名},acc={H}]}%
	{%
		\MeaningsBlock{%
	\ryumeaning%
		{食器}{ひょうたん製の野良用}{}%
	}%
	}%
	{\tailPrint[]}%
\entry
	{\headPrint[kana={がま\tl{˩˥}さーん},ipa={gamasaːɴ},pos={形},acc={R}]}%
	{%
		\MeaningsBlock{%
	\ryumeaning%
		{腕白だ}{}{}%
	\ryumeaning%
		{悪戯っぽい}{}{}%
	}%
	}%
	{\tailPrint[]}%
\entry
	{\headPrint[kana={がまん\tl{˥˩}},ipa={gamaɴ},pos={名},acc={F}]}%
	{%
		\MeaningsBlock{%
	\ryumeaning%
		{（海中の）洞穴}{陸上の洞窟は\emb{いん}}{}%
	}%
	}%
	{\tailPrint[]}%
\entry
	{\headPrint[kana={かまんた\tl{˥}},ipa={kamanta},pos={名},acc={H}]}%
	{%
		\MeaningsBlock{%
	\ryumeaning%
		{大型エイ。マンタ}{}{}%
	}%
	}%
	{\tailPrint[]}%
\entry
	{\headPrint[kana={かみ\tl{˥}},ipa={kḁmi},pos={名},acc={H}]}%
	{%
		\MeaningsBlock{%
	\ryumeaning%
		{亀}{}{}%
	}%
	}%
	{\tailPrint[]}%
\entry
	{\headPrint[kana={かみ\tl{˥}},ipa={kḁmi},pos={名},acc={H}]}%
	{%
		\MeaningsBlock{%
	\ryumeaning%
		{甕}{}{}%
	}%
	}%
	{\tailPrint[]}%
\entry
	{\headPrint[kana={かみしきるん\tl{˥˩}},ipa={kḁmiʃi̥kiruɴ},pos={動},rui={3},para={かみすくぬ},acc={F}]}%
	{%
		\MeaningsBlock{%
	\ryumeaning%
		{捕える。ひっつかむ}{}{}%
	}%
	}%
	{\tailPrint[]}%
\entry
	{\headPrint[kana={かみるん\tl{˥}},ipa={kḁmiruɴ},pos={動},rui={3},acc={H}]}%
	{%
		\MeaningsBlock{%
	\ryumeaning%
		{\ruby[g]{掴}{つか}む。捕える}{}{}%
	}%
	}%
	{\tailPrint[]}%
\entry
	{\headPrint[kana={かみん\tl{˥}},ipa={kḁmiɴ},pos={動},rui={3}]}%
	{%
		\MeaningsBlock{%
	\ryumeaning%
		{\ruby[g]{掴}{つか}む。捕える}{}{}%
	}%
	}%
	{\tailPrint[]}%
\entry
	{\headPrint[kana={かみん\tl{˥}},ipa={kḁmiɴ},pos={動},rui={3},para={かむぬ},acc={H}]}%
	{%
		\MeaningsBlock{%
	\ryumeaning%
		{頭上に乗せて運ぶ}{}{}%
	}%
	}%
	{\tailPrint[]}%
\entry
	{\headPrint[kana={かむい\tl{˥}},ipa={kḁmui},pos={名},acc={H}]}%
	{%
		\MeaningsBlock{%
	\ryumeaning%
		{梁}{}{}%
	}%
	}%
	{\tailPrint[]}%
\entry
	{\headPrint[kana={かむし\tl{˥}},ipa={kḁmuʃi},pos={名},acc={H}]}%
	{%
		\MeaningsBlock{%
	\ryumeaning%
		{ゴキブリ。油虫}{}{}%
	}%
	}%
	{\tailPrint[]}%
\entry
	{\headPrint[kana={かむん\tl{˥}},ipa={kḁmuɴ},pos={動},rui={1},para={かまぬ},acc={H}]}%
	{%
		\MeaningsBlock{%
	\ryumeaning%
		{噛む。咬む}{}{}%
	}%
	}%
	{\tailPrint[]}%
\entry
	{\headPrint[kana={がや\tl{˩˥}},ipa={gaja},pos={名},acc={R}]}%
	{%
		\MeaningsBlock{%
	\ryumeaning%
		{茅。チガヤ}{植物名}{}%
	}%
	}%
	{\tailPrint[]}%
\entry
	{\headPrint[kana={がやふきひー\tl{˩˥}},ipa={gajafu̥kihiː},pos={名},acc={R}]}%
	{%
		\MeaningsBlock{%
	\ryumeaning%
		{萱ぶきの家屋}{}{}%
	}%
	}%
	{\tailPrint[]}%
\entry
	{\headPrint[kana={かよーん\tl{˥˩}},ipa={kajoːɴ},pos={動},para={かよーわぬ},acc={F}]}%
	{%
		\MeaningsBlock{%
	\ryumeaning%
		{通う}{}{}%
	}%
	}%
	{\tailPrint[]}%
\entry
	{\headPrint[kana={がら},ipa={gara},pos={接尾}]}%
	{%
		\MeaningsBlock{%
	\ryumeaning%
		{～匹}{魚などを数える数詞。結合する数詞語根によって形が変わる。\emb{ぴとーら}「一匹」、\emb{みーから}「三匹」など}{}%
	}%
	}%
	{\tailPrint[]}%
\entry
	{\headPrint[kana={がらがら},ipa={garagara},pos={擬}]}%
	{%
		\MeaningsBlock{%
	\ryumeaning%
		{からから}{空っぽのさま}{}%
	}%
	}%
	{\tailPrint[]}%
\entry
	{\headPrint[kana={からぎるん\tl{˥}},ipa={kḁragiruɴ},pos={動},rui={3},para={からぎらぬ},acc={H}]}%
	{%
		\MeaningsBlock{%
	\ryumeaning%
		{絡げる。\ruby[g]{捲}{まく}る}{}{}%
	}%
	}%
	{\tailPrint[]}%
\entry
	{\headPrint[kana={がらし\tl{˩˥}},ipa={garaʃi},pos={名},acc={R}]}%
	{%
		\MeaningsBlock{%
	\ryumeaning%
		{カラス}{鳥の名}{}%
	}%
	}%
	{\tailPrint[]}%
\entry
	{\headPrint[kana={からすん\tl{˥˩}},ipa={kḁrasuɴ},pos={動},rui={2},para={からはぬ},acc={F}]}%
	{%
		\MeaningsBlock{%
	\ryumeaning%
		{貸す}{}{}%
	}%
	}%
	{\tailPrint[]}%
\entry
	{\headPrint[kana={から\tl{˥˩}\ws{}なるん\tl{˩˥}},ipa={kara naruɴ},pos={句},rui={1},para={から　ならぬ},acc={F R}]}%
	{%
		\MeaningsBlock{%
	\ryumeaning%
		{空になる}{全部移されてしまう}{}%
	}%
	}%
	{\tailPrint[]}%
\entry
	{\headPrint[kana={からばっさ\tl{˥}はん},ipa={kḁrabassahaɴ},pos={形},para={からばっさへぬ},acc={H}]}%
	{%
		\MeaningsBlock{%
	\ryumeaning%
		{すばしこい。身軽だ。敏捷だ}{}{}%
	}%
	}%
	{\tailPrint[]}%
\entry
	{\headPrint[kana={から\tl{˥}はん},ipa={kḁrahaɴ},pos={形},para={からへぬ},acc={H}]}%
	{%
		\MeaningsBlock{%
	\ryumeaning%
		{\ruby[g]{辛}{から}い}{}{}%
	}%
	}%
	{\tailPrint[]}%
\entry
	{\headPrint[kana={からぱん\tl{˥˩}},ipa={kḁrapaɴ},pos={名},acc={F}]}%
	{%
		\MeaningsBlock{%
	\ryumeaning%
		{\ruby[g]{裸足}{はだし}}{}{}%
	}%
	}%
	{\tailPrint[]}%
\entry
	{\headPrint[kana={からまぐん\tl{˥}},ipa={kḁramaguɴ},pos={動},rui={1},para={からまがぬ},acc={H}]}%
	{%
		\MeaningsBlock{%
	\ryumeaning%
		{絡まる。絡み付く}{}{}%
	}%
	}%
	{\tailPrint[]}%
\entry
	{\headPrint[kana={からまぐん\tl{˥}},ipa={kḁramaguɴ},pos={動},rui={1},para={からまがぬ},acc={H}]}%
	{%
		\MeaningsBlock{%
	\ryumeaning%
		{巻き付ける。絡める}{}{}%
	}%
	}%
	{\tailPrint[]}%
\entry
	{\headPrint[kana={からむん\tl{˥}},ipa={kḁramuɴ},pos={動},rui={1},para={からまぬ},acc={H}]}%
	{%
		\MeaningsBlock{%
	\ryumeaning%
		{絡み付く}{}{}%
	}%
	}%
	{\tailPrint[]}%
\entry
	{\headPrint[kana={かりー\tl{˥}},ipa={kḁriː},pos={名},acc={H}]}%
	{%
		\MeaningsBlock{%
	\ryumeaning%
		{嘉例。吉例}{}{}%
	}%
	}%
	{\tailPrint[]}%
\entry
	{\headPrint[kana={かりがすん\tl{˥˩}},ipa={kḁrigasuɴ},pos={動},rui={2},para={かりがはぬ},acc={F}]}%
	{%
		\MeaningsBlock{%
	\ryumeaning%
		{乾かす}{}{}%
	}%
	}%
	{\tailPrint[]}%
\entry
	{\headPrint[kana={かりぐん\tl{˥˩}},ipa={kariguɴ},pos={動},rui={1},para={かりがぬ},acc={F}]}%
	{%
		\MeaningsBlock{%
	\ryumeaning%
		{乾く}{}{}%
	}%
	}%
	{\tailPrint[]}%
\entry
	{\headPrint[kana={かりゆし\tl{˥}},ipa={karijuʃi},pos={名},acc={H}]}%
	{%
		\MeaningsBlock{%
	\ryumeaning%
		{航海安全。航路平安}{}{}%
	}%
	}%
	{\tailPrint[]}%
\entry
	{\headPrint[kana={かりょん\tl{˥}},ipa={kḁrjoɴ},pos={名},acc={H}]}%
	{%
		\MeaningsBlock{%
	\ryumeaning%
		{山芋}{作物}{}%
	}%
	}%
	{\tailPrint[]}%
\entry
	{\headPrint[kana={かりるん\tl{˥˩}},ipa={kḁriruɴ},pos={動},rui={1},para={からぬ},acc={F}]}%
	{%
		\MeaningsBlock{%
	\ryumeaning%
		{借りる}{}{}%
	}%
	}%
	{\tailPrint[]}%
\entry
	{\headPrint[kana={かりるん\tl{˥˩}},ipa={kḁriruɴ},pos={動},rui={3},para={かるぬ},acc={F}]}%
	{%
		\MeaningsBlock{%
	\ryumeaning%
		{枯れる}{}{}%
	\ryumeaning%
		{（声が）嗄れる}{}{}%
	}%
	}%
	{\tailPrint[]}%
\entry
	{\headPrint[kana={かるん\tl{˥˩}},ipa={kḁruɴ},pos={動},rui={1},para={からぬ},acc={F}]}%
	{%
		\MeaningsBlock{%
	\ryumeaning%
		{刈る}{刈り取る}{}%
	}%
	}%
	{\tailPrint[]}%
\entry
	{\headPrint[kana={かるん\tl{˥˩}},ipa={kḁruɴ},pos={動},rui={1},para={からぬ},acc={F}]}%
	{%
		\MeaningsBlock{%
	\ryumeaning%
		{借りる。借用する}{}{}%
	}%
	}%
	{\tailPrint[]}%
\entry
	{\headPrint[kana={かるんじるん\tl{˥˩}},ipa={kḁrundʒiruɴ},pos={動},rui={3},acc={F}]}%
	{%
		\MeaningsBlock{%
	\ryumeaning%
		{軽んじる。軽く思う。ないがしろにする}{}{}%
	}%
	}%
	{\tailPrint[]}%
\entry
	{\headPrint[kana={かれーるん\tl{˥˩}},ipa={kḁreːruɴ},pos={動},rui={1},para={かれーらぬ},acc={F}]}%
	{%
		\MeaningsBlock{%
	\ryumeaning%
		{変える。代える。替える。交換する。両替する}{}{}%
	}%
	}%
	{\tailPrint[]}%
\entry
	{\headPrint[kana={かろ\tl{˥˩}はん},ipa={kḁrohaɴ},pos={形},para={かろへぬ},acc={F}]}%
	{%
		\MeaningsBlock{%
	\ryumeaning%
		{軽い}{}{}%
	}%
	}%
	{\tailPrint[]}%
\entry
	{\headPrint[kana={かわり},ipa={kawari},pos={名}]}%
	{%
		\MeaningsBlock{%
	\ryumeaning%
		{変わり}{}{}%
	}%
	}%
	{\tailPrint[]}%
\entry
	{\headPrint[kana={かわるん\tl{˥˩}},ipa={kawaruɴ},pos={動},rui={1},para={かわらぬ},acc={F}]}%
	{%
		\MeaningsBlock{%
	\ryumeaning%
		{変わる。変化する}{}{}%
	\ryumeaning%
		{代わる}{}{}%
	}%
	}%
	{\tailPrint[]}%
\entry
	{\headPrint[kana={かん\tl{˥}},ipa={kaɴ},pos={名},acc={H}]}%
	{%
		\MeaningsBlock{%
	\ryumeaning%
		{神。神様}{}{}%
	}%
	}%
	{\tailPrint[]}%
\entry
	{\headPrint[kana={かん\tl{˥}},ipa={kaɴ},pos={名},acc={H}]}%
	{%
		\MeaningsBlock{%
	\ryumeaning%
		{寒気。寒さ}{}{}%
	}%
	}%
	{\tailPrint[]}%
\entry
	{\headPrint[kana={かん\tl{˥˩}},ipa={kaɴ},pos={名},acc={F}]}%
	{%
		\MeaningsBlock{%
	\ryumeaning%
		{蟹}{}{}%
	}%
	}%
	{\tailPrint[]}%
\entry
	{\headPrint[kana={かん\tl{˥˩}},ipa={kaɴ},pos={名},acc={F}]}%
	{%
		\MeaningsBlock{%
	\ryumeaning%
		{棺桶}{}{}%
	}%
	}%
	{\tailPrint[]}%
\entry
	{\headPrint[kana={かんがーすん\tl{˥}},ipa={kaŋgaːsuɴ},pos={動},rui={2},para={かんがーはぬ},acc={H}]}%
	{%
		\MeaningsBlock{%
	\ryumeaning%
		{（火で）焙って乾燥する}{}{}%
	}%
	}%
	{\tailPrint[]}%
\entry
	{\headPrint[kana={かんがん\tl{˥}},ipa={kaŋgaɴ},pos={名},acc={H}]}%
	{%
		\MeaningsBlock{%
	\ryumeaning%
		{鏡}{}{}%
	}%
	}%
	{\tailPrint[]}%
\entry
	{\headPrint[kana={がんがん},ipa={gaŋgaɴ},pos={擬}]}%
	{%
		\MeaningsBlock{%
	\ryumeaning%
		{早く}{早くやれと急がす様}{}%
	}%
	}%
	{\tailPrint[]}%
\entry
	{\headPrint[kana={がんきょー\tl{˥}},ipa={gaŋkjoː},pos={名},acc={H}]}%
	{%
		\MeaningsBlock{%
	\ryumeaning%
		{眼鏡。水中眼鏡}{}{}%
	}%
	}%
	{\tailPrint[]}%
\entry
	{\headPrint[kana={がんくむぬ},ipa={gankumunu},pos={名}]}%
	{%
		\MeaningsBlock{%
	\ryumeaning%
		{頑固者}{}{}%
	}%
	}%
	{\tailPrint[]}%
\entry
	{\headPrint[kana={かんげー\tl{˥}},ipa={kaŋgeː},pos={名},acc={H}]}%
	{%
		\MeaningsBlock{%
	\ryumeaning%
		{考え。思案}{}{}%
	}%
	}%
	{\tailPrint[]}%
\entry
	{\headPrint[kana={かんげーるん\tl{˥}},ipa={kaŋgeːruɴ},pos={動},rui={3},acc={H}]}%
	{%
		\MeaningsBlock{%
	\ryumeaning%
		{考える。思う}{}{}%
	}%
	}%
	{\tailPrint[]}%
\entry
	{\headPrint[kana={かんご\tl{˥}},ipa={kaŋgo},pos={名},acc={H}]}%
	{%
		\MeaningsBlock{%
	\ryumeaning%
		{肩車}{}{}%
	}%
	}%
	{\tailPrint[]}%
\entry
	{\headPrint[kana={がんざん\tl{˩˥}},ipa={gandzaɴ},pos={名},acc={R}]}%
	{%
		\MeaningsBlock{%
	\ryumeaning%
		{蚊}{}{}%
	}%
	}%
	{\tailPrint[]}%
\entry
	{\headPrint[kana={がんす\tl{˩˥}},ipa={gansu},pos={名},acc={R}]}%
	{%
		\MeaningsBlock{%
	\ryumeaning%
		{元祖。先祖}{}{}%
	}%
	}%
	{\tailPrint[]}%
\entry
	{\headPrint[kana={がんずさ\tl{˩˥}\ws{}なるん\tl{˩˥}},ipa={gandzusa naruɴ},pos={句},acc={R R}]}%
	{%
		\MeaningsBlock{%
	\ryumeaning%
		{丈夫になる。体が強くなる}{}{}%
	}%
	}%
	{\tailPrint[]}%
\entry
	{\headPrint[kana={がんずさ\tl{˩˥}はん},ipa={gandzusahaɴ},pos={形},para={がんずさへぬ},acc={R}]}%
	{%
		\MeaningsBlock{%
	\ryumeaning%
		{頑丈だ。丈夫だ。元気だ}{}{}%
	}%
	}%
	{\tailPrint[]}%
\entry
	{\headPrint[kana={かんすり\tl{˥}},ipa={kansu̥ri},pos={名},acc={H}]}%
	{%
		\MeaningsBlock{%
	\ryumeaning%
		{剃刀}{}{}%
	}%
	}%
	{\tailPrint[]}%
\entry
	{\headPrint[kana={かんだが\tl{˥}さーん},ipa={kandagasaːɴ},pos={形},acc={H}]}%
	{%
		\MeaningsBlock{%
	\ryumeaning%
		{神の霊験が高い}{}{}%
	}%
	}%
	{\tailPrint[]}%
\entry
	{\headPrint[kana={がんだるごー\tl{˥˩}},ipa={gandarugoː},pos={名},acc={F}]}%
	{%
		\MeaningsBlock{%
	\ryumeaning%
		{\ruby[g]{龕}{がん}}{葬儀で棺を運ぶ龕}{}%
	}%
	}%
	{\tailPrint[]}%
\entry
	{\headPrint[kana={かんだるん\tl{˥}},ipa={kandaruɴ},pos={動},rui={1},para={かんだらぬ},acc={H}]}%
	{%
		\MeaningsBlock{%
	\ryumeaning%
		{噛む。よく噛む}{}{}%
	}%
	}%
	{\tailPrint[]}%
\entry
	{\headPrint[kana={がんどーりん\tl{˩˥}},ipa={gandoːriɴ},pos={動},rui={3},para={がんどうーるぬ},acc={R}]}%
	{%
		\MeaningsBlock{%
	\ryumeaning%
		{疲れ果てる}{}{}%
	\ryumeaning%
		{元気をなくす}{}{}%
	}%
	}%
	{\tailPrint[]}%
\entry
	{\headPrint[kana={がんどらん\tl{˩˥}},ipa={gandoraɴ},pos={動（継）},rui={1},para={がんどらぬ},acc={R}]}%
	{%
		\MeaningsBlock{%
	\ryumeaning%
		{痩せ衰える}{}{}%
	}%
	}%
	{\tailPrint[]}%
\entry
	{\headPrint[kana={かんなり\tl{˥}},ipa={kannari},pos={名},acc={H}]}%
	{%
		\MeaningsBlock{%
	\ryumeaning%
		{雷。雷鳴}{}{}%
	}%
	}%
	{\tailPrint[]}%
\entry
	{\headPrint[kana={かんぬ\tl{˥}\ws{}ゆー\tl{˥˩}},ipa={kannu juː},pos={句},acc={H F}]}%
	{%
		\MeaningsBlock{%
	\ryumeaning%
		{神の世。豊穣の世}{}{}%
	}%
	}%
	{\tailPrint[]}%
\entry
	{\headPrint[kana={かんぱるん\tl{˥}},ipa={kampḁruɴ},pos={動},rui={1},para={かんぱらーぬ},acc={H}]}%
	{%
		\MeaningsBlock{%
	\ryumeaning%
		{噛み砕く}{}{}%
	}%
	}%
	{\tailPrint[]}%
\entry
	{\headPrint[kana={かんぽーしん\tl{˥}},ipa={kampoːʃiɴ},pos={名},acc={H}]}%
	{%
		\MeaningsBlock{%
	\ryumeaning%
		{船首材が上に出た船}{}{}%
	}%
	}%
	{\tailPrint[]}%
\entry
	{\headPrint[kana={がんまり\tl{˩˥}},ipa={gammari},pos={名},acc={R}]}%
	{%
		\MeaningsBlock{%
	\ryumeaning%
		{悪戯。悪さ}{}{}%
	}%
	}%
	{\tailPrint[]}%
\entry
	{\headPrint[kana={がんまり\ws{}しゃーん},ipa={gammari ʃaːɴ},pos={句},acc={R HL}]}%
	{%
		\MeaningsBlock{%
	\ryumeaning%
		{よく悪戯をする。悪戯っぽい}{}{}%
	}%
	}%
	{\tailPrint[]}%
\entry
	{\headPrint[kana={かんみちぃ\tl{˥}},ipa={kammitʃi},pos={名},acc={H}]}%
	{%
		\MeaningsBlock{%
	\ryumeaning%
		{神の道}{神や神司が通る道}{}%
	}%
	}%
	{\tailPrint[]}%
\LetterHead{き}
\entry
	{\headPrint[kana={きー\tl{˥}},ipa={kiː},pos={名},acc={H}]}%
	{%
		\MeaningsBlock{%
	\ryumeaning%
		{木。樹木}{}{}%
	}%
	}%
	{\tailPrint[]}%
\entry
	{\headPrint[kana={きー\tl{˥˩}},ipa={kiː},pos={名},acc={F}]}%
	{%
		\MeaningsBlock{%
	\ryumeaning%
		{毛。体毛}{}{}%
	}%
	}%
	{\tailPrint[]}%
\entry
	{\headPrint[kana={きー\tl{˥˩}},ipa={kiː},pos={名},acc={F}]}%
	{%
		\MeaningsBlock{%
	\ryumeaning%
		{気。気持}{}{}%
	}%
	}%
	{\tailPrint[]}%
\entry
	{\headPrint[kana={きーいりん\tl{˥˩}},ipa={kiːiriɴ},pos={動},rui={1},para={きーいらぬ},acc={F}]}%
	{%
		\MeaningsBlock{%
	\ryumeaning%
		{元気付ける}{}{}%
	}%
	}%
	{\tailPrint[]}%
\entry
	{\headPrint[kana={きー\tl{˥˩}\ws{}しきるん\tl{˥}},ipa={kiː ʃi̥kiruɴ},pos={句},rui={3},para={きぃすくぬ},acc={F H}]}%
	{%
		\MeaningsBlock{%
	\ryumeaning%
		{注意する。気をつける}{}{}%
	}%
	}%
	{\tailPrint[]}%
\entry
	{\headPrint[kana={きーぬ\ws{}むとぅ},ipa={kiːnu mutu},pos={句}]}%
	{%
		\MeaningsBlock{%
	\ryumeaning%
		{木の幹}{}{}%
	}%
	}%
	{\tailPrint[]}%
\entry
	{\headPrint[kana={きーる\tl{˥˩}},ipa={kiːru},pos={名},acc={F}]}%
	{%
		\MeaningsBlock{%
	\ryumeaning%
		{黄色}{}{}%
	}%
	}%
	{\tailPrint[]}%
\entry
	{\headPrint[kana={きがい\tl{˥˩}},ipa={kigai},pos={名},acc={F}]}%
	{%
		\MeaningsBlock{%
	\ryumeaning%
		{気立て。気骨}{}{}%
	}%
	}%
	{\tailPrint[]}%
\entry
	{\headPrint[kana={きかんしん\tl{˥}},ipa={ki̥kanʃiɴ},pos={名},acc={H}]}%
	{%
		\MeaningsBlock{%
	\ryumeaning%
		{発動機船}{帆船に対して言う}{}%
	}%
	}%
	{\tailPrint[]}%
\entry
	{\headPrint[kana={きさ\tl{˥}},ipa={ki̥sa},pos={名},acc={H}]}%
	{%
		\MeaningsBlock{%
	\ryumeaning%
		{既に。とっくに}{}{}%
	}%
	}%
	{\tailPrint[]}%
\entry
	{\headPrint[kana={きさり\tl{˥˩}},ipa={ki̥sari},pos={名},acc={F}]}%
	{%
		\MeaningsBlock{%
	\ryumeaning%
		{神行事}{宗教的な行事を指す}{}%
	}%
	}%
	{\tailPrint[]}%
\entry
	{\headPrint[kana={ぎし\tl{˩˥}},ipa={giʃi},pos={名},acc={R}]}%
	{%
		\MeaningsBlock{%
	\ryumeaning%
		{指図。命令}{}{}%
	}%
	}%
	{\tailPrint[]}%
\entry
	{\headPrint[kana={きしり\tl{˥}},ipa={ki̥ʃiri},pos={名},acc={H}]}%
	{%
		\MeaningsBlock{%
	\ryumeaning%
		{煙管}{}{}%
	}%
	}%
	{\tailPrint[]}%
\entry
	{\headPrint[kana={きた\tl{˥}},ipa={ki̥ta},pos={名},acc={H}]}%
	{%
		\MeaningsBlock{%
	\ryumeaning%
		{桁}{}{}%
	}%
	}%
	{\tailPrint[]}%
\entry
	{\headPrint[kana={きちぃ\tl{˥}},ipa={kitʃi},pos={名},acc={H}]}%
	{%
		\MeaningsBlock{%
	\ryumeaning%
		{垂木}{}{}%
	}%
	}%
	{\tailPrint[]}%
\entry
	{\headPrint[kana={きちぃるん\tl{˥˩}},ipa={ki̥tʃiruɴ},pos={動},rui={1},para={きちぃらぬ},acc={F}]}%
	{%
		\MeaningsBlock{%
	\ryumeaning%
		{\ruby[g]{削}{けず}る}{}{}%
	\ryumeaning%
		{（櫛で髪を）梳く}{}{}%
	}%
	}%
	{\tailPrint[]}%
\entry
	{\headPrint[kana={きちがん\tl{˥˩}},ipa={ki̥tʃigaɴ},pos={名},acc={F}]}%
	{%
		\MeaningsBlock{%
	\ryumeaning%
		{結願}{祈願の成就への感謝祭}{}%
	}%
	}%
	{\tailPrint[]}%
\entry
	{\headPrint[kana={きつぁむん\tl{˥˩}},ipa={ki̥tsamuɴ},pos={動},rui={1},para={きつぁまぬ},acc={F}]}%
	{%
		\MeaningsBlock{%
	\ryumeaning%
		{刻む}{}{}%
	}%
	}%
	{\tailPrint[]}%
\entry
	{\headPrint[kana={きっく\tl{˥}},ipa={kikku},pos={名},acc={H}]}%
	{%
		\MeaningsBlock{%
	\ryumeaning%
		{稽古。練習}{}{}%
	}%
	}%
	{\tailPrint[]}%
\entry
	{\headPrint[kana={きっしい\tl{˥˩}},ipa={kiʃʃi},pos={名},acc={F}]}%
	{%
		\MeaningsBlock{%
	\ryumeaning%
		{警察}{「警察」の転訛}{}%
	}%
	}%
	{\tailPrint[]}%
\entry
	{\headPrint[kana={きっとー\tl{˥}},ipa={kittoː},pos={名},acc={H}]}%
	{%
		\MeaningsBlock{%
	\ryumeaning%
		{毛布}{外来語の「ケット」から}{}%
	}%
	}%
	{\tailPrint[]}%
\entry
	{\headPrint[kana={きな\tl{˥˩}},ipa={ki̥na},pos={名},acc={F}]}%
	{%
		\MeaningsBlock{%
	\ryumeaning%
		{黒木。琉球黒檀}{樹木で三線の部材になる}{}%
	}%
	}%
	{\tailPrint[]}%
\entry
	{\headPrint[kana={きなし\tl{˥}},ipa={ki̥naʃi},pos={名},acc={H}]}%
	{%
		\MeaningsBlock{%
	\ryumeaning%
		{シナノキ}{植物名}{}%
	}%
	}%
	{\tailPrint[]}%
\entry
	{\headPrint[kana={ぎなむぬ\tl{˩˥}},ipa={ginamunu},pos={名},acc={R}]}%
	{%
		\MeaningsBlock{%
	\ryumeaning%
		{出しもの。芸能}{「芸の物」から}{}%
	}%
	}%
	{\tailPrint[]}%
\entry
	{\headPrint[kana={きに\tl{˥˩}},ipa={ki̥ni},pos={名},acc={F}]}%
	{%
		\MeaningsBlock{%
	\ryumeaning%
		{世帯。家庭}{}{}%
	}%
	}%
	{\tailPrint[]}%
\entry
	{\headPrint[kana={きにばがら\tl{˩˥}},ipa={ki̥nibagara},pos={名},acc={R}]}%
	{%
		\MeaningsBlock{%
	\ryumeaning%
		{分家する}{}{}%
	\ryumeaning%
		{分家}{}{}%
	}%
	}%
	{\tailPrint[]}%
\entry
	{\headPrint[kana={きぬ\tl{˥˩}\ws{}あん\tl{˥}},ipa={kinu aɴ},pos={句},acc={F H}]}%
	{%
		\MeaningsBlock{%
	\ryumeaning%
		{気がある。関心がある}{}{}%
	}%
	}%
	{\tailPrint[]}%
\entry
	{\headPrint[kana={きぬ\tl{˥}\ws{}みー\tl{˥˩}},ipa={kinu miː},pos={句},acc={H F}]}%
	{%
		\MeaningsBlock{%
	\ryumeaning%
		{林の中。\ruby[g]{藪}{やぶ}の中}{}{}%
	}%
	}%
	{\tailPrint[]}%
\entry
	{\headPrint[kana={きぬん\tl{˥˩}},ipa={ki̥nuɴ},pos={動},rui={1},para={きなぬ},acc={F}]}%
	{%
		\MeaningsBlock{%
	\ryumeaning%
		{ねだる。せがむ}{}{}%
	}%
	}%
	{\tailPrint[]}%
\entry
	{\headPrint[kana={ぎばるん\tl{˥˩}},ipa={gibaruɴ},pos={動},rui={1},para={ぎばらぬ},acc={F}]}%
	{%
		\MeaningsBlock{%
	\ryumeaning%
		{頑張る。よく働く。尽くす。努力する}{}{}%
	}%
	}%
	{\tailPrint[]}%
\entry
	{\headPrint[kana={きぷ\tl{˥˩}},ipa={ki̥pu},pos={名},acc={F}]}%
	{%
		\MeaningsBlock{%
	\ryumeaning%
		{湯気。蒸気}{}{}%
	}%
	}%
	{\tailPrint[]}%
\entry
	{\headPrint[kana={きぷさ\tl{˥˩}はーん},ipa={ki̥pusahaːɴ},pos={形},acc={F}]}%
	{%
		\MeaningsBlock{%
	\ryumeaning%
		{煙たい。煙る}{}{}%
	}%
	}%
	{\tailPrint[]}%
\entry
	{\headPrint[kana={きまるん\tl{˥}},ipa={ki̥maruɴ},pos={動},rui={1},para={きまらぬ},acc={H}]}%
	{%
		\MeaningsBlock{%
	\ryumeaning%
		{決まる。定まる}{}{}%
	}%
	}%
	{\tailPrint[]}%
\entry
	{\headPrint[kana={きみるん\tl{˥}},ipa={ki̥miruɴ},pos={動},rui={3},para={きむぬ},acc={H}]}%
	{%
		\MeaningsBlock{%
	\ryumeaning%
		{決める。決定する}{}{}%
	}%
	}%
	{\tailPrint[]}%
\entry
	{\headPrint[kana={きむ\tl{˥}},ipa={ki̥mu},pos={名},acc={H}]}%
	{%
		\MeaningsBlock{%
	\ryumeaning%
		{肝。肝臓}{}{}%
	}%
	}%
	{\tailPrint[]}%
\entry
	{\headPrint[kana={きむぐくる\tl{˥}},ipa={ki̥mugukuru},pos={名},acc={H}]}%
	{%
		\MeaningsBlock{%
	\ryumeaning%
		{心根。心情}{}{}%
	}%
	}%
	{\tailPrint[]}%
\entry
	{\headPrint[kana={きゃんぎ\tl{˥˩}},ipa={kjaŋgi},pos={名},acc={F}]}%
	{%
		\MeaningsBlock{%
	\ryumeaning%
		{イヌマキ}{樹木名。最良の建材}{}%
	}%
	}%
	{\tailPrint[]}%
\entry
	{\headPrint[kana={きゅー\tl{˥}},ipa={kjuː},pos={名},acc={H}]}%
	{%
		\MeaningsBlock{%
	\ryumeaning%
		{今日。本日}{}{}%
	}%
	}%
	{\tailPrint[]}%
\entry
	{\headPrint[kana={きゆみるん\tl{˥}},ipa={ki̥jumiruɴ},pos={動},rui={3},para={きゆむぬ},acc={H}]}%
	{%
		\MeaningsBlock{%
	\ryumeaning%
		{清める}{}{}%
	}%
	}%
	{\tailPrint[]}%
\entry
	{\headPrint[kana={ぎょーれつ\tl{˥˩}},ipa={gjoːretsu},pos={名},acc={F}]}%
	{%
		\MeaningsBlock{%
	\ryumeaning%
		{行列}{旧盆の仮装行列は\emb{みちすねー}という}{}%
	}%
	}%
	{\tailPrint[]}%
\entry
	{\headPrint[kana={ぎりかた\tl{˩˥}はん},ipa={girikḁtahaɴ},pos={形},acc={R}]}%
	{%
		\MeaningsBlock{%
	\ryumeaning%
		{義理堅い。律義だ}{}{}%
	}%
	}%
	{\tailPrint[]}%
\entry
	{\headPrint[kana={きりとーすん\tl{˥}},ipa={ki̥ritoːsuɴ},pos={動},rui={2},para={きりとはぬ},acc={H}]}%
	{%
		\MeaningsBlock{%
	\ryumeaning%
		{切り倒す}{}{}%
	}%
	}%
	{\tailPrint[]}%
\entry
	{\headPrint[kana={きりとーすん\tl{˥}},ipa={ki̥ritoːsuɴ},pos={動},rui={2},para={きりとーはぬ},acc={H}]}%
	{%
		\MeaningsBlock{%
	\ryumeaning%
		{蹴り倒す}{}{}%
	}%
	}%
	{\tailPrint[]}%
\entry
	{\headPrint[kana={きりとぅぱすん\tl{˥}},ipa={ki̥ritu̥pasuɴ},pos={動},rui={2},para={きりとぅぱはぬ},acc={H}]}%
	{%
		\MeaningsBlock{%
	\ryumeaning%
		{蹴飛ばす}{}{}%
	}%
	}%
	{\tailPrint[]}%
\entry
	{\headPrint[kana={きるん\tl{˥}},ipa={ki̥ruɴ},pos={動},rui={1},para={きらぬ},acc={H}]}%
	{%
		\MeaningsBlock{%
	\ryumeaning%
		{蹴る}{}{}%
	}%
	}%
	{\tailPrint[]}%
\entry
	{\headPrint[kana={きん\tl{˥˩}},ipa={kiɴ},pos={名},acc={F}]}%
	{%
		\MeaningsBlock{%
	\ryumeaning%
		{斤。斤数}{重さの単位}{}%
	}%
	}%
	{\tailPrint[]}%
\entry
	{\headPrint[kana={きんき\ws{}すん\tl{˥˩}},ipa={kiŋki suɴ},pos={句},acc={HL F}]}%
	{%
		\MeaningsBlock{%
	\ryumeaning%
		{黄ばむ。黄色になる}{}{}%
	}%
	}%
	{\tailPrint[]}%
\entry
	{\headPrint[kana={きんさ\tl{˥}},ipa={kinsa},pos={名},acc={H}]}%
	{%
		\MeaningsBlock{%
	\ryumeaning%
		{検査}{}{}%
	}%
	}%
	{\tailPrint[]}%
\entry
	{\headPrint[kana={きんだるん\tl{˥˩}},ipa={kindaruɴ},pos={動},rui={1},para={きんだらぬ},acc={F}]}%
	{%
		\MeaningsBlock{%
	\ryumeaning%
		{掻き乱す}{}{}%
	}%
	}%
	{\tailPrint[]}%
\entry
	{\headPrint[kana={きんつぁーらすん\tl{˥}},ipa={kintsaːrasuɴ},pos={動},rui={2},para={きんつぁーらはぬ},acc={H}]}%
	{%
		\MeaningsBlock{%
	\ryumeaning%
		{掻き混ぜる}{}{}%
	}%
	}%
	{\tailPrint[]}%
\entry
	{\headPrint[kana={ぎんみ\tl{˩˥}},ipa={gimmi},pos={名},acc={R}]}%
	{%
		\MeaningsBlock{%
	\ryumeaning%
		{吟味。検討}{}{}%
	}%
	}%
	{\tailPrint[]}%
\LetterHead{く}
\entry
	{\headPrint[kana={く\tl{˥}},ipa={ku},pos={名},acc={H}]}%
	{%
		\MeaningsBlock{%
	\ryumeaning%
		{粉。粉末}{}{}%
	}%
	}%
	{\tailPrint[]}%
\entry
	{\headPrint[kana={ぐー\tl{˩˥}},ipa={guː},pos={名},acc={R}]}%
	{%
		\MeaningsBlock{%
	\ryumeaning%
		{相棒。仲間。連れ}{}{}%
	}%
	}%
	{\tailPrint[]}%
\entry
	{\headPrint[kana={くい\tl{˥}},ipa={kui},pos={名},acc={H}]}%
	{%
		\MeaningsBlock{%
	\ryumeaning%
		{声}{}{}%
	}%
	}%
	{\tailPrint[]}%
\entry
	{\headPrint[kana={くいつぁーすん\tl{˥}},ipa={kuitsaːsuɴ},pos={動},rui={2},para={くいつぁーはぬ},acc={H}]}%
	{%
		\MeaningsBlock{%
	\ryumeaning%
		{揺する。揺らす}{}{}%
	}%
	}%
	{\tailPrint[]}%
\entry
	{\headPrint[kana={ぐいふ\tl{˩˥}},ipa={guifu},pos={名},acc={R}]}%
	{%
		\MeaningsBlock{%
	\ryumeaning%
		{御用布}{王府時代の人頭税の一つ}{}%
	}%
	}%
	{\tailPrint[]}%
\entry
	{\headPrint[kana={くいるん\tl{˥}},ipa={kuiruɴ},pos={動},acc={H}]}%
	{%
		\MeaningsBlock{%
	\ryumeaning%
		{請う。求める}{}{}%
	\ryumeaning%
		{求婚する}{嫁として申し込む}{}%
	}%
	}%
	{\tailPrint[]}%
\entry
	{\headPrint[kana={くいるん\tl{˥˩}},ipa={kuiruɴ},pos={動},rui={3},para={くいぬ},acc={F}]}%
	{%
		\MeaningsBlock{%
	\ryumeaning%
		{越える}{}{}%
	}%
	}%
	{\tailPrint[]}%
\entry
	{\headPrint[kana={くがすん\tl{˥˩}},ipa={kugusuɴ},pos={動},rui={2},para={くぐはぬ},acc={F}]}%
	{%
		\MeaningsBlock{%
	\ryumeaning%
		{焦がす}{}{}%
	}%
	}%
	{\tailPrint[]}%
\entry
	{\headPrint[kana={くがに\tl{˥˩}},ipa={kugani},pos={名},acc={F}]}%
	{%
		\MeaningsBlock{%
	\ryumeaning%
		{\ruby[g]{黄}{こ}\ruby[g]{金}{がね}}{}{}%
	}%
	}%
	{\tailPrint[]}%
\entry
	{\headPrint[kana={くがりるん\tl{˥˩}},ipa={kugariruɴ},pos={動},rui={3},para={くがるぬ},acc={F}]}%
	{%
		\MeaningsBlock{%
	\ryumeaning%
		{焦げる}{}{}%
	\ryumeaning%
		{焦がれる}{}{}%
	}%
	}%
	{\tailPrint[]}%
\entry
	{\headPrint[kana={くくぬち\tl{˥}},ipa={ku̥kunutʃi},pos={名},acc={H}]}%
	{%
		\MeaningsBlock{%
	\ryumeaning%
		{九つ}{}{}%
	}%
	}%
	{\tailPrint[note={簡略化\emb{はこな}}]}%
\entry
	{\headPrint[kana={くくる\tl{˥}},ipa={ku̥kuru},pos={名},acc={H}]}%
	{%
		\MeaningsBlock{%
	\ryumeaning%
		{心。精神}{}{}%
	}%
	}%
	{\tailPrint[]}%
\entry
	{\headPrint[kana={くくるいるん},ipa={ku̥kuruiruɴ},pos={動}]}%
	{%
		\MeaningsBlock{%
	\ryumeaning%
		{心得る。心掛ける。気をつける}{}{}%
	}%
	}%
	{\tailPrint[]}%
\entry
	{\headPrint[kana={くくるがきるん},ipa={ku̥kurugakiruɴ},pos={動},rui={3},para={くくるがくぬ}]}%
	{%
		\MeaningsBlock{%
	\ryumeaning%
		{心掛ける}{}{}%
	}%
	}%
	{\tailPrint[]}%
\entry
	{\headPrint[kana={くさむくん\tl{˥}},ipa={ku̥samukuɴ},pos={動},rui={1},para={くさむかぬ},acc={H}]}%
	{%
		\MeaningsBlock{%
	\ryumeaning%
		{憤慨する。しゃくにさわる}{}{}%
	}%
	}%
	{\tailPrint[]}%
\entry
	{\headPrint[kana={ぐさん\tl{˩˥}},ipa={gusaɴ},pos={名},acc={R}]}%
	{%
		\MeaningsBlock{%
	\ryumeaning%
		{杖}{}{}%
	}%
	}%
	{\tailPrint[]}%
\entry
	{\headPrint[kana={ぐじぃら\tl{˩˥}},ipa={gudʒira},pos={名},acc={R}]}%
	{%
		\MeaningsBlock{%
	\ryumeaning%
		{鯨}{}{}%
	}%
	}%
	{\tailPrint[]}%
\entry
	{\headPrint[kana={くしみ\tl{˥}},ipa={kuʃi̥mi},pos={名},acc={H}]}%
	{%
		\MeaningsBlock{%
	\ryumeaning%
		{甲烏賊}{}{}%
	}%
	}%
	{\tailPrint[]}%
\entry
	{\headPrint[kana={ぐじゅー\tl{˩˥}},ipa={gudʒuː},pos={名},acc={R}]}%
	{%
		\MeaningsBlock{%
	\ryumeaning%
		{九十}{}{}%
	}%
	}%
	{\tailPrint[]}%
\entry
	{\headPrint[kana={ぐしん\tl{˩˥}},ipa={guʃiɴ},pos={名},acc={R}]}%
	{%
		\MeaningsBlock{%
	\ryumeaning%
		{お酒。お神酒}{酒でお供えしたものを指す}{}%
	}%
	}%
	{\tailPrint[]}%
\entry
	{\headPrint[kana={ぐす\tl{˥˩}},ipa={gusu},pos={名},acc={F}]}%
	{%
		\MeaningsBlock{%
	\ryumeaning%
		{トウガラシ}{香辛料}{}%
	}%
	}%
	{\tailPrint[]}%
\entry
	{\headPrint[kana={ぐそー\tl{˩˥}},ipa={gusoː},pos={名},acc={R}]}%
	{%
		\MeaningsBlock{%
	\ryumeaning%
		{あの世。後生}{}{}%
	}%
	}%
	{\tailPrint[]}%
\entry
	{\headPrint[kana={ぐそーみ\tl{˩˥}},ipa={gusoːmi},pos={名},acc={R}]}%
	{%
		\MeaningsBlock{%
	\ryumeaning%
		{五勺米}{御嶽神事の徴収米}{}%
	}%
	}%
	{\tailPrint[]}%
\entry
	{\headPrint[kana={くたいるん\tl{˥}},ipa={ku̥tairuɴ},pos={動},rui={3},para={くたいるぬ},acc={H}]}%
	{%
		\MeaningsBlock{%
	\ryumeaning%
		{答える}{}{}%
	}%
	}%
	{\tailPrint[]}%
\entry
	{\headPrint[kana={くだくん},ipa={kudakuɴ},pos={動},rui={1},para={くだかぬ}]}%
	{%
		\MeaningsBlock{%
	\ryumeaning%
		{砕く。打ち砕く}{}{}%
	}%
	}%
	{\tailPrint[]}%
\entry
	{\headPrint[kana={くたんでー\tl{˥}},ipa={ku̥tandeː},pos={名},acc={H}]}%
	{%
		\MeaningsBlock{%
	\ryumeaning%
		{草臥れ。疲労}{}{}%
	}%
	}%
	{\tailPrint[]}%
\entry
	{\headPrint[kana={くたんでぃるん\tl{˥}},ipa={ku̥tandiruɴ},pos={動},rui={1},para={くたんでぃらぬ},acc={H}]}%
	{%
		\MeaningsBlock{%
	\ryumeaning%
		{疲れる。草臥れる}{}{}%
	}%
	}%
	{\tailPrint[]}%
\entry
	{\headPrint[kana={くち\tl{˥}},ipa={ku̥tʃi},pos={名},acc={H}]}%
	{%
		\MeaningsBlock{%
	\ryumeaning%
		{籐}{蔓の植物}{}%
	}%
	}%
	{\tailPrint[]}%
\entry
	{\headPrint[kana={くち\tl{˥˩}},ipa={ku̥tʃi},pos={名},acc={F}]}%
	{%
		\MeaningsBlock{%
	\ryumeaning%
		{腰}{}{}%
	}%
	}%
	{\tailPrint[]}%
\entry
	{\headPrint[kana={くちさ\tl{˥˩}はん},ipa={ku̥tʃisḁhaɴ},pos={形},para={くちさへぬ},acc={F}]}%
	{%
		\MeaningsBlock{%
	\ryumeaning%
		{苦しい。窮屈だ}{}{}%
	}%
	}%
	{\tailPrint[]}%
\entry
	{\headPrint[kana={ぐちみん\tl{˩˥}},ipa={gutʃimiɴ},pos={名},acc={R}]}%
	{%
		\MeaningsBlock{%
	\ryumeaning%
		{\ruby[g]{脇}{わき}。脇の下}{}{}%
	}%
	}%
	{\tailPrint[]}%
\entry
	{\headPrint[kana={ぐちゅぐちゅ},ipa={gutʃugutʃu},pos={擬}]}%
	{%
		\MeaningsBlock{%
	\ryumeaning%
		{こちょこちょ}{くすぐる時の形容}{}%
	}%
	}%
	{\tailPrint[]}%
\entry
	{\headPrint[kana={ぐちゅるん\tl{˥˩}},ipa={gutʃuruɴ},pos={動},rui={1},para={ぐちゅらぬ},acc={F}]}%
	{%
		\MeaningsBlock{%
	\ryumeaning%
		{\ruby[g]{擽}{くすぐ}る}{}{}%
	}%
	}%
	{\tailPrint[]}%
\entry
	{\headPrint[kana={くつぁーすん\tl{˥˩}},ipa={ku̥tsaːsuɴ},pos={動},rui={2b},para={くつぁーさぬ},acc={F}]}%
	{%
		\MeaningsBlock{%
	\ryumeaning%
		{（魚などを）拵える}{}{}%
	}%
	}%
	{\tailPrint[]}%
\entry
	{\headPrint[kana={くつぁすん\tl{˥˩}},ipa={ku̥tsasuɴ},pos={動},rui={2},para={くつぁはぬ},acc={F}]}%
	{%
		\MeaningsBlock{%
	\ryumeaning%
		{拵える。こさえる}{}{}%
	}%
	}%
	{\tailPrint[]}%
\entry
	{\headPrint[kana={ぐつぇー\tl{˩˥}},ipa={gutseː},pos={名},acc={R}]}%
	{%
		\MeaningsBlock{%
	\ryumeaning%
		{（大きな）雄牛}{}{}%
	}%
	}%
	{\tailPrint[]}%
\entry
	{\headPrint[kana={くつん\tl{˥}},ipa={ku̥tsuɴ},pos={名},acc={H}]}%
	{%
		\MeaningsBlock{%
	\ryumeaning%
		{去年。昨年}{}{}%
	}%
	}%
	{\tailPrint[]}%
\entry
	{\headPrint[kana={ぐてー\tl{˩˥}},ipa={guteː},pos={名},acc={R}]}%
	{%
		\MeaningsBlock{%
	\ryumeaning%
		{体。体格}{}{}%
	}%
	}%
	{\tailPrint[]}%
\entry
	{\headPrint[kana={くとぅしぃ\tl{˥˩}},ipa={ku̥tuʃi},pos={名},acc={F}]}%
	{%
		\MeaningsBlock{%
	\ryumeaning%
		{今年}{}{}%
	}%
	}%
	{\tailPrint[]}%
\entry
	{\headPrint[kana={くとぅしきるん\tl{˥}},ipa={ku̥tuʃi̥kiruɴ},pos={動},rui={3},para={くとすくぬ},acc={H}]}%
	{%
		\MeaningsBlock{%
	\ryumeaning%
		{言付ける。伝言する}{}{}%
	}%
	}%
	{\tailPrint[]}%
\entry
	{\headPrint[kana={くとぅば\tl{˥}},ipa={ku̥tuba},pos={名},acc={H}]}%
	{%
		\MeaningsBlock{%
	\ryumeaning%
		{言葉。言語}{}{}%
	}%
	}%
	{\tailPrint[]}%
\entry
	{\headPrint[kana={くとぅばるん\tl{˥}},ipa={ku̥tubaruɴ},pos={動},rui={1},para={くとぅばらぬ},acc={H}]}%
	{%
		\MeaningsBlock{%
	\ryumeaning%
		{断る}{}{}%
	}%
	}%
	{\tailPrint[]}%
\entry
	{\headPrint[kana={くなすん\tl{˥˩}},ipa={ku̥nasuɴ},pos={動},rui={2},para={くなはぬ},acc={F}]}%
	{%
		\MeaningsBlock{%
	\ryumeaning%
		{こなす。柔らかくする}{}{}%
	}%
	}%
	{\tailPrint[]}%
\entry
	{\headPrint[kana={くなすん\tl{˥˩}},ipa={ku̥nasuɴ},pos={動},rui={2},para={くなはぬ},acc={F}]}%
	{%
		\MeaningsBlock{%
	\ryumeaning%
		{（お腹を）下す。下痢する}{}{}%
	}%
	}%
	{\tailPrint[]}%
\entry
	{\headPrint[kana={くなちぃ\tl{˥}},ipa={kunatʃi},pos={名},acc={H}]}%
	{%
		\MeaningsBlock{%
	\ryumeaning%
		{来年の夏}{「来る夏」の意味}{}%
	}%
	}%
	{\tailPrint[]}%
\entry
	{\headPrint[kana={くぬ\tl{˥˩}},ipa={ku̥nu},pos={連体},acc={F}]}%
	{%
		\MeaningsBlock{%
	\ryumeaning%
		{この。これの}{}{}%
	}%
	}%
	{\tailPrint[]}%
\entry
	{\headPrint[kana={くぱ\tl{˥˩}},ipa={ku̥pa},pos={名},acc={F}]}%
	{%
		\MeaningsBlock{%
	\ryumeaning%
		{クバ。ビロウ}{植物名}{}%
	}%
	}%
	{\tailPrint[]}%
\entry
	{\headPrint[kana={くばかつぁ\tl{˥˩}},ipa={kubakḁtsa},pos={名},acc={F}]}%
	{%
		\MeaningsBlock{%
	\ryumeaning%
		{クバ笠}{クバを材料とした笠}{}%
	}%
	}%
	{\tailPrint[]}%
\entry
	{\headPrint[kana={くぱすん\tl{˥}},ipa={ku̥pasuɴ},pos={動},rui={2},para={くぱはぬ},acc={H}]}%
	{%
		\MeaningsBlock{%
	\ryumeaning%
		{\ruby[g]{零}{こぼ}す}{}{}%
	}%
	}%
	{\tailPrint[]}%
\entry
	{\headPrint[kana={くぱ\tl{˥}はん},ipa={ku̥pahaɴ},pos={形},para={くぱへぬ},acc={H}]}%
	{%
		\MeaningsBlock{%
	\ryumeaning%
		{下手だ。不器用だ}{}{}%
	}%
	}%
	{\tailPrint[]}%
\entry
	{\headPrint[kana={くぱるん\tl{˥}},ipa={ku̥paruɴ},pos={動},rui={1},para={くぱらぬ},acc={H}]}%
	{%
		\MeaningsBlock{%
	\ryumeaning%
		{（寒さで）凍える}{}{}%
	}%
	}%
	{\tailPrint[]}%
\entry
	{\headPrint[kana={くぱるん\tl{˥}},ipa={ku̥paruɴ},pos={動},rui={1},para={くぱらぬ},acc={H}]}%
	{%
		\MeaningsBlock{%
	\ryumeaning%
		{配る。分配する}{}{}%
	}%
	}%
	{\tailPrint[]}%
\entry
	{\headPrint[kana={くぱん\tl{˥}},ipa={ku̥paɴ},pos={名},acc={H}]}%
	{%
		\MeaningsBlock{%
	\ryumeaning%
		{神前の供物の名}{}{}%
	}%
	}%
	{\tailPrint[]}%
\entry
	{\headPrint[kana={くぷりるん\tl{˥}},ipa={ku̥puriruɴ},pos={動},rui={3},para={くぷるぬ},acc={H}]}%
	{%
		\MeaningsBlock{%
	\ryumeaning%
		{\ruby[g]{零}{こぼ}れる}{}{}%
	}%
	}%
	{\tailPrint[]}%
\entry
	{\headPrint[kana={くぷりん\tl{˥}},ipa={ku̥puriɴ},pos={動},rui={3},para={くぷるぬ},acc={H}]}%
	{%
		\MeaningsBlock{%
	\ryumeaning%
		{\ruby[g]{零}{こぼ}れる}{}{}%
	}%
	}%
	{\tailPrint[]}%
\entry
	{\headPrint[kana={ぐぼん\tl{˩˥}},ipa={guboɴ},pos={名},acc={R}]}%
	{%
		\MeaningsBlock{%
	\ryumeaning%
		{\ruby[g]{牛蒡}{ごぼう}}{野菜名}{}%
	}%
	}%
	{\tailPrint[]}%
\entry
	{\headPrint[kana={くまた\tl{˥}},ipa={ku̥mata},pos={名},acc={H}]}%
	{%
		\MeaningsBlock{%
	\ryumeaning%
		{分け前}{}{}%
	}%
	}%
	{\tailPrint[]}%
\entry
	{\headPrint[kana={くま\tl{˥}はん},ipa={ku̥mahaɴ},pos={形},para={くまへぬ},acc={H}]}%
	{%
		\MeaningsBlock{%
	\ryumeaning%
		{細かい}{}{}%
	\ryumeaning%
		{（手が）器用だ。（技が）精巧だ}{}{}%
	}%
	}%
	{\tailPrint[]}%
\entry
	{\headPrint[kana={ぐま\tl{˩˥}はん},ipa={gumahaɴ},pos={形},para={ぐまへぬ},acc={R}]}%
	{%
		\MeaningsBlock{%
	\ryumeaning%
		{小さい。小粒だ}{}{}%
	}%
	}%
	{\tailPrint[]}%
\entry
	{\headPrint[kana={くまるん\tl{˥}},ipa={ku̥maruɴ},pos={動},rui={1},para={くまらぬ},acc={H}]}%
	{%
		\MeaningsBlock{%
	\ryumeaning%
		{\ruby[g]{籠}{こも}る}{}{}%
	\ryumeaning%
		{隠れる}{}{}%
	}%
	}%
	{\tailPrint[]}%
\entry
	{\headPrint[kana={くらすん\tl{˥˩}},ipa={ku̥rasuɴ},pos={動},rui={2},para={くらはぬ},acc={F}]}%
	{%
		\MeaningsBlock{%
	\ryumeaning%
		{殺す}{}{}%
	\ryumeaning%
		{打つ。殴る}{}{}%
	}%
	}%
	{\tailPrint[]}%
\entry
	{\headPrint[kana={くらびるん\tl{˥˩}},ipa={ku̥rabiruɴ},pos={動},rui={3},acc={F}]}%
	{%
		\MeaningsBlock{%
	\ryumeaning%
		{比べる。比較する}{}{}%
	}%
	}%
	{\tailPrint[]}%
\entry
	{\headPrint[kana={ぐるっけーすん\tl{˩˥}},ipa={gurukkeːsuɴ},pos={動},rui={2b},para={ぐるっけーさぬ},acc={R}]}%
	{%
		\MeaningsBlock{%
	\ryumeaning%
		{引っくり返す}{}{}%
	}%
	}%
	{\tailPrint[]}%
\entry
	{\headPrint[kana={ぐるっけーらん\tl{˩˥}},ipa={gurukkeːraɴ},pos={動（継）},rui={1},para={ぐるっけーらぬ},acc={R}]}%
	{%
		\MeaningsBlock{%
	\ryumeaning%
		{引っくり返る}{}{}%
	}%
	}%
	{\tailPrint[]}%
\entry
	{\headPrint[kana={くるばすん\tl{˥˩}},ipa={ku̥rubasuɴ},pos={動},rui={2},para={くるばはぬ},acc={F}]}%
	{%
		\MeaningsBlock{%
	\ryumeaning%
		{転ばす。転がす}{}{}%
	\ryumeaning%
		{横たえる}{}{}%
	}%
	}%
	{\tailPrint[]}%
\entry
	{\headPrint[kana={くるぶん\tl{˥˩}},ipa={ku̥rubuɴ},pos={動},rui={1},para={くるばぬ},acc={F}]}%
	{%
		\MeaningsBlock{%
	\ryumeaning%
		{転ぶ。転がる}{}{}%
	}%
	}%
	{\tailPrint[]}%
\entry
	{\headPrint[kana={くるまぼー\tl{˥}},ipa={kurumaboː},pos={名},acc={H}]}%
	{%
		\MeaningsBlock{%
	\ryumeaning%
		{車棒。くるり棒}{豆打ちの農具}{}%
	}%
	}%
	{\tailPrint[]}%
\entry
	{\headPrint[kana={くわいるん\tl{˥}},ipa={kuwairuɴ},pos={動},acc={H}]}%
	{%
		\MeaningsBlock{%
	\ryumeaning%
		{加える}{}{}%
	}%
	}%
	{\tailPrint[]}%
\entry
	{\headPrint[kana={くん\tl{˥}},ipa={kuɴ},pos={動},rui={irr},para={くぬ},acc={H}]}%
	{%
		\MeaningsBlock{%
	\ryumeaning%
		{来る}{}{}%
	}%
	}%
	{\tailPrint[]}%
\entry
	{\headPrint[kana={くん\tl{˥}},ipa={kuɴ},pos={動},rui={1},para={くわぬ},acc={H}]}%
	{%
		\MeaningsBlock{%
	\ryumeaning%
		{漕ぐ}{}{}%
	}%
	}%
	{\tailPrint[]}%
\entry
	{\headPrint[kana={くんき\tl{˥˩}},ipa={kuŋki},pos={名},acc={F}]}%
	{%
		\MeaningsBlock{%
	\ryumeaning%
		{根気。体力}{}{}%
	}%
	}%
	{\tailPrint[]}%
\entry
	{\headPrint[kana={くんき\ws{}しきるん},ipa={kuŋki ʃi̥kiruɴ},pos={句},acc={-}]}%
	{%
		\MeaningsBlock{%
	\ryumeaning%
		{栄養をつける}{}{}%
	}%
	}%
	{\tailPrint[]}%
\entry
	{\headPrint[kana={ぐんず\tl{˩˥}},ipa={gundzu},pos={名},acc={R}]}%
	{%
		\MeaningsBlock{%
	\ryumeaning%
		{五十}{}{}%
	}%
	}%
	{\tailPrint[]}%
\entry
	{\headPrint[kana={くんぞー\tl{˥}},ipa={kundzoː},pos={名},acc={H}]}%
	{%
		\MeaningsBlock{%
	\ryumeaning%
		{怒り。怒ること}{}{}%
	}%
	}%
	{\tailPrint[]}%
\entry
	{\headPrint[kana={くんぞーむぬ},ipa={kundzoːmunu},pos={名}]}%
	{%
		\MeaningsBlock{%
	\ryumeaning%
		{怒りんぼ}{}{}%
	}%
	}%
	{\tailPrint[]}%
\entry
	{\headPrint[kana={くんた\tl{˥˩}ばり\tl{˩˥}},ipa={kuntabari},pos={名},acc={F+R}]}%
	{%
		\MeaningsBlock{%
	\ryumeaning%
		{この間。少し前}{}{}%
	}%
	}%
	{\tailPrint[]}%
\entry
	{\headPrint[kana={ぐんぼー\tl{˩˥}},ipa={gumboː},pos={名},acc={R}]}%
	{%
		\MeaningsBlock{%
	\ryumeaning%
		{私生児。隠し子}{}{}%
	}%
	}%
	{\tailPrint[]}%
\LetterHead{け}
\entry
	{\headPrint[kana={けー},ipa={keː},pos={名}]}%
	{%
		\MeaningsBlock{%
	\ryumeaning%
		{陰。日蔭}{}{}%
	}%
	}%
	{\tailPrint[]}%
\entry
	{\headPrint[kana={けー\tl{˥}},ipa={keː},pos={名},acc={H}]}%
	{%
		\MeaningsBlock{%
	\ryumeaning%
		{卵}{鶏卵は\emb{ごかけー}。魚・蟹の卵は\emb{ぱりゃん}と言う}{}%
	}%
	}%
	{\tailPrint[]}%
\entry
	{\headPrint[kana={けー\tl{˥˩}},ipa={keː},pos={名},acc={F}]}%
	{%
		\MeaningsBlock{%
	\ryumeaning%
		{井戸}{}{}%
	}%
	}%
	{\tailPrint[]}%
\entry
	{\headPrint[kana={けーし},ipa={keːʃi},pos={副}]}%
	{%
		\MeaningsBlock{%
	\ryumeaning%
		{きれいに。立派に}{}{}%
	}%
	}%
	{\tailPrint[]}%
\entry
	{\headPrint[kana={けーすん\tl{˥}},ipa={keːsuɴ},pos={動},rui={2},para={けーさぬ},acc={H}]}%
	{%
		\MeaningsBlock{%
	\ryumeaning%
		{耕す}{「返す」の義}{}%
	}%
	}%
	{\tailPrint[]}%
\entry
	{\headPrint[kana={けーら},ipa={keːra},pos={名}]}%
	{%
		\MeaningsBlock{%
	\ryumeaning%
		{\ruby[g]{皆}{みんな}。全部}{}{}%
	}%
	}%
	{\tailPrint[]}%
\entry
	{\headPrint[kana={けーら\tl{˥}ねーら\tl{˩˥}},ipa={keːraneːra},pos={名},acc={H+R}]}%
	{%
		\MeaningsBlock{%
	\ryumeaning%
		{皆さま。皆皆様}{\emb{けーら}は「皆・全部」}{}%
	}%
	}%
	{\tailPrint[]}%
\entry
	{\headPrint[kana={けーり\tl{˥}\ws{}とーすん\tl{˥}},ipa={keːri toːsuɴ},pos={句},rui={2},para={けーり　とーはぬ},acc={H H}]}%
	{%
		\MeaningsBlock{%
	\ryumeaning%
		{切り倒す。薙ぎ倒す}{}{}%
	}%
	}%
	{\tailPrint[]}%
\entry
	{\headPrint[kana={けーるん\tl{˥}},ipa={keːruɴ},pos={動},rui={1},para={けーらぬ},acc={H}]}%
	{%
		\MeaningsBlock{%
	\ryumeaning%
		{切り倒す。（固いものを）叩き切る}{}{}%
	}%
	}%
	{\tailPrint[]}%
\entry
	{\headPrint[kana={けーるん\tl{˥}},ipa={keːruɴ},pos={動},acc={H}]}%
	{%
		\MeaningsBlock{%
	\ryumeaning%
		{帰る}{}{}%
	}%
	}%
	{\tailPrint[]}%
\entry
	{\headPrint[kana={けしゃ\tl{˥}はん},ipa={keʃahaɴ},pos={形},para={けしゃへぬ},acc={H}]}%
	{%
		\MeaningsBlock{%
	\ryumeaning%
		{きれい。美しい}{物について言う}{}%
	}%
	}%
	{\tailPrint[]}%
\entry
	{\headPrint[kana={けすん\tl{˥˩}},ipa={kesuɴ},pos={動},rui={2},para={けさぬ},acc={F}]}%
	{%
		\MeaningsBlock{%
	\ryumeaning%
		{消す。消去する}{}{}%
	}%
	}%
	{\tailPrint[]}%
\entry
	{\headPrint[kana={\josho{}けっ\kako{}た},ipa={ketta},pos={副},acc={HL}]}%
	{%
		\MeaningsBlock{%
	\ryumeaning%
		{たいそう}{}{}%
	}%
	}%
	{\tailPrint[]}%
\entry
	{\headPrint[kana={けった\ws{}ぞっふぁら},ipa={ketta dzoffara},pos={句}]}%
	{%
		\MeaningsBlock{%
	\ryumeaning%
		{びっしょり}{水に濡れるさまの形容}{}%
	}%
	}%
	{\tailPrint[]}%
\entry
	{\headPrint[kana={けった\ws{}ばっさーん},ipa={ketta bassaːɴ},pos={句},acc={HL LHL}]}%
	{%
		\MeaningsBlock{%
	\ryumeaning%
		{すっかり忘れきる}{}{}%
	}%
	}%
	{\tailPrint[]}%
\entry
	{\headPrint[kana={けった\ws{}みゃーん},ipa={ketta mjaːɴ},pos={句},acc={HL HL}]}%
	{%
		\MeaningsBlock{%
	\ryumeaning%
		{十分に熟する。熟しきる}{}{}%
	}%
	}%
	{\tailPrint[]}%
\entry
	{\headPrint[kana={げん},ipa={geɴ},pos={名}]}%
	{%
		\MeaningsBlock{%
	\ryumeaning%
		{来年。来る年}{冨嘉の発音}{}%
	}%
	}%
	{\tailPrint[]}%
\LetterHead{こ}
\entry
	{\headPrint[kana={こー\tl{˥}},ipa={koː},pos={名},acc={H}]}%
	{%
		\MeaningsBlock{%
	\ryumeaning%
		{お香。線香}{}{}%
	}%
	}%
	{\tailPrint[]}%
\entry
	{\headPrint[kana={こーがき\tl{˥}},ipa={koːgaki},pos={名},acc={H}]}%
	{%
		\MeaningsBlock{%
	\ryumeaning%
		{頬かぶり。覆面}{}{}%
	}%
	}%
	{\tailPrint[]}%
\entry
	{\headPrint[kana={ごーごー},ipa={goːgoː},pos={擬}]}%
	{%
		\MeaningsBlock{%
	\ryumeaning%
		{ぐうぐう}{ぐっすり眠るさま}{}%
	\ryumeaning%
		{}{すうすうと寝入るように死んで行くさま}{}%
	}%
	}%
	{\tailPrint[]}%
\entry
	{\headPrint[kana={こーさー\tl{˥}},ipa={koːsaː},pos={名},acc={H}]}%
	{%
		\MeaningsBlock{%
	\ryumeaning%
		{拳骨}{}{}%
	}%
	}%
	{\tailPrint[]}%
\entry
	{\headPrint[kana={こーし},ipa={koːʃi},pos={副}]}%
	{%
		\MeaningsBlock{%
	\ryumeaning%
		{固く。堅実に。立派に}{}{}%
	}%
	}%
	{\tailPrint[]}%
\entry
	{\headPrint[kana={こーし\tl{˥}},ipa={koːʃi},pos={名},acc={H}]}%
	{%
		\MeaningsBlock{%
	\ryumeaning%
		{菓子}{}{}%
	}%
	}%
	{\tailPrint[]}%
\entry
	{\headPrint[kana={こーすん},ipa={koːsuɴ},pos={動},rui={2},para={こーはぬ},acc={-}]}%
	{%
		\MeaningsBlock{%
	\ryumeaning%
		{壊す}{}{}%
	}%
	}%
	{\tailPrint[]}%
\entry
	{\headPrint[kana={こーすん\tl{˥˩}},ipa={koːsuɴ},pos={動},rui={2b},para={こーさぬ},acc={F}]}%
	{%
		\MeaningsBlock{%
	\ryumeaning%
		{掘り取る。根こそぎ取る}{}{}%
	}%
	}%
	{\tailPrint[]}%
\entry
	{\headPrint[kana={こー\josho{}に\kako{}ー},ipa={koːniː},pos={名},acc={特殊}]}%
	{%
		\MeaningsBlock{%
	\ryumeaning%
		{男の児の愛称}{}{}%
	}%
	}%
	{\tailPrint[]}%
\entry
	{\headPrint[kana={こー\tl{˥}はん},ipa={koːhaɴ},pos={形},para={こーへぬ},acc={H}]}%
	{%
		\MeaningsBlock{%
	\ryumeaning%
		{固い。硬い}{}{}%
	}%
	}%
	{\tailPrint[]}%
\entry
	{\headPrint[kana={こー\tl{˥}はん},ipa={koːhaɴ},pos={形},acc={H}]}%
	{%
		\MeaningsBlock{%
	\ryumeaning%
		{歯ごたえがある}{}{}%
	}%
	}%
	{\tailPrint[]}%
\entry
	{\headPrint[kana={ごー\tl{˥}はん},ipa={goːhaɴ},pos={形},acc={H}]}%
	{%
		\MeaningsBlock{%
	\ryumeaning%
		{怖い。恐ろしい}{}{}%
	}%
	}%
	{\tailPrint[]}%
\entry
	{\headPrint[kana={こーみ\tl{˥}},ipa={koːmi},pos={名},acc={H}]}%
	{%
		\MeaningsBlock{%
	\ryumeaning%
		{お辞儀。拝む}{}{}%
	}%
	}%
	{\tailPrint[]}%
\entry
	{\headPrint[kana={こーりるん\tl{˥}},ipa={koːriruɴ},pos={動},rui={3},para={こーるぬ},acc={H}]}%
	{%
		\MeaningsBlock{%
	\ryumeaning%
		{壊れる。崩れる}{}{}%
	}%
	}%
	{\tailPrint[]}%
\entry
	{\headPrint[kana={こーるん\tl{˥}},ipa={koːruɴ},pos={動},rui={1},para={こーらぬ},acc={H}]}%
	{%
		\MeaningsBlock{%
	\ryumeaning%
		{固まる。硬くなる}{}{}%
	}%
	}%
	{\tailPrint[]}%
\entry
	{\headPrint[kana={こーろ\tl{˥}},ipa={koːro},pos={名},acc={H}]}%
	{%
		\MeaningsBlock{%
	\ryumeaning%
		{\ruby[g]{独楽}{こま}}{}{}%
	}%
	}%
	{\tailPrint[]}%
\entry
	{\headPrint[kana={こーん\tl{˥˩}},ipa={koːɴ},pos={動},para={かーぬ},acc={F}]}%
	{%
		\MeaningsBlock{%
	\ryumeaning%
		{買う}{}{}%
	}%
	}%
	{\tailPrint[]}%
\entry
	{\headPrint[kana={ごーんた},ipa={goːnta},pos={擬}]}%
	{%
		\MeaningsBlock{%
	\ryumeaning%
		{こつん}{物を軽く打つ音の形容}{}%
	}%
	}%
	{\tailPrint[]}%
\entry
	{\headPrint[kana={こい\tl{˥}},ipa={koi},pos={名},acc={H}]}%
	{%
		\MeaningsBlock{%
	\ryumeaning%
		{肥し。肥料}{}{}%
	}%
	}%
	{\tailPrint[]}%
\entry
	{\headPrint[kana={こいすぷ\tl{˥˩}},ipa={koisu̥pu},pos={名},acc={F}]}%
	{%
		\MeaningsBlock{%
	\ryumeaning%
		{肥溜}{}{}%
	}%
	}%
	{\tailPrint[]}%
\entry
	{\headPrint[kana={こいたんぐ\tl{˥˩}},ipa={koitangu},pos={名},acc={F}]}%
	{%
		\MeaningsBlock{%
	\ryumeaning%
		{水肥おけ}{}{}%
	}%
	}%
	{\tailPrint[]}%
\entry
	{\headPrint[kana={ごかけー\tl{˥}},ipa={gokakeː},pos={名},acc={H}]}%
	{%
		\MeaningsBlock{%
	\ryumeaning%
		{ニワトリの卵。鶏卵}{}{}%
	}%
	}%
	{\tailPrint[]}%
\entry
	{\headPrint[kana={こさ\tl{˥}はん},ipa={kosahaɴ},pos={形},para={こさへぬ},acc={H}]}%
	{%
		\MeaningsBlock{%
	\ryumeaning%
		{\ruby[g]{擽}{くす}ったい}{}{}%
	}%
	}%
	{\tailPrint[]}%
\entry
	{\headPrint[kana={こすぱり},ipa={kosu̥pari},pos={-}]}%
	{%
		\MeaningsBlock{%
	\ryumeaning%
		{歯ごたえがある}{}{}%
	}%
	}%
	{\tailPrint[]}%
\entry
	{\headPrint[kana={こすぱるん\tl{˥}},ipa={kosu̥paruɴ},pos={動},rui={1},para={こすぱらぬ},acc={H}]}%
	{%
		\MeaningsBlock{%
	\ryumeaning%
		{固まる。硬くなる}{}{}%
	}%
	}%
	{\tailPrint[]}%
\entry
	{\headPrint[kana={こすん\tl{˥˩}},ipa={kosuɴ},pos={動},rui={2b},para={こさぬ},acc={F}]}%
	{%
		\MeaningsBlock{%
	\ryumeaning%
		{漉す}{}{}%
	}%
	}%
	{\tailPrint[]}%
\entry
	{\headPrint[kana={こすん\tl{˥˩}},ipa={kosuɴ},pos={動},rui={2},para={こはぬ},acc={F}]}%
	{%
		\MeaningsBlock{%
	\ryumeaning%
		{越す。越える。越させる}{}{}%
	}%
	}%
	{\tailPrint[]}%
\entry
	{\headPrint[kana={ごっか\tl{˩˥}},ipa={gokka},pos={名},acc={R}]}%
	{%
		\MeaningsBlock{%
	\ryumeaning%
		{鶏}{}{}%
	}%
	}%
	{\tailPrint[]}%
\entry
	{\headPrint[kana={こっきー\tl{˥}},ipa={kokkiː},pos={名},acc={H}]}%
	{%
		\MeaningsBlock{%
	\ryumeaning%
		{ご馳走}{}{}%
	}%
	}%
	{\tailPrint[]}%
\entry
	{\headPrint[kana={こっこー\tl{˥}},ipa={kokkoː},pos={名},acc={H}]}%
	{%
		\MeaningsBlock{%
	\ryumeaning%
		{法事。法要}{}{}%
	}%
	}%
	{\tailPrint[]}%
\entry
	{\headPrint[kana={こっち},ipa={kottʃi},pos={名}]}%
	{%
		\MeaningsBlock{%
	\ryumeaning%
		{（幼児語）陰茎}{}{}%
	}%
	}%
	{\tailPrint[]}%
\entry
	{\headPrint[kana={こっぱるん\tl{˥}},ipa={kopparuɴ},pos={動},rui={1},para={こっぱらぬ},acc={H}]}%
	{%
		\MeaningsBlock{%
	\ryumeaning%
		{強張る。意地を張る}{}{}%
	\ryumeaning%
		{固くなる。硬直する}{}{}%
	}%
	}%
	{\tailPrint[]}%
\entry
	{\headPrint[kana={こっふぁすん\tl{˥˩}},ipa={koffasuɴ},pos={動},rui={2},para={こっふぁはぬ},acc={F}]}%
	{%
		\MeaningsBlock{%
	\ryumeaning%
		{崩す。壊す}{}{}%
	}%
	}%
	{\tailPrint[]}%
\entry
	{\headPrint[kana={こっふぃるん\tl{˥˩}},ipa={koffiruɴ},pos={動},rui={3},para={こっふぬ},acc={F}]}%
	{%
		\MeaningsBlock{%
	\ryumeaning%
		{崩れる}{}{}%
	}%
	}%
	{\tailPrint[]}%
\entry
	{\headPrint[kana={こむん\tl{˥}},ipa={komuɴ},pos={動},rui={1},para={こまぬ},acc={H}]}%
	{%
		\MeaningsBlock{%
	\ryumeaning%
		{お辞儀する}{}{}%
	}%
	}%
	{\tailPrint[]}%
\entry
	{\headPrint[kana={こんぎ\tl{˥}},ipa={koŋgi},pos={名},acc={H}]}%
	{%
		\MeaningsBlock{%
	\ryumeaning%
		{桑の木}{植物名}{}%
	}%
	}%
	{\tailPrint[]}%
\entry
	{\headPrint[kana={こんぎ\tl{˥}},ipa={koŋgi},pos={名},acc={H}]}%
	{%
		\MeaningsBlock{%
	\ryumeaning%
		{狂言}{}{}%
	}%
	}%
	{\tailPrint[]}%
\entry
	{\headPrint[kana={ごんた},ipa={gonta},pos={擬}]}%
	{%
		\MeaningsBlock{%
	\ryumeaning%
		{ごくんごくん}{水など、液体を勢い良く飲み下ろすさま}{}%
	}%
	}%
	{\tailPrint[]}%
\LetterHead{さ}
\entry
	{\headPrint[kana={さ},ipa={sa},pos={接尾}]}%
	{%
		\MeaningsBlock{%
	\ryumeaning%
		{～さ}{}{}%
	}%
	}%
	{\tailPrint[]}%
\entry
	{\headPrint[kana={さ\tl{˥}},ipa={sa},pos={名},acc={H}]}%
	{%
		\MeaningsBlock{%
	\ryumeaning%
		{茶}{}{}%
	}%
	}%
	{\tailPrint[]}%
\entry
	{\headPrint[kana={さ\tl{˥}},ipa={sa},pos={名},acc={H}]}%
	{%
		\MeaningsBlock{%
	\ryumeaning%
		{差}{}{}%
	}%
	}%
	{\tailPrint[]}%
\entry
	{\headPrint[kana={ざー\tl{˩˥}},ipa={dzaː},pos={名},acc={R}]}%
	{%
		\MeaningsBlock{%
	\ryumeaning%
		{\ruby[g]{何処}{どこ}。どれ}{}{}%
	}%
	}%
	{\tailPrint[]}%
\entry
	{\headPrint[kana={ざー\tl{˩˥}},ipa={dzaː},pos={名},acc={R}]}%
	{%
		\MeaningsBlock{%
	\ryumeaning%
		{座。席}{}{}%
	}%
	}%
	{\tailPrint[]}%
\entry
	{\headPrint[kana={さーだが\tl{˥˩}はん},ipa={saːdagahaɴ},pos={形},acc={F}]}%
	{%
		\MeaningsBlock{%
	\ryumeaning%
		{霊感能力が高い。霊験あらたかだ}{}{}%
	}%
	}%
	{\tailPrint[]}%
\entry
	{\headPrint[kana={ざーだぎ\tl{˩˥}},ipa={dzaːdagi},pos={名},acc={R}]}%
	{%
		\MeaningsBlock{%
	\ryumeaning%
		{苦竹}{竹の種類}{}%
	}%
	}%
	{\tailPrint[]}%
\entry
	{\headPrint[kana={さーっさーった},ipa={saːssaːtta},pos={擬}]}%
	{%
		\MeaningsBlock{%
	\ryumeaning%
		{さっさと勢いよく歩く様子}{}{}%
	}%
	}%
	{\tailPrint[]}%
\entry
	{\headPrint[kana={さーった},ipa={saːtta},pos={擬}]}%
	{%
		\MeaningsBlock{%
	\ryumeaning%
		{さっさと。素早く}{}{}%
	}%
	}%
	{\tailPrint[]}%
\entry
	{\headPrint[kana={さーに\tl{˥}},ipa={saːni},pos={名},acc={H}]}%
	{%
		\MeaningsBlock{%
	\ryumeaning%
		{\ruby[g]{月桃}{げっとう}}{植物名}{}%
	}%
	}%
	{\tailPrint[]}%
\entry
	{\headPrint[kana={ざーぬ\tl{˩˥}},ipa={dzaːnu},pos={句},acc={R}]}%
	{%
		\MeaningsBlock{%
	\ryumeaning%
		{どの。\ruby[g]{何処}{どこ}の}{}{}%
	}%
	}%
	{\tailPrint[]}%
\entry
	{\headPrint[kana={ざー\tl{˩˥}はん},ipa={dzaːhaɴ},pos={形},para={ざーへぬ},acc={R}]}%
	{%
		\MeaningsBlock{%
	\ryumeaning%
		{苦い}{}{}%
	}%
	}%
	{\tailPrint[]}%
\entry
	{\headPrint[kana={さーふーふー\ws{}すん\tl{˥˩}},ipa={saːfuːfuː suɴ},pos={句},acc={HL F}]}%
	{%
		\MeaningsBlock{%
	\ryumeaning%
		{ほろ酔いになる。上機嫌になる}{}{}%
	}%
	}%
	{\tailPrint[]}%
\entry
	{\headPrint[kana={ざーふぇー\tl{˩˥}},ipa={dzaːfeː},pos={名},acc={R}]}%
	{%
		\MeaningsBlock{%
	\ryumeaning%
		{まずいこと。困ったこと}{}{}%
	}%
	}%
	{\tailPrint[]}%
\entry
	{\headPrint[kana={さーら\tl{˥}},ipa={saːra},pos={名},acc={H}]}%
	{%
		\MeaningsBlock{%
	\ryumeaning%
		{サワラ}{魚類名}{}%
	}%
	}%
	{\tailPrint[]}%
\entry
	{\headPrint[kana={ざーりるん\tl{˩˥}},ipa={dzaːriruɴ},pos={動},rui={3},para={ざーるぬ},acc={R}]}%
	{%
		\MeaningsBlock{%
	\ryumeaning%
		{（布などが古くなって）裂ける}{}{}%
	}%
	}%
	{\tailPrint[]}%
\entry
	{\headPrint[kana={さーる\tl{˥}},ipa={saːru},pos={名},acc={H}]}%
	{%
		\MeaningsBlock{%
	\ryumeaning%
		{猿}{動物名}{}%
	}%
	}%
	{\tailPrint[]}%
\entry
	{\headPrint[kana={さーるん\tl{˥˩}},ipa={saːruɴ},pos={動},rui={1},para={さーらぬ},acc={F}]}%
	{%
		\MeaningsBlock{%
	\ryumeaning%
		{触る。触れる}{}{}%
	}%
	}%
	{\tailPrint[]}%
\entry
	{\headPrint[kana={ざいぎ\tl{˩˥}},ipa={dzaigi},pos={名},acc={R}]}%
	{%
		\MeaningsBlock{%
	\ryumeaning%
		{木材。材木}{}{}%
	}%
	}%
	{\tailPrint[]}%
\entry
	{\headPrint[kana={ざいさん\tl{˥˩}},ipa={dzaisaɴ},pos={名},acc={F}]}%
	{%
		\MeaningsBlock{%
	\ryumeaning%
		{財産}{}{}%
	}%
	}%
	{\tailPrint[]}%
\entry
	{\headPrint[kana={さうき\tl{˥}},ipa={sauki},pos={名},acc={H}]}%
	{%
		\MeaningsBlock{%
	\ryumeaning%
		{お茶うけ。茶菓子}{}{}%
	}%
	}%
	{\tailPrint[]}%
\entry
	{\headPrint[kana={さか\tl{˥˩}},ipa={sḁka},pos={名},acc={F}]}%
	{%
		\MeaningsBlock{%
	\ryumeaning%
		{上り坂}{「下り坂」は\emb{さんがり}と言う}{}%
	}%
	}%
	{\tailPrint[]}%
\entry
	{\headPrint[kana={さかいるん\tl{˥˩}},ipa={sḁkairuɴ},pos={動},rui={3},para={さかゆぬ},acc={F}]}%
	{%
		\MeaningsBlock{%
	\ryumeaning%
		{栄える。繁栄する}{}{}%
	\ryumeaning%
		{繁る}{}{}%
	}%
	}%
	{\tailPrint[]}%
\entry
	{\headPrint[kana={さかさー},ipa={sḁkasaː},pos={名}]}%
	{%
		\MeaningsBlock{%
	\ryumeaning%
		{逆さ。逆}{}{}%
	}%
	}%
	{\tailPrint[]}%
\entry
	{\headPrint[kana={さかすけー\tl{˥˩}},ipa={sḁkasu̥keː},pos={名},acc={F}]}%
	{%
		\MeaningsBlock{%
	\ryumeaning%
		{盃}{}{}%
	}%
	}%
	{\tailPrint[]}%
\entry
	{\headPrint[kana={さかなやー\tl{˥˩}},ipa={sḁkanajaː},pos={名},acc={F}]}%
	{%
		\MeaningsBlock{%
	\ryumeaning%
		{料亭。遊廓}{}{}%
	}%
	}%
	{\tailPrint[]}%
\entry
	{\headPrint[kana={ざがふなぶ\tl{˩˥}},ipa={dzagafu̥nabu},pos={名},acc={R}]}%
	{%
		\MeaningsBlock{%
	\ryumeaning%
		{ヒラミレモン}{ミカン名}{}%
	}%
	}%
	{\tailPrint[]}%
\entry
	{\headPrint[kana={さがらすん\tl{˥}},ipa={sagarasuɴ},pos={動},rui={2},para={さがらはぬ},acc={H}]}%
	{%
		\MeaningsBlock{%
	\ryumeaning%
		{掛け買いする}{掛けで買う}{}%
	}%
	}%
	{\tailPrint[]}%
\entry
	{\headPrint[kana={さがるん\tl{˥}},ipa={sagaruɴ},pos={動},rui={1},para={さがらぬ},acc={H}]}%
	{%
		\MeaningsBlock{%
	\ryumeaning%
		{（肉類が）傷む}{}{}%
	}%
	}%
	{\tailPrint[]}%
\entry
	{\headPrint[kana={さがるん\tl{˥}},ipa={sagaruɴ},pos={動},rui={1},para={さんがらぬ},acc={H}]}%
	{%
		\MeaningsBlock{%
	\ryumeaning%
		{下がる。後退する。ぶら下がる}{}{}%
	}%
	}%
	{\tailPrint[]}%
\entry
	{\headPrint[kana={さき\tl{˥˩}},ipa={sḁki},pos={名},acc={F}]}%
	{%
		\MeaningsBlock{%
	\ryumeaning%
		{先。先端。\ruby[g]{岬}{みさき}。\ruby[g]{埼}{さき}}{}{}%
	}%
	}%
	{\tailPrint[]}%
\entry
	{\headPrint[kana={さき\tl{˥˩}},ipa={sḁki},pos={名},acc={F}]}%
	{%
		\MeaningsBlock{%
	\ryumeaning%
		{酒}{}{}%
	}%
	}%
	{\tailPrint[]}%
\entry
	{\headPrint[kana={さきぐし\tl{˥˩}},ipa={sḁkiguʃi},pos={名},acc={F}]}%
	{%
		\MeaningsBlock{%
	\ryumeaning%
		{酒癖}{}{}%
	}%
	}%
	{\tailPrint[]}%
\entry
	{\headPrint[kana={さきじょーぐ\tl{˥˩}},ipa={sḁkidʒoːgu},pos={名},acc={F}]}%
	{%
		\MeaningsBlock{%
	\ryumeaning%
		{大酒飲み。酒上戸}{}{}%
	}%
	}%
	{\tailPrint[]}%
\entry
	{\headPrint[kana={さきつぁーるん\tl{˥˩}},ipa={sḁkitsaːruɴ},pos={動},acc={F}]}%
	{%
		\MeaningsBlock{%
	\ryumeaning%
		{ずたずたに裂く}{}{}%
	}%
	}%
	{\tailPrint[]}%
\entry
	{\headPrint[kana={さきふつぁ\tl{˥˩}はーん},ipa={sḁkifu̥tsahaːɴ},pos={形},acc={F}]}%
	{%
		\MeaningsBlock{%
	\ryumeaning%
		{酒臭い}{}{}%
	}%
	}%
	{\tailPrint[]}%
\entry
	{\headPrint[kana={さきまーり\tl{˥˩}},ipa={sḁkimaːri},pos={名},acc={F}]}%
	{%
		\MeaningsBlock{%
	\ryumeaning%
		{先回り}{}{}%
	}%
	}%
	{\tailPrint[]}%
\entry
	{\headPrint[kana={さきやま\tl{˥}},ipa={sḁkijama},pos={名},acc={H}]}%
	{%
		\MeaningsBlock{%
	\ryumeaning%
		{崎山}{西表島の地名}{}%
	}%
	}%
	{\tailPrint[]}%
\entry
	{\headPrint[kana={さきるん\tl{˥}},ipa={sḁkiruɴ},pos={動},rui={3},para={さくぬ},acc={H}]}%
	{%
		\MeaningsBlock{%
	\ryumeaning%
		{裂ける}{}{}%
	}%
	}%
	{\tailPrint[]}%
\entry
	{\headPrint[kana={さぎるん\tl{˥}},ipa={sagiruɴ},pos={動},rui={1},acc={H}]}%
	{%
		\MeaningsBlock{%
	\ryumeaning%
		{下痢する。通じが良くなる}{}{}%
	}%
	}%
	{\tailPrint[]}%
\entry
	{\headPrint[kana={さくほー\tl{˥}},ipa={sḁkuhoː},pos={名},acc={H}]}%
	{%
		\MeaningsBlock{%
	\ryumeaning%
		{農作。耕作}{}{}%
	}%
	}%
	{\tailPrint[]}%
\entry
	{\headPrint[kana={さくるん\tl{˥}},ipa={sḁkuruɴ},pos={動},rui={1},para={さくらぬ},acc={H}]}%
	{%
		\MeaningsBlock{%
	\ryumeaning%
		{切り開く}{肉を切り開くときに言う}{}%
	}%
	}%
	{\tailPrint[]}%
\entry
	{\headPrint[kana={さぐるん\tl{˥˩}},ipa={saguruɴ},pos={動},rui={1},para={さくらぬ},acc={F}]}%
	{%
		\MeaningsBlock{%
	\ryumeaning%
		{探る。探索する}{}{}%
	}%
	}%
	{\tailPrint[]}%
\entry
	{\headPrint[kana={さくん\tl{˥}},ipa={sḁkuɴ},pos={動},rui={1},para={さかぬ},acc={H}]}%
	{%
		\MeaningsBlock{%
	\ryumeaning%
		{裂く。割る}{}{}%
	}%
	}%
	{\tailPrint[]}%
\entry
	{\headPrint[kana={さくん\tl{˥˩}},ipa={sḁkuɴ},pos={動},rui={1},para={さかぬ},acc={F}]}%
	{%
		\MeaningsBlock{%
	\ryumeaning%
		{咲く}{}{}%
	}%
	}%
	{\tailPrint[]}%
\entry
	{\headPrint[kana={さけー\tl{˥˩}},ipa={sḁkeː},pos={名},acc={F}]}%
	{%
		\MeaningsBlock{%
	\ryumeaning%
		{境。境界}{}{}%
	}%
	}%
	{\tailPrint[]}%
\entry
	{\headPrint[kana={さこー\tl{˥}},ipa={sḁkoː},pos={名},acc={H}]}%
	{%
		\MeaningsBlock{%
	\ryumeaning%
		{咳}{}{}%
	}%
	}%
	{\tailPrint[]}%
\entry
	{\headPrint[kana={さこーし\tl{˥}},ipa={sḁkoːʃi},pos={名},acc={H}]}%
	{%
		\MeaningsBlock{%
	\ryumeaning%
		{長男}{}{}%
	}%
	}%
	{\tailPrint[]}%
\entry
	{\headPrint[kana={さこだち\tl{˥˩}},ipa={sḁkodatʃi},pos={名},acc={F}]}%
	{%
		\MeaningsBlock{%
	\ryumeaning%
		{ハスノハギリ}{樹木名}{}%
	}%
	}%
	{\tailPrint[]}%
\entry
	{\headPrint[kana={さごち},ipa={sagotʃi},pos={名}]}%
	{%
		\MeaningsBlock{%
	\ryumeaning%
		{正月}{}{}%
	}%
	}%
	{\tailPrint[]}%
\entry
	{\headPrint[kana={さこら\tl{˥}はん},ipa={sḁkorahaɴ},pos={形},para={さこらへぬ},acc={H}]}%
	{%
		\MeaningsBlock{%
	\ryumeaning%
		{塩辛い}{}{}%
	}%
	}%
	{\tailPrint[]}%
\entry
	{\headPrint[kana={ざしき\tl{˩˥}},ipa={dzaʃi̥ki},pos={名},acc={R}]}%
	{%
		\MeaningsBlock{%
	\ryumeaning%
		{座敷}{}{}%
	}%
	}%
	{\tailPrint[]}%
\entry
	{\headPrint[kana={さすん\tl{˥}},ipa={sḁsuɴ},pos={動},rui={2b},para={ささぬ},acc={H}]}%
	{%
		\MeaningsBlock{%
	\ryumeaning%
		{刺す。差す。挿す}{}{}%
	}%
	}%
	{\tailPrint[]}%
\entry
	{\headPrint[kana={さた\tl{˥}},ipa={sata},pos={名},acc={H}]}%
	{%
		\MeaningsBlock{%
	\ryumeaning%
		{噂。消息。評判}{}{}%
	}%
	}%
	{\tailPrint[]}%
\entry
	{\headPrint[kana={さた\tl{˥}},ipa={sḁta},pos={名},acc={H}]}%
	{%
		\MeaningsBlock{%
	\ryumeaning%
		{砂糖。黒糖}{}{}%
	}%
	}%
	{\tailPrint[]}%
\entry
	{\headPrint[kana={さたたくん},ipa={sḁtatḁkuɴ},pos={句},para={さたたかぬ}]}%
	{%
		\MeaningsBlock{%
	\ryumeaning%
		{製糖する}{}{}%
	}%
	}%
	{\tailPrint[]}%
\entry
	{\headPrint[kana={さたぱんべー\tl{˥}},ipa={sḁtapambeː},pos={名},acc={H}]}%
	{%
		\MeaningsBlock{%
	\ryumeaning%
		{砂糖てんぷら}{}{}%
	}%
	}%
	{\tailPrint[]}%
\entry
	{\headPrint[kana={さだみ\tl{˥}},ipa={sadami},pos={名},acc={H}]}%
	{%
		\MeaningsBlock{%
	\ryumeaning%
		{定め。規則}{}{}%
	}%
	}%
	{\tailPrint[]}%
\entry
	{\headPrint[kana={さだみるん\tl{˥}},ipa={sadamiruɴ},pos={動},rui={3},para={さだむぬ},acc={H}]}%
	{%
		\MeaningsBlock{%
	\ryumeaning%
		{定める。決める}{}{}%
	}%
	}%
	{\tailPrint[]}%
\entry
	{\headPrint[kana={さち\tl{˥˩}},ipa={sḁtʃi},pos={名},acc={F}]}%
	{%
		\MeaningsBlock{%
	\ryumeaning%
		{鉢巻}{}{}%
	}%
	}%
	{\tailPrint[]}%
\entry
	{\headPrint[kana={さってぃむ},ipa={sattimu},pos={感}]}%
	{%
		\MeaningsBlock{%
	\ryumeaning%
		{大変だ。もう大変}{}{}%
	}%
	}%
	{\tailPrint[]}%
\entry
	{\headPrint[kana={\josho{}ざっ\kako{}と},ipa={dzatto},pos={副},acc={HL}]}%
	{%
		\MeaningsBlock{%
	\ryumeaning%
		{簡単に。大雑把に}{}{}%
	}%
	}%
	{\tailPrint[]}%
\entry
	{\headPrint[kana={さっと},ipa={satto},pos={副}]}%
	{%
		\MeaningsBlock{%
	\ryumeaning%
		{てきぱきと。手早く}{}{}%
	}%
	}%
	{\tailPrint[]}%
\entry
	{\headPrint[kana={ざっふぁた},ipa={dzaffata},pos={擬}]}%
	{%
		\MeaningsBlock{%
	\ryumeaning%
		{ぐさりと}{物を突き刺す音の形容}{}%
	\ryumeaning%
		{ざぶざぶ}{}{}%
	}%
	}%
	{\tailPrint[]}%
\entry
	{\headPrint[kana={ざっふぁった},ipa={dzaffatta},pos={擬}]}%
	{%
		\MeaningsBlock{%
	\ryumeaning%
		{ざぶっと}{水に勢いよくつけるさま}{}%
	}%
	}%
	{\tailPrint[]}%
\entry
	{\headPrint[kana={さっぷん\tl{˥}},ipa={sappuɴ},pos={名},acc={H}]}%
	{%
		\MeaningsBlock{%
	\ryumeaning%
		{石鹸}{外来語の「シャボン」から}{}%
	}%
	}%
	{\tailPrint[]}%
\entry
	{\headPrint[kana={さとー\tl{˥}},ipa={sḁtoː},pos={名},acc={H}]}%
	{%
		\MeaningsBlock{%
	\ryumeaning%
		{茶湯}{仏前に供えるお茶}{}%
	}%
	}%
	{\tailPrint[]}%
\entry
	{\headPrint[kana={さとぅるん\tl{˥}},ipa={saturuɴ},pos={動},rui={1},para={さとらぬ},acc={H}]}%
	{%
		\MeaningsBlock{%
	\ryumeaning%
		{悟る。気付く}{}{}%
	}%
	}%
	{\tailPrint[]}%
\entry
	{\headPrint[kana={ざとく\tl{˩˥}},ipa={dzato̥ku},pos={名},acc={R}]}%
	{%
		\MeaningsBlock{%
	\ryumeaning%
		{床の間}{一番座の神を祭る間}{}%
	}%
	}%
	{\tailPrint[]}%
\entry
	{\headPrint[kana={さとぬし\tl{˥}},ipa={sḁtonuʃi},pos={名},acc={H}]}%
	{%
		\MeaningsBlock{%
	\ryumeaning%
		{里之子}{士族の若い男}{}%
	}%
	}%
	{\tailPrint[]}%
\entry
	{\headPrint[kana={さな\tl{˥}},ipa={sḁna},pos={名},acc={H}]}%
	{%
		\MeaningsBlock{%
	\ryumeaning%
		{傘。洋傘}{}{}%
	}%
	}%
	{\tailPrint[]}%
\entry
	{\headPrint[kana={さにさに\ws{}しゃん},ipa={sanisani ʃaɴ},pos={句}]}%
	{%
		\MeaningsBlock{%
	\ryumeaning%
		{にこにこと笑う}{}{}%
	}%
	}%
	{\tailPrint[]}%
\entry
	{\headPrint[kana={さにしゃー\ws{}すん},ipa={sḁniʃaː suɴ},pos={句},acc={-}]}%
	{%
		\MeaningsBlock{%
	\ryumeaning%
		{嬉しがる。喜ぶ。楽しむ}{}{}%
	}%
	}%
	{\tailPrint[]}%
\entry
	{\headPrint[kana={さにしゃ\tl{˥˩}はん},ipa={sḁniʃahaɴ},pos={形},para={さにしゃへぬ},acc={F}]}%
	{%
		\MeaningsBlock{%
	\ryumeaning%
		{嬉しい}{}{}%
	}%
	}%
	{\tailPrint[]}%
\entry
	{\headPrint[kana={さにち\tl{˥˩}},ipa={s̥anitʃi},pos={名},acc={F}]}%
	{%
		\MeaningsBlock{%
	\ryumeaning%
		{旧暦の三月三日}{浜下りの日}{}%
	}%
	}%
	{\tailPrint[]}%
\entry
	{\headPrint[kana={ざぬ},ipa={dzanu},pos={句}]}%
	{%
		\MeaningsBlock{%
	\ryumeaning%
		{どの。どこの}{}{}%
	}%
	}%
	{\tailPrint[]}%
\entry
	{\headPrint[kana={さねー\tl{˥˩}},ipa={sḁneː},pos={名},acc={F}]}%
	{%
		\MeaningsBlock{%
	\ryumeaning%
		{\ruby[g]{褌}{ふんどし}}{}{}%
	}%
	}%
	{\tailPrint[]}%
\entry
	{\headPrint[kana={さぱ\tl{˥˩}},ipa={sḁpa},pos={名},acc={F}]}%
	{%
		\MeaningsBlock{%
	\ryumeaning%
		{\ruby[g]{鱶}{ふか}}{}{}%
	}%
	}%
	{\tailPrint[]}%
\entry
	{\headPrint[kana={さぱぁ\tl{˥}},ipa={sḁpa},pos={名},acc={H}]}%
	{%
		\MeaningsBlock{%
	\ryumeaning%
		{相撲}{}{}%
	}%
	}%
	{\tailPrint[]}%
\entry
	{\headPrint[kana={さばくん\tl{˥}},ipa={sabakuɴ},pos={動},rui={1},para={さばかぬ},acc={H}]}%
	{%
		\MeaningsBlock{%
	\ryumeaning%
		{さばく。調停する}{}{}%
	}%
	}%
	{\tailPrint[]}%
\entry
	{\headPrint[kana={さばくん\tl{˥}},ipa={sabakuɴ},pos={動},rui={1},para={さばがぬ},acc={H}]}%
	{%
		\MeaningsBlock{%
	\ryumeaning%
		{探す。探索する}{}{}%
	}%
	}%
	{\tailPrint[]}%
\entry
	{\headPrint[kana={さぱさぱ},ipa={sḁpasḁpa},pos={擬}]}%
	{%
		\MeaningsBlock{%
	\ryumeaning%
		{粘り気のない様}{}{}%
	}%
	}%
	{\tailPrint[]}%
\entry
	{\headPrint[kana={さばに\tl{˥}},ipa={sabani},pos={名},acc={H}]}%
	{%
		\MeaningsBlock{%
	\ryumeaning%
		{くり舟}{}{}%
	}%
	}%
	{\tailPrint[]}%
\entry
	{\headPrint[kana={さぱ\tl{˥˩}はん},ipa={sḁpahaɴ},pos={形},para={さぱへぬ},acc={F}]}%
	{%
		\MeaningsBlock{%
	\ryumeaning%
		{折れやすい。裂けやすい}{}{}%
	}%
	}%
	{\tailPrint[]}%
\entry
	{\headPrint[kana={さばん\tl{˥}},ipa={sabaɴ},pos={名},acc={H}]}%
	{%
		\MeaningsBlock{%
	\ryumeaning%
		{湯呑。茶碗}{}{}%
	}%
	}%
	{\tailPrint[]}%
\entry
	{\headPrint[kana={さぱん\tl{˥˩}},ipa={sḁpaɴ},pos={名},acc={F}]}%
	{%
		\MeaningsBlock{%
	\ryumeaning%
		{草履}{}{}%
	}%
	}%
	{\tailPrint[]}%
\entry
	{\headPrint[kana={さび\tl{˥}},ipa={sabi},pos={名},acc={H}]}%
	{%
		\MeaningsBlock{%
	\ryumeaning%
		{錆}{}{}%
	}%
	}%
	{\tailPrint[]}%
\entry
	{\headPrint[kana={さびすくん},ipa={sabisukuɴ},pos={動}]}%
	{%
		\MeaningsBlock{%
	\ryumeaning%
		{錆付く}{}{}%
	}%
	}%
	{\tailPrint[]}%
\entry
	{\headPrint[kana={さふくん\tl{˥˩}},ipa={sḁfu̥kuɴ},pos={動},rui={1},para={さふかぬ},acc={F}]}%
	{%
		\MeaningsBlock{%
	\ryumeaning%
		{引きずる}{}{}%
	}%
	}%
	{\tailPrint[]}%
\entry
	{\headPrint[kana={さぶら\tl{˥}},ipa={sabura},pos={名},acc={H}]}%
	{%
		\MeaningsBlock{%
	\ryumeaning%
		{\ruby[g]{法}{ほ}\ruby[g]{螺}{ら}貝}{}{}%
	}%
	}%
	{\tailPrint[]}%
\entry
	{\headPrint[kana={さぼーりん\tl{˥˩}},ipa={saboːriɴ},pos={動},rui={3},para={さぼーるぬ},acc={F}]}%
	{%
		\MeaningsBlock{%
	\ryumeaning%
		{寂れる。落ちぶれる}{}{}%
	}%
	}%
	{\tailPrint[]}%
\entry
	{\headPrint[kana={ざぼんた},ipa={dzabonta},pos={擬}]}%
	{%
		\MeaningsBlock{%
	\ryumeaning%
		{ざぶっと}{水の中へ物を入れる音の形容}{}%
	}%
	}%
	{\tailPrint[]}%
\entry
	{\headPrint[kana={さましきるん},ipa={sḁmaʃikiruɴ},pos={句},para={さますくぬ}]}%
	{%
		\MeaningsBlock{%
	\ryumeaning%
		{縛りつける}{}{}%
	}%
	}%
	{\tailPrint[]}%
\entry
	{\headPrint[kana={さますん\tl{˥}},ipa={sḁmasuɴ},pos={動},rui={2},para={さまはぬ},acc={H}]}%
	{%
		\MeaningsBlock{%
	\ryumeaning%
		{冷ます}{}{}%
	}%
	}%
	{\tailPrint[]}%
\entry
	{\headPrint[kana={さますん\tl{˥}},ipa={sḁmasuɴ},pos={動},acc={H}]}%
	{%
		\MeaningsBlock{%
	\ryumeaning%
		{覚ます}{}{}%
	}%
	}%
	{\tailPrint[]}%
\entry
	{\headPrint[kana={さまたぎるん\tl{˥˩}},ipa={samatagiruɴ},pos={動},rui={3},para={さまたぐぬ},acc={F}]}%
	{%
		\MeaningsBlock{%
	\ryumeaning%
		{妨げる。邪魔する}{}{}%
	}%
	}%
	{\tailPrint[]}%
\entry
	{\headPrint[kana={ざまどぅらすん\tl{˩˥}},ipa={dzamadurasuɴ},pos={動},acc={R}]}%
	{%
		\MeaningsBlock{%
	\ryumeaning%
		{迷わす。惑わす}{}{}%
	\ryumeaning%
		{化かす}{}{}%
	}%
	}%
	{\tailPrint[]}%
\entry
	{\headPrint[kana={ざまどぅるん\tl{˩˥}},ipa={dzamaduruɴ},pos={動},rui={1},para={ざまどぅらぬ},acc={R}]}%
	{%
		\MeaningsBlock{%
	\ryumeaning%
		{迷う。まごつく。うろたえる}{}{}%
	}%
	}%
	{\tailPrint[]}%
\entry
	{\headPrint[kana={さまるん\tl{˥}},ipa={sḁmaruɴ},pos={動},rui={1},para={さまらぬ},acc={H}]}%
	{%
		\MeaningsBlock{%
	\ryumeaning%
		{醒める。目覚める}{}{}%
	}%
	}%
	{\tailPrint[]}%
\entry
	{\headPrint[kana={さまるん\tl{˥}},ipa={sḁmaruɴ},pos={動},rui={1},para={さまらぬ},acc={H}]}%
	{%
		\MeaningsBlock{%
	\ryumeaning%
		{\ruby[g]{縛}{しば}る}{}{}%
	\ryumeaning%
		{束ねる}{}{}%
	}%
	}%
	{\tailPrint[]}%
\entry
	{\headPrint[kana={さみるん\tl{˥}},ipa={sḁmiruɴ},pos={動},acc={H}]}%
	{%
		\MeaningsBlock{%
	\ryumeaning%
		{覚める}{}{}%
	}%
	}%
	{\tailPrint[]}%
\entry
	{\headPrint[kana={さみるん\tl{˥}},ipa={sḁmiruɴ},pos={動},acc={H}]}%
	{%
		\MeaningsBlock{%
	\ryumeaning%
		{冷める}{}{}%
	}%
	}%
	{\tailPrint[]}%
\entry
	{\headPrint[kana={さゆ\tl{˥}},ipa={saju},pos={名},acc={H}]}%
	{%
		\MeaningsBlock{%
	\ryumeaning%
		{白湯}{}{}%
	}%
	}%
	{\tailPrint[]}%
\entry
	{\headPrint[kana={さらー\tl{˥˩}},ipa={sḁraː},pos={名},acc={F}]}%
	{%
		\MeaningsBlock{%
	\ryumeaning%
		{皿}{}{}%
	}%
	}%
	{\tailPrint[]}%
\entry
	{\headPrint[kana={ざらざーら\ws{}すん},ipa={dzaradzaːra suɴ},pos={句}]}%
	{%
		\MeaningsBlock{%
	\ryumeaning%
		{ざらざらする}{}{}%
	}%
	}%
	{\tailPrint[]}%
\entry
	{\headPrint[kana={さらさら},ipa={sarasara},pos={擬}]}%
	{%
		\MeaningsBlock{%
	\ryumeaning%
		{}{水気のないさらっとしたさま}{}%
	}%
	}%
	{\tailPrint[]}%
\entry
	{\headPrint[kana={ざらざら},ipa={dzaradzara},pos={擬}]}%
	{%
		\MeaningsBlock{%
	\ryumeaning%
		{（粒など小さな物が）零れる様子}{小さな物を移したり、こぼしたりする時の音の形容}{}%
	}%
	}%
	{\tailPrint[]}%
\entry
	{\headPrint[kana={さらすん\tl{˥˩}},ipa={sḁrasuɴ},pos={動},rui={2},para={さらはぬ},acc={F}]}%
	{%
		\MeaningsBlock{%
	\ryumeaning%
		{晒す。漂白する}{}{}%
	\ryumeaning%
		{（水に）晒す}{}{}%
	}%
	}%
	{\tailPrint[]}%
\entry
	{\headPrint[kana={さり\tl{˥}},ipa={sḁri},pos={名},acc={H}]}%
	{%
		\MeaningsBlock{%
	\ryumeaning%
		{\ruby[g]{申}{さる}}{十二支の申}{}%
	}%
	}%
	{\tailPrint[]}%
\entry
	{\headPrint[kana={ざりるん\tl{˩˥}},ipa={dzariruɴ},pos={動},rui={3},para={ざるぬ},acc={R}]}%
	{%
		\MeaningsBlock{%
	\ryumeaning%
		{裂ける。破れる}{}{}%
	}%
	}%
	{\tailPrint[]}%
\entry
	{\headPrint[kana={さわがすん\tl{˥˩}},ipa={sawagasuɴ},pos={動},rui={2},para={さわがはぬ},acc={F}]}%
	{%
		\MeaningsBlock{%
	\ryumeaning%
		{騒がす。慌てさせる}{}{}%
	}%
	}%
	{\tailPrint[]}%
\entry
	{\headPrint[kana={さん\tl{˥}},ipa={saɴ},pos={名},acc={H}]}%
	{%
		\MeaningsBlock{%
	\ryumeaning%
		{桟}{戸や障子の骨}{}%
	}%
	}%
	{\tailPrint[]}%
\entry
	{\headPrint[kana={さん\tl{˥}},ipa={saɴ},pos={名},acc={H}]}%
	{%
		\MeaningsBlock{%
	\ryumeaning%
		{虱}{}{}%
	}%
	}%
	{\tailPrint[]}%
\entry
	{\headPrint[kana={さん\tl{˥}},ipa={saɴ},pos={名},acc={H}]}%
	{%
		\MeaningsBlock{%
	\ryumeaning%
		{魔除けの結び}{}{}%
	}%
	}%
	{\tailPrint[]}%
\entry
	{\headPrint[kana={ざん\tl{˩˥}},ipa={dzan},pos={名},acc={R}]}%
	{%
		\MeaningsBlock{%
	\ryumeaning%
		{ジュゴン}{海獣名}{}%
	}%
	}%
	{\tailPrint[]}%
\entry
	{\headPrint[kana={さんがり\tl{˥}},ipa={saŋgari},pos={名},acc={H}]}%
	{%
		\MeaningsBlock{%
	\ryumeaning%
		{下り坂}{}{}%
	}%
	}%
	{\tailPrint[]}%
\entry
	{\headPrint[kana={さんがるん\tl{˥}},ipa={saŋgaruɴ},pos={動},acc={H}]}%
	{%
		\MeaningsBlock{%
	\ryumeaning%
		{ぶら下がる}{}{}%
	}%
	}%
	{\tailPrint[]}%
\entry
	{\headPrint[kana={さんきら\tl{˥}},ipa={sankira},pos={名},acc={H}]}%
	{%
		\MeaningsBlock{%
	\ryumeaning%
		{山帰来}{薬草}{}%
	}%
	}%
	{\tailPrint[]}%
\entry
	{\headPrint[kana={さんぎるん\tl{˥}},ipa={saŋgiruɴ},pos={動},rui={3},para={さんぐぬ},acc={H}]}%
	{%
		\MeaningsBlock{%
	\ryumeaning%
		{吊るす。吊下げる}{}{}%
	}%
	}%
	{\tailPrint[]}%
\entry
	{\headPrint[kana={さんぎんそー\tl{˥}},ipa={saŋginsoː},pos={名},acc={H}]}%
	{%
		\MeaningsBlock{%
	\ryumeaning%
		{易者。占い師}{}{}%
	}%
	}%
	{\tailPrint[]}%
\entry
	{\headPrint[kana={さんごなー\tl{˥}},ipa={saŋgonaː},pos={名},acc={H}]}%
	{%
		\MeaningsBlock{%
	\ryumeaning%
		{浮気女。尻軽女}{}{}%
	}%
	}%
	{\tailPrint[]}%
\entry
	{\headPrint[kana={さんごなー\tl{˥}\ws{}すん\tl{˥˩}},ipa={saŋgonaː suɴ},pos={句},acc={H F}]}%
	{%
		\MeaningsBlock{%
	\ryumeaning%
		{淫らだ}{性的にだらしのない}{}%
	}%
	}%
	{\tailPrint[]}%
\entry
	{\headPrint[kana={さんざらごーなるん},ipa={sandzaragoː naruɴ},pos={名}]}%
	{%
		\MeaningsBlock{%
	\ryumeaning%
		{ばらばらに割れる}{}{}%
	}%
	}%
	{\tailPrint[]}%
\entry
	{\headPrint[kana={さんしん\tl{˥˩}},ipa={sanʃiɴ},pos={名},acc={F}]}%
	{%
		\MeaningsBlock{%
	\ryumeaning%
		{\ruby[g]{三}{さん}\ruby[g]{線}{しん}}{沖縄の三味線の呼称}{}%
	}%
	}%
	{\tailPrint[]}%
\entry
	{\headPrint[kana={さんず\tl{˥}},ipa={sandzu},pos={名},acc={H}]}%
	{%
		\MeaningsBlock{%
	\ryumeaning%
		{三十}{}{}%
	}%
	}%
	{\tailPrint[]}%
\entry
	{\headPrint[kana={さんだん},ipa={sandaɴ},pos={名}]}%
	{%
		\MeaningsBlock{%
	\ryumeaning%
		{工面}{}{}%
	}%
	}%
	{\tailPrint[]}%
\entry
	{\headPrint[kana={さんだん\ws{}すん},ipa={sandaɴ suɴ},pos={句}]}%
	{%
		\MeaningsBlock{%
	\ryumeaning%
		{工面する}{}{}%
	}%
	}%
	{\tailPrint[]}%
\entry
	{\headPrint[kana={ざんぬ\tl{˩˥}ゆ\tl{˥˩}},ipa={dzannuju},pos={名},acc={R+F}]}%
	{%
		\MeaningsBlock{%
	\ryumeaning%
		{ジュゴン}{海獣名}{}%
	}%
	}%
	{\tailPrint[]}%
\entry
	{\headPrint[kana={\josho{}さん\kako{}ば},ipa={samba},pos={名},acc={H+L}]}%
	{%
		\MeaningsBlock{%
	\ryumeaning%
		{南西方}{}{}%
	}%
	}%
	{\tailPrint[]}%
\entry
	{\headPrint[kana={\josho{}さん\kako{}ぼー},ipa={samboː},pos={名},acc={H+L}]}%
	{%
		\MeaningsBlock{%
	\ryumeaning%
		{三方}{お供えの台}{}%
	}%
	}%
	{\tailPrint[]}%
\entry
	{\headPrint[kana={さんみん\tl{˥}},ipa={sammiɴ},pos={名},acc={H}]}%
	{%
		\MeaningsBlock{%
	\ryumeaning%
		{計算}{}{}%
	}%
	}%
	{\tailPrint[]}%
\LetterHead{し}
\entry
	{\headPrint[kana={しー\tl{˥}},ipa={ʃiː},pos={名},acc={H}]}%
	{%
		\MeaningsBlock{%
	\ryumeaning%
		{手}{}{}%
	}%
	}%
	{\tailPrint[]}%
\entry
	{\headPrint[kana={しー\tl{˥}},ipa={ʃiː},pos={名},acc={H}]}%
	{%
		\MeaningsBlock{%
	\ryumeaning%
		{巣。鳥の巣}{}{}%
	}%
	}%
	{\tailPrint[]}%
\entry
	{\headPrint[kana={しー\tl{˥}あらはーん},ipa={ʃiːarahaːɴ},pos={形},acc={H+F}]}%
	{%
		\MeaningsBlock{%
	\ryumeaning%
		{手荒い。手が粗っぽい}{手ですることが粗削りである}{}%
	}%
	}%
	{\tailPrint[]}%
\entry
	{\headPrint[kana={しー\tl{˥}うさぎるん\tl{˥˩}},ipa={ʃiː.usagiruɴ},pos={句},para={しぃうさぐぬ},acc={H F}]}%
	{%
		\MeaningsBlock{%
	\ryumeaning%
		{手を合わせる。礼拝する}{}{}%
	}%
	}%
	{\tailPrint[]}%
\entry
	{\headPrint[kana={しー\tl{˥}うすくま\tl{˥˩}},ipa={ʃiː.usu̥kuma},pos={句},para={しぃうすくむぬ},acc={H F}]}%
	{%
		\MeaningsBlock{%
	\ryumeaning%
		{仕事を完了する}{}{}%
	}%
	}%
	{\tailPrint[]}%
\entry
	{\headPrint[kana={しー\ws{}うすくまん},ipa={ʃiː usu̥kumaɴ},pos={動},rui={3},para={しー　うすくむぬ}]}%
	{%
		\MeaningsBlock{%
	\ryumeaning%
		{終える。し遂げる}{}{}%
	}%
	}%
	{\tailPrint[]}%
\entry
	{\headPrint[kana={じーかたうやぐ},ipa={dʒiːkḁta.ujagu},pos={名}]}%
	{%
		\MeaningsBlock{%
	\ryumeaning%
		{母方親戚}{}{}%
	}%
	}%
	{\tailPrint[]}%
\entry
	{\headPrint[kana={しー\tl{˥}くぱ\tl{˥}はーん},ipa={ʃiːku̥pahaːɴ},pos={形},acc={H+H}]}%
	{%
		\MeaningsBlock{%
	\ryumeaning%
		{不器用だ}{}{}%
	}%
	}%
	{\tailPrint[]}%
\entry
	{\headPrint[kana={しー\tl{˥}ぐま\tl{˩˥}はん},ipa={ʃiːgumahaɴ},pos={形},acc={H+R}]}%
	{%
		\MeaningsBlock{%
	\ryumeaning%
		{器用だ}{仕事の出来がきれいだ}{}%
	}%
	}%
	{\tailPrint[]}%
\entry
	{\headPrint[kana={しーし\tl{˥}},ipa={ʃiːʃi},pos={名},acc={H}]}%
	{%
		\MeaningsBlock{%
	\ryumeaning%
		{獅子。ライオン}{}{}%
	}%
	}%
	{\tailPrint[]}%
\entry
	{\headPrint[kana={しーっぺー\tl{˥}},ipa={ʃiːppeː},pos={副},acc={H}]}%
	{%
		\MeaningsBlock{%
	\ryumeaning%
		{精一杯}{}{}%
	}%
	}%
	{\tailPrint[]}%
\entry
	{\headPrint[kana={しーなん\tl{˥}},ipa={ʃiːnaɴ},pos={名},acc={H}]}%
	{%
		\MeaningsBlock{%
	\ryumeaning%
		{帆綱}{}{}%
	}%
	}%
	{\tailPrint[]}%
\entry
	{\headPrint[kana={しーにふつはん},ipa={ʃiːnifu̥tsu̥haɴ},pos={形},acc={-}]}%
	{%
		\MeaningsBlock{%
	\ryumeaning%
		{手の仕事が遅い}{}{}%
	}%
	}%
	{\tailPrint[]}%
\entry
	{\headPrint[kana={しー\tl{˥˩}のーすん\tl{˩˥}},ipa={ʃiːnoːsuɴ},pos={動},rui={2},para={しーのーはぬ},acc={F R}]}%
	{%
		\MeaningsBlock{%
	\ryumeaning%
		{し直す}{}{}%
	}%
	}%
	{\tailPrint[]}%
\entry
	{\headPrint[kana={しーぱん\tl{˥}},ipa={ʃiːpaɴ},pos={名},acc={H}]}%
	{%
		\MeaningsBlock{%
	\ryumeaning%
		{手足}{}{}%
	}%
	}%
	{\tailPrint[]}%
\entry
	{\headPrint[kana={しーぶ\tl{˥}},ipa={ʃiːbu},pos={名},acc={H}]}%
	{%
		\MeaningsBlock{%
	\ryumeaning%
		{勝負}{}{}%
	}%
	}%
	{\tailPrint[note={勝負の転訛}]}%
\entry
	{\headPrint[kana={しーふき\tl{˥}},ipa={ʃiːfu̥ki},pos={名},acc={H}]}%
	{%
		\MeaningsBlock{%
	\ryumeaning%
		{指笛}{}{}%
	}%
	}%
	{\tailPrint[]}%
\entry
	{\headPrint[kana={しー\tl{˥}ぺしゃ\tl{˥}はん},ipa={ʃiːpeʃahaɴ},pos={形},acc={H+H}]}%
	{%
		\MeaningsBlock{%
	\ryumeaning%
		{手早い。手がすばしこい}{仕事の処理が早い}{}%
	\ryumeaning%
		{手がすぐ出る}{}{}%
	}%
	}%
	{\tailPrint[]}%
\entry
	{\headPrint[kana={しーるん\tl{˥}},ipa={ʃiːruɴ},pos={動},rui={1},para={しーらぬ},acc={H}]}%
	{%
		\MeaningsBlock{%
	\ryumeaning%
		{饐える}{食べ物が痛む}{}%
	}%
	}%
	{\tailPrint[]}%
\entry
	{\headPrint[kana={しあぎるん\tl{˥˩}},ipa={ʃi.agiruɴ},pos={動},rui={3},para={しあぐぬ},acc={F}]}%
	{%
		\MeaningsBlock{%
	\ryumeaning%
		{仕上げる}{}{}%
	}%
	}%
	{\tailPrint[]}%
\entry
	{\headPrint[kana={じぃ\tl{˥}},ipa={dʒi},pos={名},acc={H}]}%
	{%
		\MeaningsBlock{%
	\ryumeaning%
		{乳。乳房}{}{}%
	}%
	}%
	{\tailPrint[]}%
\entry
	{\headPrint[kana={じぃー\tl{˥˩}},ipa={dʒiː},pos={名},acc={F}]}%
	{%
		\MeaningsBlock{%
	\ryumeaning%
		{血。血液}{}{}%
	}%
	}%
	{\tailPrint[]}%
\entry
	{\headPrint[kana={じぃー\tl{˥˩}},ipa={dʒiː},pos={名},acc={F}]}%
	{%
		\MeaningsBlock{%
	\ryumeaning%
		{字。文字}{}{}%
	}%
	}%
	{\tailPrint[]}%
\entry
	{\headPrint[kana={じぃー\tl{˩˥}},ipa={dʒiː},pos={名},acc={R}]}%
	{%
		\MeaningsBlock{%
	\ryumeaning%
		{地。土地}{}{}%
	}%
	}%
	{\tailPrint[]}%
\entry
	{\headPrint[kana={じぃー\tl{˩˥}},ipa={dʒiː},pos={名},acc={R}]}%
	{%
		\MeaningsBlock{%
	\ryumeaning%
		{棘}{}{}%
	}%
	}%
	{\tailPrint[]}%
\entry
	{\headPrint[kana={じぃーまーみ\tl{˩˥}},ipa={dʒiːmaːmi},pos={名},acc={R}]}%
	{%
		\MeaningsBlock{%
	\ryumeaning%
		{落花生。ピーナツ}{}{}%
	}%
	}%
	{\tailPrint[]}%
\entry
	{\headPrint[kana={じぃぬ\ws{}みち\tl{˥˩}},ipa={dʒinu mitʃi},pos={名},acc={F}]}%
	{%
		\MeaningsBlock{%
	\ryumeaning%
		{血管}{}{}%
	}%
	}%
	{\tailPrint[]}%
\entry
	{\headPrint[kana={しぃぴりるん\tl{˥˩}},ipa={ʃi̥piriruɴ},pos={動},rui={3},para={しぃぴるぬ},acc={F}]}%
	{%
		\MeaningsBlock{%
	\ryumeaning%
		{萎びる。小さくなる}{}{}%
	}%
	}%
	{\tailPrint[]}%
\entry
	{\headPrint[kana={しぃぴるん\tl{˥˩}},ipa={ʃi̥piruɴ},pos={動},rui={1},para={しぃぴらぬ},acc={F}]}%
	{%
		\MeaningsBlock{%
	\ryumeaning%
		{\ruby[g]{舐}{な}める}{}{}%
	}%
	}%
	{\tailPrint[]}%
\entry
	{\headPrint[kana={しぃま\tl{˥}},ipa={ʃi̥ma},pos={名},acc={H}]}%
	{%
		\MeaningsBlock{%
	\ryumeaning%
		{手間。手間賃。報酬}{}{}%
	}%
	}%
	{\tailPrint[]}%
\entry
	{\headPrint[kana={しぃみるん\tl{˥}},ipa={ʃi̥miruɴ},pos={動},acc={H}]}%
	{%
		\MeaningsBlock{%
	\ryumeaning%
		{強める}{}{}%
	\ryumeaning%
		{（酒などを）濃くする}{}{}%
	}%
	}%
	{\tailPrint[]}%
\entry
	{\headPrint[kana={じうてー\tl{˩˥}},ipa={dʒi.uteː},pos={名},acc={R}]}%
	{%
		\MeaningsBlock{%
	\ryumeaning%
		{地謡}{}{}%
	}%
	}%
	{\tailPrint[]}%
\entry
	{\headPrint[kana={しか\tl{˥}},ipa={ʃi̥ka},pos={名},acc={H}]}%
	{%
		\MeaningsBlock{%
	\ryumeaning%
		{石垣市}{石垣市の古称。直訳「四箇」で、「四か村」の義}{}%
	}%
	}%
	{\tailPrint[]}%
\entry
	{\headPrint[kana={\josho{}しかい\kako{}とぅ},ipa={ʃi̥kaitu},pos={副},acc={HL}]}%
	{%
		\MeaningsBlock{%
	\ryumeaning%
		{しっかり。精一杯}{}{}%
	}%
	}%
	{\tailPrint[]}%
\entry
	{\headPrint[kana={しかき\tl{˥˩}},ipa={ʃi̥kaki},pos={名},acc={F}]}%
	{%
		\MeaningsBlock{%
	\ryumeaning%
		{仕掛け。工夫}{}{}%
	}%
	}%
	{\tailPrint[]}%
\entry
	{\headPrint[kana={しかきるん\tl{˥˩}},ipa={ʃi̥kakiruɴ},pos={動},rui={3},para={しかくぬ},acc={F}]}%
	{%
		\MeaningsBlock{%
	\ryumeaning%
		{仕掛ける}{}{}%
	\ryumeaning%
		{挑む}{}{}%
	}%
	}%
	{\tailPrint[]}%
\entry
	{\headPrint[kana={しかた\tl{˥˩}},ipa={ʃi̥kata},pos={名},acc={F}]}%
	{%
		\MeaningsBlock{%
	\ryumeaning%
		{仕方。方法}{}{}%
	}%
	}%
	{\tailPrint[]}%
\entry
	{\headPrint[kana={しがねー\tl{˥}},ipa={ʃiganeː},pos={名},acc={H}]}%
	{%
		\MeaningsBlock{%
	\ryumeaning%
		{手伝い。手助け}{}{}%
	}%
	}%
	{\tailPrint[]}%
\entry
	{\headPrint[kana={しかま\tl{˥}},ipa={ʃi̥kama},pos={名},acc={H}]}%
	{%
		\MeaningsBlock{%
	\ryumeaning%
		{仕事}{}{}%
	}%
	}%
	{\tailPrint[]}%
\entry
	{\headPrint[kana={\josho{}しかま\kako{}ぷち},ipa={ʃi̥kamapu̥tʃi},pos={名},acc={H+L}]}%
	{%
		\MeaningsBlock{%
	\ryumeaning%
		{宵の明星}{昔は星の明かりで野良仕事をしたことから}{}%
	}%
	}%
	{\tailPrint[]}%
\entry
	{\headPrint[kana={しがみ\tl{˥}},ipa={ʃigami},pos={名},acc={H}]}%
	{%
		\MeaningsBlock{%
	\ryumeaning%
		{手紙}{}{}%
	}%
	}%
	{\tailPrint[]}%
\entry
	{\headPrint[kana={しかむん\tl{˥}},ipa={ʃi̥kamuɴ},pos={動},acc={H}]}%
	{%
		\MeaningsBlock{%
	\ryumeaning%
		{しかむ。ひびる}{}{}%
	}%
	}%
	{\tailPrint[]}%
\entry
	{\headPrint[kana={しからさ\tl{˥}\ws{}すん\tl{˥˩}},ipa={ʃi̥karasa suɴ},pos={句},acc={H F}]}%
	{%
		\MeaningsBlock{%
	\ryumeaning%
		{悲しがる}{}{}%
	\ryumeaning%
		{寂しがる}{}{}%
	}%
	}%
	{\tailPrint[]}%
\entry
	{\headPrint[kana={しから\tl{˥}はーん},ipa={ʃi̥karahaːɴ},pos={形},acc={H}]}%
	{%
		\MeaningsBlock{%
	\ryumeaning%
		{悲しみふさぐ}{}{}%
	\ryumeaning%
		{寂しがる}{}{}%
	}%
	}%
	{\tailPrint[]}%
\entry
	{\headPrint[kana={しから\tl{˥}はん},ipa={ʃi̥karahaɴ},pos={形},para={しからへぬ},acc={H}]}%
	{%
		\MeaningsBlock{%
	\ryumeaning%
		{淋しい。悲しい}{}{}%
	}%
	}%
	{\tailPrint[]}%
\entry
	{\headPrint[kana={しきあたるん\tl{˥˩}},ipa={ʃi̥ki.ataruɴ},pos={動},rui={1},para={しきあたらぬ},acc={F}]}%
	{%
		\MeaningsBlock{%
	\ryumeaning%
		{突き当たる}{}{}%
	}%
	}%
	{\tailPrint[]}%
\entry
	{\headPrint[kana={しきあんぎるん\tl{˥˩}},ipa={ʃi̥ki.aŋgiruɴ},pos={動},rui={3},para={しきあんぐぬ},acc={F}]}%
	{%
		\MeaningsBlock{%
	\ryumeaning%
		{突き上げる}{}{}%
	}%
	}%
	{\tailPrint[]}%
\entry
	{\headPrint[kana={しきぐり\tl{˥˩}さーん},ipa={ʃi̥kigurisaːɴ},pos={形},acc={F}]}%
	{%
		\MeaningsBlock{%
	\ryumeaning%
		{聞き苦しい}{}{}%
	}%
	}%
	{\tailPrint[]}%
\entry
	{\headPrint[kana={しきだーすん\tl{˥˩}},ipa={ʃi̥kidaːsuɴ},pos={動},rui={2},acc={F}]}%
	{%
		\MeaningsBlock{%
	\ryumeaning%
		{べったり座る}{}{}%
	}%
	}%
	{\tailPrint[]}%
\entry
	{\headPrint[kana={しきだぎ\tl{˥}},ipa={ʃi̥kidagi},pos={名},acc={H}]}%
	{%
		\MeaningsBlock{%
	\ryumeaning%
		{\ruby[g]{燐寸}{マッチ}}{}{}%
	}%
	}%
	{\tailPrint[]}%
\entry
	{\headPrint[kana={しきだるん\tl{˥}},ipa={ʃi̥kidaruɴ},pos={動},rui={1},para={しきだらぬ},acc={H}]}%
	{%
		\MeaningsBlock{%
	\ryumeaning%
		{叩く。砕く}{}{}%
	}%
	}%
	{\tailPrint[]}%
\entry
	{\headPrint[kana={しきたん\tl{˥}},ipa={ʃi̥kitaɴ},pos={名},acc={H}]}%
	{%
		\MeaningsBlock{%
	\ryumeaning%
		{石炭}{}{}%
	}%
	}%
	{\tailPrint[note={石炭の転訛}]}%
\entry
	{\headPrint[kana={しきつぁーすん\tl{˥˩}},ipa={ʃi̥kitsaːsuɴ},pos={動},rui={2b},para={しきつぁーさぬ},acc={F}]}%
	{%
		\MeaningsBlock{%
	\ryumeaning%
		{にじり寄る。いざる}{}{}%
	}%
	}%
	{\tailPrint[]}%
\entry
	{\headPrint[kana={しきっつぁーるん},ipa={ʃi̥kittsaːruɴ},pos={動}]}%
	{%
		\MeaningsBlock{%
	\ryumeaning%
		{いざって進む。膝行する}{}{}%
	}%
	}%
	{\tailPrint[]}%
\entry
	{\headPrint[kana={しきとぅぱすん\tl{˥˩}},ipa={ʃi̥kitu̥pasuɴ},pos={動},rui={2},para={しきとぱはぬ},acc={F}]}%
	{%
		\MeaningsBlock{%
	\ryumeaning%
		{突き飛ばす}{}{}%
	}%
	}%
	{\tailPrint[]}%
\entry
	{\headPrint[kana={しきにんむつん},ipa={ʃikinimmutsuɴ}]}%
	{%
		\MeaningsBlock{%
	\ryumeaning%
		{責任を負う。責任を持つ}{}{}%
	}%
	}%
	{\tailPrint[]}%
\entry
	{\headPrint[kana={しきべー},ipa={ʃi̥kibeː},pos={名}]}%
	{%
		\MeaningsBlock{%
	\ryumeaning%
		{色気違い。多淫の女子。好色の女子}{}{}%
	}%
	}%
	{\tailPrint[]}%
\entry
	{\headPrint[kana={しきぽーるん\tl{˥˩}},ipa={ʃi̥kipoːruɴ},pos={動},rui={1},para={しきぽーらぬ}]}%
	{%
		\MeaningsBlock{%
	\ryumeaning%
		{散乱させる。散らかす}{}{}%
	}%
	}%
	{\tailPrint[]}%
\entry
	{\headPrint[kana={しきほーん\tl{˥˩}},ipa={ʃi̥kihoːɴ},pos={動},acc={F}]}%
	{%
		\MeaningsBlock{%
	\ryumeaning%
		{ついばむ}{}{}%
	}%
	}%
	{\tailPrint[]}%
\entry
	{\headPrint[kana={しきゆ\tl{˥}},ipa={ʃi̥kiju},pos={名},acc={H}]}%
	{%
		\MeaningsBlock{%
	\ryumeaning%
		{石油}{初めは\emb{しきたんゆ}と言っていた}{}%
	}%
	}%
	{\tailPrint[]}%
\entry
	{\headPrint[kana={しきり\tl{˥˩}},ipa={ʃi̥kiri},pos={名},acc={F}]}%
	{%
		\MeaningsBlock{%
	\ryumeaning%
		{\ruby[g]{海鼠}{なまこ}}{}{}%
	}%
	}%
	{\tailPrint[]}%
\entry
	{\headPrint[kana={しきるん\tl{˥˩}},ipa={ʃi̥kiruɴ},pos={動},rui={3},para={すくーぬ},acc={F}]}%
	{%
		\MeaningsBlock{%
	\ryumeaning%
		{漬ける}{}{}%
	}%
	}%
	{\tailPrint[]}%
\entry
	{\headPrint[kana={しきるん\tl{˥˩}},ipa={ʃi̥kiruɴ},pos={動},rui={3},para={しくぬ},acc={F}]}%
	{%
		\MeaningsBlock{%
	\ryumeaning%
		{お供えする}{}{}%
	}%
	}%
	{\tailPrint[]}%
\entry
	{\headPrint[kana={しきるん\tl{˥˩}},ipa={ʃi̥kiruɴ},pos={動},rui={3},acc={F}]}%
	{%
		\MeaningsBlock{%
	\ryumeaning%
		{（灯りや火を）点ける}{}{}%
	}%
	}%
	{\tailPrint[]}%
\entry
	{\headPrint[kana={しぎるん\tl{˥}},ipa={ʃigiruɴ},pos={動},rui={3},para={すぐぬ},acc={H}]}%
	{%
		\MeaningsBlock{%
	\ryumeaning%
		{過ぎ去る}{}{}%
	}%
	}%
	{\tailPrint[]}%
\entry
	{\headPrint[kana={しきん\tl{˥}},ipa={ʃi̥kiɴ},pos={名},acc={H}]}%
	{%
		\MeaningsBlock{%
	\ryumeaning%
		{世間。世の中}{}{}%
	}%
	}%
	{\tailPrint[]}%
\entry
	{\headPrint[kana={しき\tl{˥}\ws{}んぐん\tl{˥˩}},ipa={ʃi̥ki ŋguɴ},pos={句},para={しきんがぬ},acc={H F}]}%
	{%
		\MeaningsBlock{%
	\ryumeaning%
		{付いて行く}{}{}%
	}%
	}%
	{\tailPrint[]}%
\entry
	{\headPrint[kana={しき\tl{˥˩}んじるん\tl{˥}},ipa={ʃi̥kindʒiruɴ},pos={動},rui={3},para={しきんどぅぬ},acc={F H}]}%
	{%
		\MeaningsBlock{%
	\ryumeaning%
		{突き出る}{}{}%
	}%
	}%
	{\tailPrint[]}%
\entry
	{\headPrint[kana={しきんだすん},ipa={ʃi̥kindasuɴ},pos={動},rui={2},para={しきんだはぬ}]}%
	{%
		\MeaningsBlock{%
	\ryumeaning%
		{突き出す}{}{}%
	\ryumeaning%
		{突き出させる}{突き出るようにする}{}%
	}%
	}%
	{\tailPrint[]}%
\entry
	{\headPrint[kana={しぐ\tl{˥}},ipa={ʃigu},pos={名},acc={H}]}%
	{%
		\MeaningsBlock{%
	\ryumeaning%
		{小刀。ナイフ}{}{}%
	}%
	}%
	{\tailPrint[]}%
\entry
	{\headPrint[kana={しぐとぅ\tl{˥˩}},ipa={ʃigutu},pos={名},acc={F}]}%
	{%
		\MeaningsBlock{%
	\ryumeaning%
		{仕事。業務}{}{}%
	}%
	}%
	{\tailPrint[]}%
\entry
	{\headPrint[kana={しくむん\tl{˥˩}},ipa={ʃi̥kumuɴ},pos={動},rui={1},para={しくまぬ},acc={F}]}%
	{%
		\MeaningsBlock{%
	\ryumeaning%
		{仕込む。準備する}{}{}%
	}%
	}%
	{\tailPrint[]}%
\entry
	{\headPrint[kana={じぐや\tl{˩˥}},ipa={dʒiguja},pos={名},acc={R}]}%
	{%
		\MeaningsBlock{%
	\ryumeaning%
		{十五夜}{中秋の名月}{}%
	}%
	}%
	{\tailPrint[]}%
\entry
	{\headPrint[kana={しくん\tl{˥˩}},ipa={ʃi̥kuɴ},pos={動},rui={1},para={しかぬ},acc={F}]}%
	{%
		\MeaningsBlock{%
	\ryumeaning%
		{突く}{}{}%
	}%
	}%
	{\tailPrint[]}%
\entry
	{\headPrint[kana={しけーすん\tl{˥˩}},ipa={ʃi̥keːsuɴ},pos={動},rui={2},para={しけーさぬ},acc={F}]}%
	{%
		\MeaningsBlock{%
	\ryumeaning%
		{案内する。お迎えする}{}{}%
	}%
	}%
	{\tailPrint[]}%
\entry
	{\headPrint[kana={しけーるん\tl{˥˩}},ipa={ʃikeːruɴ},pos={動},acc={F}]}%
	{%
		\MeaningsBlock{%
	\ryumeaning%
		{支える}{}{}%
	}%
	}%
	{\tailPrint[]}%
\entry
	{\headPrint[kana={しけん\tl{˥}},ipa={ʃi̥keɴ},pos={名},acc={H}]}%
	{%
		\MeaningsBlock{%
	\ryumeaning%
		{（天体の）月}{}{}%
	}%
	}%
	{\tailPrint[]}%
\entry
	{\headPrint[kana={しけんゆー},ipa={ʃi̥keɴjuː},pos={名}]}%
	{%
		\MeaningsBlock{%
	\ryumeaning%
		{月夜}{}{}%
	}%
	}%
	{\tailPrint[]}%
\entry
	{\headPrint[kana={しこーるん\tl{˥˩}},ipa={ʃi̥koːruɴ},pos={動},rui={1},para={すこーらぬ},acc={F}]}%
	{%
		\MeaningsBlock{%
	\ryumeaning%
		{準備する。用意する}{}{}%
	\ryumeaning%
		{作る。拵える}{}{}%
	}%
	}%
	{\tailPrint[]}%
\entry
	{\headPrint[kana={しさ\tl{˥}はん},ipa={ʃi̥sahaɴ},pos={形},para={しさへぬ},acc={H}]}%
	{%
		\MeaningsBlock{%
	\ryumeaning%
		{\ruby[g]{酸}{す}っぱい}{}{}%
	}%
	}%
	{\tailPrint[]}%
\entry
	{\headPrint[kana={ししー\tl{˥}},ipa={ʃiʃiː},pos={名},acc={H}]}%
	{%
		\MeaningsBlock{%
	\ryumeaning%
		{煤}{}{}%
	}%
	}%
	{\tailPrint[]}%
\entry
	{\headPrint[kana={ししき\tl{˥}},ipa={ʃiʃi̥ki},pos={名},acc={H}]}%
	{%
		\MeaningsBlock{%
	\ryumeaning%
		{（頭上運搬の）当て物}{頭に乗せる輪。頭上運搬に使う丸い輪}{}%
	}%
	}%
	{\tailPrint[]}%
\entry
	{\headPrint[kana={ししきるん},ipa={ʃiʃi̥kiruɴ},pos={句},para={しぃすくぬ}]}%
	{%
		\MeaningsBlock{%
	\ryumeaning%
		{手がける。手をつける}{}{}%
	}%
	}%
	{\tailPrint[]}%
\entry
	{\headPrint[kana={しじゃー\tl{˥}},ipa={ʃidʒaː},pos={名},acc={H}]}%
	{%
		\MeaningsBlock{%
	\ryumeaning%
		{年上。年長}{}{}%
	}%
	}%
	{\tailPrint[]}%
\entry
	{\headPrint[kana={ししゃん\tl{˥˩}},ipa={ʃʃaɴ},pos={動（継）},rui={1},para={しさぬ},acc={F}]}%
	{%
		\MeaningsBlock{%
	\ryumeaning%
		{知っている}{}{}%
	}%
	}%
	{\tailPrint[]}%
\entry
	{\headPrint[kana={ししらすん\tl{˥}},ipa={ʃʃirasuɴ},pos={動},acc={H}]}%
	{%
		\MeaningsBlock{%
	\ryumeaning%
		{滑らせる}{}{}%
	}%
	}%
	{\tailPrint[]}%
\entry
	{\headPrint[kana={しじり\tl{˥}},ipa={ʃidʒiri},pos={名},acc={H}]}%
	{%
		\MeaningsBlock{%
	\ryumeaning%
		{\ruby[g]{硯}{すずり}}{}{}%
	}%
	}%
	{\tailPrint[]}%
\entry
	{\headPrint[kana={ししりん\tl{˥}},ipa={ʃʃiriɴ},pos={動},rui={3},para={ししぃるぬ},acc={H}]}%
	{%
		\MeaningsBlock{%
	\ryumeaning%
		{滑る}{}{}%
	}%
	}%
	{\tailPrint[]}%
\entry
	{\headPrint[kana={ししん\tl{˥}},ipa={ʃʃiɴ},pos={名},acc={H}]}%
	{%
		\MeaningsBlock{%
	\ryumeaning%
		{節祭}{}{}%
	}%
	}%
	{\tailPrint[]}%
\entry
	{\headPrint[kana={ししん\tl{˥˩}},ipa={ʃʃiɴ},pos={動},rui={3},acc={F}]}%
	{%
		\MeaningsBlock{%
	\ryumeaning%
		{捨てる}{}{}%
	}%
	}%
	{\tailPrint[]}%
\entry
	{\headPrint[kana={した\tl{˥}},ipa={ʃi̥ta},pos={名},acc={H}]}%
	{%
		\MeaningsBlock{%
	\ryumeaning%
		{舌}{}{}%
	}%
	}%
	{\tailPrint[]}%
\entry
	{\headPrint[kana={したい},ipa={ʃi̥tai},pos={感}]}%
	{%
		\MeaningsBlock{%
	\ryumeaning%
		{やったー。でかした}{}{}%
	}%
	}%
	{\tailPrint[]}%
\entry
	{\headPrint[kana={したいが\tl{˥˩}},ipa={ʃi̥taiga},pos={句},acc={F}]}%
	{%
		\MeaningsBlock{%
	\ryumeaning%
		{下へ。下方へ}{}{}%
	}%
	}%
	{\tailPrint[]}%
\entry
	{\headPrint[kana={したく\ws{}すん},ipa={si̥taku suɴ},pos={句},para={とぅりむたぬ}]}%
	{%
		\MeaningsBlock{%
	\ryumeaning%
		{支度する}{}{}%
	}%
	}%
	{\tailPrint[]}%
\entry
	{\headPrint[kana={したくぱるん\tl{˥}},ipa={ʃi̥taku̥paruɴ},pos={動},rui={1},para={したくぱらぬ},acc={H}]}%
	{%
		\MeaningsBlock{%
	\ryumeaning%
		{\ruby[g]{吃}{ども}る}{}{}%
	}%
	}%
	{\tailPrint[]}%
\entry
	{\headPrint[kana={しだすん\tl{˥˩}},ipa={ʃidasuɴ},pos={動},rui={2},para={しだはぬ},acc={F}]}%
	{%
		\MeaningsBlock{%
	\ryumeaning%
		{着飾る}{}{}%
	}%
	}%
	{\tailPrint[]}%
\entry
	{\headPrint[kana={したち\tl{˥}},ipa={ʃi̥tatʃi},pos={名},acc={H}]}%
	{%
		\MeaningsBlock{%
	\ryumeaning%
		{醤油}{}{}%
	}%
	}%
	{\tailPrint[]}%
\entry
	{\headPrint[kana={したてぃるん\tl{˥˩}},ipa={ʃi̥tatiruɴ},pos={動},acc={F}]}%
	{%
		\MeaningsBlock{%
	\ryumeaning%
		{仕立てる}{}{}%
	}%
	}%
	{\tailPrint[]}%
\entry
	{\headPrint[kana={しち\tl{˥}},ipa={ʃi̥tʃi},pos={名},acc={H}]}%
	{%
		\MeaningsBlock{%
	\ryumeaning%
		{\ruby[g]{節}{せつ}。季節}{}{}%
	}%
	}%
	{\tailPrint[]}%
\entry
	{\headPrint[kana={しっさん\tl{˥}},ipa={ssaɴ},pos={動（継）},rui={3},para={しっすぬ},acc={H}]}%
	{%
		\MeaningsBlock{%
	\ryumeaning%
		{絶やす}{}{}%
	}%
	}%
	{\tailPrint[]}%
\entry
	{\headPrint[kana={しっし\tl{˥˩}},ipa={ʃiʃʃi},pos={名},acc={F}]}%
	{%
		\MeaningsBlock{%
	\ryumeaning%
		{手拭い。タオル}{}{}%
	}%
	}%
	{\tailPrint[]}%
\entry
	{\headPrint[kana={しっしとぅるん\tl{˥}},ipa={ʃʃituruɴ},pos={動},rui={1},para={しっしとぅらぬ},acc={H}]}%
	{%
		\MeaningsBlock{%
	\ryumeaning%
		{切り取る}{}{}%
	}%
	}%
	{\tailPrint[]}%
\entry
	{\headPrint[kana={しっしりるん\tl{˥}},ipa={ʃʃiriruɴ},pos={動},rui={3},para={しっしるぬ},acc={H}]}%
	{%
		\MeaningsBlock{%
	\ryumeaning%
		{ずり落ちる}{}{}%
	\ryumeaning%
		{滑る}{}{}%
	}%
	}%
	{\tailPrint[]}%
\entry
	{\headPrint[kana={しっしるん\tl{˥}},ipa={ʃʃiruɴ},pos={動},acc={H}]}%
	{%
		\MeaningsBlock{%
	\ryumeaning%
		{絶える。途切れる}{}{}%
	}%
	}%
	{\tailPrint[]}%
\entry
	{\headPrint[kana={しっしんふに\tl{˥}},ipa={ʃʃinfu̥ni},pos={名},acc={H}]}%
	{%
		\MeaningsBlock{%
	\ryumeaning%
		{節祭で漕ぐくり舟}{}{}%
	}%
	}%
	{\tailPrint[]}%
\entry
	{\headPrint[kana={しっすするん\tl{˥˩}},ipa={ssuruɴ},pos={動},rui={1},para={しっすらぬ},acc={F}]}%
	{%
		\MeaningsBlock{%
	\ryumeaning%
		{\ruby[g]{擦}{す}る。\ruby[g]{擦}{こす}る}{}{}%
	}%
	}%
	{\tailPrint[]}%
\entry
	{\headPrint[kana={しっすぴかり\tl{˥}},ipa={ssupi̥kari},pos={名},acc={H}]}%
	{%
		\MeaningsBlock{%
	\ryumeaning%
		{\ruby[g]{蛍}{ほたる}}{}{}%
	}%
	}%
	{\tailPrint[]}%
\entry
	{\headPrint[kana={しっすみん\tl{˥}},ipa={ssumiɴ},pos={名},acc={H}]}%
	{%
		\MeaningsBlock{%
	\ryumeaning%
		{白いキノコ}{}{}%
	}%
	}%
	{\tailPrint[]}%
\entry
	{\headPrint[kana={しっするん\tl{˥˩}},ipa={ssuruɴ},pos={動},rui={1},para={しっすらぬ},acc={F}]}%
	{%
		\MeaningsBlock{%
	\ryumeaning%
		{拭く。拭き取る}{}{}%
	}%
	}%
	{\tailPrint[]}%
\entry
	{\headPrint[kana={しっすん\tl{˥}},ipa={ssuɴ},pos={動},rui={2},para={しっさぬ},acc={H}]}%
	{%
		\MeaningsBlock{%
	\ryumeaning%
		{切る}{}{}%
	}%
	}%
	{\tailPrint[]}%
\entry
	{\headPrint[kana={しっすん\tl{˥˩}},ipa={ssuɴ},pos={動},rui={2b},para={しっさぬ},acc={F}]}%
	{%
		\MeaningsBlock{%
	\ryumeaning%
		{着る}{}{}%
	}%
	}%
	{\tailPrint[]}%
\entry
	{\headPrint[kana={しっすん\tl{˥˩}},ipa={ssuɴ},pos={動},acc={F}]}%
	{%
		\MeaningsBlock{%
	\ryumeaning%
		{注ぐ。差す}{湯，茶，水だけに使う}{}%
	}%
	}%
	{\tailPrint[]}%
\entry
	{\headPrint[kana={しっすん\tl{˥˩}},ipa={ssuɴ},pos={動},acc={F}]}%
	{%
		\MeaningsBlock{%
	\ryumeaning%
		{知る}{}{}%
	}%
	}%
	{\tailPrint[]}%
\entry
	{\headPrint[kana={しっせー\tl{˥}},ipa={sseː},pos={名},acc={H}]}%
	{%
		\MeaningsBlock{%
	\ryumeaning%
		{白髪}{}{}%
	}%
	}%
	{\tailPrint[]}%
\entry
	{\headPrint[kana={じっち},ipa={dʒittʃi},pos={名}]}%
	{%
		\MeaningsBlock{%
	\ryumeaning%
		{乳房}{}{}%
	}%
	}%
	{\tailPrint[]}%
\entry
	{\headPrint[kana={しっつぁーすん\tl{˥}},ipa={ʃi̥ttsaːsuɴ},pos={動},rui={2},para={しっつぁはぬ},acc={H}]}%
	{%
		\MeaningsBlock{%
	\ryumeaning%
		{撫でる}{}{}%
	}%
	}%
	{\tailPrint[]}%
\entry
	{\headPrint[kana={しな\tl{˥}},ipa={ʃi̥na},pos={名},acc={H}]}%
	{%
		\MeaningsBlock{%
	\ryumeaning%
		{太陽。日}{}{}%
	}%
	}%
	{\tailPrint[]}%
\entry
	{\headPrint[kana={しな\tl{˥}},ipa={ʃi̥na},pos={名},acc={H}]}%
	{%
		\MeaningsBlock{%
	\ryumeaning%
		{品。種類}{}{}%
	}%
	}%
	{\tailPrint[]}%
\entry
	{\headPrint[kana={しな\tl{˥}あがるん\tl{˥˩}},ipa={ʃi̥na.agaruɴ},pos={句},para={しぃなあがらぬ},acc={H F}]}%
	{%
		\MeaningsBlock{%
	\ryumeaning%
		{日が昇る。日の出になる}{}{}%
	}%
	}%
	{\tailPrint[]}%
\entry
	{\headPrint[kana={しな\tl{˥}いるん\tl{˥˩}},ipa={ʃi̥na.iruɴ},pos={動},rui={1},para={しないらぬ},acc={H+F}]}%
	{%
		\MeaningsBlock{%
	\ryumeaning%
		{日が沈む。日没になる}{}{}%
	}%
	}%
	{\tailPrint[]}%
\entry
	{\headPrint[kana={しな\tl{˥}しっすん\tl{˥}},ipa={ʃi̥nassuɴ},pos={動},rui={2},para={しなしっさぬ},acc={H+H}]}%
	{%
		\MeaningsBlock{%
	\ryumeaning%
		{日が照る}{}{}%
	}%
	}%
	{\tailPrint[]}%
\entry
	{\headPrint[kana={しなた\tl{˥˩}},ipa={ʃi̥nata},pos={名},acc={F}]}%
	{%
		\MeaningsBlock{%
	\ryumeaning%
		{後ろ。後方}{}{}%
	}%
	}%
	{\tailPrint[]}%
\entry
	{\headPrint[kana={しなぬ\tl{˥}\ws{}みー\tl{˥˩}},ipa={ʃi̥nanu miː},pos={句},acc={H F}]}%
	{%
		\MeaningsBlock{%
	\ryumeaning%
		{日向}{}{}%
	}%
	}%
	{\tailPrint[]}%
\entry
	{\headPrint[kana={しなばく\tl{˥˩}},ipa={ʃinabaku},pos={名},acc={F}]}%
	{%
		\MeaningsBlock{%
	\ryumeaning%
		{首里大屋子}{王府時代の役職}{}%
	}%
	}%
	{\tailPrint[]}%
\entry
	{\headPrint[kana={しにぶたん},ipa={ʃi̥nibutaɴ},pos={動}]}%
	{%
		\MeaningsBlock{%
	\ryumeaning%
		{死に果てる}{}{}%
	}%
	}%
	{\tailPrint[]}%
\entry
	{\headPrint[kana={しぬ\tl{˥˩}},ipa={ʃi̥nu},pos={名},acc={F}]}%
	{%
		\MeaningsBlock{%
	\ryumeaning%
		{昨日}{}{}%
	}%
	}%
	{\tailPrint[]}%
\entry
	{\headPrint[kana={\josho{}しぬ\kako{}\ws{}ふき},ipa={ʃi̥nu fu̥ki},pos={句},acc={H L}]}%
	{%
		\MeaningsBlock{%
	\ryumeaning%
		{手首}{}{}%
	}%
	}%
	{\tailPrint[]}%
\entry
	{\headPrint[kana={しのー\tl{˥}},ipa={ʃi̥noː},pos={名},acc={H}]}%
	{%
		\MeaningsBlock{%
	\ryumeaning%
		{\ruby[g]{角}{つの}}{}{}%
	}%
	}%
	{\tailPrint[]}%
\entry
	{\headPrint[kana={しのー\tl{˥}},ipa={ʃi̥noː},pos={名},acc={H}]}%
	{%
		\MeaningsBlock{%
	\ryumeaning%
		{\ruby[g]{篩}{ふるい}}{}{}%
	}%
	}%
	{\tailPrint[]}%
\entry
	{\headPrint[kana={じばぐ\tl{˩˥}},ipa={dʒibagu},pos={名},acc={R}]}%
	{%
		\MeaningsBlock{%
	\ryumeaning%
		{重箱}{}{}%
	}%
	}%
	{\tailPrint[note={重箱の転訛}]}%
\entry
	{\headPrint[kana={しぱみるん\tl{˥˩}},ipa={ʃi̥pamiruɴ},pos={動},acc={F}]}%
	{%
		\MeaningsBlock{%
	\ryumeaning%
		{狭める。狭くする}{}{}%
	}%
	}%
	{\tailPrint[]}%
\entry
	{\headPrint[kana={じばり\tl{˥˩}},ipa={dʒibari},pos={名},acc={F}]}%
	{%
		\MeaningsBlock{%
	\ryumeaning%
		{釣針}{}{}%
	}%
	}%
	{\tailPrint[]}%
\entry
	{\headPrint[kana={しび\tl{˥˩}},ipa={ʃibi},pos={名},acc={F}]}%
	{%
		\MeaningsBlock{%
	\ryumeaning%
		{鮪}{}{}%
	}%
	}%
	{\tailPrint[]}%
\entry
	{\headPrint[kana={しぴ\tl{˥˩}},ipa={ʃi̥pi},pos={名},acc={F}]}%
	{%
		\MeaningsBlock{%
	\ryumeaning%
		{タカラガイ。子安貝}{}{}%
	}%
	}%
	{\tailPrint[]}%
\entry
	{\headPrint[kana={しぴ\tl{˥˩}},ipa={ʃi̥pi},pos={名},acc={F}]}%
	{%
		\MeaningsBlock{%
	\ryumeaning%
		{尻}{}{}%
	\ryumeaning%
		{後ろ。後}{}{}%
	}%
	}%
	{\tailPrint[]}%
\entry
	{\headPrint[kana={しぴかろ\tl{˥˩}はん},ipa={ʃi̥pik̥arohaɴ},pos={形},acc={F}]}%
	{%
		\MeaningsBlock{%
	\ryumeaning%
		{尻が軽い}{さっさと仕事を片づける}{}%
	}%
	}%
	{\tailPrint[]}%
\entry
	{\headPrint[kana={しぴぬ\tl{˥˩}みん\tl{˩˥}},ipa={ʃi̥pinumiɴ},pos={名},acc={F R}]}%
	{%
		\MeaningsBlock{%
	\ryumeaning%
		{尻の穴。肛門}{}{}%
	}%
	}%
	{\tailPrint[]}%
\entry
	{\headPrint[kana={しぴゃた\tl{˥˩}},ipa={ʃi̥pjata},pos={名},acc={F}]}%
	{%
		\MeaningsBlock{%
	\ryumeaning%
		{後ろ。後。後方}{}{}%
	}%
	}%
	{\tailPrint[]}%
\entry
	{\headPrint[kana={しぴりん\tl{˥˩}},ipa={ʃi̥piriɴ},pos={動},rui={3},para={しぴるぬ},acc={F}]}%
	{%
		\MeaningsBlock{%
	\ryumeaning%
		{萎びる。腫れがひく}{}{}%
	}%
	}%
	{\tailPrint[]}%
\entry
	{\headPrint[kana={しびる\tl{˥˩}},ipa={ʃibiru},pos={名},acc={F}]}%
	{%
		\MeaningsBlock{%
	\ryumeaning%
		{\ruby[g]{葱}{ねぎ}}{}{}%
	}%
	}%
	{\tailPrint[]}%
\entry
	{\headPrint[kana={しぴんさ\tl{˥˩}はーん},ipa={ʃi̥pinsahaːɴ},pos={形},acc={F}]}%
	{%
		\MeaningsBlock{%
	\ryumeaning%
		{尻が重い。不精だ}{なかなか動こうとしない}{}%
	}%
	}%
	{\tailPrint[]}%
\entry
	{\headPrint[kana={しふく\tl{˥}},ipa={ʃifu̥ku},pos={名},acc={H}]}%
	{%
		\MeaningsBlock{%
	\ryumeaning%
		{\ruby[g]{拳}{こぶし}}{}{}%
	}%
	}%
	{\tailPrint[]}%
\entry
	{\headPrint[kana={じふねー\tl{˩˥}},ipa={dʒifu̥neː},pos={名},acc={R}]}%
	{%
		\MeaningsBlock{%
	\ryumeaning%
		{陸酔い}{下船後も酔うこと}{}%
	}%
	}%
	{\tailPrint[]}%
\entry
	{\headPrint[kana={じぶん\tl{˥˩}},ipa={dʒibuɴ},pos={名},acc={F}]}%
	{%
		\MeaningsBlock{%
	\ryumeaning%
		{時分。ころ}{}{}%
	}%
	}%
	{\tailPrint[]}%
\entry
	{\headPrint[kana={じぼー\tl{˥˩}},ipa={dʒiboː},pos={名},acc={F}]}%
	{%
		\MeaningsBlock{%
	\ryumeaning%
		{釣り竿}{}{}%
	}%
	}%
	{\tailPrint[]}%
\entry
	{\headPrint[kana={しぼへぬ},ipa={ʃi̥bohenu},pos={句},para={しぼはん}]}%
	{%
		\MeaningsBlock{%
	\ryumeaning%
		{したくない。やりたくない}{}{}%
	}%
	}%
	{\tailPrint[]}%
\entry
	{\headPrint[kana={しまいるん\tl{˥˩}},ipa={ʃimairuɴ},pos={動},acc={F}]}%
	{%
		\MeaningsBlock{%
	\ryumeaning%
		{済む。終わる}{}{}%
	}%
	}%
	{\tailPrint[]}%
\entry
	{\headPrint[kana={しまはかるん\tl{˥}},ipa={si̥mahḁkaruɴ},pos={動},rui={1},para={しまはからぬ},acc={H}]}%
	{%
		\MeaningsBlock{%
	\ryumeaning%
		{手間取る}{}{}%
	}%
	}%
	{\tailPrint[]}%
\entry
	{\headPrint[kana={しまふさら\tl{˥}},ipa={ʃi̥mafu̥sara},pos={名},acc={H}]}%
	{%
		\MeaningsBlock{%
	\ryumeaning%
		{疫病除けの行事}{}{}%
	}%
	}%
	{\tailPrint[]}%
\entry
	{\headPrint[kana={しみ\tl{˥˩}},ipa={ʃi̥mi},pos={名},acc={F}]}%
	{%
		\MeaningsBlock{%
	\ryumeaning%
		{\ruby[g]{爪}{つめ}}{}{}%
	}%
	}%
	{\tailPrint[]}%
\entry
	{\headPrint[kana={しみしきるん\tl{˥}},ipa={ʃi̥miʃikiruɴ},pos={動},rui={3},para={しみすくぬ},acc={H}]}%
	{%
		\MeaningsBlock{%
	\ryumeaning%
		{責めつける}{}{}%
	}%
	}%
	{\tailPrint[]}%
\entry
	{\headPrint[kana={しみしきるん\tl{˥}},ipa={ʃi̥miʃi̥kiruɴ},pos={動},rui={3},para={しみすくぬ},acc={H}]}%
	{%
		\MeaningsBlock{%
	\ryumeaning%
		{締めつける}{}{}%
	}%
	}%
	{\tailPrint[]}%
\entry
	{\headPrint[kana={しみらすん\tl{˥}},ipa={ʃi̥mirasuɴ},pos={動},rui={2},para={しみらはぬ},acc={H}]}%
	{%
		\MeaningsBlock{%
	\ryumeaning%
		{湿らせる}{水気を与える}{}%
	}%
	}%
	{\tailPrint[]}%
\entry
	{\headPrint[kana={しみるん\tl{˥}},ipa={ʃi̥miruɴ},pos={動},rui={3},para={すむぬ},acc={H}]}%
	{%
		\MeaningsBlock{%
	\ryumeaning%
		{締める。締めつける}{}{}%
	}%
	}%
	{\tailPrint[]}%
\entry
	{\headPrint[kana={しみるん\tl{˥}},ipa={ʃi̥miruɴ},pos={動},rui={1},para={しみらぬ},acc={H}]}%
	{%
		\MeaningsBlock{%
	\ryumeaning%
		{湿る。湿っぽくなる}{}{}%
	}%
	}%
	{\tailPrint[]}%
\entry
	{\headPrint[kana={しみるん\tl{˥}},ipa={ʃi̥miruɴ},pos={動},acc={H}]}%
	{%
		\MeaningsBlock{%
	\ryumeaning%
		{責める。非難する}{}{}%
	\ryumeaning%
		{攻める}{}{}%
	}%
	}%
	{\tailPrint[]}%
\entry
	{\headPrint[kana={しみん\tl{˥˩}},ipa={ʃi̥miɴ},pos={動},rui={3},para={すむぬ},acc={F}]}%
	{%
		\MeaningsBlock{%
	\ryumeaning%
		{染める}{}{}%
	}%
	}%
	{\tailPrint[]}%
\entry
	{\headPrint[kana={しむぬ\tl{˥˩}},ipa={ʃi̥munu},pos={名},acc={F}]}%
	{%
		\MeaningsBlock{%
	\ryumeaning%
		{吸い物}{料理名}{}%
	}%
	}%
	{\tailPrint[]}%
\entry
	{\headPrint[kana={しゃー\tl{˥}},ipa={ʃaː},pos={名},acc={H}]}%
	{%
		\MeaningsBlock{%
	\ryumeaning%
		{\ruby[g]{枡}{ます}。\ruby[g]{一升}{いっしょう}\ruby[g]{枡}{ます}}{}{}%
	}%
	}%
	{\tailPrint[]}%
\entry
	{\headPrint[kana={しゃー\tl{˥}},ipa={ʃaː},pos={副},acc={H}]}%
	{%
		\MeaningsBlock{%
	\ryumeaning%
		{何時も。常に}{}{}%
	}%
	}%
	{\tailPrint[]}%
\entry
	{\headPrint[kana={しゃーみ\tl{˥}},ipa={ʃaːmi},pos={名},acc={H}]}%
	{%
		\MeaningsBlock{%
	\ryumeaning%
		{つもり。予定}{}{}%
	}%
	}%
	{\tailPrint[]}%
\entry
	{\headPrint[kana={しゃーみるん\tl{˥˩}},ipa={ʃaːmiruɴ},pos={動},rui={3},para={しゃーむぬ},acc={F}]}%
	{%
		\MeaningsBlock{%
	\ryumeaning%
		{やっつける。やり込める。𠮟りつける}{}{}%
	}%
	}%
	{\tailPrint[]}%
\entry
	{\headPrint[kana={しゃー\tl{˥}\ws{}んぐん\tl{˥˩}},ipa={ʃaː ŋguɴ},pos={句},acc={H F}]}%
	{%
		\MeaningsBlock{%
	\ryumeaning%
		{よく行く。しばしば行く}{}{}%
	}%
	}%
	{\tailPrint[]}%
\entry
	{\headPrint[kana={じゃく\tl{˩˥}},ipa={dʒaku},pos={名},acc={R}]}%
	{%
		\MeaningsBlock{%
	\ryumeaning%
		{\ruby[g]{鰹}{かつお}の餌の小魚}{「雑魚」から}{}%
	}%
	}%
	{\tailPrint[]}%
\entry
	{\headPrint[kana={しゃしゃびら\tl{˥}},ipa={ʃaʃabira},pos={名},acc={H}]}%
	{%
		\MeaningsBlock{%
	\ryumeaning%
		{\ruby[g]{杓文字}{しゃもじ}}{}{}%
	}%
	}%
	{\tailPrint[]}%
\entry
	{\headPrint[kana={しゃま\tl{˥}},ipa={ʃama},pos={名},acc={H}]}%
	{%
		\MeaningsBlock{%
	\ryumeaning%
		{兄。年上}{「兄方」の義}{}%
	}%
	}%
	{\tailPrint[]}%
\entry
	{\headPrint[kana={しゃまかた\tl{˥}},ipa={ʃamakḁta},pos={名},acc={H}]}%
	{%
		\MeaningsBlock{%
	\ryumeaning%
		{先輩}{}{}%
	}%
	}%
	{\tailPrint[]}%
\entry
	{\headPrint[kana={しゃみしきるん},ipa={ʃamiʃi̥kiruɴ},pos={動},rui={3},para={しゃみすくぬ}]}%
	{%
		\MeaningsBlock{%
	\ryumeaning%
		{叱りつける。どやしつける}{}{}%
	}%
	}%
	{\tailPrint[]}%
\entry
	{\headPrint[kana={しゃみしきん\tl{˥˩}},ipa={ʃamiʃi̥kiɴ},pos={動},rui={3},para={しゃみすくぬ},acc={F}]}%
	{%
		\MeaningsBlock{%
	\ryumeaning%
		{叱りつける}{}{}%
	}%
	}%
	{\tailPrint[]}%
\entry
	{\headPrint[kana={しゃんしゃん\tl{˥}},ipa={ʃanʃaɴ},pos={名},acc={H}]}%
	{%
		\MeaningsBlock{%
	\ryumeaning%
		{セミ。アブラゼミ}{}{}%
	}%
	}%
	{\tailPrint[]}%
\entry
	{\headPrint[kana={しゅくぶん\tl{˩˥}},ipa={ʃu̥kubuɴ},pos={名},acc={R}]}%
	{%
		\MeaningsBlock{%
	\ryumeaning%
		{職分。任務}{}{}%
	}%
	}%
	{\tailPrint[]}%
\entry
	{\headPrint[kana={しゅぷ\tl{˥˩}},ipa={ʃu̥pu},pos={名},acc={F}]}%
	{%
		\MeaningsBlock{%
	\ryumeaning%
		{鉄砲}{}{}%
	}%
	}%
	{\tailPrint[]}%
\entry
	{\headPrint[kana={しゅる\tl{˥˩}},ipa={ʃuru},pos={名},acc={F}]}%
	{%
		\MeaningsBlock{%
	\ryumeaning%
		{棕櫚}{}{}%
	}%
	}%
	{\tailPrint[]}%
\entry
	{\headPrint[kana={しゅるん\tl{˥}},ipa={ʃuruɴ},pos={動},rui={1},para={しゅらぬ},acc={H}]}%
	{%
		\MeaningsBlock{%
	\ryumeaning%
		{（稲を）扱く。そぐ。脱穀する。しごく。しごき落とす}{}{}%
	}%
	}%
	{\tailPrint[]}%
\entry
	{\headPrint[kana={じょーぎ\tl{˩˥}},ipa={dʒoːgi},pos={名},acc={R}]}%
	{%
		\MeaningsBlock{%
	\ryumeaning%
		{定規。物差し}{}{}%
	}%
	}%
	{\tailPrint[]}%
\entry
	{\headPrint[kana={じょーぐ\tl{˩˥}},ipa={dʒoːgu},pos={名},acc={R}]}%
	{%
		\MeaningsBlock{%
	\ryumeaning%
		{\ruby[g]{上戸}{じょうご}}{酒好き}{}%
	}%
	}%
	{\tailPrint[]}%
\entry
	{\headPrint[kana={じょー\tl{˥˩}ふか\tl{˥˩}はーん},ipa={dʒoːfu̥kahaːɴ},pos={形},acc={F F}]}%
	{%
		\MeaningsBlock{%
	\ryumeaning%
		{情が深い}{}{}%
	}%
	}%
	{\tailPrint[]}%
\entry
	{\headPrint[kana={じょーぶくろ\tl{˥˩}},ipa={dʒoːbukuro},pos={名},acc={F}]}%
	{%
		\MeaningsBlock{%
	\ryumeaning%
		{封筒}{}{}%
	}%
	}%
	{\tailPrint[]}%
\entry
	{\headPrint[kana={しょーむぬ\tl{˥˩}},ipa={ʃoːmunu},pos={名},acc={F}]}%
	{%
		\MeaningsBlock{%
	\ryumeaning%
		{良い物。立派な物}{}{}%
	}%
	}%
	{\tailPrint[]}%
\entry
	{\headPrint[kana={しょっこー\tl{˥}},ipa={ʃokkoː},pos={名},acc={H}]}%
	{%
		\MeaningsBlock{%
	\ryumeaning%
		{法事。法要}{直訳「焼香」。普通は\emb{こっこー}と言う}{}%
	}%
	}%
	{\tailPrint[]}%
\entry
	{\headPrint[kana={しら\tl{˥˩}},ipa={ʃi̥ra},pos={名},acc={F}]}%
	{%
		\MeaningsBlock{%
	\ryumeaning%
		{\ruby[g]{稲}{いな}むら}{}{}%
	}%
	}%
	{\tailPrint[]}%
\entry
	{\headPrint[kana={しらいん\tl{˥˩}},ipa={ʃi̥raiɴ},pos={動},rui={3},para={しらるぬ},acc={F}]}%
	{%
		\MeaningsBlock{%
	\ryumeaning%
		{してやられる}{}{}%
	}%
	}%
	{\tailPrint[]}%
\entry
	{\headPrint[kana={しらす\tl{˥˩}},ipa={ʃi̥rasu},pos={名},acc={F}]}%
	{%
		\MeaningsBlock{%
	\ryumeaning%
		{予行演習。リハーサル}{}{}%
	}%
	}%
	{\tailPrint[]}%
\entry
	{\headPrint[kana={じらば\tl{˩˥}},ipa={dʒiraba},pos={名},acc={R}]}%
	{%
		\MeaningsBlock{%
	\ryumeaning%
		{古謡の労働歌}{}{}%
	}%
	}%
	{\tailPrint[]}%
\entry
	{\headPrint[kana={じり\tl{˩˥}},ipa={dʒiri},pos={名},acc={R}]}%
	{%
		\MeaningsBlock{%
	\ryumeaning%
		{暗礁。リーフ}{}{}%
	}%
	}%
	{\tailPrint[]}%
\entry
	{\headPrint[kana={じりく\tl{˩˥}にち\tl{˥˩}},ipa={dʒiriku̥nitʃi},pos={名},acc={R+F}]}%
	{%
		\MeaningsBlock{%
	\ryumeaning%
		{十六日祭}{旧暦一月十六日の墓前での先祖や死者供養}{}%
	}%
	}%
	{\tailPrint[]}%
\entry
	{\headPrint[kana={しるし\tl{˥˩}},ipa={ʃiruʃi},pos={名},acc={F}]}%
	{%
		\MeaningsBlock{%
	\ryumeaning%
		{\ruby[g]{印}{しるし}。兆候}{}{}%
	}%
	}%
	{\tailPrint[]}%
\entry
	{\headPrint[kana={しるん\tl{˥}},ipa={ʃi̥ruɴ},pos={動},rui={1},para={しらぬ},acc={H}]}%
	{%
		\MeaningsBlock{%
	\ryumeaning%
		{（太陽が）照る}{}{}%
	}%
	}%
	{\tailPrint[]}%
\entry
	{\headPrint[kana={じろーし\tl{˥˩}},ipa={dʒiroːʃi},pos={名},acc={F}]}%
	{%
		\MeaningsBlock{%
	\ryumeaning%
		{サルカケミカン}{鋭い棘の植物}{}%
	}%
	}%
	{\tailPrint[]}%
\entry
	{\headPrint[kana={しわー\tl{˥˩}},ipa={ʃi̥waː},pos={名},acc={F}]}%
	{%
		\MeaningsBlock{%
	\ryumeaning%
		{心配。不安}{}{}%
	}%
	}%
	{\tailPrint[]}%
\entry
	{\headPrint[kana={しわざ\tl{˥˩}},ipa={ʃi̥wadza},pos={名},acc={F}]}%
	{%
		\MeaningsBlock{%
	\ryumeaning%
		{仕業。行為}{}{}%
	}%
	}%
	{\tailPrint[]}%
\entry
	{\headPrint[kana={しん\tl{˥}},ipa={ʃiɴ},pos={名},acc={H}]}%
	{%
		\MeaningsBlock{%
	\ryumeaning%
		{墨。墨汁}{}{}%
	}%
	}%
	{\tailPrint[]}%
\entry
	{\headPrint[kana={しん\tl{˥}},ipa={siɴ},pos={名},acc={H}]}%
	{%
		\MeaningsBlock{%
	\ryumeaning%
		{\ruby[g]{唾}{つば}。唾液}{}{}%
	}%
	}%
	{\tailPrint[]}%
\entry
	{\headPrint[kana={しん\tl{˥}},ipa={ʃiɴ},pos={名},acc={H}]}%
	{%
		\MeaningsBlock{%
	\ryumeaning%
		{\ruby[g]{稗}{ひえ}}{}{}%
	}%
	}%
	{\tailPrint[]}%
\entry
	{\headPrint[kana={しん\tl{˥}},ipa={ʃiɴ},pos={名},acc={H}]}%
	{%
		\MeaningsBlock{%
	\ryumeaning%
		{招待客。お客}{}{}%
	}%
	}%
	{\tailPrint[]}%
\entry
	{\headPrint[kana={しん\tl{˥˩}},ipa={ʃin},pos={名},acc={F}]}%
	{%
		\MeaningsBlock{%
	\ryumeaning%
		{千}{}{}%
	}%
	}%
	{\tailPrint[]}%
\entry
	{\headPrint[kana={じん\tl{˥˩}},ipa={dʒiɴ},pos={名},acc={F}]}%
	{%
		\MeaningsBlock{%
	\ryumeaning%
		{天。空}{}{}%
	}%
	}%
	{\tailPrint[]}%
\entry
	{\headPrint[kana={じん\tl{˥˩}},ipa={dʒiɴ},pos={名},acc={F}]}%
	{%
		\MeaningsBlock{%
	\ryumeaning%
		{お膳。膳}{}{}%
	}%
	}%
	{\tailPrint[]}%
\entry
	{\headPrint[kana={じん\tl{˩˥}},ipa={dʒiɴ},pos={名},acc={R}]}%
	{%
		\MeaningsBlock{%
	\ryumeaning%
		{銭。お金}{}{}%
	}%
	}%
	{\tailPrint[]}%
\entry
	{\headPrint[kana={しんか\tl{˥}},ipa={ʃinka},pos={名},acc={H}]}%
	{%
		\MeaningsBlock{%
	\ryumeaning%
		{従業員。部下}{「臣下」すなわち「家来」から}{}%
	}%
	}%
	{\tailPrint[note={「臣下」の転}]}%
\entry
	{\headPrint[kana={しんがま\tl{˥}},ipa={ʃiŋgama},pos={名},acc={H}]}%
	{%
		\MeaningsBlock{%
	\ryumeaning%
		{\ruby[g]{蠍}{さそり}}{}{}%
	}%
	}%
	{\tailPrint[]}%
\entry
	{\headPrint[kana={しんがら\tl{˥}},ipa={ʃiŋgara},pos={名},acc={H}]}%
	{%
		\MeaningsBlock{%
	\ryumeaning%
		{大きな金梃子}{}{}%
	}%
	}%
	{\tailPrint[]}%
\entry
	{\headPrint[kana={じんぎり\tl{˩˥}},ipa={dʒiŋgiri},pos={名},acc={R}]}%
	{%
		\MeaningsBlock{%
	\ryumeaning%
		{茶筒}{}{}%
	}%
	}%
	{\tailPrint[]}%
\entry
	{\headPrint[kana={しんくつぇー\tl{˥}},ipa={ʃiŋkutseː},pos={名},acc={H}]}%
	{%
		\MeaningsBlock{%
	\ryumeaning%
		{洗骨}{法事の一つ}{}%
	}%
	}%
	{\tailPrint[]}%
\entry
	{\headPrint[kana={しんけー\tl{˥}},ipa={ʃinkeː},pos={名},acc={H}]}%
	{%
		\MeaningsBlock{%
	\ryumeaning%
		{気違い}{「神経」から}{}%
	}%
	}%
	{\tailPrint[]}%
\entry
	{\headPrint[kana={\josho{}しんじ\kako{}ふちり},ipa={ʃindʒifu̥tʃiri},pos={名},acc={H+L}]}%
	{%
		\MeaningsBlock{%
	\ryumeaning%
		{煎じ薬}{}{}%
	}%
	}%
	{\tailPrint[]}%
\entry
	{\headPrint[kana={しんじょー\tl{˥}},ipa={ʃindʒoː},pos={名},acc={H}]}%
	{%
		\MeaningsBlock{%
	\ryumeaning%
		{天井}{}{}%
	}%
	}%
	{\tailPrint[]}%
\entry
	{\headPrint[kana={しんじるん\tl{˥}},ipa={ʃindʒiruɴ},pos={動},rui={3},para={しんじぬ},acc={H}]}%
	{%
		\MeaningsBlock{%
	\ryumeaning%
		{煎じる}{}{}%
	}%
	}%
	{\tailPrint[]}%
\entry
	{\headPrint[kana={しんじるん\tl{˥}},ipa={ʃindʒiruɴ},pos={動},rui={3},para={しんずぬ},acc={H}]}%
	{%
		\MeaningsBlock{%
	\ryumeaning%
		{信じる。信用する。信仰する}{}{}%
	}%
	}%
	{\tailPrint[]}%
\entry
	{\headPrint[kana={しんしん\tl{˥}},ipa={ʃinʃiɴ},pos={名},acc={H}]}%
	{%
		\MeaningsBlock{%
	\ryumeaning%
		{先生}{}{}%
	}%
	}%
	{\tailPrint[note={先生の転訛}]}%
\entry
	{\headPrint[kana={しんず\tl{˥}},ipa={ʃindzu},pos={名},acc={H}]}%
	{%
		\MeaningsBlock{%
	\ryumeaning%
		{四十}{}{}%
	}%
	}%
	{\tailPrint[]}%
\entry
	{\headPrint[kana={じん\tl{˩˥}たみるん\tl{˥˩}},ipa={dʒintḁmiruɴ},pos={動},rui={3},para={じんたむぬ},acc={R+F}]}%
	{%
		\MeaningsBlock{%
	\ryumeaning%
		{金を貯める。貯金する}{}{}%
	}%
	}%
	{\tailPrint[]}%
\entry
	{\headPrint[kana={しんだら\tl{˥}さーん},ipa={ʃindarasaːɴ},pos={形},acc={H}]}%
	{%
		\MeaningsBlock{%
	\ryumeaning%
		{可愛い。可愛らしい}{}{}%
	}%
	}%
	{\tailPrint[]}%
\entry
	{\headPrint[kana={しんどー\tl{˥}},ipa={ʃindoː},pos={名},acc={H}]}%
	{%
		\MeaningsBlock{%
	\ryumeaning%
		{船頭}{}{}%
	}%
	}%
	{\tailPrint[]}%
\entry
	{\headPrint[kana={じんぬ\tl{˥˩}\ws{}ふかー\tl{˥}},ipa={dʒinnu fu̥kaː},pos={句},acc={F H}]}%
	{%
		\MeaningsBlock{%
	\ryumeaning%
		{天の川。銀河}{}{}%
	}%
	}%
	{\tailPrint[]}%
\entry
	{\headPrint[kana={じんぱるん},ipa={dʒimpḁruɴ},pos={動}]}%
	{%
		\MeaningsBlock{%
	\ryumeaning%
		{\ruby[g]{抓}{つね}る}{}{}%
	}%
	}%
	{\tailPrint[]}%
\entry
	{\headPrint[kana={じんぱるん\tl{˥˩}},ipa={dʒimpḁruɴ},pos={動},rui={1},para={じんぱらぬ},acc={F}]}%
	{%
		\MeaningsBlock{%
	\ryumeaning%
		{噛み切る}{}{}%
	}%
	}%
	{\tailPrint[]}%
\entry
	{\headPrint[kana={しんび\tl{˥}},ipa={ʃimbi},pos={名},acc={H}]}%
	{%
		\MeaningsBlock{%
	\ryumeaning%
		{手の指}{}{}%
	}%
	}%
	{\tailPrint[]}%
\entry
	{\headPrint[kana={じんふくる\tl{˩˥}},ipa={dʒinfu̥kuru},pos={名},acc={R}]}%
	{%
		\MeaningsBlock{%
	\ryumeaning%
		{財布。お金入れ}{直訳「銭袋」}{}%
	}%
	}%
	{\tailPrint[]}%
\entry
	{\headPrint[kana={しんぷら\tl{˥}},ipa={ʃimpura},pos={名},acc={H}]}%
	{%
		\MeaningsBlock{%
	\ryumeaning%
		{天ぷら}{}{}%
	}%
	}%
	{\tailPrint[note={「てんぷら」の転訛}]}%
\entry
	{\headPrint[kana={じんぶん\tl{˩˥}},ipa={dʒimbuɴ},pos={名},acc={R}]}%
	{%
		\MeaningsBlock{%
	\ryumeaning%
		{知恵。分別}{}{}%
	}%
	}%
	{\tailPrint[]}%
\entry
	{\headPrint[kana={じんぶん\tl{˩˥}\ws{}とーらん\tl{˥˩}},ipa={dʒimbuɴ toːraɴ},pos={句},acc={R H}]}%
	{%
		\MeaningsBlock{%
	\ryumeaning%
		{気が利かなくなる}{}{}%
	}%
	}%
	{\tailPrint[]}%
\entry
	{\headPrint[kana={じんむちゃー\tl{˩˥}},ipa={dʒimmutʃaː},pos={名},acc={R}]}%
	{%
		\MeaningsBlock{%
	\ryumeaning%
		{金持ち}{}{}%
	}%
	}%
	{\tailPrint[]}%
\entry
	{\headPrint[kana={じんもーき\tl{˩˥}},ipa={dʒimmoːki},pos={名},acc={R}]}%
	{%
		\MeaningsBlock{%
	\ryumeaning%
		{金儲け}{}{}%
	}%
	}%
	{\tailPrint[]}%
\LetterHead{す}
\entry
	{\headPrint[kana={すー\tl{˥}},ipa={suː},pos={名},acc={H}]}%
	{%
		\MeaningsBlock{%
	\ryumeaning%
		{潮。海水}{}{}%
	}%
	}%
	{\tailPrint[]}%
\entry
	{\headPrint[kana={すー\tl{˥}},ipa={suː},pos={名},acc={H}]}%
	{%
		\MeaningsBlock{%
	\ryumeaning%
		{おつゆ。汁}{}{}%
	}%
	}%
	{\tailPrint[]}%
\entry
	{\headPrint[kana={ずーし\tl{˩˥}},ipa={dzuːʃi},pos={名},acc={R}]}%
	{%
		\MeaningsBlock{%
	\ryumeaning%
		{おじや。雑炊}{}{}%
	}%
	}%
	{\tailPrint[]}%
\entry
	{\headPrint[kana={すーだが\tl{˥}},ipa={suːdaga},pos={名},acc={H}]}%
	{%
		\MeaningsBlock{%
	\ryumeaning%
		{合計。総額}{}{}%
	}%
	}%
	{\tailPrint[]}%
\entry
	{\headPrint[kana={すーぬ\tl{˥}\ws{}ふち\tl{˥˩}},ipa={suːnu fu̥tʃi},pos={名},acc={H F}]}%
	{%
		\MeaningsBlock{%
	\ryumeaning%
		{なぎさ。波打ち際}{}{}%
	}%
	}%
	{\tailPrint[]}%
\entry
	{\headPrint[kana={すーぴき\tl{˥˩}},ipa={suːpi̥ki},pos={名},acc={F}]}%
	{%
		\MeaningsBlock{%
	\ryumeaning%
		{潮流}{}{}%
	}%
	}%
	{\tailPrint[]}%
\entry
	{\headPrint[kana={すーぴすん\tl{˥}},ipa={suːpi̥ʃi},pos={動},acc={H}]}%
	{%
		\MeaningsBlock{%
	\ryumeaning%
		{潮が引く。干潮になる}{}{}%
	}%
	}%
	{\tailPrint[]}%
\entry
	{\headPrint[kana={ずーぶるん\tl{˥}},ipa={dzuːburuɴ},pos={動},rui={1},para={ずーぶらぬ},acc={H}]}%
	{%
		\MeaningsBlock{%
	\ryumeaning%
		{発情する}{}{}%
	}%
	}%
	{\tailPrint[]}%
\entry
	{\headPrint[kana={ずーぶん\tl{˥}},ipa={dzuːbuɴ},pos={動},rui={1},para={ずーばぬ},acc={H}]}%
	{%
		\MeaningsBlock{%
	\ryumeaning%
		{交尾する}{}{}%
	}%
	}%
	{\tailPrint[]}%
\entry
	{\headPrint[kana={すーぺー\tl{˥}},ipa={suːpeː},pos={名},acc={H}]}%
	{%
		\MeaningsBlock{%
	\ryumeaning%
		{杓子。お玉}{}{}%
	}%
	}%
	{\tailPrint[]}%
\entry
	{\headPrint[kana={すーまん\tl{˥}},ipa={suːmam},pos={名},acc={H}]}%
	{%
		\MeaningsBlock{%
	\ryumeaning%
		{小満}{季節の名。沖縄では梅雨時期にあたる}{}%
	}%
	}%
	{\tailPrint[]}%
\entry
	{\headPrint[kana={すーむに\tl{˥}},ipa={suːmuni},pos={名},acc={H}]}%
	{%
		\MeaningsBlock{%
	\ryumeaning%
		{苦言。忠言}{}{}%
	}%
	}%
	{\tailPrint[]}%
\entry
	{\headPrint[kana={すーる\tl{˥˩}},ipa={su̥ru},pos={名},acc={F}]}%
	{%
		\MeaningsBlock{%
	\ryumeaning%
		{\ruby[g]{弦}{げん}}{}{}%
	}%
	}%
	{\tailPrint[]}%
\entry
	{\headPrint[kana={すー\tl{˥}んつん\tl{˩˥}},ipa={suːntsuɴ},pos={動},acc={H+R}]}%
	{%
		\MeaningsBlock{%
	\ryumeaning%
		{潮が満ちる。満潮になる}{}{}%
	}%
	}%
	{\tailPrint[]}%
\entry
	{\headPrint[kana={すぅとぅ\tl{˥}},ipa={su̥tu},pos={名},acc={H}]}%
	{%
		\MeaningsBlock{%
	\ryumeaning%
		{お\ruby[g]{土産}{みやげ}}{}{}%
	}%
	}%
	{\tailPrint[]}%
\entry
	{\headPrint[kana={すぅな\tl{˥}はん},ipa={su̥nahaɴ},pos={形},para={しぃなへぬ？},acc={H}]}%
	{%
		\MeaningsBlock{%
	\ryumeaning%
		{\ruby[g]{拙}{つたな}い。幼稚}{}{}%
	}%
	}%
	{\tailPrint[]}%
\entry
	{\headPrint[kana={すぅぬ},ipa={su̥nu},pos={名}]}%
	{%
		\MeaningsBlock{%
	\ryumeaning%
		{昨日}{}{}%
	}%
	}%
	{\tailPrint[]}%
\entry
	{\headPrint[kana={すぅぬ\tl{˥}},ipa={su̥nu},pos={名},acc={H}]}%
	{%
		\MeaningsBlock{%
	\ryumeaning%
		{着物}{}{}%
	}%
	}%
	{\tailPrint[]}%
\entry
	{\headPrint[kana={すぅぬ\tl{˥}\ws{}あらすん\tl{˥˩}},ipa={su̥nu. arasuɴ},pos={句},para={すぬあらはぬ},acc={H F}]}%
	{%
		\MeaningsBlock{%
	\ryumeaning%
		{洗濯する}{}{}%
	}%
	}%
	{\tailPrint[]}%
\entry
	{\headPrint[kana={すぅぬ\tl{˥}\ws{}かれーるん\tl{˥˩}},ipa={su̥nu kareːruɴ},pos={句},para={すぬかれるぬ},acc={H F}]}%
	{%
		\MeaningsBlock{%
	\ryumeaning%
		{着替える}{}{}%
	}%
	}%
	{\tailPrint[]}%
\entry
	{\headPrint[kana={すぅぬ\tl{˥}\ws{}すっすん\tl{˥˩}},ipa={su̥nu ssuɴ},pos={句},acc={H F}]}%
	{%
		\MeaningsBlock{%
	\ryumeaning%
		{着物を着る}{}{}%
	}%
	}%
	{\tailPrint[]}%
\entry
	{\headPrint[kana={すぅぬん\tl{˥˩}},ipa={su̥nuɴ},pos={動},rui={1},para={すなぬ},acc={F}]}%
	{%
		\MeaningsBlock{%
	\ryumeaning%
		{死ぬ}{敬語は\emb{まらすん}}{}%
	}%
	}%
	{\tailPrint[]}%
\entry
	{\headPrint[kana={すぅぱ\tl{˥}},ipa={su̥pa},pos={名},acc={H}]}%
	{%
		\MeaningsBlock{%
	\ryumeaning%
		{\ruby[g]{唇}{くちびる}}{}{}%
	}%
	}%
	{\tailPrint[]}%
\entry
	{\headPrint[kana={すぅぷ\tl{˥˩}},ipa={su̥pu},pos={名},acc={F}]}%
	{%
		\MeaningsBlock{%
	\ryumeaning%
		{\ruby[g]{壺}{つぼ}}{}{}%
	\ryumeaning%
		{急所}{}{}%
	}%
	}%
	{\tailPrint[]}%
\entry
	{\headPrint[kana={すぅふく\tl{˥}},ipa={sufu̥ku},pos={名},acc={H}]}%
	{%
		\MeaningsBlock{%
	\ryumeaning%
		{クモ。クモの糸}{}{}%
	}%
	}%
	{\tailPrint[]}%
\entry
	{\headPrint[kana={すぅま\tl{˥}},ipa={su̥ma},pos={名},acc={H}]}%
	{%
		\MeaningsBlock{%
	\ryumeaning%
		{島。郷里}{我が島(べーすま)など}{}%
	}%
	}%
	{\tailPrint[]}%
\entry
	{\headPrint[kana={すぅまるん\tl{˥}},ipa={su̥maruɴ},pos={動},rui={1},para={すまらぬ},acc={H}]}%
	{%
		\MeaningsBlock{%
	\ryumeaning%
		{\ruby[g]{詰}{つ}まる}{}{}%
	}%
	}%
	{\tailPrint[]}%
\entry
	{\headPrint[kana={すか\tl{˥}},ipa={su̥ka},pos={名},acc={H}]}%
	{%
		\MeaningsBlock{%
	\ryumeaning%
		{囲炉裏。炉}{}{}%
	}%
	}%
	{\tailPrint[]}%
\entry
	{\headPrint[kana={すか\tl{˥}},ipa={su̥ka},pos={名},acc={H}]}%
	{%
		\MeaningsBlock{%
	\ryumeaning%
		{柄}{道具の柄}{}%
	}%
	}%
	{\tailPrint[]}%
\entry
	{\headPrint[kana={すかーるぬ},ipa={su̥kaːrunu},pos={句}]}%
	{%
		\MeaningsBlock{%
	\ryumeaning%
		{聞こえない}{}{}%
	}%
	}%
	{\tailPrint[]}%
\entry
	{\headPrint[kana={すがい\tl{˥}},ipa={sugai},pos={名},acc={H}]}%
	{%
		\MeaningsBlock{%
	\ryumeaning%
		{身なり。身だしなみ}{}{}%
	}%
	}%
	{\tailPrint[]}%
\entry
	{\headPrint[kana={すかいるん\tl{˥˩}},ipa={su̥kairuɴ},pos={動},rui={3},para={すかるぬ},acc={F}]}%
	{%
		\MeaningsBlock{%
	\ryumeaning%
		{聞こえる。聞かれる}{}{}%
	\ryumeaning%
		{聞ける}{聞くことが出来る}{}%
	}%
	}%
	{\tailPrint[]}%
\entry
	{\headPrint[kana={すかうちるん\tl{˥}},ipa={su̥ka.utʃiruɴ},pos={動},rui={3},para={すかうとぬ},acc={H}]}%
	{%
		\MeaningsBlock{%
	\ryumeaning%
		{落っこちる}{「落ちる」の強調}{}%
	}%
	}%
	{\tailPrint[]}%
\entry
	{\headPrint[kana={すかさ\tl{˥}},ipa={su̥kasa},pos={名},acc={F}]}%
	{%
		\MeaningsBlock{%
	\ryumeaning%
		{司。神司}{}{}%
	}%
	}%
	{\tailPrint[]}%
\entry
	{\headPrint[kana={すかすん\tl{˥˩}},ipa={su̥kasuɴ},pos={動},rui={2},para={すかはぬ},acc={F}]}%
	{%
		\MeaningsBlock{%
	\ryumeaning%
		{聞かせる}{}{}%
	\ryumeaning%
		{知らせる}{}{}%
	}%
	}%
	{\tailPrint[]}%
\entry
	{\headPrint[kana={すかすん\tl{˥˩}},ipa={su̥kasuɴ},pos={動},rui={2},para={すかはぬ},acc={F}]}%
	{%
		\MeaningsBlock{%
	\ryumeaning%
		{鋤く}{}{}%
	}%
	}%
	{\tailPrint[]}%
\entry
	{\headPrint[kana={すか\ws{}っしりん},ipa={su̥ka ʃʃiriɴ},pos={句},rui={3},para={すか　っすぬ}]}%
	{%
		\MeaningsBlock{%
	\ryumeaning%
		{投げ捨てる}{捨てるの強調}{}%
	}%
	}%
	{\tailPrint[]}%
\entry
	{\headPrint[kana={すかなーうや\tl{˥}},ipa={su̥kanaː.uja},pos={名},acc={H}]}%
	{%
		\MeaningsBlock{%
	\ryumeaning%
		{養い親。養父母}{}{}%
	}%
	}%
	{\tailPrint[]}%
\entry
	{\headPrint[kana={すかなすん\tl{˥}},ipa={su̥kanasuɴ},pos={動},rui={2},para={すかなはぬ},acc={H}]}%
	{%
		\MeaningsBlock{%
	\ryumeaning%
		{養育する。飼育する}{}{}%
	}%
	}%
	{\tailPrint[]}%
\entry
	{\headPrint[kana={すかのーすん\tl{˥}},ipa={su̥kanoːsuɴ},pos={動},rui={2},para={すかのーはぬ},acc={H}]}%
	{%
		\MeaningsBlock{%
	\ryumeaning%
		{養う。飼う}{}{}%
	}%
	}%
	{\tailPrint[]}%
\entry
	{\headPrint[kana={すかは\tl{˥}\ws{}なるん\tl{˩˥}},ipa={su̥kaha naruɴ},pos={句},acc={H R}]}%
	{%
		\MeaningsBlock{%
	\ryumeaning%
		{近づく}{ただし，時間的なものにのみいう}{}%
	}%
	}%
	{\tailPrint[]}%
\entry
	{\headPrint[kana={すか\tl{˥}はん},ipa={su̥kahaɴ},pos={形},acc={H}]}%
	{%
		\MeaningsBlock{%
	\ryumeaning%
		{近い。近距離}{}{}%
	}%
	}%
	{\tailPrint[]}%
\entry
	{\headPrint[kana={すかふちるん\tl{˥}},ipa={su̥kafu̥tʃiruɴ},pos={動},rui={3},para={すかふとぅぬ},acc={H}]}%
	{%
		\MeaningsBlock{%
	\ryumeaning%
		{おっかぶせる}{}{}%
	}%
	}%
	{\tailPrint[]}%
\entry
	{\headPrint[kana={すかま\tl{˥}},ipa={su̥kama},pos={名},acc={H}]}%
	{%
		\MeaningsBlock{%
	\ryumeaning%
		{仕事。労役。業務}{}{}%
	}%
	}%
	{\tailPrint[]}%
\entry
	{\headPrint[kana={すかみち\tl{˥}},ipa={su̥kamitʃi},pos={名},acc={H}]}%
	{%
		\MeaningsBlock{%
	\ryumeaning%
		{近道}{}{}%
	}%
	}%
	{\tailPrint[]}%
\entry
	{\headPrint[kana={すかみんだすん\tl{˥}},ipa={su̥kamindasuɴ},pos={動},rui={2},para={すかみんだはぬ},acc={H}]}%
	{%
		\MeaningsBlock{%
	\ryumeaning%
		{掴み出す}{}{}%
	}%
	}%
	{\tailPrint[]}%
\entry
	{\headPrint[kana={すかむん\tl{˥}},ipa={su̥kamuɴ},pos={動},rui={1},para={すかまぬ},acc={H}]}%
	{%
		\MeaningsBlock{%
	\ryumeaning%
		{\ruby[g]{掴}{つか}む。捕える}{}{}%
	}%
	}%
	{\tailPrint[]}%
\entry
	{\headPrint[kana={すかり\tl{˥˩}すくん\tl{˥}},ipa={su̥karisu̥kuɴ},pos={動},rui={3},para={すかりすくぬ},acc={F+H}]}%
	{%
		\MeaningsBlock{%
	\ryumeaning%
		{叱りつける。戒める。躾ける}{}{}%
	}%
	}%
	{\tailPrint[]}%
\entry
	{\headPrint[kana={すかるん\tl{˥˩}},ipa={su̥karuɴ},pos={動},rui={1},para={すからぬ},acc={F}]}%
	{%
		\MeaningsBlock{%
	\ryumeaning%
		{漬かる。浸る}{}{}%
	}%
	}%
	{\tailPrint[]}%
\entry
	{\headPrint[kana={すぎゃん\tl{˥}},ipa={sugjaɴ},pos={動（継）},rui={3},para={すぐぬ},acc={H}]}%
	{%
		\MeaningsBlock{%
	\ryumeaning%
		{過ぎ去る}{}{}%
	}%
	}%
	{\tailPrint[]}%
\entry
	{\headPrint[kana={すく\tl{˥˩}},ipa={su̥ku},pos={名},acc={F}]}%
	{%
		\MeaningsBlock{%
	\ryumeaning%
		{底。奥}{}{}%
	}%
	}%
	{\tailPrint[]}%
\entry
	{\headPrint[kana={すくく\tl{˥}},ipa={su̥kuku},pos={名},acc={H}]}%
	{%
		\MeaningsBlock{%
	\ryumeaning%
		{梟}{}{}%
	}%
	}%
	{\tailPrint[]}%
\entry
	{\headPrint[kana={すくっふゎすん\tl{˥˩}},ipa={su̥kuffwasuɴ},pos={動},rui={2b},para={すくっふゎさぬ},acc={F}]}%
	{%
		\MeaningsBlock{%
	\ryumeaning%
		{ぶち壊す}{}{}%
	}%
	}%
	{\tailPrint[]}%
\entry
	{\headPrint[kana={すくび\tl{˥˩}},ipa={su̥kubi},pos={名},acc={F}]}%
	{%
		\MeaningsBlock{%
	\ryumeaning%
		{\ruby[g]{帯}{おび}}{}{}%
	}%
	}%
	{\tailPrint[]}%
\entry
	{\headPrint[kana={すくぶ\tl{˥}},ipa={su̥kubu},pos={名},acc={H}]}%
	{%
		\MeaningsBlock{%
	\ryumeaning%
		{\ruby[g]{籾}{もみ}\ruby[g]{殻}{がら}}{}{}%
	}%
	}%
	{\tailPrint[]}%
\entry
	{\headPrint[kana={すくふぁすん\tl{˥˩}},ipa={su̥kufasuɴ},pos={動},rui={2},para={すくふぁさぬ},acc={F}]}%
	{%
		\MeaningsBlock{%
	\ryumeaning%
		{強く振る。揺らす}{}{}%
	}%
	}%
	{\tailPrint[]}%
\entry
	{\headPrint[kana={すくまるん\tl{˥}},ipa={su̥kumaruɴ},pos={動},acc={H}]}%
	{%
		\MeaningsBlock{%
	\ryumeaning%
		{縮こまる。小さくなる}{}{}%
	\ryumeaning%
		{\ruby[g]{屈}{かが}む}{}{}%
	}%
	}%
	{\tailPrint[]}%
\entry
	{\headPrint[kana={すくまん\tl{˥}},ipa={su̥kumaɴ},pos={名},acc={H}]}%
	{%
		\MeaningsBlock{%
	\ryumeaning%
		{初穂祭}{}{}%
	}%
	}%
	{\tailPrint[]}%
\entry
	{\headPrint[kana={すくむん},ipa={su̥kumuɴ},pos={動},rui={1},para={すかまぬ}]}%
	{%
		\MeaningsBlock{%
	\ryumeaning%
		{すくむ}{身動きが出来なくなる}{}%
	}%
	}%
	{\tailPrint[]}%
\entry
	{\headPrint[kana={すくり\tl{˥}},ipa={su̥kuri},pos={名},acc={H}]}%
	{%
		\MeaningsBlock{%
	\ryumeaning%
		{造り。構造}{}{}%
	}%
	}%
	{\tailPrint[]}%
\entry
	{\headPrint[kana={すぐりん\tl{˥˩}},ipa={suguriɴ},pos={動},rui={3},acc={F}]}%
	{%
		\MeaningsBlock{%
	\ryumeaning%
		{優れる。勝る}{}{}%
	}%
	}%
	{\tailPrint[]}%
\entry
	{\headPrint[kana={すくるん\tl{˥}},ipa={su̥kuruɴ},pos={動},rui={1},para={すくらぬ},acc={H}]}%
	{%
		\MeaningsBlock{%
	\ryumeaning%
		{作る。作成する}{}{}%
	}%
	}%
	{\tailPrint[]}%
\entry
	{\headPrint[kana={すくん\tl{˥}},ipa={su̥kuɴ},pos={動},rui={1},para={すかぬ},acc={H}]}%
	{%
		\MeaningsBlock{%
	\ryumeaning%
		{着く。到着する}{}{}%
	}%
	}%
	{\tailPrint[]}%
\entry
	{\headPrint[kana={すくん\tl{˥}},ipa={su̥kuɴ},pos={動},rui={1},para={すかぬ},acc={H}]}%
	{%
		\MeaningsBlock{%
	\ryumeaning%
		{置く。置いておく}{}{}%
	}%
	}%
	{\tailPrint[]}%
\entry
	{\headPrint[kana={すくん\tl{˥}},ipa={su̥kuɴ},pos={動},rui={1},para={すかぬ},acc={H}]}%
	{%
		\MeaningsBlock{%
	\ryumeaning%
		{搗く}{玄米をついて白米にする}{}%
	}%
	}%
	{\tailPrint[]}%
\entry
	{\headPrint[kana={すくん\tl{˥}},ipa={su̥kuɴ},pos={動},rui={1},para={すかぬ},acc={H}]}%
	{%
		\MeaningsBlock{%
	\ryumeaning%
		{好く。好む。好きだ}{}{}%
	}%
	}%
	{\tailPrint[]}%
\entry
	{\headPrint[kana={すくん\tl{˥˩}},ipa={su̥kuɴ},pos={動},rui={1},para={すかぬ},acc={F}]}%
	{%
		\MeaningsBlock{%
	\ryumeaning%
		{聞く}{}{}%
	}%
	}%
	{\tailPrint[]}%
\entry
	{\headPrint[kana={すくん\tl{˥˩}},ipa={su̥kuɴ},pos={動},rui={1},para={すかぬ},acc={F}]}%
	{%
		\MeaningsBlock{%
	\ryumeaning%
		{敷く}{}{}%
	}%
	}%
	{\tailPrint[]}%
\entry
	{\headPrint[kana={すくん\tl{˥˩}},ipa={su̥kuɴ},pos={動},rui={1},para={すかぬ},acc={F}]}%
	{%
		\MeaningsBlock{%
	\ryumeaning%
		{突く}{}{}%
	}%
	}%
	{\tailPrint[]}%
\entry
	{\headPrint[kana={すけーとー\tl{˥˩}\ws{}あん\tl{˥}},ipa={su̥keːtoː aɴ},pos={句},acc={F H}]}%
	{%
		\MeaningsBlock{%
	\ryumeaning%
		{使い出がある}{}{}%
	}%
	}%
	{\tailPrint[]}%
\entry
	{\headPrint[kana={すこーすん\tl{˥˩}},ipa={su̥koːsuɴ},pos={動},rui={3},para={すこーはぬ},acc={F}]}%
	{%
		\MeaningsBlock{%
	\ryumeaning%
		{\ruby[g]{掬}{すく}う}{}{}%
	}%
	}%
	{\tailPrint[]}%
\entry
	{\headPrint[kana={すこーん\tl{˥˩}},ipa={su̥koːɴ},pos={動},rui={1},para={すかーぬ},acc={F}]}%
	{%
		\MeaningsBlock{%
	\ryumeaning%
		{使う。使用する。費やす}{}{}%
	}%
	}%
	{\tailPrint[]}%
\entry
	{\headPrint[kana={すさーらすん\tl{˥}},ipa={ssaːrasuɴ},pos={動},rui={2},para={すさーらはぬ},acc={H}]}%
	{%
		\MeaningsBlock{%
	\ryumeaning%
		{垂らす。ぶら下げる}{}{}%
	}%
	}%
	{\tailPrint[]}%
\entry
	{\headPrint[kana={すさーり\tl{˥}},ipa={ssaːri},pos={名},acc={H}]}%
	{%
		\MeaningsBlock{%
	\ryumeaning%
		{白蟻}{}{}%
	}%
	}%
	{\tailPrint[]}%
\entry
	{\headPrint[kana={すさーるん},ipa={ssaruɴ},pos={動},rui={1},para={すさーらぬ},acc={-}]}%
	{%
		\MeaningsBlock{%
	\ryumeaning%
		{垂れる。垂れ下がる}{}{}%
	}%
	}%
	{\tailPrint[]}%
\entry
	{\headPrint[kana={すさぎるん\tl{˥}},ipa={ssagiruɴ},pos={動},rui={3},para={すさぐぬ},acc={H}]}%
	{%
		\MeaningsBlock{%
	\ryumeaning%
		{（臼・杵で）搗く}{}{}%
	}%
	}%
	{\tailPrint[]}%
\entry
	{\headPrint[kana={すさざぎ\tl{˥}},ipa={ssadzagi},pos={名},acc={H}]}%
	{%
		\MeaningsBlock{%
	\ryumeaning%
		{モンパノキ}{海浜の植物}{}%
	}%
	}%
	{\tailPrint[]}%
\entry
	{\headPrint[kana={すさ\tl{˥}はん},ipa={ssahaɴ},pos={形},para={すさへぬ},acc={H}]}%
	{%
		\MeaningsBlock{%
	\ryumeaning%
		{強い。強力だ。優勢だ}{}{}%
	}%
	}%
	{\tailPrint[]}%
\entry
	{\headPrint[kana={すさびるん\tl{˥}},ipa={ssabiruɴ},pos={動},rui={3},para={すさぶぬ},acc={H}]}%
	{%
		\MeaningsBlock{%
	\ryumeaning%
		{調べる}{}{}%
	}%
	}%
	{\tailPrint[]}%
\entry
	{\headPrint[kana={すさぶ\tl{˥}},ipa={ssabu},pos={名},acc={H}]}%
	{%
		\MeaningsBlock{%
	\ryumeaning%
		{\ruby[g]{白保}{しらほ}}{石垣島の集落名。海岸に位置し、1771年の津波で集落がほとんど壊滅し、その後、波照間島からの人で再建された。現在の八重山諸方言の中で白保方言が波照間方言にもっと近い}{}%
	}%
	}%
	{\tailPrint[]}%
\entry
	{\headPrint[kana={すさむん\tl{˥}},ipa={ssamuɴ},pos={動},rui={1},para={すさまぬ},acc={H}]}%
	{%
		\MeaningsBlock{%
	\ryumeaning%
		{白む。白くなる。白みを帯びる。色があせる}{}{}%
	}%
	}%
	{\tailPrint[]}%
\entry
	{\headPrint[kana={すさりぶち\tl{˥˩}},ipa={ssaributʃi},pos={名},acc={F}]}%
	{%
		\MeaningsBlock{%
	\ryumeaning%
		{挨拶。口上}{}{}%
	}%
	}%
	{\tailPrint[]}%
\entry
	{\headPrint[kana={すさりん\tl{˥˩}},ipa={ssariɴ},pos={動},rui={3},para={すさるぬ},acc={F}]}%
	{%
		\MeaningsBlock{%
	\ryumeaning%
		{申し上げる}{謙譲語}{}%
	}%
	}%
	{\tailPrint[]}%
\entry
	{\headPrint[kana={すされー},ipa={ssareː},pos={感}]}%
	{%
		\MeaningsBlock{%
	\ryumeaning%
		{ご免下さい}{他家訪問時の声かけ}{}%
	}%
	}%
	{\tailPrint[]}%
\entry
	{\headPrint[kana={すさんち\tl{˥}},ipa={ssantʃi},pos={名},acc={H}]}%
	{%
		\MeaningsBlock{%
	\ryumeaning%
		{秋。秋季}{}{}%
	}%
	}%
	{\tailPrint[]}%
\entry
	{\headPrint[kana={すするん\tl{˥˩}},ipa={ssuruɴ},pos={動},rui={1},para={すすらぬ},acc={F}]}%
	{%
		\MeaningsBlock{%
	\ryumeaning%
		{拭く。\ruby[g]{擦}{こす}る}{}{}%
	}%
	}%
	{\tailPrint[]}%
\entry
	{\headPrint[kana={\josho{}すそー\kako{}し},ipa={ssoːʃi},pos={副},acc={HL}]}%
	{%
		\MeaningsBlock{%
	\ryumeaning%
		{白い。白っぽい}{}{}%
	}%
	}%
	{\tailPrint[]}%
\entry
	{\headPrint[kana={す\josho{}たー\kako{}すたー},ipa={su̥taːsu̥taː},pos={副},acc={重複}]}%
	{%
		\MeaningsBlock{%
	\ryumeaning%
		{度々}{}{}%
	}%
	}%
	{\tailPrint[]}%
\entry
	{\headPrint[kana={すたうい},ipa={su̥ta.ui},pos={名}]}%
	{%
		\MeaningsBlock{%
	\ryumeaning%
		{上下逆}{}{}%
	}%
	}%
	{\tailPrint[]}%
\entry
	{\headPrint[kana={すだちるん\tl{˥}},ipa={sudatʃiruɴ},pos={動},acc={H}]}%
	{%
		\MeaningsBlock{%
	\ryumeaning%
		{育てる}{}{}%
	}%
	}%
	{\tailPrint[]}%
\entry
	{\headPrint[kana={すだつん\tl{˥}},ipa={sudatsuɴ},pos={動},para={すだたぬ},acc={H}]}%
	{%
		\MeaningsBlock{%
	\ryumeaning%
		{育つ}{}{}%
	}%
	}%
	{\tailPrint[]}%
\entry
	{\headPrint[kana={すたはこち\tl{˥˩}},ipa={su̥tahḁkotʃi},pos={名},acc={F}]}%
	{%
		\MeaningsBlock{%
	\ryumeaning%
		{下あご}{}{}%
	}%
	}%
	{\tailPrint[]}%
\entry
	{\headPrint[kana={すたみん\tl{˥}},ipa={su̥tamiɴ},pos={名},acc={H}]}%
	{%
		\MeaningsBlock{%
	\ryumeaning%
		{カタツムリ}{}{}%
	}%
	}%
	{\tailPrint[]}%
\entry
	{\headPrint[kana={すたらすん\tl{˥}},ipa={su̥tarasuɴ},pos={動},rui={2},para={すたらはぬ},acc={H}]}%
	{%
		\MeaningsBlock{%
	\ryumeaning%
		{滴らせる。水を切る}{}{}%
	}%
	}%
	{\tailPrint[]}%
\entry
	{\headPrint[kana={すたるん\tl{˥}},ipa={su̥taruɴ},pos={動},rui={1},para={すたらぬ},acc={H}]}%
	{%
		\MeaningsBlock{%
	\ryumeaning%
		{（雫が）垂れる。滴る。湿る}{}{}%
	}%
	}%
	{\tailPrint[]}%
\entry
	{\headPrint[kana={すっす\tl{˥}},ipa={sussu},pos={名},acc={H}]}%
	{%
		\MeaningsBlock{%
	\ryumeaning%
		{\ruby[g]{裾}{すそ}}{}{}%
	}%
	}%
	{\tailPrint[]}%
\entry
	{\headPrint[kana={すっち\tl{˥}},ipa={suttʃi},pos={名},acc={H}]}%
	{%
		\MeaningsBlock{%
	\ryumeaning%
		{潮時。干潮時}{}{}%
	}%
	}%
	{\tailPrint[]}%
\entry
	{\headPrint[kana={すとーま\tl{˥˩}},ipa={su̥toːma},pos={名},acc={F}]}%
	{%
		\MeaningsBlock{%
	\ryumeaning%
		{\ruby[g]{風}{かざ}\ruby[g]{下}{しも}}{}{}%
	}%
	}%
	{\tailPrint[]}%
\entry
	{\headPrint[kana={すとぅち\tl{˥}},ipa={su̥tutʃi},pos={名},acc={H}]}%
	{%
		\MeaningsBlock{%
	\ryumeaning%
		{\ruby[g]{蘇鉄}{そてつ}}{備荒植物}{}%
	}%
	}%
	{\tailPrint[]}%
\entry
	{\headPrint[kana={すとぅみるん\tl{˥}},ipa={su̥tumiruɴ},pos={動},rui={3},para={すとぅむぬ},acc={H}]}%
	{%
		\MeaningsBlock{%
	\ryumeaning%
		{勤める。勤務する}{}{}%
	}%
	}%
	{\tailPrint[]}%
\entry
	{\headPrint[kana={すとぅむち\tl{˥}},ipa={su̥tumutʃi},pos={名},acc={H}]}%
	{%
		\MeaningsBlock{%
	\ryumeaning%
		{朝}{}{}%
	}%
	}%
	{\tailPrint[]}%
\entry
	{\headPrint[kana={すとぅりすとぅり},ipa={su̥turisu̥turi},pos={擬}]}%
	{%
		\MeaningsBlock{%
	\ryumeaning%
		{ぽたぽた。しとしと}{雨の静かに降るさま}{}%
	}%
	}%
	{\tailPrint[]}%
\entry
	{\headPrint[kana={すとぅる\tl{˥}さーん},ipa={su̥turusaːɴ},pos={形},acc={H}]}%
	{%
		\MeaningsBlock{%
	\ryumeaning%
		{うら寂しい}{}{}%
	}%
	}%
	{\tailPrint[]}%
\entry
	{\headPrint[kana={すな\tl{˥}},ipa={su̥na},pos={名},acc={H}]}%
	{%
		\MeaningsBlock{%
	\ryumeaning%
		{綱。縄}{}{}%
	}%
	}%
	{\tailPrint[]}%
\entry
	{\headPrint[kana={すながらすん\tl{˥˩}},ipa={su̥nagarasuɴ},pos={動},rui={2},para={すながらはぬ},acc={F}]}%
	{%
		\MeaningsBlock{%
	\ryumeaning%
		{続かせる。連ねる}{}{}%
	}%
	}%
	{\tailPrint[]}%
\entry
	{\headPrint[kana={すながるん\tl{˥˩}},ipa={su̥nagaruɴ},pos={動},acc={F}]}%
	{%
		\MeaningsBlock{%
	\ryumeaning%
		{連ねる}{}{}%
	\ryumeaning%
		{並べる}{}{}%
	}%
	}%
	{\tailPrint[]}%
\entry
	{\headPrint[kana={すなすん\tl{˥˩}},ipa={su̥nasuɴ},pos={動},rui={2},para={すなはぬ},acc={F}]}%
	{%
		\MeaningsBlock{%
	\ryumeaning%
		{殺す}{「死なす」の意}{}%
	\ryumeaning%
		{殴る}{}{}%
	}%
	}%
	{\tailPrint[]}%
\entry
	{\headPrint[kana={すなちるん\tl{˥}},ipa={su̥natʃiruɴ},pos={動},rui={3},para={すなとはぬ},acc={H}]}%
	{%
		\MeaningsBlock{%
	\ryumeaning%
		{育てる}{}{}%
	}%
	}%
	{\tailPrint[]}%
\entry
	{\headPrint[kana={すねー\tl{˥}},ipa={su̥neː},pos={名},acc={H}]}%
	{%
		\MeaningsBlock{%
	\ryumeaning%
		{曽根。浅堆。バンク}{}{}%
	}%
	}%
	{\tailPrint[]}%
\entry
	{\headPrint[kana={すぱしきるん\tl{˥˩}},ipa={su̥paʃi̥kiruɴ},pos={動},rui={1?},para={すぱすかぬ},acc={F}]}%
	{%
		\MeaningsBlock{%
	\ryumeaning%
		{くっ付ける。接着する}{}{}%
	}%
	}%
	{\tailPrint[]}%
\entry
	{\headPrint[kana={すぱに\tl{˥˩}},ipa={su̥pani},pos={名},acc={F}]}%
	{%
		\MeaningsBlock{%
	\ryumeaning%
		{岩礒。岩岸壁}{}{}%
	}%
	}%
	{\tailPrint[]}%
\entry
	{\headPrint[kana={すぱみるん\tl{˥}},ipa={su̥pamiruɴ},pos={動},rui={3},acc={H}]}%
	{%
		\MeaningsBlock{%
	\ryumeaning%
		{狭める}{}{}%
	}%
	}%
	{\tailPrint[]}%
\entry
	{\headPrint[kana={すぱんこっち},ipa={su̥paŋkottʃi},pos={名}]}%
	{%
		\MeaningsBlock{%
	\ryumeaning%
		{ミズガンピ。浜紫檀}{植物名}{}%
	}%
	}%
	{\tailPrint[]}%
\entry
	{\headPrint[kana={すび\tl{˥˩}},ipa={subi},pos={名},acc={F}]}%
	{%
		\MeaningsBlock{%
	\ryumeaning%
		{礒。荒磯}{}{}%
	}%
	}%
	{\tailPrint[]}%
\entry
	{\headPrint[kana={すぷしん\tl{˥˩}},ipa={su̥puʃiɴ},pos={名},acc={F}]}%
	{%
		\MeaningsBlock{%
	\ryumeaning%
		{\ruby[g]{膝}{ひざ}}{}{}%
	}%
	}%
	{\tailPrint[]}%
\entry
	{\headPrint[kana={すぷた\tl{˥}はん},ipa={su̥putahaɴ},pos={形},para={すぷたへぬ},acc={H}]}%
	{%
		\MeaningsBlock{%
	\ryumeaning%
		{不潔だ。汚い}{}{}%
	}%
	}%
	{\tailPrint[]}%
\entry
	{\headPrint[kana={すぷ\tl{˥}つぁーん},ipa={suputsaːɴ},pos={形},acc={H}]}%
	{%
		\MeaningsBlock{%
	\ryumeaning%
		{不潔だ}{}{}%
	}%
	}%
	{\tailPrint[]}%
\entry
	{\headPrint[kana={すぶりやみ},ipa={su̥burijami},pos={動}]}%
	{%
		\MeaningsBlock{%
	\ryumeaning%
		{腹がしぼるように痛む}{}{}%
	}%
	}%
	{\tailPrint[]}%
\entry
	{\headPrint[kana={すぷりん\tl{˥}},ipa={su̥puriɴ},pos={名},acc={H}]}%
	{%
		\MeaningsBlock{%
	\ryumeaning%
		{\ruby[g]{冬瓜}{とうがん}}{}{}%
	}%
	}%
	{\tailPrint[]}%
\entry
	{\headPrint[kana={すぷるいし\tl{˥}},ipa={su̥puru.iʃi},pos={名},acc={H}]}%
	{%
		\MeaningsBlock{%
	\ryumeaning%
		{アザミサンゴ}{}{}%
	}%
	}%
	{\tailPrint[]}%
\entry
	{\headPrint[kana={すぷるん\tl{˥}},ipa={su̥puruɴ},pos={動},rui={1},acc={H}]}%
	{%
		\MeaningsBlock{%
	\ryumeaning%
		{\ruby[g]{搾}{しぼ}る}{}{}%
	}%
	}%
	{\tailPrint[]}%
\entry
	{\headPrint[kana={すぷるん\tl{˥˩}},ipa={su̥puruɴ},pos={動},rui={1},para={すぷらぬ},acc={F}]}%
	{%
		\MeaningsBlock{%
	\ryumeaning%
		{すする。\ruby[g]{舐}{な}める}{}{}%
	}%
	}%
	{\tailPrint[]}%
\entry
	{\headPrint[kana={すぷるん\tl{˥˩}},ipa={su̥puruɴ},pos={動},rui={1},para={すぷらぬ},acc={F}]}%
	{%
		\MeaningsBlock{%
	\ryumeaning%
		{吸う。吸い取る}{}{}%
	}%
	}%
	{\tailPrint[]}%
\entry
	{\headPrint[kana={すぽ\tl{˥}はん},ipa={su̥pohaɴ},pos={形},para={すぽへぬ},acc={H}]}%
	{%
		\MeaningsBlock{%
	\ryumeaning%
		{渋い。渋みがある}{}{}%
	}%
	}%
	{\tailPrint[]}%
\entry
	{\headPrint[kana={すまじしゃ\tl{˥}はん},ipa={su̥madʒiʃahaɴ},pos={形},para={しまじへぬ},acc={H}]}%
	{%
		\MeaningsBlock{%
	\ryumeaning%
		{コクのある旨味}{}{}%
	}%
	}%
	{\tailPrint[]}%
\entry
	{\headPrint[kana={すますん\tl{˥}},ipa={su̥masuɴ},pos={動},rui={2},para={すまはぬ},acc={H}]}%
	{%
		\MeaningsBlock{%
	\ryumeaning%
		{済ます。済ませる。終える。し遂げる}{}{}%
	}%
	}%
	{\tailPrint[]}%
\entry
	{\headPrint[kana={すますん\tl{˥}},ipa={su̥masuɴ},pos={動},rui={2},para={すまはぬ},acc={H}]}%
	{%
		\MeaningsBlock{%
	\ryumeaning%
		{潜らせる。沈ませる}{}{}%
	}%
	}%
	{\tailPrint[]}%
\entry
	{\headPrint[kana={すますん\tl{˥}},ipa={su̥masuɴ},pos={動},rui={2},para={すまはぬ},acc={H}]}%
	{%
		\MeaningsBlock{%
	\ryumeaning%
		{澄ます}{}{}%
	}%
	}%
	{\tailPrint[]}%
\entry
	{\headPrint[kana={すまっち\ws{}すん},ipa={su̥mattʃi suɴ},pos={句}]}%
	{%
		\MeaningsBlock{%
	\ryumeaning%
		{粗末にする。ぞんざいに扱う}{}{}%
	}%
	}%
	{\tailPrint[]}%
\entry
	{\headPrint[kana={すまどぅるん\tl{˥˩}},ipa={su̥maburuɴ},pos={動},rui={1},para={すまどぅらぬ},acc={F}]}%
	{%
		\MeaningsBlock{%
	\ryumeaning%
		{迷う。とまどう}{}{}%
	}%
	}%
	{\tailPrint[]}%
\entry
	{\headPrint[kana={すま\tl{˥}むに\tl{˩˥}},ipa={su̥mamuni},pos={名},acc={H+R}]}%
	{%
		\MeaningsBlock{%
	\ryumeaning%
		{島言葉。方言}{}{}%
	}%
	}%
	{\tailPrint[]}%
\entry
	{\headPrint[kana={すまるん\tl{˥˩}},ipa={su̥maruɴ},pos={動},rui={1},para={すまらぬ},acc={F}]}%
	{%
		\MeaningsBlock{%
	\ryumeaning%
		{染まる}{}{}%
	}%
	}%
	{\tailPrint[]}%
\entry
	{\headPrint[kana={すむ\tl{˥}},ipa={su̥mu},pos={名},acc={H}]}%
	{%
		\MeaningsBlock{%
	\ryumeaning%
		{肝。心}{}{}%
	}%
	}%
	{\tailPrint[]}%
\entry
	{\headPrint[kana={すむあーり},ipa={su̥mu.aːri},pos={名}]}%
	{%
		\MeaningsBlock{%
	\ryumeaning%
		{胸騒ぎ。落ち着かない}{}{}%
	}%
	}%
	{\tailPrint[]}%
\entry
	{\headPrint[kana={すむ\tl{˥}いしゃが\tl{˥}はーん},ipa={su̥muiʃagahaːɴ},pos={形},acc={H+H}]}%
	{%
		\MeaningsBlock{%
	\ryumeaning%
		{小胆だ。気が小さい}{}{}%
	}%
	}%
	{\tailPrint[]}%
\entry
	{\headPrint[kana={すむ\tl{˥}いた\tl{˥}はん},ipa={su̥muitahaɴ},pos={形},para={すむいたへぬ},acc={H+H}]}%
	{%
		\MeaningsBlock{%
	\ryumeaning%
		{気の毒だ。可哀そうだ}{心が痛いの義}{}%
	}%
	}%
	{\tailPrint[]}%
\entry
	{\headPrint[kana={すむ\tl{˥}いり\tl{˥˩}},ipa={su̥muiri},pos={名},acc={H+F}]}%
	{%
		\MeaningsBlock{%
	\ryumeaning%
		{肝入り。熱心}{}{}%
	}%
	}%
	{\tailPrint[]}%
\entry
	{\headPrint[kana={すむぐくる\tl{˥}},ipa={su̥mugukuru},pos={名},acc={H}]}%
	{%
		\MeaningsBlock{%
	\ryumeaning%
		{心情。精神}{}{}%
	}%
	}%
	{\tailPrint[]}%
\entry
	{\headPrint[kana={すむ\tl{˥}ぐり\tl{˥}さん},ipa={su̥mugurisaɴ},pos={形},acc={H+H}]}%
	{%
		\MeaningsBlock{%
	\ryumeaning%
		{可哀そうだ。気の毒だ}{}{}%
	}%
	}%
	{\tailPrint[]}%
\entry
	{\headPrint[kana={すむ\tl{˥}ぐりしゃ\tl{˥}はーん},ipa={su̥muguriʃahaːɴ},pos={形},acc={H+H}]}%
	{%
		\MeaningsBlock{%
	\ryumeaning%
		{可哀そうそうな思いがする。みじめな思いがする}{}{}%
	}%
	}%
	{\tailPrint[]}%
\entry
	{\headPrint[kana={すむくん\tl{˥}},ipa={su̥mukuɴ},pos={動},acc={H}]}%
	{%
		\MeaningsBlock{%
	\ryumeaning%
		{背く}{}{}%
	}%
	}%
	{\tailPrint[]}%
\entry
	{\headPrint[kana={すむ\tl{˥}けしゃ\tl{˥}はーん},ipa={su̥mukeʃahaːɴ},pos={形},acc={H+H}]}%
	{%
		\MeaningsBlock{%
	\ryumeaning%
		{心が美しい。心が優しい}{}{}%
	}%
	}%
	{\tailPrint[]}%
\entry
	{\headPrint[kana={すむ\tl{˥}ずさ\tl{˥}はーん},ipa={su̥mudzusahaːɴ},pos={形},acc={H+H}]}%
	{%
		\MeaningsBlock{%
	\ryumeaning%
		{心強い。頼もしく思う}{}{}%
	}%
	}%
	{\tailPrint[]}%
\entry
	{\headPrint[kana={すむ\tl{˥}すさ\tl{˥}はん},ipa={su̥mususahaɴ},pos={形},acc={H+H}]}%
	{%
		\MeaningsBlock{%
	\ryumeaning%
		{心強い}{}{}%
	}%
	}%
	{\tailPrint[]}%
\entry
	{\headPrint[kana={すむだーりるん\tl{˥˩}},ipa={su̥mudaːruɴ},pos={動},acc={F}]}%
	{%
		\MeaningsBlock{%
	\ryumeaning%
		{悄げる。気落ちする}{}{}%
	}%
	}%
	{\tailPrint[]}%
\entry
	{\headPrint[kana={すむ\tl{˥}だが\tl{˥}はん},ipa={su̥mudagahaɴ},pos={形},para={すむだがへぬ},acc={H+H}]}%
	{%
		\MeaningsBlock{%
	\ryumeaning%
		{気位が高い}{}{}%
	}%
	}%
	{\tailPrint[]}%
\entry
	{\headPrint[kana={すむち\tl{˥˩}},ipa={su̥mutʃi},pos={名},acc={F}]}%
	{%
		\MeaningsBlock{%
	\ryumeaning%
		{書物。本}{}{}%
	}%
	}%
	{\tailPrint[]}%
\entry
	{\headPrint[kana={すむとぅ\tl{˥}},ipa={su̥mutu},pos={名},acc={H}]}%
	{%
		\MeaningsBlock{%
	\ryumeaning%
		{\ruby[g]{鞭}{むち}}{}{}%
	}%
	}%
	{\tailPrint[]}%
\entry
	{\headPrint[kana={すむ\tl{˥}\ws{}ねーぬ\tl{˩˥}},ipa={su̥mu neːnu},pos={句},acc={H R}]}%
	{%
		\MeaningsBlock{%
	\ryumeaning%
		{人情がない。薄情}{}{}%
	}%
	}%
	{\tailPrint[]}%
\entry
	{\headPrint[kana={すむぴらすん\tl{˥}},ipa={su̥mupi̥rasuɴ},pos={動},rui={2},para={すむぴらはぬ},acc={H}]}%
	{%
		\MeaningsBlock{%
	\ryumeaning%
		{打ち解ける}{}{}%
	\ryumeaning%
		{仲直りする}{}{}%
	}%
	}%
	{\tailPrint[]}%
\entry
	{\headPrint[kana={すむやむーん},ipa={su̥mujamuːɴ},pos={句}]}%
	{%
		\MeaningsBlock{%
	\ryumeaning%
		{心が痛む。胸の痛む思いがする。見るに忍びない。哀れだ}{}{}%
	}%
	}%
	{\tailPrint[]}%
\entry
	{\headPrint[kana={すむ\tl{˥}やむん\tl{˩˥}},ipa={su̥mujamuɴ},pos={動},rui={1},para={すむやまぬ},acc={H+R}]}%
	{%
		\MeaningsBlock{%
	\ryumeaning%
		{心配する。悩む}{}{}%
	}%
	}%
	{\tailPrint[]}%
\entry
	{\headPrint[kana={すむん\tl{˥}},ipa={su̥muɴ},pos={動},rui={1},para={すまぬ},acc={H}]}%
	{%
		\MeaningsBlock{%
	\ryumeaning%
		{潜る。沈む}{}{}%
	}%
	}%
	{\tailPrint[]}%
\entry
	{\headPrint[kana={すむん\tl{˥}},ipa={su̥muɴ},pos={動},rui={1},para={すまぬ},acc={H}]}%
	{%
		\MeaningsBlock{%
	\ryumeaning%
		{澄む}{}{}%
	}%
	}%
	{\tailPrint[]}%
\entry
	{\headPrint[kana={すむん\tl{˥}},ipa={su̥muɴ},pos={動},rui={1},para={すまぬ},acc={H}]}%
	{%
		\MeaningsBlock{%
	\ryumeaning%
		{済む。終わる}{}{}%
	}%
	}%
	{\tailPrint[]}%
\entry
	{\headPrint[kana={すむん\tl{˥}},ipa={su̥muɴ},pos={動},rui={1},para={すまぬ},acc={H}]}%
	{%
		\MeaningsBlock{%
	\ryumeaning%
		{包む}{}{}%
	}%
	}%
	{\tailPrint[]}%
\entry
	{\headPrint[kana={すむん\tl{˥}},ipa={su̥muɴ},pos={動},rui={1},para={すまぬ},acc={H}]}%
	{%
		\MeaningsBlock{%
	\ryumeaning%
		{摘む。摘み切る。千切る}{}{}%
	\ryumeaning%
		{つねる}{}{}%
	}%
	}%
	{\tailPrint[]}%
\entry
	{\headPrint[kana={すむん\tl{˥˩}},ipa={su̥muɴ},pos={動},rui={1},para={すまぬ},acc={F}]}%
	{%
		\MeaningsBlock{%
	\ryumeaning%
		{積む。積み重ねる}{}{}%
	}%
	}%
	{\tailPrint[]}%
\entry
	{\headPrint[kana={すら\tl{˥}},ipa={sura},pos={名},acc={H}]}%
	{%
		\MeaningsBlock{%
	\ryumeaning%
		{先。先端}{}{}%
	}%
	}%
	{\tailPrint[]}%
\entry
	{\headPrint[kana={するいるん\tl{˥}},ipa={su̥ruiruɴ},pos={動},acc={H}]}%
	{%
		\MeaningsBlock{%
	\ryumeaning%
		{揃える。準備する}{}{}%
	\ryumeaning%
		{列を整える。行列をつくる}{}{}%
	\ryumeaning%
		{集まらせる。集合させる}{}{}%
	}%
	}%
	{\tailPrint[]}%
\entry
	{\headPrint[kana={するいん\tl{˥}},ipa={suruiɴ},pos={動},rui={1},acc={H}]}%
	{%
		\MeaningsBlock{%
	\ryumeaning%
		{揃える}{}{}%
	}%
	}%
	{\tailPrint[]}%
\entry
	{\headPrint[kana={するん\tl{˥}},ipa={su̥ruɴ},pos={動},rui={1},para={すらぬ},acc={H}]}%
	{%
		\MeaningsBlock{%
	\ryumeaning%
		{剃る}{}{}%
	}%
	}%
	{\tailPrint[]}%
\entry
	{\headPrint[kana={すん\tl{˥}},ipa={suɴ},pos={名},acc={H}]}%
	{%
		\MeaningsBlock{%
	\ryumeaning%
		{\ruby[g]{損}{そん}}{}{}%
	}%
	}%
	{\tailPrint[]}%
\entry
	{\headPrint[kana={すん\tl{˥˩}},ipa={suɴ},pos={動},rui={2b},para={さぬ},acc={F}]}%
	{%
		\MeaningsBlock{%
	\ryumeaning%
		{する。やる}{}{}%
	}%
	}%
	{\tailPrint[]}%
\entry
	{\headPrint[kana={すんぐるん\tl{˥}},ipa={suŋguruɴ},pos={動},rui={1},para={すんぐらぬ},acc={H}]}%
	{%
		\MeaningsBlock{%
	\ryumeaning%
		{\ruby[g]{鞭}{むち}打つ}{}{}%
	}%
	}%
	{\tailPrint[]}%
\entry
	{\headPrint[kana={すんぽー\tl{˥}},ipa={sumpoː},pos={名},acc={H}]}%
	{%
		\MeaningsBlock{%
	\ryumeaning%
		{寸法}{}{}%
	}%
	}%
	{\tailPrint[]}%
\LetterHead{せ}
\entry
	{\headPrint[kana={せー},ipa={seː},pos={感}]}%
	{%
		\MeaningsBlock{%
	\ryumeaning%
		{さあー}{行動始めの掛け声}{}%
	}%
	}%
	{\tailPrint[]}%
\entry
	{\headPrint[kana={せー\tl{˥}},ipa={seː},pos={名},acc={H}]}%
	{%
		\MeaningsBlock{%
	\ryumeaning%
		{小エビ}{}{}%
	}%
	}%
	{\tailPrint[]}%
\entry
	{\headPrint[kana={せーぐ\tl{˥}},ipa={seːgu},pos={名},acc={H}]}%
	{%
		\MeaningsBlock{%
	\ryumeaning%
		{大工}{}{}%
	}%
	}%
	{\tailPrint[note={「細工」に対応}]}%
\entry
	{\headPrint[kana={せいろー\tl{˥}},ipa={seiroː},pos={名},acc={H}]}%
	{%
		\MeaningsBlock{%
	\ryumeaning%
		{蒸籠}{蒸し器のこと}{}%
	}%
	}%
	{\tailPrint[]}%
\LetterHead{そ}
\entry
	{\headPrint[kana={そー\tl{˥}},ipa={soː},pos={名},acc={H}]}%
	{%
		\MeaningsBlock{%
	\ryumeaning%
		{正気。意識}{}{}%
	}%
	}%
	{\tailPrint[]}%
\entry
	{\headPrint[kana={そー\tl{˥}},ipa={soː},pos={名},acc={H}]}%
	{%
		\MeaningsBlock{%
	\ryumeaning%
		{\ruby[g]{竿}{さお}}{}{}%
	}%
	}%
	{\tailPrint[]}%
\entry
	{\headPrint[kana={ぞー\tl{˩˥}},ipa={dzoː},pos={名},acc={R}]}%
	{%
		\MeaningsBlock{%
	\ryumeaning%
		{\ruby[g]{門}{もん}}{}{}%
	}%
	}%
	{\tailPrint[]}%
\entry
	{\headPrint[kana={そーいり\tl{˥˩}},ipa={soː.iri},pos={名},acc={F}]}%
	{%
		\MeaningsBlock{%
	\ryumeaning%
		{利口。賢い}{}{}%
	}%
	}%
	{\tailPrint[]}%
\entry
	{\headPrint[kana={そーぎ\tl{˥}},ipa={soːgi},pos={名},acc={H}]}%
	{%
		\MeaningsBlock{%
	\ryumeaning%
		{\ruby[g]{箕}{み}}{}{}%
	}%
	}%
	{\tailPrint[]}%
\entry
	{\headPrint[kana={そーじ\tl{˥}},ipa={soːdʒi},pos={名},acc={H}]}%
	{%
		\MeaningsBlock{%
	\ryumeaning%
		{掃除}{}{}%
	}%
	}%
	{\tailPrint[]}%
\entry
	{\headPrint[kana={ぞーじ\tl{˩˥}},ipa={dzoːdʒi},pos={名},acc={R}]}%
	{%
		\MeaningsBlock{%
	\ryumeaning%
		{上手。上手い}{}{}%
	}%
	}%
	{\tailPrint[]}%
\entry
	{\headPrint[kana={そーしき\tl{˥˩}},ipa={soːʃi̥ki},pos={名},acc={F}]}%
	{%
		\MeaningsBlock{%
	\ryumeaning%
		{葬式}{}{}%
	}%
	}%
	{\tailPrint[]}%
\entry
	{\headPrint[kana={ぞーしき\tl{˥˩}},ipa={dzoːʃi̥ki},pos={名},acc={F}]}%
	{%
		\MeaningsBlock{%
	\ryumeaning%
		{炊事}{}{}%
	}%
	}%
	{\tailPrint[]}%
\entry
	{\headPrint[kana={そーどー\tl{˥˩}},ipa={soːdoː},pos={名},acc={F}]}%
	{%
		\MeaningsBlock{%
	\ryumeaning%
		{騒動。騒ぎ}{}{}%
	}%
	}%
	{\tailPrint[]}%
\entry
	{\headPrint[kana={ぞーとぅ\tl{˩˥}},ipa={dzoːtu},pos={名},acc={R}]}%
	{%
		\MeaningsBlock{%
	\ryumeaning%
		{上等。立派}{}{}%
	}%
	}%
	{\tailPrint[]}%
\entry
	{\headPrint[kana={ぞーぬ\tl{˩˥}},ipa={dzoːnu},pos={名},acc={R}]}%
	{%
		\MeaningsBlock{%
	\ryumeaning%
		{租税。年貢}{上納の義}{}%
	}%
	}%
	{\tailPrint[]}%
\entry
	{\headPrint[kana={そー\ws{}ねーぬ},ipa={soː neːnu},pos={句}]}%
	{%
		\MeaningsBlock{%
	\ryumeaning%
		{\ruby[g]{耄}{もう}\ruby[g]{碌}{ろく}している}{}{}%
	}%
	}%
	{\tailPrint[]}%
\entry
	{\headPrint[kana={そーべー\tl{˥}},ipa={soːbeː},pos={名},acc={H}]}%
	{%
		\MeaningsBlock{%
	\ryumeaning%
		{安っぽい。粗悪な}{}{}%
	}%
	}%
	{\tailPrint[]}%
\entry
	{\headPrint[kana={そーみん\tl{˥}},ipa={soːmiɴ},pos={名},acc={H}]}%
	{%
		\MeaningsBlock{%
	\ryumeaning%
		{\ruby[g]{素}{そう}\ruby[g]{麵}{めん}}{}{}%
	}%
	}%
	{\tailPrint[note={そうめんの転訛}]}%
\entry
	{\headPrint[kana={そーり\tl{˥˩}\ws{}くん\tl{˥}},ipa={soːri kuɴ},pos={句},acc={F H}]}%
	{%
		\MeaningsBlock{%
	\ryumeaning%
		{一緒に来る}{}{}%
	}%
	}%
	{\tailPrint[]}%
\entry
	{\headPrint[kana={ぞーりぞーり},ipa={dzoːridzoːri},pos={擬}]}%
	{%
		\MeaningsBlock{%
	\ryumeaning%
		{ザァザァ}{激しい雨のさま}{}%
	}%
	}%
	{\tailPrint[]}%
\entry
	{\headPrint[kana={そーりん\tl{˥}},ipa={soːriɴ},pos={名},acc={H}]}%
	{%
		\MeaningsBlock{%
	\ryumeaning%
		{お盆。旧盆}{}{}%
	}%
	}%
	{\tailPrint[]}%
\entry
	{\headPrint[kana={そーり\tl{˥˩}\ws{}んぐん\tl{˥˩}},ipa={soːri ŋguɴ},pos={句},acc={F F}]}%
	{%
		\MeaningsBlock{%
	\ryumeaning%
		{一緒に行く}{}{}%
	}%
	}%
	{\tailPrint[]}%
\entry
	{\headPrint[kana={そーるん\tl{˥˩}},ipa={soːruɴ},pos={動},rui={1},para={そーらぬ},acc={F}]}%
	{%
		\MeaningsBlock{%
	\ryumeaning%
		{連れる}{}{}%
	\ryumeaning%
		{娶る。妻にする}{}{}%
	}%
	}%
	{\tailPrint[]}%
\entry
	{\headPrint[kana={そーるん\tl{˥˩}},ipa={soːruɴ},pos={動},rui={1},para={そーらぬ},acc={F}]}%
	{%
		\MeaningsBlock{%
	\ryumeaning%
		{（ヘラで）雑草を根こそぎに取る}{}{}%
	}%
	}%
	{\tailPrint[]}%
\entry
	{\headPrint[kana={ぞっか\tl{˩˥}},ipa={dzokka},pos={名},acc={R}]}%
	{%
		\MeaningsBlock{%
	\ryumeaning%
		{\ruby[g]{急須}{きゅうす}}{}{}%
	}%
	}%
	{\tailPrint[]}%
\entry
	{\headPrint[kana={ぞっふゎらすん\tl{˥˩}},ipa={dzoffwarasuɴ},pos={動},rui={2},para={ぞっふぁらはぬ},acc={F}]}%
	{%
		\MeaningsBlock{%
	\ryumeaning%
		{濡らす}{}{}%
	}%
	}%
	{\tailPrint[]}%
\entry
	{\headPrint[kana={ぞっふゎりん\tl{˥˩}},ipa={dzoffwariɴ},pos={動},rui={3},para={ぞっふぁるぬ},acc={F}]}%
	{%
		\MeaningsBlock{%
	\ryumeaning%
		{濡れる}{}{}%
	}%
	}%
	{\tailPrint[]}%
\entry
	{\headPrint[kana={ぞっふゎるん\tl{˥˩}},ipa={dzoffwaruɴ},pos={動},rui={3},acc={F}]}%
	{%
		\MeaningsBlock{%
	\ryumeaning%
		{濡らす}{}{}%
	}%
	}%
	{\tailPrint[]}%
\entry
	{\headPrint[kana={そり\tl{˥˩}\ws{}んぐん\tl{˥˩}},ipa={sori ŋguɴ},pos={句},acc={F F}]}%
	{%
		\MeaningsBlock{%
	\ryumeaning%
		{連れていく}{}{}%
	}%
	}%
	{\tailPrint[]}%
\entry
	{\headPrint[kana={そるん\tl{˥˩}},ipa={soruɴ},pos={動},rui={1},para={そらぬ},acc={F}]}%
	{%
		\MeaningsBlock{%
	\ryumeaning%
		{（妻を）娶る}{}{}%
	}%
	}%
	{\tailPrint[]}%
\entry
	{\headPrint[kana={そるん\tl{˥˩}},ipa={soruɴ},pos={動},rui={1},para={そらぬ},acc={F}]}%
	{%
		\MeaningsBlock{%
	\ryumeaning%
		{（へらで）除草する}{}{}%
	}%
	}%
	{\tailPrint[]}%
\entry
	{\headPrint[kana={そんぐん\tl{˥˩}},ipa={soŋguɴ},pos={動},rui={1},para={そんがぬ},acc={F}]}%
	{%
		\MeaningsBlock{%
	\ryumeaning%
		{引きずる。連れていく}{}{}%
	}%
	}%
	{\tailPrint[]}%
\entry
	{\headPrint[kana={そん\tl{˥}さみゃ\tl{˩˥}},ipa={sonsamja},pos={名},acc={H+R}]}%
	{%
		\MeaningsBlock{%
	\ryumeaning%
		{白鷺}{鳥の名}{}%
	}%
	}%
	{\tailPrint[]}%
\entry
	{\headPrint[kana={そん\tl{˥}\ws{}すん\tl{˥˩}},ipa={soɴ suɴ},pos={句},acc={H F}]}%
	{%
		\MeaningsBlock{%
	\ryumeaning%
		{損する。欠損する}{}{}%
	}%
	}%
	{\tailPrint[]}%
\entry
	{\headPrint[kana={そんだん\tl{˥}},ipa={sondaɴ},pos={名},acc={H}]}%
	{%
		\MeaningsBlock{%
	\ryumeaning%
		{相談}{}{}%
	}%
	}%
	{\tailPrint[note={相談の転訛}]}%
\entry
	{\headPrint[kana={そんだん\tl{˥}\ws{}すん\tl{˥˩}},ipa={sondaɴ suɴ},pos={句},acc={H F}]}%
	{%
		\MeaningsBlock{%
	\ryumeaning%
		{相談する。言い合わせる}{}{}%
	}%
	}%
	{\tailPrint[]}%
\LetterHead{た}
\entry
	{\headPrint[kana={た},ipa={ta},pos={接尾}]}%
	{%
		\MeaningsBlock{%
	\ryumeaning%
		{～方}{\emb{あらた}「東方」、\emb{いらた}「西方」など方向を示す}{}%
	}%
	}%
	{\tailPrint[]}%
\entry
	{\headPrint[kana={たー\tl{˥˩}},ipa={taː},pos={名},acc={F}]}%
	{%
		\MeaningsBlock{%
	\ryumeaning%
		{どなた。誰}{敬語}{}%
	}%
	}%
	{\tailPrint[]}%
\entry
	{\headPrint[kana={だー\tl{˥˩}},ipa={daː},pos={名},acc={F}]}%
	{%
		\MeaningsBlock{%
	\ryumeaning%
		{お前。君}{年下へ使う}{}%
	}%
	}%
	{\tailPrint[]}%
\entry
	{\headPrint[kana={だーぐ\tl{˩˥}},ipa={daːgu},pos={名},acc={R}]}%
	{%
		\MeaningsBlock{%
	\ryumeaning%
		{道具}{}{}%
	}%
	}%
	{\tailPrint[]}%
\entry
	{\headPrint[kana={\josho{}たー\kako{}こ},ipa={taːko},pos={名},acc={H+L}]}%
	{%
		\MeaningsBlock{%
	\ryumeaning%
		{\ruby[g]{凧}{たこ}}{「たこ」の転訛}{}%
	}%
	}%
	{\tailPrint[]}%
\entry
	{\headPrint[kana={だーさ\tl{˩˥}はん},ipa={daːsahaɴ},pos={形},para={だーさへぬ},acc={R}]}%
	{%
		\MeaningsBlock{%
	\ryumeaning%
		{立派だ。獲物が多い}{}{}%
	}%
	}%
	{\tailPrint[]}%
\entry
	{\headPrint[kana={だーし\tl{˩˥}},ipa={daːʃi},pos={名},acc={R}]}%
	{%
		\MeaningsBlock{%
	\ryumeaning%
		{\ruby[g]{出汁}{だし}}{}{}%
	}%
	}%
	{\tailPrint[]}%
\entry
	{\headPrint[kana={だーびむぬ\tl{˩˥}},ipa={daːbimunu},pos={名},acc={R}]}%
	{%
		\MeaningsBlock{%
	\ryumeaning%
		{\ruby[g]{玩具}{おもちゃ}}{遊び道具}{}%
	}%
	}%
	{\tailPrint[]}%
\entry
	{\headPrint[kana={だーぶん\tl{˩˥}},ipa={daːbuɴ},pos={動},rui={1},para={だーばぬ},acc={R}]}%
	{%
		\MeaningsBlock{%
	\ryumeaning%
		{\ruby[g]{弄}{いじ}る。\ruby[g]{弄}{もてあそ}ぶ}{}{}%
	}%
	}%
	{\tailPrint[]}%
\entry
	{\headPrint[kana={たーむじぃ\tl{˥}},ipa={taːmudʒi},pos={名},acc={H}]}%
	{%
		\MeaningsBlock{%
	\ryumeaning%
		{田芋}{}{}%
	}%
	}%
	{\tailPrint[]}%
\entry
	{\headPrint[kana={たーら\tl{˥}},ipa={taːra},pos={名},acc={H}]}%
	{%
		\MeaningsBlock{%
	\ryumeaning%
		{\ruby[g]{俵}{たわら}}{}{}%
	}%
	}%
	{\tailPrint[]}%
\entry
	{\headPrint[kana={だーり},ipa={daːri},pos={-}]}%
	{%
		\MeaningsBlock{%
	\ryumeaning%
		{ぐったり。ぐだっと}{}{}%
	}%
	}%
	{\tailPrint[]}%
\entry
	{\headPrint[kana={だーりゃん\tl{˩˥}},ipa={daːrjaɴ},pos={動（継）},rui={3},para={だーるぬ},acc={R}]}%
	{%
		\MeaningsBlock{%
	\ryumeaning%
		{}{ぐたっとなって元気がなくなる}{}%
	}%
	}%
	{\tailPrint[]}%
\entry
	{\headPrint[kana={だーりるん\tl{˩˥}},ipa={daːriruɴ},pos={動},rui={3},para={だーるぬ},acc={R}]}%
	{%
		\MeaningsBlock{%
	\ryumeaning%
		{だれる。ぐったりする}{}{}%
	\ryumeaning%
		{萎える。（花などが）萎む。（植物の勢いが）衰える}{}{}%
	}%
	}%
	{\tailPrint[]}%
\entry
	{\headPrint[kana={たーりん\tl{˥˩}},ipa={taːriɴ},pos={動},rui={3},para={たーるぬ},acc={F}]}%
	{%
		\MeaningsBlock{%
	\ryumeaning%
		{寝入る。熟睡する}{}{}%
	}%
	}%
	{\tailPrint[]}%
\entry
	{\headPrint[kana={だーりん\tl{˩˥}},ipa={daːriɴ},pos={動},rui={3},para={だーるぬ},acc={R}]}%
	{%
		\MeaningsBlock{%
	\ryumeaning%
		{だれる。疲れる。（植物が）萎れる}{}{}%
	}%
	}%
	{\tailPrint[]}%
\entry
	{\headPrint[kana={だーんた},ipa={daːnta},pos={擬}]}%
	{%
		\MeaningsBlock{%
	\ryumeaning%
		{どんと。しっかりと。どすん}{物事をしっかり決める。\emb{だーんた\ws{}きみるん}「しっかり決める」のように使う}{}%
	}%
	}%
	{\tailPrint[]}%
\entry
	{\headPrint[kana={たい\tl{˥}\ws{}すん\tl{˥˩}},ipa={tai suɴ},pos={句},acc={H F}]}%
	{%
		\MeaningsBlock{%
	\ryumeaning%
		{対抗する。反抗する。敵対する}{}{}%
	\ryumeaning%
		{嫉妬する}{}{}%
	}%
	}%
	{\tailPrint[]}%
\entry
	{\headPrint[kana={だいばん\tl{˩˥}},ipa={daibaɴ},pos={名},acc={R}]}%
	{%
		\MeaningsBlock{%
	\ryumeaning%
		{大型のカツオ}{大判鰹}{}%
	}%
	}%
	{\tailPrint[]}%
\entry
	{\headPrint[kana={だいま\tl{˥˩}},ipa={daima},pos={名},acc={F}]}%
	{%
		\MeaningsBlock{%
	\ryumeaning%
		{君たち。お前ら}{}{}%
	}%
	}%
	{\tailPrint[]}%
\entry
	{\headPrint[kana={たいらぎるん\tl{˥}},ipa={tairagiruɴ},pos={動},acc={H}]}%
	{%
		\MeaningsBlock{%
	\ryumeaning%
		{（飲食物を残らず）飲食し尽くす}{}{}%
	}%
	}%
	{\tailPrint[]}%
\entry
	{\headPrint[kana={たか\tl{˥˩}},ipa={tḁka},pos={名},acc={F}]}%
	{%
		\MeaningsBlock{%
	\ryumeaning%
		{\ruby[g]{鸇}{サシバ}}{}{}%
	}%
	}%
	{\tailPrint[]}%
\entry
	{\headPrint[kana={た\josho{}か\kako{}た},ipa={tḁkata},pos={名},acc={特殊}]}%
	{%
		\MeaningsBlock{%
	\ryumeaning%
		{高地。台地}{}{}%
	}%
	}%
	{\tailPrint[]}%
\entry
	{\headPrint[kana={た\josho{}か\kako{}たかし},ipa={takatḁkaʃi̥},pos={副},acc={重複}]}%
	{%
		\MeaningsBlock{%
	\ryumeaning%
		{高々と。うず高く}{}{}%
	}%
	}%
	{\tailPrint[]}%
\entry
	{\headPrint[kana={たか\tl{˥}はん},ipa={tḁkahaɴ},pos={形},para={たかへぬ},acc={H}]}%
	{%
		\MeaningsBlock{%
	\ryumeaning%
		{高い}{}{}%
	}%
	}%
	{\tailPrint[]}%
\entry
	{\headPrint[kana={たかぴぃとぅ\tl{˥}},ipa={tḁkapi̥tu},pos={名},acc={H}]}%
	{%
		\MeaningsBlock{%
	\ryumeaning%
		{背の高い人。偉い人}{}{}%
	}%
	}%
	{\tailPrint[]}%
\entry
	{\headPrint[kana={たかぶ\tl{˥}},ipa={tḁkabu},pos={名},acc={H}]}%
	{%
		\MeaningsBlock{%
	\ryumeaning%
		{\ruby[g]{煙草}{たばこ}}{}{}%
	}%
	}%
	{\tailPrint[]}%
\entry
	{\headPrint[kana={たかふもん\tl{˥}},ipa={tḁkafu̥moɴ},pos={名},acc={H}]}%
	{%
		\MeaningsBlock{%
	\ryumeaning%
		{入道雲}{}{}%
	}%
	}%
	{\tailPrint[]}%
\entry
	{\headPrint[kana={たかぶるん\tl{˥}},ipa={tḁkaburuɴ},pos={動},acc={H}]}%
	{%
		\MeaningsBlock{%
	\ryumeaning%
		{高ぶる。威張る。高く止まる}{}{}%
	}%
	}%
	{\tailPrint[]}%
\entry
	{\headPrint[kana={たかぼん},ipa={tḁkaboɴ},pos={名}]}%
	{%
		\MeaningsBlock{%
	\ryumeaning%
		{年寄りが煙草を吸う道具}{}{}%
	}%
	}%
	{\tailPrint[]}%
\entry
	{\headPrint[kana={たから\tl{˥}},ipa={tḁkara},pos={名},acc={H}]}%
	{%
		\MeaningsBlock{%
	\ryumeaning%
		{宝}{}{}%
	}%
	}%
	{\tailPrint[]}%
\entry
	{\headPrint[kana={たき\tl{˥}},ipa={tḁki},pos={名},acc={H}]}%
	{%
		\MeaningsBlock{%
	\ryumeaning%
		{丈。身長}{}{}%
	}%
	}%
	{\tailPrint[]}%
\entry
	{\headPrint[kana={たき\tl{˥˩}},ipa={tḁki},pos={名},acc={F}]}%
	{%
		\MeaningsBlock{%
	\ryumeaning%
		{竹}{}{}%
	}%
	}%
	{\tailPrint[]}%
\entry
	{\headPrint[kana={だき\tl{˥˩}},ipa={daki},pos={名},acc={F}]}%
	{%
		\MeaningsBlock{%
	\ryumeaning%
		{岳。山}{}{}%
	}%
	}%
	{\tailPrint[]}%
\entry
	{\headPrint[kana={だぎ},ipa={dagi},pos={副}]}%
	{%
		\MeaningsBlock{%
	\ryumeaning%
		{その程度}{}{}%
	}%
	}%
	{\tailPrint[]}%
\entry
	{\headPrint[kana={だきすん\tl{˩˥}},ipa={daki̥suɴ},pos={動},rui={2},para={だきさぬ},acc={R}]}%
	{%
		\MeaningsBlock{%
	\ryumeaning%
		{（固い物、刃物で）叩き切る}{「切る」の強調}{}%
	}%
	}%
	{\tailPrint[]}%
\entry
	{\headPrint[kana={たきどぅん\tl{˥}},ipa={tḁkiduɴ},pos={名},acc={H}]}%
	{%
		\MeaningsBlock{%
	\ryumeaning%
		{竹富}{竹富島。八重山諸島の島の名}{}%
	}%
	}%
	{\tailPrint[]}%
\entry
	{\headPrint[kana={たきふんた},ipa={tḁkifunta},pos={名}]}%
	{%
		\MeaningsBlock{%
	\ryumeaning%
		{竹床}{}{}%
	}%
	}%
	{\tailPrint[]}%
\entry
	{\headPrint[kana={たぎらすん\tl{˥˩}},ipa={tagirasuɴ},pos={動},rui={2},para={たぎらはぬ},acc={F}]}%
	{%
		\MeaningsBlock{%
	\ryumeaning%
		{たぎらせる。煮えたぎらせる。沸騰させる}{}{}%
	}%
	}%
	{\tailPrint[]}%
\entry
	{\headPrint[kana={たぎるん\tl{˥˩}},ipa={tagiruɴ},pos={動},acc={F}]}%
	{%
		\MeaningsBlock{%
	\ryumeaning%
		{たぎる。煮えたぎる。沸騰する}{}{}%
	}%
	}%
	{\tailPrint[]}%
\entry
	{\headPrint[kana={たく\tl{˥}},ipa={tḁku},pos={名},acc={H}]}%
	{%
		\MeaningsBlock{%
	\ryumeaning%
		{蛸}{}{}%
	}%
	}%
	{\tailPrint[]}%
\entry
	{\headPrint[kana={だくむん\tl{˩˥}},ipa={daku̥muɴ},pos={動},rui={1},para={だくまぬ},acc={R}]}%
	{%
		\MeaningsBlock{%
	\ryumeaning%
		{ぶち込む。投げ入れる。放り込む}{}{}%
	}%
	}%
	{\tailPrint[]}%
\entry
	{\headPrint[kana={たくらむん},ipa={takuramuɴ},pos={動}]}%
	{%
		\MeaningsBlock{%
	\ryumeaning%
		{企む}{}{}%
	}%
	}%
	{\tailPrint[]}%
\entry
	{\headPrint[kana={たぐるん\tl{˥}},ipa={taguruɴ},pos={動},rui={1},para={たぐらぬ},acc={H}]}%
	{%
		\MeaningsBlock{%
	\ryumeaning%
		{\ruby[g]{手繰}{たぐ}る}{}{}%
	}%
	}%
	{\tailPrint[]}%
\entry
	{\headPrint[kana={たくん\tl{˥˩}},ipa={tḁkuɴ},pos={動},rui={1},para={たかぬ},acc={F}]}%
	{%
		\MeaningsBlock{%
	\ryumeaning%
		{炊く。炊事する}{}{}%
	}%
	}%
	{\tailPrint[]}%
\entry
	{\headPrint[kana={だぐん\tl{˥˩}},ipa={daguɴ},pos={動},rui={1},para={だがぬ},acc={F}]}%
	{%
		\MeaningsBlock{%
	\ryumeaning%
		{抱く}{}{}%
	}%
	}%
	{\tailPrint[]}%
\entry
	{\headPrint[kana={たげー\tl{˥}},ipa={tageː},pos={副},acc={H}]}%
	{%
		\MeaningsBlock{%
	\ryumeaning%
		{だいぶ。大概}{}{}%
	}%
	}%
	{\tailPrint[]}%
\entry
	{\headPrint[kana={たけーるん\tl{˥}},ipa={tḁkeːruɴ},pos={動},rui={1},para={たけーらぬ},acc={H}]}%
	{%
		\MeaningsBlock{%
	\ryumeaning%
		{叫ぶ。怒鳴る。唸る}{}{}%
	}%
	}%
	{\tailPrint[]}%
\entry
	{\headPrint[kana={たこーすん\tl{˥˩}},ipa={tḁkoːsuɴ},pos={動},rui={2},para={たこはぬ},acc={F}]}%
	{%
		\MeaningsBlock{%
	\ryumeaning%
		{貯える。貯蔵する}{}{}%
	}%
	}%
	{\tailPrint[]}%
\entry
	{\headPrint[kana={たじぃにるん\tl{˥}},ipa={tadʒiniruɴ},pos={動},rui={3},para={たずぬぬ},acc={H}]}%
	{%
		\MeaningsBlock{%
	\ryumeaning%
		{尋ねる。問う。探す}{}{}%
	}%
	}%
	{\tailPrint[]}%
\entry
	{\headPrint[kana={たしかみるん\tl{˥}},ipa={taʃi̥kamiruɴ},pos={動},rui={3},para={たしかむぬ},acc={H}]}%
	{%
		\MeaningsBlock{%
	\ryumeaning%
		{確かめる。確認する}{}{}%
	\ryumeaning%
		{捜し求める}{}{}%
	}%
	}%
	{\tailPrint[]}%
\entry
	{\headPrint[kana={たしきるん\tl{˥}},ipa={taʃi̥kiruɴ},pos={動},rui={3},para={たすくぬ},acc={H}]}%
	{%
		\MeaningsBlock{%
	\ryumeaning%
		{助ける。救う}{}{}%
	}%
	}%
	{\tailPrint[]}%
\entry
	{\headPrint[kana={だすきん\tl{˩˥}},ipa={dasu̥kiɴ},pos={動},rui={3},para={だすくぬ},acc={R}]}%
	{%
		\MeaningsBlock{%
	\ryumeaning%
		{叩きつける}{}{}%
	}%
	}%
	{\tailPrint[]}%
\entry
	{\headPrint[kana={たすたがー\tl{˥˩}},ipa={tasu̥tagaː},pos={句},acc={F}]}%
	{%
		\MeaningsBlock{%
	\ryumeaning%
		{知ったことか}{}{}%
	}%
	}%
	{\tailPrint[]}%
\entry
	{\headPrint[kana={だすたすくん\tl{˩˥}},ipa={dasu̥tasu̥kuɴ},pos={動},rui={1},para={だすたすかぬ},acc={R}]}%
	{%
		\MeaningsBlock{%
	\ryumeaning%
		{放置する。構わない。捨てておく。放っておく}{}{}%
	}%
	}%
	{\tailPrint[]}%
\entry
	{\headPrint[kana={だすますん\tl{˩˥}},ipa={dasu̥masuɴ},pos={動},rui={2},para={だすまはぬ},acc={R}]}%
	{%
		\MeaningsBlock{%
	\ryumeaning%
		{（鞭で）ひっぱたく}{}{}%
	}%
	}%
	{\tailPrint[]}%
\entry
	{\headPrint[kana={たたーるん\tl{˥}},ipa={tḁtaːruɴ},pos={動},rui={1},para={たたーらぬ},acc={H}]}%
	{%
		\MeaningsBlock{%
	\ryumeaning%
		{祟る}{}{}%
	}%
	}%
	{\tailPrint[]}%
\entry
	{\headPrint[kana={たたぎばるん\tl{˥}},ipa={tḁtagibaruɴ},pos={動},rui={1},para={たたぎばらぬ},acc={H}]}%
	{%
		\MeaningsBlock{%
	\ryumeaning%
		{叩き割る}{}{}%
	}%
	}%
	{\tailPrint[]}%
\entry
	{\headPrint[kana={たたぐん\tl{˥}},ipa={tḁtaguɴ},pos={動},rui={1},para={たたがぬ},acc={H}]}%
	{%
		\MeaningsBlock{%
	\ryumeaning%
		{\ruby[g]{叩}{たた}く}{}{}%
	}%
	}%
	{\tailPrint[]}%
\entry
	{\headPrint[kana={たたしきるん\tl{˥}},ipa={tḁtaʃi̥kiruɴ},pos={動},rui={3},para={たたすくぬ},acc={H}]}%
	{%
		\MeaningsBlock{%
	\ryumeaning%
		{投げて叩きつける}{}{}%
	}%
	}%
	{\tailPrint[]}%
\entry
	{\headPrint[kana={ただすん\tl{˥}},ipa={tadasuɴ},pos={動},rui={2},para={ただはぬ},acc={H}]}%
	{%
		\MeaningsBlock{%
	\ryumeaning%
		{質す。確かめる}{質問して確かめる}{}%
	}%
	}%
	{\tailPrint[]}%
\entry
	{\headPrint[kana={たたっくむん\tl{˥}},ipa={tḁtakku̥muɴ},pos={動},rui={1},para={たたっくまぬ},acc={H}]}%
	{%
		\MeaningsBlock{%
	\ryumeaning%
		{放り込む}{}{}%
	}%
	}%
	{\tailPrint[]}%
\entry
	{\headPrint[kana={たたむん},ipa={tḁtamuɴ},pos={動},rui={1},para={たたまぬ}]}%
	{%
		\MeaningsBlock{%
	\ryumeaning%
		{畳む。折り畳む}{}{}%
	}%
	}%
	{\tailPrint[]}%
\entry
	{\headPrint[kana={たためー\tl{˥˩}},ipa={tḁtameː},pos={名},acc={F}]}%
	{%
		\MeaningsBlock{%
	\ryumeaning%
		{\ruby[g]{畳}{たたみ}}{}{}%
	}%
	}%
	{\tailPrint[]}%
\entry
	{\headPrint[kana={ただりん},ipa={tadariɴ},pos={動},rui={3}]}%
	{%
		\MeaningsBlock{%
	\ryumeaning%
		{\ruby[g]{爛}{ただ}れる}{}{}%
	}%
	}%
	{\tailPrint[]}%
\entry
	{\headPrint[kana={たち\tl{˥˩}},ipa={tatʃi},pos={名},acc={F}]}%
	{%
		\MeaningsBlock{%
	\ryumeaning%
		{\ruby[g]{辰}{たつ}}{十二支の辰}{}%
	}%
	}%
	{\tailPrint[]}%
\entry
	{\headPrint[kana={だちご\tl{˥˩}},ipa={datʃigo},pos={名},acc={F}]}%
	{%
		\MeaningsBlock{%
	\ryumeaning%
		{ダンチク}{竹科の植物}{}%
	}%
	}%
	{\tailPrint[]}%
\entry
	{\headPrint[kana={たちのーるん\tl{˥}},ipa={tḁtʃinoːruɴ},pos={動},rui={1},para={たちのーらぬ},acc={H}]}%
	{%
		\MeaningsBlock{%
	\ryumeaning%
		{立ち直る}{}{}%
	}%
	}%
	{\tailPrint[]}%
\entry
	{\headPrint[kana={たちるん\tl{˥}},ipa={tḁtʃiruɴ},pos={動},rui={3},para={たとぬ},acc={H}]}%
	{%
		\MeaningsBlock{%
	\ryumeaning%
		{立てる。建てる}{}{}%
	}%
	}%
	{\tailPrint[]}%
\entry
	{\headPrint[kana={だっきょん\tl{˩˥}},ipa={dakkjoɴ},pos={名},acc={R}]}%
	{%
		\MeaningsBlock{%
	\ryumeaning%
		{ラッキョウ}{野菜}{}%
	}%
	}%
	{\tailPrint[]}%
\entry
	{\headPrint[kana={だっくゎーるん\tl{˥}},ipa={dakkwaːruɴ},pos={動},rui={3},acc={H}]}%
	{%
		\MeaningsBlock{%
	\ryumeaning%
		{くっ付く。ひっ付く}{}{}%
	}%
	}%
	{\tailPrint[]}%
\entry
	{\headPrint[kana={だっくゎすん},ipa={dakkwasuɴ},pos={動},rui={2},para={だっくゎーさぬ}]}%
	{%
		\MeaningsBlock{%
	\ryumeaning%
		{引っつける。塗りつける}{}{}%
	}%
	}%
	{\tailPrint[]}%
\entry
	{\headPrint[kana={だっさが},ipa={dassaga},pos={-}]}%
	{%
		\MeaningsBlock{%
	\ryumeaning%
		{ぶち当たる。ぶち当てる}{}{}%
	}%
	}%
	{\tailPrint[]}%
\entry
	{\headPrint[kana={だった},ipa={datta},pos={擬}]}%
	{%
		\MeaningsBlock{%
	\ryumeaning%
		{急に凪ぐ様。急に小さくなる様}{風が凪ぐとき勢いが急に弱くなる。\emb{だった\ws{}とりるん}「風が急に凪いだ」のように使う}{}%
	}%
	}%
	{\tailPrint[]}%
\entry
	{\headPrint[kana={だっち},ipa={dattʃi},pos={-}]}%
	{%
		\MeaningsBlock{%
	\ryumeaning%
		{突っ立つ。立つ}{}{}%
	}%
	}%
	{\tailPrint[]}%
\entry
	{\headPrint[kana={だっつぁぐん\tl{˥˩}},ipa={dattsaguɴ},pos={動},rui={1},para={だっつぁがぬ},acc={F}]}%
	{%
		\MeaningsBlock{%
	\ryumeaning%
		{はかどる。能率が上がる}{}{}%
	}%
	}%
	{\tailPrint[]}%
\entry
	{\headPrint[kana={だっふぁだっふぁ},ipa={daffadaffa},pos={擬}]}%
	{%
		\MeaningsBlock{%
	\ryumeaning%
		{堂々と歩くさま}{}{}%
	}%
	}%
	{\tailPrint[]}%
\entry
	{\headPrint[kana={たつん\tl{˥}},ipa={tḁtsuɴ},pos={動},rui={2},para={たたぬ},acc={H}]}%
	{%
		\MeaningsBlock{%
	\ryumeaning%
		{立つ。建つ}{}{}%
	}%
	}%
	{\tailPrint[]}%
\entry
	{\headPrint[kana={たつん\tl{˥}},ipa={tḁtsuɴ},pos={動},rui={2},para={たたぬ},acc={H}]}%
	{%
		\MeaningsBlock{%
	\ryumeaning%
		{経つ。経過する}{}{}%
	}%
	}%
	{\tailPrint[]}%
\entry
	{\headPrint[kana={たでーま},ipa={tadeːma},pos={副}]}%
	{%
		\MeaningsBlock{%
	\ryumeaning%
		{すぐに。瞬時に。即座に}{}{}%
	}%
	}%
	{\tailPrint[]}%
\entry
	{\headPrint[kana={たとぅいん\tl{˥}},ipa={tḁtuiɴ},pos={動},acc={H}]}%
	{%
		\MeaningsBlock{%
	\ryumeaning%
		{例える}{}{}%
	}%
	}%
	{\tailPrint[]}%
\entry
	{\headPrint[kana={たどぅるん\tl{˥}},ipa={taduruɴ},pos={動},rui={1},para={たどぅらぬ},acc={H}]}%
	{%
		\MeaningsBlock{%
	\ryumeaning%
		{\ruby[g]{辿}{たど}る}{}{}%
	}%
	}%
	{\tailPrint[]}%
\entry
	{\headPrint[kana={たな\tl{˥˩}},ipa={tḁna},pos={名},acc={F}]}%
	{%
		\MeaningsBlock{%
	\ryumeaning%
		{\ruby[g]{棚}{たな}}{}{}%
	}%
	}%
	{\tailPrint[]}%
\entry
	{\headPrint[kana={たなー\tl{˥}},ipa={tḁnaː},pos={名},acc={H}]}%
	{%
		\MeaningsBlock{%
	\ryumeaning%
		{田んぼ。水田}{}{}%
	}%
	}%
	{\tailPrint[]}%
\entry
	{\headPrint[kana={たな\tl{˥}いびるん\tl{˥˩}},ipa={tḁna.ibiruɴ},pos={句},acc={H F}]}%
	{%
		\MeaningsBlock{%
	\ryumeaning%
		{田植えをする}{}{}%
	}%
	}%
	{\tailPrint[]}%
\entry
	{\headPrint[kana={たなが\tl{˥}},ipa={tḁnaga},pos={名},acc={H}]}%
	{%
		\MeaningsBlock{%
	\ryumeaning%
		{田んぼ。水田}{}{}%
	}%
	}%
	{\tailPrint[]}%
\entry
	{\headPrint[kana={たなぎるん},ipa={tḁnagiruɴ},pos={動},rui={3},para={たなぐぬ}]}%
	{%
		\MeaningsBlock{%
	\ryumeaning%
		{頼りにする}{}{}%
	}%
	}%
	{\tailPrint[]}%
\entry
	{\headPrint[kana={たなぐん},ipa={tḁnaguɴ},pos={動}]}%
	{%
		\MeaningsBlock{%
	\ryumeaning%
		{頼る}{}{}%
	}%
	}%
	{\tailPrint[]}%
\entry
	{\headPrint[kana={たに\tl{˥}},ipa={tḁni},pos={名},acc={H}]}%
	{%
		\MeaningsBlock{%
	\ryumeaning%
		{種。種子}{}{}%
	}%
	}%
	{\tailPrint[]}%
\entry
	{\headPrint[kana={たにかたうやぐ\tl{˥}},ipa={tḁnikḁta.ujagu},pos={名},acc={H}]}%
	{%
		\MeaningsBlock{%
	\ryumeaning%
		{父方の親戚}{}{}%
	}%
	}%
	{\tailPrint[]}%
\entry
	{\headPrint[kana={たに\ws{}しさん},ipa={tḁni ʃi̥saɴ},pos={句}]}%
	{%
		\MeaningsBlock{%
	\ryumeaning%
		{種がなくなる。絶滅する。（作物が）根絶やしになる}{}{}%
	}%
	}%
	{\tailPrint[]}%
\entry
	{\headPrint[kana={たにどぅり\tl{˥˩}},ipa={tḁniduri},pos={名},acc={F}]}%
	{%
		\MeaningsBlock{%
	\ryumeaning%
		{種取祭}{}{}%
	}%
	}%
	{\tailPrint[note={\emb{たに}「種」が平進型であるため、\emb{たにどぅり}も平進型に所属すると期待されるが、下降型である。八重山の他の方言を見ると、石垣四箇や宮良方言も「種取」が下降型になっている}]}%
\entry
	{\headPrint[kana={たに\tl{˥}\ws{}とぅるん\tl{˥}},ipa={tḁni tu̥ruɴ},pos={句},para={たにとらぬ},acc={H H}]}%
	{%
		\MeaningsBlock{%
	\ryumeaning%
		{採種する。去勢する}{}{}%
	}%
	}%
	{\tailPrint[]}%
\entry
	{\headPrint[kana={たぬみ\tl{˥}},ipa={tḁnumi},pos={名},acc={H}]}%
	{%
		\MeaningsBlock{%
	\ryumeaning%
		{頼み。頼り}{}{}%
	}%
	}%
	{\tailPrint[]}%
\entry
	{\headPrint[kana={たぬむん\tl{˥}},ipa={tḁnumuɴ},pos={動},rui={1},para={たぬまぬ},acc={H}]}%
	{%
		\MeaningsBlock{%
	\ryumeaning%
		{頼む}{}{}%
	}%
	}%
	{\tailPrint[]}%
\entry
	{\headPrint[kana={たぱ\tl{˥}},ipa={tḁpa},pos={名},acc={H}]}%
	{%
		\MeaningsBlock{%
	\ryumeaning%
		{\ruby[g]{束}{たば}}{}{}%
	}%
	}%
	{\tailPrint[]}%
\entry
	{\headPrint[kana={たぴ\tl{˥˩}},ipa={tḁpi},pos={名},acc={F}]}%
	{%
		\MeaningsBlock{%
	\ryumeaning%
		{旅。旅行}{}{}%
	}%
	}%
	{\tailPrint[]}%
\entry
	{\headPrint[kana={たぶ\tl{˥}},ipa={tabu},pos={名},acc={H}]}%
	{%
		\MeaningsBlock{%
	\ryumeaning%
		{たも網}{}{}%
	}%
	}%
	{\tailPrint[]}%
\entry
	{\headPrint[kana={だふつり\tl{˩˥}わっきん\tl{˥˩}},ipa={dafu̥tsuriwakkiɴ},pos={動},rui={3},para={だふつりわっくぬ},acc={R F}]}%
	{%
		\MeaningsBlock{%
	\ryumeaning%
		{追い散らす}{}{}%
	}%
	}%
	{\tailPrint[]}%
\entry
	{\headPrint[kana={だふつるん\tl{˩˥}},ipa={dafu̥tsuruɴ},pos={動},rui={1},para={だふつらぬ},acc={R}]}%
	{%
		\MeaningsBlock{%
	\ryumeaning%
		{追い立てる}{}{}%
	}%
	}%
	{\tailPrint[]}%
\entry
	{\headPrint[kana={たぼーらりん\tl{˥}},ipa={taboːrariɴ},pos={動},para={たぼーらるぬ},acc={H}]}%
	{%
		\MeaningsBlock{%
	\ryumeaning%
		{賜る。頂戴する}{}{}%
	}%
	}%
	{\tailPrint[]}%
\entry
	{\headPrint[kana={たぼらいるん\tl{˥}},ipa={taborairuɴ},pos={動},para={たぼらるぬ},acc={H}]}%
	{%
		\MeaningsBlock{%
	\ryumeaning%
		{賜る。頂く}{}{}%
	\ryumeaning%
		{授ける。下さる}{}{}%
	}%
	}%
	{\tailPrint[]}%
\entry
	{\headPrint[kana={\josho{}たぼらり\kako{}むぬ},ipa={taborarimunu},pos={名},acc={H+L}]}%
	{%
		\MeaningsBlock{%
	\ryumeaning%
		{頂き物。貰い物}{}{}%
	}%
	}%
	{\tailPrint[]}%
\entry
	{\headPrint[kana={たぼるん\tl{˥}},ipa={taboruɴ},pos={動},rui={1},acc={H}]}%
	{%
		\MeaningsBlock{%
	\ryumeaning%
		{下さる}{}{}%
	\ryumeaning%
		{頂く}{}{}%
	}%
	}%
	{\tailPrint[]}%
\entry
	{\headPrint[kana={だぼんがすん\tl{˩˥}},ipa={daboŋgasuɴ},pos={動},rui={2},para={だぼんがはぬ},acc={R}]}%
	{%
		\MeaningsBlock{%
	\ryumeaning%
		{（液体の中に物を）強く投げ入れる}{}{}%
	}%
	}%
	{\tailPrint[]}%
\entry
	{\headPrint[kana={だぼんた},ipa={dabonta},pos={擬}]}%
	{%
		\MeaningsBlock{%
	\ryumeaning%
		{どぼんと。ざぶんと}{}{}%
	}%
	}%
	{\tailPrint[]}%
\entry
	{\headPrint[kana={たま},ipa={tama},pos={接尾}]}%
	{%
		\MeaningsBlock{%
	\ryumeaning%
		{指小辞}{小さいものの愛称。\emb{うたま}「子供」、\emb{ゆんたま}「小魚」、\emb{いしんたま}「小石」など}{}%
	}%
	}%
	{\tailPrint[]}%
\entry
	{\headPrint[kana={た\josho{}まー\kako{}たま},ipa={tḁmaːtḁma},pos={副},acc={重複}]}%
	{%
		\MeaningsBlock{%
	\ryumeaning%
		{たまに。たまたま}{}{}%
	}%
	}%
	{\tailPrint[]}%
\entry
	{\headPrint[kana={たましぃ\tl{˥}},ipa={tḁmaʃi},pos={名},acc={H}]}%
	{%
		\MeaningsBlock{%
	\ryumeaning%
		{魂。霊魂}{}{}%
	}%
	}%
	{\tailPrint[]}%
\entry
	{\headPrint[kana={たまち\tl{˥}},ipa={tḁmatʃi},pos={名},acc={H}]}%
	{%
		\MeaningsBlock{%
	\ryumeaning%
		{魂。霊魂}{}{}%
	}%
	}%
	{\tailPrint[]}%
\entry
	{\headPrint[kana={たまち\tl{˥}},ipa={tḁmatʃi},pos={名},acc={H}]}%
	{%
		\MeaningsBlock{%
	\ryumeaning%
		{割当て。持分}{\emb{くまた}と同義}{}%
	}%
	}%
	{\tailPrint[]}%
\entry
	{\headPrint[kana={たまるん\tl{˥}},ipa={tḁmaruɴ},pos={動},rui={1},para={たまらぬ},acc={H}]}%
	{%
		\MeaningsBlock{%
	\ryumeaning%
		{（曲がっている物が）真っ直ぐになる}{日本語の「矯める」と同源の\emb{たみるん}「まっすぐにする」に対応する自動詞}{}%
	}%
	}%
	{\tailPrint[]}%
\entry
	{\headPrint[kana={たまるん\tl{˥˩}},ipa={tḁmaruɴ},pos={動},rui={1},para={たまらぬ},acc={F}]}%
	{%
		\MeaningsBlock{%
	\ryumeaning%
		{溜まる。貯まる}{}{}%
	}%
	}%
	{\tailPrint[]}%
\entry
	{\headPrint[kana={たまん\tl{˥}},ipa={tḁmaɴ},pos={名},acc={H}]}%
	{%
		\MeaningsBlock{%
	\ryumeaning%
		{玉。弾}{}{}%
	}%
	}%
	{\tailPrint[]}%
\entry
	{\headPrint[kana={たみ\tl{˥}},ipa={tḁmi},pos={名},acc={H}]}%
	{%
		\MeaningsBlock{%
	\ryumeaning%
		{\ruby[g]{為}{ため}}{}{}%
	}%
	}%
	{\tailPrint[]}%
\entry
	{\headPrint[kana={たみすん\tl{˥}},ipa={tḁmisuɴ},pos={動},rui={2b},para={たみさぬ},acc={H}]}%
	{%
		\MeaningsBlock{%
	\ryumeaning%
		{試す}{}{}%
	}%
	}%
	{\tailPrint[]}%
\entry
	{\headPrint[kana={たみらすん},ipa={tḁmirasuɴ},pos={動},rui={2},para={たみらはぬ}]}%
	{%
		\MeaningsBlock{%
	\ryumeaning%
		{真っ直ぐにする}{}{}%
	}%
	}%
	{\tailPrint[]}%
\entry
	{\headPrint[kana={たみるん\tl{˥}},ipa={tḁmiruɴ},pos={動},acc={H}]}%
	{%
		\MeaningsBlock{%
	\ryumeaning%
		{真っ直ぐにする}{}{}%
	}%
	}%
	{\tailPrint[]}%
\entry
	{\headPrint[kana={たみるん\tl{˥˩}},ipa={tḁmiruɴ},pos={動},rui={1},acc={F}]}%
	{%
		\MeaningsBlock{%
	\ryumeaning%
		{貯める。蓄える}{}{}%
	\ryumeaning%
		{溜める}{}{}%
	}%
	}%
	{\tailPrint[]}%
\entry
	{\headPrint[kana={たむき\tl{˥}},ipa={tamuki},pos={名},acc={H}]}%
	{%
		\MeaningsBlock{%
	\ryumeaning%
		{\ruby[g]{風}{かざ}\ruby[g]{向}{むき}}{}{}%
	}%
	}%
	{\tailPrint[]}%
\entry
	{\headPrint[kana={たむつん\tl{˥}},ipa={tamutsuɴ},pos={動},para={たむたぬ},acc={H}]}%
	{%
		\MeaningsBlock{%
	\ryumeaning%
		{保つ。長持ちする}{}{}%
	}%
	}%
	{\tailPrint[]}%
\entry
	{\headPrint[kana={たむぬ\tl{˥}},ipa={tamunu},pos={名},acc={H}]}%
	{%
		\MeaningsBlock{%
	\ryumeaning%
		{\ruby[g]{薪}{たきぎ}}{}{}%
	}%
	}%
	{\tailPrint[]}%
\entry
	{\headPrint[kana={たむん},ipa={tamun},pos={名}]}%
	{%
		\MeaningsBlock{%
	\ryumeaning%
		{\ruby[g]{薪}{たきぎ}}{}{}%
	}%
	}%
	{\tailPrint[]}%
\entry
	{\headPrint[kana={たゆるん\tl{˥}},ipa={tajuruɴ},pos={動},rui={1},para={たゆらぬ},acc={H}]}%
	{%
		\MeaningsBlock{%
	\ryumeaning%
		{頼る}{}{}%
	}%
	}%
	{\tailPrint[]}%
\entry
	{\headPrint[kana={たらーすん\tl{˥˩}},ipa={tḁraːsuɴ},pos={動},rui={2},para={たらーはぬ},acc={F}]}%
	{%
		\MeaningsBlock{%
	\ryumeaning%
		{物を補う。埋め合わせる}{}{}%
	}%
	}%
	{\tailPrint[]}%
\entry
	{\headPrint[kana={たらーぬ\tl{˥˩}},ipa={tḁraːnu},pos={句},acc={F}]}%
	{%
		\MeaningsBlock{%
	\ryumeaning%
		{足りない。不足する}{}{}%
	}%
	}%
	{\tailPrint[]}%
\entry
	{\headPrint[kana={たらぐ\tl{˥}},ipa={taragu},pos={名},acc={H}]}%
	{%
		\MeaningsBlock{%
	\ryumeaning%
		{\ruby[g]{俵}{たわら}}{}{}%
	}%
	}%
	{\tailPrint[]}%
\entry
	{\headPrint[kana={たらぐ\tl{˥}},ipa={taragu},pos={名},acc={H}]}%
	{%
		\MeaningsBlock{%
	\ryumeaning%
		{オタマジャクシ}{}{}%
	}%
	}%
	{\tailPrint[]}%
\entry
	{\headPrint[kana={たらすん\tl{˥˩}},ipa={tḁrasuɴ},pos={動},rui={2},para={たらはぬ},acc={F}]}%
	{%
		\MeaningsBlock{%
	\ryumeaning%
		{足らす。補う}{}{}%
	}%
	}%
	{\tailPrint[]}%
\entry
	{\headPrint[kana={たらすん\tl{˥˩}},ipa={tḁrasuɴ},pos={動},rui={2},para={たらはぬ},acc={F}]}%
	{%
		\MeaningsBlock{%
	\ryumeaning%
		{溶かす。とろかす}{}{}%
	}%
	}%
	{\tailPrint[]}%
\entry
	{\headPrint[kana={だらだら},ipa={daradara},pos={擬}]}%
	{%
		\MeaningsBlock{%
	\ryumeaning%
		{だらだら}{動作の鈍いさま}{}%
	}%
	}%
	{\tailPrint[]}%
\entry
	{\headPrint[kana={たりるん\tl{˥}},ipa={tḁriruɴ},pos={動},rui={3},acc={H}]}%
	{%
		\MeaningsBlock{%
	\ryumeaning%
		{溶ける。垂れる}{}{}%
	}%
	}%
	{\tailPrint[]}%
\entry
	{\headPrint[kana={たりるん\tl{˥˩}},ipa={tariruɴ},pos={動},acc={F}]}%
	{%
		\MeaningsBlock{%
	\ryumeaning%
		{足りる}{十分に足りる}{}%
	}%
	}%
	{\tailPrint[]}%
\entry
	{\headPrint[kana={だりるん\tl{˩˥}},ipa={dariruɴ},pos={動},rui={3},para={だるぬ},acc={R}]}%
	{%
		\MeaningsBlock{%
	\ryumeaning%
		{\ruby[g]{萎}{しお}れる}{}{}%
	\ryumeaning%
		{元気がなくなる。悄げかえる}{}{}%
	}%
	}%
	{\tailPrint[]}%
\entry
	{\headPrint[kana={たるざー\tl{˥}},ipa={tarudzaː},pos={名},acc={H}]}%
	{%
		\MeaningsBlock{%
	\ryumeaning%
		{どいつが。誰が}{}{}%
	}%
	}%
	{\tailPrint[]}%
\entry
	{\headPrint[kana={だる\tl{˩˥}さーん},ipa={darusaːɴ},pos={形},para={だるさーへぬ},acc={R}]}%
	{%
		\MeaningsBlock{%
	\ryumeaning%
		{だるい。疲労ぎみだ}{}{}%
	}%
	}%
	{\tailPrint[]}%
\entry
	{\headPrint[kana={たれー\tl{˥}},ipa={tḁreː},pos={名},acc={H}]}%
	{%
		\MeaningsBlock{%
	\ryumeaning%
		{\ruby[g]{盥}{たらい}}{}{}%
	}%
	}%
	{\tailPrint[]}%
\entry
	{\headPrint[kana={たん\tl{˥˩}},ipa={taɴ},pos={名},acc={F}]}%
	{%
		\MeaningsBlock{%
	\ryumeaning%
		{炭。木炭}{}{}%
	}%
	}%
	{\tailPrint[]}%
\entry
	{\headPrint[kana={たん\tl{˥˩}},ipa={taɴ},pos={名},acc={F}]}%
	{%
		\MeaningsBlock{%
	\ryumeaning%
		{反}{田畑の面積や反物の数}{}%
	}%
	}%
	{\tailPrint[]}%
\entry
	{\headPrint[kana={だん\tl{˩˥}},ipa={daɴ},pos={名},acc={R}]}%
	{%
		\MeaningsBlock{%
	\ryumeaning%
		{段。壇}{}{}%
	}%
	}%
	{\tailPrint[]}%
\entry
	{\headPrint[kana={たんか\tl{˥˩}},ipa={taŋka},pos={名},acc={F}]}%
	{%
		\MeaningsBlock{%
	\ryumeaning%
		{真向かい。正面}{}{}%
	}%
	}%
	{\tailPrint[]}%
\entry
	{\headPrint[kana={たんかー\tl{˥}},ipa={taŋkaː},pos={名},acc={H}]}%
	{%
		\MeaningsBlock{%
	\ryumeaning%
		{満一歳の誕生日}{}{}%
	}%
	}%
	{\tailPrint[]}%
\entry
	{\headPrint[kana={たんがーむぬ\tl{˥}},ipa={taŋgaːmunu},pos={名},acc={H}]}%
	{%
		\MeaningsBlock{%
	\ryumeaning%
		{独り者。独身}{}{}%
	}%
	}%
	{\tailPrint[]}%
\entry
	{\headPrint[kana={だんがさ\tl{˥}},ipa={daŋgasa},pos={名},acc={H}]}%
	{%
		\MeaningsBlock{%
	\ryumeaning%
		{洋傘}{}{}%
	}%
	}%
	{\tailPrint[note={蘭傘の転訛}]}%
\entry
	{\headPrint[kana={たんかにげー\tl{˩˥}},ipa={taŋkanigeː},pos={名},acc={R}]}%
	{%
		\MeaningsBlock{%
	\ryumeaning%
		{遥拝}{遠隔地からの礼拝}{}%
	}%
	}%
	{\tailPrint[]}%
\entry
	{\headPrint[kana={たんきむぬ},ipa={taŋkimunu},pos={名}]}%
	{%
		\MeaningsBlock{%
	\ryumeaning%
		{短気者}{}{}%
	}%
	}%
	{\tailPrint[]}%
\entry
	{\headPrint[kana={たんきるん\tl{˥}},ipa={taŋkiruɴ},pos={動},rui={3},para={たんくぬ},acc={H}]}%
	{%
		\MeaningsBlock{%
	\ryumeaning%
		{ひるむ。怖気づく}{}{}%
	}%
	}%
	{\tailPrint[]}%
\entry
	{\headPrint[kana={たんぐ\tl{˥}},ipa={taŋgu},pos={名},acc={H}]}%
	{%
		\MeaningsBlock{%
	\ryumeaning%
		{桶。水桶}{}{}%
	}%
	}%
	{\tailPrint[]}%
\entry
	{\headPrint[kana={だんざー\tl{˥}},ipa={dandzaː},pos={名},acc={H}]}%
	{%
		\MeaningsBlock{%
	\ryumeaning%
		{お前は。きさまは}{}{}%
	}%
	}%
	{\tailPrint[]}%
\entry
	{\headPrint[kana={だ\josho{}ん\kako{}た},ipa={danta},pos={副},acc={LHL}]}%
	{%
		\MeaningsBlock{%
	\ryumeaning%
		{しっかり。がっちり}{}{}%
	}%
	}%
	{\tailPrint[]}%
\entry
	{\headPrint[kana={\josho{}だん\kako{}だん},ipa={dandan},pos={名},acc={H+L}]}%
	{%
		\MeaningsBlock{%
	\ryumeaning%
		{階段}{}{}%
	}%
	}%
	{\tailPrint[]}%
\entry
	{\headPrint[kana={たんでー\tl{˥}},ipa={tandeː},pos={感},acc={H}]}%
	{%
		\MeaningsBlock{%
	\ryumeaning%
		{どうか。是非とも}{}{}%
	}%
	}%
	{\tailPrint[]}%
\entry
	{\headPrint[kana={だんぱち\tl{˥}},ipa={dampḁtʃi},pos={名},acc={H}]}%
	{%
		\MeaningsBlock{%
	\ryumeaning%
		{断髪。理髪}{}{}%
	}%
	}%
	{\tailPrint[]}%
\entry
	{\headPrint[kana={だんぱちーやー},ipa={dampatʃiːjaː},pos={名}]}%
	{%
		\MeaningsBlock{%
	\ryumeaning%
		{床屋。理髪店}{}{}%
	}%
	}%
	{\tailPrint[]}%
\entry
	{\headPrint[kana={だんぱん\tl{˥˩}},ipa={dampaɴ},pos={名},acc={F}]}%
	{%
		\MeaningsBlock{%
	\ryumeaning%
		{談判。抗議}{}{}%
	}%
	}%
	{\tailPrint[]}%
\LetterHead{ち}
\entry
	{\headPrint[kana={ちー},ipa={tʃiː},pos={接尾}]}%
	{%
		\MeaningsBlock{%
	\ryumeaning%
		{～個}{\emb{ぴとちぃ}「一個)」、\emb{はたちぃ}「二十歳」}{}%
	\ryumeaning%
		{～歳}{\emb{ぴとちぃ}「一個)」、\emb{はたちぃ}「二十歳」}{}%
	}%
	}%
	{\tailPrint[]}%
\entry
	{\headPrint[kana={ちーかたうやぐ\tl{˥}},ipa={tʃiːkḁta.ujagu},pos={名},acc={H}]}%
	{%
		\MeaningsBlock{%
	\ryumeaning%
		{母方親戚}{}{}%
	}%
	}%
	{\tailPrint[]}%
\entry
	{\headPrint[kana={ちき\tl{˥}},ipa={tʃi̥ki},pos={名},acc={H}]}%
	{%
		\MeaningsBlock{%
	\ryumeaning%
		{（月日の）月}{}{}%
	}%
	}%
	{\tailPrint[]}%
\entry
	{\headPrint[kana={ちくどぅん\tl{˥}},ipa={tʃi̥kuduɴ},pos={名},acc={H}]}%
	{%
		\MeaningsBlock{%
	\ryumeaning%
		{筑登之}{王府時代の位階}{}%
	}%
	}%
	{\tailPrint[]}%
\entry
	{\headPrint[kana={ちごーん\tl{˥˩}},ipa={tʃigoːɴ},pos={動},para={ちがわぬ},acc={F}]}%
	{%
		\MeaningsBlock{%
	\ryumeaning%
		{違う}{}{}%
	\ryumeaning%
		{変わる}{}{}%
	}%
	}%
	{\tailPrint[]}%
\entry
	{\headPrint[kana={ちじ\tl{˥}},ipa={tʃidʒi},pos={名},acc={H}]}%
	{%
		\MeaningsBlock{%
	\ryumeaning%
		{頂。頂上}{}{}%
	}%
	}%
	{\tailPrint[]}%
\entry
	{\headPrint[kana={ちじきるん\tl{˥˩}},ipa={tʃidʒikiruɴ},pos={動},rui={3},para={ちじくぬ},acc={F}]}%
	{%
		\MeaningsBlock{%
	\ryumeaning%
		{続ける}{}{}%
	}%
	}%
	{\tailPrint[]}%
\entry
	{\headPrint[kana={ちじくん\tl{˥˩}},ipa={tʃidʒikuɴ},pos={動},rui={1},para={ちじかぬ},acc={F}]}%
	{%
		\MeaningsBlock{%
	\ryumeaning%
		{続く}{}{}%
	}%
	}%
	{\tailPrint[]}%
\entry
	{\headPrint[kana={ちちすむん\tl{˥˩}},ipa={tʃitʃisu̥muɴ},pos={動},rui={1},para={ちちしぃまぬ},acc={F}]}%
	{%
		\MeaningsBlock{%
	\ryumeaning%
		{慎む}{}{}%
	}%
	}%
	{\tailPrint[]}%
\entry
	{\headPrint[kana={ちぶ\tl{˥˩}},ipa={tʃibu},pos={名},acc={F}]}%
	{%
		\MeaningsBlock{%
	\ryumeaning%
		{坪}{面積の単位}{}%
	}%
	}%
	{\tailPrint[]}%
\entry
	{\headPrint[kana={\josho{}ちゃー\kako{}が},ipa={tʃaːga},pos={句},acc={HL}]}%
	{%
		\MeaningsBlock{%
	\ryumeaning%
		{どうか。どうだい}{}{}%
	}%
	}%
	{\tailPrint[]}%
\entry
	{\headPrint[kana={ちゃんぷるー\tl{˥}},ipa={tʃampuruː},pos={名},acc={H}]}%
	{%
		\MeaningsBlock{%
	\ryumeaning%
		{}{油でいためたおかず類}{}%
	}%
	}%
	{\tailPrint[]}%
\entry
	{\headPrint[kana={ちゅー},ipa={tʃuː},pos={助}]}%
	{%
		\MeaningsBlock{%
	\ryumeaning%
		{～だそうだ}{}{}%
	}%
	}%
	{\tailPrint[]}%
\entry
	{\headPrint[kana={ちょーみー\tl{˥˩}},ipa={tʃoːmiː},pos={名},acc={F}]}%
	{%
		\MeaningsBlock{%
	\ryumeaning%
		{長命。長寿}{}{}%
	}%
	}%
	{\tailPrint[]}%
\entry
	{\headPrint[kana={ちんだみ},ipa={tʃindami},pos={名}]}%
	{%
		\MeaningsBlock{%
	\ryumeaning%
		{調弦}{三線用語}{}%
	}%
	}%
	{\tailPrint[]}%
\LetterHead{つ}
\entry
	{\headPrint[kana={つぐん\tl{˥˩}},ipa={tsuguɴ},pos={動},rui={1},para={つがぬ},acc={F}]}%
	{%
		\MeaningsBlock{%
	\ryumeaning%
		{継ぐ}{}{}%
	}%
	}%
	{\tailPrint[]}%
\LetterHead{て}
\entry
	{\headPrint[kana={てー\tl{˥}},ipa={teː},pos={名},acc={H}]}%
	{%
		\MeaningsBlock{%
	\ryumeaning%
		{魚の干物}{}{}%
	}%
	}%
	{\tailPrint[]}%
\entry
	{\headPrint[kana={てー\tl{˥}},ipa={teː},pos={名},acc={H}]}%
	{%
		\MeaningsBlock{%
	\ryumeaning%
		{支える力}{}{}%
	}%
	}%
	{\tailPrint[]}%
\entry
	{\headPrint[kana={てー\tl{˥}},ipa={teː},pos={名},acc={H}]}%
	{%
		\MeaningsBlock{%
	\ryumeaning%
		{\ruby[g]{松明}{たいまつ}}{}{}%
	}%
	}%
	{\tailPrint[]}%
\entry
	{\headPrint[kana={てーく\tl{˥}},ipa={teːku},pos={名},acc={H}]}%
	{%
		\MeaningsBlock{%
	\ryumeaning%
		{太鼓}{}{}%
	}%
	}%
	{\tailPrint[]}%
\entry
	{\headPrint[kana={でーぐ\tl{˩˥}},ipa={deːgu},pos={名},acc={R}]}%
	{%
		\MeaningsBlock{%
	\ryumeaning%
		{大工}{}{}%
	}%
	}%
	{\tailPrint[]}%
\entry
	{\headPrint[kana={でーぐに\tl{˩˥}},ipa={deːguni},pos={名},acc={R}]}%
	{%
		\MeaningsBlock{%
	\ryumeaning%
		{大根}{}{}%
	}%
	}%
	{\tailPrint[]}%
\entry
	{\headPrint[kana={てーし\tl{˥}},ipa={teːʃi},pos={名},acc={H}]}%
	{%
		\MeaningsBlock{%
	\ryumeaning%
		{野いちご。ナワシロイチゴ}{}{}%
	}%
	}%
	{\tailPrint[note={多良間方言の\emb{たぎす}に対応}]}%
\entry
	{\headPrint[kana={でーじ\tl{˩˥}},ipa={deːdʒi},pos={名},acc={R}]}%
	{%
		\MeaningsBlock{%
	\ryumeaning%
		{大事。大変}{}{}%
	}%
	}%
	{\tailPrint[]}%
\entry
	{\headPrint[kana={でーだが\tl{˥˩}},ipa={deːdaga},pos={名},acc={F}]}%
	{%
		\MeaningsBlock{%
	\ryumeaning%
		{高価な}{}{}%
	}%
	}%
	{\tailPrint[]}%
\entry
	{\headPrint[kana={てーに\tl{˥}},ipa={teːni},pos={名},acc={H}]}%
	{%
		\MeaningsBlock{%
	\ryumeaning%
		{\ruby[g]{舵}{かじ}\ruby[g]{棒}{ぼう}}{}{}%
	}%
	}%
	{\tailPrint[]}%
\entry
	{\headPrint[kana={てぃがら\tl{˥}},ipa={tigara},pos={名},acc={H}]}%
	{%
		\MeaningsBlock{%
	\ryumeaning%
		{手柄。獲物}{}{}%
	}%
	}%
	{\tailPrint[]}%
\entry
	{\headPrint[kana={てぃんま\tl{˥˩}},ipa={timma},pos={名},acc={F}]}%
	{%
		\MeaningsBlock{%
	\ryumeaning%
		{伝馬船}{}{}%
	}%
	}%
	{\tailPrint[]}%
\entry
	{\headPrint[kana={てしきるん\tl{˥˩}},ipa={teʃi̥kiruɴ},pos={動},rui={3},para={てすくぬ},acc={F}]}%
	{%
		\MeaningsBlock{%
	\ryumeaning%
		{（火を）熾す。焚き付ける。（薪に火を）点ける。（薪などを）くべる}{}{}%
	}%
	}%
	{\tailPrint[]}%
\entry
	{\headPrint[kana={てすかるん\tl{˥˩}},ipa={tesu̥karuɴ},pos={動},rui={1},para={てすからぬ},acc={F}]}%
	{%
		\MeaningsBlock{%
	\ryumeaning%
		{（火が）燃えだす}{}{}%
	}%
	}%
	{\tailPrint[]}%
\entry
	{\headPrint[kana={でん\tl{˥˩}},ipa={deɴ},pos={名},acc={F}]}%
	{%
		\MeaningsBlock{%
	\ryumeaning%
		{値段。代金}{}{}%
	}%
	}%
	{\tailPrint[]}%
\LetterHead{と}
\entry
	{\headPrint[kana={とー},ipa={toː},pos={感}]}%
	{%
		\MeaningsBlock{%
	\ryumeaning%
		{もう}{}{}%
	}%
	}%
	{\tailPrint[]}%
\entry
	{\headPrint[kana={とー\tl{˥}},ipa={toː},pos={名},acc={H}]}%
	{%
		\MeaningsBlock{%
	\ryumeaning%
		{唐。中国}{}{}%
	}%
	}%
	{\tailPrint[]}%
\entry
	{\headPrint[kana={とー\tl{˥˩}},ipa={toː},pos={名},acc={F}]}%
	{%
		\MeaningsBlock{%
	\ryumeaning%
		{沖。大海}{}{}%
	}%
	}%
	{\tailPrint[]}%
\entry
	{\headPrint[kana={とー\tl{˥˩}},ipa={toː},pos={名},acc={F}]}%
	{%
		\MeaningsBlock{%
	\ryumeaning%
		{\ruby[g]{十}{とお}}{}{}%
	}%
	}%
	{\tailPrint[note={簡略化\emb{とぅ}}]}%
\entry
	{\headPrint[kana={とー\tl{˥˩}},ipa={toː},pos={名},acc={F}]}%
	{%
		\MeaningsBlock{%
	\ryumeaning%
		{低平地。窪地}{}{}%
	}%
	}%
	{\tailPrint[]}%
\entry
	{\headPrint[kana={とーかち\tl{˥}},ipa={toːkḁtʃi},pos={名},acc={H}]}%
	{%
		\MeaningsBlock{%
	\ryumeaning%
		{米寿}{八十八歳の祝い}{}%
	}%
	}%
	{\tailPrint[]}%
\entry
	{\headPrint[kana={とーさ\tl{˥}},ipa={toːsa},pos={名},acc={H}]}%
	{%
		\MeaningsBlock{%
	\ryumeaning%
		{田草}{}{}%
	}%
	}%
	{\tailPrint[]}%
\entry
	{\headPrint[kana={とーしんべー\tl{˥}},ipa={toːʃimbeː},pos={名},acc={H}]}%
	{%
		\MeaningsBlock{%
	\ryumeaning%
		{おたふくかぜ。風疹}{}{}%
	}%
	}%
	{\tailPrint[]}%
\entry
	{\headPrint[kana={とーすん\tl{˥}},ipa={toːsuɴ},pos={動},rui={2},para={とーはぬ},acc={H}]}%
	{%
		\MeaningsBlock{%
	\ryumeaning%
		{倒す}{サトウキビの収穫にも言う}{}%
	}%
	}%
	{\tailPrint[]}%
\entry
	{\headPrint[kana={とーすん\tl{˥}},ipa={toːsuɴ},pos={動},rui={2},para={とーはぬ},acc={H}]}%
	{%
		\MeaningsBlock{%
	\ryumeaning%
		{通す}{}{}%
	}%
	}%
	{\tailPrint[]}%
\entry
	{\headPrint[kana={\josho{}どー\kako{}でぃん},ipa={doːdiɴ},pos={副},acc={HL}]}%
	{%
		\MeaningsBlock{%
	\ryumeaning%
		{どうぞ。是非とも}{}{}%
	}%
	}%
	{\tailPrint[]}%
\entry
	{\headPrint[kana={とーに\tl{˥}},ipa={toːni},pos={名},acc={H}]}%
	{%
		\MeaningsBlock{%
	\ryumeaning%
		{豚の餌入れ}{}{}%
	}%
	}%
	{\tailPrint[]}%
\entry
	{\headPrint[kana={とーぴくん\tl{˥}},ipa={toːpi̥kuɴ},pos={動},rui={1},para={とーぴかぬ},acc={H}]}%
	{%
		\MeaningsBlock{%
	\ryumeaning%
		{\ruby[g]{擦}{こす}る。磨く}{}{}%
	}%
	}%
	{\tailPrint[]}%
\entry
	{\headPrint[kana={とーら\tl{˥˩}},ipa={toːra},pos={名},acc={F}]}%
	{%
		\MeaningsBlock{%
	\ryumeaning%
		{炊事小屋。福屋}{作業小屋を兼ねる}{}%
	}%
	}%
	{\tailPrint[note={「唐倉」の転。白保方言も下降型である}]}%
\entry
	{\headPrint[kana={どーり\tl{˩˥}},ipa={doːri},pos={名},acc={R}]}%
	{%
		\MeaningsBlock{%
	\ryumeaning%
		{道理。ことわり}{}{}%
	}%
	}%
	{\tailPrint[]}%
\entry
	{\headPrint[kana={とーりるん\tl{˥}},ipa={toːriruɴ},pos={動},rui={3},para={とーるぬ},acc={H}]}%
	{%
		\MeaningsBlock{%
	\ryumeaning%
		{倒れる}{}{}%
	\ryumeaning%
		{滅ぶ。滅亡する}{}{}%
	}%
	}%
	{\tailPrint[]}%
\entry
	{\headPrint[kana={とーるん\tl{˥}},ipa={toːruɴ},pos={動},rui={1},para={とーらぬ},acc={H}]}%
	{%
		\MeaningsBlock{%
	\ryumeaning%
		{通る}{}{}%
	\ryumeaning%
		{通用する。認められる}{}{}%
	}%
	}%
	{\tailPrint[]}%
\entry
	{\headPrint[kana={とーん\tl{˥˩}},ipa={toːɴ},pos={動},para={とわぬ},acc={F}]}%
	{%
		\MeaningsBlock{%
	\ryumeaning%
		{問う。尋ねる}{}{}%
	}%
	}%
	{\tailPrint[]}%
\entry
	{\headPrint[kana={どぅー\tl{˩˥}},ipa={duː},pos={名},acc={R}]}%
	{%
		\MeaningsBlock{%
	\ryumeaning%
		{自分。体。胴体}{}{}%
	}%
	}%
	{\tailPrint[]}%
\entry
	{\headPrint[kana={どぅー\tl{˩˥}かってぃ\tl{˥˩}},ipa={duːkatti},pos={名},acc={R+F}]}%
	{%
		\MeaningsBlock{%
	\ryumeaning%
		{自分勝手}{}{}%
	}%
	}%
	{\tailPrint[]}%
\entry
	{\headPrint[kana={どぅーかってぃ\ws{}しゃーん},ipa={duːkatti ʃaːɴ},pos={句}]}%
	{%
		\MeaningsBlock{%
	\ryumeaning%
		{自分勝手している。自分勝手だ}{}{}%
	}%
	}%
	{\tailPrint[]}%
\entry
	{\headPrint[kana={どぅー\tl{˩˥}がんじゅー\tl{˩˥}さーん},ipa={duːgandʒuːsaːɴ},pos={形},acc={R+R}]}%
	{%
		\MeaningsBlock{%
	\ryumeaning%
		{体が健康だ}{}{}%
	}%
	}%
	{\tailPrint[]}%
\entry
	{\headPrint[kana={\josho{}どぅー\kako{}し},ipa={duːʃi},pos={句},acc={HL}]}%
	{%
		\MeaningsBlock{%
	\ryumeaning%
		{自分で}{}{}%
	}%
	}%
	{\tailPrint[]}%
\entry
	{\headPrint[kana={どぅー\tl{˩˥}ずさ\tl{˩˥}はん},ipa={duːdzusahaɴ},pos={形},para={どうずさへぬ},acc={R+H}]}%
	{%
		\MeaningsBlock{%
	\ryumeaning%
		{体が丈夫だ。健康だ}{}{}%
	}%
	}%
	{\tailPrint[]}%
\entry
	{\headPrint[kana={どぅー\tl{˩˥}だるさ\tl{˩˥}はーん},ipa={duːdarusahaːɴ},pos={形},acc={R+R}]}%
	{%
		\MeaningsBlock{%
	\ryumeaning%
		{}{体がだるくて重い感じだ}{}%
	}%
	}%
	{\tailPrint[]}%
\entry
	{\headPrint[kana={どぅー\tl{˩˥}たんき\tl{˥}},ipa={duːtaŋki},pos={名},acc={R+H}]}%
	{%
		\MeaningsBlock{%
	\ryumeaning%
		{骨惜しみ}{}{}%
	}%
	}%
	{\tailPrint[]}%
\entry
	{\headPrint[kana={とぅーぴとぅち},ipa={tuːpi̥tutʃi},pos={名}]}%
	{%
		\MeaningsBlock{%
	\ryumeaning%
		{十一。十一個}{}{}%
	}%
	}%
	{\tailPrint[]}%
\entry
	{\headPrint[kana={どぅー\tl{˩˥}よがすん\tl{˩˥}},ipa={duːjogasuɴ},pos={動},rui={2},para={どぅよがはぬ},acc={R+R}]}%
	{%
		\MeaningsBlock{%
	\ryumeaning%
		{休む。休憩する}{}{}%
	}%
	}%
	{\tailPrint[]}%
\entry
	{\headPrint[kana={どぅーる\tl{˩˥}},ipa={duːru},pos={名},acc={R}]}%
	{%
		\MeaningsBlock{%
	\ryumeaning%
		{\ruby[g]{泥}{どろ}}{}{}%
	}%
	}%
	{\tailPrint[]}%
\entry
	{\headPrint[kana={とぅいしみるん\tl{˥˩}},ipa={tuiʃimiruɴ},pos={動},rui={3},para={とぅいすむぬ},acc={F}]}%
	{%
		\MeaningsBlock{%
	\ryumeaning%
		{問い詰める}{}{}%
	}%
	}%
	{\tailPrint[]}%
\entry
	{\headPrint[kana={とぅいすくん\tl{˥˩}},ipa={tuisu̥kuɴ},pos={動},rui={1},para={とぅいすかぬ},acc={F}]}%
	{%
		\MeaningsBlock{%
	\ryumeaning%
		{問い聞く}{}{}%
	}%
	}%
	{\tailPrint[]}%
\entry
	{\headPrint[kana={とぅが\tl{˥}},ipa={tuga},pos={名},acc={H}]}%
	{%
		\MeaningsBlock{%
	\ryumeaning%
		{罪。罰}{}{}%
	}%
	}%
	{\tailPrint[]}%
\entry
	{\headPrint[kana={とぅがにん},ipa={tuganin},pos={名}]}%
	{%
		\MeaningsBlock{%
	\ryumeaning%
		{罪人}{}{}%
	}%
	}%
	{\tailPrint[]}%
\entry
	{\headPrint[kana={とぅがみるん\tl{˥}},ipa={tugamiruɴ},pos={動},rui={3},para={とぅがむぬ},acc={H}]}%
	{%
		\MeaningsBlock{%
	\ryumeaning%
		{咎める。罰する}{}{}%
	}%
	}%
	{\tailPrint[]}%
\entry
	{\headPrint[kana={とぅきるん\tl{˥}},ipa={tu̥kiruɴ},pos={動},rui={3},para={とぅくぬ},acc={H}]}%
	{%
		\MeaningsBlock{%
	\ryumeaning%
		{溶ける。融ける}{}{}%
	}%
	}%
	{\tailPrint[]}%
\entry
	{\headPrint[kana={とぅく\tl{˥}},ipa={tu̥ku},pos={名},acc={H}]}%
	{%
		\MeaningsBlock{%
	\ryumeaning%
		{得。利得}{}{}%
	}%
	}%
	{\tailPrint[]}%
\entry
	{\headPrint[kana={とぅく\tl{˥}},ipa={tu̥ku},pos={名},acc={H}]}%
	{%
		\MeaningsBlock{%
	\ryumeaning%
		{徳}{}{}%
	}%
	}%
	{\tailPrint[]}%
\entry
	{\headPrint[kana={とぅく\tl{˥˩}},ipa={tu̥ku},pos={名},acc={F}]}%
	{%
		\MeaningsBlock{%
	\ryumeaning%
		{仏壇。床の間}{}{}%
	}%
	}%
	{\tailPrint[]}%
\entry
	{\headPrint[kana={とぅく\tl{˥˩}},ipa={tu̥ku},pos={名},acc={F}]}%
	{%
		\MeaningsBlock{%
	\ryumeaning%
		{寝床}{}{}%
	}%
	}%
	{\tailPrint[]}%
\entry
	{\headPrint[kana={どぅく\tl{˩˥}},ipa={duku},pos={名},acc={R}]}%
	{%
		\MeaningsBlock{%
	\ryumeaning%
		{\ruby[g]{毒}{どく}}{}{}%
	}%
	}%
	{\tailPrint[]}%
\entry
	{\headPrint[kana={どぅぐ},ipa={dugu},pos={副}]}%
	{%
		\MeaningsBlock{%
	\ryumeaning%
		{\ruby[g]{酷}{ひど}い}{}{}%
	}%
	}%
	{\tailPrint[]}%
\entry
	{\headPrint[kana={とぅくっとぅ\tl{˥}},ipa={tu̥kuttu},pos={名},acc={H}]}%
	{%
		\MeaningsBlock{%
	\ryumeaning%
		{安心。安堵}{}{}%
	}%
	}%
	{\tailPrint[]}%
\entry
	{\headPrint[kana={どぅ\josho{}ぐ\kako{}どぅぐ},ipa={dugudugu},pos={副},acc={重複}]}%
	{%
		\MeaningsBlock{%
	\ryumeaning%
		{あんまりだ}{\emb{どぅぐ}の強調}{}%
	}%
	}%
	{\tailPrint[]}%
\entry
	{\headPrint[kana={どぅぐりしゃ\tl{˩˥}はーん},ipa={duguriʃahaːɴ},pos={形},para={どくりしゃへぬ},acc={R}]}%
	{%
		\MeaningsBlock{%
	\ryumeaning%
		{気の毒だ。気まずい。可哀そうだ。恐縮だ}{}{}%
	}%
	}%
	{\tailPrint[]}%
\entry
	{\headPrint[kana={とぅくる\tl{˥}},ipa={tu̥kuru},pos={名},acc={H}]}%
	{%
		\MeaningsBlock{%
	\ryumeaning%
		{所。場所}{}{}%
	}%
	}%
	{\tailPrint[]}%
\entry
	{\headPrint[kana={とぅぐん\tl{˥}},ipa={tuguɴ},pos={動},rui={1},para={とがぬ},acc={H}]}%
	{%
		\MeaningsBlock{%
	\ryumeaning%
		{研ぐ。研磨する}{}{}%
	}%
	}%
	{\tailPrint[]}%
\entry
	{\headPrint[kana={とぅけー\tl{˥˩}},ipa={tu̥keː},pos={名},acc={F}]}%
	{%
		\MeaningsBlock{%
	\ryumeaning%
		{遠い所。遠方}{}{}%
	}%
	}%
	{\tailPrint[]}%
\entry
	{\headPrint[kana={どぅげーるん\tl{˩˥}},ipa={dugeːruɴ},pos={動},rui={1},para={どげーらぬ},acc={R}]}%
	{%
		\MeaningsBlock{%
	\ryumeaning%
		{\ruby[g]{怒鳴}{どな}る}{}{}%
	}%
	}%
	{\tailPrint[]}%
\entry
	{\headPrint[kana={とぅさ\tl{˥˩}はん},ipa={tu̥sahaɴ},pos={形},para={とうさへぬ},acc={F}]}%
	{%
		\MeaningsBlock{%
	\ryumeaning%
		{遠い。遠方だ}{}{}%
	}%
	}%
	{\tailPrint[]}%
\entry
	{\headPrint[kana={とぅし\tl{˥}},ipa={tuʃi},pos={名},acc={H}]}%
	{%
		\MeaningsBlock{%
	\ryumeaning%
		{\ruby[g]{砥石}{といし}}{}{}%
	}%
	}%
	{\tailPrint[]}%
\entry
	{\headPrint[kana={どぅし\tl{˥˩}},ipa={duʃi},pos={名},acc={F}]}%
	{%
		\MeaningsBlock{%
	\ryumeaning%
		{友人。友達}{}{}%
	}%
	}%
	{\tailPrint[]}%
\entry
	{\headPrint[kana={とぅしぃぬ\tl{˥}\ws{}ゆー\tl{˩˥}},ipa={tu̥ʃinu juː},pos={句},acc={H R}]}%
	{%
		\MeaningsBlock{%
	\ryumeaning%
		{大晦日の夜}{}{}%
	}%
	}%
	{\tailPrint[]}%
\entry
	{\headPrint[kana={とぅしぃぬ\tl{˥}\ws{}ゆる\tl{˩˥}},ipa={tu̥ʃinu juru},pos={句},acc={H R}]}%
	{%
		\MeaningsBlock{%
	\ryumeaning%
		{大晦日の夜}{}{}%
	}%
	}%
	{\tailPrint[]}%
\entry
	{\headPrint[kana={とぅしきるん\tl{˥}},ipa={tu̥ʃikiruɴ},pos={動},rui={3},para={とぅすくぬ},acc={H}]}%
	{%
		\MeaningsBlock{%
	\ryumeaning%
		{説得する}{}{}%
	}%
	}%
	{\tailPrint[]}%
\entry
	{\headPrint[kana={どぅしけな},ipa={duʃi̥kena},pos={句}]}%
	{%
		\MeaningsBlock{%
	\ryumeaning%
		{各自で。自分で}{自分でやるなど}{}%
	}%
	}%
	{\tailPrint[]}%
\entry
	{\headPrint[kana={とぅしみるん\tl{˥}},ipa={tuʃi̥miruɴ},pos={動},rui={3},para={とぅしむぬ},acc={H}]}%
	{%
		\MeaningsBlock{%
	\ryumeaning%
		{問い責める。問い詰める}{}{}%
	}%
	}%
	{\tailPrint[]}%
\entry
	{\headPrint[kana={とぅじみるん\tl{˥˩}},ipa={tudʒimiruɴ},pos={動},rui={3},acc={F}]}%
	{%
		\MeaningsBlock{%
	\ryumeaning%
		{終える}{}{}%
	}%
	}%
	{\tailPrint[]}%
\entry
	{\headPrint[kana={とぅちぃ\tl{˥}},ipa={tu̥tʃi},pos={名},acc={H}]}%
	{%
		\MeaningsBlock{%
	\ryumeaning%
		{年。歳。年齢}{}{}%
	}%
	}%
	{\tailPrint[]}%
\entry
	{\headPrint[kana={とぅちぃとぅるん\tl{˥}},ipa={tu̥tʃitu̥ruɴ},pos={動},rui={1},para={とぅちとぅらぬ},acc={H}]}%
	{%
		\MeaningsBlock{%
	\ryumeaning%
		{年を取る}{}{}%
	}%
	}%
	{\tailPrint[]}%
\entry
	{\headPrint[kana={とぅとぅぎるん\tl{˥}},ipa={tu̥tugiruɴ},pos={動},rui={1},acc={H}]}%
	{%
		\MeaningsBlock{%
	\ryumeaning%
		{届ける}{}{}%
	}%
	}%
	{\tailPrint[]}%
\entry
	{\headPrint[kana={とぅとぅぐん\tl{˥}},ipa={tu̥tuguɴ},pos={動},rui={1},para={ととがぬ},acc={H}]}%
	{%
		\MeaningsBlock{%
	\ryumeaning%
		{届く}{}{}%
	}%
	}%
	{\tailPrint[]}%
\entry
	{\headPrint[kana={とぅとぅのーん\tl{˥}},ipa={tu̥tunoːɴ},pos={動},para={とぅとぅのーあぬ},acc={H}]}%
	{%
		\MeaningsBlock{%
	\ryumeaning%
		{整う。きちんとなる。うまくまとまる}{}{}%
	}%
	}%
	{\tailPrint[]}%
\entry
	{\headPrint[kana={どぅなー\tl{˩˥}さーん},ipa={dunaːsaːɴ},pos={形},acc={R}]}%
	{%
		\MeaningsBlock{%
	\ryumeaning%
		{（動作が）緩慢である。のろのろしている}{}{}%
	}%
	}%
	{\tailPrint[]}%
\entry
	{\headPrint[kana={とぅなばるん\tl{˥˩}},ipa={tu̥nabaruɴ},pos={動},rui={1},para={とぅなばらぬ},acc={F}]}%
	{%
		\MeaningsBlock{%
	\ryumeaning%
		{黙る。押し黙る}{}{}%
	}%
	}%
	{\tailPrint[]}%
\entry
	{\headPrint[kana={どぅなはん},ipa={dunahaɴ},pos={形},acc={-}]}%
	{%
		\MeaningsBlock{%
	\ryumeaning%
		{\ruby[g]{鈍}{のろ}い}{}{}%
	}%
	}%
	{\tailPrint[]}%
\entry
	{\headPrint[kana={とぅぬすく\tl{˥˩}},ipa={tu̥nusu̥ku},pos={名},acc={F}]}%
	{%
		\MeaningsBlock{%
	\ryumeaning%
		{\ruby[g]{登}{と}\ruby[g]{野}{の}\ruby[g]{城}{しろ}}{石垣島の地名}{}%
	}%
	}%
	{\tailPrint[]}%
\entry
	{\headPrint[kana={とぅのー\tl{˥}},ipa={tu̥noː},pos={名},acc={H}]}%
	{%
		\MeaningsBlock{%
	\ryumeaning%
		{アカテツ}{樹木名}{}%
	}%
	}%
	{\tailPrint[]}%
\entry
	{\headPrint[kana={\josho{}どぅばだ\kako{}にげー},ipa={dubadanigeː},pos={名},acc={H+L}]}%
	{%
		\MeaningsBlock{%
	\ryumeaning%
		{健康祈願}{\emb{どぅ―}「体」と\emb{ぱだ}「肌」の「願い」からか}{}%
	}%
	}%
	{\tailPrint[]}%
\entry
	{\headPrint[kana={とぅぴくいるん\tl{˥˩}},ipa={tu̥pikuiruɴ},pos={動},acc={F}]}%
	{%
		\MeaningsBlock{%
	\ryumeaning%
		{飛び越える}{}{}%
	}%
	}%
	{\tailPrint[]}%
\entry
	{\headPrint[kana={とぅぴゆ\tl{˥˩}},ipa={tu̥piju},pos={名},acc={F}]}%
	{%
		\MeaningsBlock{%
	\ryumeaning%
		{飛魚}{}{}%
	}%
	}%
	{\tailPrint[]}%
\entry
	{\headPrint[kana={とぅぷちん\tl{˥˩}},ipa={tu̥putʃiɴ},pos={動},rui={3},para={とぅぷとぅぬ},acc={F}]}%
	{%
		\MeaningsBlock{%
	\ryumeaning%
		{踏み外す。飛び込む}{}{}%
	}%
	}%
	{\tailPrint[]}%
\entry
	{\headPrint[kana={とぅぷん\tl{˥˩}},ipa={tu̥puɴ},pos={動},rui={1},para={とぅぱぬ},acc={F}]}%
	{%
		\MeaningsBlock{%
	\ryumeaning%
		{飛ぶ。飛行する}{}{}%
	}%
	}%
	{\tailPrint[]}%
\entry
	{\headPrint[kana={とぅまらすん\tl{˥˩}},ipa={tu̥marasuɴ},pos={動},rui={2},para={とぅまらはぬ},acc={F}]}%
	{%
		\MeaningsBlock{%
	\ryumeaning%
		{泊まらせる}{}{}%
	}%
	}%
	{\tailPrint[]}%
\entry
	{\headPrint[kana={とぅまるん\tl{˥˩}},ipa={tu̥maruɴ},pos={動},rui={1},para={とぅまらぬ},acc={F}]}%
	{%
		\MeaningsBlock{%
	\ryumeaning%
		{泊まる。宿る。宿泊する}{}{}%
	}%
	}%
	{\tailPrint[]}%
\entry
	{\headPrint[kana={とぅまるん\tl{˥˩}},ipa={tu̥maruɴ},pos={動},rui={1},para={とーまらぬ},acc={F}]}%
	{%
		\MeaningsBlock{%
	\ryumeaning%
		{窪む。へこむ}{}{}%
	}%
	}%
	{\tailPrint[]}%
\entry
	{\headPrint[kana={とぅみるん\tl{˥}},ipa={tu̥miruɴ},pos={動},rui={3}]}%
	{%
		\MeaningsBlock{%
	\ryumeaning%
		{止める}{}{}%
	}%
	}%
	{\tailPrint[]}%
\entry
	{\headPrint[kana={とぅみるん\tl{˥}},ipa={tu̥miruɴ},pos={動},rui={3},para={とぅむぬ},acc={H}]}%
	{%
		\MeaningsBlock{%
	\ryumeaning%
		{探す。見つける}{}{}%
	\ryumeaning%
		{拾う}{}{}%
	\ryumeaning%
		{求める}{}{}%
	}%
	}%
	{\tailPrint[]}%
\entry
	{\headPrint[kana={とぅみん\tl{˥}},ipa={tu̥miɴ},pos={動},rui={3},para={とぅむぬ},acc={H}]}%
	{%
		\MeaningsBlock{%
	\ryumeaning%
		{探す。見つける}{}{}%
	\ryumeaning%
		{（嫁を）探す。（妻を）娶る}{}{}%
	}%
	}%
	{\tailPrint[]}%
\entry
	{\headPrint[kana={とぅむ\tl{˥}},ipa={tu̥mu},pos={名},acc={H}]}%
	{%
		\MeaningsBlock{%
	\ryumeaning%
		{\ruby[g]{艫}{とも}}{船の船尾}{}%
	}%
	}%
	{\tailPrint[]}%
\entry
	{\headPrint[kana={とぅゆますん\tl{˥}},ipa={tujumasuɴ},pos={動},rui={2},para={とぅゆまはぬ},acc={H}]}%
	{%
		\MeaningsBlock{%
	\ryumeaning%
		{轟かせる。とよます}{}{}%
	}%
	}%
	{\tailPrint[]}%
\entry
	{\headPrint[kana={とぅゆまりん\tl{˥}},ipa={tujumariɴ},pos={動},rui={1},para={とぅゆまらぬ},acc={H}]}%
	{%
		\MeaningsBlock{%
	\ryumeaning%
		{とよむ。世に鳴り響く}{}{}%
	}%
	}%
	{\tailPrint[]}%
\entry
	{\headPrint[kana={とぅら\tl{˥˩}},ipa={tura},pos={名},acc={F}]}%
	{%
		\MeaningsBlock{%
	\ryumeaning%
		{\ruby[g]{虎}{とら}}{十二支の寅}{}%
	}%
	}%
	{\tailPrint[]}%
\entry
	{\headPrint[kana={とぅり\tl{˥}},ipa={turi},pos={名},acc={H}]}%
	{%
		\MeaningsBlock{%
	\ryumeaning%
		{灯り。ランプ}{}{}%
	}%
	}%
	{\tailPrint[]}%
\entry
	{\headPrint[kana={とぅり\tl{˥˩}},ipa={tu̥ri},pos={名},acc={F}]}%
	{%
		\MeaningsBlock{%
	\ryumeaning%
		{鳥}{}{}%
	}%
	}%
	{\tailPrint[]}%
\entry
	{\headPrint[kana={とぅり\tl{˥˩}},ipa={tu̥ri},pos={名},acc={F}]}%
	{%
		\MeaningsBlock{%
	\ryumeaning%
		{\ruby[g]{酉}{とり}}{十二支の酉}{}%
	}%
	}%
	{\tailPrint[]}%
\entry
	{\headPrint[kana={とぅりかいすん\tl{˥}},ipa={tu̥rikaisuɴ},pos={動},rui={2b},para={とぅりかいさぬ},acc={H}]}%
	{%
		\MeaningsBlock{%
	\ryumeaning%
		{取り返す。取り戻す}{}{}%
	}%
	}%
	{\tailPrint[]}%
\entry
	{\headPrint[kana={とぅりかいるん\tl{˥}},ipa={tu̥rikairuɴ},pos={動},rui={1},para={とぅりかいらぬ},acc={H}]}%
	{%
		\MeaningsBlock{%
	\ryumeaning%
		{取り替える}{}{}%
	}%
	}%
	{\tailPrint[]}%
\entry
	{\headPrint[kana={とぅりくむん},ipa={tu̥riku̥muɴ},pos={動},rui={1},para={とぅりくまぬ}]}%
	{%
		\MeaningsBlock{%
	\ryumeaning%
		{取り込む}{取ってしまい込む}{}%
	}%
	}%
	{\tailPrint[]}%
\entry
	{\headPrint[kana={とぅりくむん\tl{˥}},ipa={tu̥rikumuɴ},pos={動},rui={1},para={とぅりくまぬ},acc={H}]}%
	{%
		\MeaningsBlock{%
	\ryumeaning%
		{取り組む}{}{}%
	}%
	}%
	{\tailPrint[]}%
\entry
	{\headPrint[kana={とぅりけーすん\tl{˥}},ipa={tu̥rikeːsuɴ},pos={動},rui={2b},para={とぅりけーさぬ},acc={H}]}%
	{%
		\MeaningsBlock{%
	\ryumeaning%
		{取り返す}{}{}%
	}%
	}%
	{\tailPrint[]}%
\entry
	{\headPrint[kana={とぅりしまるん\tl{˥}},ipa={tu̥riʃimaruɴ},pos={動},rui={1},para={とぅりしまらぬ},acc={H}]}%
	{%
		\MeaningsBlock{%
	\ryumeaning%
		{取り締まる}{}{}%
	}%
	}%
	{\tailPrint[]}%
\entry
	{\headPrint[kana={とぅりたちん\tl{˥}},ipa={tu̥ritḁtʃiɴ},pos={動},rui={3},para={とぅりたとぅぬ},acc={H}]}%
	{%
		\MeaningsBlock{%
	\ryumeaning%
		{取り立てる}{}{}%
	}%
	}%
	{\tailPrint[]}%
\entry
	{\headPrint[kana={とぅり\tl{˥}\ws{}っしるん\tl{˥˩}},ipa={tu̥ri ʃʃiruɴ},pos={句},acc={H F}]}%
	{%
		\MeaningsBlock{%
	\ryumeaning%
		{取り除く。取って捨てる}{}{}%
	}%
	}%
	{\tailPrint[]}%
\entry
	{\headPrint[kana={\josho{}とぅり\kako{}\ws{}っしん\tl{˥˩}},ipa={tu̥ri ʃʃiɴ},pos={句},para={とりすとぬ},acc={H L}]}%
	{%
		\MeaningsBlock{%
	\ryumeaning%
		{取り去る。取って捨てる}{}{}%
	}%
	}%
	{\tailPrint[]}%
\entry
	{\headPrint[kana={とぅりはかろーん\tl{˥}},ipa={tu̥rihakaroːɴ},pos={動},rui={1},para={とぅりはからぬ},acc={H}]}%
	{%
		\MeaningsBlock{%
	\ryumeaning%
		{取り計らう}{}{}%
	}%
	}%
	{\tailPrint[]}%
\entry
	{\headPrint[kana={とぅり\tl{˥}ぱんつん\tl{˥˩}},ipa={tu̥ripantsuɴ},pos={動},para={とぅりぱんつぁぬ},acc={H+F}]}%
	{%
		\MeaningsBlock{%
	\ryumeaning%
		{取り外す}{}{}%
	}%
	}%
	{\tailPrint[]}%
\entry
	{\headPrint[kana={とぅりぴんがすん\tl{˥}},ipa={tu̥ripiŋgasuɴ},pos={動},rui={2},para={とぅりぴんがはぬ},acc={H}]}%
	{%
		\MeaningsBlock{%
	\ryumeaning%
		{取り逃がす}{}{}%
	}%
	}%
	{\tailPrint[]}%
\entry
	{\headPrint[kana={とぅりむつん\tl{˥}},ipa={tu̥rimutsuɴ},pos={動},para={とぅりむたぬ},acc={H}]}%
	{%
		\MeaningsBlock{%
	\ryumeaning%
		{持てなす。接待する}{}{}%
	}%
	}%
	{\tailPrint[]}%
\entry
	{\headPrint[kana={とぅりむどぅすん\tl{˥}},ipa={tu̥rimudusuɴ},pos={動},rui={2b},para={とぅりむどぅさぬ},acc={H}]}%
	{%
		\MeaningsBlock{%
	\ryumeaning%
		{取り戻す}{}{}%
	}%
	}%
	{\tailPrint[]}%
\entry
	{\headPrint[kana={とぅりゆしるん},ipa={turijuʃiruɴ},pos={動}]}%
	{%
		\MeaningsBlock{%
	\ryumeaning%
		{取り寄せる}{}{}%
	}%
	}%
	{\tailPrint[]}%
\entry
	{\headPrint[kana={とぅりるん\tl{˥˩}},ipa={tu̥riruɴ},pos={動},rui={3},para={とぅるぬ},acc={F}]}%
	{%
		\MeaningsBlock{%
	\ryumeaning%
		{凪ぐ。静かになる}{}{}%
	}%
	}%
	{\tailPrint[]}%
\entry
	{\headPrint[kana={とぅりんだすん\tl{˥}},ipa={tu̥rindasuɴ},pos={動},rui={2},para={とぅりんだはぬ},acc={H}]}%
	{%
		\MeaningsBlock{%
	\ryumeaning%
		{取り出す}{}{}%
	}%
	}%
	{\tailPrint[]}%
\entry
	{\headPrint[kana={どぅるびちゃー\tl{˩˥}\ws{}なるん\tl{˩˥}},ipa={durubitʃaː naruɴ},pos={句},acc={R R}]}%
	{%
		\MeaningsBlock{%
	\ryumeaning%
		{泥だらけになる}{}{}%
	}%
	}%
	{\tailPrint[]}%
\entry
	{\headPrint[kana={どぅるみち\tl{˩˥}},ipa={durumitʃi},pos={名},acc={R}]}%
	{%
		\MeaningsBlock{%
	\ryumeaning%
		{泥道}{}{}%
	}%
	}%
	{\tailPrint[]}%
\entry
	{\headPrint[kana={とぅるん\tl{˥}},ipa={tu̥ruɴ},pos={動},rui={1},para={とぅらぬ},acc={H}]}%
	{%
		\MeaningsBlock{%
	\ryumeaning%
		{取る。奪う}{}{}%
	}%
	}%
	{\tailPrint[]}%
\entry
	{\headPrint[kana={とぅん\tl{˥˩}},ipa={tuɴ},pos={名},acc={F}]}%
	{%
		\MeaningsBlock{%
	\ryumeaning%
		{妻。嫁}{}{}%
	}%
	}%
	{\tailPrint[]}%
\entry
	{\headPrint[kana={とぅんけーるん},ipa={tuŋkeːriɴ},pos={動}]}%
	{%
		\MeaningsBlock{%
	\ryumeaning%
		{振り返る。振り向く}{}{}%
	}%
	}%
	{\tailPrint[]}%
\entry
	{\headPrint[kana={とぅんじー\tl{˥}},ipa={tundʒiː},pos={名},acc={H}]}%
	{%
		\MeaningsBlock{%
	\ryumeaning%
		{冬至}{}{}%
	}%
	}%
	{\tailPrint[]}%
\entry
	{\headPrint[kana={とぅんじるん\tl{˥˩}},ipa={tundʒiruɴ},pos={動},acc={F}]}%
	{%
		\MeaningsBlock{%
	\ryumeaning%
		{飛び出る}{}{}%
	}%
	}%
	{\tailPrint[]}%
\entry
	{\headPrint[kana={とぅんそるん\tl{˥˩}},ipa={tunsoruɴ},pos={動},rui={1},para={とぅんそらぬ},acc={F}]}%
	{%
		\MeaningsBlock{%
	\ryumeaning%
		{娶る}{嫁として縁組する}{}%
	}%
	}%
	{\tailPrint[]}%
\entry
	{\headPrint[kana={とぅんたち\tl{˥˩}},ipa={tuntḁtʃi},pos={名},acc={F}]}%
	{%
		\MeaningsBlock{%
	\ryumeaning%
		{つま先立ち}{}{}%
	}%
	}%
	{\tailPrint[]}%
\entry
	{\headPrint[kana={とぅんたちびるん\tl{˥˩}},ipa={tuntḁtʃibiruɴ},pos={動},acc={F}]}%
	{%
		\MeaningsBlock{%
	\ryumeaning%
		{爪立つ}{爪先で立つ}{}%
	}%
	}%
	{\tailPrint[]}%
\entry
	{\headPrint[kana={とぅんちるん\tl{˥˩}},ipa={tuntʃiruɴ},pos={動},acc={F}]}%
	{%
		\MeaningsBlock{%
	\ryumeaning%
		{突き出る}{}{}%
	}%
	}%
	{\tailPrint[]}%
\entry
	{\headPrint[kana={とぅんぶとぅ\tl{˥˩}},ipa={tumbutu},pos={名},acc={F}]}%
	{%
		\MeaningsBlock{%
	\ryumeaning%
		{\ruby[g]{夫婦}{めおと}。\ruby[g]{夫婦}{ふうふ}}{}{}%
	}%
	}%
	{\tailPrint[]}%
\entry
	{\headPrint[kana={とっきん},ipa={tokkin},pos={名}]}%
	{%
		\MeaningsBlock{%
	\ryumeaning%
		{（野生の）バンザクロ}{}{}%
	}%
	}%
	{\tailPrint[]}%
\entry
	{\headPrint[kana={どっふぁった},ipa={doffatta},pos={擬}]}%
	{%
		\MeaningsBlock{%
	\ryumeaning%
		{どかっと。どっかり}{重いものを置いたり，大きな人が座ったりするさま}{}%
	}%
	}%
	{\tailPrint[]}%
\entry
	{\headPrint[kana={とふかすん\tl{˥˩}},ipa={tofu̥kasuɴ},pos={動},rui={2},para={とふかはぬ},acc={F}]}%
	{%
		\MeaningsBlock{%
	\ryumeaning%
		{突き通す}{}{}%
	}%
	}%
	{\tailPrint[]}%
\entry
	{\headPrint[kana={とまらすん},ipa={tomarasuɴ},pos={動}]}%
	{%
		\MeaningsBlock{%
	\ryumeaning%
		{窪ませる}{}{}%
	}%
	}%
	{\tailPrint[]}%
\entry
	{\headPrint[kana={どまんぐるん\tl{˩˥}},ipa={domaŋguruɴ},pos={動},rui={1},para={どまんぐらぬ},acc={R}]}%
	{%
		\MeaningsBlock{%
	\ryumeaning%
		{うろたえる。正気を失う}{}{}%
	}%
	}%
	{\tailPrint[]}%
\entry
	{\headPrint[kana={どみがすん\tl{˩˥}},ipa={domigasuɴ},pos={動},rui={2},para={どみがはぬ},acc={R}]}%
	{%
		\MeaningsBlock{%
	\ryumeaning%
		{突き崩す}{}{}%
	}%
	}%
	{\tailPrint[]}%
\entry
	{\headPrint[kana={どみんがすん\tl{˩˥}},ipa={domiŋgasuɴ},pos={動},rui={2},para={どみんがはぬ},acc={R}]}%
	{%
		\MeaningsBlock{%
	\ryumeaning%
		{叩きのめす}{}{}%
	}%
	}%
	{\tailPrint[]}%
\entry
	{\headPrint[kana={とらしみん},ipa={toraʃimiɴ},pos={動}]}%
	{%
		\MeaningsBlock{%
	\ryumeaning%
		{与える。取らせる}{}{}%
	}%
	}%
	{\tailPrint[]}%
\entry
	{\headPrint[kana={どらん\tl{˥˩}},ipa={doraɴ},pos={名},acc={F}]}%
	{%
		\MeaningsBlock{%
	\ryumeaning%
		{\ruby[g]{銅鑼}{どら}}{}{}%
	}%
	}%
	{\tailPrint[]}%
\entry
	{\headPrint[kana={とろふかすん\tl{˥}},ipa={torofu̥kasuɴ},pos={動},rui={2},para={とろふかぬ},acc={H}]}%
	{%
		\MeaningsBlock{%
	\ryumeaning%
		{下痢する}{}{}%
	}%
	}%
	{\tailPrint[]}%
\entry
	{\headPrint[kana={どんどん},ipa={dondoɴ},pos={擬}]}%
	{%
		\MeaningsBlock{%
	\ryumeaning%
		{どきどき}{心臓が鼓動する音の形容}{}%
	}%
	}%
	{\tailPrint[]}%
\entry
	{\headPrint[kana={どんぶり\tl{˩˥}},ipa={domburi},pos={名},acc={R}]}%
	{%
		\MeaningsBlock{%
	\ryumeaning%
		{\ruby[g]{丼}{どんぶり}}{食器名}{}%
	}%
	}%
	{\tailPrint[]}%
\LetterHead{な}
\entry
	{\headPrint[kana={なー\tl{˥}},ipa={naː},pos={名},acc={H}]}%
	{%
		\MeaningsBlock{%
	\ryumeaning%
		{ここ。此処に}{}{}%
	}%
	}%
	{\tailPrint[]}%
\entry
	{\headPrint[kana={なー\tl{˥˩}},ipa={naː},pos={名},acc={F}]}%
	{%
		\MeaningsBlock{%
	\ryumeaning%
		{名。名前}{}{}%
	}%
	}%
	{\tailPrint[]}%
\entry
	{\headPrint[kana={なーがすん\tl{˥˩}},ipa={naːgasuɴ},pos={動},rui={2},para={なーがはぬ},acc={F}]}%
	{%
		\MeaningsBlock{%
	\ryumeaning%
		{泣かす}{}{}%
	}%
	}%
	{\tailPrint[]}%
\entry
	{\headPrint[kana={なーぐん\tl{˥˩}},ipa={naːguɴ},pos={動},rui={1},para={なーがぬ},acc={F}]}%
	{%
		\MeaningsBlock{%
	\ryumeaning%
		{泣く}{}{}%
	\ryumeaning%
		{鳴く。（鳥が）囀る}{}{}%
	}%
	}%
	{\tailPrint[]}%
\entry
	{\headPrint[kana={なーしきうや\tl{˩˥}},ipa={naːʃi̥ki.uja},pos={名},acc={R}]}%
	{%
		\MeaningsBlock{%
	\ryumeaning%
		{名付け親}{}{}%
	}%
	}%
	{\tailPrint[]}%
\entry
	{\headPrint[kana={なーしきん\tl{˩˥}},ipa={naːʃi̥kiɴ},pos={動},rui={3},para={なんすくぬ},acc={R}]}%
	{%
		\MeaningsBlock{%
	\ryumeaning%
		{名を付ける。命名する}{}{}%
	}%
	}%
	{\tailPrint[]}%
\entry
	{\headPrint[kana={なーすん\tl{˩˥}},ipa={naːsuɴ},pos={動},rui={2},para={なーはぬ},acc={R}]}%
	{%
		\MeaningsBlock{%
	\ryumeaning%
		{\ruby[g]{綯}{な}う}{縄を綯う}{}%
	}%
	}%
	{\tailPrint[]}%
\entry
	{\headPrint[kana={なー\tl{˥˩}たか\tl{˥}はん},ipa={naːtakḁhaɴ},pos={形},para={なーたかへぬ},acc={F+H}]}%
	{%
		\MeaningsBlock{%
	\ryumeaning%
		{名高い。有名だ}{}{}%
	}%
	}%
	{\tailPrint[]}%
\entry
	{\headPrint[kana={なーどぅーどぅ},ipa={naːduːdu},pos={句}]}%
	{%
		\MeaningsBlock{%
	\ryumeaning%
		{各自。自分}{}{}%
	}%
	}%
	{\tailPrint[]}%
\entry
	{\headPrint[kana={な\josho{}ー\kako{}なー\ws{}しゃん\tl{˥˩}},ipa={naːnaː ʃaɴ},pos={句},acc={重複 L}]}%
	{%
		\MeaningsBlock{%
	\ryumeaning%
		{長々とした}{}{}%
	}%
	}%
	{\tailPrint[]}%
\entry
	{\headPrint[kana={なーばぐ\tl{˩˥}},ipa={naːbagu},pos={名},acc={R}]}%
	{%
		\MeaningsBlock{%
	\ryumeaning%
		{長持ち。衣類入れ}{}{}%
	}%
	}%
	{\tailPrint[]}%
\entry
	{\headPrint[kana={なー\tl{˩˥}はん},ipa={naːhaɴ},pos={形},para={まろはん},acc={R}]}%
	{%
		\MeaningsBlock{%
	\ryumeaning%
		{長い。永い}{}{}%
	}%
	}%
	{\tailPrint[]}%
\entry
	{\headPrint[kana={なーびら\tl{˩˥}},ipa={naːbira},pos={名},acc={R}]}%
	{%
		\MeaningsBlock{%
	\ryumeaning%
		{糸瓜}{}{}%
	}%
	}%
	{\tailPrint[]}%
\entry
	{\headPrint[kana={なーふく\tl{˩˥}},ipa={naːfu̥ku},pos={名},acc={R}]}%
	{%
		\MeaningsBlock{%
	\ryumeaning%
		{目隠しの石垣}{中の石垣の義。沖縄で言うヒンプン}{}%
	}%
	}%
	{\tailPrint[]}%
\entry
	{\headPrint[kana={なーぶに\tl{˩˥}},ipa={naːbuni},pos={名},acc={R}]}%
	{%
		\MeaningsBlock{%
	\ryumeaning%
		{背骨}{}{}%
	}%
	}%
	{\tailPrint[]}%
\entry
	{\headPrint[kana={なー\tl{˩˥}むら\tl{˥˩}},ipa={naːmura},pos={名},acc={R+F}]}%
	{%
		\MeaningsBlock{%
	\ryumeaning%
		{前村}{波照間島の集落名}{}%
	}%
	}%
	{\tailPrint[]}%
\entry
	{\headPrint[kana={なーらぼん\tl{˥}},ipa={naːrabon},pos={名},acc={H}]}%
	{%
		\MeaningsBlock{%
	\ryumeaning%
		{ナーラボン}{山芋の種類}{}%
	}%
	}%
	{\tailPrint[]}%
\entry
	{\headPrint[kana={なーり\tl{˩˥}},ipa={naːri},pos={名},acc={R}]}%
	{%
		\MeaningsBlock{%
	\ryumeaning%
		{実。果物}{}{}%
	}%
	}%
	{\tailPrint[]}%
\entry
	{\headPrint[kana={なーりむぬ\tl{˩˥}},ipa={naːrimunu},pos={名},acc={R}]}%
	{%
		\MeaningsBlock{%
	\ryumeaning%
		{果物。果樹}{}{}%
	}%
	}%
	{\tailPrint[]}%
\entry
	{\headPrint[kana={なーりん\tl{˩˥}},ipa={naːriɴ},pos={動},rui={3},para={なーるぬ},acc={R}]}%
	{%
		\MeaningsBlock{%
	\ryumeaning%
		{流れる。漂流する}{}{}%
	}%
	}%
	{\tailPrint[]}%
\entry
	{\headPrint[kana={なーるん\tl{˥˩}},ipa={naːruɴ},pos={動},rui={1},para={なーらぬ},acc={F}]}%
	{%
		\MeaningsBlock{%
	\ryumeaning%
		{鳴る。騒ぐ}{}{}%
	}%
	}%
	{\tailPrint[]}%
\entry
	{\headPrint[kana={なーるん\tl{˩˥}},ipa={naːruɴ},pos={動},rui={1},para={なーらぬ},acc={R}]}%
	{%
		\MeaningsBlock{%
	\ryumeaning%
		{出来る。実がなる}{}{}%
	}%
	}%
	{\tailPrint[]}%
\entry
	{\headPrint[kana={なーれるん\tl{˩˥}},ipa={naːreruɴ},pos={動},rui={3},para={なーるぬ},acc={R}]}%
	{%
		\MeaningsBlock{%
	\ryumeaning%
		{慣れる}{}{}%
	}%
	}%
	{\tailPrint[]}%
\entry
	{\headPrint[kana={ないぐん\tl{˩˥}},ipa={naiguɴ},pos={動},acc={R}]}%
	{%
		\MeaningsBlock{%
	\ryumeaning%
		{びっこをひく。びっこになる}{}{}%
	\ryumeaning%
		{}{足腰が立たなくなる}{}%
	}%
	}%
	{\tailPrint[]}%
\entry
	{\headPrint[kana={なが\tl{˥˩}},ipa={naga},pos={名},acc={F}]}%
	{%
		\MeaningsBlock{%
	\ryumeaning%
		{中。中に}{}{}%
	}%
	}%
	{\tailPrint[]}%
\entry
	{\headPrint[kana={ながあみ\tl{˥˩}},ipa={naga.ami},pos={名},acc={F}]}%
	{%
		\MeaningsBlock{%
	\ryumeaning%
		{長雨}{}{}%
	}%
	}%
	{\tailPrint[]}%
\entry
	{\headPrint[kana={ながいき\tl{˥˩}},ipa={naga.iki},pos={名},acc={F}]}%
	{%
		\MeaningsBlock{%
	\ryumeaning%
		{長生き。長寿}{}{}%
	}%
	}%
	{\tailPrint[]}%
\entry
	{\headPrint[kana={ながいぎ\tl{˥˩}\ws{}すん\tl{˥˩}},ipa={naga.igi suɴ},pos={句},acc={F F}]}%
	{%
		\MeaningsBlock{%
	\ryumeaning%
		{長生きする}{}{}%
	}%
	}%
	{\tailPrint[]}%
\entry
	{\headPrint[kana={なかざら\tl{˩˥}},ipa={nakadzara},pos={名},acc={R}]}%
	{%
		\MeaningsBlock{%
	\ryumeaning%
		{中皿}{}{}%
	}%
	}%
	{\tailPrint[]}%
\entry
	{\headPrint[kana={なかだち\tl{˩˥}},ipa={nakadatʃi},pos={名},acc={R}]}%
	{%
		\MeaningsBlock{%
	\ryumeaning%
		{仲人。媒酌人}{}{}%
	}%
	}%
	{\tailPrint[]}%
\entry
	{\headPrint[kana={ながびくん\tl{˥˩}},ipa={nagabikuɴ},pos={動},rui={1},para={ながびかぬ},acc={F}]}%
	{%
		\MeaningsBlock{%
	\ryumeaning%
		{長引く}{}{}%
	}%
	}%
	{\tailPrint[]}%
\entry
	{\headPrint[kana={ながびり\tl{˥˩}},ipa={nagabiri},pos={名},acc={F}]}%
	{%
		\MeaningsBlock{%
	\ryumeaning%
		{長居}{}{}%
	}%
	}%
	{\tailPrint[]}%
\entry
	{\headPrint[kana={ながみん\tl{˥}},ipa={nagamiɴ},pos={動},rui={3},para={ながむぬ},acc={H}]}%
	{%
		\MeaningsBlock{%
	\ryumeaning%
		{眺める}{}{}%
	}%
	}%
	{\tailPrint[]}%
\entry
	{\headPrint[kana={ながやどぅ\tl{˩˥}},ipa={nagajadu},pos={名},acc={R}]}%
	{%
		\MeaningsBlock{%
	\ryumeaning%
		{中戸}{}{}%
	}%
	}%
	{\tailPrint[]}%
\entry
	{\headPrint[kana={ながやみ\tl{˥˩}},ipa={nagajami},pos={名},acc={F}]}%
	{%
		\MeaningsBlock{%
	\ryumeaning%
		{長患い}{}{}%
	}%
	}%
	{\tailPrint[]}%
\entry
	{\headPrint[kana={なぎ\tl{˩˥}},ipa={nagi},pos={名},acc={R}]}%
	{%
		\MeaningsBlock{%
	\ryumeaning%
		{長さ。丈}{}{}%
	}%
	}%
	{\tailPrint[]}%
\entry
	{\headPrint[kana={なきくがりん},ipa={nakikugariɴ},pos={動}]}%
	{%
		\MeaningsBlock{%
	\ryumeaning%
		{泣き焦がれる}{}{}%
	}%
	}%
	{\tailPrint[]}%
\entry
	{\headPrint[kana={なぎすかり\tl{˥˩}},ipa={nagisu̥kari},pos={動},rui={3},para={なぎすかるぬ},acc={F}]}%
	{%
		\MeaningsBlock{%
	\ryumeaning%
		{泣きかかる。泣きつく}{}{}%
	}%
	}%
	{\tailPrint[]}%
\entry
	{\headPrint[kana={なぎとーすん\tl{˩˥}},ipa={naigitoːsuɴ},pos={動},rui={2},para={なぎとーはぬ},acc={R}]}%
	{%
		\MeaningsBlock{%
	\ryumeaning%
		{薙ぎ倒す。切り倒す}{}{}%
	}%
	}%
	{\tailPrint[]}%
\entry
	{\headPrint[kana={なぎぼたりるん\tl{˩˥}},ipa={nagibotariruɴ},pos={動},rui={3},para={なぎぼたるぬ},acc={R}]}%
	{%
		\MeaningsBlock{%
	\ryumeaning%
		{泣き疲れる。泣きくたびれる。泣きしおれる}{}{}%
	}%
	}%
	{\tailPrint[]}%
\entry
	{\headPrint[kana={なぎまーべ\tl{˥˩}},ipa={nagimaːbe},pos={名},acc={F}]}%
	{%
		\MeaningsBlock{%
	\ryumeaning%
		{泣き真似。うそ泣き}{}{}%
	}%
	}%
	{\tailPrint[]}%
\entry
	{\headPrint[kana={なぐさみん\tl{˥˩}},ipa={nagusamiɴ},pos={動},rui={3},para={なぐさむぬ},acc={F}]}%
	{%
		\MeaningsBlock{%
	\ryumeaning%
		{慰める}{}{}%
	}%
	}%
	{\tailPrint[]}%
\entry
	{\headPrint[kana={なぐりしゃ\tl{˩˥}はん},ipa={naguriʃahaɴ},pos={形},para={なぐりしゃへぬ},acc={R}]}%
	{%
		\MeaningsBlock{%
	\ryumeaning%
		{名残惜しい。悲しい}{}{}%
	}%
	}%
	{\tailPrint[]}%
\entry
	{\headPrint[kana={なげーくとぅ},ipa={nageːkutu},pos={句}]}%
	{%
		\MeaningsBlock{%
	\ryumeaning%
		{長いこと。長い間}{}{}%
	}%
	}%
	{\tailPrint[]}%
\entry
	{\headPrint[kana={なさき\tl{˩˥}},ipa={nasaki},pos={名},acc={R}]}%
	{%
		\MeaningsBlock{%
	\ryumeaning%
		{情け。思いやり}{}{}%
	}%
	}%
	{\tailPrint[]}%
\entry
	{\headPrint[kana={なざぎ\tl{˩˥}},ipa={nadzagi},pos={名},acc={R}]}%
	{%
		\MeaningsBlock{%
	\ryumeaning%
		{ハエキビ}{雑草名}{}%
	}%
	}%
	{\tailPrint[]}%
\entry
	{\headPrint[kana={なしきるん\tl{˥}},ipa={naʃikiruɴ},pos={動},acc={H}]}%
	{%
		\MeaningsBlock{%
	\ryumeaning%
		{\ruby[g]{託}{かこ}ける。口実にする}{}{}%
	}%
	}%
	{\tailPrint[]}%
\entry
	{\headPrint[kana={なしぴ\tl{˩˥}},ipa={naʃi̥pi},pos={名},acc={R}]}%
	{%
		\MeaningsBlock{%
	\ryumeaning%
		{\ruby[g]{茄子}{なすび}}{}{}%
	}%
	}%
	{\tailPrint[]}%
\entry
	{\headPrint[kana={なすぅくん\tl{˥˩}},ipa={nasu̥kuɴ},pos={動},rui={1},para={なすかぬ},acc={F}]}%
	{%
		\MeaningsBlock{%
	\ryumeaning%
		{懐く。なじむ}{}{}%
	}%
	}%
	{\tailPrint[]}%
\entry
	{\headPrint[kana={なすかっさ\tl{˩˥}はん},ipa={nasu̥kassahaɴ},pos={形},para={なすかへぬ},acc={R}]}%
	{%
		\MeaningsBlock{%
	\ryumeaning%
		{懐かしい}{}{}%
	}%
	}%
	{\tailPrint[]}%
\entry
	{\headPrint[kana={なすむら\tl{˥˩}},ipa={nasu̥mura},pos={名},acc={F}]}%
	{%
		\MeaningsBlock{%
	\ryumeaning%
		{\ruby[g]{長}{な}\ruby[g]{石}{いし}村}{波照間島の集落名}{}%
	}%
	}%
	{\tailPrint[]}%
\entry
	{\headPrint[kana={なすん\tl{˩˥}},ipa={nasuɴ},pos={動},rui={2},para={なはぬ},acc={R}]}%
	{%
		\MeaningsBlock{%
	\ryumeaning%
		{\ruby[g]{綯}{な}う}{縄を綯う}{}%
	}%
	}%
	{\tailPrint[]}%
\entry
	{\headPrint[kana={なすん\tl{˩˥}},ipa={nasuɴ},pos={動},rui={2},para={なはぬ},acc={R}]}%
	{%
		\MeaningsBlock{%
	\ryumeaning%
		{産む}{}{}%
	}%
	}%
	{\tailPrint[]}%
\entry
	{\headPrint[kana={なだ\tl{˩˥}},ipa={nada},pos={名},acc={R}]}%
	{%
		\MeaningsBlock{%
	\ryumeaning%
		{\ruby[g]{涙}{なみだ}}{普通\emb{なんだ}と言う}{}%
	}%
	}%
	{\tailPrint[]}%
\entry
	{\headPrint[kana={なだぎるん\tl{˥˩}},ipa={nadagiruɴ},pos={動},rui={3},para={なだぐぬ},acc={F}]}%
	{%
		\MeaningsBlock{%
	\ryumeaning%
		{均す。平坦にする}{}{}%
	}%
	}%
	{\tailPrint[]}%
\entry
	{\headPrint[kana={なだみるん\tl{˥˩}},ipa={nadamiruɴ},pos={動},rui={3},para={なだむぬ},acc={F}]}%
	{%
		\MeaningsBlock{%
	\ryumeaning%
		{宥める}{}{}%
	}%
	}%
	{\tailPrint[]}%
\entry
	{\headPrint[kana={なだらがすん\tl{˥˩}},ipa={nadaragasuɴ},pos={動},rui={2},para={なだらがはぬ},acc={F}]}%
	{%
		\MeaningsBlock{%
	\ryumeaning%
		{平坦にする。均す}{}{}%
	}%
	}%
	{\tailPrint[]}%
\entry
	{\headPrint[kana={なだらが\tl{˥˩}はん},ipa={nadaragahaɴ},pos={形},acc={F}]}%
	{%
		\MeaningsBlock{%
	\ryumeaning%
		{平坦だ}{}{}%
	}%
	}%
	{\tailPrint[]}%
\entry
	{\headPrint[kana={なだらぎゃーん},ipa={nadaragjaːɴ},pos={動（継）}]}%
	{%
		\MeaningsBlock{%
	\ryumeaning%
		{なだらかだ}{}{}%
	}%
	}%
	{\tailPrint[]}%
\entry
	{\headPrint[kana={なだらぎるん\tl{˥˩}},ipa={nadaragiruɴ},pos={動},rui={3},para={なだらぐぬ},acc={F}]}%
	{%
		\MeaningsBlock{%
	\ryumeaning%
		{なだらかにする。平らにならす}{}{}%
	}%
	}%
	{\tailPrint[]}%
\entry
	{\headPrint[kana={なだらぐん\tl{˥˩}},ipa={nadaraguɴ},pos={動},rui={1},para={なだらがぬ},acc={F}]}%
	{%
		\MeaningsBlock{%
	\ryumeaning%
		{なだらかになる。平坦になる}{}{}%
	\ryumeaning%
		{凪ぐ。穏やかになる}{}{}%
	}%
	}%
	{\tailPrint[]}%
\entry
	{\headPrint[kana={なちぃ\tl{˥˩}},ipa={natʃi},pos={名},acc={F}]}%
	{%
		\MeaningsBlock{%
	\ryumeaning%
		{夏。夏季}{「初夏」は\emb{ばがなちぃ}}{}%
	}%
	}%
	{\tailPrint[]}%
\entry
	{\headPrint[kana={なちぃまき\tl{˥˩}},ipa={natʃi̥maki},pos={名},acc={F}]}%
	{%
		\MeaningsBlock{%
	\ryumeaning%
		{夏負け}{}{}%
	}%
	}%
	{\tailPrint[]}%
\entry
	{\headPrint[kana={なちすぅぬ\tl{˥˩}},ipa={natʃisu̥nu},pos={名},acc={F}]}%
	{%
		\MeaningsBlock{%
	\ryumeaning%
		{夏着}{}{}%
	}%
	}%
	{\tailPrint[]}%
\entry
	{\headPrint[kana={なちぶさー\tl{˥˩}},ipa={natʃi̥busaː},pos={形},acc={F}]}%
	{%
		\MeaningsBlock{%
	\ryumeaning%
		{泣き虫だ}{}{}%
	}%
	}%
	{\tailPrint[]}%
\entry
	{\headPrint[kana={なっす\tl{˩˥}},ipa={nassu},pos={名},acc={R}]}%
	{%
		\MeaningsBlock{%
	\ryumeaning%
		{苗代}{}{}%
	}%
	}%
	{\tailPrint[]}%
\entry
	{\headPrint[kana={なっすだー},ipa={nassudaː},pos={名}]}%
	{%
		\MeaningsBlock{%
	\ryumeaning%
		{苗代田}{}{}%
	}%
	}%
	{\tailPrint[]}%
\entry
	{\headPrint[kana={ななず\tl{˩˥}},ipa={nanadzu},pos={名},acc={R}]}%
	{%
		\MeaningsBlock{%
	\ryumeaning%
		{七十}{}{}%
	}%
	}%
	{\tailPrint[]}%
\entry
	{\headPrint[kana={ななち\tl{˩˥}},ipa={nanatʃi},pos={名},acc={R}]}%
	{%
		\MeaningsBlock{%
	\ryumeaning%
		{七つ}{簡略化：なな}{}%
	}%
	}%
	{\tailPrint[]}%
\entry
	{\headPrint[kana={なば\tl{˩˥}},ipa={naba},pos={名},acc={R}]}%
	{%
		\MeaningsBlock{%
	\ryumeaning%
		{\ruby[g]{茸}{きのこ}}{}{}%
	}%
	}%
	{\tailPrint[]}%
\entry
	{\headPrint[kana={なび\tl{˩˥}},ipa={nabi},pos={名},acc={R}]}%
	{%
		\MeaningsBlock{%
	\ryumeaning%
		{\ruby[g]{鍋}{なべ}}{}{}%
	}%
	}%
	{\tailPrint[]}%
\entry
	{\headPrint[kana={なびしき\tl{˩˥}},ipa={nabiʃi̥ki},pos={名},acc={R}]}%
	{%
		\MeaningsBlock{%
	\ryumeaning%
		{おこげ}{}{}%
	}%
	}%
	{\tailPrint[]}%
\entry
	{\headPrint[kana={なびふつぁらすん\tl{˩˥}},ipa={nabifu̥tsarasuɴ},pos={動},acc={R}]}%
	{%
		\MeaningsBlock{%
	\ryumeaning%
		{鍋を焦がす}{}{}%
	}%
	}%
	{\tailPrint[]}%
\entry
	{\headPrint[kana={なびふつぁりん\tl{˩˥}},ipa={nabifu̥tsariɴ},pos={動},rui={3},para={なびふつぁるぬ},acc={R}]}%
	{%
		\MeaningsBlock{%
	\ryumeaning%
		{焦げる}{}{}%
	}%
	}%
	{\tailPrint[]}%
\entry
	{\headPrint[kana={なふこ\tl{˩˥}はん},ipa={nafu̥kohaɴ},pos={形},para={なふこはへぬ},acc={R}]}%
	{%
		\MeaningsBlock{%
	\ryumeaning%
		{なめらかだ。滑りやすい}{}{}%
	}%
	}%
	{\tailPrint[]}%
\entry
	{\headPrint[kana={なま\tl{˥}},ipa={nama},pos={名},acc={H}]}%
	{%
		\MeaningsBlock{%
	\ryumeaning%
		{今。現在}{}{}%
	}%
	}%
	{\tailPrint[]}%
\entry
	{\headPrint[kana={なま\tl{˥}},ipa={nama},pos={名},acc={H}]}%
	{%
		\MeaningsBlock{%
	\ryumeaning%
		{生}{}{}%
	}%
	}%
	{\tailPrint[]}%
\entry
	{\headPrint[kana={なまき\tl{˩˥}},ipa={namaki},pos={名},acc={R}]}%
	{%
		\MeaningsBlock{%
	\ryumeaning%
		{生木}{枯れていない薪など}{}%
	}%
	}%
	{\tailPrint[]}%
\entry
	{\headPrint[kana={なまぐみ\tl{˩˥}},ipa={namagumi},pos={名},acc={R}]}%
	{%
		\MeaningsBlock{%
	\ryumeaning%
		{生米}{半煮のご飯}{}%
	}%
	}%
	{\tailPrint[]}%
\entry
	{\headPrint[kana={なましぃ\tl{˩˥}},ipa={namaʃi},pos={名},acc={R}]}%
	{%
		\MeaningsBlock{%
	\ryumeaning%
		{\ruby[g]{刺身}{さしみ}}{}{}%
	}%
	}%
	{\tailPrint[]}%
\entry
	{\headPrint[kana={なまにー\tl{˩˥}},ipa={namaniː},pos={名},acc={R}]}%
	{%
		\MeaningsBlock{%
	\ryumeaning%
		{半煮え}{}{}%
	}%
	}%
	{\tailPrint[]}%
\entry
	{\headPrint[kana={なまふささーん},ipa={namafu̥sasaːɴ},pos={形},acc={-}]}%
	{%
		\MeaningsBlock{%
	\ryumeaning%
		{生臭い}{生物，野菜，魚，肉類の煮てない物，または，枯れていないものの臭いがする}{}%
	}%
	}%
	{\tailPrint[]}%
\entry
	{\headPrint[kana={なまふつぁ\tl{˩˥}はん},ipa={namafu̥tsahaɴ},pos={形},para={なまふつぁへぬ},acc={R}]}%
	{%
		\MeaningsBlock{%
	\ryumeaning%
		{生臭い}{}{}%
	}%
	}%
	{\tailPrint[]}%
\entry
	{\headPrint[kana={なまふつぁりむぬ\tl{˩˥}},ipa={namafu̥tsarimunu},pos={名},acc={R}]}%
	{%
		\MeaningsBlock{%
	\ryumeaning%
		{怠け者。横着な者}{}{}%
	}%
	}%
	{\tailPrint[]}%
\entry
	{\headPrint[kana={なまむぬ\tl{˩˥}},ipa={namamunu},pos={名},acc={R}]}%
	{%
		\MeaningsBlock{%
	\ryumeaning%
		{生物}{生の物}{}%
	}%
	}%
	{\tailPrint[]}%
\entry
	{\headPrint[kana={なまるん\tl{˩˥}},ipa={namaruɴ},pos={動},rui={1},para={なまらぬ},acc={R}]}%
	{%
		\MeaningsBlock{%
	\ryumeaning%
		{\ruby[g]{鈍}{なま}る。（切れ味が）鈍る}{}{}%
	}%
	}%
	{\tailPrint[]}%
\entry
	{\headPrint[kana={なもり\tl{˩˥}},ipa={namori},pos={名},acc={R}]}%
	{%
		\MeaningsBlock{%
	\ryumeaning%
		{一合枡}{}{}%
	}%
	}%
	{\tailPrint[]}%
\entry
	{\headPrint[kana={なや\tl{˥˩}},ipa={naja},pos={名},acc={F}]}%
	{%
		\MeaningsBlock{%
	\ryumeaning%
		{\ruby[g]{鰹}{かつお}工場}{「納屋」すなわち「漁具小屋」から}{}%
	}%
	}%
	{\tailPrint[]}%
\entry
	{\headPrint[kana={ならすん\tl{˥˩}},ipa={narasuɴ},pos={動},rui={2},para={ならはぬ},acc={F}]}%
	{%
		\MeaningsBlock{%
	\ryumeaning%
		{鳴らす}{}{}%
	}%
	}%
	{\tailPrint[]}%
\entry
	{\headPrint[kana={ならすん\tl{˩˥}},ipa={narasuɴ},pos={動},rui={2},para={ならはぬ},acc={R}]}%
	{%
		\MeaningsBlock{%
	\ryumeaning%
		{教える}{}{}%
	\ryumeaning%
		{習う。学ぶ}{}{}%
	}%
	}%
	{\tailPrint[]}%
\entry
	{\headPrint[kana={ならすん\tl{˩˥}},ipa={narasuɴ},pos={動},rui={2},para={ならはぬ},acc={R}]}%
	{%
		\MeaningsBlock{%
	\ryumeaning%
		{慣れる}{}{}%
	}%
	}%
	{\tailPrint[]}%
\entry
	{\headPrint[kana={ならびるん\tl{˥˩}},ipa={narabiruɴ},pos={動},rui={3},para={ならぶぬ},acc={F}]}%
	{%
		\MeaningsBlock{%
	\ryumeaning%
		{並べる}{}{}%
	}%
	}%
	{\tailPrint[]}%
\entry
	{\headPrint[kana={ならぶん\tl{˥˩}},ipa={narabuɴ},pos={動},rui={1},para={ならばぬ},acc={F}]}%
	{%
		\MeaningsBlock{%
	\ryumeaning%
		{並ぶ。整列する}{}{}%
	}%
	}%
	{\tailPrint[]}%
\entry
	{\headPrint[kana={な\josho{}り\kako{}さ},ipa={narisa},pos={名},acc={特殊}]}%
	{%
		\MeaningsBlock{%
	\ryumeaning%
		{浜砂利}{サンゴの砕けた砂利}{}%
	}%
	}%
	{\tailPrint[]}%
\entry
	{\headPrint[kana={なれー\tl{˩˥}},ipa={nareː},pos={名},acc={R}]}%
	{%
		\MeaningsBlock{%
	\ryumeaning%
		{ならい。習慣}{}{}%
	}%
	}%
	{\tailPrint[]}%
\entry
	{\headPrint[kana={なん\tl{˥˩}},ipa={naɴ},pos={名},acc={F}]}%
	{%
		\MeaningsBlock{%
	\ryumeaning%
		{名。名前}{}{}%
	}%
	}%
	{\tailPrint[]}%
\entry
	{\headPrint[kana={なん\tl{˥˩}},ipa={naɴ},pos={名},acc={F}]}%
	{%
		\MeaningsBlock{%
	\ryumeaning%
		{波。波浪}{}{}%
	}%
	}%
	{\tailPrint[]}%
\entry
	{\headPrint[kana={なん\tl{˩˥}},ipa={naɴ},pos={名},acc={R}]}%
	{%
		\MeaningsBlock{%
	\ryumeaning%
		{釣り糸}{}{}%
	}%
	}%
	{\tailPrint[]}%
\entry
	{\headPrint[kana={なん\tl{˩˥}},ipa={naɴ},pos={名},acc={R}]}%
	{%
		\MeaningsBlock{%
	\ryumeaning%
		{菜}{}{}%
	}%
	}%
	{\tailPrint[]}%
\entry
	{\headPrint[kana={なん\tl{˥˩}\ws{}えぬん\tl{˥˩}},ipa={naɴ enuɴ},pos={句},acc={F F}]}%
	{%
		\MeaningsBlock{%
	\ryumeaning%
		{名乗る}{名前を告げる}{}%
	}%
	}%
	{\tailPrint[]}%
\entry
	{\headPrint[kana={なんが\tl{˩˥}},ipa={naŋga},pos={名},acc={R}]}%
	{%
		\MeaningsBlock{%
	\ryumeaning%
		{七日。一周忌}{}{}%
	}%
	}%
	{\tailPrint[]}%
\entry
	{\headPrint[kana={なんが\tl{˩˥}\ws{}ししるん\tl{˥˩}},ipa={naŋga ʃʃiruɴ},pos={句},acc={R F}]}%
	{%
		\MeaningsBlock{%
	\ryumeaning%
		{投げ捨てる}{}{}%
	}%
	}%
	{\tailPrint[]}%
\entry
	{\headPrint[kana={なんがそーりん\tl{˩˥}},ipa={naŋgasoːriɴ},pos={名},acc={R}]}%
	{%
		\MeaningsBlock{%
	\ryumeaning%
		{旧暦の七夕}{}{}%
	}%
	}%
	{\tailPrint[]}%
\entry
	{\headPrint[kana={なんが\tl{˩˥}\ws{}っしん\tl{˩}},ipa={naŋga ʃʃiɴ},pos={句},para={なんがすとぬ},acc={R L}]}%
	{%
		\MeaningsBlock{%
	\ryumeaning%
		{捨てる。投げ捨てる}{}{}%
	}%
	}%
	{\tailPrint[]}%
\entry
	{\headPrint[kana={なんぎ\tl{˩˥}},ipa={naŋgi},pos={名},acc={R}]}%
	{%
		\MeaningsBlock{%
	\ryumeaning%
		{難儀}{}{}%
	}%
	}%
	{\tailPrint[]}%
\entry
	{\headPrint[kana={なんぎるん\tl{˩˥}},ipa={naŋgiruɴ},pos={動},rui={3},para={なんぐぬ},acc={R}]}%
	{%
		\MeaningsBlock{%
	\ryumeaning%
		{投げる}{}{}%
	\ryumeaning%
		{叩きつける}{}{}%
	}%
	}%
	{\tailPrint[]}%
\entry
	{\headPrint[kana={なんざ\tl{˩˥}},ipa={nandza},pos={名},acc={R}]}%
	{%
		\MeaningsBlock{%
	\ryumeaning%
		{銀}{}{}%
	}%
	}%
	{\tailPrint[]}%
\entry
	{\headPrint[kana={なんだ\tl{˩˥}},ipa={nanda},pos={名},acc={R}]}%
	{%
		\MeaningsBlock{%
	\ryumeaning%
		{\ruby[g]{涙}{なみだ}}{}{}%
	}%
	}%
	{\tailPrint[]}%
\entry
	{\headPrint[kana={なんだら\tl{˩˥}},ipa={nandara},pos={名},acc={R}]}%
	{%
		\MeaningsBlock{%
	\ryumeaning%
		{ビーチロック}{浜に出来る砂岩}{}%
	}%
	}%
	{\tailPrint[]}%
\entry
	{\headPrint[kana={なんだ\tl{˩˥}\ws{}んじるん\tl{˥˩}},ipa={nanda ndʒiruɴ},pos={句},acc={R H}]}%
	{%
		\MeaningsBlock{%
	\ryumeaning%
		{涙が出る。涙ぐむ}{}{}%
	}%
	}%
	{\tailPrint[]}%
\entry
	{\headPrint[kana={なんぶらすん\tl{˩˥}},ipa={namburasuɴ},pos={動},rui={2},para={なんぶらはぬ},acc={R}]}%
	{%
		\MeaningsBlock{%
	\ryumeaning%
		{流す}{}{}%
	}%
	}%
	{\tailPrint[]}%
\entry
	{\headPrint[kana={なんぶらん},ipa={namburaɴ},pos={名}]}%
	{%
		\MeaningsBlock{%
	\ryumeaning%
		{荒海}{}{}%
	}%
	}%
	{\tailPrint[]}%
\entry
	{\headPrint[kana={なんぶりるん},ipa={namburirun},pos={動}]}%
	{%
		\MeaningsBlock{%
	\ryumeaning%
		{漂わす}{}{}%
	}%
	}%
	{\tailPrint[]}%
\entry
	{\headPrint[kana={なんぶりん\tl{˩˥}},ipa={namburiɴ},pos={動},rui={3},para={なんぶるぬ},acc={R}]}%
	{%
		\MeaningsBlock{%
	\ryumeaning%
		{流れ出す。流失する}{}{}%
	}%
	}%
	{\tailPrint[]}%
\LetterHead{に}
\entry
	{\headPrint[kana={にー\tl{˥}},ipa={niː},pos={名},acc={H}]}%
	{%
		\MeaningsBlock{%
	\ryumeaning%
		{\ruby[g]{子}{ね}}{十二支の子}{}%
	}%
	}%
	{\tailPrint[]}%
\entry
	{\headPrint[kana={にーはい},ipa={niːhai},pos={感}]}%
	{%
		\MeaningsBlock{%
	\ryumeaning%
		{ありがとう}{}{}%
	}%
	}%
	{\tailPrint[]}%
\entry
	{\headPrint[kana={にーはいゆー},ipa={niːhaijuː},pos={感}]}%
	{%
		\MeaningsBlock{%
	\ryumeaning%
		{ありがとうございます}{敬語}{}%
	}%
	}%
	{\tailPrint[]}%
\entry
	{\headPrint[kana={にーびち\tl{˩˥}},ipa={niːbitʃi},pos={名},acc={R}]}%
	{%
		\MeaningsBlock{%
	\ryumeaning%
		{結婚式}{}{}%
	}%
	}%
	{\tailPrint[]}%
\entry
	{\headPrint[kana={にーぶた\tl{˩˥}},ipa={niːbuta},pos={名},acc={R}]}%
	{%
		\MeaningsBlock{%
	\ryumeaning%
		{腫れ物。吹き出物}{}{}%
	}%
	}%
	{\tailPrint[]}%
\entry
	{\headPrint[kana={にーぶやー\tl{˥˩}},ipa={niːbujaː},pos={名},acc={F}]}%
	{%
		\MeaningsBlock{%
	\ryumeaning%
		{寝坊助。寝坊}{}{}%
	}%
	}%
	{\tailPrint[]}%
\entry
	{\headPrint[kana={にーまらん\tl{˥˩}},ipa={niːmaraɴ},pos={動},acc={F}]}%
	{%
		\MeaningsBlock{%
	\ryumeaning%
		{（食物が）腐る。腐りかかる}{}{}%
	}%
	}%
	{\tailPrint[]}%
\entry
	{\headPrint[kana={にーらすん},ipa={niːrasuɴ},pos={動}]}%
	{%
		\MeaningsBlock{%
	\ryumeaning%
		{似せる}{}{}%
	}%
	}%
	{\tailPrint[]}%
\entry
	{\headPrint[kana={にが\tl{˩˥}},ipa={niga},pos={名},acc={R}]}%
	{%
		\MeaningsBlock{%
	\ryumeaning%
		{今晩}{今夜}{}%
	}%
	}%
	{\tailPrint[]}%
\entry
	{\headPrint[kana={にく\tl{˩˥}},ipa={niku},pos={名},acc={R}]}%
	{%
		\MeaningsBlock{%
	\ryumeaning%
		{\ruby[g]{肉}{にく}}{}{}%
	}%
	}%
	{\tailPrint[]}%
\entry
	{\headPrint[kana={にげー\tl{˩˥}},ipa={nigeː},pos={名},acc={R}]}%
	{%
		\MeaningsBlock{%
	\ryumeaning%
		{願い。祈願}{}{}%
	}%
	}%
	{\tailPrint[]}%
\entry
	{\headPrint[kana={にげー\tl{˩˥}しきるん\tl{˥}},ipa={nigeːʃi̥kiruɴ},pos={動},rui={3},para={にげーすくぬ},acc={R+H}]}%
	{%
		\MeaningsBlock{%
	\ryumeaning%
		{祈る}{}{}%
	}%
	}%
	{\tailPrint[]}%
\entry
	{\headPrint[kana={にげーふちぃ\tl{˩˥}},ipa={nigeːfu̥tʃi},pos={名},acc={R}]}%
	{%
		\MeaningsBlock{%
	\ryumeaning%
		{祈願の言葉}{}{}%
	}%
	}%
	{\tailPrint[]}%
\entry
	{\headPrint[kana={にげー\tl{˩˥}\ws{}んじるん\tl{˥}},ipa={nigeː ndʒiruɴ},pos={句},rui={3},para={にげーんどぅぬ},acc={R+H}]}%
	{%
		\MeaningsBlock{%
	\ryumeaning%
		{願い出る。申し入れる}{}{}%
	}%
	}%
	{\tailPrint[]}%
\entry
	{\headPrint[kana={にげるん\tl{˩˥}},ipa={nigeruɴ},pos={動},rui={1},para={にげーらぬ},acc={R}]}%
	{%
		\MeaningsBlock{%
	\ryumeaning%
		{祈願する。祈る}{}{}%
	}%
	}%
	{\tailPrint[]}%
\entry
	{\headPrint[kana={にし\tl{˥˩}},ipa={niʃi},pos={名},acc={F}]}%
	{%
		\MeaningsBlock{%
	\ryumeaning%
		{北。北方}{}{}%
	}%
	}%
	{\tailPrint[]}%
\entry
	{\headPrint[kana={にしかち\tl{˥˩}},ipa={niʃikḁtʃi},pos={名},acc={F}]}%
	{%
		\MeaningsBlock{%
	\ryumeaning%
		{北風}{}{}%
	}%
	}%
	{\tailPrint[]}%
\entry
	{\headPrint[kana={に\josho{}じ\kako{}きょー},ipa={nidʒikjoː},pos={名},acc={特殊}]}%
	{%
		\MeaningsBlock{%
	\ryumeaning%
		{ウイキョウ}{ハーブの一種}{}%
	}%
	}%
	{\tailPrint[]}%
\entry
	{\headPrint[kana={にししゃ\tl{˥˩}はん},ipa={niʃiʃahaɴ},pos={形},para={にししゃへぬ},acc={F}]}%
	{%
		\MeaningsBlock{%
	\ryumeaning%
		{似る。似ている}{}{}%
	}%
	}%
	{\tailPrint[]}%
\entry
	{\headPrint[kana={にした\tl{˥˩}},ipa={niʃi̥ta},pos={名},acc={F}]}%
	{%
		\MeaningsBlock{%
	\ryumeaning%
		{北方。北側}{}{}%
	}%
	}%
	{\tailPrint[]}%
\entry
	{\headPrint[kana={にしななち\tl{˩˥}},ipa={niʃi̥nanatʃi},pos={名},acc={R}]}%
	{%
		\MeaningsBlock{%
	\ryumeaning%
		{北斗星。北斗七星}{}{}%
	}%
	}%
	{\tailPrint[]}%
\entry
	{\headPrint[kana={にしむら\tl{˥˩}},ipa={niʃi̥mura},pos={名},acc={F}]}%
	{%
		\MeaningsBlock{%
	\ryumeaning%
		{北村}{集落名}{}%
	}%
	}%
	{\tailPrint[]}%
\entry
	{\headPrint[kana={にしゃ\tl{˥˩}はん},ipa={niʃahaɴ},pos={形},acc={F}]}%
	{%
		\MeaningsBlock{%
	\ryumeaning%
		{不味い。旨くない}{}{}%
	}%
	}%
	{\tailPrint[]}%
\entry
	{\headPrint[kana={にた\tl{˩˥}はん},ipa={nitahaɴ},pos={形},para={にたへぬ},acc={R}]}%
	{%
		\MeaningsBlock{%
	\ryumeaning%
		{憎い。妬ましい}{}{}%
	}%
	}%
	{\tailPrint[]}%
\entry
	{\headPrint[kana={にたむん\tl{˩˥}},ipa={nitamuɴ},pos={動},rui={1},para={にたまぬ},acc={R}]}%
	{%
		\MeaningsBlock{%
	\ryumeaning%
		{\ruby[g]{妬}{ねた}む}{}{}%
	\ryumeaning%
		{恨む}{}{}%
	}%
	}%
	{\tailPrint[]}%
\entry
	{\headPrint[kana={にちぃ\tl{˥˩}},ipa={nitʃi},pos={名},acc={F}]}%
	{%
		\MeaningsBlock{%
	\ryumeaning%
		{\ruby[g]{胸}{むね}}{}{}%
	}%
	}%
	{\tailPrint[]}%
\entry
	{\headPrint[kana={にちぃ\tl{˩˥}},ipa={nitʃi},pos={名},acc={R}]}%
	{%
		\MeaningsBlock{%
	\ryumeaning%
		{熱。体温}{}{}%
	}%
	}%
	{\tailPrint[]}%
\entry
	{\headPrint[kana={にち\tl{˩˥}\ws{}んじるん\tl{˥}},ipa={nitʃi ndʒiruɴ},pos={句},para={にちんどうぬ},acc={R H}]}%
	{%
		\MeaningsBlock{%
	\ryumeaning%
		{熱が出る。発熱する}{}{}%
	}%
	}%
	{\tailPrint[]}%
\entry
	{\headPrint[kana={にばり\tl{˩˥}},ipa={nibari},pos={名},acc={R}]}%
	{%
		\MeaningsBlock{%
	\ryumeaning%
		{ハタ}{ハタ類の総称。魚類名}{}%
	}%
	}%
	{\tailPrint[]}%
\entry
	{\headPrint[kana={にぶ\tl{˩˥}},ipa={nibu},pos={名},acc={R}]}%
	{%
		\MeaningsBlock{%
	\ryumeaning%
		{\ruby[g]{柄杓}{ひしゃく}}{}{}%
	}%
	}%
	{\tailPrint[]}%
\entry
	{\headPrint[kana={にふつぁ\tl{˥˩}はん},ipa={nifu̥tsahaɴ},pos={形},para={にふつぁへぬ},acc={F}]}%
	{%
		\MeaningsBlock{%
	\ryumeaning%
		{遅い。遅れる}{}{}%
	}%
	}%
	{\tailPrint[]}%
\entry
	{\headPrint[kana={にむちぃ\tl{˩˥}},ipa={nimutʃi},pos={名},acc={R}]}%
	{%
		\MeaningsBlock{%
	\ryumeaning%
		{荷物}{}{}%
	}%
	}%
	{\tailPrint[]}%
\entry
	{\headPrint[kana={にん\tl{˥˩}},ipa={niɴ},pos={名},acc={F}]}%
	{%
		\MeaningsBlock{%
	\ryumeaning%
		{\ruby[g]{年}{ねん}}{}{}%
	}%
	}%
	{\tailPrint[]}%
\entry
	{\headPrint[kana={にん\tl{˩˥}},ipa={niɴ},pos={名},acc={R}]}%
	{%
		\MeaningsBlock{%
	\ryumeaning%
		{\ruby[g]{根}{ね}}{植物の根}{}%
	}%
	}%
	{\tailPrint[]}%
\entry
	{\headPrint[kana={にん\tl{˩˥}},ipa={niɴ},pos={名},acc={R}]}%
	{%
		\MeaningsBlock{%
	\ryumeaning%
		{地震}{}{}%
	}%
	}%
	{\tailPrint[]}%
\entry
	{\headPrint[kana={にんうりん},ipa={nin.uriɴ},pos={句}]}%
	{%
		\MeaningsBlock{%
	\ryumeaning%
		{根付く}{}{}%
	}%
	}%
	{\tailPrint[]}%
\entry
	{\headPrint[kana={にんぎん\tl{˩˥}},ipa={niŋgiɴ},pos={名},acc={R}]}%
	{%
		\MeaningsBlock{%
	\ryumeaning%
		{人。人間}{}{}%
	}%
	}%
	{\tailPrint[]}%
\entry
	{\headPrint[kana={にんぐ\tl{˩˥}},ipa={niŋgu},pos={名},acc={R}]}%
	{%
		\MeaningsBlock{%
	\ryumeaning%
		{年貢}{}{}%
	}%
	}%
	{\tailPrint[]}%
\entry
	{\headPrint[kana={にんじるん\tl{˥˩}},ipa={nindʒiruɴ},pos={動},rui={3},para={にんずぬ},acc={F}]}%
	{%
		\MeaningsBlock{%
	\ryumeaning%
		{念ずる。祈る}{}{}%
	}%
	}%
	{\tailPrint[]}%
\entry
	{\headPrint[kana={にんず},ipa={nindzu},pos={名},acc={-}]}%
	{%
		\MeaningsBlock{%
	\ryumeaning%
		{二十}{年齢の「二十歳」は\emb{ぱたち}という}{}%
	}%
	}%
	{\tailPrint[]}%
\entry
	{\headPrint[kana={にんずー\tl{˩˥}},ipa={nindzuː},pos={名},acc={R}]}%
	{%
		\MeaningsBlock{%
	\ryumeaning%
		{人数}{}{}%
	}%
	}%
	{\tailPrint[]}%
\entry
	{\headPrint[kana={にん\tl{˩˥}すくん\tl{˥}},ipa={ninsu̥kuɴ},pos={動},rui={1},para={にんすかぬ},acc={R+H}]}%
	{%
		\MeaningsBlock{%
	\ryumeaning%
		{根付く}{}{}%
	}%
	}%
	{\tailPrint[]}%
\entry
	{\headPrint[kana={にんどーしぃ\tl{˩˥}},ipa={nindoːʃi},pos={名},acc={R}]}%
	{%
		\MeaningsBlock{%
	\ryumeaning%
		{ハイキビ}{雑草名}{}%
	}%
	}%
	{\tailPrint[]}%
\entry
	{\headPrint[kana={にんにん},ipa={ninniɴ},pos={名}]}%
	{%
		\MeaningsBlock{%
	\ryumeaning%
		{年々。毎年}{}{}%
	}%
	}%
	{\tailPrint[]}%
\entry
	{\headPrint[kana={にんぬ\ws{}んさはん},ipa={ninnu nsahaɴ},pos={句}]}%
	{%
		\MeaningsBlock{%
	\ryumeaning%
		{荷が重い}{}{}%
	}%
	}%
	{\tailPrint[]}%
\entry
	{\headPrint[kana={にんばぎるん\tl{˩˥}},ipa={nimbagiruɴ},pos={句},para={にんばぐぬ},acc={R}]}%
	{%
		\MeaningsBlock{%
	\ryumeaning%
		{根分けする。株分けする}{}{}%
	}%
	}%
	{\tailPrint[]}%
\entry
	{\headPrint[kana={にんばぷち\tl{˩˥}},ipa={nimbapu̥tʃi},pos={名},acc={R}]}%
	{%
		\MeaningsBlock{%
	\ryumeaning%
		{北極星}{}{}%
	}%
	}%
	{\tailPrint[]}%
\entry
	{\headPrint[kana={にんぶち\tl{˩˥}},ipa={nimbutʃi},pos={名},acc={R}]}%
	{%
		\MeaningsBlock{%
	\ryumeaning%
		{念仏者。司祭}{}{}%
	}%
	}%
	{\tailPrint[]}%
\entry
	{\headPrint[kana={にんぶり\tl{˥˩}},ipa={nimburi},pos={名},acc={F}]}%
	{%
		\MeaningsBlock{%
	\ryumeaning%
		{居眠り}{}{}%
	}%
	}%
	{\tailPrint[]}%
\entry
	{\headPrint[kana={にんまるん\tl{˥˩}},ipa={nimmaruɴ},pos={動},para={にんまるぬ},acc={F}]}%
	{%
		\MeaningsBlock{%
	\ryumeaning%
		{腐りかかる}{腐敗にまではいたっていない}{}%
	}%
	}%
	{\tailPrint[]}%
\entry
	{\headPrint[kana={にんむとぅ\tl{˩˥}},ipa={nimmutu},pos={名},acc={R}]}%
	{%
		\MeaningsBlock{%
	\ryumeaning%
		{根元}{}{}%
	}%
	}%
	{\tailPrint[]}%
\LetterHead{ぬ}
\entry
	{\headPrint[kana={ぬー\tl{˩˥}},ipa={nuː},pos={名},acc={R}]}%
	{%
		\MeaningsBlock{%
	\ryumeaning%
		{野。野原}{}{}%
	}%
	}%
	{\tailPrint[]}%
\entry
	{\headPrint[kana={ぬー\tl{˩˥}},ipa={nuː},pos={名},acc={R}]}%
	{%
		\MeaningsBlock{%
	\ryumeaning%
		{何}{}{}%
	}%
	}%
	{\tailPrint[]}%
\entry
	{\headPrint[kana={ぬーぐとぅ},ipa={nuːgutu},pos={名}]}%
	{%
		\MeaningsBlock{%
	\ryumeaning%
		{何事}{}{}%
	}%
	}%
	{\tailPrint[]}%
\entry
	{\headPrint[kana={ぬ\josho{}ー\kako{}しゃる},ipa={nuːʃaru},pos={句},acc={LHL}]}%
	{%
		\MeaningsBlock{%
	\ryumeaning%
		{どんな。どのような}{}{}%
	}%
	}%
	{\tailPrint[]}%
\entry
	{\headPrint[kana={ぬーしん},ipa={nuːʃiɴ},pos={副}]}%
	{%
		\MeaningsBlock{%
	\ryumeaning%
		{どうしても}{}{}%
	}%
	}%
	{\tailPrint[]}%
\entry
	{\headPrint[kana={ぬ\josho{}ー\kako{}た},ipa={nuːta},pos={句},acc={LHL}]}%
	{%
		\MeaningsBlock{%
	\ryumeaning%
		{なんと。何と言ってるか}{}{}%
	}%
	}%
	{\tailPrint[]}%
\entry
	{\headPrint[kana={ぬーたる},ipa={nuːtaru},pos={句}]}%
	{%
		\MeaningsBlock{%
	\ryumeaning%
		{なんと。何と言ってるか}{}{}%
	}%
	}%
	{\tailPrint[]}%
\entry
	{\headPrint[kana={ぬーとぅ},ipa={nuːtu},pos={名}]}%
	{%
		\MeaningsBlock{%
	\ryumeaning%
		{謎々}{「なんだろう」の意味}{}%
	}%
	}%
	{\tailPrint[]}%
\entry
	{\headPrint[kana={ぬーび\tl{˩˥}},ipa={nuːbi},pos={名},acc={R}]}%
	{%
		\MeaningsBlock{%
	\ryumeaning%
		{背伸び}{}{}%
	}%
	}%
	{\tailPrint[]}%
\entry
	{\headPrint[kana={ぬ\josho{}ー\kako{}や},ipa={nuːja},pos={句},acc={LHL}]}%
	{%
		\MeaningsBlock{%
	\ryumeaning%
		{何が}{}{}%
	}%
	}%
	{\tailPrint[]}%
\entry
	{\headPrint[kana={ぬ\josho{}ー\ws{}やばん\tl{˩}},ipa={nuː jabaɴ},pos={句},acc={LH L}]}%
	{%
		\MeaningsBlock{%
	\ryumeaning%
		{なんでも}{}{}%
	}%
	}%
	{\tailPrint[]}%
\entry
	{\headPrint[kana={ぬ\josho{}ー\kako{}ん},ipa={nuːɴ},pos={副},acc={LHL}]}%
	{%
		\MeaningsBlock{%
	\ryumeaning%
		{全く。何も}{}{}%
	}%
	}%
	{\tailPrint[]}%
\entry
	{\headPrint[kana={ぬーん\tl{˩˥}},ipa={nuːɴ},pos={動},rui={1},para={ぬわぬ},acc={R}]}%
	{%
		\MeaningsBlock{%
	\ryumeaning%
		{縫う。裁縫}{}{}%
	}%
	}%
	{\tailPrint[]}%
\entry
	{\headPrint[kana={ぬいむぬ\tl{˩˥}},ipa={nuimunu},pos={名},acc={R}]}%
	{%
		\MeaningsBlock{%
	\ryumeaning%
		{縫い物。裁縫}{}{}%
	}%
	}%
	{\tailPrint[]}%
\entry
	{\headPrint[kana={ぬが\tl{˩˥}},ipa={nuga},pos={名},acc={R}]}%
	{%
		\MeaningsBlock{%
	\ryumeaning%
		{\ruby[g]{糠}{ぬか}。\ruby[g]{米}{こめ}\ruby[g]{糠}{ぬか}}{}{}%
	}%
	}%
	{\tailPrint[]}%
\entry
	{\headPrint[kana={ぬがーらすん},ipa={nugaːrasuɴ},pos={動},acc={R}]}%
	{%
		\MeaningsBlock{%
	\ryumeaning%
		{許す。免れさせる}{}{}%
	}%
	}%
	{\tailPrint[]}%
\entry
	{\headPrint[kana={ぬがーるん\tl{˩˥}},ipa={nugaːruɴ},pos={動},rui={1},acc={R}]}%
	{%
		\MeaningsBlock{%
	\ryumeaning%
		{許される。免れる}{}{}%
	\ryumeaning%
		{卒業する}{義務教育を「免れた」から}{}%
	}%
	}%
	{\tailPrint[]}%
\entry
	{\headPrint[kana={ぬがりゃん\tl{˩˥}},ipa={nugarjaɴ},pos={動（継）},para={ぬがるぬ},acc={R}]}%
	{%
		\MeaningsBlock{%
	\ryumeaning%
		{脱した}{}{}%
	}%
	}%
	{\tailPrint[]}%
\entry
	{\headPrint[kana={ぬきひー},ipa={nukihiː},pos={名}]}%
	{%
		\MeaningsBlock{%
	\ryumeaning%
		{貫き家。本建築}{掘っ立て小屋に対する}{}%
	}%
	}%
	{\tailPrint[]}%
\entry
	{\headPrint[kana={ぬぎり\tl{˩˥}},ipa={nugiri},pos={名},acc={R}]}%
	{%
		\MeaningsBlock{%
	\ryumeaning%
		{\ruby[g]{鋸}{のこぎり}}{}{}%
	}%
	}%
	{\tailPrint[]}%
\entry
	{\headPrint[kana={ぬぎるん\tl{˩˥}},ipa={nugiruɴ},pos={動},rui={3},para={ぬぐぬ},acc={R}]}%
	{%
		\MeaningsBlock{%
	\ryumeaning%
		{抜ける。外れる}{}{}%
	}%
	}%
	{\tailPrint[]}%
\entry
	{\headPrint[kana={ぬぎるん\tl{˩˥}},ipa={nugiruɴ},pos={動},acc={R}]}%
	{%
		\MeaningsBlock{%
	\ryumeaning%
		{退く。遠のく}{}{}%
	\ryumeaning%
		{離れる}{}{}%
	}%
	}%
	{\tailPrint[]}%
\entry
	{\headPrint[kana={ぬぎんだん},ipa={nugindaɴ},pos={動},rui={3},para={ぬぎんどぅぬ}]}%
	{%
		\MeaningsBlock{%
	\ryumeaning%
		{抜きんでる}{}{}%
	}%
	}%
	{\tailPrint[]}%
\entry
	{\headPrint[kana={ぬぐじま\tl{˥˩}},ipa={nugudʒima},pos={名},acc={F}]}%
	{%
		\MeaningsBlock{%
	\ryumeaning%
		{山のない平坦な島}{}{}%
	}%
	}%
	{\tailPrint[]}%
\entry
	{\headPrint[kana={ぬぐり\tl{˩˥}},ipa={nuguri},pos={名},acc={R}]}%
	{%
		\MeaningsBlock{%
	\ryumeaning%
		{残り。残部}{}{}%
	}%
	}%
	{\tailPrint[]}%
\entry
	{\headPrint[kana={ぬぐりしゃ\tl{˩˥}はん},ipa={nuguriʃahaɴ},pos={形},para={ぬぐりしゃへぬ},acc={R}]}%
	{%
		\MeaningsBlock{%
	\ryumeaning%
		{恐ろしい。危険だ}{}{}%
	}%
	}%
	{\tailPrint[]}%
\entry
	{\headPrint[kana={ぬぐん\tl{˩˥}},ipa={nuguɴ},pos={動},rui={1},para={ぬがぬ},acc={R}]}%
	{%
		\MeaningsBlock{%
	\ryumeaning%
		{抜く。貫く。脱ぐ}{}{}%
	\ryumeaning%
		{刺す。突き刺す}{}{}%
	}%
	}%
	{\tailPrint[]}%
\entry
	{\headPrint[kana={ぬしん\tl{˥˩}},ipa={nuʃiɴ},pos={動},rui={3},para={ぬすぬ},acc={F}]}%
	{%
		\MeaningsBlock{%
	\ryumeaning%
		{乗せる。積む}{}{}%
	}%
	}%
	{\tailPrint[]}%
\entry
	{\headPrint[kana={ぬすかるん\tl{˥˩}},ipa={nusu̥karuɴ},pos={動},rui={1},para={ぬすからぬ},acc={F}]}%
	{%
		\MeaningsBlock{%
	\ryumeaning%
		{近づく。（間近に、目前に）せまる}{}{}%
	}%
	}%
	{\tailPrint[]}%
\entry
	{\headPrint[kana={ぬすとぅり\tl{˩˥}},ipa={nusu̥turi},pos={名},acc={R}]}%
	{%
		\MeaningsBlock{%
	\ryumeaning%
		{泥棒。盗人}{「盗み取り」の略}{}%
	}%
	}%
	{\tailPrint[]}%
\entry
	{\headPrint[kana={ぬずみ\tl{˥˩}},ipa={nudzumi},pos={名},acc={F}]}%
	{%
		\MeaningsBlock{%
	\ryumeaning%
		{望み。希望}{}{}%
	}%
	}%
	{\tailPrint[]}%
\entry
	{\headPrint[kana={ぬすむん\tl{˩˥}},ipa={nusu̥muɴ},pos={動},rui={1},para={ぬすまぬ},acc={R}]}%
	{%
		\MeaningsBlock{%
	\ryumeaning%
		{盗む}{}{}%
	}%
	}%
	{\tailPrint[]}%
\entry
	{\headPrint[kana={ぬずむん},ipa={nudzumuɴ},pos={動},rui={1},para={ぬずまぬ}]}%
	{%
		\MeaningsBlock{%
	\ryumeaning%
		{望む。所望する。欲する}{}{}%
	}%
	}%
	{\tailPrint[]}%
\entry
	{\headPrint[kana={ぬちぃ\tl{˩˥}},ipa={nutʃi},pos={名},acc={R}]}%
	{%
		\MeaningsBlock{%
	\ryumeaning%
		{命。生命}{}{}%
	}%
	}%
	{\tailPrint[]}%
\entry
	{\headPrint[kana={ぬちぃ\ws{}しさん},ipa={nutʃi ssaɴ},pos={句}]}%
	{%
		\MeaningsBlock{%
	\ryumeaning%
		{こと切れる。死亡する}{}{}%
	}%
	}%
	{\tailPrint[]}%
\entry
	{\headPrint[kana={ぬちぃずー\tl{˩˥}さーん},ipa={nutʃidzuːsaːɴ},pos={形},acc={R}]}%
	{%
		\MeaningsBlock{%
	\ryumeaning%
		{命強い}{生命力が強い。命が永い}{}%
	}%
	}%
	{\tailPrint[]}%
\entry
	{\headPrint[kana={ぬちぃぬうや\tl{˩˥}},ipa={nutʃinu.uja},pos={名},acc={R}]}%
	{%
		\MeaningsBlock{%
	\ryumeaning%
		{命の恩人}{}{}%
	}%
	}%
	{\tailPrint[]}%
\entry
	{\headPrint[kana={ぬちぃまろ\tl{˩˥}はん},ipa={nutʃimarohaɴ},pos={形},para={ぬちまろへぬ},acc={R}]}%
	{%
		\MeaningsBlock{%
	\ryumeaning%
		{短命。薄命}{}{}%
	}%
	}%
	{\tailPrint[]}%
\entry
	{\headPrint[kana={ぬちぃむやん},ipa={nutʃimujaɴ},pos={句},para={ぬちむぬ}]}%
	{%
		\MeaningsBlock{%
	\ryumeaning%
		{助かる。生き返る}{}{}%
	}%
	}%
	{\tailPrint[]}%
\entry
	{\headPrint[kana={ぬっふぇーるん\tl{˥˩}},ipa={nuffeːruɴ},pos={動},acc={F}]}%
	{%
		\MeaningsBlock{%
	\ryumeaning%
		{\ruby[g]{舐}{な}める}{}{}%
	}%
	}%
	{\tailPrint[]}%
\entry
	{\headPrint[kana={ぬっふたりるん},ipa={nuffutariruɴ},pos={動}]}%
	{%
		\MeaningsBlock{%
	\ryumeaning%
		{寝入る。よく眠る}{}{}%
	}%
	}%
	{\tailPrint[]}%
\entry
	{\headPrint[kana={ぬっふん\tl{˥˩}},ipa={nuffuɴ},pos={動},rui={1},para={ぬっふわぬ},acc={F}]}%
	{%
		\MeaningsBlock{%
	\ryumeaning%
		{寝る。眠る}{}{}%
	}%
	}%
	{\tailPrint[]}%
\entry
	{\headPrint[kana={ぬどぅ\tl{˩˥}},ipa={nudu},pos={名},acc={R}]}%
	{%
		\MeaningsBlock{%
	\ryumeaning%
		{\ruby[g]{喉}{のど}}{}{}%
	}%
	}%
	{\tailPrint[]}%
\entry
	{\headPrint[kana={ぬどぅ\tl{˩˥}\ws{}かりるん\tl{˥˩}},ipa={nudu kḁriruɴ},pos={句},acc={R F}]}%
	{%
		\MeaningsBlock{%
	\ryumeaning%
		{喉が渇く}{}{}%
	}%
	}%
	{\tailPrint[]}%
\entry
	{\headPrint[kana={ぬぬ\tl{˥˩}},ipa={nunu},pos={名},acc={F}]}%
	{%
		\MeaningsBlock{%
	\ryumeaning%
		{布。織物}{}{}%
	}%
	}%
	{\tailPrint[]}%
\entry
	{\headPrint[kana={ぬのさらし\tl{˥˩}},ipa={nunosaraʃi},pos={名},acc={F}]}%
	{%
		\MeaningsBlock{%
	\ryumeaning%
		{布晒し}{}{}%
	}%
	}%
	{\tailPrint[]}%
\entry
	{\headPrint[kana={ぬばすん\tl{˩˥}},ipa={nubasuɴ},pos={動},rui={2},para={ぬばはぬ},acc={R}]}%
	{%
		\MeaningsBlock{%
	\ryumeaning%
		{延ばす。延期する}{}{}%
	}%
	}%
	{\tailPrint[]}%
\entry
	{\headPrint[kana={ぬびるん\tl{˩˥}},ipa={nubiruɴ},pos={動},rui={2},para={ぬばはぬ},acc={R}]}%
	{%
		\MeaningsBlock{%
	\ryumeaning%
		{延ばす。伸べる}{}{}%
	}%
	}%
	{\tailPrint[]}%
\entry
	{\headPrint[kana={ぬぶしるん\tl{˥˩}},ipa={nubuʃiɴ},pos={動},rui={3},para={ぬぶすぬ},acc={F}]}%
	{%
		\MeaningsBlock{%
	\ryumeaning%
		{\ruby[g]{逆上}{のぼ}せる}{}{}%
	}%
	}%
	{\tailPrint[]}%
\entry
	{\headPrint[kana={ぬぶしん\tl{˩˥}},ipa={nubuʃin},pos={名},acc={R}]}%
	{%
		\MeaningsBlock{%
	\ryumeaning%
		{首。喉元}{}{}%
	}%
	}%
	{\tailPrint[]}%
\entry
	{\headPrint[kana={ぬふた\tl{˩˥}はーん},ipa={nufu̥tahaːɴ},pos={形},acc={R}]}%
	{%
		\MeaningsBlock{%
	\ryumeaning%
		{眠たい}{}{}%
	}%
	}%
	{\tailPrint[]}%
\entry
	{\headPrint[kana={ぬぶたん\tl{˩˥}},ipa={nubutaɴ},pos={動（継）},rui={3},para={ぬぶとぅぬ},acc={R}]}%
	{%
		\MeaningsBlock{%
	\ryumeaning%
		{飲み尽くす}{}{}%
	}%
	}%
	{\tailPrint[]}%
\entry
	{\headPrint[kana={ぬふちるん},ipa={nufu̥tʃiruɴ},pos={句},para={ぬふとぅぬ}]}%
	{%
		\MeaningsBlock{%
	\ryumeaning%
		{寝かせる。寝かしつける}{}{}%
	}%
	}%
	{\tailPrint[]}%
\entry
	{\headPrint[kana={ぬぶるん\tl{˥˩}},ipa={nuburuɴ},pos={動},rui={1},para={ぬぶらぬ},acc={F}]}%
	{%
		\MeaningsBlock{%
	\ryumeaning%
		{登る。上がる}{}{}%
	\ryumeaning%
		{乗る}{}{}%
	}%
	}%
	{\tailPrint[]}%
\entry
	{\headPrint[kana={ぬぶん\tl{˥˩}},ipa={nubuɴ},pos={動},rui={1},para={ぬばぬ},acc={F}]}%
	{%
		\MeaningsBlock{%
	\ryumeaning%
		{伸びる}{}{}%
	}%
	}%
	{\tailPrint[]}%
\entry
	{\headPrint[kana={ぬみ\ws{}うらへ},ipa={numi urahe},pos={句}]}%
	{%
		\MeaningsBlock{%
	\ryumeaning%
		{飲み下ろす}{}{}%
	}%
	}%
	{\tailPrint[]}%
\entry
	{\headPrint[kana={ぬみくむん\tl{˩˥}},ipa={numiku̥muɴ},pos={動},rui={1},para={ぬみくまぬ},acc={R}]}%
	{%
		\MeaningsBlock{%
	\ryumeaning%
		{飲み込む}{}{}%
	}%
	}%
	{\tailPrint[]}%
\entry
	{\headPrint[kana={ぬみのがすん\tl{˩˥}},ipa={numinogasuɴ},pos={動},rui={2},para={ぬみのがはぬ},acc={R}]}%
	{%
		\MeaningsBlock{%
	\ryumeaning%
		{飲み残す}{}{}%
	}%
	}%
	{\tailPrint[]}%
\entry
	{\headPrint[kana={ぬむん\tl{˩˥}},ipa={numuɴ},pos={動},rui={1},para={ぬまぬ},acc={R}]}%
	{%
		\MeaningsBlock{%
	\ryumeaning%
		{飲む}{}{}%
	}%
	}%
	{\tailPrint[]}%
\entry
	{\headPrint[kana={ぬり\tl{˩˥}},ipa={nuri},pos={名},acc={R}]}%
	{%
		\MeaningsBlock{%
	\ryumeaning%
		{苔}{}{}%
	}%
	}%
	{\tailPrint[]}%
\entry
	{\headPrint[kana={ぬり\tl{˩˥}},ipa={nuri},pos={名},acc={R}]}%
	{%
		\MeaningsBlock{%
	\ryumeaning%
		{\ruby[g]{糊}{のり}}{}{}%
	}%
	}%
	{\tailPrint[]}%
\entry
	{\headPrint[kana={ぬりうくりるん},ipa={nuri.ukuriruɴ},pos={動},rui={3},para={ぬりうくるぬ}]}%
	{%
		\MeaningsBlock{%
	\ryumeaning%
		{乗り遅れる。乗りはずす}{}{}%
	}%
	}%
	{\tailPrint[]}%
\entry
	{\headPrint[kana={ぬりだっくゎーすん},ipa={nuridakkwaːsuɴ},pos={動}]}%
	{%
		\MeaningsBlock{%
	\ryumeaning%
		{塗りたくる}{}{}%
	}%
	}%
	{\tailPrint[]}%
\entry
	{\headPrint[kana={ぬりふもん\tl{˩˥}},ipa={nurifu̥mon},pos={名},acc={R}]}%
	{%
		\MeaningsBlock{%
	\ryumeaning%
		{積雲。積乱雲}{}{}%
	}%
	}%
	{\tailPrint[]}%
\entry
	{\headPrint[kana={ぬりむぬ\tl{˥˩}},ipa={nurimunu},pos={名},acc={F}]}%
	{%
		\MeaningsBlock{%
	\ryumeaning%
		{乗り物。交通機関}{}{}%
	}%
	}%
	{\tailPrint[]}%
\entry
	{\headPrint[kana={ぬるさ\tl{˩˥}はん},ipa={nurusahaɴ},pos={形},para={ぬるさへぬ},acc={R}]}%
	{%
		\MeaningsBlock{%
	\ryumeaning%
		{\ruby[g]{温}{ぬる}い}{}{}%
	}%
	}%
	{\tailPrint[]}%
\entry
	{\headPrint[kana={ぬるみゃん\tl{˩˥}},ipa={nurumjaɴ},pos={動（継）},rui={1},para={ぬるまぬ},acc={R}]}%
	{%
		\MeaningsBlock{%
	\ryumeaning%
		{\ruby[g]{温}{ぬる}まる}{}{}%
	}%
	}%
	{\tailPrint[]}%
\entry
	{\headPrint[kana={ぬるみるん\tl{˩˥}},ipa={nurumiruɴ},pos={動},acc={R}]}%
	{%
		\MeaningsBlock{%
	\ryumeaning%
		{\ruby[g]{温}{ぬる}くする}{温める}{}%
	}%
	}%
	{\tailPrint[]}%
\entry
	{\headPrint[kana={ぬるん\tl{˥˩}},ipa={nuruɴ},pos={動},rui={1},para={ぬらぬ},acc={F}]}%
	{%
		\MeaningsBlock{%
	\ryumeaning%
		{乗る。載る}{}{}%
	}%
	}%
	{\tailPrint[]}%
\entry
	{\headPrint[kana={ぬるん\tl{˥˩}},ipa={nuruɴ},pos={動},rui={1},para={ぬらぬ},acc={F}]}%
	{%
		\MeaningsBlock{%
	\ryumeaning%
		{塗る}{}{}%
	}%
	}%
	{\tailPrint[]}%
\entry
	{\headPrint[kana={ぬん\tl{˩˥}},ipa={nun},pos={名},acc={R}]}%
	{%
		\MeaningsBlock{%
	\ryumeaning%
		{\ruby[g]{蚤}{のみ}}{}{}%
	}%
	}%
	{\tailPrint[]}%
\entry
	{\headPrint[kana={ぬん\tl{˩˥}},ipa={nuɴ},pos={名},acc={R}]}%
	{%
		\MeaningsBlock{%
	\ryumeaning%
		{\ruby[g]{鑿}{のみ}}{}{}%
	}%
	}%
	{\tailPrint[]}%
\LetterHead{ね}
\entry
	{\headPrint[kana={ねー\tl{˩˥}},ipa={neː},pos={副},acc={R}]}%
	{%
		\MeaningsBlock{%
	\ryumeaning%
		{どう。どんな}{}{}%
	}%
	}%
	{\tailPrint[]}%
\entry
	{\headPrint[kana={ねーき\tl{˩˥}},ipa={neːki},pos={句},acc={R}]}%
	{%
		\MeaningsBlock{%
	\ryumeaning%
		{どうして。何故}{}{}%
	}%
	}%
	{\tailPrint[]}%
\entry
	{\headPrint[kana={ねーきる},ipa={neːkiru},pos={句}]}%
	{%
		\MeaningsBlock{%
	\ryumeaning%
		{どうして。何故}{}{}%
	}%
	}%
	{\tailPrint[]}%
\entry
	{\headPrint[kana={ねーしゃるん},ipa={neːʃaruɴ},pos={句}]}%
	{%
		\MeaningsBlock{%
	\ryumeaning%
		{どんな。如何なる}{}{}%
	}%
	}%
	{\tailPrint[]}%
\entry
	{\headPrint[kana={ねーな\tl{˩˥}さん\tl{˥˩}},ipa={neːnasaɴ},pos={動（継）},acc={R F}]}%
	{%
		\MeaningsBlock{%
	\ryumeaning%
		{なくす。失う。なくなった}{}{}%
	}%
	}%
	{\tailPrint[]}%
\entry
	{\headPrint[kana={ねー\tl{˩˥}\ws{}なるん\tl{˩˥}},ipa={neː naruɴ},pos={句},acc={R R}]}%
	{%
		\MeaningsBlock{%
	\ryumeaning%
		{なくなる}{}{}%
	}%
	}%
	{\tailPrint[]}%
\entry
	{\headPrint[kana={ねーぬ\tl{˩˥}},ipa={neːnu},pos={動},rui={irr},acc={R}]}%
	{%
		\MeaningsBlock{%
	\ryumeaning%
		{無い}{}{}%
	}%
	}%
	{\tailPrint[]}%
\entry
	{\headPrint[kana={\josho{}ねー\kako{}や},ipa={neːja},pos={句},acc={HL}]}%
	{%
		\MeaningsBlock{%
	\ryumeaning%
		{どう。どうかね}{}{}%
	}%
	}%
	{\tailPrint[]}%
\entry
	{\headPrint[kana={ねーり\tl{˥˩}},ipa={neːri},pos={名},acc={F}]}%
	{%
		\MeaningsBlock{%
	\ryumeaning%
		{右。右方。右側}{}{}%
	}%
	}%
	{\tailPrint[]}%
\entry
	{\headPrint[kana={ねーる\tl{˩˥}},ipa={neːru},pos={指示様態},acc={R}]}%
	{%
		\MeaningsBlock{%
	\ryumeaning%
		{どのように。どんな風に}{}{}%
	}%
	}%
	{\tailPrint[]}%
\entry
	{\headPrint[kana={ねーるん\tl{˥˩}},ipa={neːruɴ},pos={動},rui={1},para={ねーらぬ},acc={F}]}%
	{%
		\MeaningsBlock{%
	\ryumeaning%
		{煮える}{}{}%
	}%
	}%
	{\tailPrint[]}%
\entry
	{\headPrint[kana={ねっすん\tl{˥˩}},ipa={nessuɴ},pos={動},rui={2},acc={F}]}%
	{%
		\MeaningsBlock{%
	\ryumeaning%
		{煮る}{}{}%
	}%
	}%
	{\tailPrint[]}%
\LetterHead{の}
\entry
	{\headPrint[kana={のー\tl{˥˩}},ipa={noː},pos={名},acc={F}]}%
	{%
		\MeaningsBlock{%
	\ryumeaning%
		{脳。大脳}{}{}%
	}%
	}%
	{\tailPrint[]}%
\entry
	{\headPrint[kana={のーさ\tl{˩˥}\ws{}なるん\tl{˩˥}},ipa={noːsa naruɴ},pos={句},acc={R R}]}%
	{%
		\MeaningsBlock{%
	\ryumeaning%
		{暖かくなる。暖まる}{}{}%
	}%
	}%
	{\tailPrint[]}%
\entry
	{\headPrint[kana={のーさ\tl{˩˥}はん},ipa={noːsahaɴ},pos={形},para={のーさへぬ},acc={R}]}%
	{%
		\MeaningsBlock{%
	\ryumeaning%
		{暖かい}{}{}%
	}%
	}%
	{\tailPrint[]}%
\entry
	{\headPrint[kana={\josho{}のー\kako{}しん},ipa={noːʃiɴ},pos={副},acc={HL}]}%
	{%
		\MeaningsBlock{%
	\ryumeaning%
		{どうしても}{}{}%
	}%
	}%
	{\tailPrint[]}%
\entry
	{\headPrint[kana={のーじん\tl{˥˩}},ipa={noːdʒiɴ},pos={名},acc={F}]}%
	{%
		\MeaningsBlock{%
	\ryumeaning%
		{\ruby[g]{虹}{にじ}}{}{}%
	}%
	}%
	{\tailPrint[]}%
\entry
	{\headPrint[kana={のーすん\tl{˩˥}},ipa={noːsuɴ},pos={動},rui={2},para={のーはぬ},acc={R}]}%
	{%
		\MeaningsBlock{%
	\ryumeaning%
		{治す}{病気を治す}{}%
	\ryumeaning%
		{直す}{物を直す}{}%
	}%
	}%
	{\tailPrint[]}%
\entry
	{\headPrint[kana={のーなさん},ipa={noːnasaɴ},pos={動}]}%
	{%
		\MeaningsBlock{%
	\ryumeaning%
		{居なくなる}{}{}%
	}%
	}%
	{\tailPrint[]}%
\entry
	{\headPrint[kana={のーり\tl{˩˥}},ipa={noːri},pos={名},acc={R}]}%
	{%
		\MeaningsBlock{%
	\ryumeaning%
		{\ruby[g]{実}{みの}り}{}{}%
	}%
	}%
	{\tailPrint[]}%
\entry
	{\headPrint[kana={のーり\tl{˩˥}ゆ\tl{˥˩}},ipa={noːriju},pos={名},acc={R+F}]}%
	{%
		\MeaningsBlock{%
	\ryumeaning%
		{豊年}{「稔り世」の意}{}%
	}%
	}%
	{\tailPrint[]}%
\entry
	{\headPrint[kana={のーるん},ipa={noːrun},pos={動}]}%
	{%
		\MeaningsBlock{%
	\ryumeaning%
		{実る}{}{}%
	}%
	}%
	{\tailPrint[]}%
\entry
	{\headPrint[kana={のーるん\tl{˩˥}},ipa={noːruɴ},pos={動},rui={1},para={のーらぬ},acc={R}]}%
	{%
		\MeaningsBlock{%
	\ryumeaning%
		{治る}{体の病気が治る}{}%
	\ryumeaning%
		{直る}{物が直る}{}%
	}%
	}%
	{\tailPrint[]}%
\entry
	{\headPrint[kana={のがすん\tl{˥˩}},ipa={nogasuɴ},pos={動},rui={2},para={のがはぬ},acc={F}]}%
	{%
		\MeaningsBlock{%
	\ryumeaning%
		{残す}{}{}%
	}%
	}%
	{\tailPrint[]}%
\entry
	{\headPrint[kana={のがるん\tl{˩˥}},ipa={nogaruɴ},pos={動},rui={1},para={のがらぬ},acc={R}]}%
	{%
		\MeaningsBlock{%
	\ryumeaning%
		{残る。余る}{}{}%
	}%
	}%
	{\tailPrint[]}%
\entry
	{\headPrint[kana={のごり},ipa={nogori},pos={名}]}%
	{%
		\MeaningsBlock{%
	\ryumeaning%
		{残り。残部}{}{}%
	}%
	}%
	{\tailPrint[]}%
\LetterHead{は}
\entry
	{\headPrint[kana={はー},ipa={haː},pos={名},acc={-}]}%
	{%
		\MeaningsBlock{%
	\ryumeaning%
		{あそこ。あちら}{}{}%
	}%
	}%
	{\tailPrint[]}%
\entry
	{\headPrint[kana={ばー\tl{˩˥}},ipa={baː},pos={名},acc={R}]}%
	{%
		\MeaningsBlock{%
	\ryumeaning%
		{私。自分。私の}{}{}%
	}%
	}%
	{\tailPrint[]}%
\entry
	{\headPrint[kana={ぱー\tl{˥}},ipa={paː},pos={名},acc={H}]}%
	{%
		\MeaningsBlock{%
	\ryumeaning%
		{祖母。おばあさん}{}{}%
	}%
	}%
	{\tailPrint[]}%
\entry
	{\headPrint[kana={ぱー\tl{˥˩}},ipa={paː},pos={名},acc={F}]}%
	{%
		\MeaningsBlock{%
	\ryumeaning%
		{葉}{}{}%
	}%
	}%
	{\tailPrint[]}%
\entry
	{\headPrint[kana={ばー\ws{}うたま},ipa={baː utama},pos={句}]}%
	{%
		\MeaningsBlock{%
	\ryumeaning%
		{私の子}{}{}%
	}%
	}%
	{\tailPrint[]}%
\entry
	{\headPrint[kana={はーがに\tl{˥}},ipa={haːgani},pos={名},acc={H}]}%
	{%
		\MeaningsBlock{%
	\ryumeaning%
		{\ruby[g]{鋼}{はがね}}{「刃金」から}{}%
	}%
	}%
	{\tailPrint[]}%
\entry
	{\headPrint[kana={ばーき\tl{˩˥}},ipa={baːki},pos={名},acc={R}]}%
	{%
		\MeaningsBlock{%
	\ryumeaning%
		{ざる。竹ざる}{}{}%
	}%
	}%
	{\tailPrint[]}%
\entry
	{\headPrint[kana={ばーぐん\tl{˩˥}},ipa={baːguɴ},pos={動},rui={1},para={ばーがぬ},acc={R}]}%
	{%
		\MeaningsBlock{%
	\ryumeaning%
		{湧き出る。酒などが発酵する}{}{}%
	}%
	}%
	{\tailPrint[]}%
\entry
	{\headPrint[kana={ばーさ\tl{˩˥}},ipa={baːsa},pos={名},acc={R}]}%
	{%
		\MeaningsBlock{%
	\ryumeaning%
		{芭蕉。バナナ}{}{}%
	}%
	}%
	{\tailPrint[]}%
\entry
	{\headPrint[kana={ばーさぬ\tl{˩˥}\ws{}なーり\tl{˩˥}},ipa={baːsanu naːri},pos={句},acc={R R}]}%
	{%
		\MeaningsBlock{%
	\ryumeaning%
		{バナナの実}{}{}%
	}%
	}%
	{\tailPrint[]}%
\entry
	{\headPrint[kana={ばーしゃ\tl{˩˥}はん},ipa={baːʃahaɴ},pos={形},para={ばーしゃへぬ},acc={R}]}%
	{%
		\MeaningsBlock{%
	\ryumeaning%
		{\ruby[g]{可笑}{おか}しい}{}{}%
	}%
	}%
	{\tailPrint[]}%
\entry
	{\headPrint[kana={ばーすん\tl{˥˩}},ipa={baːsuɴ},pos={動},rui={2},para={ばーはぬ},acc={F}]}%
	{%
		\MeaningsBlock{%
	\ryumeaning%
		{（水などで）薄める}{}{}%
	}%
	}%
	{\tailPrint[]}%
\entry
	{\headPrint[kana={ぱーすん\tl{˩˥}},ipa={paːsuɴ},pos={動},acc={R}]}%
	{%
		\MeaningsBlock{%
	\ryumeaning%
		{囃す}{}{}%
	}%
	}%
	{\tailPrint[]}%
\entry
	{\headPrint[kana={ぱーち\tl{˥˩}},ipa={paːtʃi̥},pos={名},acc={F}]}%
	{%
		\MeaningsBlock{%
	\ryumeaning%
		{蜂}{}{}%
	}%
	}%
	{\tailPrint[]}%
\entry
	{\headPrint[kana={ぱーった},ipa={paːtta},pos={擬}]}%
	{%
		\MeaningsBlock{%
	\ryumeaning%
		{ぱあっと}{急に明るくなるさま}{}%
	}%
	}%
	{\tailPrint[]}%
\entry
	{\headPrint[kana={はーはー},ipa={haːhaː},pos={擬}]}%
	{%
		\MeaningsBlock{%
	\ryumeaning%
		{はあはあ}{あえぐさま}{}%
	}%
	}%
	{\tailPrint[]}%
\entry
	{\headPrint[kana={ばーばー},ipa={baːbaː},pos={擬}]}%
	{%
		\MeaningsBlock{%
	\ryumeaning%
		{びゅうびゅう}{風の強く吹くさま}{}%
	\ryumeaning%
		{ぼうぼう}{火の勢いよく燃えるさま}{}%
	}%
	}%
	{\tailPrint[]}%
\entry
	{\headPrint[kana={ばーらいるん\tl{˥˩}},ipa={baːrairuɴ},pos={動},rui={3},para={ばらーるぬ},acc={F}]}%
	{%
		\MeaningsBlock{%
	\ryumeaning%
		{笑われる。物笑いになる}{}{}%
	}%
	}%
	{\tailPrint[]}%
\entry
	{\headPrint[kana={ばーり\tl{˥˩}},ipa={baːri},pos={名},acc={F}]}%
	{%
		\MeaningsBlock{%
	\ryumeaning%
		{割れ目}{}{}%
	}%
	}%
	{\tailPrint[]}%
\entry
	{\headPrint[kana={ばーりゃん\tl{˥˩}},ipa={baːrjaɴ},pos={動},rui={1},para={ばーらぬ},acc={F}]}%
	{%
		\MeaningsBlock{%
	\ryumeaning%
		{微笑む}{}{}%
	}%
	}%
	{\tailPrint[]}%
\entry
	{\headPrint[kana={ばーるん\tl{˥˩}},ipa={baːruɴ},pos={動},rui={1},para={ばーらぬ},acc={F}]}%
	{%
		\MeaningsBlock{%
	\ryumeaning%
		{笑う}{}{}%
	}%
	}%
	{\tailPrint[]}%
\entry
	{\headPrint[kana={はいから\tl{˥}さーん},ipa={haikarasaːɴ},pos={形},acc={H}]}%
	{%
		\MeaningsBlock{%
	\ryumeaning%
		{ハイカラだ。おしゃれだ}{}{}%
	}%
	}%
	{\tailPrint[]}%
\entry
	{\headPrint[kana={ばいま\tl{˩˥}},ipa={baima},pos={名},acc={R}]}%
	{%
		\MeaningsBlock{%
	\ryumeaning%
		{私たち}{}{}%
	}%
	}%
	{\tailPrint[]}%
\entry
	{\headPrint[kana={ばいるん\tl{˥˩}},ipa={bairuɴ},pos={動},acc={F}]}%
	{%
		\MeaningsBlock{%
	\ryumeaning%
		{薄める}{}{}%
	}%
	}%
	{\tailPrint[]}%
\entry
	{\headPrint[kana={ぱいるん\tl{˥}},ipa={pairuɴ},pos={動},rui={1},para={ぱいらぬ},acc={H}]}%
	{%
		\MeaningsBlock{%
	\ryumeaning%
		{映える}{}{}%
	\ryumeaning%
		{似合う}{}{}%
	}%
	}%
	{\tailPrint[]}%
\entry
	{\headPrint[kana={ぱか\tl{˥}},ipa={pḁka},pos={名},acc={H}]}%
	{%
		\MeaningsBlock{%
	\ryumeaning%
		{区画。畝}{}{}%
	}%
	}%
	{\tailPrint[]}%
\entry
	{\headPrint[kana={ぱか\tl{˥˩}},ipa={pḁka},pos={名},acc={F}]}%
	{%
		\MeaningsBlock{%
	\ryumeaning%
		{墓}{}{}%
	}%
	}%
	{\tailPrint[]}%
\entry
	{\headPrint[kana={ばがいとぅるん\tl{˩˥}},ipa={bagaitu̥ruɴ},pos={動},acc={R}]}%
	{%
		\MeaningsBlock{%
	\ryumeaning%
		{奪い取る}{}{}%
	}%
	}%
	{\tailPrint[]}%
\entry
	{\headPrint[kana={ばがけーる\tl{˩˥}},ipa={bagakeːru},pos={句},para={ばがけーらぬ},acc={R}]}%
	{%
		\MeaningsBlock{%
	\ryumeaning%
		{若返る}{}{}%
	}%
	}%
	{\tailPrint[]}%
\entry
	{\headPrint[kana={はかじ\tl{˥}\ws{}とぅるん\tl{˥}},ipa={hḁkadʒi tu̥ruɴ},pos={句},acc={H H}]}%
	{%
		\MeaningsBlock{%
	\ryumeaning%
		{掻き集めて取る}{}{}%
	}%
	}%
	{\tailPrint[]}%
\entry
	{\headPrint[kana={はかじるん\tl{˥}},ipa={hḁkadʒiruɴ},pos={動},rui={1},para={はかじらぬ},acc={H}]}%
	{%
		\MeaningsBlock{%
	\ryumeaning%
		{引っ掻く}{}{}%
	}%
	}%
	{\tailPrint[]}%
\entry
	{\headPrint[kana={ばがすけん\tl{˩˥}},ipa={bagasu̥keɴ},pos={名},acc={R}]}%
	{%
		\MeaningsBlock{%
	\ryumeaning%
		{若月}{}{}%
	}%
	}%
	{\tailPrint[]}%
\entry
	{\headPrint[kana={はかずるん\tl{˥}},ipa={hḁkadzuruɴ},pos={動},rui={1},para={はかずらぬ},acc={H}]}%
	{%
		\MeaningsBlock{%
	\ryumeaning%
		{（爪で）引っ掻く}{}{}%
	}%
	}%
	{\tailPrint[]}%
\entry
	{\headPrint[kana={はかすん\tl{˥˩}},ipa={hḁkasuɴ},pos={動},rui={2},para={はかはぬ},acc={F}]}%
	{%
		\MeaningsBlock{%
	\ryumeaning%
		{弁償させる}{}{}%
	}%
	}%
	{\tailPrint[]}%
\entry
	{\headPrint[kana={ばがすん\tl{˥˩}},ipa={bagasuɴ},pos={動},rui={2},para={ばがはぬ},acc={F}]}%
	{%
		\MeaningsBlock{%
	\ryumeaning%
		{煮る}{}{}%
	}%
	}%
	{\tailPrint[]}%
\entry
	{\headPrint[kana={ばがすん\tl{˩˥}},ipa={bagasuɴ},pos={動},rui={2},para={ばがはぬ},acc={R}]}%
	{%
		\MeaningsBlock{%
	\ryumeaning%
		{奪う}{}{}%
	}%
	}%
	{\tailPrint[]}%
\entry
	{\headPrint[kana={ばがすん\tl{˩˥}},ipa={bagasuɴ},pos={動},rui={2},para={ばがはぬ},acc={R}]}%
	{%
		\MeaningsBlock{%
	\ryumeaning%
		{引き離す。離す}{}{}%
	\ryumeaning%
		{剥がす}{}{}%
	}%
	}%
	{\tailPrint[]}%
\entry
	{\headPrint[kana={ぱかすん\tl{˥}},ipa={pḁkasuɴ},pos={動},rui={2},para={ぱかはぬ},acc={H}]}%
	{%
		\MeaningsBlock{%
	\ryumeaning%
		{吐かせる}{}{}%
	}%
	}%
	{\tailPrint[]}%
\entry
	{\headPrint[kana={ぱかすん\tl{˥˩}},ipa={pḁkasuɴ},pos={動},rui={2},para={ぱかはぬ},acc={F}]}%
	{%
		\MeaningsBlock{%
	\ryumeaning%
		{（全責任を）負わせる}{}{}%
	}%
	}%
	{\tailPrint[]}%
\entry
	{\headPrint[kana={ばがっせー\tl{˩˥}},ipa={bagasseː},pos={名},acc={R}]}%
	{%
		\MeaningsBlock{%
	\ryumeaning%
		{若白髪}{}{}%
	}%
	}%
	{\tailPrint[]}%
\entry
	{\headPrint[kana={ばがなち\tl{˩˥}},ipa={baganatʃi},pos={名},acc={R}]}%
	{%
		\MeaningsBlock{%
	\ryumeaning%
		{若夏。初夏}{「若夏」の義}{}%
	}%
	}%
	{\tailPrint[]}%
\entry
	{\headPrint[kana={はがに},ipa={hagani},pos={名}]}%
	{%
		\MeaningsBlock{%
	\ryumeaning%
		{\ruby[g]{鋼}{はがね}}{「刃金」から}{}%
	}%
	}%
	{\tailPrint[]}%
\entry
	{\headPrint[kana={ばがぱ\tl{˩˥}},ipa={bagapa},pos={名},acc={R}]}%
	{%
		\MeaningsBlock{%
	\ryumeaning%
		{若葉。新緑}{}{}%
	}%
	}%
	{\tailPrint[]}%
\entry
	{\headPrint[kana={ばが\tl{˩˥}はん},ipa={bagahaɴ},pos={形},para={ばがへぬ},acc={R}]}%
	{%
		\MeaningsBlock{%
	\ryumeaning%
		{若い}{}{}%
	}%
	}%
	{\tailPrint[]}%
\entry
	{\headPrint[kana={ばがへ\ws{}とぅるん\tl{˥}},ipa={bagahe tu̥ruɴ},pos={句},acc={R H}]}%
	{%
		\MeaningsBlock{%
	\ryumeaning%
		{奪い取る}{}{}%
	}%
	}%
	{\tailPrint[]}%
\entry
	{\headPrint[kana={ぱかま\tl{˥˩}},ipa={pḁkama},pos={名},acc={F}]}%
	{%
		\MeaningsBlock{%
	\ryumeaning%
		{\ruby[g]{袴}{はかま}}{}{}%
	}%
	}%
	{\tailPrint[]}%
\entry
	{\headPrint[kana={ばがむぬ\tl{˩˥}},ipa={bagamunu},pos={名},acc={R}]}%
	{%
		\MeaningsBlock{%
	\ryumeaning%
		{若者。青年}{}{}%
	}%
	}%
	{\tailPrint[]}%
\entry
	{\headPrint[kana={はからすん\tl{˥}},ipa={hḁkarasuɴ},pos={動},rui={2},para={はからはぬ},acc={H}]}%
	{%
		\MeaningsBlock{%
	\ryumeaning%
		{（魚を）網で捕る}{}{}%
	}%
	}%
	{\tailPrint[]}%
\entry
	{\headPrint[kana={ぱからっ\tl{˥}さ},ipa={pḁkarassa},pos={形},acc={H}]}%
	{%
		\MeaningsBlock{%
	\ryumeaning%
		{かんばしい。立派}{}{}%
	}%
	}%
	{\tailPrint[]}%
\entry
	{\headPrint[kana={ぱからっさー\tl{˥}\ws{}ねーぬ\tl{˩˥}},ipa={pḁkarassaː neːnu},pos={句},acc={H R}]}%
	{%
		\MeaningsBlock{%
	\ryumeaning%
		{芳しくない。良くない}{\emb{まぱからさねーぬ}は強調語}{}%
	}%
	}%
	{\tailPrint[]}%
\entry
	{\headPrint[kana={ばがらぬ\tl{˩˥}},ipa={bagaranu},pos={句},acc={R}]}%
	{%
		\MeaningsBlock{%
	\ryumeaning%
		{知らない}{}{}%
	}%
	}%
	{\tailPrint[]}%
\entry
	{\headPrint[kana={ぱがりるん\tl{˥}},ipa={pagariruɴ},pos={動},rui={3},para={ぱがるぬ},acc={H}]}%
	{%
		\MeaningsBlock{%
	\ryumeaning%
		{剥がれる}{}{}%
	}%
	}%
	{\tailPrint[]}%
\entry
	{\headPrint[kana={ばがりん\tl{˩˥}},ipa={bagariɴ},pos={動},rui={3},para={ばがるぬ},acc={R}]}%
	{%
		\MeaningsBlock{%
	\ryumeaning%
		{別れる。離婚する}{}{}%
	}%
	}%
	{\tailPrint[]}%
\entry
	{\headPrint[kana={はかるん\tl{˥}},ipa={hḁkaruɴ},pos={動},rui={1},para={はからぬ},acc={H}]}%
	{%
		\MeaningsBlock{%
	\ryumeaning%
		{掛かる。引っ掛かる}{}{}%
	\ryumeaning%
		{（禁忌に）かかる}{}{}%
	\ryumeaning%
		{（病気に）罹る。患う}{}{}%
	}%
	}%
	{\tailPrint[]}%
\entry
	{\headPrint[kana={はかるん\tl{˥}},ipa={hḁkaruɴ},pos={動},rui={1},para={はからぬ},acc={H}]}%
	{%
		\MeaningsBlock{%
	\ryumeaning%
		{謀る。企む}{}{}%
	}%
	}%
	{\tailPrint[]}%
\entry
	{\headPrint[kana={ばがるん\tl{˩˥}},ipa={bagaruɴ},pos={動},rui={1},para={ばがらぬ},acc={R}]}%
	{%
		\MeaningsBlock{%
	\ryumeaning%
		{分かる。理解する}{}{}%
	\ryumeaning%
		{推し量る}{}{}%
	}%
	}%
	{\tailPrint[]}%
\entry
	{\headPrint[kana={ぱかるん\tl{˥}},ipa={pḁkaruɴ},pos={動},rui={1},para={ぱからぬ},acc={H}]}%
	{%
		\MeaningsBlock{%
	\ryumeaning%
		{計る。計測する}{}{}%
	}%
	}%
	{\tailPrint[]}%
\entry
	{\headPrint[kana={はき\tl{˥˩}},ipa={hḁki},pos={名},acc={F}]}%
	{%
		\MeaningsBlock{%
	\ryumeaning%
		{欠片。破片}{}{}%
	}%
	}%
	{\tailPrint[]}%
\entry
	{\headPrint[kana={ばぎ},ipa={bagi},pos={助}]}%
	{%
		\MeaningsBlock{%
	\ryumeaning%
		{～まで}{}{}%
	}%
	}%
	{\tailPrint[]}%
\entry
	{\headPrint[kana={はきあしみるん\tl{˥}},ipa={hḁki.aʃi̥miruɴ},pos={動},rui={3},para={はきあすむぬ},acc={H}]}%
	{%
		\MeaningsBlock{%
	\ryumeaning%
		{掻き集める}{}{}%
	}%
	}%
	{\tailPrint[]}%
\entry
	{\headPrint[kana={はき\tl{˥}いりるん\tl{˥˩}},ipa={hḁki.iriruɴ},pos={動},rui={3},para={はきいるぬ},acc={H+F}]}%
	{%
		\MeaningsBlock{%
	\ryumeaning%
		{書き入れる}{}{}%
	}%
	}%
	{\tailPrint[]}%
\entry
	{\headPrint[kana={はき\tl{˥}いりるん\tl{˥˩}},ipa={haki.iriruɴ},pos={動},acc={H+F}]}%
	{%
		\MeaningsBlock{%
	\ryumeaning%
		{口にかき込む}{}{}%
	}%
	}%
	{\tailPrint[]}%
\entry
	{\headPrint[kana={ばぎし\tl{˩˥}},ipa={bagiʃi},pos={名},acc={R}]}%
	{%
		\MeaningsBlock{%
	\ryumeaning%
		{バケツ}{}{}%
	}%
	}%
	{\tailPrint[]}%
\entry
	{\headPrint[kana={ぱぎじ\tl{˥}},ipa={pagidʒi},pos={名},acc={H}]}%
	{%
		\MeaningsBlock{%
	\ryumeaning%
		{痩せ地}{}{}%
	}%
	}%
	{\tailPrint[]}%
\entry
	{\headPrint[kana={はきしぃきるん\tl{˥}},ipa={hḁkiʃi̥kiruɴ},pos={動},rui={3},para={はきすくぬ},acc={H}]}%
	{%
		\MeaningsBlock{%
	\ryumeaning%
		{駆ける。走る}{}{}%
	}%
	}%
	{\tailPrint[note={「駆け付ける」に対応}]}%
\entry
	{\headPrint[kana={はきじゃー\tl{˥}},ipa={hakidʒaː},pos={名},acc={H}]}%
	{%
		\MeaningsBlock{%
	\ryumeaning%
		{\ruby[g]{雨}{あま}\ruby[g]{樋}{どい}}{}{}%
	}%
	}%
	{\tailPrint[]}%
\entry
	{\headPrint[kana={ぱぎすぷる\tl{˥}},ipa={pagisu̥puru},pos={名},acc={H}]}%
	{%
		\MeaningsBlock{%
	\ryumeaning%
		{\ruby[g]{禿}{はげ}\ruby[g]{頭}{あたま}}{}{}%
	}%
	}%
	{\tailPrint[]}%
\entry
	{\headPrint[kana={はきたすん\tl{˥}},ipa={hḁkitasuɴ},pos={動},rui={2b},para={はきたさぬ},acc={H}]}%
	{%
		\MeaningsBlock{%
	\ryumeaning%
		{書き添える。書き足す}{}{}%
	}%
	}%
	{\tailPrint[]}%
\entry
	{\headPrint[kana={ばぎだま\tl{˩˥}},ipa={bagidama},pos={名},acc={R}]}%
	{%
		\MeaningsBlock{%
	\ryumeaning%
		{分け前}{}{}%
	}%
	}%
	{\tailPrint[]}%
\entry
	{\headPrint[kana={はきとぅみるん\tl{˥}},ipa={hḁkitu̥miruɴ},pos={動},rui={3},para={はきとぅむぬ},acc={H}]}%
	{%
		\MeaningsBlock{%
	\ryumeaning%
		{書き留める}{}{}%
	}%
	}%
	{\tailPrint[]}%
\entry
	{\headPrint[kana={ばぎとぅるん\tl{˥}},ipa={pagituruɴ},pos={動},rui={1},para={ぱぎとぅらぬ},acc={H}]}%
	{%
		\MeaningsBlock{%
	\ryumeaning%
		{剥ぎ取る}{}{}%
	\ryumeaning%
		{かっぱらう}{}{}%
	}%
	}%
	{\tailPrint[]}%
\entry
	{\headPrint[kana={はきぽーるん\tl{˥}},ipa={hḁkipoːruɴ},pos={動},rui={1},para={はきぽーらぬ},acc={H}]}%
	{%
		\MeaningsBlock{%
	\ryumeaning%
		{掻き散らかす}{}{}%
	}%
	}%
	{\tailPrint[]}%
\entry
	{\headPrint[kana={はきまーすん\tl{˥}},ipa={hḁkimaːsuɴ},pos={動},para={はきまはぬ},acc={H}]}%
	{%
		\MeaningsBlock{%
	\ryumeaning%
		{掻き回す}{}{}%
	}%
	}%
	{\tailPrint[]}%
\entry
	{\headPrint[kana={はきまーるん\tl{˥}},ipa={hakimaːruɴ},pos={動},acc={H}]}%
	{%
		\MeaningsBlock{%
	\ryumeaning%
		{駆け回る}{}{}%
	}%
	}%
	{\tailPrint[]}%
\entry
	{\headPrint[kana={はき\tl{˥}まんじるん\tl{˩˥}},ipa={hḁkimandʒiruɴ},pos={動},rui={3},para={はきまんずぬ},acc={H+R}]}%
	{%
		\MeaningsBlock{%
	\ryumeaning%
		{掻き乱す。掻き混ぜる}{}{}%
	\ryumeaning%
		{いびる}{}{}%
	}%
	}%
	{\tailPrint[]}%
\entry
	{\headPrint[kana={ばぎみじ\tl{˥˩}},ipa={bagimidʒi},pos={名},acc={F}]}%
	{%
		\MeaningsBlock{%
	\ryumeaning%
		{湧水。泉}{}{}%
	}%
	}%
	{\tailPrint[]}%
\entry
	{\headPrint[kana={はきもらすん\tl{˥}},ipa={hḁkimorasuɴ},pos={動},rui={2},para={はきもらはぬ},acc={H}]}%
	{%
		\MeaningsBlock{%
	\ryumeaning%
		{書き落とす。書き漏らす}{}{}%
	}%
	}%
	{\tailPrint[]}%
\entry
	{\headPrint[kana={はきるん\tl{˥}},ipa={hḁkiruɴ},pos={動},rui={3},para={はくぬ},acc={H}]}%
	{%
		\MeaningsBlock{%
	\ryumeaning%
		{掛ける}{}{}%
	}%
	}%
	{\tailPrint[]}%
\entry
	{\headPrint[kana={はきるん\tl{˥˩}},ipa={hḁkiruɴ},pos={動},rui={3},para={はくぬ},acc={F}]}%
	{%
		\MeaningsBlock{%
	\ryumeaning%
		{欠ける}{}{}%
	}%
	}%
	{\tailPrint[]}%
\entry
	{\headPrint[kana={ばぎるん\tl{˩˥}},ipa={bagiruɴ},pos={動},rui={3},para={ばぐぬ},acc={R}]}%
	{%
		\MeaningsBlock{%
	\ryumeaning%
		{分ける。和解させる}{}{}%
	}%
	}%
	{\tailPrint[]}%
\entry
	{\headPrint[kana={ぱぎるん\tl{˥}},ipa={pagiruɴ},pos={動},rui={3},para={ぱぐぬ},acc={H}]}%
	{%
		\MeaningsBlock{%
	\ryumeaning%
		{\ruby[g]{色}{いろ}\ruby[g]{褪}{あ}せる}{}{}%
	}%
	}%
	{\tailPrint[]}%
\entry
	{\headPrint[kana={はきん\tl{˥˩}},ipa={hḁkiɴ},pos={動},rui={3},para={はくぬ},acc={F}]}%
	{%
		\MeaningsBlock{%
	\ryumeaning%
		{欠ける}{}{}%
	}%
	}%
	{\tailPrint[]}%
\entry
	{\headPrint[kana={ばぎん\tl{˩˥}},ipa={bagiɴ},pos={動},rui={3},para={ばーぐぬ},acc={R}]}%
	{%
		\MeaningsBlock{%
	\ryumeaning%
		{分ける。分別する}{}{}%
	}%
	}%
	{\tailPrint[]}%
\entry
	{\headPrint[kana={ぱぎん\tl{˥}},ipa={pagiɴ},pos={動},rui={3},para={ぱぐぬ},acc={H}]}%
	{%
		\MeaningsBlock{%
	\ryumeaning%
		{禿げる}{}{}%
	}%
	}%
	{\tailPrint[]}%
\entry
	{\headPrint[kana={はきんだすん\tl{˥}},ipa={hḁkindasuɴ},pos={動},rui={2},para={はきんだはぬ},acc={H}]}%
	{%
		\MeaningsBlock{%
	\ryumeaning%
		{掻き出す}{}{}%
	}%
	}%
	{\tailPrint[]}%
\entry
	{\headPrint[kana={ぱきんだすん\tl{˥}},ipa={pḁkindasuɴ},pos={動},rui={2},para={ぱきんだはぬ},acc={H}]}%
	{%
		\MeaningsBlock{%
	\ryumeaning%
		{吐き出す}{}{}%
	}%
	}%
	{\tailPrint[]}%
\entry
	{\headPrint[kana={ぱく\tl{˥˩}},ipa={pḁku},pos={名},acc={F}]}%
	{%
		\MeaningsBlock{%
	\ryumeaning%
		{箱}{}{}%
	}%
	}%
	{\tailPrint[]}%
\entry
	{\headPrint[kana={ぱく\tl{˥˩}},ipa={pḁku},pos={名},acc={F}]}%
	{%
		\MeaningsBlock{%
	\ryumeaning%
		{蛇}{}{}%
	}%
	}%
	{\tailPrint[]}%
\entry
	{\headPrint[kana={ばくるん\tl{˩˥}},ipa={baku̥ruɴ},pos={動},rui={1},para={ばくらぬ},acc={R}]}%
	{%
		\MeaningsBlock{%
	\ryumeaning%
		{侮る。嘲る}{}{}%
	}%
	}%
	{\tailPrint[]}%
\entry
	{\headPrint[kana={ぱくるん\tl{˥}},ipa={pḁkuruɴ},pos={動},rui={1},acc={H}]}%
	{%
		\MeaningsBlock{%
	\ryumeaning%
		{からかう。ひやかす}{}{}%
	}%
	}%
	{\tailPrint[]}%
\entry
	{\headPrint[kana={はくん\tl{˥}},ipa={hḁkuɴ},pos={動},rui={1},para={はかぬ},acc={H}]}%
	{%
		\MeaningsBlock{%
	\ryumeaning%
		{書く。描く}{}{}%
	}%
	}%
	{\tailPrint[]}%
\entry
	{\headPrint[kana={はくん\tl{˥}},ipa={hḁkuɴ},pos={動},rui={1},para={はかぬ},acc={H}]}%
	{%
		\MeaningsBlock{%
	\ryumeaning%
		{掃く}{}{}%
	}%
	}%
	{\tailPrint[]}%
\entry
	{\headPrint[kana={はくん\tl{˥}},ipa={hḁkuɴ},pos={動},rui={1},para={はかぬ},acc={H}]}%
	{%
		\MeaningsBlock{%
	\ryumeaning%
		{搔く}{}{}%
	}%
	}%
	{\tailPrint[]}%
\entry
	{\headPrint[kana={ばぐん\tl{˥}},ipa={baguɴ},pos={動},rui={1},acc={H}]}%
	{%
		\MeaningsBlock{%
	\ryumeaning%
		{（木を）挽く}{}{}%
	}%
	}%
	{\tailPrint[]}%
\entry
	{\headPrint[kana={ばぐん\tl{˥˩}},ipa={baguɴ},pos={動},rui={1},para={ばがぬ},acc={F}]}%
	{%
		\MeaningsBlock{%
	\ryumeaning%
		{（水が）湧く。（酒が）醸される}{}{}%
	}%
	}%
	{\tailPrint[]}%
\entry
	{\headPrint[kana={ぱくん\tl{˥˩}},ipa={pḁkuɴ},pos={動},rui={1},para={ぱかぬ},acc={F}]}%
	{%
		\MeaningsBlock{%
	\ryumeaning%
		{佩く}{}{}%
	\ryumeaning%
		{ひっかぶる}{}{}%
	\ryumeaning%
		{償う。責任を負う。弁償する}{}{}%
	}%
	}%
	{\tailPrint[]}%
\entry
	{\headPrint[kana={ぱくん\tl{˥˩}},ipa={hḁkuɴ},pos={動},rui={1},para={はかぬ},acc={F}]}%
	{%
		\MeaningsBlock{%
	\ryumeaning%
		{償う。責任を負う。弁償する}{}{}%
	}%
	}%
	{\tailPrint[]}%
\entry
	{\headPrint[kana={ぱくん\tl{˥˩}},ipa={pḁkuɴ},pos={動},rui={1},para={ぱかぬ}]}%
	{%
		\MeaningsBlock{%
	\ryumeaning%
		{吐く}{}{}%
	}%
	}%
	{\tailPrint[]}%
\entry
	{\headPrint[kana={ぱぐん\tl{˥}},ipa={paguɴ},pos={動},rui={1},para={ぱがぬ},acc={H}]}%
	{%
		\MeaningsBlock{%
	\ryumeaning%
		{剥ぐ。剥ぎ取る}{}{}%
	}%
	}%
	{\tailPrint[]}%
\entry
	{\headPrint[kana={はこー\tl{˥˩}},ipa={hḁkoː},pos={名},acc={F}]}%
	{%
		\MeaningsBlock{%
	\ryumeaning%
		{水夫。船員}{}{}%
	}%
	}%
	{\tailPrint[note={古語「水夫（かこ）」に対応}]}%
\entry
	{\headPrint[kana={はこー\tl{˥˩}},ipa={hḁkoː},pos={名},acc={F}]}%
	{%
		\MeaningsBlock{%
	\ryumeaning%
		{屋敷}{}{}%
	}%
	}%
	{\tailPrint[]}%
\entry
	{\headPrint[kana={ばごー\tl{˩˥}},ipa={bagoː},pos={名},acc={R}]}%
	{%
		\MeaningsBlock{%
	\ryumeaning%
		{ノカラムシ}{雑草名}{}%
	}%
	}%
	{\tailPrint[]}%
\entry
	{\headPrint[kana={はこすん\tl{˥}},ipa={hḁkosuɴ},pos={動},rui={2},para={はこはぬ},acc={H}]}%
	{%
		\MeaningsBlock{%
	\ryumeaning%
		{隠す}{}{}%
	}%
	}%
	{\tailPrint[]}%
\entry
	{\headPrint[kana={はこち\tl{˥˩}},ipa={hḁkotʃi},pos={名},acc={F}]}%
	{%
		\MeaningsBlock{%
	\ryumeaning%
		{\ruby[g]{顎}{あご}}{}{}%
	}%
	}%
	{\tailPrint[]}%
\entry
	{\headPrint[kana={はこますん\tl{˥˩}},ipa={hḁkomasuɴ},pos={動},rui={2},para={はこまはぬ},acc={F}]}%
	{%
		\MeaningsBlock{%
	\ryumeaning%
		{囲む。囲ませる}{}{}%
	}%
	}%
	{\tailPrint[]}%
\entry
	{\headPrint[kana={はこむん\tl{˥˩}},ipa={hḁkomuɴ},pos={動},acc={F}]}%
	{%
		\MeaningsBlock{%
	\ryumeaning%
		{囲む}{}{}%
	}%
	}%
	{\tailPrint[]}%
\entry
	{\headPrint[kana={はこりん\tl{˥}},ipa={hḁkoriɴ},pos={動},rui={3},para={はこるぬ},acc={H}]}%
	{%
		\MeaningsBlock{%
	\ryumeaning%
		{隠れる}{}{}%
	}%
	}%
	{\tailPrint[]}%
\entry
	{\headPrint[kana={ばざーるん\tl{˩˥}},ipa={badzaːruɴ},pos={動},rui={1},para={ばざーらぬ},acc={R}]}%
	{%
		\MeaningsBlock{%
	\ryumeaning%
		{勢いづく}{勢いが強くなる}{}%
	}%
	}%
	{\tailPrint[]}%
\entry
	{\headPrint[kana={ばさすぅぬ\tl{˩˥}},ipa={basasu̥nu},pos={名},acc={R}]}%
	{%
		\MeaningsBlock{%
	\ryumeaning%
		{芭蕉布の着物}{}{}%
	}%
	}%
	{\tailPrint[]}%
\entry
	{\headPrint[kana={ばざるん\tl{˥˩}},ipa={badzaruɴ},pos={動},rui={1},para={ばざらーぬ},acc={F}]}%
	{%
		\MeaningsBlock{%
	\ryumeaning%
		{はしゃぐ。気が高まる}{}{}%
	}%
	}%
	{\tailPrint[]}%
\entry
	{\headPrint[kana={ばしぃ\tl{˥˩}},ipa={baʃi},pos={名},acc={F}]}%
	{%
		\MeaningsBlock{%
	\ryumeaning%
		{\ruby[g]{鷲}{わし}}{鳥の名}{}%
	}%
	}%
	{\tailPrint[]}%
\entry
	{\headPrint[kana={はじまるん\tl{˥˩}},ipa={hadʒimaruɴ},pos={動},rui={1},para={はじまらぬ},acc={F}]}%
	{%
		\MeaningsBlock{%
	\ryumeaning%
		{始まる}{}{}%
	}%
	}%
	{\tailPrint[]}%
\entry
	{\headPrint[kana={はじみるん\tl{˥˩}},ipa={hadʒimiruɴ},pos={動},rui={3},para={はじむぬ},acc={F}]}%
	{%
		\MeaningsBlock{%
	\ryumeaning%
		{始める}{}{}%
	}%
	}%
	{\tailPrint[]}%
\entry
	{\headPrint[kana={ばしゅ\tl{˩˥}},ipa={baʃu},pos={名},acc={R}]}%
	{%
		\MeaningsBlock{%
	\ryumeaning%
		{場所。場面}{}{}%
	}%
	}%
	{\tailPrint[]}%
\entry
	{\headPrint[kana={ぱしゅくりん\tl{˥}},ipa={pḁʃukuriɴ},pos={動},rui={3},para={ぱしゅくるぬ},acc={H}]}%
	{%
		\MeaningsBlock{%
	\ryumeaning%
		{\ruby[g]{弾}{はじ}ける。破裂する}{}{}%
	}%
	}%
	{\tailPrint[]}%
\entry
	{\headPrint[kana={ぱすぅか\tl{˥}さーん},ipa={pasu̥kasaːɴ},pos={形},acc={H}]}%
	{%
		\MeaningsBlock{%
	\ryumeaning%
		{（人に対して）恥ずかしい}{世間体が悪い。きまりが悪い。面目がない}{}%
	}%
	}%
	{\tailPrint[]}%
\entry
	{\headPrint[kana={ぱすぅこー\tl{˥}はん},ipa={pasu̥koːhaɴ},pos={形},acc={H}]}%
	{%
		\MeaningsBlock{%
	\ryumeaning%
		{（稲、麦などの芒が皮膚をつきさすような）痛くて痒い感じがする}{}{}%
	}%
	}%
	{\tailPrint[]}%
\entry
	{\headPrint[kana={ぱすかはん},ipa={pasu̥kahaɴ},pos={形},para={ぱすかへぬ},acc={-}]}%
	{%
		\MeaningsBlock{%
	\ryumeaning%
		{恥ずかしい}{}{}%
	}%
	}%
	{\tailPrint[]}%
\entry
	{\headPrint[kana={ばずら\tl{˩˥}},ipa={badzura},pos={名},acc={R}]}%
	{%
		\MeaningsBlock{%
	\ryumeaning%
		{トカゲ}{動物名}{}%
	}%
	}%
	{\tailPrint[]}%
\entry
	{\headPrint[kana={ばずらすん\tl{˥˩}},ipa={badzurasuɴ},pos={動},rui={2},para={ばずらはぬ},acc={F}]}%
	{%
		\MeaningsBlock{%
	\ryumeaning%
		{屠殺する。解体する}{}{}%
	}%
	}%
	{\tailPrint[]}%
\entry
	{\headPrint[kana={ばた\tl{˩˥}},ipa={bata},pos={名},acc={R}]}%
	{%
		\MeaningsBlock{%
	\ryumeaning%
		{\ruby[g]{餡}{あん}。餡子}{}{}%
	}%
	}%
	{\tailPrint[]}%
\entry
	{\headPrint[kana={ばた\tl{˩˥}},ipa={bata},pos={名},acc={R}]}%
	{%
		\MeaningsBlock{%
	\ryumeaning%
		{腹。腹わた。内臓}{}{}%
	}%
	}%
	{\tailPrint[]}%
\entry
	{\headPrint[kana={ばた\tl{˩˥}},ipa={bata},pos={名},acc={R}]}%
	{%
		\MeaningsBlock{%
	\ryumeaning%
		{綿。木綿}{}{}%
	}%
	}%
	{\tailPrint[]}%
\entry
	{\headPrint[kana={ぱた\tl{˥˩}},ipa={pḁta},pos={名},acc={F}]}%
	{%
		\MeaningsBlock{%
	\ryumeaning%
		{側。端。傍ら}{}{}%
	}%
	}%
	{\tailPrint[]}%
\entry
	{\headPrint[kana={ぱた\tl{˥˩}},ipa={pḁta},pos={名},acc={F}]}%
	{%
		\MeaningsBlock{%
	\ryumeaning%
		{\ruby[g]{旗}{はた}}{}{}%
	}%
	}%
	{\tailPrint[]}%
\entry
	{\headPrint[kana={ばたさ\tl{˩˥}},ipa={batasa},pos={名},acc={R}]}%
	{%
		\MeaningsBlock{%
	\ryumeaning%
		{下男。小使い}{}{}%
	}%
	}%
	{\tailPrint[]}%
\entry
	{\headPrint[kana={ばたすん\tl{˥˩}},ipa={batasuɴ},pos={動},rui={2},para={ばたはぬ},acc={F}]}%
	{%
		\MeaningsBlock{%
	\ryumeaning%
		{渡す}{}{}%
	}%
	}%
	{\tailPrint[]}%
\entry
	{\headPrint[kana={ぱたち\tl{˥˩}},ipa={pḁtatʃi},pos={名},acc={F}]}%
	{%
		\MeaningsBlock{%
	\ryumeaning%
		{二十歳}{}{}%
	}%
	}%
	{\tailPrint[]}%
\entry
	{\headPrint[kana={ばたぬ\tl{˩˥}\ws{}むしぃ\tl{˥˩}},ipa={batanu muʃi},pos={句},acc={R F}]}%
	{%
		\MeaningsBlock{%
	\ryumeaning%
		{寄生虫。回虫}{}{}%
	}%
	}%
	{\tailPrint[]}%
\entry
	{\headPrint[kana={ばたふくら},ipa={batafu̥kura},pos={句}]}%
	{%
		\MeaningsBlock{%
	\ryumeaning%
		{腹がもたれる。消化不良}{}{}%
	}%
	}%
	{\tailPrint[]}%
\entry
	{\headPrint[kana={ばたふくりるん\tl{˩˥}},ipa={batafu̥kuriruɴ},pos={動},acc={R}]}%
	{%
		\MeaningsBlock{%
	\ryumeaning%
		{（消化不良で）腹が膨れる}{}{}%
	}%
	}%
	{\tailPrint[]}%
\entry
	{\headPrint[kana={ばたぶたー\tl{˩˥}},ipa={batabutaː},pos={名},acc={R}]}%
	{%
		\MeaningsBlock{%
	\ryumeaning%
		{でぶ。肥満者}{}{}%
	}%
	}%
	{\tailPrint[]}%
\entry
	{\headPrint[kana={ばたふちり\tl{˩˥}},ipa={batafu̥tʃiri},pos={名},acc={R}]}%
	{%
		\MeaningsBlock{%
	\ryumeaning%
		{悪口を言うこと}{}{}%
	}%
	}%
	{\tailPrint[]}%
\entry
	{\headPrint[kana={ばたふちりむぬ\tl{˩˥}},ipa={batafu̥tʃirimunu},pos={名},acc={R}]}%
	{%
		\MeaningsBlock{%
	\ryumeaning%
		{憎まれ口}{}{}%
	}%
	}%
	{\tailPrint[]}%
\entry
	{\headPrint[kana={ぱたむん\tl{˥˩}},ipa={pḁtamuɴ},pos={名},acc={F}]}%
	{%
		\MeaningsBlock{%
	\ryumeaning%
		{機織り機}{}{}%
	}%
	}%
	{\tailPrint[]}%
\entry
	{\headPrint[kana={ばた\tl{˩˥}やみ\tl{˩˥}},ipa={batajami},pos={名},acc={R+R}]}%
	{%
		\MeaningsBlock{%
	\ryumeaning%
		{腹痛}{}{}%
	}%
	}%
	{\tailPrint[]}%
\entry
	{\headPrint[kana={ばた\tl{˩˥}よー\tl{˩˥}さーん},ipa={batajoːsaːɴ},pos={形},acc={R+R}]}%
	{%
		\MeaningsBlock{%
	\ryumeaning%
		{お腹をこわしやすい}{}{}%
	}%
	}%
	{\tailPrint[]}%
\entry
	{\headPrint[kana={ぱたらぐん\tl{˥˩}},ipa={pḁtaraguɴ},pos={動},rui={1},para={ぱたらがぬ},acc={F}]}%
	{%
		\MeaningsBlock{%
	\ryumeaning%
		{働く。仕事する}{}{}%
	}%
	}%
	{\tailPrint[]}%
\entry
	{\headPrint[kana={ぱたらしやみどぅむ},ipa={pataraʃijamidumu},pos={名}]}%
	{%
		\MeaningsBlock{%
	\ryumeaning%
		{お転婆}{}{}%
	}%
	}%
	{\tailPrint[]}%
\entry
	{\headPrint[kana={ばたるん\tl{˥˩}},ipa={bataruɴ},pos={動},rui={1},para={ばたらぬ},acc={F}]}%
	{%
		\MeaningsBlock{%
	\ryumeaning%
		{渡る}{}{}%
	}%
	}%
	{\tailPrint[]}%
\entry
	{\headPrint[kana={ばた\tl{˩˥}んち\tl{˩˥}},ipa={batantʃi},pos={句},acc={R+R}]}%
	{%
		\MeaningsBlock{%
	\ryumeaning%
		{満腹になる。腹一杯になる}{}{}%
	}%
	}%
	{\tailPrint[]}%
\entry
	{\headPrint[kana={ばち},ipa={batʃi},pos={名},acc={F}]}%
	{%
		\MeaningsBlock{%
	\ryumeaning%
		{（大鼓の）ばち}{}{}%
	}%
	}%
	{\tailPrint[]}%
\entry
	{\headPrint[kana={ばち\tl{˥˩}},ipa={batʃi},pos={名},acc={F}]}%
	{%
		\MeaningsBlock{%
	\ryumeaning%
		{罰}{}{}%
	}%
	}%
	{\tailPrint[]}%
\entry
	{\headPrint[kana={ぱち\tl{˥˩}},ipa={pḁtʃi},pos={名},acc={F}]}%
	{%
		\MeaningsBlock{%
	\ryumeaning%
		{橋。梯子}{}{}%
	}%
	}%
	{\tailPrint[]}%
\entry
	{\headPrint[kana={ぱち\tl{˥˩}},ipa={pḁtʃi},pos={名},acc={F}]}%
	{%
		\MeaningsBlock{%
	\ryumeaning%
		{お初。初物}{}{}%
	}%
	}%
	{\tailPrint[]}%
\entry
	{\headPrint[kana={ばちさはん},ipa={batʃi̥sahaɴ},pos={形},para={ばちさへぬ},acc={-}]}%
	{%
		\MeaningsBlock{%
	\ryumeaning%
		{\ruby[g]{細}{ほそ}い}{}{}%
	}%
	}%
	{\tailPrint[]}%
\entry
	{\headPrint[kana={ぱちず\tl{˥}},ipa={pḁtʃidzu},pos={名},acc={H}]}%
	{%
		\MeaningsBlock{%
	\ryumeaning%
		{八十}{}{}%
	}%
	}%
	{\tailPrint[]}%
\entry
	{\headPrint[kana={ばちばた\tl{˩˥}},ipa={batʃibata},pos={名},acc={R}]}%
	{%
		\MeaningsBlock{%
	\ryumeaning%
		{小腸}{「細い腸」の義}{}%
	}%
	}%
	{\tailPrint[]}%
\entry
	{\headPrint[kana={ぱちま\tl{˥˩}はん},ipa={pḁtʃimahaɴ},pos={形},para={ぱちまへぬ},acc={F}]}%
	{%
		\MeaningsBlock{%
	\ryumeaning%
		{\ruby[g]{眩}{まぶ}しい}{}{}%
	}%
	}%
	{\tailPrint[]}%
\entry
	{\headPrint[kana={ぱちまはん},ipa={pḁtʃimahaɴ},pos={形}]}%
	{%
		\MeaningsBlock{%
	\ryumeaning%
		{眩しい}{}{}%
	}%
	}%
	{\tailPrint[]}%
\entry
	{\headPrint[kana={ぱちみがすん\tl{˥}},ipa={pḁtʃimigasuɴ},pos={動},rui={2},para={ぱちみがはぬ},acc={H}]}%
	{%
		\MeaningsBlock{%
	\ryumeaning%
		{平手で打つ}{平手打ちを食わせる}{}%
	}%
	}%
	{\tailPrint[]}%
\entry
	{\headPrint[kana={ぱちるま},ipa={pḁtʃiruma},pos={名}]}%
	{%
		\MeaningsBlock{%
	\ryumeaning%
		{波照間}{波照間島。波照間・白保での呼び方}{}%
	}%
	}%
	{\tailPrint[]}%
\entry
	{\headPrint[kana={ぱちるん\tl{˥}},ipa={pḁtʃiruɴ},pos={動},rui={1},para={ぱちらぬ},acc={H}]}%
	{%
		\MeaningsBlock{%
	\ryumeaning%
		{（不吉を）祓う}{}{}%
	}%
	}%
	{\tailPrint[]}%
\entry
	{\headPrint[kana={ぱちるん\tl{˥}},ipa={pḁtʃiruɴ},pos={動},rui={1},para={ぱちらぬ},acc={H}]}%
	{%
		\MeaningsBlock{%
	\ryumeaning%
		{（着物を）脱ぐ}{}{}%
	}%
	}%
	{\tailPrint[]}%
\entry
	{\headPrint[kana={ばつぁ\tl{˩˥}},ipa={batsḁ},pos={名},acc={R}]}%
	{%
		\MeaningsBlock{%
	\ryumeaning%
		{罰。罪。咎}{}{}%
	}%
	}%
	{\tailPrint[]}%
\entry
	{\headPrint[kana={ばつぁかぷん},ipa={batsakḁpuɴ},pos={句},para={ばつぁかぱぬ}]}%
	{%
		\MeaningsBlock{%
	\ryumeaning%
		{罰を被る。咎を被る}{}{}%
	}%
	}%
	{\tailPrint[]}%
\entry
	{\headPrint[kana={ばつぁみるん\tl{˩˥}},ipa={batsamiruɴ},pos={動},rui={3},para={ばつぁむぬ},acc={R}]}%
	{%
		\MeaningsBlock{%
	\ryumeaning%
		{（牛馬などの動物を）繋ぐ。繋ぎとめる}{}{}%
	}%
	}%
	{\tailPrint[]}%
\entry
	{\headPrint[kana={ぱつぁむん\tl{˥}},ipa={pḁtsamuɴ},pos={動},rui={1},para={ぱつぁまぬ},acc={H}]}%
	{%
		\MeaningsBlock{%
	\ryumeaning%
		{挟む。（箸などを）掴む}{}{}%
	}%
	}%
	{\tailPrint[]}%
\entry
	{\headPrint[kana={ぱつぁん\tl{˥}},ipa={patsaɴ},pos={名},acc={H}]}%
	{%
		\MeaningsBlock{%
	\ryumeaning%
		{\ruby[g]{鋏}{はさみ}}{}{}%
	}%
	}%
	{\tailPrint[]}%
\entry
	{\headPrint[kana={ばっくるん\tl{˩˥}},ipa={bakkuruɴ},pos={動},rui={1},para={ばっくらぬ},acc={R}]}%
	{%
		\MeaningsBlock{%
	\ryumeaning%
		{からかう。馬鹿にする}{}{}%
	}%
	}%
	{\tailPrint[]}%
\entry
	{\headPrint[kana={はっさみよー},ipa={hassamijoː},pos={感}]}%
	{%
		\MeaningsBlock{%
	\ryumeaning%
		{なんたることか}{}{}%
	}%
	}%
	{\tailPrint[]}%
\entry
	{\headPrint[kana={ばっしむぬ},ipa={baʃʃimunu},pos={名}]}%
	{%
		\MeaningsBlock{%
	\ryumeaning%
		{忘れ物}{}{}%
	}%
	}%
	{\tailPrint[]}%
\entry
	{\headPrint[kana={ばっし\tl{˥˩}めさ\tl{˩˥}はん},ipa={baʃʃimesahaɴ},pos={形},acc={F+R}]}%
	{%
		\MeaningsBlock{%
	\ryumeaning%
		{忘れっぽい}{}{}%
	}%
	}%
	{\tailPrint[]}%
\entry
	{\headPrint[kana={ばっしるん\tl{˥˩}},ipa={baʃʃiruɴ},pos={動},rui={3},para={ばっすぬ},acc={F}]}%
	{%
		\MeaningsBlock{%
	\ryumeaning%
		{忘れる}{}{}%
	}%
	}%
	{\tailPrint[]}%
\entry
	{\headPrint[kana={ばっしん\tl{˥˩}},ipa={baʃʃiɴ},pos={動},rui={3},para={ばっすぬ},acc={F}]}%
	{%
		\MeaningsBlock{%
	\ryumeaning%
		{忘れる}{}{}%
	}%
	}%
	{\tailPrint[]}%
\entry
	{\headPrint[kana={ぱったらげー},ipa={pattarageː},pos={名}]}%
	{%
		\MeaningsBlock{%
	\ryumeaning%
		{てんてこ舞い}{}{}%
	}%
	}%
	{\tailPrint[]}%
\entry
	{\headPrint[kana={ばっぺー\tl{˩˥}},ipa={bappeː},pos={名},acc={R}]}%
	{%
		\MeaningsBlock{%
	\ryumeaning%
		{間違い}{}{}%
	}%
	}%
	{\tailPrint[]}%
\entry
	{\headPrint[kana={ばっぺー\ws{}すん},ipa={bappeː suɴ},pos={句}]}%
	{%
		\MeaningsBlock{%
	\ryumeaning%
		{間違う。誤る}{}{}%
	}%
	}%
	{\tailPrint[]}%
\entry
	{\headPrint[kana={はてろー\tl{˥˩}},ipa={hateroː},pos={名},acc={F}]}%
	{%
		\MeaningsBlock{%
	\ryumeaning%
		{波照間。波照間島}{石垣市での呼び方}{}%
	}%
	}%
	{\tailPrint[]}%
\entry
	{\headPrint[kana={ぱてろーま\tl{˥˩}},ipa={pateroːma},pos={名},acc={F}]}%
	{%
		\MeaningsBlock{%
	\ryumeaning%
		{波照間。波照間島}{石垣市や民謡での呼び方}{}%
	}%
	}%
	{\tailPrint[]}%
\entry
	{\headPrint[kana={ぱとーま\tl{˥}},ipa={patoːma},pos={名},acc={H}]}%
	{%
		\MeaningsBlock{%
	\ryumeaning%
		{\ruby[g]{鳩間}{はとま}。鳩間島}{八重山諸島の島の名}{}%
	}%
	}%
	{\tailPrint[]}%
\entry
	{\headPrint[kana={ぱとん\tl{˥}},ipa={pḁtoɴ},pos={名},acc={H}]}%
	{%
		\MeaningsBlock{%
	\ryumeaning%
		{鳩}{}{}%
	}%
	}%
	{\tailPrint[]}%
\entry
	{\headPrint[kana={はな\tl{˥}},ipa={hana},pos={句},acc={H}]}%
	{%
		\MeaningsBlock{%
	\ryumeaning%
		{あそこに。あちらに}{}{}%
	}%
	}%
	{\tailPrint[]}%
\entry
	{\headPrint[kana={ぱな\tl{˥}},ipa={pḁna},pos={名},acc={H}]}%
	{%
		\MeaningsBlock{%
	\ryumeaning%
		{花}{}{}%
	}%
	}%
	{\tailPrint[]}%
\entry
	{\headPrint[kana={ぱな\tl{˥˩}},ipa={pḁna},pos={名},acc={F}]}%
	{%
		\MeaningsBlock{%
	\ryumeaning%
		{\ruby[g]{鼻}{はな}}{}{}%
	}%
	}%
	{\tailPrint[]}%
\entry
	{\headPrint[kana={ぱな\tl{˥˩}},ipa={pḁna},pos={名},acc={F}]}%
	{%
		\MeaningsBlock{%
	\ryumeaning%
		{先。崎。岬}{}{}%
	}%
	}%
	{\tailPrint[]}%
\entry
	{\headPrint[kana={ぱないぎ\tl{˥}},ipa={pḁna.igi},pos={名},acc={H}]}%
	{%
		\MeaningsBlock{%
	\ryumeaning%
		{花生け。花瓶}{}{}%
	}%
	}%
	{\tailPrint[]}%
\entry
	{\headPrint[kana={ぱなぐみ\tl{˥}},ipa={pḁnagumi},pos={名},acc={H}]}%
	{%
		\MeaningsBlock{%
	\ryumeaning%
		{供米}{神前に供える米}{}%
	}%
	}%
	{\tailPrint[note={直訳「花米」}]}%
\entry
	{\headPrint[kana={ぱなし\tl{˥}},ipa={pḁnaʃi},pos={名},acc={H}]}%
	{%
		\MeaningsBlock{%
	\ryumeaning%
		{話}{}{}%
	}%
	}%
	{\tailPrint[]}%
\entry
	{\headPrint[kana={ぱなしき\tl{˥˩}},ipa={pḁnaʃi̥ki},pos={名},acc={F}]}%
	{%
		\MeaningsBlock{%
	\ryumeaning%
		{\ruby[g]{病}{やまい}。病気}{}{}%
	}%
	}%
	{\tailPrint[]}%
\entry
	{\headPrint[kana={ぱなしきぬ\ws{}はやるん},ipa={panaʃi̥kinu hajaruɴ},pos={句}]}%
	{%
		\MeaningsBlock{%
	\ryumeaning%
		{風邪が流行る}{}{}%
	}%
	}%
	{\tailPrint[]}%
\entry
	{\headPrint[kana={ぱなじな\tl{˥}},ipa={pḁnadʒina},pos={名},acc={H}]}%
	{%
		\MeaningsBlock{%
	\ryumeaning%
		{（牛の）鼻綱}{}{}%
	}%
	}%
	{\tailPrint[]}%
\entry
	{\headPrint[kana={ぱなすん\tl{˥}},ipa={pḁnasuɴ},pos={動},rui={2},para={ぱなはぬ},acc={H}]}%
	{%
		\MeaningsBlock{%
	\ryumeaning%
		{離す}{}{}%
	}%
	}%
	{\tailPrint[]}%
\entry
	{\headPrint[kana={ぱなすん\tl{˥}},ipa={pḁnasuɴ},pos={動},rui={2},para={ぱなさぬ},acc={H}]}%
	{%
		\MeaningsBlock{%
	\ryumeaning%
		{話す。語る}{}{}%
	}%
	}%
	{\tailPrint[]}%
\entry
	{\headPrint[kana={ぱなた\tl{˥}},ipa={pḁnata},pos={名},acc={H}]}%
	{%
		\MeaningsBlock{%
	\ryumeaning%
		{先。先端}{}{}%
	}%
	}%
	{\tailPrint[]}%
\entry
	{\headPrint[kana={ぱなだり\tl{˥˩}},ipa={pḁnadari},pos={名},acc={F}]}%
	{%
		\MeaningsBlock{%
	\ryumeaning%
		{鼻汁。\ruby[g]{青}{あお}\ruby[g]{洟}{ばな}}{}{}%
	}%
	}%
	{\tailPrint[]}%
\entry
	{\headPrint[kana={ぱなたるん\tl{˥}},ipa={pḁnataruɴ},pos={動},rui={1},para={ぱなたらぬ},acc={H}]}%
	{%
		\MeaningsBlock{%
	\ryumeaning%
		{太る。肥える。肥満になる}{}{}%
	}%
	}%
	{\tailPrint[]}%
\entry
	{\headPrint[kana={\josho{}ぱなぬ\kako{}\ws{}ふぁー},ipa={pḁnanu faː},pos={句},acc={H L}]}%
	{%
		\MeaningsBlock{%
	\ryumeaning%
		{（御嶽の）\ruby[g]{氏子}{うじこ}}{お嶽の祭祀にかかわる血族の人々}{}%
	}%
	}%
	{\tailPrint[]}%
\entry
	{\headPrint[kana={ぱな\tl{˥˩}ぴぃすん\tl{˥}},ipa={pḁnapi̥suɴ},pos={句},para={ぱなぴさぬ},acc={F+H}]}%
	{%
		\MeaningsBlock{%
	\ryumeaning%
		{くしゃみをする}{}{}%
	}%
	}%
	{\tailPrint[]}%
\entry
	{\headPrint[kana={ぱな\tl{˥˩}ふき\tl{˩˥}},ipa={pḁnafu̥ki},pos={名},acc={F+R}]}%
	{%
		\MeaningsBlock{%
	\ryumeaning%
		{\ruby[g]{鼾}{ひびき}}{}{}%
	}%
	}%
	{\tailPrint[]}%
\entry
	{\headPrint[kana={はなやかすん\tl{˥}},ipa={hanajakasuɴ},pos={動},rui={2b},para={はなやかさぬ},acc={H}]}%
	{%
		\MeaningsBlock{%
	\ryumeaning%
		{賑やかに騒ぐ。賑わす}{}{}%
	}%
	}%
	{\tailPrint[]}%
\entry
	{\headPrint[kana={ぱなり\tl{˥}},ipa={pḁnari},pos={名},acc={H}]}%
	{%
		\MeaningsBlock{%
	\ryumeaning%
		{離れ。新城島}{新城島の略称}{}%
	}%
	}%
	{\tailPrint[]}%
\entry
	{\headPrint[kana={ぱなりぬ\ws{}すま\tl{˥}},ipa={pḁnarinu su̥ma},pos={名},acc={H}]}%
	{%
		\MeaningsBlock{%
	\ryumeaning%
		{離れの島。新城島}{八重山諸島の島の名。直訳「離島」}{}%
	}%
	}%
	{\tailPrint[]}%
\entry
	{\headPrint[kana={ぱなりるん\tl{˥}},ipa={pḁnariruɴ},pos={動},rui={3},para={ぱなるぬ},acc={H}]}%
	{%
		\MeaningsBlock{%
	\ryumeaning%
		{離れる}{}{}%
	\ryumeaning%
		{別れる}{}{}%
	\ryumeaning%
		{剥がれる}{}{}%
	}%
	}%
	{\tailPrint[]}%
\entry
	{\headPrint[kana={ぱなりん\tl{˥}},ipa={pḁnariɴ},pos={動},rui={3},para={ぱなるぬ},acc={H}]}%
	{%
		\MeaningsBlock{%
	\ryumeaning%
		{離れる}{}{}%
	}%
	}%
	{\tailPrint[]}%
\entry
	{\headPrint[kana={ぱに\tl{˥˩}},ipa={pḁni},pos={名},acc={F}]}%
	{%
		\MeaningsBlock{%
	\ryumeaning%
		{羽。翼。鰭}{}{}%
	}%
	}%
	{\tailPrint[]}%
\entry
	{\headPrint[kana={はにかいすん\tl{˥}},ipa={hanikaisuɴ},pos={動},acc={H}]}%
	{%
		\MeaningsBlock{%
	\ryumeaning%
		{（物に突き当たって）跳ね返す}{}{}%
	}%
	}%
	{\tailPrint[]}%
\entry
	{\headPrint[kana={はにかいるん\tl{˥}},ipa={hanikairuɴ},pos={動},rui={1},para={はにかいらぬ},acc={H}]}%
	{%
		\MeaningsBlock{%
	\ryumeaning%
		{跳ね返る}{}{}%
	}%
	}%
	{\tailPrint[]}%
\entry
	{\headPrint[kana={ばぬ\tl{˩˥}},ipa={banu},pos={名},acc={R}]}%
	{%
		\MeaningsBlock{%
	\ryumeaning%
		{\ruby[g]{我}{われ}。私。\ruby[g]{俺}{おれ}}{}{}%
	}%
	}%
	{\tailPrint[]}%
\entry
	{\headPrint[kana={ば\tl{˥}ひー\tl{˩˥}},ipa={bahiː},pos={名},acc={H+R}]}%
	{%
		\MeaningsBlock{%
	\ryumeaning%
		{私の家}{}{}%
	}%
	}%
	{\tailPrint[]}%
\entry
	{\headPrint[kana={ぱぴる\tl{˥}},ipa={pḁpiru},pos={名},acc={H}]}%
	{%
		\MeaningsBlock{%
	\ryumeaning%
		{蝶。蛾}{}{}%
	}%
	}%
	{\tailPrint[]}%
\entry
	{\headPrint[kana={ぱま\tl{˥}},ipa={pḁma},pos={名},acc={H}]}%
	{%
		\MeaningsBlock{%
	\ryumeaning%
		{浜。砂浜}{}{}%
	}%
	}%
	{\tailPrint[]}%
\entry
	{\headPrint[kana={ぱまかん\tl{˥}},ipa={pḁmakaɴ},pos={名},acc={H}]}%
	{%
		\MeaningsBlock{%
	\ryumeaning%
		{浜蟹}{}{}%
	}%
	}%
	{\tailPrint[]}%
\entry
	{\headPrint[kana={\josho{}ぱま\kako{}ふちぃ},ipa={pḁmafu̥tʃi},pos={名},acc={H+L}]}%
	{%
		\MeaningsBlock{%
	\ryumeaning%
		{なぎさ。礒辺}{}{}%
	}%
	}%
	{\tailPrint[]}%
\entry
	{\headPrint[kana={ぱまるん\tl{˥}},ipa={pḁmaruɴ},pos={動},rui={1},para={ぱまらぬ},acc={H}]}%
	{%
		\MeaningsBlock{%
	\ryumeaning%
		{はまる}{}{}%
	}%
	}%
	{\tailPrint[]}%
\entry
	{\headPrint[kana={ばみがすん\tl{˥˩}},ipa={bamigasuɴ},pos={動},rui={2},para={ばみがはぬ},acc={F}]}%
	{%
		\MeaningsBlock{%
	\ryumeaning%
		{殴る。ぶん殴る}{}{}%
	}%
	}%
	{\tailPrint[]}%
\entry
	{\headPrint[kana={ぱみるん\tl{˥}},ipa={pḁmiruɴ},pos={動},rui={1},acc={H}]}%
	{%
		\MeaningsBlock{%
	\ryumeaning%
		{はめる。はめ込む。当てはめる}{}{}%
	}%
	}%
	{\tailPrint[]}%
\entry
	{\headPrint[kana={ぱめー\tl{˥˩}},ipa={pameː},pos={名},acc={F}]}%
	{%
		\MeaningsBlock{%
	\ryumeaning%
		{食糧}{}{}%
	}%
	}%
	{\tailPrint[note={飯米からか}]}%
\entry
	{\headPrint[kana={はやーるん\tl{˥}},ipa={hajaːruɴ},pos={動},rui={1},para={はやーらぬ},acc={H}]}%
	{%
		\MeaningsBlock{%
	\ryumeaning%
		{流行る。流行する}{}{}%
	}%
	}%
	{\tailPrint[]}%
\entry
	{\headPrint[kana={はやじに\tl{˥}\ws{}すん\tl{˥˩}},ipa={hajadʒini suɴ},pos={句},acc={H F}]}%
	{%
		\MeaningsBlock{%
	\ryumeaning%
		{早世する。早死する}{}{}%
	}%
	}%
	{\tailPrint[]}%
\entry
	{\headPrint[kana={ぱやすん\tl{˩˥}},ipa={pajasuɴ},pos={動},acc={R}]}%
	{%
		\MeaningsBlock{%
	\ryumeaning%
		{囃す}{}{}%
	}%
	}%
	{\tailPrint[]}%
\entry
	{\headPrint[kana={はやまり\tl{˥}},ipa={hajamari},pos={名},acc={H}]}%
	{%
		\MeaningsBlock{%
	\ryumeaning%
		{早生まれ}{同年でも早く生まれた者}{}%
	}%
	}%
	{\tailPrint[]}%
\entry
	{\headPrint[kana={はやまるん\tl{˥}},ipa={hajamaruɴ},pos={動},rui={1},para={はやまらぬ},acc={H}]}%
	{%
		\MeaningsBlock{%
	\ryumeaning%
		{早まる。早くなる}{}{}%
	}%
	}%
	{\tailPrint[]}%
\entry
	{\headPrint[kana={はやみるん\tl{˥}},ipa={hajamiruɴ},pos={動},acc={H}]}%
	{%
		\MeaningsBlock{%
	\ryumeaning%
		{早める。早くする}{}{}%
	}%
	}%
	{\tailPrint[]}%
\entry
	{\headPrint[kana={はやり\tl{˥}},ipa={hajari},pos={名},acc={H}]}%
	{%
		\MeaningsBlock{%
	\ryumeaning%
		{流行り。流行}{}{}%
	}%
	}%
	{\tailPrint[]}%
\entry
	{\headPrint[kana={はやりやん\tl{˥}},ipa={hajarijaɴ},pos={名},acc={H}]}%
	{%
		\MeaningsBlock{%
	\ryumeaning%
		{伝染病}{}{}%
	}%
	}%
	{\tailPrint[]}%
\entry
	{\headPrint[kana={ばら\tl{˩˥}},ipa={bara},pos={名},acc={R}]}%
	{%
		\MeaningsBlock{%
	\ryumeaning%
		{\ruby[g]{藁}{わら}}{}{}%
	}%
	}%
	{\tailPrint[]}%
\entry
	{\headPrint[kana={ぱら\tl{˥}},ipa={pḁra},pos={名},acc={H}]}%
	{%
		\MeaningsBlock{%
	\ryumeaning%
		{\ruby[g]{柱}{はしら}}{}{}%
	}%
	}%
	{\tailPrint[]}%
\entry
	{\headPrint[kana={はらごー\tl{˥}},ipa={haragoː},pos={名},acc={H}]}%
	{%
		\MeaningsBlock{%
	\ryumeaning%
		{魚の腹肉}{カツオの腹肉を指す}{}%
	}%
	}%
	{\tailPrint[]}%
\entry
	{\headPrint[kana={ばらざん\tl{˩˥}},ipa={baradzaɴ},pos={名},acc={R}]}%
	{%
		\MeaningsBlock{%
	\ryumeaning%
		{\ruby[g]{藁}{わら}\ruby[g]{算}{ざん}}{結縄記録。王府時代の記録法の一種}{}%
	}%
	}%
	{\tailPrint[]}%
\entry
	{\headPrint[kana={ばらじぃな\tl{˩˥}},ipa={baradʒina},pos={名},acc={R}]}%
	{%
		\MeaningsBlock{%
	\ryumeaning%
		{\ruby[g]{藁}{わら}縄}{藁束は\emb{ばらふた}}{}%
	}%
	}%
	{\tailPrint[]}%
\entry
	{\headPrint[kana={ぱらすん\tl{˥}},ipa={pḁrasuɴ},pos={動},rui={2},para={ぱらはぬ},acc={H}]}%
	{%
		\MeaningsBlock{%
	\ryumeaning%
		{払う。支払う}{}{}%
	}%
	}%
	{\tailPrint[]}%
\entry
	{\headPrint[kana={ぱらすん\tl{˥}},ipa={pḁrasuɴ},pos={動},rui={2},para={ぱらはぬ},acc={H}]}%
	{%
		\MeaningsBlock{%
	\ryumeaning%
		{走らせる}{}{}%
	\ryumeaning%
		{（血、汗などを）流す}{}{}%
	\ryumeaning%
		{注ぐ}{}{}%
	}%
	}%
	{\tailPrint[]}%
\entry
	{\headPrint[kana={ばらふた},ipa={barafu̥ta},pos={名}]}%
	{%
		\MeaningsBlock{%
	\ryumeaning%
		{\ruby[g]{藁}{わら}束}{}{}%
	}%
	}%
	{\tailPrint[]}%
\entry
	{\headPrint[kana={ばらふちぃ},ipa={barafu̥tʃi},pos={名}]}%
	{%
		\MeaningsBlock{%
	\ryumeaning%
		{\ruby[g]{藁}{わら}製の草鞋}{}{}%
	}%
	}%
	{\tailPrint[]}%
\entry
	{\headPrint[kana={はらみ\tl{˥}},ipa={harami},pos={名},acc={H}]}%
	{%
		\MeaningsBlock{%
	\ryumeaning%
		{（鰹など魚の）卵巣}{}{}%
	}%
	}%
	{\tailPrint[]}%
\entry
	{\headPrint[kana={ぱり\tl{˥}},ipa={pḁri},pos={名},acc={H}]}%
	{%
		\MeaningsBlock{%
	\ryumeaning%
		{針}{釣針は\emb{じばり}}{}%
	}%
	}%
	{\tailPrint[]}%
\entry
	{\headPrint[kana={ばりあてぃん\tl{˥˩}},ipa={bari.atiɴ},pos={動},acc={F}]}%
	{%
		\MeaningsBlock{%
	\ryumeaning%
		{割り当てる。振り当てる}{}{}%
	}%
	}%
	{\tailPrint[]}%
\entry
	{\headPrint[kana={ばりくだくん\tl{˩˥}},ipa={bariku̥dakuɴ},pos={動},rui={1},para={ばりくだかぬ},acc={R}]}%
	{%
		\MeaningsBlock{%
	\ryumeaning%
		{割り砕く}{}{}%
	}%
	}%
	{\tailPrint[]}%
\entry
	{\headPrint[kana={ばり\ws{}っしん},ipa={bari ʃʃiɴ},pos={句}]}%
	{%
		\MeaningsBlock{%
	\ryumeaning%
		{割ってしまう}{}{}%
	}%
	}%
	{\tailPrint[]}%
\entry
	{\headPrint[kana={ぱりみじぃ\tl{˥}},ipa={pḁrimidʒi},pos={名},acc={H}]}%
	{%
		\MeaningsBlock{%
	\ryumeaning%
		{湧水。泉}{}{}%
	}%
	}%
	{\tailPrint[]}%
\entry
	{\headPrint[kana={ぱりゃん\tl{˥}},ipa={pḁrjaɴ},pos={名},acc={H}]}%
	{%
		\MeaningsBlock{%
	\ryumeaning%
		{（魚の）卵巣。（蟹の）卵}{}{}%
	}%
	}%
	{\tailPrint[]}%
\entry
	{\headPrint[kana={ぱりん\tl{˥}},ipa={pḁriɴ},pos={動},rui={3},para={ぱるぬ},acc={H}]}%
	{%
		\MeaningsBlock{%
	\ryumeaning%
		{（天気が）晴れる}{}{}%
	}%
	}%
	{\tailPrint[]}%
\entry
	{\headPrint[kana={ぱるむん\tl{˥}},ipa={pḁrumuɴ},pos={動},rui={1},para={ぱるまぬ},acc={H}]}%
	{%
		\MeaningsBlock{%
	\ryumeaning%
		{孕む。妊娠する}{}{}%
	}%
	}%
	{\tailPrint[]}%
\entry
	{\headPrint[kana={ぱるん\tl{˥}},ipa={pḁruɴ},pos={動},rui={1},para={ぱらぬ},acc={H}]}%
	{%
		\MeaningsBlock{%
	\ryumeaning%
		{走る。行く。発つ}{}{}%
	}%
	}%
	{\tailPrint[]}%
\entry
	{\headPrint[kana={ぱるん\tl{˥˩}},ipa={pḁruɴ},pos={動},rui={1},para={ぱらぬ},acc={F}]}%
	{%
		\MeaningsBlock{%
	\ryumeaning%
		{張る}{}{}%
	}%
	}%
	{\tailPrint[]}%
\entry
	{\headPrint[kana={ばん\tl{˩˥}},ipa={baɴ},pos={名},acc={R}]}%
	{%
		\MeaningsBlock{%
	\ryumeaning%
		{番。順番。見張りする}{}{}%
	}%
	}%
	{\tailPrint[]}%
\entry
	{\headPrint[kana={ぱん\tl{˥}},ipa={paɴ},pos={名},acc={H}]}%
	{%
		\MeaningsBlock{%
	\ryumeaning%
		{祈詞。呪詞}{祈祷などの祈りの言葉}{}%
	}%
	}%
	{\tailPrint[]}%
\entry
	{\headPrint[kana={ぱん\tl{˥}},ipa={paɴ},pos={名},acc={H}]}%
	{%
		\MeaningsBlock{%
	\ryumeaning%
		{刃。刃先}{}{}%
	}%
	}%
	{\tailPrint[]}%
\entry
	{\headPrint[kana={ぱん\tl{˥}},ipa={paɴ},pos={名},acc={H}]}%
	{%
		\MeaningsBlock{%
	\ryumeaning%
		{足}{}{}%
	}%
	}%
	{\tailPrint[]}%
\entry
	{\headPrint[kana={ぱん\tl{˥}},ipa={paɴ},pos={名},acc={H}]}%
	{%
		\MeaningsBlock{%
	\ryumeaning%
		{歯}{}{}%
	}%
	}%
	{\tailPrint[note={単独では多く\emb{ふちぃぬ\ws{}ぱん}を使う。複合語の構成要素となる}]}%
\entry
	{\headPrint[kana={ぱん\tl{˥˩}},ipa={paɴ},pos={名},acc={F}]}%
	{%
		\MeaningsBlock{%
	\ryumeaning%
		{印。印鑑}{}{}%
	}%
	}%
	{\tailPrint[]}%
\entry
	{\headPrint[kana={ぱんがま\tl{˥˩}},ipa={pamgama},pos={名},acc={F}]}%
	{%
		\MeaningsBlock{%
	\ryumeaning%
		{\ruby[g]{羽釜}{はがま}}{}{}%
	}%
	}%
	{\tailPrint[]}%
\entry
	{\headPrint[kana={ぱんきん\tl{˥}},ipa={paŋkiɴ},pos={動},rui={3},para={ぱんくぬ},acc={H}]}%
	{%
		\MeaningsBlock{%
	\ryumeaning%
		{捲れる}{}{}%
	}%
	}%
	{\tailPrint[]}%
\entry
	{\headPrint[kana={ぱんくん\tl{˥}},ipa={paŋkuɴ},pos={動},rui={1},para={ぱんかぬ},acc={H}]}%
	{%
		\MeaningsBlock{%
	\ryumeaning%
		{\ruby[g]{弾}{はじ}く}{}{}%
	}%
	}%
	{\tailPrint[]}%
\entry
	{\headPrint[kana={ぱんじ\tl{˥˩}},ipa={pandʒi},pos={名},acc={F}]}%
	{%
		\MeaningsBlock{%
	\ryumeaning%
		{ハゼノキ}{\emb{はじまき}「ハゼノキにかぶれる」}{}%
	}%
	}%
	{\tailPrint[]}%
\entry
	{\headPrint[kana={ばんじぃん\tl{˥˩}},ipa={bandʒiɴ},pos={名},acc={F}]}%
	{%
		\MeaningsBlock{%
	\ryumeaning%
		{盛り。真っ盛り}{}{}%
	}%
	}%
	{\tailPrint[]}%
\entry
	{\headPrint[kana={ばんしゅる\tl{˩˥}},ipa={banʃuru},pos={名},acc={R}]}%
	{%
		\MeaningsBlock{%
	\ryumeaning%
		{バンザクロ}{野生の果物。\emb{とっきん}ともいう}{}%
	}%
	}%
	{\tailPrint[]}%
\entry
	{\headPrint[kana={はんじょー\tl{˥}},ipa={handʒoː},pos={名},acc={H}]}%
	{%
		\MeaningsBlock{%
	\ryumeaning%
		{繁昌。繁栄}{}{}%
	}%
	}%
	{\tailPrint[]}%
\entry
	{\headPrint[kana={ばんぞーがに\tl{˩˥}},ipa={bandzoːgani},pos={名},acc={R}]}%
	{%
		\MeaningsBlock{%
	\ryumeaning%
		{\ruby[g]{曲尺}{かねじゃく}}{}{}%
	}%
	}%
	{\tailPrint[]}%
\entry
	{\headPrint[kana={ぱんた\tl{˥˩}},ipa={panta},pos={名},acc={F}]}%
	{%
		\MeaningsBlock{%
	\ryumeaning%
		{端。端っこ}{}{}%
	}%
	}%
	{\tailPrint[]}%
\entry
	{\headPrint[kana={ぱんだー\tl{˥}},ipa={pandaː},pos={名},acc={H}]}%
	{%
		\MeaningsBlock{%
	\ryumeaning%
		{おばさん達}{神司の集団を指すことも}{}%
	}%
	}%
	{\tailPrint[]}%
\entry
	{\headPrint[kana={ぱんたっさ\tl{˥}はん},ipa={pantassahaɴ},pos={形},para={ぱんたっさへぬ},acc={H}]}%
	{%
		\MeaningsBlock{%
	\ryumeaning%
		{忙しい。多忙だ}{}{}%
	}%
	}%
	{\tailPrint[]}%
\entry
	{\headPrint[kana={ぱんたるん\tl{˥}},ipa={pantaruɴ},pos={動},rui={3},para={ぱんたるぬ},acc={H}]}%
	{%
		\MeaningsBlock{%
	\ryumeaning%
		{太る。肥える。肥満になる}{}{}%
	}%
	}%
	{\tailPrint[]}%
\entry
	{\headPrint[kana={ぱんちるん\tl{˥˩}},ipa={pantʃiruɴ},pos={動},rui={3},para={ぱんつぬ},acc={F}]}%
	{%
		\MeaningsBlock{%
	\ryumeaning%
		{\ruby[g]{解}{ほど}ける}{}{}%
	}%
	}%
	{\tailPrint[]}%
\entry
	{\headPrint[kana={ぱんちん\tl{˥˩}},ipa={pantʃiɴ},pos={動},rui={3},para={ぱんつぬ},acc={F}]}%
	{%
		\MeaningsBlock{%
	\ryumeaning%
		{外れる}{}{}%
	}%
	}%
	{\tailPrint[]}%
\entry
	{\headPrint[kana={ぱんつん},ipa={pantsuɴ},pos={動}]}%
	{%
		\MeaningsBlock{%
	\ryumeaning%
		{外す}{}{}%
	}%
	}%
	{\tailPrint[]}%
\entry
	{\headPrint[kana={ばんどー\tl{˥˩}},ipa={bandoː},pos={名},acc={F}]}%
	{%
		\MeaningsBlock{%
	\ryumeaning%
		{口の大きな水甕}{}{}%
	}%
	}%
	{\tailPrint[]}%
\entry
	{\headPrint[kana={ぱんにち\tl{˥˩}},ipa={pannitʃi},pos={名},acc={F}]}%
	{%
		\MeaningsBlock{%
	\ryumeaning%
		{半日}{}{}%
	}%
	}%
	{\tailPrint[]}%
\entry
	{\headPrint[kana={ぱんぬ\ws{}こー},ipa={pannu koː},pos={名}]}%
	{%
		\MeaningsBlock{%
	\ryumeaning%
		{足の甲}{}{}%
	}%
	}%
	{\tailPrint[]}%
\entry
	{\headPrint[kana={ぱんぬ\ws{}にふつぁー},ipa={pannu nifu̥tsaː},pos={句}]}%
	{%
		\MeaningsBlock{%
	\ryumeaning%
		{足が遅い}{}{}%
	}%
	}%
	{\tailPrint[]}%
\entry
	{\headPrint[kana={ぱんぬ\ws{}ぴぃさ},ipa={pannu pi̥sa},pos={名}]}%
	{%
		\MeaningsBlock{%
	\ryumeaning%
		{足の甲}{}{}%
	}%
	}%
	{\tailPrint[]}%
\entry
	{\headPrint[kana={ぱんぬ\ws{}ふき},ipa={pannu fu̥ki},pos={名}]}%
	{%
		\MeaningsBlock{%
	\ryumeaning%
		{足首}{}{}%
	}%
	}%
	{\tailPrint[]}%
\entry
	{\headPrint[kana={ぱんぶん\tl{˥}},ipa={pambuɴ},pos={名},acc={H}]}%
	{%
		\MeaningsBlock{%
	\ryumeaning%
		{半分}{}{}%
	}%
	}%
	{\tailPrint[]}%
\entry
	{\headPrint[kana={ぱんぶんばぎ\tl{˥}},ipa={pambumbagi},pos={名},acc={H}]}%
	{%
		\MeaningsBlock{%
	\ryumeaning%
		{半分分け}{}{}%
	}%
	}%
	{\tailPrint[]}%
\entry
	{\headPrint[kana={ぱんべー\tl{˥}},ipa={pambeː},pos={名},acc={H}]}%
	{%
		\MeaningsBlock{%
	\ryumeaning%
		{天ぷら}{}{}%
	}%
	}%
	{\tailPrint[]}%
\LetterHead{ひ}
\entry
	{\headPrint[kana={ひー\tl{˩˥}},ipa={hiː},pos={名},acc={R}]}%
	{%
		\MeaningsBlock{%
	\ryumeaning%
		{家。家屋}{}{}%
	}%
	}%
	{\tailPrint[]}%
\entry
	{\headPrint[kana={びー\tl{˩˥}},ipa={biː},pos={名},acc={R}]}%
	{%
		\MeaningsBlock{%
	\ryumeaning%
		{\ruby[g]{亥}{い}}{十二支の亥}{}%
	}%
	}%
	{\tailPrint[]}%
\entry
	{\headPrint[kana={ぴー\tl{˥}},ipa={piː},pos={名},acc={H}]}%
	{%
		\MeaningsBlock{%
	\ryumeaning%
		{火。火事}{}{}%
	}%
	}%
	{\tailPrint[]}%
\entry
	{\headPrint[kana={ぴー\tl{˥}},ipa={piː},pos={名},acc={H}]}%
	{%
		\MeaningsBlock{%
	\ryumeaning%
		{\ruby[g]{干}{ひ}\ruby[g]{瀬}{せ}}{干潮時に干上がる瀬}{}%
	}%
	}%
	{\tailPrint[]}%
\entry
	{\headPrint[kana={ぴー\tl{˥}},ipa={piː},pos={名},acc={H}]}%
	{%
		\MeaningsBlock{%
	\ryumeaning%
		{女陰}{}{}%
	}%
	}%
	{\tailPrint[]}%
\entry
	{\headPrint[kana={びーうしぃ\tl{˥˩}},ipa={biː.uʃi},pos={名},acc={F}]}%
	{%
		\MeaningsBlock{%
	\ryumeaning%
		{雄牛}{大きい雄牛は\emb{ぐちぇー}と言う}{}%
	}%
	}%
	{\tailPrint[]}%
\entry
	{\headPrint[kana={びーがーら\tl{˥˩}},ipa={biːgaːra},pos={名},acc={F}]}%
	{%
		\MeaningsBlock{%
	\ryumeaning%
		{雄瓦}{}{}%
	}%
	}%
	{\tailPrint[]}%
\entry
	{\headPrint[kana={ぴーかち\tl{˥}},ipa={piːkḁtʃi},pos={名},acc={H}]}%
	{%
		\MeaningsBlock{%
	\ryumeaning%
		{火風}{塩害のひどい台風}{}%
	}%
	}%
	{\tailPrint[]}%
\entry
	{\headPrint[kana={ひーすぅぬ\tl{˩˥}},ipa={hiːsu̥nu},pos={名},acc={R}]}%
	{%
		\MeaningsBlock{%
	\ryumeaning%
		{普段着}{}{}%
	}%
	}%
	{\tailPrint[]}%
\entry
	{\headPrint[kana={びーたこりるん\tl{˩˥}},ipa={biːtḁkoriruɴ},pos={動},rui={3},para={びーたこるぬ},acc={R}]}%
	{%
		\MeaningsBlock{%
	\ryumeaning%
		{酔いつぶれる}{}{}%
	}%
	}%
	{\tailPrint[]}%
\entry
	{\headPrint[kana={びーちゃー\tl{˩˥}},ipa={biːtʃaː},pos={名},acc={R}]}%
	{%
		\MeaningsBlock{%
	\ryumeaning%
		{酔っ払い}{}{}%
	}%
	}%
	{\tailPrint[]}%
\entry
	{\headPrint[kana={ぴーなん\tl{˥}},ipa={piːnaɴ},pos={名},acc={H}]}%
	{%
		\MeaningsBlock{%
	\ryumeaning%
		{火縄}{火種にする縄}{}%
	}%
	}%
	{\tailPrint[]}%
\entry
	{\headPrint[kana={\josho{}ぴーぬ\kako{}\ws{}ばり},ipa={piːnu bari},pos={句},acc={H L}]}%
	{%
		\MeaningsBlock{%
	\ryumeaning%
		{\ruby[g]{干}{ひ}\ruby[g]{瀬}{せ}の割れ目}{リーフの割れ目}{}%
	}%
	}%
	{\tailPrint[]}%
\entry
	{\headPrint[kana={ぴーみじぃ\tl{˥}},ipa={piːmidʒi},pos={名},acc={H}]}%
	{%
		\MeaningsBlock{%
	\ryumeaning%
		{冷や水}{}{}%
	}%
	}%
	{\tailPrint[]}%
\entry
	{\headPrint[kana={びーむぬ\tl{˥˩}},ipa={biːmunu},pos={名},acc={F}]}%
	{%
		\MeaningsBlock{%
	\ryumeaning%
		{雄}{}{}%
	}%
	}%
	{\tailPrint[]}%
\entry
	{\headPrint[kana={びーら},ipa={biːra},pos={名}]}%
	{%
		\MeaningsBlock{%
	\ryumeaning%
		{ひ弱。虚弱}{}{}%
	}%
	}%
	{\tailPrint[]}%
\entry
	{\headPrint[kana={びーらー},ipa={biːraː},pos={名}]}%
	{%
		\MeaningsBlock{%
	\ryumeaning%
		{病弱な子供}{}{}%
	}%
	}%
	{\tailPrint[]}%
\entry
	{\headPrint[kana={ぴーる},ipa={piːru},pos={名}]}%
	{%
		\MeaningsBlock{%
	\ryumeaning%
		{吉日。日選び}{}{}%
	}%
	}%
	{\tailPrint[]}%
\entry
	{\headPrint[kana={びーるん\tl{˩˥}},ipa={biːruɴ},pos={動},rui={1},para={びーらぬ},acc={R}]}%
	{%
		\MeaningsBlock{%
	\ryumeaning%
		{中毒する}{}{}%
	}%
	}%
	{\tailPrint[]}%
\entry
	{\headPrint[kana={ぴーるん\tl{˥}},ipa={piːruɴ},pos={動},rui={1},para={ぴらぬ},acc={H}]}%
	{%
		\MeaningsBlock{%
	\ryumeaning%
		{冷える}{}{}%
	}%
	}%
	{\tailPrint[]}%
\entry
	{\headPrint[kana={ぴぃさ\tl{˥˩}},ipa={pi̥sa},pos={名},acc={F}]}%
	{%
		\MeaningsBlock{%
	\ryumeaning%
		{足の甲}{}{}%
	}%
	}%
	{\tailPrint[]}%
\entry
	{\headPrint[kana={ぴぃさ\tl{˥˩}},ipa={pi̥sa},pos={名},acc={F}]}%
	{%
		\MeaningsBlock{%
	\ryumeaning%
		{屋根}{}{}%
	}%
	}%
	{\tailPrint[]}%
\entry
	{\headPrint[kana={ぴぃささしー\tl{˥}},ipa={pi̥sasḁʃiː},pos={名},acc={H}]}%
	{%
		\MeaningsBlock{%
	\ryumeaning%
		{人差し指}{}{}%
	}%
	}%
	{\tailPrint[]}%
\entry
	{\headPrint[kana={ぴぃさじま\tl{˥˩}},ipa={pi̥sadʒima},pos={名},acc={F}]}%
	{%
		\MeaningsBlock{%
	\ryumeaning%
		{平坦な島}{}{}%
	}%
	}%
	{\tailPrint[]}%
\entry
	{\headPrint[kana={ぴぃさすん\tl{˥}},ipa={pi̥sasuɴ},pos={動},rui={2},para={ぴさはぬ},acc={H}]}%
	{%
		\MeaningsBlock{%
	\ryumeaning%
		{（米飯を）炊く}{}{}%
	}%
	}%
	{\tailPrint[]}%
\entry
	{\headPrint[kana={ぴぃさ\tl{˥˩}はん},ipa={pi̥sahaɴ},pos={形},para={ぴしゃへぬ},acc={F}]}%
	{%
		\MeaningsBlock{%
	\ryumeaning%
		{薄い。薄めだ}{}{}%
	}%
	}%
	{\tailPrint[]}%
\entry
	{\headPrint[kana={ぴぃさんたり\tl{˥˩}},ipa={pi̥santari},pos={名},acc={F}]}%
	{%
		\MeaningsBlock{%
	\ryumeaning%
		{平たい。平べったぃ}{}{}%
	}%
	}%
	{\tailPrint[]}%
\entry
	{\headPrint[kana={\josho{}ぴぃしゃー\kako{}むぬ},ipa={pi̥ʃaːmunu},pos={名},acc={H+L}]}%
	{%
		\MeaningsBlock{%
	\ryumeaning%
		{拾い物}{}{}%
	}%
	}%
	{\tailPrint[]}%
\entry
	{\headPrint[kana={ぴぃす\tl{˥˩}},ipa={pi̥su},pos={名},acc={F}]}%
	{%
		\MeaningsBlock{%
	\ryumeaning%
		{昼間。日中}{}{}%
	}%
	}%
	{\tailPrint[]}%
\entry
	{\headPrint[kana={ぴぃすあとぅ},ipa={pi̥su.atu},pos={名}]}%
	{%
		\MeaningsBlock{%
	\ryumeaning%
		{昼後。午後}{}{}%
	}%
	}%
	{\tailPrint[]}%
\entry
	{\headPrint[kana={ぴぃすまぬっふぃ\tl{˥˩}},ipa={pi̥su̥manuffi},pos={名},acc={F}]}%
	{%
		\MeaningsBlock{%
	\ryumeaning%
		{昼寝}{}{}%
	}%
	}%
	{\tailPrint[]}%
\entry
	{\headPrint[kana={ぴぃすまり\tl{˥˩}},ipa={pi̥sumari},pos={名},acc={F}]}%
	{%
		\MeaningsBlock{%
	\ryumeaning%
		{真昼。正午}{}{}%
	}%
	}%
	{\tailPrint[]}%
\entry
	{\headPrint[kana={ぴぃすまりむぬ\tl{˩˥}},ipa={pi̥sumarimunu},pos={名},acc={R}]}%
	{%
		\MeaningsBlock{%
	\ryumeaning%
		{昼食。昼飯}{}{}%
	}%
	}%
	{\tailPrint[]}%
\entry
	{\headPrint[kana={ぴぃすまり\tl{˥˩}やしみ\tl{˩˥}},ipa={pi̥sumarijaʃi̥mi},pos={名},acc={F+R}]}%
	{%
		\MeaningsBlock{%
	\ryumeaning%
		{昼休み}{}{}%
	}%
	}%
	{\tailPrint[]}%
\entry
	{\headPrint[kana={ぴぃすん\tl{˥}},ipa={pi̥suɴ},pos={動},rui={2},para={ぴさぬ},acc={H}]}%
	{%
		\MeaningsBlock{%
	\ryumeaning%
		{（潮が）引く}{干潮になる}{}%
	}%
	}%
	{\tailPrint[]}%
\entry
	{\headPrint[kana={ぴぃすん\tl{˥}},ipa={pi̥suɴ},pos={動},rui={2},para={ぴさぬ},acc={H}]}%
	{%
		\MeaningsBlock{%
	\ryumeaning%
		{（屁を）ひる。（屁を）こく}{}{}%
	}%
	}%
	{\tailPrint[]}%
\entry
	{\headPrint[kana={ぴぃすん\tl{˥˩}},ipa={pi̥suɴ},pos={動},rui={2},para={ぴさぬ},acc={F}]}%
	{%
		\MeaningsBlock{%
	\ryumeaning%
		{拾う}{}{}%
	}%
	}%
	{\tailPrint[]}%
\entry
	{\headPrint[kana={ぴぃせー\tl{˥˩}},ipa={pi̥seː},pos={名},acc={F}]}%
	{%
		\MeaningsBlock{%
	\ryumeaning%
		{台地上の平地}{}{}%
	}%
	}%
	{\tailPrint[]}%
\entry
	{\headPrint[kana={ぴぃせー\tl{˥˩}},ipa={pi̥seː},pos={名},acc={F}]}%
	{%
		\MeaningsBlock{%
	\ryumeaning%
		{\ruby[g]{平}{ひら}\ruby[g]{得}{え}}{石垣島の集落名}{}%
	}%
	}%
	{\tailPrint[]}%
\entry
	{\headPrint[kana={ぴぃそ\tl{˥}はん},ipa={pi̥sohaɴ},pos={形},para={ぴそへぬ},acc={H}]}%
	{%
		\MeaningsBlock{%
	\ryumeaning%
		{広い}{}{}%
	}%
	}%
	{\tailPrint[]}%
\entry
	{\headPrint[kana={ぴぃとぅ\tl{˥˩}},ipa={pi̥tu},pos={名},acc={F}]}%
	{%
		\MeaningsBlock{%
	\ryumeaning%
		{人。人間}{}{}%
	\ryumeaning%
		{他人}{}{}%
	}%
	}%
	{\tailPrint[]}%
\entry
	{\headPrint[kana={ぴぃとぅいし\tl{˥}},ipa={pi̥tu.iʃi},pos={名},acc={H}]}%
	{%
		\MeaningsBlock{%
	\ryumeaning%
		{一息}{}{}%
	}%
	}%
	{\tailPrint[]}%
\entry
	{\headPrint[kana={ぴぃとぅかた\tl{˥}},ipa={pi̥tukḁta},pos={名},acc={H}]}%
	{%
		\MeaningsBlock{%
	\ryumeaning%
		{一方。片側}{}{}%
	}%
	}%
	{\tailPrint[]}%
\entry
	{\headPrint[kana={ぴぃとぅきぱ\tl{˥}},ipa={pi̥tuki̥pa},pos={名},acc={H}]}%
	{%
		\MeaningsBlock{%
	\ryumeaning%
		{一畝}{}{}%
	}%
	}%
	{\tailPrint[]}%
\entry
	{\headPrint[kana={ぴぃとぅごはすん},ipa={pi̥tugohasuɴ},pos={句},para={ひとぅごはさぬ}]}%
	{%
		\MeaningsBlock{%
	\ryumeaning%
		{人怖じする。人見知りをする}{}{}%
	}%
	}%
	{\tailPrint[]}%
\entry
	{\headPrint[kana={ぴぃとぅしぐとぅ\tl{˥}},ipa={pi̥tuʃigutu},pos={名},acc={H}]}%
	{%
		\MeaningsBlock{%
	\ryumeaning%
		{一仕事}{}{}%
	}%
	}%
	{\tailPrint[]}%
\entry
	{\headPrint[kana={ぴぃとぅすこーん},ipa={pi̥tusu̥koːɴ},pos={句},para={ぴとぅすかぬ}]}%
	{%
		\MeaningsBlock{%
	\ryumeaning%
		{人を使う}{}{}%
	}%
	}%
	{\tailPrint[]}%
\entry
	{\headPrint[kana={ぴぃとぅすな\tl{˥}},ipa={pi̥tusu̥na},pos={名},acc={H}]}%
	{%
		\MeaningsBlock{%
	\ryumeaning%
		{一品}{}{}%
	}%
	}%
	{\tailPrint[]}%
\entry
	{\headPrint[kana={ぴぃとぅだぎ\tl{˥˩}},ipa={pi̥tudagi},pos={名},acc={F}]}%
	{%
		\MeaningsBlock{%
	\ryumeaning%
		{人並に}{}{}%
	}%
	}%
	{\tailPrint[]}%
\entry
	{\headPrint[kana={ぴぃとぅたなぐん},ipa={pi̥tutḁnaguɴ},pos={句},para={ひとぅたながぬ}]}%
	{%
		\MeaningsBlock{%
	\ryumeaning%
		{人に頼る。人頼み}{}{}%
	}%
	}%
	{\tailPrint[]}%
\entry
	{\headPrint[kana={ぴぃとぅっし\tl{˥}},ipa={pi̥tuʃʃi},pos={名},acc={H}]}%
	{%
		\MeaningsBlock{%
	\ryumeaning%
		{\ruby[g]{一}{ひと}\ruby[g]{切}{き}れ}{}{}%
	}%
	}%
	{\tailPrint[]}%
\entry
	{\headPrint[kana={ぴぃとぅぴれー\tl{˥}},ipa={pi̥tupi̥reː},pos={名},acc={H}]}%
	{%
		\MeaningsBlock{%
	\ryumeaning%
		{付き合い。交際}{}{}%
	}%
	}%
	{\tailPrint[]}%
\entry
	{\headPrint[kana={ぴぃとぅむし\tl{˥}},ipa={pi̥tumuʃi},pos={名},acc={H}]}%
	{%
		\MeaningsBlock{%
	\ryumeaning%
		{一度。一回}{}{}%
	}%
	}%
	{\tailPrint[]}%
\entry
	{\headPrint[kana={ぴぃとぅゆー\tl{˥}},ipa={pi̥tujuː},pos={名},acc={H}]}%
	{%
		\MeaningsBlock{%
	\ryumeaning%
		{一晩。一夜}{}{}%
	}%
	}%
	{\tailPrint[]}%
\entry
	{\headPrint[kana={ぴぃとぅり\tl{˥}},ipa={pi̥turi},pos={名},acc={H}]}%
	{%
		\MeaningsBlock{%
	\ryumeaning%
		{独り。一人}{}{}%
	}%
	}%
	{\tailPrint[]}%
\entry
	{\headPrint[kana={ぴぃとぅりたま\tl{˥}},ipa={pi̥turitama},pos={名},acc={H}]}%
	{%
		\MeaningsBlock{%
	\ryumeaning%
		{一人っ子}{}{}%
	}%
	}%
	{\tailPrint[]}%
\entry
	{\headPrint[kana={ぴぃとぅりふち\tl{˥}},ipa={pi̥turifu̥tʃi},pos={名},acc={H}]}%
	{%
		\MeaningsBlock{%
	\ryumeaning%
		{独り言}{}{}%
	}%
	}%
	{\tailPrint[]}%
\entry
	{\headPrint[kana={ぴぃな\tl{˥}},ipa={pi̥na},pos={名},acc={H}]}%
	{%
		\MeaningsBlock{%
	\ryumeaning%
		{\ruby[g]{皺}{しわ}}{}{}%
	}%
	}%
	{\tailPrint[note={日本語の「襞（ひだ）」に対応}]}%
\entry
	{\headPrint[kana={ぴからすん\tl{˥}},ipa={pi̥karasuɴ},pos={動},rui={2},para={ぴからはぬ},acc={H}]}%
	{%
		\MeaningsBlock{%
	\ryumeaning%
		{光らせる}{磨いて光らせる}{}%
	}%
	}%
	{\tailPrint[]}%
\entry
	{\headPrint[kana={ぴかるん\tl{˥}},ipa={pi̥karuɴ},pos={動},rui={1},para={ぴからぬ},acc={H}]}%
	{%
		\MeaningsBlock{%
	\ryumeaning%
		{光る}{}{}%
	}%
	}%
	{\tailPrint[]}%
\entry
	{\headPrint[kana={ぴき\tl{˥˩}},ipa={pi̥ki},pos={名},acc={F}]}%
	{%
		\MeaningsBlock{%
	\ryumeaning%
		{竿秤}{}{}%
	}%
	}%
	{\tailPrint[]}%
\entry
	{\headPrint[kana={ぴき\tl{˥˩}},ipa={pi̥ki},pos={名},acc={F}]}%
	{%
		\MeaningsBlock{%
	\ryumeaning%
		{血筋。血統}{}{}%
	}%
	}%
	{\tailPrint[]}%
\entry
	{\headPrint[kana={ぴき\tl{˥˩}},ipa={pi̥ki},pos={名},acc={F}]}%
	{%
		\MeaningsBlock{%
	\ryumeaning%
		{引き}{引くこと}{}%
	}%
	}%
	{\tailPrint[]}%
\entry
	{\headPrint[kana={ひきあうん\tl{˥˩}},ipa={hi̥ki.auɴ},pos={動},para={ひきあわぬ},acc={F}]}%
	{%
		\MeaningsBlock{%
	\ryumeaning%
		{引き合う。相当する}{}{}%
	}%
	}%
	{\tailPrint[]}%
\entry
	{\headPrint[kana={ひきあんぎるん\tl{˥˩}},ipa={hi̥ki.angiruɴ},pos={句},para={ひきあんぐぬ},acc={F}]}%
	{%
		\MeaningsBlock{%
	\ryumeaning%
		{引き揚げる}{}{}%
	}%
	}%
	{\tailPrint[]}%
\entry
	{\headPrint[kana={ひき\tl{˥˩}うきるん\tl{˥}},ipa={hi̥ki.ukiruɴ},pos={動},rui={3},para={ひきうくぬ},acc={F+H}]}%
	{%
		\MeaningsBlock{%
	\ryumeaning%
		{引き受ける}{}{}%
	}%
	}%
	{\tailPrint[]}%
\entry
	{\headPrint[kana={ぴきうし\tl{˥˩}},ipa={pi̥ki.uʃi},pos={名},acc={F}]}%
	{%
		\MeaningsBlock{%
	\ryumeaning%
		{引き臼}{}{}%
	}%
	}%
	{\tailPrint[]}%
\entry
	{\headPrint[kana={ぴきうらすん},ipa={pi̥ki.urasuɴ},pos={動},rui={2},para={ぴきうらはぬ}]}%
	{%
		\MeaningsBlock{%
	\ryumeaning%
		{引き下ろす}{}{}%
	}%
	}%
	{\tailPrint[]}%
\entry
	{\headPrint[kana={ぴきうるすん},ipa={pi̥ki.urusuɴ},pos={動}]}%
	{%
		\MeaningsBlock{%
	\ryumeaning%
		{引き下ろす}{}{}%
	}%
	}%
	{\tailPrint[]}%
\entry
	{\headPrint[kana={ぴきけーすん\tl{˥˩}},ipa={pi̥kikeːsuɴ},pos={句},para={ぴきけーさぬ},acc={F}]}%
	{%
		\MeaningsBlock{%
	\ryumeaning%
		{引き返す}{}{}%
	}%
	}%
	{\tailPrint[]}%
\entry
	{\headPrint[kana={ぴきさくん\tl{˥˩}},ipa={pi̥kisḁkuɴ},pos={動},para={ぴきさかぬ},acc={F}]}%
	{%
		\MeaningsBlock{%
	\ryumeaning%
		{引き裂く}{}{}%
	}%
	}%
	{\tailPrint[]}%
\entry
	{\headPrint[kana={ぴきさんぎるん\tl{˥˩}},ipa={pi̥kisaŋgiruɴ},pos={動},rui={3},para={ぴきさんぐぬ},acc={F}]}%
	{%
		\MeaningsBlock{%
	\ryumeaning%
		{\ruby[g]{引}{ひ}っ\ruby[g]{提}{さ}げる}{}{}%
	}%
	}%
	{\tailPrint[]}%
\entry
	{\headPrint[kana={ぴき\tl{˥˩}さんぎん\tl{˩˥}},ipa={pi̥kisaŋgiɴ},pos={句},para={ぴきさんぐぬ},acc={F+R}]}%
	{%
		\MeaningsBlock{%
	\ryumeaning%
		{ぶら下げる}{}{}%
	}%
	}%
	{\tailPrint[]}%
\entry
	{\headPrint[kana={ぴきしぃきん\tl{˥˩}},ipa={pi̥kiʃi̥kiɴ},pos={句},para={ぴきすくぬ},acc={F}]}%
	{%
		\MeaningsBlock{%
	\ryumeaning%
		{引き付ける。引っ張る}{}{}%
	\ryumeaning%
		{痙攣する}{}{}%
	}%
	}%
	{\tailPrint[]}%
\entry
	{\headPrint[kana={ぴきすー\tl{˥˩}},ipa={pi̥kisuː},pos={名},acc={F}]}%
	{%
		\MeaningsBlock{%
	\ryumeaning%
		{引き潮}{}{}%
	}%
	}%
	{\tailPrint[]}%
\entry
	{\headPrint[kana={ぴきそーるん\tl{˥˩}},ipa={pi̥kisoːruɴ},pos={句},para={ぴきそーらぬ},acc={F}]}%
	{%
		\MeaningsBlock{%
	\ryumeaning%
		{引き連れる}{}{}%
	}%
	}%
	{\tailPrint[]}%
\entry
	{\headPrint[kana={ぴきだーすん},ipa={pi̥kidaːsuɴ},pos={句},para={ぴきんたはぬ}]}%
	{%
		\MeaningsBlock{%
	\ryumeaning%
		{引き出す}{}{}%
	}%
	}%
	{\tailPrint[]}%
\entry
	{\headPrint[kana={ぴきとーすん},ipa={pi̥kitoːsuɴ},pos={句},para={ぴきとはぬ}]}%
	{%
		\MeaningsBlock{%
	\ryumeaning%
		{引き倒す}{}{}%
	}%
	}%
	{\tailPrint[]}%
\entry
	{\headPrint[kana={ぴきとぅるん\tl{˥˩}},ipa={pi̥kituruɴ},pos={動},rui={1},para={ぴきとぅらぬ},acc={F}]}%
	{%
		\MeaningsBlock{%
	\ryumeaning%
		{引き取る}{}{}%
	}%
	}%
	{\tailPrint[]}%
\entry
	{\headPrint[kana={ぴきなん\tl{˩˥}},ipa={pi̥kinaɴ},pos={名},acc={R}]}%
	{%
		\MeaningsBlock{%
	\ryumeaning%
		{曳き縄漁}{}{}%
	}%
	}%
	{\tailPrint[]}%
\entry
	{\headPrint[kana={ぴき\tl{˥˩}ぬぐん\tl{˩˥}},ipa={pi̥kinuguɴ},pos={動},rui={1},para={ぴきぬがぬ},acc={F+R}]}%
	{%
		\MeaningsBlock{%
	\ryumeaning%
		{引き抜く。抜き出す}{}{}%
	}%
	}%
	{\tailPrint[]}%
\entry
	{\headPrint[kana={ぴき\tl{˥˩}ぬばすん\tl{˩˥}},ipa={pi̥kinubasuɴ},pos={動},rui={2},para={ぴきぬばはぬ},acc={F+R}]}%
	{%
		\MeaningsBlock{%
	\ryumeaning%
		{引き伸ばす}{}{}%
	}%
	}%
	{\tailPrint[]}%
\entry
	{\headPrint[kana={ぴきぬ\ws{}ふり},ipa={pi̥kinu furi},pos={名}]}%
	{%
		\MeaningsBlock{%
	\ryumeaning%
		{分銅}{}{}%
	}%
	}%
	{\tailPrint[]}%
\entry
	{\headPrint[kana={ぴきはかるん},ipa={pi̥kihakaruɴ},pos={句},para={ぴきはからぬ}]}%
	{%
		\MeaningsBlock{%
	\ryumeaning%
		{引っかかる}{}{}%
	}%
	}%
	{\tailPrint[]}%
\entry
	{\headPrint[kana={ぴきやどぅ\tl{˩˥}},ipa={pi̥kijadu},pos={名},acc={R}]}%
	{%
		\MeaningsBlock{%
	\ryumeaning%
		{引き戸}{}{}%
	}%
	}%
	{\tailPrint[]}%
\entry
	{\headPrint[kana={ぴき\tl{˥˩}やぶるん\tl{˩˥}},ipa={pi̥kijaburuɴ},pos={動},rui={1},para={ぴきやぶらぬ},acc={F+R}]}%
	{%
		\MeaningsBlock{%
	\ryumeaning%
		{引き破る}{}{}%
	}%
	}%
	{\tailPrint[]}%
\entry
	{\headPrint[kana={ぴきゆしるん\tl{˥˩}},ipa={pi̥kijuʃiruɴ},pos={動},para={ぴきゆするぬ},acc={F}]}%
	{%
		\MeaningsBlock{%
	\ryumeaning%
		{引き寄せる。寄せ集める}{}{}%
	}%
	}%
	{\tailPrint[]}%
\entry
	{\headPrint[kana={びぎり\tl{˥˩}},ipa={bigiri},pos={名},acc={F}]}%
	{%
		\MeaningsBlock{%
	\ryumeaning%
		{男の兄弟}{「女の兄弟」は\emb{ぶなり}}{}%
	}%
	}%
	{\tailPrint[]}%
\entry
	{\headPrint[kana={ぴきんだすん},ipa={pi̥kindasuɴ},pos={動},rui={2},para={ぴきんだはぬ}]}%
	{%
		\MeaningsBlock{%
	\ryumeaning%
		{引き出す}{}{}%
	}%
	}%
	{\tailPrint[]}%
\entry
	{\headPrint[kana={ぴく},ipa={pi̥ku},pos={副}]}%
	{%
		\MeaningsBlock{%
	\ryumeaning%
		{早く。すぐに}{}{}%
	}%
	}%
	{\tailPrint[]}%
\entry
	{\headPrint[kana={ぴくだり\tl{˥˩}},ipa={pi̥kudari},pos={副},acc={F}]}%
	{%
		\MeaningsBlock{%
	\ryumeaning%
		{早めに。早く}{}{}%
	}%
	}%
	{\tailPrint[]}%
\entry
	{\headPrint[kana={ぴくら\tl{˥˩}やむん\tl{˩˥}},ipa={pi̥kurajamuɴ},pos={動},rui={1},para={ぴくらやまぬ},acc={F+R}]}%
	{%
		\MeaningsBlock{%
	\ryumeaning%
		{\ruby[g]{痺}{しび}れる}{}{}%
	}%
	}%
	{\tailPrint[]}%
\entry
	{\headPrint[kana={ぴくん\tl{˥}},ipa={pi̥kuɴ},pos={動},rui={1},para={ぴかぬ},acc={H}]}%
	{%
		\MeaningsBlock{%
	\ryumeaning%
		{引く。弾く}{}{}%
	}%
	}%
	{\tailPrint[]}%
\entry
	{\headPrint[kana={ぴくん\tl{˥}},ipa={pi̥kuɴ},pos={動},rui={1},acc={H}]}%
	{%
		\MeaningsBlock{%
	\ryumeaning%
		{穴をあける}{}{}%
	}%
	}%
	{\tailPrint[]}%
\entry
	{\headPrint[kana={ぴくん\tl{˥˩}},ipa={pi̥kuɴ},pos={動},rui={1},para={ぴかぬ},acc={F}]}%
	{%
		\MeaningsBlock{%
	\ryumeaning%
		{（臼を）挽く}{}{}%
	}%
	}%
	{\tailPrint[]}%
\entry
	{\headPrint[kana={ぴぐん\tl{˥}},ipa={piguɴ},pos={動},rui={1},para={ぴがぬ},acc={H}]}%
	{%
		\MeaningsBlock{%
	\ryumeaning%
		{\ruby[g]{削}{けず}る}{}{}%
	}%
	}%
	{\tailPrint[]}%
\entry
	{\headPrint[kana={ぴこ\tl{˥˩}はん},ipa={pi̥kohaɴ},pos={形},para={ぴこへぬ},acc={F}]}%
	{%
		\MeaningsBlock{%
	\ryumeaning%
		{危ない。危険だ}{}{}%
	}%
	}%
	{\tailPrint[]}%
\entry
	{\headPrint[kana={ぴさんたらすん\tl{˥˩}},ipa={pi̥santarasuɴ},pos={動},rui={2},para={ぴさんたらはぬ},acc={F}]}%
	{%
		\MeaningsBlock{%
	\ryumeaning%
		{押して平たくする。潰れる。平たくなる。ひしゃげる}{}{}%
	}%
	}%
	{\tailPrint[]}%
\entry
	{\headPrint[kana={ぴじがら\tl{˥}},ipa={pidʒigara},pos={名},acc={H}]}%
	{%
		\MeaningsBlock{%
	\ryumeaning%
		{\ruby[g]{鰹}{かつお}の削りかす}{}{}%
	}%
	}%
	{\tailPrint[]}%
\entry
	{\headPrint[kana={びしじ\tl{˥˩}},ipa={biʃidʒi},pos={名},acc={F}]}%
	{%
		\MeaningsBlock{%
	\ryumeaning%
		{礎石。土台}{}{}%
	}%
	}%
	{\tailPrint[]}%
\entry
	{\headPrint[kana={ぴしゃ\tl{˥}はん},ipa={pi̥ʃahaɴ},pos={形},para={ぴしゃへぬ},acc={H}]}%
	{%
		\MeaningsBlock{%
	\ryumeaning%
		{寒い。冷える}{}{}%
	}%
	}%
	{\tailPrint[]}%
\entry
	{\headPrint[kana={ぴじゅる\tl{˥}さーん},ipa={pidʒurusaːɴ},pos={形},acc={H}]}%
	{%
		\MeaningsBlock{%
	\ryumeaning%
		{冷たい}{}{}%
	}%
	}%
	{\tailPrint[]}%
\entry
	{\headPrint[kana={びしょー\tl{˥˩}},ipa={biʃoː},pos={名},acc={F}]}%
	{%
		\MeaningsBlock{%
	\ryumeaning%
		{（磯に続く）浅瀬}{}{}%
	}%
	}%
	{\tailPrint[]}%
\entry
	{\headPrint[kana={びしるん\tl{˥˩}},ipa={biʃiruɴ},pos={動},rui={3},para={びすぬ},acc={F}]}%
	{%
		\MeaningsBlock{%
	\ryumeaning%
		{据える}{}{}%
	}%
	}%
	{\tailPrint[]}%
\entry
	{\headPrint[kana={ひすくり\tl{˩˥}},ipa={hisu̥kuri},pos={名},acc={R}]}%
	{%
		\MeaningsBlock{%
	\ryumeaning%
		{家造り}{}{}%
	}%
	}%
	{\tailPrint[]}%
\entry
	{\headPrint[kana={ひすくりよい\tl{˩˥}},ipa={hisu̥kurijoi},pos={名},acc={R}]}%
	{%
		\MeaningsBlock{%
	\ryumeaning%
		{家の落成祝い}{}{}%
	}%
	}%
	{\tailPrint[]}%
\entry
	{\headPrint[kana={ぴそーぎるん},ipa={pi̥soːgiruɴ},pos={動},rui={3},para={ぴそーぐぬ}]}%
	{%
		\MeaningsBlock{%
	\ryumeaning%
		{広げる。拡大する}{広げて敷く}{}%
	\ryumeaning%
		{延べる。延ばす}{}{}%
	}%
	}%
	{\tailPrint[]}%
\entry
	{\headPrint[kana={びたこら\tl{˩˥}},ipa={bitḁkora},pos={名},acc={R}]}%
	{%
		\MeaningsBlock{%
	\ryumeaning%
		{酒に酔う}{}{}%
	}%
	}%
	{\tailPrint[]}%
\entry
	{\headPrint[kana={ひだみるん\tl{˥}},ipa={hidamiruɴ},pos={動},para={ひだみぬ},acc={H}]}%
	{%
		\MeaningsBlock{%
	\ryumeaning%
		{隔てる。離れる}{}{}%
	}%
	}%
	{\tailPrint[]}%
\entry
	{\headPrint[kana={ひだみん\tl{˥}},ipa={hidamiɴ},pos={動},rui={3},acc={H}]}%
	{%
		\MeaningsBlock{%
	\ryumeaning%
		{隔てる。離れる}{}{}%
	}%
	}%
	{\tailPrint[]}%
\entry
	{\headPrint[kana={びちぃ\tl{˩˥}},ipa={bitʃi},pos={名},acc={R}]}%
	{%
		\MeaningsBlock{%
	\ryumeaning%
		{別}{}{}%
	}%
	}%
	{\tailPrint[]}%
\entry
	{\headPrint[kana={ぴちぃ\tl{˥˩}},ipa={pi̥tʃi},pos={名},acc={F}]}%
	{%
		\MeaningsBlock{%
	\ryumeaning%
		{\ruby[g]{未}{ひつじ}}{十二支の未}{}%
	}%
	}%
	{\tailPrint[]}%
\entry
	{\headPrint[kana={びちぃぬ\tl{˩˥}},ipa={bitʃinu},pos={句},acc={R}]}%
	{%
		\MeaningsBlock{%
	\ryumeaning%
		{別の。他の}{}{}%
	}%
	}%
	{\tailPrint[]}%
\entry
	{\headPrint[kana={びちぃびちぃ},ipa={bitʃibitʃi},pos={副}]}%
	{%
		\MeaningsBlock{%
	\ryumeaning%
		{別々}{}{}%
	}%
	}%
	{\tailPrint[]}%
\entry
	{\headPrint[kana={びちゃびちゃ},ipa={bitʃabitʃa},pos={擬}]}%
	{%
		\MeaningsBlock{%
	\ryumeaning%
		{じめじめ}{湿っているさま}{}%
	}%
	}%
	{\tailPrint[]}%
\entry
	{\headPrint[kana={ぴっしゃ\tl{˥}},ipa={piʃʃa},pos={名},acc={H}]}%
	{%
		\MeaningsBlock{%
	\ryumeaning%
		{筆者。書記}{王府時代の役人}{}%
	}%
	}%
	{\tailPrint[]}%
\entry
	{\headPrint[kana={びっちゅる\tl{˩˥}},ipa={bittʃuru},pos={名},acc={R}]}%
	{%
		\MeaningsBlock{%
	\ryumeaning%
		{陽石}{男根をかたどった石}{}%
	}%
	}%
	{\tailPrint[]}%
\entry
	{\headPrint[kana={ぴてー\tl{˥}},ipa={pi̥teː},pos={名},acc={H}]}%
	{%
		\MeaningsBlock{%
	\ryumeaning%
		{畑}{}{}%
	}%
	}%
	{\tailPrint[]}%
\entry
	{\headPrint[kana={ぴてーぎ\tl{˥}},ipa={pi̥teːgi},pos={名},acc={H}]}%
	{%
		\MeaningsBlock{%
	\ryumeaning%
		{畑}{}{}%
	}%
	}%
	{\tailPrint[]}%
\entry
	{\headPrint[kana={ぴてーすぅぬ\tl{˥}},ipa={pi̥teːsu̥nu},pos={名},acc={H}]}%
	{%
		\MeaningsBlock{%
	\ryumeaning%
		{野良着。作業着}{}{}%
	}%
	}%
	{\tailPrint[]}%
\entry
	{\headPrint[kana={ぴてー\tl{˥}すくるん\tl{˥}},ipa={pi̥teːsu̥kuruɴ},pos={句},acc={H+H}]}%
	{%
		\MeaningsBlock{%
	\ryumeaning%
		{畑仕事。農作業}{}{}%
	}%
	}%
	{\tailPrint[]}%
\entry
	{\headPrint[kana={ぴてーひ\tl{˥}},ipa={pi̥teːhi},pos={名},acc={H}]}%
	{%
		\MeaningsBlock{%
	\ryumeaning%
		{畑小屋}{}{}%
	}%
	}%
	{\tailPrint[]}%
\entry
	{\headPrint[kana={ぴてん\tl{˥}},ipa={pi̥teɴ},pos={名},acc={H}]}%
	{%
		\MeaningsBlock{%
	\ryumeaning%
		{\ruby[g]{一}{いち}\ruby[g]{日}{にち}}{}{}%
	}%
	}%
	{\tailPrint[]}%
\entry
	{\headPrint[kana={ぴてんぴじゅ},ipa={pi̥tempidʒu},pos={名}]}%
	{%
		\MeaningsBlock{%
	\ryumeaning%
		{一日中。終日}{}{}%
	}%
	}%
	{\tailPrint[]}%
\entry
	{\headPrint[kana={ぴとぅち\tl{˥}},ipa={pi̥tutʃi},pos={名},acc={H}]}%
	{%
		\MeaningsBlock{%
	\ryumeaning%
		{一つ}{簡略化：てぃ}{}%
	}%
	}%
	{\tailPrint[]}%
\entry
	{\headPrint[kana={ぴとぅみしり\ws{}すん},ipa={pi̥tumiʃiri suɴ},pos={句}]}%
	{%
		\MeaningsBlock{%
	\ryumeaning%
		{人見知りをする}{}{}%
	}%
	}%
	{\tailPrint[]}%
\entry
	{\headPrint[kana={びどぅむ\tl{˥˩}},ipa={bidumu},pos={名},acc={F}]}%
	{%
		\MeaningsBlock{%
	\ryumeaning%
		{男。男性}{}{}%
	}%
	}%
	{\tailPrint[]}%
\entry
	{\headPrint[kana={ぴなかん\tl{˥}},ipa={pi̥nakaɴ},pos={名},acc={H}]}%
	{%
		\MeaningsBlock{%
	\ryumeaning%
		{火の神}{台所の守護神}{}%
	}%
	}%
	{\tailPrint[]}%
\entry
	{\headPrint[kana={ぴながん\tl{˥}},ipa={pi̥nagaɴ},pos={名},acc={H}]}%
	{%
		\MeaningsBlock{%
	\ryumeaning%
		{太陽の\ruby[g]{暈}{かさ}。\ruby[g]{暈}{かさ}}{}{}%
	}%
	}%
	{\tailPrint[]}%
\entry
	{\headPrint[kana={びな\tl{˩˥}さーん},ipa={binasaːɴ},pos={形},acc={R}]}%
	{%
		\MeaningsBlock{%
	\ryumeaning%
		{病弱だ}{}{}%
	}%
	}%
	{\tailPrint[]}%
\entry
	{\headPrint[kana={ひな\tl{˩˥}\ws{}のーぬ\tl{˩˥}},ipa={hina noːnu},pos={句},acc={R R}]}%
	{%
		\MeaningsBlock{%
	\ryumeaning%
		{留守}{「家に居ない」の意}{}%
	}%
	}%
	{\tailPrint[]}%
\entry
	{\headPrint[kana={ぴな\tl{˥}はん},ipa={pi̥nahaɴ},pos={形},para={びなへぬ},acc={H}]}%
	{%
		\MeaningsBlock{%
	\ryumeaning%
		{劣る。貧弱だ}{}{}%
	}%
	}%
	{\tailPrint[]}%
\entry
	{\headPrint[kana={ぴなり\tl{˥}},ipa={pi̥nari},pos={名},acc={H}]}%
	{%
		\MeaningsBlock{%
	\ryumeaning%
		{左。左方}{}{}%
	}%
	}%
	{\tailPrint[]}%
\entry
	{\headPrint[kana={ぴなるん\tl{˥˩}},ipa={pi̥naruɴ},pos={動},rui={1},para={ぴならぬ},acc={F}]}%
	{%
		\MeaningsBlock{%
	\ryumeaning%
		{減る。減少する}{}{}%
	}%
	}%
	{\tailPrint[]}%
\entry
	{\headPrint[kana={ぴに\tl{˥˩}},ipa={pi̥ni},pos={名},acc={F}]}%
	{%
		\MeaningsBlock{%
	\ryumeaning%
		{\ruby[g]{鬚}{ひげ}}{}{}%
	}%
	}%
	{\tailPrint[]}%
\entry
	{\headPrint[kana={ぴにるん\tl{˥}},ipa={pi̥niruɴ},pos={動},rui={3},para={ぴぬぬ},acc={H}]}%
	{%
		\MeaningsBlock{%
	\ryumeaning%
		{\ruby[g]{捻}{ひね}る}{}{}%
	}%
	}%
	{\tailPrint[]}%
\entry
	{\headPrint[kana={ひぬばん\tl{˩˥}},ipa={hinubaɴ},pos={名},acc={R}]}%
	{%
		\MeaningsBlock{%
	\ryumeaning%
		{留守番}{}{}%
	}%
	}%
	{\tailPrint[]}%
\entry
	{\headPrint[kana={ぴぱちぃ\tl{˥}},ipa={pi̥patʃi},pos={名},acc={H}]}%
	{%
		\MeaningsBlock{%
	\ryumeaning%
		{ヒハツモドキ}{香辛料の植物}{}%
	}%
	}%
	{\tailPrint[]}%
\entry
	{\headPrint[kana={ぴぱり\tl{˥˩}},ipa={pi̥pari},pos={名},acc={F}]}%
	{%
		\MeaningsBlock{%
	\ryumeaning%
		{地割れ。旱害}{旱魃による田畑の地割れ}{}%
	}%
	}%
	{\tailPrint[]}%
\entry
	{\headPrint[kana={\josho{}ひばん\kako{}むり},ipa={hibammuri},pos={名},acc={H+L}]}%
	{%
		\MeaningsBlock{%
	\ryumeaning%
		{火番盛}{王府時代の烽火台}{}%
	}%
	}%
	{\tailPrint[]}%
\entry
	{\headPrint[kana={びびんたま\tl{˥}},ipa={bibintama},pos={名},acc={H}]}%
	{%
		\MeaningsBlock{%
	\ryumeaning%
		{小指}{}{}%
	}%
	}%
	{\tailPrint[]}%
\entry
	{\headPrint[kana={ぴま\tl{˥˩}},ipa={pi̥ma},pos={名},acc={F}]}%
	{%
		\MeaningsBlock{%
	\ryumeaning%
		{\ruby[g]{暇}{ひま}}{}{}%
	}%
	}%
	{\tailPrint[]}%
\entry
	{\headPrint[kana={ぴみざ\tl{˥}},ipa={pi̥midza},pos={名},acc={H}]}%
	{%
		\MeaningsBlock{%
	\ryumeaning%
		{\ruby[g]{山羊}{やぎ}}{}{}%
	}%
	}%
	{\tailPrint[]}%
\entry
	{\headPrint[kana={ぴむ\tl{˥}},ipa={pi̥mu},pos={名},acc={H}]}%
	{%
		\MeaningsBlock{%
	\ryumeaning%
		{（魚の）えら}{}{}%
	}%
	}%
	{\tailPrint[]}%
\entry
	{\headPrint[kana={ぴゃーぐ\tl{˥}},ipa={pjaːgu},pos={名},acc={H}]}%
	{%
		\MeaningsBlock{%
	\ryumeaning%
		{百}{}{}%
	}%
	}%
	{\tailPrint[]}%
\entry
	{\headPrint[kana={ぴゃっかるん\tl{˥}},ipa={pjakkaruɴ},pos={動},rui={1},para={ぴゃっからぬ},acc={H}]}%
	{%
		\MeaningsBlock{%
	\ryumeaning%
		{這って進む}{}{}%
	}%
	}%
	{\tailPrint[]}%
\entry
	{\headPrint[kana={ぴゃんが\tl{˥}},ipa={pjaŋga},pos={名},acc={H}]}%
	{%
		\MeaningsBlock{%
	\ryumeaning%
		{お転婆。お転婆娘}{}{}%
	}%
	}%
	{\tailPrint[]}%
\entry
	{\headPrint[kana={びゅーびゅー},ipa={bjuːbjuː},pos={擬}]}%
	{%
		\MeaningsBlock{%
	\ryumeaning%
		{ひゅうひゅう}{穴やすきま風が吹きぬける音の形容}{}%
	}%
	}%
	{\tailPrint[]}%
\entry
	{\headPrint[kana={ひよー\tl{˥}},ipa={hijoː},pos={名},acc={H}]}%
	{%
		\MeaningsBlock{%
	\ryumeaning%
		{日雇い}{}{}%
	}%
	}%
	{\tailPrint[]}%
\entry
	{\headPrint[kana={びょー\tl{˩˥}はん},ipa={bjoːhaɴ},pos={形},para={びょーへぬ？},acc={R}]}%
	{%
		\MeaningsBlock{%
	\ryumeaning%
		{\ruby[g]{痒}{かゆ}い}{}{}%
	}%
	}%
	{\tailPrint[]}%
\entry
	{\headPrint[kana={びょすた},ipa={bjosu̥ta},pos={副}]}%
	{%
		\MeaningsBlock{%
	\ryumeaning%
		{一目散に}{}{}%
	}%
	}%
	{\tailPrint[]}%
\entry
	{\headPrint[kana={ぴょん\tl{˥˩}},ipa={pjoɴ},pos={名},acc={F}]}%
	{%
		\MeaningsBlock{%
	\ryumeaning%
		{足跡}{}{}%
	}%
	}%
	{\tailPrint[]}%
\entry
	{\headPrint[kana={びら\tl{˩˥}},ipa={bira},pos={名},acc={R}]}%
	{%
		\MeaningsBlock{%
	\ryumeaning%
		{\ruby[g]{韮}{にら}}{}{}%
	}%
	}%
	{\tailPrint[]}%
\entry
	{\headPrint[kana={ぴら\tl{˥}},ipa={pi̥ra},pos={名},acc={H}]}%
	{%
		\MeaningsBlock{%
	\ryumeaning%
		{\ruby[g]{箆}{へら}}{}{}%
	}%
	}%
	{\tailPrint[]}%
\entry
	{\headPrint[kana={びらーすん\tl{˩˥}},ipa={biraːsuɴ},pos={動},rui={2},para={びらはぬ},acc={R}]}%
	{%
		\MeaningsBlock{%
	\ryumeaning%
		{中毒させる}{}{}%
	}%
	}%
	{\tailPrint[]}%
\entry
	{\headPrint[kana={ぴらい\tl{˥}},ipa={pi̥rai},pos={名},acc={H}]}%
	{%
		\MeaningsBlock{%
	\ryumeaning%
		{付き合い。交際}{}{}%
	}%
	}%
	{\tailPrint[]}%
\entry
	{\headPrint[kana={びらがすん\tl{˥˩}},ipa={biragasuɴ},pos={動},rui={2},para={びらがはぬ},acc={F}]}%
	{%
		\MeaningsBlock{%
	\ryumeaning%
		{押しつぶす。押して平たくする}{}{}%
	}%
	}%
	{\tailPrint[]}%
\entry
	{\headPrint[kana={ひらきるん\tl{˥}},ipa={hirakiruɴ},pos={動},rui={3},para={ひらくぬ},acc={H}]}%
	{%
		\MeaningsBlock{%
	\ryumeaning%
		{開ける}{}{}%
	\ryumeaning%
		{発展する}{}{}%
	}%
	}%
	{\tailPrint[]}%
\entry
	{\headPrint[kana={びらぐ\tl{˩˥}},ipa={biragu},pos={名},acc={R}]}%
	{%
		\MeaningsBlock{%
	\ryumeaning%
		{竹製の食べ物籠}{}{}%
	}%
	}%
	{\tailPrint[]}%
\entry
	{\headPrint[kana={ぴらぐ\tl{˥}},ipa={pi̥ragu},pos={名},acc={H}]}%
	{%
		\MeaningsBlock{%
	\ryumeaning%
		{寒気。寒波}{}{}%
	}%
	}%
	{\tailPrint[]}%
\entry
	{\headPrint[kana={ぴらす},ipa={pi̥rasu},pos={名}]}%
	{%
		\MeaningsBlock{%
	\ryumeaning%
		{付き合い。交際}{}{}%
	}%
	}%
	{\tailPrint[]}%
\entry
	{\headPrint[kana={ぴらすか\tl{˥}},ipa={pi̥rasu̥ka},pos={名},acc={H}]}%
	{%
		\MeaningsBlock{%
	\ryumeaning%
		{怠惰な農夫。下手な農夫}{}{}%
	}%
	}%
	{\tailPrint[]}%
\entry
	{\headPrint[kana={ぴらすん\tl{˥}},ipa={pi̥rasuɴ},pos={動},rui={2},para={ぴらはぬ},acc={H}]}%
	{%
		\MeaningsBlock{%
	\ryumeaning%
		{冷やす。冷ます}{}{}%
	}%
	}%
	{\tailPrint[]}%
\entry
	{\headPrint[kana={ぴらすん\tl{˥}},ipa={pi̥rasuɴ},pos={動},rui={2},para={ぴらはぬ},acc={H}]}%
	{%
		\MeaningsBlock{%
	\ryumeaning%
		{おどける}{}{}%
	}%
	}%
	{\tailPrint[]}%
\entry
	{\headPrint[kana={びらま\tl{˩˥}},ipa={birama},pos={名},acc={R}]}%
	{%
		\MeaningsBlock{%
	\ryumeaning%
		{兄さん}{年上の男性を指す}{}%
	}%
	}%
	{\tailPrint[]}%
\entry
	{\headPrint[kana={ぴり\tl{˥}},ipa={piri},pos={-},acc={H}]}%
	{%
		\MeaningsBlock{%
	\ryumeaning%
		{冷える。冷たい}{}{}%
	}%
	}%
	{\tailPrint[]}%
\entry
	{\headPrint[kana={ぴりかち\tl{˥}},ipa={pirikḁtʃi},pos={名},acc={H}]}%
	{%
		\MeaningsBlock{%
	\ryumeaning%
		{冷たい風。涼風}{}{}%
	}%
	}%
	{\tailPrint[]}%
\entry
	{\headPrint[kana={びりかなぱ\tl{˩˥}},ipa={birikḁnapa},pos={名},acc={R}]}%
	{%
		\MeaningsBlock{%
	\ryumeaning%
		{クワズイモ}{毒のある植物}{}%
	}%
	}%
	{\tailPrint[]}%
\entry
	{\headPrint[kana={ぴりがるん\tl{˥}},ipa={pirigaruɴ},pos={動},rui={1},para={ぴりがらぬ},acc={H}]}%
	{%
		\MeaningsBlock{%
	\ryumeaning%
		{澄む。沈澱する}{}{}%
	}%
	}%
	{\tailPrint[]}%
\entry
	{\headPrint[kana={ぴりしゃー\ws{}すん},ipa={piriʃaː suɴ},pos={句}]}%
	{%
		\MeaningsBlock{%
	\ryumeaning%
		{涼む}{}{}%
	}%
	}%
	{\tailPrint[]}%
\entry
	{\headPrint[kana={ぴりしゃ\tl{˥}はん},ipa={piriʃahaɴ},pos={形},para={ぴりしゃへぬ},acc={H}]}%
	{%
		\MeaningsBlock{%
	\ryumeaning%
		{涼しい}{}{}%
	}%
	}%
	{\tailPrint[]}%
\entry
	{\headPrint[kana={ぴりふつぁりむぬ},ipa={pirifu̥tsarimunu},pos={名}]}%
	{%
		\MeaningsBlock{%
	\ryumeaning%
		{冷や飯}{冷えた食事}{}%
	}%
	}%
	{\tailPrint[]}%
\entry
	{\headPrint[kana={ぴる\tl{˥˩}},ipa={pi̥ru},pos={名},acc={F}]}%
	{%
		\MeaningsBlock{%
	\ryumeaning%
		{昼。昼間}{}{}%
	}%
	}%
	{\tailPrint[]}%
\entry
	{\headPrint[kana={ぴる\tl{˥˩}},ipa={pi̥ru},pos={名},acc={F}]}%
	{%
		\MeaningsBlock{%
	\ryumeaning%
		{\ruby[g]{大蒜}{にんにく}}{}{}%
	}%
	}%
	{\tailPrint[]}%
\entry
	{\headPrint[kana={ぴるぬ\ws{}うち},ipa={pi̥runu utʃi},pos={句}]}%
	{%
		\MeaningsBlock{%
	\ryumeaning%
		{昼間。昼の内に}{}{}%
	}%
	}%
	{\tailPrint[]}%
\entry
	{\headPrint[kana={ぴるまさん},ipa={pi̥rumasaɴ},pos={形},para={ぴるまへぬ},acc={H}]}%
	{%
		\MeaningsBlock{%
	\ryumeaning%
		{不思議だ。珍しい}{}{}%
	}%
	}%
	{\tailPrint[]}%
\entry
	{\headPrint[kana={ぴるま\tl{˥}はん},ipa={pi̥rumahaɴ},pos={形},para={ぴるまへぬ},acc={H}]}%
	{%
		\MeaningsBlock{%
	\ryumeaning%
		{不思議だ。珍しい}{}{}%
	}%
	}%
	{\tailPrint[]}%
\entry
	{\headPrint[kana={ひるまるん\tl{˥}},ipa={hirumaruɴ},pos={動},rui={1},para={ひるまらぬ},acc={H}]}%
	{%
		\MeaningsBlock{%
	\ryumeaning%
		{広がる。広まる}{噂などについていう}{}%
	}%
	}%
	{\tailPrint[]}%
\entry
	{\headPrint[kana={ぴるまるん\tl{˥}},ipa={pi̥rumaruɴ},pos={動},rui={1},para={ぴるまらぬ},acc={H}]}%
	{%
		\MeaningsBlock{%
	\ryumeaning%
		{広まる。広がる}{}{}%
	}%
	}%
	{\tailPrint[]}%
\entry
	{\headPrint[kana={ひるみるん\tl{˥}},ipa={hirumiruɴ},pos={動},rui={3},para={ひるむぬ},acc={H}]}%
	{%
		\MeaningsBlock{%
	\ryumeaning%
		{広める。広げる。流行らせる}{}{}%
	}%
	}%
	{\tailPrint[]}%
\entry
	{\headPrint[kana={ひるみん\tl{˥}},ipa={hirumiɴ},pos={動},rui={3},para={ひるむぬ},acc={H}]}%
	{%
		\MeaningsBlock{%
	\ryumeaning%
		{広める。広げる。流行らせる}{}{}%
	}%
	}%
	{\tailPrint[]}%
\entry
	{\headPrint[kana={ひるん\tl{˥}},ipa={hiruɴ},pos={動},rui={3},para={ふーぬ},acc={H}]}%
	{%
		\MeaningsBlock{%
	\ryumeaning%
		{あげる。与える}{}{}%
	}%
	}%
	{\tailPrint[]}%
\entry
	{\headPrint[kana={びるん\tl{˥˩}},ipa={biruɴ},pos={動},rui={1},para={びらぬ},acc={F}]}%
	{%
		\MeaningsBlock{%
	\ryumeaning%
		{座る}{}{}%
	}%
	}%
	{\tailPrint[]}%
\entry
	{\headPrint[kana={ぴん\tl{˥}},ipa={piɴ},pos={名},acc={H}]}%
	{%
		\MeaningsBlock{%
	\ryumeaning%
		{屁}{}{}%
	}%
	}%
	{\tailPrint[]}%
\entry
	{\headPrint[kana={ぴん\tl{˥˩}},ipa={piɴ},pos={名},acc={F}]}%
	{%
		\MeaningsBlock{%
	\ryumeaning%
		{笛。竹笛}{}{}%
	}%
	}%
	{\tailPrint[]}%
\entry
	{\headPrint[kana={ぴん\tl{˥˩}},ipa={piɴ},pos={名},acc={F}]}%
	{%
		\MeaningsBlock{%
	\ryumeaning%
		{\ruby[g]{稗}{ひえ}}{}{}%
	}%
	}%
	{\tailPrint[]}%
\entry
	{\headPrint[kana={ぴん\tl{˥˩}},ipa={piɴ},pos={名},acc={F}]}%
	{%
		\MeaningsBlock{%
	\ryumeaning%
		{\ruby[g]{日}{ひ}。日にち}{}{}%
	}%
	}%
	{\tailPrint[]}%
\entry
	{\headPrint[kana={ぴんがすん\tl{˥}},ipa={piŋgasuɴ},pos={動},rui={2},para={ぴんがはぬ},acc={H}]}%
	{%
		\MeaningsBlock{%
	\ryumeaning%
		{逃がす}{}{}%
	}%
	}%
	{\tailPrint[]}%
\entry
	{\headPrint[kana={ぴんがん\tl{˥}},ipa={piŋgaɴ},pos={名},acc={H}]}%
	{%
		\MeaningsBlock{%
	\ryumeaning%
		{彼岸}{}{}%
	}%
	}%
	{\tailPrint[]}%
\entry
	{\headPrint[kana={ぴんぎまーるん\tl{˥}},ipa={piŋgimaːruɴ},pos={動},rui={1},para={ぴんぎまーらぬ},acc={H}]}%
	{%
		\MeaningsBlock{%
	\ryumeaning%
		{逃げ回る}{}{}%
	}%
	}%
	{\tailPrint[]}%
\entry
	{\headPrint[kana={ぴんぎるん\tl{˥}},ipa={piŋgiruɴ},pos={動},rui={3},para={ぴんぐぬ},acc={H}]}%
	{%
		\MeaningsBlock{%
	\ryumeaning%
		{逃げる。逃亡する}{}{}%
	}%
	}%
	{\tailPrint[]}%
\entry
	{\headPrint[kana={ぴんくるん\tl{˥}},ipa={piŋkuruɴ},pos={動},acc={H}]}%
	{%
		\MeaningsBlock{%
	\ryumeaning%
		{抉る}{}{}%
	}%
	}%
	{\tailPrint[]}%
\entry
	{\headPrint[kana={ぴんぐるん\tl{˥}},ipa={piŋguruɴ},pos={動},rui={1},para={ぴんぐらぬ},acc={H}]}%
	{%
		\MeaningsBlock{%
	\ryumeaning%
		{冷える。凍える}{}{}%
	}%
	}%
	{\tailPrint[]}%
\entry
	{\headPrint[kana={ひんすん\tl{˥˩}},ipa={hinsuɴ},pos={動},rui={2b},para={ひんさぬ},acc={F}]}%
	{%
		\MeaningsBlock{%
	\ryumeaning%
		{すねる}{}{}%
	}%
	}%
	{\tailPrint[]}%
\entry
	{\headPrint[kana={ぴんそー\tl{˥}},ipa={pinsoː},pos={名},acc={H}]}%
	{%
		\MeaningsBlock{%
	\ryumeaning%
		{貧乏。貧しい}{}{}%
	}%
	}%
	{\tailPrint[]}%
\entry
	{\headPrint[kana={ぴんだま\tl{˥}},ipa={pindama},pos={名},acc={H}]}%
	{%
		\MeaningsBlock{%
	\ryumeaning%
		{火の球。人魂}{}{}%
	}%
	}%
	{\tailPrint[]}%
\entry
	{\headPrint[kana={ぴんちぃ\tl{˥˩}},ipa={pintʃi},pos={名},acc={F}]}%
	{%
		\MeaningsBlock{%
	\ryumeaning%
		{日付。日和。日数}{}{}%
	}%
	}%
	{\tailPrint[]}%
\entry
	{\headPrint[kana={ひんとー\tl{˥˩}},ipa={hintoː},pos={名},acc={F}]}%
	{%
		\MeaningsBlock{%
	\ryumeaning%
		{返事。返答}{}{}%
	}%
	}%
	{\tailPrint[]}%
\entry
	{\headPrint[kana={びんとー\tl{˥˩}},ipa={bintoː},pos={名},acc={F}]}%
	{%
		\MeaningsBlock{%
	\ryumeaning%
		{弁当}{}{}%
	}%
	}%
	{\tailPrint[]}%
\entry
	{\headPrint[kana={ぴんとまらすん\tl{˥}},ipa={pintomarasuɴ},pos={動},rui={2},para={ぴんとまらはぬ},acc={H}]}%
	{%
		\MeaningsBlock{%
	\ryumeaning%
		{尖らせる}{}{}%
	}%
	}%
	{\tailPrint[]}%
\entry
	{\headPrint[kana={ぴんとまり},ipa={pinto̥mari},pos={名}]}%
	{%
		\MeaningsBlock{%
	\ryumeaning%
		{鋭く尖った}{}{}%
	}%
	}%
	{\tailPrint[]}%
\entry
	{\headPrint[kana={ぴんとまるん\tl{˥}},ipa={pintomaruɴ},pos={動},rui={1},para={ぴんとまらぬ},acc={H}]}%
	{%
		\MeaningsBlock{%
	\ryumeaning%
		{尖る}{}{}%
	}%
	}%
	{\tailPrint[]}%
\entry
	{\headPrint[kana={ぴんふつぁ\tl{˥}はーん},ipa={pimfu̥tsahaːɴ},pos={形},acc={H}]}%
	{%
		\MeaningsBlock{%
	\ryumeaning%
		{屁臭い}{屁の臭いがして臭い}{}%
	}%
	}%
	{\tailPrint[]}%
\LetterHead{ふ}
\entry
	{\headPrint[kana={ぶー},ipa={buː},pos={接頭}]}%
	{%
		\MeaningsBlock{%
	\ryumeaning%
		{\ruby[g]{大}{おお}。\ruby[g]{大}{だい}。大きい}{\emb{ぶーいし}「大石」、\emb{ぶーあみ}「大雨」など}{}%
	}%
	}%
	{\tailPrint[]}%
\entry
	{\headPrint[kana={ぶー\tl{˩˥}},ipa={buː},pos={名},acc={R}]}%
	{%
		\MeaningsBlock{%
	\ryumeaning%
		{\ruby[g]{紐}{ひも}}{}{}%
	}%
	}%
	{\tailPrint[]}%
\entry
	{\headPrint[kana={ぶー\tl{˩˥}},ipa={buː},pos={名},acc={R}]}%
	{%
		\MeaningsBlock{%
	\ryumeaning%
		{\ruby[g]{苧}{ちょ}\ruby[g]{麻}{ま}}{上布の原料}{}%
	}%
	}%
	{\tailPrint[]}%
\entry
	{\headPrint[kana={ぷー\tl{˥}},ipa={puː},pos={名},acc={H}]}%
	{%
		\MeaningsBlock{%
	\ryumeaning%
		{\ruby[g]{穂}{ほ}}{植物の穂}{}%
	}%
	}%
	{\tailPrint[]}%
\entry
	{\headPrint[kana={ふーがー\tl{˥}},ipa={fuːgaː},pos={名},acc={H}]}%
	{%
		\MeaningsBlock{%
	\ryumeaning%
		{大川}{地名。石垣市の字名}{}%
	}%
	}%
	{\tailPrint[]}%
\entry
	{\headPrint[kana={ぶーかち\tl{˥}},ipa={buːkḁtʃi},pos={名},acc={H}]}%
	{%
		\MeaningsBlock{%
	\ryumeaning%
		{大風。台風}{}{}%
	}%
	}%
	{\tailPrint[]}%
\entry
	{\headPrint[kana={ぷーき},ipa={puːki},pos={名}]}%
	{%
		\MeaningsBlock{%
	\ryumeaning%
		{暑気。猛暑}{}{}%
	}%
	}%
	{\tailPrint[]}%
\entry
	{\headPrint[kana={ぷーき\tl{˥˩}},ipa={puːki},pos={名},acc={F}]}%
	{%
		\MeaningsBlock{%
	\ryumeaning%
		{マラリア。熱病}{}{}%
	}%
	}%
	{\tailPrint[]}%
\entry
	{\headPrint[kana={ふーく\tl{˥}},ipa={fuːku},pos={名},acc={H}]}%
	{%
		\MeaningsBlock{%
	\ryumeaning%
		{奉公}{}{}%
	}%
	}%
	{\tailPrint[]}%
\entry
	{\headPrint[kana={ふーさた},ipa={fuːsḁta},pos={名}]}%
	{%
		\MeaningsBlock{%
	\ryumeaning%
		{黒糖。黒砂糖}{}{}%
	}%
	}%
	{\tailPrint[]}%
\entry
	{\headPrint[kana={ぶーさ\tl{˥}\ws{}なすん\tl{˩˥}},ipa={buːsa nasuɴ},pos={句},acc={H R}]}%
	{%
		\MeaningsBlock{%
	\ryumeaning%
		{多くする。増やす。増す}{}{}%
	}%
	}%
	{\tailPrint[]}%
\entry
	{\headPrint[kana={ぶーさ\tl{˥}\ws{}なるん\tl{˩˥}},ipa={buːsa naruɴ},pos={句},acc={H R}]}%
	{%
		\MeaningsBlock{%
	\ryumeaning%
		{多くなる。増える}{}{}%
	\ryumeaning%
		{成長する}{}{}%
	}%
	}%
	{\tailPrint[]}%
\entry
	{\headPrint[kana={ぶーさら},ipa={buːsḁra},pos={名}]}%
	{%
		\MeaningsBlock{%
	\ryumeaning%
		{大皿}{}{}%
	}%
	}%
	{\tailPrint[]}%
\entry
	{\headPrint[kana={\josho{}ふー\kako{}し},ipa={fuːʃi},pos={副},acc={HL}]}%
	{%
		\MeaningsBlock{%
	\ryumeaning%
		{黒く。黒ずむ}{}{}%
	}%
	}%
	{\tailPrint[]}%
\entry
	{\headPrint[kana={ぶーしけん},ipa={buːʃi̥keɴ},pos={名}]}%
	{%
		\MeaningsBlock{%
	\ryumeaning%
		{満月}{}{}%
	}%
	}%
	{\tailPrint[]}%
\entry
	{\headPrint[kana={\josho{}ふー\kako{}し\ws{}なるん\tl{˩˥}},ipa={fuːʃi naruɴ},pos={句},acc={HL R}]}%
	{%
		\MeaningsBlock{%
	\ryumeaning%
		{黒む。黒くなる。黒ずむ}{黒みを帯びる}{}%
	}%
	}%
	{\tailPrint[]}%
\entry
	{\headPrint[kana={ぶーすー\tl{˥}},ipa={buːsuː},pos={名},acc={H}]}%
	{%
		\MeaningsBlock{%
	\ryumeaning%
		{大潮}{}{}%
	}%
	}%
	{\tailPrint[]}%
\entry
	{\headPrint[kana={ふーすぅぬ\tl{˥}},ipa={fuːsu̥nu},pos={名},acc={H}]}%
	{%
		\MeaningsBlock{%
	\ryumeaning%
		{古着}{}{}%
	}%
	}%
	{\tailPrint[]}%
\entry
	{\headPrint[kana={ふーたい\tl{˥˩}},ipa={fuːtai},pos={名},acc={F}]}%
	{%
		\MeaningsBlock{%
	\ryumeaning%
		{風袋}{軽量のときの入れ物の重さ}{}%
	}%
	}%
	{\tailPrint[]}%
\entry
	{\headPrint[kana={ぶーどぅり\tl{˩˥}},ipa={buːduri},pos={名},acc={R}]}%
	{%
		\MeaningsBlock{%
	\ryumeaning%
		{居所。住所}{}{}%
	}%
	}%
	{\tailPrint[]}%
\entry
	{\headPrint[kana={ふー\tl{˥}\ws{}なすん\tl{˩˥}},ipa={fuː nasuɴ},pos={句},acc={H R}]}%
	{%
		\MeaningsBlock{%
	\ryumeaning%
		{古くする}{}{}%
	}%
	}%
	{\tailPrint[]}%
\entry
	{\headPrint[kana={ふー\tl{˥}\ws{}なるん\tl{˩˥}},ipa={fuː naruɴ},pos={句},para={ふーならぬ},acc={H R}]}%
	{%
		\MeaningsBlock{%
	\ryumeaning%
		{古くなる。古む。古びる}{}{}%
	}%
	}%
	{\tailPrint[]}%
\entry
	{\headPrint[kana={ぶーなん\tl{˥}},ipa={buːnaɴ},pos={名},acc={H}]}%
	{%
		\MeaningsBlock{%
	\ryumeaning%
		{大波。津波}{}{}%
	}%
	}%
	{\tailPrint[]}%
\entry
	{\headPrint[kana={\josho{}ぷーぬ\kako{}\ws{}むぬ},ipa={puːnu munu},pos={句},acc={H L}]}%
	{%
		\MeaningsBlock{%
	\ryumeaning%
		{穀物}{穂の物の義}{}%
	}%
	}%
	{\tailPrint[]}%
\entry
	{\headPrint[kana={ぶーばた\tl{˥}},ipa={buːbata},pos={名},acc={H}]}%
	{%
		\MeaningsBlock{%
	\ryumeaning%
		{胃。胃袋}{}{}%
	}%
	}%
	{\tailPrint[]}%
\entry
	{\headPrint[kana={ぶーび\tl{˥}},ipa={buːbi},pos={名},acc={H}]}%
	{%
		\MeaningsBlock{%
	\ryumeaning%
		{親指}{}{}%
	}%
	}%
	{\tailPrint[]}%
\entry
	{\headPrint[kana={ぶー\tl{˥}ひー\tl{˩˥}},ipa={buːhiː},pos={名},acc={H+R}]}%
	{%
		\MeaningsBlock{%
	\ryumeaning%
		{母屋}{}{}%
	}%
	}%
	{\tailPrint[]}%
\entry
	{\headPrint[kana={ぶーぴぃとぅ},ipa={buːpi̥tu},pos={名}]}%
	{%
		\MeaningsBlock{%
	\ryumeaning%
		{大きな人。偉い人}{}{}%
	}%
	}%
	{\tailPrint[]}%
\entry
	{\headPrint[kana={ぷーりん\tl{˥}},ipa={puːriɴ},pos={名},acc={H}]}%
	{%
		\MeaningsBlock{%
	\ryumeaning%
		{豊年祭}{}{}%
	}%
	}%
	{\tailPrint[]}%
\entry
	{\headPrint[kana={ふーる\tl{˥}},ipa={fuːru},pos={名},acc={H}]}%
	{%
		\MeaningsBlock{%
	\ryumeaning%
		{\ruby[g]{厠}{かわや}。便所}{}{}%
	}%
	}%
	{\tailPrint[]}%
\entry
	{\headPrint[kana={ふーん\tl{˥˩}},ipa={fuːɴ},pos={動},rui={1},para={ふはぬ},acc={F}]}%
	{%
		\MeaningsBlock{%
	\ryumeaning%
		{閉じる。閉める}{}{}%
	\ryumeaning%
		{蓋をする}{}{}%
	}%
	}%
	{\tailPrint[]}%
\entry
	{\headPrint[kana={ふーん\tl{˩˥}},ipa={fuːɴ},pos={感},acc={R}]}%
	{%
		\MeaningsBlock{%
	\ryumeaning%
		{}{相手の返事が納得できない時の気のない返事}{}%
	}%
	}%
	{\tailPrint[]}%
\entry
	{\headPrint[kana={ふぁどぅまるん\tl{˥˩}},ipa={fadumaruɴ},pos={動},rui={1},para={ふあどまらぬ},acc={F}]}%
	{%
		\MeaningsBlock{%
	\ryumeaning%
		{薄暗くなる}{}{}%
	}%
	}%
	{\tailPrint[]}%
\entry
	{\headPrint[kana={ぶい\tl{˥˩}},ipa={bui},pos={名},acc={F}]}%
	{%
		\MeaningsBlock{%
	\ryumeaning%
		{甥。姪}{}{}%
	}%
	}%
	{\tailPrint[]}%
\entry
	{\headPrint[kana={ふぃーつぁるん},ipa={fiːtsaruɴ},pos={動},rui={1},para={ふぃーつぁらぬ}]}%
	{%
		\MeaningsBlock{%
	\ryumeaning%
		{身震いする。寒気がする}{}{}%
	}%
	}%
	{\tailPrint[]}%
\entry
	{\headPrint[kana={ふい\tl{˥˩}\ws{}っしん\tl{˩}},ipa={fu̥i ʃʃiɴ},pos={句},acc={F L}]}%
	{%
		\MeaningsBlock{%
	\ryumeaning%
		{振って落とす}{}{}%
	}%
	}%
	{\tailPrint[]}%
\entry
	{\headPrint[kana={ぶいふぁー\tl{˥˩}},ipa={buifaː},pos={名},acc={F}]}%
	{%
		\MeaningsBlock{%
	\ryumeaning%
		{甥。姪}{}{}%
	}%
	}%
	{\tailPrint[]}%
\entry
	{\headPrint[kana={ふか\tl{˥}},ipa={fu̥ka},pos={名},acc={H}]}%
	{%
		\MeaningsBlock{%
	\ryumeaning%
		{\ruby[g]{外}{そと}。\ruby[g]{他}{ほか}}{}{}%
	}%
	}%
	{\tailPrint[]}%
\entry
	{\headPrint[kana={ぶ\josho{}が\kako{}き},ipa={bugaki},pos={名},acc={特殊}]}%
	{%
		\MeaningsBlock{%
	\ryumeaning%
		{ヤエヤマアオキ}{植物名}{}%
	}%
	}%
	{\tailPrint[]}%
\entry
	{\headPrint[kana={ふがじ\tl{˥}},ipa={fugadʒi},pos={名},acc={H}]}%
	{%
		\MeaningsBlock{%
	\ryumeaning%
		{\ruby[g]{青}{あお}\ruby[g]{雁}{がん}\ruby[g]{皮}{ぴ}}{和紙の原料}{}%
	}%
	}%
	{\tailPrint[]}%
\entry
	{\headPrint[kana={ふかすん\tl{˥}},ipa={fu̥kasuɴ},pos={動},rui={2},para={ふかはぬ},acc={H}]}%
	{%
		\MeaningsBlock{%
	\ryumeaning%
		{（湯を）沸かす}{}{}%
	}%
	}%
	{\tailPrint[]}%
\entry
	{\headPrint[kana={ふかすん\tl{˥˩}},ipa={fu̥kasuɴ},pos={動},rui={2},para={ふかはぬ},acc={F}]}%
	{%
		\MeaningsBlock{%
	\ryumeaning%
		{（篩で）\ruby[g]{濾}{こ}す}{}{}%
	}%
	}%
	{\tailPrint[]}%
\entry
	{\headPrint[kana={ふかな},ipa={fu̥kana},pos={句}]}%
	{%
		\MeaningsBlock{%
	\ryumeaning%
		{別に。外に}{}{}%
	}%
	}%
	{\tailPrint[]}%
\entry
	{\headPrint[kana={ふか\tl{˥˩}はん},ipa={fu̥kahaɴ},pos={形},para={ふかへぬ},acc={F}]}%
	{%
		\MeaningsBlock{%
	\ryumeaning%
		{深い}{}{}%
	}%
	}%
	{\tailPrint[]}%
\entry
	{\headPrint[kana={ぶがま},ipa={bugama},pos={名}]}%
	{%
		\MeaningsBlock{%
	\ryumeaning%
		{クロヨナ}{植物名}{}%
	}%
	}%
	{\tailPrint[]}%
\entry
	{\headPrint[kana={ふかまー\tl{˥}},ipa={fu̥kamaː},pos={名},acc={H}]}%
	{%
		\MeaningsBlock{%
	\ryumeaning%
		{外孫}{}{}%
	}%
	}%
	{\tailPrint[]}%
\entry
	{\headPrint[kana={ふかまーり\tl{˥}},ipa={fu̥kamaːri},pos={名},acc={H}]}%
	{%
		\MeaningsBlock{%
	\ryumeaning%
		{外回り。外出}{}{}%
	}%
	}%
	{\tailPrint[]}%
\entry
	{\headPrint[kana={ふかむら\tl{˥˩}},ipa={fu̥kamura},pos={名},acc={F}]}%
	{%
		\MeaningsBlock{%
	\ryumeaning%
		{\ruby[g]{冨嘉}{ふか}村}{波照間島の集落名。「外村」とも書く}{}%
	}%
	}%
	{\tailPrint[]}%
\entry
	{\headPrint[kana={ふがら\tl{˥}},ipa={fugara},pos={名},acc={H}]}%
	{%
		\MeaningsBlock{%
	\ryumeaning%
		{クロツグの繊維}{黒縄の材料}{}%
	}%
	}%
	{\tailPrint[]}%
\entry
	{\headPrint[kana={ふかりるん\tl{˥}},ipa={fu̥kariruɴ},pos={動},rui={3},para={ふかーるぬ},acc={H}]}%
	{%
		\MeaningsBlock{%
	\ryumeaning%
		{（風に）吹かれる。（風に）当たる}{}{}%
	}%
	}%
	{\tailPrint[]}%
\entry
	{\headPrint[kana={ぶがりるん\tl{˩˥}},ipa={bugariruɴ},pos={動},rui={3},para={ぶがるぬ},acc={R}]}%
	{%
		\MeaningsBlock{%
	\ryumeaning%
		{疲れる。疲労する}{}{}%
	}%
	}%
	{\tailPrint[]}%
\entry
	{\headPrint[kana={ふき\tl{˥}},ipa={fu̥ki},pos={名},acc={H}]}%
	{%
		\MeaningsBlock{%
	\ryumeaning%
		{茎。芽}{}{}%
	}%
	}%
	{\tailPrint[]}%
\entry
	{\headPrint[kana={ふきーるん\tl{˥˩}},ipa={fu̥kiːruɴ},pos={動},acc={F}]}%
	{%
		\MeaningsBlock{%
	\ryumeaning%
		{気絶する。失神する}{}{}%
	}%
	}%
	{\tailPrint[]}%
\entry
	{\headPrint[kana={ふきありるん\tl{˥}},ipa={fu̥ki.ariruɴ},pos={動},rui={3},para={ふきあるぬ},acc={H}]}%
	{%
		\MeaningsBlock{%
	\ryumeaning%
		{吹き荒れる。吹きまくる}{}{}%
	}%
	}%
	{\tailPrint[]}%
\entry
	{\headPrint[kana={ふきけーし\tl{˥˩}},ipa={fu̥kikeːʃi},pos={名},acc={F}]}%
	{%
		\MeaningsBlock{%
	\ryumeaning%
		{台風の吹き返し}{}{}%
	}%
	}%
	{\tailPrint[]}%
\entry
	{\headPrint[kana={ふきつぁーるん\tl{˥˩}},ipa={fu̥kitsaːruɴ},pos={動},rui={3},acc={F}]}%
	{%
		\MeaningsBlock{%
	\ryumeaning%
		{（帯が下に下がって）だらしない姿をする}{}{}%
	}%
	}%
	{\tailPrint[]}%
\entry
	{\headPrint[kana={ふきっつぁーるん\tl{˥˩}},ipa={fu̥kittsaːruɴ},pos={動},rui={1},acc={F}]}%
	{%
		\MeaningsBlock{%
	\ryumeaning%
		{緩む}{緩くなる}{}%
	}%
	}%
	{\tailPrint[]}%
\entry
	{\headPrint[kana={ふきるん\tl{˥˩}},ipa={fu̥kiruɴ},pos={動},acc={F}]}%
	{%
		\MeaningsBlock{%
	\ryumeaning%
		{\ruby[g]{潜}{くぐ}る。\ruby[g]{潜}{くぐ}り抜ける}{}{}%
	}%
	}%
	{\tailPrint[]}%
\entry
	{\headPrint[kana={ふきん\tl{˥}},ipa={fu̥kiɴ},pos={動},rui={3},acc={H}]}%
	{%
		\MeaningsBlock{%
	\ryumeaning%
		{囀る。高鳴きする}{}{}%
	}%
	}%
	{\tailPrint[]}%
\entry
	{\headPrint[kana={ふきん\tl{˥˩}},ipa={fu̥kiɴ},pos={動},rui={3},acc={F}]}%
	{%
		\MeaningsBlock{%
	\ryumeaning%
		{卒倒する}{}{}%
	}%
	}%
	{\tailPrint[]}%
\entry
	{\headPrint[kana={ふきんじるん\tl{˥}},ipa={fu̥kindʒiruɴ},pos={動},rui={3},para={ふきんどぅぬ},acc={H}]}%
	{%
		\MeaningsBlock{%
	\ryumeaning%
		{勢いよく溢れる}{}{}%
	}%
	}%
	{\tailPrint[]}%
\entry
	{\headPrint[kana={ふく\tl{˥}},ipa={fu̥ku},pos={名},acc={H}]}%
	{%
		\MeaningsBlock{%
	\ryumeaning%
		{石垣}{}{}%
	}%
	}%
	{\tailPrint[]}%
\entry
	{\headPrint[kana={ふく\tl{˥}},ipa={fu̥ku},pos={名},acc={H}]}%
	{%
		\MeaningsBlock{%
	\ryumeaning%
		{粉。粉末}{}{}%
	}%
	}%
	{\tailPrint[]}%
\entry
	{\headPrint[kana={ふく\tl{˥}},ipa={fu̥ku},pos={名},acc={H}]}%
	{%
		\MeaningsBlock{%
	\ryumeaning%
		{肺臓}{}{}%
	}%
	}%
	{\tailPrint[]}%
\entry
	{\headPrint[kana={ふくじぃ\tl{˥}},ipa={fu̥kudʒi},pos={名},acc={H}]}%
	{%
		\MeaningsBlock{%
	\ryumeaning%
		{\ruby[g]{埃}{ほこり}}{}{}%
	}%
	}%
	{\tailPrint[]}%
\entry
	{\headPrint[kana={ふくた\tl{˥}},ipa={fu̥kuta},pos={名},acc={H}]}%
	{%
		\MeaningsBlock{%
	\ryumeaning%
		{ぼろの着物。防寒着}{}{}%
	}%
	}%
	{\tailPrint[]}%
\entry
	{\headPrint[kana={ふくな\tl{˥˩}},ipa={fu̥kuna},pos={名},acc={F}]}%
	{%
		\MeaningsBlock{%
	\ryumeaning%
		{ハルノノゲシ}{雑草名}{}%
	}%
	}%
	{\tailPrint[]}%
\entry
	{\headPrint[kana={ふくなすん\tl{˥}},ipa={fu̥kunasuɴ},pos={動},rui={2},para={ふくなはぬ},acc={H}]}%
	{%
		\MeaningsBlock{%
	\ryumeaning%
		{砕く。粉々にする。粉末にする}{}{}%
	}%
	}%
	{\tailPrint[]}%
\entry
	{\headPrint[kana={ふく\ws{}なるん},ipa={fuku naruɴ},pos={句}]}%
	{%
		\MeaningsBlock{%
	\ryumeaning%
		{粉々に割れる}{}{}%
	}%
	}%
	{\tailPrint[]}%
\entry
	{\headPrint[kana={ふくなるん\tl{˩˥}},ipa={fu̥kunaruɴ},pos={動},rui={1},para={ふくならぬ},acc={R}]}%
	{%
		\MeaningsBlock{%
	\ryumeaning%
		{粉々になる}{}{}%
	}%
	}%
	{\tailPrint[]}%
\entry
	{\headPrint[kana={ぶくぶく},ipa={bukubuku},pos={擬}]}%
	{%
		\MeaningsBlock{%
	\ryumeaning%
		{ブクブク}{泡が出るさま}{}%
	}%
	}%
	{\tailPrint[]}%
\entry
	{\headPrint[kana={ふくますん\tl{˥}},ipa={fu̥kumasuɴ},pos={動},rui={2},para={ふくまはぬ},acc={H}]}%
	{%
		\MeaningsBlock{%
	\ryumeaning%
		{含ませる}{}{}%
	}%
	}%
	{\tailPrint[]}%
\entry
	{\headPrint[kana={ふくむん},ipa={fu̥kumuɴ},pos={動}]}%
	{%
		\MeaningsBlock{%
	\ryumeaning%
		{（味などが）しみ込む}{}{}%
	}%
	}%
	{\tailPrint[]}%
\entry
	{\headPrint[kana={ふくむん\tl{˥}},ipa={fu̥kumuɴ},pos={動},rui={1},para={ふくまぬ},acc={H}]}%
	{%
		\MeaningsBlock{%
	\ryumeaning%
		{\ruby[g]{含}{ふく}む}{}{}%
	}%
	}%
	{\tailPrint[]}%
\entry
	{\headPrint[kana={ぶくむん\tl{˥}},ipa={buku̥muɴ},pos={動},rui={1},para={ぶくまぬ},acc={H}]}%
	{%
		\MeaningsBlock{%
	\ryumeaning%
		{取り混ぜる}{}{}%
	\ryumeaning%
		{ぶち込む。放り込む}{}{}%
	}%
	}%
	{\tailPrint[]}%
\entry
	{\headPrint[kana={ふくら\tl{˥˩}さーん},ipa={fu̥kurasaːɴ},pos={形},acc={F}]}%
	{%
		\MeaningsBlock{%
	\ryumeaning%
		{ふんわりする}{}{}%
	}%
	}%
	{\tailPrint[]}%
\entry
	{\headPrint[kana={ふくらすん\tl{˥˩}},ipa={fu̥kurasuɴ},pos={動},rui={2},acc={F}]}%
	{%
		\MeaningsBlock{%
	\ryumeaning%
		{膨らます}{}{}%
	}%
	}%
	{\tailPrint[]}%
\entry
	{\headPrint[kana={ふくら\tl{˥˩}はん},ipa={fu̥kurahaɴ},pos={形},para={ふくらへぬ},acc={F}]}%
	{%
		\MeaningsBlock{%
	\ryumeaning%
		{ありがとう。ありがたい}{}{}%
	}%
	}%
	{\tailPrint[]}%
\entry
	{\headPrint[kana={ふくらび\tl{˥}},ipa={fu̥kurabi},pos={名},acc={H}]}%
	{%
		\MeaningsBlock{%
	\ryumeaning%
		{カワハギ}{魚類名}{}%
	}%
	}%
	{\tailPrint[]}%
\entry
	{\headPrint[kana={ふくらますん\tl{˥˩}},ipa={fu̥kuramasuɴ},pos={動},rui={2},acc={F}]}%
	{%
		\MeaningsBlock{%
	\ryumeaning%
		{膨らませる}{}{}%
	}%
	}%
	{\tailPrint[]}%
\entry
	{\headPrint[kana={ふくりるん\tl{˥˩}},ipa={fu̥kuriruɴ},pos={動},rui={3},para={ふくるぬ},acc={F}]}%
	{%
		\MeaningsBlock{%
	\ryumeaning%
		{膨れる。腫れる}{}{}%
	}%
	}%
	{\tailPrint[]}%
\entry
	{\headPrint[kana={ふくりん\tl{˥˩}},ipa={fu̥kuriɴ},pos={動},rui={3},acc={F}]}%
	{%
		\MeaningsBlock{%
	\ryumeaning%
		{\ruby[g]{膨}{ふく}らむ}{}{}%
	\ryumeaning%
		{（水を含んで）膨れる}{}{}%
	}%
	}%
	{\tailPrint[]}%
\entry
	{\headPrint[kana={ふくる\tl{˥}},ipa={fu̥kuru},pos={名},acc={H}]}%
	{%
		\MeaningsBlock{%
	\ryumeaning%
		{\ruby[g]{袋}{ふくろ}}{}{}%
	}%
	}%
	{\tailPrint[]}%
\entry
	{\headPrint[kana={ふくるさ\tl{˥}はん},ipa={fu̥kurusahaɴ},pos={形},para={ふくらへぬ},acc={F}]}%
	{%
		\MeaningsBlock{%
	\ryumeaning%
		{ふんわりする}{}{}%
	}%
	}%
	{\tailPrint[]}%
\entry
	{\headPrint[kana={ふくるふくる},ipa={fu̥kurufu̥kuru},pos={擬}]}%
	{%
		\MeaningsBlock{%
	\ryumeaning%
		{ふっくらとした様}{}{}%
	}%
	}%
	{\tailPrint[]}%
\entry
	{\headPrint[kana={ふくるん\tl{˥˩}},ipa={fu̥kuruɴ},pos={動},rui={1},para={ふくらぬ},acc={F}]}%
	{%
		\MeaningsBlock{%
	\ryumeaning%
		{括る。縛る}{}{}%
	}%
	}%
	{\tailPrint[]}%
\entry
	{\headPrint[kana={ふくん\tl{˥˩}},ipa={fu̥kuɴ},pos={動},rui={1},para={ふかぬ},acc={F}]}%
	{%
		\MeaningsBlock{%
	\ryumeaning%
		{（屋根を）葺く}{}{}%
	}%
	}%
	{\tailPrint[]}%
\entry
	{\headPrint[kana={ふごー\tl{˥˩}},ipa={fugoː},pos={名},acc={F}]}%
	{%
		\MeaningsBlock{%
	\ryumeaning%
		{不合格}{不合格の短略}{}%
	}%
	}%
	{\tailPrint[]}%
\entry
	{\headPrint[kana={ふこーら\tl{˥˩}はーん},ipa={fukoːrahaːɴ},pos={形},acc={F}]}%
	{%
		\MeaningsBlock{%
	\ryumeaning%
		{ありがたい}{同等以下にしか用いない}{}%
	}%
	}%
	{\tailPrint[]}%
\entry
	{\headPrint[kana={ふこ\tl{˥˩}はん},ipa={fu̥kohaɴ},pos={形},para={ふこへぬ},acc={F}]}%
	{%
		\MeaningsBlock{%
	\ryumeaning%
		{（畑が）じめじめする}{}{}%
	}%
	}%
	{\tailPrint[]}%
\entry
	{\headPrint[kana={ふこん\tl{˥}},ipa={fu̥koɴ},pos={名},acc={H}]}%
	{%
		\MeaningsBlock{%
	\ryumeaning%
		{福木}{樹木名。暴風林になる}{}%
	}%
	}%
	{\tailPrint[]}%
\entry
	{\headPrint[kana={ぶざ\tl{˥}},ipa={budza},pos={名},acc={H}]}%
	{%
		\MeaningsBlock{%
	\ryumeaning%
		{百姓。平民}{王府時代の下層身分}{}%
	}%
	}%
	{\tailPrint[]}%
\entry
	{\headPrint[kana={ふさがるん\tl{˥˩}},ipa={fu̥sagaruɴ},pos={動},rui={1},para={ふさがらぬ},acc={F}]}%
	{%
		\MeaningsBlock{%
	\ryumeaning%
		{塞がる}{}{}%
	\ryumeaning%
		{邪魔になる}{}{}%
	}%
	}%
	{\tailPrint[]}%
\entry
	{\headPrint[kana={ふさぐん\tl{˥˩}},ipa={fu̥saguɴ},pos={動},acc={F}]}%
	{%
		\MeaningsBlock{%
	\ryumeaning%
		{塞ぐ。閉じる}{}{}%
	}%
	}%
	{\tailPrint[]}%
\entry
	{\headPrint[kana={ぶざすけー\tl{˥}},ipa={budzasu̥keː},pos={名},acc={H}]}%
	{%
		\MeaningsBlock{%
	\ryumeaning%
		{床の間の神棚}{}{}%
	}%
	}%
	{\tailPrint[]}%
\entry
	{\headPrint[kana={ぶさは\tl{˥}\ws{}なるん\tl{˩˥}},ipa={busaha naruɴ},pos={句},acc={H R}]}%
	{%
		\MeaningsBlock{%
	\ryumeaning%
		{大きくなる。成長する}{}{}%
	}%
	}%
	{\tailPrint[]}%
\entry
	{\headPrint[kana={ぶさ\tl{˥}はん},ipa={busahaɴ},pos={形},para={ぶさへぬ},acc={H}]}%
	{%
		\MeaningsBlock{%
	\ryumeaning%
		{大きい}{}{}%
	\ryumeaning%
		{多い}{}{}%
	}%
	}%
	{\tailPrint[]}%
\entry
	{\headPrint[kana={ぶざま\tl{˩˥}},ipa={budzama},pos={名},acc={R}]}%
	{%
		\MeaningsBlock{%
	\ryumeaning%
		{叔父。伯父}{}{}%
	}%
	}%
	{\tailPrint[]}%
\entry
	{\headPrint[kana={ふし\tl{˥}},ipa={fu̥ʃi},pos={名},acc={H}]}%
	{%
		\MeaningsBlock{%
	\ryumeaning%
		{\ruby[g]{節}{ふし}}{}{}%
	}%
	}%
	{\tailPrint[]}%
\entry
	{\headPrint[kana={ぶし\tl{˥˩}},ipa={buʃi},pos={名},acc={F}]}%
	{%
		\MeaningsBlock{%
	\ryumeaning%
		{武士。強力者}{}{}%
	}%
	}%
	{\tailPrint[]}%
\entry
	{\headPrint[kana={ふしうた\tl{˥}},ipa={fu̥ʃi.uta},pos={名},acc={H}]}%
	{%
		\MeaningsBlock{%
	\ryumeaning%
		{節歌}{}{}%
	}%
	}%
	{\tailPrint[]}%
\entry
	{\headPrint[kana={ふしがらぬ},ipa={fu̥ʃigaranu},pos={句}]}%
	{%
		\MeaningsBlock{%
	\ryumeaning%
		{我慢できない}{}{}%
	}%
	}%
	{\tailPrint[]}%
\entry
	{\headPrint[kana={ふしきるん\tl{˥}},ipa={fu̥ʃikiruɴ},pos={動},rui={3},para={ふしくぬ},acc={H}]}%
	{%
		\MeaningsBlock{%
	\ryumeaning%
		{（口に）くわえる。食いつく。（歯で）噛む。食い切る}{}{}%
	}%
	}%
	{\tailPrint[]}%
\entry
	{\headPrint[kana={ふしぐん\tl{˥}},ipa={fu̥ʃiguɴ},pos={動},rui={1},para={ふしがぬ},acc={H}]}%
	{%
		\MeaningsBlock{%
	\ryumeaning%
		{防ぐ}{}{}%
	}%
	}%
	{\tailPrint[]}%
\entry
	{\headPrint[kana={ぶしゃ\tl{˥}},ipa={buʃa},pos={名},acc={H}]}%
	{%
		\MeaningsBlock{%
	\ryumeaning%
		{長兄}{}{}%
	}%
	}%
	{\tailPrint[]}%
\entry
	{\headPrint[kana={ぶす\tl{˥}},ipa={busu},pos={名},acc={H}]}%
	{%
		\MeaningsBlock{%
	\ryumeaning%
		{潮水}{}{}%
	}%
	}%
	{\tailPrint[]}%
\entry
	{\headPrint[kana={ふすぅま\tl{˥}},ipa={fu̥suma},pos={名},acc={H}]}%
	{%
		\MeaningsBlock{%
	\ryumeaning%
		{黒島}{八重山諸島の島の名}{}%
	}%
	}%
	{\tailPrint[]}%
\entry
	{\headPrint[kana={ふすく\tl{˥˩}},ipa={fu̥suku},pos={名},acc={F}]}%
	{%
		\MeaningsBlock{%
	\ryumeaning%
		{不足}{}{}%
	}%
	}%
	{\tailPrint[]}%
\entry
	{\headPrint[kana={ぶすとぅ\ws{}なるん},ipa={busu̥tu naruɴ},pos={句},para={ふすとならぬ}]}%
	{%
		\MeaningsBlock{%
	\ryumeaning%
		{育つ。成長する}{}{}%
	}%
	}%
	{\tailPrint[]}%
\entry
	{\headPrint[kana={ぶすぷ\tl{˩˥}},ipa={busu̥pu},pos={名},acc={R}]}%
	{%
		\MeaningsBlock{%
	\ryumeaning%
		{\ruby[g]{尾}{お}。尻尾}{}{}%
	}%
	}%
	{\tailPrint[]}%
\entry
	{\headPrint[kana={ふた\tl{˥˩}},ipa={fu̥ta},pos={名},acc={F}]}%
	{%
		\MeaningsBlock{%
	\ryumeaning%
		{\ruby[g]{蓋}{ふた}}{}{}%
	}%
	}%
	{\tailPrint[]}%
\entry
	{\headPrint[kana={ふだー\tl{˥˩}},ipa={fudaː},pos={名},acc={F}]}%
	{%
		\MeaningsBlock{%
	\ryumeaning%
		{\ruby[g]{札}{ふだ}}{}{}%
	}%
	}%
	{\tailPrint[]}%
\entry
	{\headPrint[kana={ふたーち\tl{˥˩}},ipa={fu̥taːtʃi},pos={名},acc={F}]}%
	{%
		\MeaningsBlock{%
	\ryumeaning%
		{二つ}{}{}%
	}%
	}%
	{\tailPrint[]}%
\entry
	{\headPrint[kana={ふたーちぃ\tl{˥˩}ばぎ\tl{˩˥}},ipa={fu̥taːtsu̥bagi},pos={句},acc={F+R}]}%
	{%
		\MeaningsBlock{%
	\ryumeaning%
		{両方とも。二つとも}{}{}%
	}%
	}%
	{\tailPrint[]}%
\entry
	{\headPrint[kana={ふだいり\tl{˥˩}},ipa={fuda.iri},pos={名},acc={F}]}%
	{%
		\MeaningsBlock{%
	\ryumeaning%
		{札入れ。投票}{}{}%
	}%
	}%
	{\tailPrint[]}%
\entry
	{\headPrint[kana={ふたうや\tl{˥˩}},ipa={fu̥ta.uja},pos={名},acc={F}]}%
	{%
		\MeaningsBlock{%
	\ryumeaning%
		{両親}{}{}%
	}%
	}%
	{\tailPrint[]}%
\entry
	{\headPrint[kana={ふたかた\tl{˥˩}},ipa={fu̥takḁta},pos={名},acc={F}]}%
	{%
		\MeaningsBlock{%
	\ryumeaning%
		{両方}{}{}%
	}%
	}%
	{\tailPrint[]}%
\entry
	{\headPrint[kana={ふたぎな\tl{˥˩}},ipa={fu̥tagina},pos={副},acc={F}]}%
	{%
		\MeaningsBlock{%
	\ryumeaning%
		{今すぐ。直ちに}{}{}%
	}%
	}%
	{\tailPrint[]}%
\entry
	{\headPrint[kana={ふたし\tl{˥˩}},ipa={fu̥taʃi},pos={名},acc={F}]}%
	{%
		\MeaningsBlock{%
	\ryumeaning%
		{両手。両手で}{}{}%
	}%
	}%
	{\tailPrint[]}%
\entry
	{\headPrint[kana={ふたしみ\tl{˥}},ipa={fu̥taʃi̥mi},pos={名},acc={H}]}%
	{%
		\MeaningsBlock{%
	\ryumeaning%
		{\ruby[g]{守宮}{やもり}}{}{}%
	}%
	}%
	{\tailPrint[]}%
\entry
	{\headPrint[kana={ふたじろー\tl{˥˩}},ipa={fu̥tadʒiroː},pos={名},acc={F}]}%
	{%
		\MeaningsBlock{%
	\ryumeaning%
		{食べ物を入れる竹籠}{}{}%
	}%
	}%
	{\tailPrint[]}%
\entry
	{\headPrint[kana={ふたち\tl{˥˩}},ipa={fu̥tatʃi},pos={名},acc={F}]}%
	{%
		\MeaningsBlock{%
	\ryumeaning%
		{二つ}{簡略化\emb{たー}}{}%
	}%
	}%
	{\tailPrint[]}%
\entry
	{\headPrint[kana={ふたなんが\tl{˥˩}},ipa={fu̥tanaŋga},pos={名},acc={F}]}%
	{%
		\MeaningsBlock{%
	\ryumeaning%
		{二週忌}{}{}%
	}%
	}%
	{\tailPrint[]}%
\entry
	{\headPrint[kana={ふだにん\tl{˥}},ipa={fudaniɴ},pos={名},acc={H}]}%
	{%
		\MeaningsBlock{%
	\ryumeaning%
		{札人。正丁}{王府時代の課税対象（15～50歳）の成人}{}%
	}%
	}%
	{\tailPrint[]}%
\entry
	{\headPrint[kana={ふたばぎ\tl{˥˩}},ipa={fu̥tabagi},pos={句},acc={F}]}%
	{%
		\MeaningsBlock{%
	\ryumeaning%
		{二等分。折半}{}{}%
	}%
	}%
	{\tailPrint[]}%
\entry
	{\headPrint[kana={ふたびし\tl{˥˩}},ipa={fu̥tabiʃi},pos={名},acc={F}]}%
	{%
		\MeaningsBlock{%
	\ryumeaning%
		{間。合間}{}{}%
	}%
	}%
	{\tailPrint[]}%
\entry
	{\headPrint[kana={ふたまた\tl{˩˥}},ipa={fu̥tamata},pos={名},acc={R}]}%
	{%
		\MeaningsBlock{%
	\ryumeaning%
		{二股。二又}{}{}%
	}%
	}%
	{\tailPrint[]}%
\entry
	{\headPrint[kana={ふたみん\tl{˥˩}},ipa={fu̥tamiɴ},pos={名},acc={F}]}%
	{%
		\MeaningsBlock{%
	\ryumeaning%
		{両目}{}{}%
	}%
	}%
	{\tailPrint[]}%
\entry
	{\headPrint[kana={ふたり\tl{˥}},ipa={fu̥tari},pos={名},acc={H}]}%
	{%
		\MeaningsBlock{%
	\ryumeaning%
		{垢。汚れ}{}{}%
	}%
	}%
	{\tailPrint[]}%
\entry
	{\headPrint[kana={ふち\tl{˥}},ipa={fu̥tʃi},pos={名},acc={H}]}%
	{%
		\MeaningsBlock{%
	\ryumeaning%
		{櫛}{}{}%
	}%
	}%
	{\tailPrint[]}%
\entry
	{\headPrint[kana={ふち\tl{˥˩}},ipa={fu̥tʃi},pos={名},acc={F}]}%
	{%
		\MeaningsBlock{%
	\ryumeaning%
		{\ruby[g]{口}{くち}}{}{}%
	}%
	}%
	{\tailPrint[]}%
\entry
	{\headPrint[kana={ふち\tl{˥˩}},ipa={fu̥tʃi},pos={名},acc={F}]}%
	{%
		\MeaningsBlock{%
	\ryumeaning%
		{筆。毛筆}{}{}%
	}%
	}%
	{\tailPrint[]}%
\entry
	{\headPrint[kana={ぷち\tl{˥˩}},ipa={pu̥tʃi},pos={名},acc={F}]}%
	{%
		\MeaningsBlock{%
	\ryumeaning%
		{星}{}{}%
	}%
	}%
	{\tailPrint[]}%
\entry
	{\headPrint[kana={ふちぃ\tl{˥}},ipa={fu̥tʃi},pos={名},acc={H}]}%
	{%
		\MeaningsBlock{%
	\ryumeaning%
		{\ruby[g]{草鞋}{わらじ}。靴}{藁製は\emb{ばら-ふちぃ}と言う}{}%
	}%
	}%
	{\tailPrint[]}%
\entry
	{\headPrint[kana={ふちぃ\tl{˥˩}},ipa={fu̥tʃi},pos={名},acc={F}]}%
	{%
		\MeaningsBlock{%
	\ryumeaning%
		{\ruby[g]{縁}{ふち}}{}{}%
	}%
	}%
	{\tailPrint[]}%
\entry
	{\headPrint[kana={ふちぃかろ\tl{˥˩}はん},ipa={fu̥tʃikḁrohaɴ},pos={形},acc={F}]}%
	{%
		\MeaningsBlock{%
	\ryumeaning%
		{口が軽い}{}{}%
	}%
	}%
	{\tailPrint[]}%
\entry
	{\headPrint[kana={ふちぃくぱはん},ipa={fu̥tʃiku̥pahaɴ},pos={形}]}%
	{%
		\MeaningsBlock{%
	\ryumeaning%
		{口下手だ。訥弁だ}{}{}%
	\ryumeaning%
		{口が重い。吃りがちだ}{}{}%
	}%
	}%
	{\tailPrint[]}%
\entry
	{\headPrint[kana={ふちぃぬ\tl{˥˩}\ws{}ぱん\tl{˥}},ipa={fu̥tʃinu paɴ},pos={句},acc={F H}]}%
	{%
		\MeaningsBlock{%
	\ryumeaning%
		{歯}{}{}%
	}%
	}%
	{\tailPrint[note={直訳「口の歯」。同音同アクセントの\emb{ぱん}「足」と区別するためか}]}%
\entry
	{\headPrint[kana={ふちぃやにっしゃーん},ipa={fu̥tʃijaniʃʃaːɴ},pos={形}]}%
	{%
		\MeaningsBlock{%
	\ryumeaning%
		{口汚い}{言うことが汚い}{}%
	}%
	}%
	{\tailPrint[]}%
\entry
	{\headPrint[kana={ふちぃ\tl{˥˩}やば\tl{˩˥}はん},ipa={fu̥tʃijabahaɴ},pos={形},acc={F+R}]}%
	{%
		\MeaningsBlock{%
	\ryumeaning%
		{口達者だ。能弁だ}{}{}%
	}%
	}%
	{\tailPrint[]}%
\entry
	{\headPrint[kana={ふちくる},ipa={fu̥tʃikuru},pos={名}]}%
	{%
		\MeaningsBlock{%
	\ryumeaning%
		{\ruby[g]{懐}{ふところ}}{}{}%
	}%
	}%
	{\tailPrint[]}%
\entry
	{\headPrint[kana={ふちぬ\tl{˥˩}\ws{}すぱ\tl{˥}},ipa={fu̥tʃinu su̥pa},pos={句},acc={F H}]}%
	{%
		\MeaningsBlock{%
	\ryumeaning%
		{\ruby[g]{唇}{くちびる}}{直訳「口の唇」}{}%
	}%
	}%
	{\tailPrint[]}%
\entry
	{\headPrint[kana={ふちめー\tl{˥}},ipa={fu̥tʃimeː},pos={名},acc={H}]}%
	{%
		\MeaningsBlock{%
	\ryumeaning%
		{台所。炊事場}{}{}%
	}%
	}%
	{\tailPrint[]}%
\entry
	{\headPrint[kana={ふち\tl{˥˩}やぶら\tl{˩˥}},ipa={fu̥tʃijabura},pos={名},acc={F+R}]}%
	{%
		\MeaningsBlock{%
	\ryumeaning%
		{食欲がないこと}{}{}%
	}%
	}%
	{\tailPrint[]}%
\entry
	{\headPrint[kana={ふちり\tl{˥}},ipa={fu̥tʃiri},pos={名},acc={H}]}%
	{%
		\MeaningsBlock{%
	\ryumeaning%
		{薬。薬品}{}{}%
	}%
	}%
	{\tailPrint[]}%
\entry
	{\headPrint[kana={ぷちり\tl{˥}},ipa={pu̥tʃiri},pos={名},acc={H}]}%
	{%
		\MeaningsBlock{%
	\ryumeaning%
		{稲光。稲妻}{}{}%
	}%
	}%
	{\tailPrint[]}%
\entry
	{\headPrint[kana={ぷちんた},ipa={pu̥tʃinta},pos={擬}]}%
	{%
		\MeaningsBlock{%
	\ryumeaning%
		{ぷつんと。ぷつりと。さくりと}{糸や縄が切れるさま}{}%
	}%
	}%
	{\tailPrint[]}%
\entry
	{\headPrint[kana={ぷちんた},ipa={putʃinta},pos={擬}]}%
	{%
		\MeaningsBlock{%
	\ryumeaning%
		{ぽきりと。ぽきっと}{軽くものが折れるさま}{}%
	}%
	}%
	{\tailPrint[]}%
\entry
	{\headPrint[kana={ふつ\tl{˥}},ipa={fu̥tsu},pos={名},acc={H}]}%
	{%
		\MeaningsBlock{%
	\ryumeaning%
		{\ruby[g]{糞}{くそ}}{}{}%
	}%
	}%
	{\tailPrint[]}%
\entry
	{\headPrint[kana={ぷつ\tl{˥˩}},ipa={pu̥tsu},pos={名},acc={F}]}%
	{%
		\MeaningsBlock{%
	\ryumeaning%
		{\ruby[g]{臍}{へそ}}{}{}%
	}%
	}%
	{\tailPrint[]}%
\entry
	{\headPrint[kana={ふつぁ\tl{˥}},ipa={fu̥tsa},pos={名},acc={H}]}%
	{%
		\MeaningsBlock{%
	\ryumeaning%
		{草。雑草}{}{}%
	}%
	}%
	{\tailPrint[]}%
\entry
	{\headPrint[kana={ふつぁーすん\tl{˥}},ipa={fu̥tsaːsuɴ},pos={動},rui={2},para={ふつぁはぬ},acc={H}]}%
	{%
		\MeaningsBlock{%
	\ryumeaning%
		{煮返す。再び沸騰させる}{}{}%
	}%
	}%
	{\tailPrint[]}%
\entry
	{\headPrint[kana={ふつぁあん},ipa={fu̥tsa.aɴ},pos={名}]}%
	{%
		\MeaningsBlock{%
	\ryumeaning%
		{建て網の一種}{}{}%
	}%
	}%
	{\tailPrint[]}%
\entry
	{\headPrint[kana={ふつぁ\tl{˥}はん},ipa={fu̥tsahaɴ},pos={形},para={ふつぁへぬ},acc={H}]}%
	{%
		\MeaningsBlock{%
	\ryumeaning%
		{臭い}{悪臭がする}{}%
	}%
	}%
	{\tailPrint[]}%
\entry
	{\headPrint[kana={ふつぁはん},ipa={futsahaɴ},pos={形},para={ふつぁへぬ}]}%
	{%
		\MeaningsBlock{%
	\ryumeaning%
		{欲しい。欲しがる}{}{}%
	}%
	}%
	{\tailPrint[]}%
\entry
	{\headPrint[kana={ふつぁび},ipa={fu̥tsabi},pos={名},acc={-}]}%
	{%
		\MeaningsBlock{%
	\ryumeaning%
		{\ruby[g]{楔}{くさび}}{}{}%
	}%
	}%
	{\tailPrint[]}%
\entry
	{\headPrint[kana={ふつぁび\tl{˥}},ipa={fu̥tsabi},pos={名},acc={H}]}%
	{%
		\MeaningsBlock{%
	\ryumeaning%
		{ベラ}{魚類名}{}%
	}%
	}%
	{\tailPrint[]}%
\entry
	{\headPrint[kana={ふつぁまら\tl{˥}},ipa={fu̥tsamara},pos={名},acc={H}]}%
	{%
		\MeaningsBlock{%
	\ryumeaning%
		{雨乞い祈願の仮面神}{}{}%
	}%
	}%
	{\tailPrint[]}%
\entry
	{\headPrint[kana={ふつぁらし\tl{˥}},ipa={fu̥tsaraʃi},pos={名},acc={H}]}%
	{%
		\MeaningsBlock{%
	\ryumeaning%
		{草むら}{}{}%
	}%
	}%
	{\tailPrint[]}%
\entry
	{\headPrint[kana={ふつぁらすん\tl{˥}},ipa={fu̥tsarasuɴ},pos={動},rui={2},para={ふつぁらはぬ},acc={H}]}%
	{%
		\MeaningsBlock{%
	\ryumeaning%
		{腐らせる}{}{}%
	}%
	}%
	{\tailPrint[]}%
\entry
	{\headPrint[kana={ふつぁりかー},ipa={fu̥tsarikaː},pos={名}]}%
	{%
		\MeaningsBlock{%
	\ryumeaning%
		{腐臭。悪臭}{}{}%
	}%
	}%
	{\tailPrint[]}%
\entry
	{\headPrint[kana={ふつぁりん\tl{˥}},ipa={fu̥tsariɴ},pos={動},rui={3},acc={H}]}%
	{%
		\MeaningsBlock{%
	\ryumeaning%
		{腐る。腐敗する}{}{}%
	}%
	}%
	{\tailPrint[]}%
\entry
	{\headPrint[kana={ふつぁるん\tl{˥˩}},ipa={fu̥tsaruɴ},pos={動},para={ふつわらぬ},acc={F}]}%
	{%
		\MeaningsBlock{%
	\ryumeaning%
		{塞がる}{}{}%
	\ryumeaning%
		{（目を）閉じる}{}{}%
	}%
	}%
	{\tailPrint[]}%
\entry
	{\headPrint[kana={ふつぁんだに\tl{˥}},ipa={fu̥tsandani},pos={名},acc={H}]}%
	{%
		\MeaningsBlock{%
	\ryumeaning%
		{草。雑草}{}{}%
	}%
	}%
	{\tailPrint[]}%
\entry
	{\headPrint[kana={ぷつぉー\tl{˥}},ipa={pu̥tsoː},pos={名},acc={H}]}%
	{%
		\MeaningsBlock{%
	\ryumeaning%
		{煙草入れ}{}{}%
	}%
	}%
	{\tailPrint[]}%
\entry
	{\headPrint[kana={ふっかばり\tl{˥}},ipa={fukkabari},pos={-},acc={H}]}%
	{%
		\MeaningsBlock{%
	\ryumeaning%
		{黒い。黒ずむ}{}{}%
	}%
	}%
	{\tailPrint[]}%
\entry
	{\headPrint[kana={ふつがら\tl{˥}},ipa={fu̥tsugara},pos={名},acc={H}]}%
	{%
		\MeaningsBlock{%
	\ryumeaning%
		{オオタニワタリ}{植物名}{}%
	}%
	}%
	{\tailPrint[]}%
\entry
	{\headPrint[kana={ふつくれー},ipa={fu̥tsukureː},pos={感}]}%
	{%
		\MeaningsBlock{%
	\ryumeaning%
		{\ruby[g]{嫌}{いや}だ。不承知だ}{「くそくらえ」の義}{}%
	}%
	}%
	{\tailPrint[]}%
\entry
	{\headPrint[kana={ぶった},ipa={butta},pos={助}]}%
	{%
		\MeaningsBlock{%
	\ryumeaning%
		{だらけ}{}{}%
	}%
	}%
	{\tailPrint[]}%
\entry
	{\headPrint[kana={ふっひ\tl{˥˩}},ipa={fu̥hi},pos={名},acc={F}]}%
	{%
		\MeaningsBlock{%
	\ryumeaning%
		{冬。冬季}{}{}%
	}%
	}%
	{\tailPrint[]}%
\entry
	{\headPrint[kana={ふっひ\ws{}ししるん},ipa={fu̥hi ʃʃiruɴ},pos={動}]}%
	{%
		\MeaningsBlock{%
	\ryumeaning%
		{振り捨てる。振り払う}{}{}%
	}%
	}%
	{\tailPrint[]}%
\entry
	{\headPrint[kana={ふっふぁ\tl{˥}},ipa={ffa},pos={名},acc={H}]}%
	{%
		\MeaningsBlock{%
	\ryumeaning%
		{倉}{穀物倉}{}%
	}%
	}%
	{\tailPrint[]}%
\entry
	{\headPrint[kana={ふっふぁ\tl{˥}},ipa={ffa},pos={名},acc={H}]}%
	{%
		\MeaningsBlock{%
	\ryumeaning%
		{\ruby[g]{鞍}{くら}}{}{}%
	}%
	}%
	{\tailPrint[]}%
\entry
	{\headPrint[kana={ふつ\tl{˥}ふつぁ\tl{˥}はーん},ipa={fu̥tsufu̥tsahaːɴ},pos={形},acc={H+H}]}%
	{%
		\MeaningsBlock{%
	\ryumeaning%
		{糞臭い}{糞の臭いがする}{}%
	}%
	}%
	{\tailPrint[]}%
\entry
	{\headPrint[kana={ふっふん\tl{˥}},ipa={ffuɴ},pos={動},rui={1},para={ふっふぁぬ},acc={H}]}%
	{%
		\MeaningsBlock{%
	\ryumeaning%
		{（雨が）降る}{}{}%
	}%
	}%
	{\tailPrint[]}%
\entry
	{\headPrint[kana={ふっふん\tl{˥˩}},ipa={fu̥fuɴ},pos={動},rui={1},para={ふっわぬ},acc={F}]}%
	{%
		\MeaningsBlock{%
	\ryumeaning%
		{振る}{}{}%
	}%
	}%
	{\tailPrint[]}%
\entry
	{\headPrint[kana={ふっふん\tl{˥˩}},ipa={ffuɴ},pos={動},rui={1},para={ふっわぬ},acc={F}]}%
	{%
		\MeaningsBlock{%
	\ryumeaning%
		{揺れる。振動する}{}{}%
	}%
	}%
	{\tailPrint[]}%
\entry
	{\headPrint[kana={ふつ\tl{˥}\ws{}まりぽーるん\tl{˩˥}},ipa={fu̥tsu maripoːruɴ},pos={句},acc={H R}]}%
	{%
		\MeaningsBlock{%
	\ryumeaning%
		{（大便を）排泄し散らかす}{}{}%
	}%
	}%
	{\tailPrint[]}%
\entry
	{\headPrint[kana={ふつ\tl{˥}まるん\tl{˩˥}},ipa={fu̥tsumaruɴ},pos={句},para={ふつまらぬ},acc={H R}]}%
	{%
		\MeaningsBlock{%
	\ryumeaning%
		{糞を垂れる。排便する}{}{}%
	}%
	}%
	{\tailPrint[]}%
\entry
	{\headPrint[kana={ふつるん\tl{˥}},ipa={fu̥tsuruɴ},pos={動},rui={1},para={ふつらぬ},acc={H}]}%
	{%
		\MeaningsBlock{%
	\ryumeaning%
		{漁る。探し回る}{}{}%
	}%
	}%
	{\tailPrint[]}%
\entry
	{\headPrint[kana={ふつん\tl{˥}},ipa={fu̥tsuɴ},pos={動},rui={2},para={ふつあぬ},acc={H}]}%
	{%
		\MeaningsBlock{%
	\ryumeaning%
		{沸騰する。煮立つ}{}{}%
	}%
	}%
	{\tailPrint[]}%
\entry
	{\headPrint[kana={ぷつん\tl{˥}},ipa={pu̥tsuɴ},pos={動},rui={2},para={ぷつぁぬ},acc={H}]}%
	{%
		\MeaningsBlock{%
	\ryumeaning%
		{干す。乾かす}{}{}%
	}%
	}%
	{\tailPrint[]}%
\entry
	{\headPrint[kana={ふてー\tl{˥˩}},ipa={futeː},pos={名},acc={F}]}%
	{%
		\MeaningsBlock{%
	\ryumeaning%
		{\ruby[g]{額}{ひたい}}{}{}%
	}%
	}%
	{\tailPrint[]}%
\entry
	{\headPrint[kana={ぶとぅ\tl{˥˩}},ipa={butu},pos={名},acc={F}]}%
	{%
		\MeaningsBlock{%
	\ryumeaning%
		{夫}{}{}%
	}%
	}%
	{\tailPrint[]}%
\entry
	{\headPrint[kana={ぷとぅぎるん},ipa={pu̥tugiruɴ},pos={動},rui={3},para={ぷとぅぐぬ}]}%
	{%
		\MeaningsBlock{%
	\ryumeaning%
		{\ruby[g]{解}{ほど}ける}{糸や結んだ物などが解ける}{}%
	}%
	}%
	{\tailPrint[]}%
\entry
	{\headPrint[kana={ぷとぅぎん\tl{˥}},ipa={pu̥tugiɴ},pos={名},acc={H}]}%
	{%
		\MeaningsBlock{%
	\ryumeaning%
		{仏}{仏様}{}%
	}%
	}%
	{\tailPrint[]}%
\entry
	{\headPrint[kana={ぷとぅぎんぬ\tl{˥}\ws{}めー\tl{˩˥}},ipa={pu̥tuginnu meː},pos={句},acc={H R}]}%
	{%
		\MeaningsBlock{%
	\ryumeaning%
		{仏の前。仏前}{}{}%
	}%
	}%
	{\tailPrint[]}%
\entry
	{\headPrint[kana={ぷとぅぐん\tl{˥}},ipa={pu̥tuguɴ},pos={動},rui={1},para={ぷとぅがぬ},acc={H}]}%
	{%
		\MeaningsBlock{%
	\ryumeaning%
		{\ruby[g]{解}{ほど}く}{}{}%
	}%
	}%
	{\tailPrint[]}%
\entry
	{\headPrint[kana={ぶとぅち\tl{˩˥}},ipa={bututʃi},pos={名},acc={R}]}%
	{%
		\MeaningsBlock{%
	\ryumeaning%
		{\ruby[g]{一昨日}{おととい}}{}{}%
	}%
	}%
	{\tailPrint[]}%
\entry
	{\headPrint[kana={ふとぅちぱん\tl{˥}},ipa={fu̥tutʃipaɴ},pos={名},acc={H}]}%
	{%
		\MeaningsBlock{%
	\ryumeaning%
		{虫歯}{}{}%
	}%
	}%
	{\tailPrint[]}%
\entry
	{\headPrint[kana={ふとぅちるん\tl{˥}},ipa={fu̥tutʃiruɴ},pos={動},rui={3},para={ふとぅとぅぬ},acc={H}]}%
	{%
		\MeaningsBlock{%
	\ryumeaning%
		{\ruby[g]{朽}{く}ちる}{}{}%
	}%
	}%
	{\tailPrint[]}%
\entry
	{\headPrint[kana={ふとぅふとぅ},ipa={futufutu},pos={擬}]}%
	{%
		\MeaningsBlock{%
	\ryumeaning%
		{どきどき。ぶるぶる}{不安、心配、怒り、恐怖などで興奮する様}{}%
	}%
	}%
	{\tailPrint[]}%
\entry
	{\headPrint[kana={ぶとぅぶとぅ},ipa={butubutu},pos={擬}]}%
	{%
		\MeaningsBlock{%
	\ryumeaning%
		{沸騰する様}{物が煮えたぎる様}{}%
	}%
	}%
	{\tailPrint[]}%
\entry
	{\headPrint[kana={ぶとぅむち},ipa={butumutʃi},pos={名}]}%
	{%
		\MeaningsBlock{%
	\ryumeaning%
		{結婚する。嫁に行く}{「夫を持つ」の義}{}%
	}%
	}%
	{\tailPrint[]}%
\entry
	{\headPrint[kana={ぶどぅり\tl{˥˩}},ipa={buduri},pos={名},acc={F}]}%
	{%
		\MeaningsBlock{%
	\ryumeaning%
		{踊り。舞踊}{}{}%
	}%
	}%
	{\tailPrint[]}%
\entry
	{\headPrint[kana={ぶどぅるん\tl{˥˩}},ipa={buduruɴ},pos={動},acc={F}]}%
	{%
		\MeaningsBlock{%
	\ryumeaning%
		{踊る}{}{}%
	}%
	}%
	{\tailPrint[]}%
\entry
	{\headPrint[kana={ぶなが\tl{˥}},ipa={bunaga},pos={名},acc={H}]}%
	{%
		\MeaningsBlock{%
	\ryumeaning%
		{祭。祭礼}{}{}%
	}%
	}%
	{\tailPrint[]}%
\entry
	{\headPrint[kana={ふなぐん\tl{˥}},ipa={fu̥naguɴ},pos={動},rui={1},para={ふながぬ},acc={H}]}%
	{%
		\MeaningsBlock{%
	\ryumeaning%
		{叩く。砕く}{}{}%
	}%
	}%
	{\tailPrint[]}%
\entry
	{\headPrint[kana={ふなたぴ\tl{˥}},ipa={fu̥natḁpi},pos={名},acc={H}]}%
	{%
		\MeaningsBlock{%
	\ryumeaning%
		{船旅}{}{}%
	}%
	}%
	{\tailPrint[]}%
\entry
	{\headPrint[kana={ふなだま\tl{˥}},ipa={fu̥nadama},pos={名},acc={H}]}%
	{%
		\MeaningsBlock{%
	\ryumeaning%
		{船霊}{}{}%
	}%
	}%
	{\tailPrint[]}%
\entry
	{\headPrint[kana={ふなだまり\tl{˥}},ipa={fu̥nadamari},pos={名},acc={H}]}%
	{%
		\MeaningsBlock{%
	\ryumeaning%
		{船の停泊場所。船溜り}{}{}%
	}%
	}%
	{\tailPrint[]}%
\entry
	{\headPrint[kana={ふなでーぐ},ipa={fu̥nadeːgu},pos={名}]}%
	{%
		\MeaningsBlock{%
	\ryumeaning%
		{船大工}{}{}%
	}%
	}%
	{\tailPrint[]}%
\entry
	{\headPrint[kana={ふなどぅり\tl{˥˩}},ipa={fu̥naduri},pos={名},acc={F}]}%
	{%
		\MeaningsBlock{%
	\ryumeaning%
		{メジロ}{鳥の名}{}%
	}%
	}%
	{\tailPrint[]}%
\entry
	{\headPrint[kana={ふなはこー},ipa={fu̥nahḁkoː},pos={名}]}%
	{%
		\MeaningsBlock{%
	\ryumeaning%
		{水夫。船員}{}{}%
	}%
	}%
	{\tailPrint[]}%
\entry
	{\headPrint[kana={ふなばに\tl{˥}},ipa={fu̥nabani},pos={名},acc={H}]}%
	{%
		\MeaningsBlock{%
	\ryumeaning%
		{船の横板。横舷}{}{}%
	}%
	}%
	{\tailPrint[]}%
\entry
	{\headPrint[kana={ぶなび},ipa={bunabi},pos={名}]}%
	{%
		\MeaningsBlock{%
	\ryumeaning%
		{よなべ。夜業}{}{}%
	}%
	}%
	{\tailPrint[]}%
\entry
	{\headPrint[kana={ぶなびやー\tl{˥}},ipa={bunabijaː},pos={名},acc={H}]}%
	{%
		\MeaningsBlock{%
	\ryumeaning%
		{よなべ小屋}{王府時代の夜業小屋}{}%
	}%
	}%
	{\tailPrint[]}%
\entry
	{\headPrint[kana={ふなぶ\tl{˥}},ipa={fu̥nabu},pos={名},acc={H}]}%
	{%
		\MeaningsBlock{%
	\ryumeaning%
		{蜜柑}{}{}%
	}%
	}%
	{\tailPrint[]}%
\entry
	{\headPrint[kana={ぶなり\tl{˥˩}},ipa={bunari},pos={名},acc={F}]}%
	{%
		\MeaningsBlock{%
	\ryumeaning%
		{姉妹}{}{}%
	}%
	}%
	{\tailPrint[]}%
\entry
	{\headPrint[kana={ぶなりかん\tl{˥˩}},ipa={bunarikaɴ},pos={名},acc={F}]}%
	{%
		\MeaningsBlock{%
	\ryumeaning%
		{姉妹神}{}{}%
	}%
	}%
	{\tailPrint[]}%
\entry
	{\headPrint[kana={ふに\tl{˥}},ipa={fu̥ni},pos={名},acc={H}]}%
	{%
		\MeaningsBlock{%
	\ryumeaning%
		{船。船舶}{}{}%
	}%
	}%
	{\tailPrint[]}%
\entry
	{\headPrint[kana={ぷに\tl{˥}},ipa={pu̥ni},pos={名},acc={H}]}%
	{%
		\MeaningsBlock{%
	\ryumeaning%
		{骨。骨格}{}{}%
	}%
	}%
	{\tailPrint[]}%
\entry
	{\headPrint[kana={ぷにやしみ\tl{˥}},ipa={pu̥nijaʃi̥mi},pos={名},acc={H}]}%
	{%
		\MeaningsBlock{%
	\ryumeaning%
		{骨休め。休養}{}{}%
	}%
	}%
	{\tailPrint[]}%
\entry
	{\headPrint[kana={ふにんじるん},ipa={fu̥nindʒiruɴ},pos={句},para={ふにんどぬ},acc={H}]}%
	{%
		\MeaningsBlock{%
	\ryumeaning%
		{出港する}{船が出るの義}{}%
	}%
	}%
	{\tailPrint[]}%
\entry
	{\headPrint[kana={ぶぬ\tl{˩˥}},ipa={bunu},pos={名},acc={R}]}%
	{%
		\MeaningsBlock{%
	\ryumeaning%
		{\ruby[g]{斧}{おの}}{}{}%
	}%
	}%
	{\tailPrint[]}%
\entry
	{\headPrint[kana={ふねー\tl{˥}},ipa={fu̥neː},pos={名},acc={H}]}%
	{%
		\MeaningsBlock{%
	\ryumeaning%
		{船酔い}{}{}%
	}%
	}%
	{\tailPrint[]}%
\entry
	{\headPrint[kana={ふのーら\tl{˥˩}},ipa={fu̥noːra},pos={名},acc={F}]}%
	{%
		\MeaningsBlock{%
	\ryumeaning%
		{船浦}{船の停泊地}{}%
	}%
	}%
	{\tailPrint[]}%
\entry
	{\headPrint[kana={ぶば\tl{˥}},ipa={buba},pos={名},acc={H}]}%
	{%
		\MeaningsBlock{%
	\ryumeaning%
		{伯母。叔母}{}{}%
	}%
	}%
	{\tailPrint[]}%
\entry
	{\headPrint[kana={ふばむに\tl{˥}},ipa={fubamuni},pos={名},acc={H}]}%
	{%
		\MeaningsBlock{%
	\ryumeaning%
		{冗談}{}{}%
	}%
	}%
	{\tailPrint[]}%
\entry
	{\headPrint[kana={ふひつぁるん},ipa={fuhi̥tsaruɴ},pos={動},rui={1},para={ふひつぁらぬ},acc={-}]}%
	{%
		\MeaningsBlock{%
	\ryumeaning%
		{寒くて震える}{}{}%
	}%
	}%
	{\tailPrint[]}%
\entry
	{\headPrint[kana={ふみあぎるん\tl{˥}},ipa={fu̥mi.agiruɴ},pos={動},acc={H}]}%
	{%
		\MeaningsBlock{%
	\ryumeaning%
		{褒めあげる。褒め讃える}{}{}%
	}%
	}%
	{\tailPrint[]}%
\entry
	{\headPrint[kana={ふみるん\tl{˥}},ipa={fu̥miruɴ},pos={動},rui={3},para={ふむぬ},acc={H}]}%
	{%
		\MeaningsBlock{%
	\ryumeaning%
		{褒める}{}{}%
	}%
	}%
	{\tailPrint[]}%
\entry
	{\headPrint[kana={ふむむぬ\tl{˥˩}},ipa={fu̥mumunu},pos={名},acc={F}]}%
	{%
		\MeaningsBlock{%
	\ryumeaning%
		{履物}{}{}%
	}%
	}%
	{\tailPrint[]}%
\entry
	{\headPrint[kana={ふむん},ipa={fu̥muɴ},pos={動},rui={1},para={ふまぬ}]}%
	{%
		\MeaningsBlock{%
	\ryumeaning%
		{（足で）踏む}{}{}%
	}%
	}%
	{\tailPrint[]}%
\entry
	{\headPrint[kana={ふむん\tl{˥}},ipa={fu̥muɴ},pos={動},rui={1},para={ふまぬ},acc={H}]}%
	{%
		\MeaningsBlock{%
	\ryumeaning%
		{編む}{竹などで編む時にいう}{}%
	}%
	}%
	{\tailPrint[]}%
\entry
	{\headPrint[kana={ふむん\tl{˥˩}},ipa={fu̥muɴ},pos={動},rui={1},para={ふまぬ},acc={F}]}%
	{%
		\MeaningsBlock{%
	\ryumeaning%
		{（履物を）はく}{}{}%
	}%
	}%
	{\tailPrint[]}%
\entry
	{\headPrint[kana={ふむん\tl{˥˩}},ipa={fu̥muɴ},pos={動},rui={1},para={ふまぬ},acc={F}]}%
	{%
		\MeaningsBlock{%
	\ryumeaning%
		{（水を）汲む}{}{}%
	}%
	}%
	{\tailPrint[]}%
\entry
	{\headPrint[kana={ふもーらすん\tl{˥}},ipa={fu̥moːrasuɴ},pos={動},rui={2},para={ふもーらはぬ},acc={H}]}%
	{%
		\MeaningsBlock{%
	\ryumeaning%
		{煙を立てる}{}{}%
	}%
	}%
	{\tailPrint[]}%
\entry
	{\headPrint[kana={ふもん\tl{˥}},ipa={fu̥moɴ},pos={名},acc={H}]}%
	{%
		\MeaningsBlock{%
	\ryumeaning%
		{雲}{}{}%
	}%
	}%
	{\tailPrint[]}%
\entry
	{\headPrint[kana={ぶや\tl{˥}},ipa={buja},pos={名},acc={H}]}%
	{%
		\MeaningsBlock{%
	\ryumeaning%
		{祖父。おじいさん}{}{}%
	}%
	}%
	{\tailPrint[]}%
\entry
	{\headPrint[kana={ふゆー\tl{˥}},ipa={fujuː},pos={名},acc={H}]}%
	{%
		\MeaningsBlock{%
	\ryumeaning%
		{怠け。不精}{}{}%
	}%
	}%
	{\tailPrint[]}%
\entry
	{\headPrint[kana={ふゆーむぬ\tl{˥}},ipa={fujuːmunu},pos={名},acc={H}]}%
	{%
		\MeaningsBlock{%
	\ryumeaning%
		{怠け者。不精者}{}{}%
	}%
	}%
	{\tailPrint[]}%
\entry
	{\headPrint[kana={ふゆな\ws{}むぬ},ipa={fujuna munu},pos={句}]}%
	{%
		\MeaningsBlock{%
	\ryumeaning%
		{不精者}{}{}%
	}%
	}%
	{\tailPrint[]}%
\entry
	{\headPrint[kana={ぶら},ipa={bura},pos={名}]}%
	{%
		\MeaningsBlock{%
	\ryumeaning%
		{（楽器の）法螺}{}{}%
	}%
	}%
	{\tailPrint[]}%
\entry
	{\headPrint[kana={ぶら\tl{˥˩}},ipa={bura},pos={名},acc={F}]}%
	{%
		\MeaningsBlock{%
	\ryumeaning%
		{\ruby[g]{釣}{つる}\ruby[g]{瓶}{べ}}{}{}%
	}%
	}%
	{\tailPrint[]}%
\entry
	{\headPrint[kana={ぶら\tl{˩˥}},ipa={bura},pos={名},acc={R}]}%
	{%
		\MeaningsBlock{%
	\ryumeaning%
		{\ruby[g]{法}{ほ}\ruby[g]{螺}{ら}貝}{楽器の法螺}{}%
	}%
	}%
	{\tailPrint[]}%
\entry
	{\headPrint[kana={ふらー\tl{˥˩}},ipa={furaː},pos={名},acc={F}]}%
	{%
		\MeaningsBlock{%
	\ryumeaning%
		{馬鹿。馬鹿者}{}{}%
	}%
	}%
	{\tailPrint[]}%
\entry
	{\headPrint[kana={ふらふら},ipa={furafura},pos={擬}]}%
	{%
		\MeaningsBlock{%
	\ryumeaning%
		{ふらふらっと}{目眩がする様。気が遠くなる様}{}%
	}%
	}%
	{\tailPrint[]}%
\entry
	{\headPrint[kana={ぶりー\tl{˩˥}},ipa={buriː},pos={名},acc={R}]}%
	{%
		\MeaningsBlock{%
	\ryumeaning%
		{失礼。無礼}{}{}%
	}%
	}%
	{\tailPrint[]}%
\entry
	{\headPrint[kana={ぷりばーり\tl{˥˩}},ipa={pu̥ribaːri},pos={名},acc={F}]}%
	{%
		\MeaningsBlock{%
	\ryumeaning%
		{馬鹿笑い。大笑い}{}{}%
	}%
	}%
	{\tailPrint[]}%
\entry
	{\headPrint[kana={ふりはきん\tl{˥˩}},ipa={fu̥rihakiɴ},pos={動},rui={3},para={ふりはくぬ},acc={F}]}%
	{%
		\MeaningsBlock{%
	\ryumeaning%
		{振りかける}{}{}%
	}%
	}%
	{\tailPrint[]}%
\entry
	{\headPrint[kana={ぷ\josho{}り\kako{}ぷり\ws{}すん\tl{˥˩}},ipa={pu̥ripu̥ri suɴ},pos={句},rui={2},para={ぷりぷりさぬ},acc={重複}]}%
	{%
		\MeaningsBlock{%
	\ryumeaning%
		{ぼんやりする}{}{}%
	}%
	}%
	{\tailPrint[]}%
\entry
	{\headPrint[kana={ふりまーすん\tl{˥˩}},ipa={fu̥rimaːsuɴ},pos={動},rui={2b},para={ふりまーさぬ},acc={F}]}%
	{%
		\MeaningsBlock{%
	\ryumeaning%
		{振り回す}{}{}%
	}%
	}%
	{\tailPrint[]}%
\entry
	{\headPrint[kana={ぷりむに\tl{˥˩}},ipa={pu̥rimuni},pos={名},acc={F}]}%
	{%
		\MeaningsBlock{%
	\ryumeaning%
		{虚言。嘘}{}{}%
	}%
	}%
	{\tailPrint[]}%
\entry
	{\headPrint[kana={ぷりむぬ\tl{˥˩}},ipa={pu̥rimunu},pos={名},acc={F}]}%
	{%
		\MeaningsBlock{%
	\ryumeaning%
		{気違い。馬鹿者}{}{}%
	}%
	}%
	{\tailPrint[]}%
\entry
	{\headPrint[kana={ぶりるん\tl{˩˥}},ipa={buriruɴ},pos={動},rui={3},para={ぶるぬ},acc={R}]}%
	{%
		\MeaningsBlock{%
	\ryumeaning%
		{（波が）荒れる}{}{}%
	}%
	}%
	{\tailPrint[]}%
\entry
	{\headPrint[kana={ぶりるん\tl{˩˥}},ipa={buriruɴ},pos={動},rui={3},para={ぶるぬ},acc={R}]}%
	{%
		\MeaningsBlock{%
	\ryumeaning%
		{折れる}{}{}%
	}%
	}%
	{\tailPrint[]}%
\entry
	{\headPrint[kana={ぷりるん\tl{˥˩}},ipa={pu̥riruɴ},pos={動},rui={3},para={ぷるぬ},acc={F}]}%
	{%
		\MeaningsBlock{%
	\ryumeaning%
		{惚れる}{}{}%
	\ryumeaning%
		{気が狂う}{}{}%
	}%
	}%
	{\tailPrint[]}%
\entry
	{\headPrint[kana={ふるばすん\tl{˥˩}},ipa={fu̥rubasuɴ},pos={動},rui={2b},para={ふるばさぬ},acc={F}]}%
	{%
		\MeaningsBlock{%
	\ryumeaning%
		{滅ぼす}{}{}%
	}%
	}%
	{\tailPrint[]}%
\entry
	{\headPrint[kana={ふるぶん\tl{˥˩}},ipa={fu̥rubuɴ},pos={動},rui={1},para={ふるばぬ},acc={F}]}%
	{%
		\MeaningsBlock{%
	\ryumeaning%
		{滅ぶ}{}{}%
	}%
	}%
	{\tailPrint[]}%
\entry
	{\headPrint[kana={ふるや\tl{˥}},ipa={furuja},pos={名},acc={H}]}%
	{%
		\MeaningsBlock{%
	\ryumeaning%
		{\ruby[g]{厠}{かわや}。便所}{}{}%
	}%
	}%
	{\tailPrint[]}%
\entry
	{\headPrint[kana={ぶるん\tl{˥˩}},ipa={buruɴ},pos={動},rui={1},para={ぶらぬ},acc={F}]}%
	{%
		\MeaningsBlock{%
	\ryumeaning%
		{（木の実などを）摘む}{}{}%
	}%
	}%
	{\tailPrint[]}%
\entry
	{\headPrint[kana={ぶるん\tl{˩˥}},ipa={buruɴ},pos={動},rui={1},para={ぶらぬ},acc={R}]}%
	{%
		\MeaningsBlock{%
	\ryumeaning%
		{折る。収穫する}{}{}%
	}%
	}%
	{\tailPrint[]}%
\entry
	{\headPrint[kana={ぶるん\tl{˩˥}},ipa={buruɴ},pos={動},rui={1},para={ぶらぬ},acc={R}]}%
	{%
		\MeaningsBlock{%
	\ryumeaning%
		{摘み折る。もぎ取る}{}{}%
	}%
	}%
	{\tailPrint[]}%
\entry
	{\headPrint[kana={ぷるん\tl{˥}},ipa={pu̥ruɴ},pos={動},rui={1},para={ぷらぬ},acc={H}]}%
	{%
		\MeaningsBlock{%
	\ryumeaning%
		{掘る。彫る}{}{}%
	}%
	}%
	{\tailPrint[]}%
\entry
	{\headPrint[kana={ふゎどぅまり\tl{˥˩}\ws{}なるん\tl{˩˥}},ipa={ffwadumari naruɴ},pos={句},acc={F R}]}%
	{%
		\MeaningsBlock{%
	\ryumeaning%
		{暗くなる}{}{}%
	\ryumeaning%
		{曇る}{}{}%
	}%
	}%
	{\tailPrint[]}%
\entry
	{\headPrint[kana={ふわ\tl{˥˩}はん},ipa={fu̥wahaɴ},pos={形},para={ふわぁへぬ},acc={F}]}%
	{%
		\MeaningsBlock{%
	\ryumeaning%
		{暗い}{}{}%
	}%
	}%
	{\tailPrint[]}%
\entry
	{\headPrint[kana={ふゎむん\tl{˥˩}},ipa={fwamuɴ},pos={動},rui={1},para={ふゎーまぬ},acc={F}]}%
	{%
		\MeaningsBlock{%
	\ryumeaning%
		{暗くなる。（日が）暮れる}{}{}%
	\ryumeaning%
		{曇る}{}{}%
	}%
	}%
	{\tailPrint[]}%
\entry
	{\headPrint[kana={ふん\tl{˥}},ipa={fuɴ},pos={名},acc={H}]}%
	{%
		\MeaningsBlock{%
	\ryumeaning%
		{運}{}{}%
	}%
	}%
	{\tailPrint[]}%
\entry
	{\headPrint[kana={ふん\tl{˥˩}},ipa={fuɴ},pos={名},acc={F}]}%
	{%
		\MeaningsBlock{%
	\ryumeaning%
		{\ruby[g]{釘}{くぎ}}{}{}%
	}%
	}%
	{\tailPrint[]}%
\entry
	{\headPrint[kana={ぶん\tl{˥˩}},ipa={buɴ},pos={動},rui={irr},para={ぶらぬ},acc={F}]}%
	{%
		\MeaningsBlock{%
	\ryumeaning%
		{居る}{\emb{ぶなー}「居るか」}{}%
	}%
	}%
	{\tailPrint[]}%
\entry
	{\headPrint[kana={ぶん\tl{˩˥}},ipa={buɴ},pos={名},acc={R}]}%
	{%
		\MeaningsBlock{%
	\ryumeaning%
		{恩。恩義}{}{}%
	}%
	}%
	{\tailPrint[note={「恩」の転訛}]}%
\entry
	{\headPrint[kana={ぶん\tl{˩˥}かいすん\tl{˥}},ipa={bunkaisuɴ},pos={動},acc={R+H}]}%
	{%
		\MeaningsBlock{%
	\ryumeaning%
		{恩を返す}{}{}%
	}%
	}%
	{\tailPrint[]}%
\entry
	{\headPrint[kana={ぶんぎ\tl{˩˥}},ipa={buŋgi},pos={名},acc={R}]}%
	{%
		\MeaningsBlock{%
	\ryumeaning%
		{恩。恩義}{}{}%
	}%
	}%
	{\tailPrint[note={「恩義」の転訛}]}%
\entry
	{\headPrint[kana={ふんさまるん\tl{˥˩}},ipa={funsamaruɴ},pos={動},acc={F}]}%
	{%
		\MeaningsBlock{%
	\ryumeaning%
		{ふん縛る}{}{}%
	}%
	}%
	{\tailPrint[]}%
\entry
	{\headPrint[kana={ふんしー\tl{˥}},ipa={funʃiː},pos={名},acc={H}]}%
	{%
		\MeaningsBlock{%
	\ryumeaning%
		{風水。風水学}{}{}%
	}%
	}%
	{\tailPrint[]}%
\entry
	{\headPrint[kana={ふんしきるん\tl{˥˩}},ipa={funʃi̥kiruɴ},pos={動},rui={3},para={ふんすくぬ},acc={F}]}%
	{%
		\MeaningsBlock{%
	\ryumeaning%
		{踏みつける}{}{}%
	}%
	}%
	{\tailPrint[]}%
\entry
	{\headPrint[kana={ふんた\tl{˥˩}},ipa={funta},pos={名},acc={F}]}%
	{%
		\MeaningsBlock{%
	\ryumeaning%
		{床。板床}{}{}%
	}%
	}%
	{\tailPrint[]}%
\entry
	{\headPrint[kana={ふんだやー},ipa={fundajaː},pos={名}]}%
	{%
		\MeaningsBlock{%
	\ryumeaning%
		{我がままだ。勝手放題だ}{}{}%
	}%
	}%
	{\tailPrint[]}%
\entry
	{\headPrint[kana={ぶんちうちん\tl{˩˥}},ipa={buntʃi.utʃiɴ},pos={動},rui={3},para={ぶんちうとぅぬ},acc={R}]}%
	{%
		\MeaningsBlock{%
	\ryumeaning%
		{跳び下りる}{}{}%
	}%
	}%
	{\tailPrint[]}%
\entry
	{\headPrint[kana={ふんつぁーすん\tl{˥˩}},ipa={funtsaːsuɴ},pos={動},rui={2b},para={ふんつぁーさぬ},acc={F}]}%
	{%
		\MeaningsBlock{%
	\ryumeaning%
		{踏みにじる。踏みつける}{しきりに踏む}{}%
	}%
	}%
	{\tailPrint[]}%
\entry
	{\headPrint[kana={ふんつん\tl{˥˩}},ipa={funtsuɴ},pos={動},rui={2},para={ぶんつぁぬ},acc={F}]}%
	{%
		\MeaningsBlock{%
	\ryumeaning%
		{跳ぶ}{命令形は\emb{ふんち}}{}%
	}%
	}%
	{\tailPrint[]}%
\entry
	{\headPrint[kana={ふんでー\tl{˥˩}},ipa={fundeː},pos={名},acc={F}]}%
	{%
		\MeaningsBlock{%
	\ryumeaning%
		{我がまま。あまったれ}{}{}%
	}%
	}%
	{\tailPrint[]}%
\entry
	{\headPrint[kana={ぷんでー\ws{}すむん},ipa={pundeː su̥muɴ},pos={動},rui={2},para={ふくらはぬ}]}%
	{%
		\MeaningsBlock{%
	\ryumeaning%
		{甘やかす}{}{}%
	}%
	}%
	{\tailPrint[]}%
\entry
	{\headPrint[kana={ふんとー\tl{˥}},ipa={funtoː},pos={名},acc={H}]}%
	{%
		\MeaningsBlock{%
	\ryumeaning%
		{本当。真実}{}{}%
	}%
	}%
	{\tailPrint[]}%
\entry
	{\headPrint[kana={ぶんどぅり},ipa={bunduri},pos={名}]}%
	{%
		\MeaningsBlock{%
	\ryumeaning%
		{踊り。舞踊}{}{}%
	}%
	}%
	{\tailPrint[]}%
\entry
	{\headPrint[kana={ぶんどぅるん\tl{˥˩}},ipa={bunduruɴ},pos={動},rui={1},para={ぶんどぅらぬ},acc={F}]}%
	{%
		\MeaningsBlock{%
	\ryumeaning%
		{踊る。舞う}{}{}%
	}%
	}%
	{\tailPrint[]}%
\entry
	{\headPrint[kana={ふんぱんちるん\tl{˥˩}},ipa={fumpantʃiruɴ},pos={動},rui={1},acc={F}]}%
	{%
		\MeaningsBlock{%
	\ryumeaning%
		{踏み外す}{}{}%
	}%
	}%
	{\tailPrint[]}%
\entry
	{\headPrint[kana={ふんぱんつぁすん\tl{˥˩}},ipa={fumpantsasuɴ},pos={動},rui={2},para={ふんぱんつぁはぬ},acc={F}]}%
	{%
		\MeaningsBlock{%
	\ryumeaning%
		{踏み外す}{}{}%
	}%
	}%
	{\tailPrint[]}%
\entry
	{\headPrint[kana={ふんびらがすん\tl{˥˩}},ipa={fumbiragasuɴ},pos={動},rui={2},para={ふんびらがはぬ},acc={F}]}%
	{%
		\MeaningsBlock{%
	\ryumeaning%
		{踏み潰す}{}{}%
	}%
	}%
	{\tailPrint[]}%
\entry
	{\headPrint[kana={ふんまるぐん\tl{˥˩}},ipa={fummaruguɴ},pos={動},rui={1},para={ふんまるがぬ},acc={F}]}%
	{%
		\MeaningsBlock{%
	\ryumeaning%
		{強く束ねる}{}{}%
	}%
	}%
	{\tailPrint[]}%
\LetterHead{へ}
\entry
	{\headPrint[kana={べー},ipa={beː},pos={感}]}%
	{%
		\MeaningsBlock{%
	\ryumeaning%
		{わぁー。へぇー}{}{}%
	}%
	}%
	{\tailPrint[]}%
\entry
	{\headPrint[kana={べー\tl{˩˥}},ipa={beː},pos={名},acc={R}]}%
	{%
		\MeaningsBlock{%
	\ryumeaning%
		{我。我が}{\emb{べーま}「我ら」、\emb{べーひー}「我が家」、\emb{べーすま}「我が島」など}{}%
	}%
	}%
	{\tailPrint[]}%
\entry
	{\headPrint[kana={べー\tl{˩˥}},ipa={beː},pos={名},acc={R}]}%
	{%
		\MeaningsBlock{%
	\ryumeaning%
		{芽}{}{}%
	}%
	}%
	{\tailPrint[]}%
\entry
	{\headPrint[kana={べー\tl{˩˥}},ipa={beː},pos={副},acc={R}]}%
	{%
		\MeaningsBlock{%
	\ryumeaning%
		{少し。わずか}{}{}%
	}%
	}%
	{\tailPrint[]}%
\entry
	{\headPrint[kana={ぺー\tl{˥}},ipa={peː},pos={名},acc={H}]}%
	{%
		\MeaningsBlock{%
	\ryumeaning%
		{南。南方}{}{}%
	}%
	}%
	{\tailPrint[]}%
\entry
	{\headPrint[kana={ぺー\tl{˥˩}},ipa={peː},pos={名},acc={F}]}%
	{%
		\MeaningsBlock{%
	\ryumeaning%
		{\ruby[g]{鍬}{くわ}}{}{}%
	}%
	}%
	{\tailPrint[]}%
\entry
	{\headPrint[kana={ぺー\tl{˥˩}},ipa={peː},pos={名},acc={F}]}%
	{%
		\MeaningsBlock{%
	\ryumeaning%
		{\ruby[g]{蠅}{はえ}}{}{}%
	}%
	}%
	{\tailPrint[]}%
\entry
	{\headPrint[kana={へー\tl{˥˩}\ws{}あますん\tl{˥}},ipa={heː amasuɴ},pos={句},acc={F H}]}%
	{%
		\MeaningsBlock{%
	\ryumeaning%
		{食い残す。食い余す}{}{}%
	}%
	}%
	{\tailPrint[]}%
\entry
	{\headPrint[kana={ぺーかち\tl{˥}},ipa={peːkḁtʃi},pos={名},acc={H}]}%
	{%
		\MeaningsBlock{%
	\ryumeaning%
		{南風}{}{}%
	}%
	}%
	{\tailPrint[]}%
\entry
	{\headPrint[kana={へーじぶん\tl{˥˩}},ipa={heːdʒibuɴ},pos={名},acc={F}]}%
	{%
		\MeaningsBlock{%
	\ryumeaning%
		{食べ頃}{}{}%
	}%
	}%
	{\tailPrint[]}%
\entry
	{\headPrint[kana={ぺーしゃな\tl{˥}},ipa={peːʃana},pos={副},acc={H}]}%
	{%
		\MeaningsBlock{%
	\ryumeaning%
		{早く。急いで}{}{}%
	}%
	}%
	{\tailPrint[]}%
\entry
	{\headPrint[kana={ぺーしゅら\tl{˥}},ipa={peːʃura},pos={名},acc={H}]}%
	{%
		\MeaningsBlock{%
	\ryumeaning%
		{夕立ち}{}{}%
	}%
	}%
	{\tailPrint[]}%
\entry
	{\headPrint[kana={べーすま},ipa={beːsu̥ma},pos={名}]}%
	{%
		\MeaningsBlock{%
	\ryumeaning%
		{我が島。（島民間で）波照間島}{島民間では波照間島を指す。島外の人を除いた「我々の島」という表現なので、波照間島の島民のみが使う名称である}{}%
	}%
	}%
	{\tailPrint[]}%
\entry
	{\headPrint[kana={ぺーすん\tl{˥}},ipa={peːsuɴ},pos={動},rui={2},para={ぺーはぬ},acc={H}]}%
	{%
		\MeaningsBlock{%
	\ryumeaning%
		{（釣り糸などを）這わす}{}{}%
	}%
	}%
	{\tailPrint[]}%
\entry
	{\headPrint[kana={べーちゃー},ipa={beːtʃaː},pos={感}]}%
	{%
		\MeaningsBlock{%
	\ryumeaning%
		{羨ましい}{感嘆詞}{}%
	}%
	}%
	{\tailPrint[]}%
\entry
	{\headPrint[kana={ぺーつぁーるん\tl{˥}},ipa={peːtsaːruɴ},pos={動},rui={1},para={ぺーつぁーらぬ},acc={H}]}%
	{%
		\MeaningsBlock{%
	\ryumeaning%
		{這う。這いつくばう}{}{}%
	}%
	}%
	{\tailPrint[]}%
\entry
	{\headPrint[kana={べ\josho{}ー\kako{}な},ipa={beːna},pos={副},acc={LHL}]}%
	{%
		\MeaningsBlock{%
	\ryumeaning%
		{少し。わずか}{}{}%
	}%
	}%
	{\tailPrint[]}%
\entry
	{\headPrint[kana={\josho{}ぺーぬ\kako{}\ws{}かた},ipa={peːnu kḁta},pos={句},acc={H L}]}%
	{%
		\MeaningsBlock{%
	\ryumeaning%
		{南方}{}{}%
	}%
	}%
	{\tailPrint[]}%
\entry
	{\headPrint[kana={へーぱんちるん\tl{˥˩}},ipa={heːpantʃiruɴ},pos={動},rui={3},para={へーぱんつぬ},acc={F}]}%
	{%
		\MeaningsBlock{%
	\ryumeaning%
		{食い外す}{}{}%
	}%
	}%
	{\tailPrint[]}%
\entry
	{\headPrint[kana={\josho{}べー\kako{}び},ipa={beːbi},pos={副},acc={HL}]}%
	{%
		\MeaningsBlock{%
	\ryumeaning%
		{少し。わずか}{}{}%
	}%
	}%
	{\tailPrint[]}%
\entry
	{\headPrint[kana={ぺーまるん\tl{˥}},ipa={peːmaruɴ},pos={動},rui={1},para={ぺーまらぬ},acc={H}]}%
	{%
		\MeaningsBlock{%
	\ryumeaning%
		{早まる}{}{}%
	}%
	}%
	{\tailPrint[]}%
\entry
	{\headPrint[kana={ぺーむら\tl{˥˩}},ipa={peːmura},pos={名},acc={F}]}%
	{%
		\MeaningsBlock{%
	\ryumeaning%
		{南村}{波照間の村落名}{}%
	}%
	}%
	{\tailPrint[]}%
\entry
	{\headPrint[kana={ぺーり\tl{˥˩}},ipa={peːri},pos={名},acc={F}]}%
	{%
		\MeaningsBlock{%
	\ryumeaning%
		{日照り。\ruby[g]{旱}{かん}\ruby[g]{魃}{ばつ}}{}{}%
	}%
	}%
	{\tailPrint[]}%
\entry
	{\headPrint[kana={ぺーりくむん\tl{˥}},ipa={peːriku̥mun},pos={動},rui={1},para={ぺーりくまぬ},acc={H}]}%
	{%
		\MeaningsBlock{%
	\ryumeaning%
		{入り込む}{}{}%
	}%
	}%
	{\tailPrint[]}%
\entry
	{\headPrint[kana={ぺーり\ws{}しゃーん},ipa={peːri ʃaːɴ},pos={句},acc={F HL}]}%
	{%
		\MeaningsBlock{%
	\ryumeaning%
		{\ruby[g]{旱}{かん}\ruby[g]{魃}{ばつ}だ}{}{}%
	}%
	}%
	{\tailPrint[]}%
\entry
	{\headPrint[kana={ぺーりまぎ\tl{˥˩}},ipa={peːrimagi},pos={名},acc={F}]}%
	{%
		\MeaningsBlock{%
	\ryumeaning%
		{旱害}{}{}%
	}%
	}%
	{\tailPrint[]}%
\entry
	{\headPrint[kana={ぺーる\tl{˩˥}},ipa={peːru},pos={名},acc={R}]}%
	{%
		\MeaningsBlock{%
	\ryumeaning%
		{酢。食酢}{}{}%
	}%
	}%
	{\tailPrint[]}%
\entry
	{\headPrint[kana={ぺーるん\tl{˥}},ipa={peːruɴ},pos={動},rui={1},para={ぺーらぬ},acc={H}]}%
	{%
		\MeaningsBlock{%
	\ryumeaning%
		{入る}{}{}%
	}%
	}%
	{\tailPrint[]}%
\entry
	{\headPrint[kana={ぺしきびり\tl{˥˩}},ipa={peʃi̥kibiri},pos={名},acc={F}]}%
	{%
		\MeaningsBlock{%
	\ryumeaning%
		{ひざまずき。正座}{}{}%
	}%
	}%
	{\tailPrint[]}%
\entry
	{\headPrint[kana={ぺしくらやむん\tl{˩˥}},ipa={peʃi̥kurajamuɴ},pos={動},rui={1},para={ぺしくらやまぬ},acc={R}]}%
	{%
		\MeaningsBlock{%
	\ryumeaning%
		{\ruby[g]{痺}{しび}れる}{}{}%
	}%
	}%
	{\tailPrint[]}%
\entry
	{\headPrint[kana={へだま\tl{˥˩}},ipa={hedama},pos={名},acc={F}]}%
	{%
		\MeaningsBlock{%
	\ryumeaning%
		{食いしん坊}{}{}%
	}%
	}%
	{\tailPrint[]}%
\entry
	{\headPrint[kana={へだま\tl{˥˩}さーん},ipa={hedamasaːɴ},pos={形},acc={F}]}%
	{%
		\MeaningsBlock{%
	\ryumeaning%
		{食い意地が汚い}{}{}%
	}%
	}%
	{\tailPrint[]}%
\entry
	{\headPrint[kana={ぺっかるん\tl{˥}},ipa={pekkaruɴ},pos={動},rui={1},para={ぺっからぬ},acc={H}]}%
	{%
		\MeaningsBlock{%
	\ryumeaning%
		{這う}{}{}%
	}%
	}%
	{\tailPrint[]}%
\entry
	{\headPrint[kana={\josho{}ぺったぬ\kako{}\ws{}かた},ipa={pettanu kḁta},pos={句},acc={H L}]}%
	{%
		\MeaningsBlock{%
	\ryumeaning%
		{南方}{}{}%
	}%
	}%
	{\tailPrint[]}%
\entry
	{\headPrint[kana={ぺっつぁるん\tl{˥}},ipa={pettsaruɴ},pos={動},rui={1},para={ぺっつぁらぬ},acc={H}]}%
	{%
		\MeaningsBlock{%
	\ryumeaning%
		{這いつくばう}{}{}%
	}%
	}%
	{\tailPrint[]}%
\entry
	{\headPrint[kana={へぶたん},ipa={hebutaɴ},pos={句},para={へぶとぅぬ}]}%
	{%
		\MeaningsBlock{%
	\ryumeaning%
		{食べ尽くす}{}{}%
	}%
	}%
	{\tailPrint[]}%
\entry
	{\headPrint[kana={ぺふつぁ\tl{˥}},ipa={pefu̥tsa},pos={名},acc={H}]}%
	{%
		\MeaningsBlock{%
	\ryumeaning%
		{\ruby[g]{鳶}{とび}}{}{}%
	}%
	}%
	{\tailPrint[]}%
\entry
	{\headPrint[kana={ぺりふちぃ\tl{˥}},ipa={perifu̥tʃi},pos={名},acc={H}]}%
	{%
		\MeaningsBlock{%
	\ryumeaning%
		{入口}{}{}%
	}%
	}%
	{\tailPrint[]}%
\entry
	{\headPrint[kana={ぺるん\tl{˥}},ipa={peruɴ},pos={動},rui={1},para={ぺらぬ},acc={H}]}%
	{%
		\MeaningsBlock{%
	\ryumeaning%
		{入る}{}{}%
	}%
	}%
	{\tailPrint[]}%
\LetterHead{ほ}
\entry
	{\headPrint[kana={ぼー\tl{˩˥}},ipa={boː},pos={名},acc={R}]}%
	{%
		\MeaningsBlock{%
	\ryumeaning%
		{棒。棒術}{}{}%
	}%
	}%
	{\tailPrint[]}%
\entry
	{\headPrint[kana={ぼー\tl{˩˥}},ipa={boː},pos={名},acc={R}]}%
	{%
		\MeaningsBlock{%
	\ryumeaning%
		{共同作業}{}{}%
	}%
	}%
	{\tailPrint[note={意味的に八重山の他の方言の\emb{ゆい}に対応する}]}%
\entry
	{\headPrint[kana={ぽー\tl{˥˩}},ipa={poː},pos={名},acc={F}]}%
	{%
		\MeaningsBlock{%
	\ryumeaning%
		{帆}{}{}%
	}%
	}%
	{\tailPrint[]}%
\entry
	{\headPrint[kana={ぼーがり\tl{˩˥}},ipa={boːgari},pos={名},acc={R}]}%
	{%
		\MeaningsBlock{%
	\ryumeaning%
		{疲れ。疲労}{}{}%
	}%
	}%
	{\tailPrint[]}%
\entry
	{\headPrint[kana={ぽーぎぶしぃ\tl{˥}},ipa={poːgibuʃi},pos={名},acc={H}]}%
	{%
		\MeaningsBlock{%
	\ryumeaning%
		{ほうき星。彗星}{}{}%
	}%
	}%
	{\tailPrint[]}%
\entry
	{\headPrint[kana={ぽーぐん\tl{˥}},ipa={poːguɴ},pos={動},rui={1},para={ぽーがぬ},acc={H}]}%
	{%
		\MeaningsBlock{%
	\ryumeaning%
		{掃く。掃除する}{}{}%
	}%
	}%
	{\tailPrint[]}%
\entry
	{\headPrint[kana={ぽーしん\tl{˥˩}},ipa={poːʃiɴ},pos={名},acc={F}]}%
	{%
		\MeaningsBlock{%
	\ryumeaning%
		{帆船}{}{}%
	}%
	}%
	{\tailPrint[]}%
\entry
	{\headPrint[kana={ぼーす\tl{˩˥}},ipa={boːsu},pos={名},acc={R}]}%
	{%
		\MeaningsBlock{%
	\ryumeaning%
		{芒種}{季節名。沖縄は梅雨時期に当たる}{}%
	}%
	}%
	{\tailPrint[]}%
\entry
	{\headPrint[kana={ぼーず\tl{˩˥}},ipa={boːdzu},pos={名},acc={R}]}%
	{%
		\MeaningsBlock{%
	\ryumeaning%
		{坊主。僧侶}{}{}%
	}%
	}%
	{\tailPrint[]}%
\entry
	{\headPrint[kana={ぽーち\tl{˥}},ipa={poːtʃi},pos={名},acc={H}]}%
	{%
		\MeaningsBlock{%
	\ryumeaning%
		{\ruby[g]{箒}{ほうき}}{}{}%
	}%
	}%
	{\tailPrint[]}%
\entry
	{\headPrint[kana={ぼーま\tl{˥}},ipa={boːma},pos={名},acc={H}]}%
	{%
		\MeaningsBlock{%
	\ryumeaning%
		{長姉}{}{}%
	}%
	}%
	{\tailPrint[]}%
\entry
	{\headPrint[kana={ぼーま\tl{˥}},ipa={boːma},pos={名},acc={H}]}%
	{%
		\MeaningsBlock{%
	\ryumeaning%
		{\ruby[g]{大浜}{おおはま}}{石垣島の集落名}{}%
	}%
	}%
	{\tailPrint[]}%
\entry
	{\headPrint[kana={ほーむぬ\tl{˥˩}},ipa={hoːmunu},pos={名},acc={F}]}%
	{%
		\MeaningsBlock{%
	\ryumeaning%
		{食べ物。食糧}{}{}%
	}%
	}%
	{\tailPrint[]}%
\entry
	{\headPrint[kana={ほーらきし\tl{˥}},ipa={hoːraki̥ʃi},pos={名},acc={H}]}%
	{%
		\MeaningsBlock{%
	\ryumeaning%
		{あほらしい。間抜け}{}{}%
	}%
	}%
	{\tailPrint[]}%
\entry
	{\headPrint[kana={ほーらし},ipa={hoːraʃi},pos={-}]}%
	{%
		\MeaningsBlock{%
	\ryumeaning%
		{威張る。自惚れる}{}{}%
	}%
	}%
	{\tailPrint[]}%
\entry
	{\headPrint[kana={ぽーりすくん},ipa={poːrisu̥kuɴ},pos={動},rui={3},para={ぽーりしぃくな}]}%
	{%
		\MeaningsBlock{%
	\ryumeaning%
		{蒔き散らかす}{}{}%
	}%
	}%
	{\tailPrint[]}%
\entry
	{\headPrint[kana={ぼーり\tl{˩˥}のがすん\tl{˩˥}},ipa={boːrinogasuɴ},pos={句},para={ぼーりのがはぬ},acc={R R}]}%
	{%
		\MeaningsBlock{%
	\ryumeaning%
		{疲れをとる}{}{}%
	}%
	}%
	{\tailPrint[]}%
\entry
	{\headPrint[kana={ぼーりぼーり},ipa={boːriboːri},pos={感}]}%
	{%
		\MeaningsBlock{%
	\ryumeaning%
		{お利口さん}{}{}%
	}%
	}%
	{\tailPrint[]}%
\entry
	{\headPrint[kana={ぽーりゃん\tl{˥˩}},ipa={poːrjaɴ},pos={動（継）},rui={1},para={ぽーらぬ},acc={F}]}%
	{%
		\MeaningsBlock{%
	\ryumeaning%
		{散らかる。散らばる}{}{}%
	}%
	}%
	{\tailPrint[]}%
\entry
	{\headPrint[kana={ぽーるん\tl{˥˩}},ipa={poːruɴ},pos={動},rui={1},para={ぽーらぬ},acc={F}]}%
	{%
		\MeaningsBlock{%
	\ryumeaning%
		{散らす}{}{}%
	}%
	}%
	{\tailPrint[]}%
\entry
	{\headPrint[kana={ほーん\tl{˥˩}},ipa={hoːɴ},pos={動},rui={1},para={はーぬ},acc={F}]}%
	{%
		\MeaningsBlock{%
	\ryumeaning%
		{食べる}{}{}%
	}%
	}%
	{\tailPrint[]}%
\entry
	{\headPrint[kana={ほすん\tl{˥}},ipa={hosuɴ},pos={動},rui={2},para={ほはぬ},acc={H}]}%
	{%
		\MeaningsBlock{%
	\ryumeaning%
		{釣る}{}{}%
	}%
	}%
	{\tailPrint[]}%
\entry
	{\headPrint[kana={ぼたりるん\tl{˩˥}},ipa={botariruɴ},pos={動},rui={3},para={ぼたるぬ},acc={R}]}%
	{%
		\MeaningsBlock{%
	\ryumeaning%
		{疲れる。疲労する}{}{}%
	}%
	}%
	{\tailPrint[]}%
\entry
	{\headPrint[kana={ぼたりん\tl{˩˥}},ipa={botariɴ},pos={動},rui={3},para={ぼたるぬ},acc={R}]}%
	{%
		\MeaningsBlock{%
	\ryumeaning%
		{疲れる}{}{}%
	}%
	}%
	{\tailPrint[]}%
\entry
	{\headPrint[kana={ほっか\tl{˥}},ipa={hokka},pos={名},acc={H}]}%
	{%
		\MeaningsBlock{%
	\ryumeaning%
		{手品}{}{}%
	}%
	}%
	{\tailPrint[]}%
\entry
	{\headPrint[kana={ぽっつぁ\tl{˥}},ipa={pottsa},pos={名},acc={H}]}%
	{%
		\MeaningsBlock{%
	\ryumeaning%
		{包丁}{}{}%
	}%
	}%
	{\tailPrint[]}%
\entry
	{\headPrint[kana={ぽっつぁーるん\tl{˥}},ipa={pottsaːruɴ},pos={動},acc={H}]}%
	{%
		\MeaningsBlock{%
	\ryumeaning%
		{ばらばらになる}{}{}%
	}%
	}%
	{\tailPrint[]}%
\entry
	{\headPrint[kana={ぽっつぁらすん},ipa={pottsarasuɴ},pos={動}]}%
	{%
		\MeaningsBlock{%
	\ryumeaning%
		{ばらす。ばらばらにする}{}{}%
	}%
	}%
	{\tailPrint[]}%
\entry
	{\headPrint[kana={ぼふぁった},ipa={bofatta},pos={擬}]}%
	{%
		\MeaningsBlock{%
	\ryumeaning%
		{さっと貫通するさま}{}{}%
	}%
	}%
	{\tailPrint[]}%
\entry
	{\headPrint[kana={ぼふた},ipa={bofu̥ta},pos={名}]}%
	{%
		\MeaningsBlock{%
	\ryumeaning%
		{全部}{}{}%
	}%
	}%
	{\tailPrint[]}%
\entry
	{\headPrint[kana={ぼふった},ipa={bofutta},pos={擬}]}%
	{%
		\MeaningsBlock{%
	\ryumeaning%
		{さっと貫通するさま}{}{}%
	}%
	}%
	{\tailPrint[]}%
\entry
	{\headPrint[kana={ぼへぬ},ipa={bohenu},pos={句},para={ぼはん}]}%
	{%
		\MeaningsBlock{%
	\ryumeaning%
		{～たくない}{\emb{へぼへぬ}「食べたくない」など}{}%
	}%
	}%
	{\tailPrint[]}%
\entry
	{\headPrint[kana={ぼんぎるん\tl{˩˥}},ipa={boŋgiruɴ},pos={動},rui={3},para={ぼんぐぬ},acc={R}]}%
	{%
		\MeaningsBlock{%
	\ryumeaning%
		{放り投げる。投げ捨てる}{}{}%
	}%
	}%
	{\tailPrint[]}%
\LetterHead{ま}
\entry
	{\headPrint[kana={まー},ipa={maː},pos={接頭}]}%
	{%
		\MeaningsBlock{%
	\ryumeaning%
		{真～。本～}{}{}%
	}%
	}%
	{\tailPrint[note={\emb{まー}が付く語は下降型に所属する}]}%
\entry
	{\headPrint[kana={まー\tl{˥}},ipa={maː},pos={名},acc={H}]}%
	{%
		\MeaningsBlock{%
	\ryumeaning%
		{孫}{\emb{またまー}は「曽孫(ひまご)」}{}%
	}%
	}%
	{\tailPrint[]}%
\entry
	{\headPrint[kana={まー\tl{˥˩}},ipa={maː},pos={名},acc={F}]}%
	{%
		\MeaningsBlock{%
	\ryumeaning%
		{間}{場所を表す語}{}%
	}%
	}%
	{\tailPrint[]}%
\entry
	{\headPrint[kana={まーおーび},ipa={maː.oːbi},pos={名}]}%
	{%
		\MeaningsBlock{%
	\ryumeaning%
		{倍。二倍}{}{}%
	}%
	}%
	{\tailPrint[]}%
\entry
	{\headPrint[kana={まーぎ},ipa={maːgi},pos={形}]}%
	{%
		\MeaningsBlock{%
	\ryumeaning%
		{大きい}{}{}%
	}%
	}%
	{\tailPrint[]}%
\entry
	{\headPrint[kana={まーぎ\tl{˩˥}くい\tl{˥}},ipa={maːgikui},pos={名},acc={R+H}]}%
	{%
		\MeaningsBlock{%
	\ryumeaning%
		{大声}{}{}%
	}%
	}%
	{\tailPrint[]}%
\entry
	{\headPrint[kana={まーぐん\tl{˥˩}},ipa={maːguɴ},pos={動},rui={1},para={まーがぬ},acc={F}]}%
	{%
		\MeaningsBlock{%
	\ryumeaning%
		{播く。蒔く}{}{}%
	}%
	}%
	{\tailPrint[]}%
\entry
	{\headPrint[kana={まーぐん\tl{˥˩}},ipa={maːguɴ},pos={動},rui={1},para={まーがぬ},acc={F}]}%
	{%
		\MeaningsBlock{%
	\ryumeaning%
		{巻く}{}{}%
	}%
	}%
	{\tailPrint[]}%
\entry
	{\headPrint[kana={まーさかさー\ws{}ねーぬ},ipa={maːsḁkasaː neːnu},pos={形}]}%
	{%
		\MeaningsBlock{%
	\ryumeaning%
		{体調が悪い}{}{}%
	}%
	}%
	{\tailPrint[]}%
\entry
	{\headPrint[kana={まーし\tl{˥}},ipa={maːʃi},pos={名},acc={H}]}%
	{%
		\MeaningsBlock{%
	\ryumeaning%
		{箸}{}{}%
	}%
	}%
	{\tailPrint[note={\emb{おみぱし}「御御箸」の転。アクセントより本来母音始まりであることが分かる}]}%
\entry
	{\headPrint[kana={まーし\tl{˩˥}},ipa={maːʃi},pos={名},acc={R}]}%
	{%
		\MeaningsBlock{%
	\ryumeaning%
		{（～より）まし。良い}{}{}%
	}%
	}%
	{\tailPrint[]}%
\entry
	{\headPrint[kana={まーしぃ\tl{˥˩}},ipa={maːʃi},pos={名},acc={F}]}%
	{%
		\MeaningsBlock{%
	\ryumeaning%
		{（田の）区画}{}{}%
	}%
	}%
	{\tailPrint[]}%
\entry
	{\headPrint[kana={まーじぃ\tl{˥˩}},ipa={maːdʒi},pos={名},acc={F}]}%
	{%
		\MeaningsBlock{%
	\ryumeaning%
		{一緒。同期生}{}{}%
	}%
	}%
	{\tailPrint[]}%
\entry
	{\headPrint[kana={まーしとぅ\tl{˥˩}},ipa={maːʃitu},pos={名},acc={F}]}%
	{%
		\MeaningsBlock{%
	\ryumeaning%
		{開き戸}{}{}%
	}%
	}%
	{\tailPrint[]}%
\entry
	{\headPrint[kana={まーじゃかち\tl{˥˩}},ipa={maːdʒakḁtʃi},pos={名},acc={F}]}%
	{%
		\MeaningsBlock{%
	\ryumeaning%
		{つむじ風}{}{}%
	}%
	}%
	{\tailPrint[]}%
\entry
	{\headPrint[kana={まーす\tl{˩˥}},ipa={maːsu},pos={名},acc={R}]}%
	{%
		\MeaningsBlock{%
	\ryumeaning%
		{塩}{}{}%
	}%
	}%
	{\tailPrint[]}%
\entry
	{\headPrint[kana={まーすにー\tl{˩˥}},ipa={maːsuniː},pos={名},acc={R}]}%
	{%
		\MeaningsBlock{%
	\ryumeaning%
		{塩煮}{}{}%
	}%
	}%
	{\tailPrint[]}%
\entry
	{\headPrint[kana={まーすむ\tl{˥˩}},ipa={maːsu̥mu},pos={名},acc={F}]}%
	{%
		\MeaningsBlock{%
	\ryumeaning%
		{本心。本気}{}{}%
	}%
	}%
	{\tailPrint[]}%
\entry
	{\headPrint[kana={まーすん\tl{˥˩}},ipa={maːsuɴ},pos={動},rui={2},para={まーはぬ},acc={F}]}%
	{%
		\MeaningsBlock{%
	\ryumeaning%
		{回す。回転させる}{}{}%
	}%
	}%
	{\tailPrint[]}%
\entry
	{\headPrint[kana={まーたき},ipa={maːtḁki},pos={名}]}%
	{%
		\MeaningsBlock{%
	\ryumeaning%
		{同等。対等。公平}{}{}%
	}%
	}%
	{\tailPrint[]}%
\entry
	{\headPrint[kana={まーち\tl{˩˥}},ipa={maːtʃi},pos={名},acc={R}]}%
	{%
		\MeaningsBlock{%
	\ryumeaning%
		{松}{}{}%
	}%
	}%
	{\tailPrint[]}%
\entry
	{\headPrint[kana={まーに\tl{˩˥}},ipa={maːni},pos={名},acc={R}]}%
	{%
		\MeaningsBlock{%
	\ryumeaning%
		{クロツグ}{植物名}{}%
	}%
	}%
	{\tailPrint[]}%
\entry
	{\headPrint[kana={まーぱからさー\ws{}ねーぬ},ipa={maːpakarasaː neːnu},pos={句}]}%
	{%
		\MeaningsBlock{%
	\ryumeaning%
		{全く利口でない。まるで愚かだ}{}{}%
	}%
	}%
	{\tailPrint[]}%
\entry
	{\headPrint[kana={まー\tl{˥}はん},ipa={maːhaɴ},pos={形},para={まーへぬ},acc={H}]}%
	{%
		\MeaningsBlock{%
	\ryumeaning%
		{旨い。美味しい}{}{}%
	}%
	}%
	{\tailPrint[]}%
\entry
	{\headPrint[kana={まーび},ipa={maːbi},pos={副}]}%
	{%
		\MeaningsBlock{%
	\ryumeaning%
		{もっと}{}{}%
	}%
	}%
	{\tailPrint[]}%
\entry
	{\headPrint[kana={まー\ws{}ぴぃとぅむし},ipa={maː pi̥tumuʃi},pos={句}]}%
	{%
		\MeaningsBlock{%
	\ryumeaning%
		{もう一度}{}{}%
	}%
	}%
	{\tailPrint[]}%
\entry
	{\headPrint[kana={まーびる\tl{˥˩}},ipa={maːbiru},pos={名},acc={F}]}%
	{%
		\MeaningsBlock{%
	\ryumeaning%
		{真昼。日中}{}{}%
	}%
	}%
	{\tailPrint[]}%
\entry
	{\headPrint[kana={まーぶり\tl{˩˥}},ipa={maːburi},pos={名},acc={R}]}%
	{%
		\MeaningsBlock{%
	\ryumeaning%
		{魂。霊魂}{}{}%
	}%
	}%
	{\tailPrint[]}%
\entry
	{\headPrint[kana={まーぷり\tl{˥˩}},ipa={maːpu̥ri},pos={名},acc={F}]}%
	{%
		\MeaningsBlock{%
	\ryumeaning%
		{べた惚れ。熱中}{}{}%
	}%
	}%
	{\tailPrint[]}%
\entry
	{\headPrint[kana={まーぷりむん\tl{˥˩}},ipa={maːpu̥rimuɴ},pos={名},acc={F}]}%
	{%
		\MeaningsBlock{%
	\ryumeaning%
		{大馬鹿者}{}{}%
	}%
	}%
	{\tailPrint[]}%
\entry
	{\headPrint[kana={まーべー\tl{˥˩}},ipa={maːbeː},pos={名},acc={F}]}%
	{%
		\MeaningsBlock{%
	\ryumeaning%
		{真似}{}{}%
	}%
	}%
	{\tailPrint[note={真似すること}]}%
\entry
	{\headPrint[kana={まーべー\tl{˥˩}\ws{}すん\tl{˥˩}},ipa={maːbeː suɴ},pos={句},acc={F F}]}%
	{%
		\MeaningsBlock{%
	\ryumeaning%
		{真似する。真似る}{}{}%
	}%
	}%
	{\tailPrint[]}%
\entry
	{\headPrint[kana={まーましぃ\tl{˥˩}},ipa={maːmaʃi},pos={名},acc={F}]}%
	{%
		\MeaningsBlock{%
	\ryumeaning%
		{まわり。周囲}{}{}%
	}%
	}%
	{\tailPrint[]}%
\entry
	{\headPrint[kana={まーみ\tl{˩˥}},ipa={maːmi},pos={名},acc={R}]}%
	{%
		\MeaningsBlock{%
	\ryumeaning%
		{心臓}{魚などの心臓}{}%
	}%
	}%
	{\tailPrint[]}%
\entry
	{\headPrint[kana={まーみ\tl{˩˥}},ipa={maːmi},pos={名},acc={R}]}%
	{%
		\MeaningsBlock{%
	\ryumeaning%
		{豆}{}{}%
	}%
	}%
	{\tailPrint[]}%
\entry
	{\headPrint[kana={まー\ws{}みちん},ipa={maː mitʃiɴ},pos={句}]}%
	{%
		\MeaningsBlock{%
	\ryumeaning%
		{再来年}{}{}%
	}%
	}%
	{\tailPrint[]}%
\entry
	{\headPrint[kana={まーみなん\tl{˩˥}},ipa={maːminaɴ},pos={名},acc={R}]}%
	{%
		\MeaningsBlock{%
	\ryumeaning%
		{モヤシ}{「豆菜」から}{}%
	}%
	}%
	{\tailPrint[]}%
\entry
	{\headPrint[kana={まーむしぃぴ\tl{˥˩}},ipa={maːmuʃi̥pi},pos={名},acc={F}]}%
	{%
		\MeaningsBlock{%
	\ryumeaning%
		{真結び。本結び}{}{}%
	}%
	}%
	{\tailPrint[]}%
\entry
	{\headPrint[kana={まーむぬ\tl{˥˩}},ipa={maːmunu},pos={名},acc={F}]}%
	{%
		\MeaningsBlock{%
	\ryumeaning%
		{本物}{\emb{ハブ}「毒蛇」にも言う}{}%
	}%
	}%
	{\tailPrint[]}%
\entry
	{\headPrint[kana={まーむん},ipa={maːmun},pos={名}]}%
	{%
		\MeaningsBlock{%
	\ryumeaning%
		{小麦}{}{}%
	}%
	}%
	{\tailPrint[]}%
\entry
	{\headPrint[kana={まー\josho{}らん\kako{}ぶに},ipa={maːrambuni},pos={名},acc={特殊}]}%
	{%
		\MeaningsBlock{%
	\ryumeaning%
		{\ruby[g]{馬}{まー}\ruby[g]{艦}{らん}\ruby[g]{船}{せん}}{近世中期以降の沖縄の帆船}{}%
	}%
	}%
	{\tailPrint[]}%
\entry
	{\headPrint[kana={まーり\tl{˥˩}},ipa={maːri},pos={名},acc={F}]}%
	{%
		\MeaningsBlock{%
	\ryumeaning%
		{順番}{}{}%
	}%
	}%
	{\tailPrint[]}%
\entry
	{\headPrint[kana={まーり\tl{˩˥}},ipa={maːri},pos={名},acc={R}]}%
	{%
		\MeaningsBlock{%
	\ryumeaning%
		{碗}{}{}%
	}%
	}%
	{\tailPrint[]}%
\entry
	{\headPrint[kana={まーり\tl{˩˥}},ipa={maːri},pos={名},acc={R}]}%
	{%
		\MeaningsBlock{%
	\ryumeaning%
		{周り。周囲}{}{}%
	}%
	}%
	{\tailPrint[]}%
\entry
	{\headPrint[kana={まーり\tl{˥˩}\ws{}みるん\tl{˩˥}},ipa={maːri miruɴ},pos={句},acc={F R}]}%
	{%
		\MeaningsBlock{%
	\ryumeaning%
		{見回る。巡回する}{}{}%
	}%
	}%
	{\tailPrint[]}%
\entry
	{\headPrint[kana={まーるん\tl{˥˩}},ipa={maːruɴ},pos={動},rui={1},para={まーらぬ},acc={F}]}%
	{%
		\MeaningsBlock{%
	\ryumeaning%
		{回る。回転する}{}{}%
	}%
	}%
	{\tailPrint[]}%
\entry
	{\headPrint[kana={まかすん\tl{˥˩}},ipa={makasuɴ},pos={動},rui={2},para={まかはぬ},acc={F}]}%
	{%
		\MeaningsBlock{%
	\ryumeaning%
		{負かす}{}{}%
	}%
	}%
	{\tailPrint[]}%
\entry
	{\headPrint[kana={まかすん\tl{˥˩}},ipa={makasuɴ},pos={動},rui={2},para={まかはぬ},acc={F}]}%
	{%
		\MeaningsBlock{%
	\ryumeaning%
		{任せる}{}{}%
	}%
	}%
	{\tailPrint[]}%
\entry
	{\headPrint[kana={まがすん\tl{˩˥}},ipa={magasuɴ},pos={動},rui={2},para={まがはぬ},acc={R}]}%
	{%
		\MeaningsBlock{%
	\ryumeaning%
		{炊事する}{}{}%
	}%
	}%
	{\tailPrint[]}%
\entry
	{\headPrint[kana={まぎ\tl{˥˩}},ipa={magi},pos={名},acc={F}]}%
	{%
		\MeaningsBlock{%
	\ryumeaning%
		{かぶれ。まけ}{}{}%
	}%
	}%
	{\tailPrint[]}%
\entry
	{\headPrint[kana={まきしょー},ipa={maki̥ʃoː},pos={感}]}%
	{%
		\MeaningsBlock{%
	\ryumeaning%
		{もういやだ。やるもんか}{}{}%
	}%
	}%
	{\tailPrint[]}%
\entry
	{\headPrint[kana={まぎぽーるん\tl{˩˥}},ipa={magipoːruɴ},pos={動},rui={1},para={まぎぽーらぬ},acc={R}]}%
	{%
		\MeaningsBlock{%
	\ryumeaning%
		{撒き散らす}{}{}%
	}%
	}%
	{\tailPrint[]}%
\entry
	{\headPrint[kana={まぎり\tl{˥˩}},ipa={magiri},pos={名},acc={F}]}%
	{%
		\MeaningsBlock{%
	\ryumeaning%
		{間切}{王府時代の行政区画}{}%
	}%
	}%
	{\tailPrint[]}%
\entry
	{\headPrint[kana={まきるん\tl{˥˩}},ipa={makiruɴ},pos={動},rui={3},para={まくぬ},acc={F}]}%
	{%
		\MeaningsBlock{%
	\ryumeaning%
		{負ける。降参する}{}{}%
	}%
	}%
	{\tailPrint[]}%
\entry
	{\headPrint[kana={まく\tl{˩˥}},ipa={maku},pos={名},acc={R}]}%
	{%
		\MeaningsBlock{%
	\ryumeaning%
		{幕}{}{}%
	}%
	}%
	{\tailPrint[]}%
\entry
	{\headPrint[kana={まぐくる\tl{˥˩}},ipa={magukuru},pos={名},acc={F}]}%
	{%
		\MeaningsBlock{%
	\ryumeaning%
		{真心}{}{}%
	}%
	}%
	{\tailPrint[]}%
\entry
	{\headPrint[kana={まくとぅ\tl{˩˥}},ipa={maku̥tu},pos={名},acc={R}]}%
	{%
		\MeaningsBlock{%
	\ryumeaning%
		{誠}{}{}%
	}%
	}%
	{\tailPrint[]}%
\entry
	{\headPrint[kana={まさむぬ\tl{˥}},ipa={masamunu},pos={名},acc={H}]}%
	{%
		\MeaningsBlock{%
	\ryumeaning%
		{ご馳走}{}{}%
	}%
	}%
	{\tailPrint[]}%
\entry
	{\headPrint[kana={まざむん\tl{˥˩}},ipa={madzamuɴ},pos={名},acc={F}]}%
	{%
		\MeaningsBlock{%
	\ryumeaning%
		{魔物。妖怪。幽霊}{}{}%
	}%
	}%
	{\tailPrint[]}%
\entry
	{\headPrint[kana={まさるん\tl{˥˩}},ipa={masaruɴ},pos={動},rui={1},para={まさらぬ},acc={F}]}%
	{%
		\MeaningsBlock{%
	\ryumeaning%
		{優れる。勝る}{}{}%
	}%
	}%
	{\tailPrint[]}%
\entry
	{\headPrint[kana={ましかく\tl{˥˩}},ipa={maʃi̥kaku},pos={名},acc={F}]}%
	{%
		\MeaningsBlock{%
	\ryumeaning%
		{ミミズク}{フクロウの種類}{}%
	}%
	}%
	{\tailPrint[]}%
\entry
	{\headPrint[kana={ましん\tl{˥˩}},ipa={maʃiɴ},pos={名},acc={F}]}%
	{%
		\MeaningsBlock{%
	\ryumeaning%
		{本当だ}{}{}%
	}%
	}%
	{\tailPrint[]}%
\entry
	{\headPrint[kana={ますぶどぅり\tl{˥˩}},ipa={masu̥buduri},pos={名},acc={F}]}%
	{%
		\MeaningsBlock{%
	\ryumeaning%
		{巻き踊り}{豊年祭の時の集団円舞}{}%
	}%
	}%
	{\tailPrint[]}%
\entry
	{\headPrint[kana={まそーみ\tl{˥˩}},ipa={masoːmi},pos={名},acc={F}]}%
	{%
		\MeaningsBlock{%
	\ryumeaning%
		{真正面}{御嶽の遥拝所の祭壇}{}%
	}%
	}%
	{\tailPrint[]}%
\entry
	{\headPrint[kana={また\tl{˥˩}},ipa={mata},pos={名},acc={F}]}%
	{%
		\MeaningsBlock{%
	\ryumeaning%
		{股。叉}{}{}%
	}%
	}%
	{\tailPrint[]}%
\entry
	{\headPrint[kana={また\tl{˥˩}},ipa={mata},pos={副},acc={F}]}%
	{%
		\MeaningsBlock{%
	\ryumeaning%
		{又。再び}{}{}%
	}%
	}%
	{\tailPrint[]}%
\entry
	{\headPrint[kana={またいちふ\tl{˥˩}},ipa={mata.itʃifu},pos={名},acc={F}]}%
	{%
		\MeaningsBlock{%
	\ryumeaning%
		{\ruby[g]{又}{また}\ruby[g]{従兄弟}{いとこ}}{}{}%
	}%
	}%
	{\tailPrint[]}%
\entry
	{\headPrint[kana={まだぎるん\tl{˩˥}},ipa={madagiruɴ},pos={動},rui={1?},para={まだぎらぬ},acc={R}]}%
	{%
		\MeaningsBlock{%
	\ryumeaning%
		{片付ける。整理する}{}{}%
	}%
	}%
	{\tailPrint[]}%
\entry
	{\headPrint[kana={またぐい\tl{˥˩}},ipa={matagui},pos={名},acc={F}]}%
	{%
		\MeaningsBlock{%
	\ryumeaning%
		{追肥}{}{}%
	}%
	}%
	{\tailPrint[]}%
\entry
	{\headPrint[kana={またさ\tl{˩˥}},ipa={matasa},pos={名},acc={R}]}%
	{%
		\MeaningsBlock{%
	\ryumeaning%
		{燕}{}{}%
	}%
	}%
	{\tailPrint[]}%
\entry
	{\headPrint[kana={またそーり\tl{˥˩}},ipa={matasoːri},pos={名},acc={F}]}%
	{%
		\MeaningsBlock{%
	\ryumeaning%
		{男の再婚}{「また嫁を連れる」の義}{}%
	}%
	}%
	{\tailPrint[]}%
\entry
	{\headPrint[kana={まだ\tl{˥˩}はん},ipa={madahaɴ},pos={形},para={まだへぬ},acc={F}]}%
	{%
		\MeaningsBlock{%
	\ryumeaning%
		{立派だ。獲物が多い}{}{}%
	}%
	}%
	{\tailPrint[]}%
\entry
	{\headPrint[kana={またびし\tl{˩˥}},ipa={matabiʃi},pos={名},acc={R}]}%
	{%
		\MeaningsBlock{%
	\ryumeaning%
		{股下。股間}{}{}%
	}%
	}%
	{\tailPrint[]}%
\entry
	{\headPrint[kana={またべー\tl{˩˥}},ipa={matabeː},pos={名},acc={R}]}%
	{%
		\MeaningsBlock{%
	\ryumeaning%
		{脇芽}{}{}%
	}%
	}%
	{\tailPrint[]}%
\entry
	{\headPrint[kana={またまー\tl{˥˩}},ipa={matamaː},pos={名},acc={F}]}%
	{%
		\MeaningsBlock{%
	\ryumeaning%
		{\ruby[g]{曽}{ひ}孫}{}{}%
	}%
	}%
	{\tailPrint[]}%
\entry
	{\headPrint[kana={またむち\tl{˥˩}},ipa={matamutʃi},pos={名},acc={F}]}%
	{%
		\MeaningsBlock{%
	\ryumeaning%
		{女の再婚}{「また夫を持つ」の義}{}%
	}%
	}%
	{\tailPrint[]}%
\entry
	{\headPrint[kana={またんがるん\tl{˩˥}},ipa={mataŋgaruɴ},pos={動},rui={1},para={またんがらぬ},acc={R}]}%
	{%
		\MeaningsBlock{%
	\ryumeaning%
		{\ruby[g]{跨}{またが}る}{}{}%
	}%
	}%
	{\tailPrint[]}%
\entry
	{\headPrint[kana={まちがいん\tl{˥}},ipa={matʃigaiɴ},pos={動},acc={H}]}%
	{%
		\MeaningsBlock{%
	\ryumeaning%
		{間違える。取り違える}{}{}%
	}%
	}%
	{\tailPrint[]}%
\entry
	{\headPrint[kana={まちくがりん\tl{˩˥}},ipa={matʃikugariɴ},pos={動},rui={3},para={まちくがるぬ},acc={R}]}%
	{%
		\MeaningsBlock{%
	\ryumeaning%
		{待ち焦がれる}{}{}%
	}%
	}%
	{\tailPrint[]}%
\entry
	{\headPrint[kana={まちげー\tl{˩˥}},ipa={matʃigeː},pos={名},acc={R}]}%
	{%
		\MeaningsBlock{%
	\ryumeaning%
		{間違い。誤り}{}{}%
	}%
	}%
	{\tailPrint[]}%
\entry
	{\headPrint[kana={まちげーるん\tl{˥}},ipa={matʃigeːruɴ},pos={動},rui={1},para={まちげーらぬ},acc={H}]}%
	{%
		\MeaningsBlock{%
	\ryumeaning%
		{間違える}{}{}%
	}%
	}%
	{\tailPrint[]}%
\entry
	{\headPrint[kana={まちぼさーん},ipa={matʃibosaːɴ},pos={形}]}%
	{%
		\MeaningsBlock{%
	\ryumeaning%
		{待ち遠しく思う。待ちわびる}{}{}%
	}%
	}%
	{\tailPrint[]}%
\entry
	{\headPrint[kana={まちやー\tl{˥˩}},ipa={matʃijaː},pos={名},acc={F}]}%
	{%
		\MeaningsBlock{%
	\ryumeaning%
		{店。商店}{}{}%
	}%
	}%
	{\tailPrint[]}%
\entry
	{\headPrint[kana={まちゅ\tl{˩˥}},ipa={matʃu},pos={名},acc={R}]}%
	{%
		\MeaningsBlock{%
	\ryumeaning%
		{眉。睫毛}{}{}%
	}%
	}%
	{\tailPrint[]}%
\entry
	{\headPrint[kana={まちり\tl{˥˩}},ipa={matʃiri},pos={名},acc={F}]}%
	{%
		\MeaningsBlock{%
	\ryumeaning%
		{祭}{}{}%
	}%
	}%
	{\tailPrint[]}%
\entry
	{\headPrint[kana={まっふぁ\tl{˩˥}},ipa={maffa},pos={名},acc={R}]}%
	{%
		\MeaningsBlock{%
	\ryumeaning%
		{\ruby[g]{枕}{まくら}}{}{}%
	}%
	}%
	{\tailPrint[]}%
\entry
	{\headPrint[kana={まつら\tl{˥}},ipa={matsura},pos={名},acc={H}]}%
	{%
		\MeaningsBlock{%
	\ryumeaning%
		{船の竜骨。キール}{}{}%
	}%
	}%
	{\tailPrint[]}%
\entry
	{\headPrint[kana={まつん\tl{˩˥}},ipa={matsuɴ},pos={動},rui={2},para={またぬ},acc={R}]}%
	{%
		\MeaningsBlock{%
	\ryumeaning%
		{待つ}{}{}%
	}%
	}%
	{\tailPrint[]}%
\entry
	{\headPrint[kana={まとーば\tl{˥˩}},ipa={matoːba},pos={-},acc={F}]}%
	{%
		\MeaningsBlock{%
	\ryumeaning%
		{正直だ}{}{}%
	}%
	}%
	{\tailPrint[]}%
\entry
	{\headPrint[kana={まどぅぎるん\tl{˩˥}},ipa={madugiruɴ},pos={動},rui={3},para={まどぅぐぬ},acc={R}]}%
	{%
		\MeaningsBlock{%
	\ryumeaning%
		{片付ける}{}{}%
	}%
	}%
	{\tailPrint[]}%
\entry
	{\headPrint[kana={まとぅまるん\tl{˥˩}},ipa={matumaruɴ},pos={動},rui={1},para={まとぅまらぬ},acc={F}]}%
	{%
		\MeaningsBlock{%
	\ryumeaning%
		{まとまる}{}{}%
	}%
	}%
	{\tailPrint[]}%
\entry
	{\headPrint[kana={まとぅみん\tl{˥˩}},ipa={matumiɴ},pos={動},rui={3},para={まとぅむぬ},acc={F}]}%
	{%
		\MeaningsBlock{%
	\ryumeaning%
		{まとめる}{}{}%
	}%
	}%
	{\tailPrint[]}%
\entry
	{\headPrint[kana={まとむ\tl{˥˩}},ipa={matomu},pos={名},acc={F}]}%
	{%
		\MeaningsBlock{%
	\ryumeaning%
		{まとも。真後ろ}{}{}%
	}%
	}%
	{\tailPrint[]}%
\entry
	{\headPrint[kana={まな\tl{˥}},ipa={mana},pos={名},acc={H}]}%
	{%
		\MeaningsBlock{%
	\ryumeaning%
		{今。現在}{}{}%
	}%
	}%
	{\tailPrint[]}%
\entry
	{\headPrint[kana={まなぬ\tl{˥}\ws{}ゆー\tl{˥˩}},ipa={mananu juː},pos={句},acc={H F}]}%
	{%
		\MeaningsBlock{%
	\ryumeaning%
		{今の世。現代}{}{}%
	}%
	}%
	{\tailPrint[]}%
\entry
	{\headPrint[kana={まなばら\tl{˥˩}},ipa={manabara},pos={名},acc={F}]}%
	{%
		\MeaningsBlock{%
	\ryumeaning%
		{今頃}{}{}%
	}%
	}%
	{\tailPrint[]}%
\entry
	{\headPrint[kana={まなまな},ipa={manamana},pos={名}]}%
	{%
		\MeaningsBlock{%
	\ryumeaning%
		{今すぐ}{}{}%
	}%
	}%
	{\tailPrint[]}%
\entry
	{\headPrint[kana={まにあーすん\tl{˥}},ipa={mani.aːsuɴ},pos={動},rui={2b},para={まにあーさぬ},acc={H}]}%
	{%
		\MeaningsBlock{%
	\ryumeaning%
		{（時間に）間に合わせる}{}{}%
	}%
	}%
	{\tailPrint[]}%
\entry
	{\headPrint[kana={まにあうん},ipa={mani.auɴ},pos={動},para={まにあーぬ}]}%
	{%
		\MeaningsBlock{%
	\ryumeaning%
		{間に合う}{時間にも物についてもつかう}{}%
	}%
	}%
	{\tailPrint[]}%
\entry
	{\headPrint[kana={ま\josho{}べ\kako{}ー},ipa={mabeː},pos={副},acc={LHL}]}%
	{%
		\MeaningsBlock{%
	\ryumeaning%
		{もう少し。もうちょっと}{}{}%
	}%
	}%
	{\tailPrint[]}%
\entry
	{\headPrint[kana={まむるん\tl{˥˩}},ipa={mamuruɴ},pos={動},rui={1},para={まむらぬ},acc={F}]}%
	{%
		\MeaningsBlock{%
	\ryumeaning%
		{守る}{}{}%
	}%
	}%
	{\tailPrint[]}%
\entry
	{\headPrint[kana={まゆ\tl{˥˩}},ipa={maju},pos={名},acc={F}]}%
	{%
		\MeaningsBlock{%
	\ryumeaning%
		{猫}{\emb{まーゆ}とも言う}{}%
	}%
	}%
	{\tailPrint[]}%
\entry
	{\headPrint[kana={まゆ\tl{˩˥}},ipa={maju},pos={名},acc={R}]}%
	{%
		\MeaningsBlock{%
	\ryumeaning%
		{\ruby[g]{繭}{まゆ}}{}{}%
	}%
	}%
	{\tailPrint[]}%
\entry
	{\headPrint[kana={まら\tl{˩˥}},ipa={mara},pos={名},acc={R}]}%
	{%
		\MeaningsBlock{%
	\ryumeaning%
		{陰茎}{小児語で\emb{こっち}と言う}{}%
	}%
	}%
	{\tailPrint[]}%
\entry
	{\headPrint[kana={まらすん\tl{˥˩}},ipa={marasuɴ},pos={動},rui={2},para={まらはぬ},acc={F}]}%
	{%
		\MeaningsBlock{%
	\ryumeaning%
		{亡くなる}{「死ぬ」の尊敬語}{}%
	}%
	}%
	{\tailPrint[]}%
\entry
	{\headPrint[kana={まり\tl{˥˩}},ipa={mari},pos={名},acc={F}]}%
	{%
		\MeaningsBlock{%
	\ryumeaning%
		{生まれ。生地。血統}{}{}%
	}%
	}%
	{\tailPrint[]}%
\entry
	{\headPrint[kana={まりぐん\tl{˥˩}},ipa={mariguɴ},pos={動},rui={1},para={まりがぬ},acc={F}]}%
	{%
		\MeaningsBlock{%
	\ryumeaning%
		{束ねる}{}{}%
	}%
	}%
	{\tailPrint[]}%
\entry
	{\headPrint[kana={まりしき\tl{˥˩}},ipa={mariʃi̥ki},pos={名},acc={F}]}%
	{%
		\MeaningsBlock{%
	\ryumeaning%
		{臨月}{}{}%
	}%
	}%
	{\tailPrint[]}%
\entry
	{\headPrint[kana={まりじま\tl{˥˩}},ipa={maridʒima},pos={名},acc={F}]}%
	{%
		\MeaningsBlock{%
	\ryumeaning%
		{生まれ島。故郷。古里}{}{}%
	}%
	}%
	{\tailPrint[]}%
\entry
	{\headPrint[kana={まりどぅしぃ\tl{˥˩}},ipa={mariduːʃi},pos={名},acc={F}]}%
	{%
		\MeaningsBlock{%
	\ryumeaning%
		{生まれ年。生年}{}{}%
	}%
	}%
	{\tailPrint[]}%
\entry
	{\headPrint[kana={まりぴん\tl{˥˩}},ipa={maripiɴ},pos={名},acc={F}]}%
	{%
		\MeaningsBlock{%
	\ryumeaning%
		{誕生日}{}{}%
	}%
	}%
	{\tailPrint[]}%
\entry
	{\headPrint[kana={まりぶなー\tl{˥˩}},ipa={maribunaː},pos={名},acc={F}]}%
	{%
		\MeaningsBlock{%
	\ryumeaning%
		{生年祝い。誕生祝い}{}{}%
	}%
	}%
	{\tailPrint[]}%
\entry
	{\headPrint[kana={まりぽーるん\tl{˥˩}},ipa={maripoːruɴ},pos={動},rui={1},para={まりぽーらぬ},acc={F}]}%
	{%
		\MeaningsBlock{%
	\ryumeaning%
		{（大小便を）排泄し散らす}{}{}%
	}%
	}%
	{\tailPrint[]}%
\entry
	{\headPrint[kana={まりるん\tl{˥˩}},ipa={mariruɴ},pos={動},rui={3},para={まるぬ},acc={F}]}%
	{%
		\MeaningsBlock{%
	\ryumeaning%
		{生まれる。誕生する}{}{}%
	}%
	}%
	{\tailPrint[]}%
\entry
	{\headPrint[kana={まるばい\tl{˥˩}},ipa={marubai},pos={名},acc={F}]}%
	{%
		\MeaningsBlock{%
	\ryumeaning%
		{丸裸。全裸}{}{}%
	}%
	}%
	{\tailPrint[]}%
\entry
	{\headPrint[kana={まるばすん\tl{˥˩}},ipa={marubasuɴ},pos={動},rui={2},para={まりばはぬ},acc={F}]}%
	{%
		\MeaningsBlock{%
	\ryumeaning%
		{強く叩く}{}{}%
	}%
	}%
	{\tailPrint[]}%
\entry
	{\headPrint[kana={まるぶさ\tl{˥˩}},ipa={marubusa},pos={名},acc={F}]}%
	{%
		\MeaningsBlock{%
	\ryumeaning%
		{ダンドク}{南方系の草花}{}%
	}%
	}%
	{\tailPrint[]}%
\entry
	{\headPrint[kana={まるぶちん\tl{˥˩}},ipa={marubutʃiɴ},pos={動},rui={3},para={まるぶとぅぬ},acc={F}]}%
	{%
		\MeaningsBlock{%
	\ryumeaning%
		{転ぶ。転がる。転倒する}{}{}%
	}%
	}%
	{\tailPrint[]}%
\entry
	{\headPrint[kana={まるぶん\tl{˥˩}},ipa={marubuɴ},pos={動},rui={1},para={まるばぬ},acc={F}]}%
	{%
		\MeaningsBlock{%
	\ryumeaning%
		{転ぶ。まろぶ}{}{}%
	}%
	}%
	{\tailPrint[]}%
\entry
	{\headPrint[kana={まるん\tl{˩˥}},ipa={maruɴ},pos={動},rui={1},para={まらぬ},acc={R}]}%
	{%
		\MeaningsBlock{%
	\ryumeaning%
		{（糞を）放つ。出す}{}{}%
	}%
	}%
	{\tailPrint[]}%
\entry
	{\headPrint[kana={まろ\tl{˩˥}はん},ipa={marohaɴ},pos={形},para={まろへぬ},acc={R}]}%
	{%
		\MeaningsBlock{%
	\ryumeaning%
		{短い。低い}{}{}%
	}%
	}%
	{\tailPrint[]}%
\entry
	{\headPrint[kana={まん\tl{˥˩}},ipa={maɴ},pos={名},acc={F}]}%
	{%
		\MeaningsBlock{%
	\ryumeaning%
		{万}{数字の単位}{}%
	}%
	}%
	{\tailPrint[]}%
\entry
	{\headPrint[kana={まんいち\tl{˥˩}},ipa={maɴ.itʃi},pos={副},acc={F}]}%
	{%
		\MeaningsBlock{%
	\ryumeaning%
		{万一。もしも}{}{}%
	}%
	}%
	{\tailPrint[]}%
\entry
	{\headPrint[kana={まんが\tl{˩˥}},ipa={maŋga},pos={副},acc={R}]}%
	{%
		\MeaningsBlock{%
	\ryumeaning%
		{真っ直ぐ}{}{}%
	}%
	}%
	{\tailPrint[]}%
\entry
	{\headPrint[kana={まんがたんが},ipa={maŋgataŋga},pos={擬}]}%
	{%
		\MeaningsBlock{%
	\ryumeaning%
		{率直に}{}{}%
	}%
	}%
	{\tailPrint[]}%
\entry
	{\headPrint[kana={まんかりかんかり},ipa={maŋkarikaŋkari},pos={擬}]}%
	{%
		\MeaningsBlock{%
	\ryumeaning%
		{曲がりくねって}{}{}%
	}%
	}%
	{\tailPrint[]}%
\entry
	{\headPrint[kana={まんかるん\tl{˥˩}},ipa={maŋkaruɴ},pos={動},rui={1},para={まんからぬ},acc={F}]}%
	{%
		\MeaningsBlock{%
	\ryumeaning%
		{曲がる。たわむ}{}{}%
	}%
	}%
	{\tailPrint[]}%
\entry
	{\headPrint[kana={まんきるん\tl{˥˩}},ipa={maŋkiruɴ},pos={動},rui={3},para={まんくぬ},acc={F}]}%
	{%
		\MeaningsBlock{%
	\ryumeaning%
		{曲げる}{}{}%
	}%
	}%
	{\tailPrint[]}%
\entry
	{\headPrint[kana={まんざーすん},ipa={mandzaːsuɴ},pos={動}]}%
	{%
		\MeaningsBlock{%
	\ryumeaning%
		{混ぜ合わせる}{}{}%
	\ryumeaning%
		{加える}{}{}%
	\ryumeaning%
		{一緒にする}{}{}%
	}%
	}%
	{\tailPrint[]}%
\entry
	{\headPrint[kana={まんざるん\tl{˩˥}},ipa={mandzaruɴ},pos={動},rui={1},para={まんざらぬ},acc={R}]}%
	{%
		\MeaningsBlock{%
	\ryumeaning%
		{混ざる。混合する}{}{}%
	}%
	}%
	{\tailPrint[]}%
\entry
	{\headPrint[kana={まんじょー\tl{˥}},ipa={mandʒoː},pos={名},acc={H}]}%
	{%
		\MeaningsBlock{%
	\ryumeaning%
		{パパイア}{果樹名}{}%
	}%
	}%
	{\tailPrint[]}%
\entry
	{\headPrint[kana={まんじるん\tl{˩˥}},ipa={mandʒiruɴ},pos={動},rui={3},para={まんずぬ},acc={R}]}%
	{%
		\MeaningsBlock{%
	\ryumeaning%
		{混ぜる。混合する}{}{}%
	}%
	}%
	{\tailPrint[]}%
\entry
	{\headPrint[kana={まんしん\tl{˥˩}},ipa={manʃiɴ},pos={名},acc={F}]}%
	{%
		\MeaningsBlock{%
	\ryumeaning%
		{一杯。満載}{満船からか}{}%
	}%
	}%
	{\tailPrint[]}%
\entry
	{\headPrint[kana={まんしん\tl{˩˥}},ipa={manʃiɴ},pos={名},acc={R}]}%
	{%
		\MeaningsBlock{%
	\ryumeaning%
		{\ruby[g]{首}{くび}}{}{}%
	}%
	}%
	{\tailPrint[]}%
\entry
	{\headPrint[kana={まんしん\ws{}なるん},ipa={manʃiɴ naruɴ},pos={句}]}%
	{%
		\MeaningsBlock{%
	\ryumeaning%
		{満杯になる}{}{}%
	}%
	}%
	{\tailPrint[]}%
\entry
	{\headPrint[kana={まんず},ipa={mandzu},pos={名}]}%
	{%
		\MeaningsBlock{%
	\ryumeaning%
		{和え物}{}{}%
	}%
	}%
	{\tailPrint[]}%
\entry
	{\headPrint[kana={まんつぁ\tl{˩˥}},ipa={mantsa},pos={名},acc={R}]}%
	{%
		\MeaningsBlock{%
	\ryumeaning%
		{まな板}{}{}%
	}%
	}%
	{\tailPrint[]}%
\entry
	{\headPrint[kana={まんどん\tl{˩˥}},ipa={mandoɴ},pos={副},acc={R}]}%
	{%
		\MeaningsBlock{%
	\ryumeaning%
		{一杯。沢山}{}{}%
	}%
	}%
	{\tailPrint[]}%
\entry
	{\headPrint[kana={ま\josho{}んなー\kako{}が},ipa={mannaːga},pos={名},acc={特殊}]}%
	{%
		\MeaningsBlock{%
	\ryumeaning%
		{真ん中}{}{}%
	}%
	}%
	{\tailPrint[]}%
\LetterHead{み}
\entry
	{\headPrint[kana={みー\tl{˥}},ipa={miː},pos={名},acc={H}]}%
	{%
		\MeaningsBlock{%
	\ryumeaning%
		{\ruby[g]{巳}{み}}{十二支の巳}{}%
	}%
	}%
	{\tailPrint[]}%
\entry
	{\headPrint[kana={みー\tl{˥˩}},ipa={miː},pos={接頭},acc={F}]}%
	{%
		\MeaningsBlock{%
	\ryumeaning%
		{新しい}{}{}%
	}%
	}%
	{\tailPrint[]}%
\entry
	{\headPrint[kana={みーうとすん},ipa={miː.utosuɴ},pos={動}]}%
	{%
		\MeaningsBlock{%
	\ryumeaning%
		{見落とす}{}{}%
	}%
	}%
	{\tailPrint[]}%
\entry
	{\headPrint[kana={みーかんがん\tl{˩˥}},ipa={miːkaŋgaɴ},pos={名},acc={R}]}%
	{%
		\MeaningsBlock{%
	\ryumeaning%
		{眼鏡。水中眼鏡}{}{}%
	}%
	}%
	{\tailPrint[]}%
\entry
	{\headPrint[kana={みーぐりしゃーん},ipa={miːguriʃaːɴ},pos={形},acc={-}]}%
	{%
		\MeaningsBlock{%
	\ryumeaning%
		{見にくい。見苦しい}{}{}%
	}%
	}%
	{\tailPrint[]}%
\entry
	{\headPrint[kana={みーしきるん},ipa={miːʃi̥kiruɴ},pos={句},para={みーすくぬ？}]}%
	{%
		\MeaningsBlock{%
	\ryumeaning%
		{目を付ける。睨む}{}{}%
	}%
	}%
	{\tailPrint[]}%
\entry
	{\headPrint[kana={みー\tl{˩˥}\ws{}しちん\tl{˥˩}},ipa={miː ʃi̥tʃiɴ},pos={句},acc={R F}]}%
	{%
		\MeaningsBlock{%
	\ryumeaning%
		{見捨てる}{}{}%
	}%
	}%
	{\tailPrint[]}%
\entry
	{\headPrint[kana={みーち\tl{˥˩}},ipa={miːtʃi},pos={名},acc={F}]}%
	{%
		\MeaningsBlock{%
	\ryumeaning%
		{三つ}{}{}%
	}%
	}%
	{\tailPrint[note={簡略化\emb{みぃ}}]}%
\entry
	{\headPrint[kana={みーとーすん\tl{˩˥}},ipa={miːtoːsuɴ},pos={動},rui={2},para={みーとーはぬ},acc={R}]}%
	{%
		\MeaningsBlock{%
	\ryumeaning%
		{見通す。続けて見る}{}{}%
	}%
	}%
	{\tailPrint[]}%
\entry
	{\headPrint[kana={みーどぅー\tl{˩˥}さーん},ipa={miːduːsaːɴ},pos={形},acc={R}]}%
	{%
		\MeaningsBlock{%
	\ryumeaning%
		{久しく会わない。中々見ることのない}{}{}%
	}%
	}%
	{\tailPrint[]}%
\entry
	{\headPrint[kana={みーとぅどぅきるん\tl{˩˥}},ipa={miːtu̥dukiɴ},pos={動},rui={3},para={みーとぅどぅくぬ},acc={R}]}%
	{%
		\MeaningsBlock{%
	\ryumeaning%
		{見届ける}{}{}%
	}%
	}%
	{\tailPrint[]}%
\entry
	{\headPrint[kana={みーならすん\tl{˩˥}},ipa={miːnarasuɴ},pos={動},rui={2},para={みーならはぬ},acc={R}]}%
	{%
		\MeaningsBlock{%
	\ryumeaning%
		{見習う}{}{}%
	}%
	}%
	{\tailPrint[]}%
\entry
	{\headPrint[kana={みーなりゃん},ipa={miːnarjaɴ},pos={動（継）}]}%
	{%
		\MeaningsBlock{%
	\ryumeaning%
		{見慣れる}{}{}%
	}%
	}%
	{\tailPrint[]}%
\entry
	{\headPrint[kana={みーにち\tl{˥˩}},ipa={miːnitʃi},pos={名},acc={F}]}%
	{%
		\MeaningsBlock{%
	\ryumeaning%
		{命日}{}{}%
	}%
	}%
	{\tailPrint[]}%
\entry
	{\headPrint[kana={みーぬがすん\tl{˩˥}},ipa={miːnugasuɴ},pos={動},rui={2},para={みーぬがはぬ},acc={R}]}%
	{%
		\MeaningsBlock{%
	\ryumeaning%
		{見逃す}{}{}%
	}%
	}%
	{\tailPrint[]}%
\entry
	{\headPrint[kana={みーばっぺー\tl{˩˥}},ipa={miːbappeː},pos={名},acc={R}]}%
	{%
		\MeaningsBlock{%
	\ryumeaning%
		{見間違い。見誤り}{}{}%
	}%
	}%
	{\tailPrint[]}%
\entry
	{\headPrint[kana={みーふかさりん\tl{˩˥}},ipa={miːfu̥kasariɴ},pos={動},para={みーふかさるぬ},acc={R}]}%
	{%
		\MeaningsBlock{%
	\ryumeaning%
		{見透かされる。透き通って見える}{}{}%
	}%
	}%
	{\tailPrint[]}%
\entry
	{\headPrint[kana={みーふかすん\tl{˩˥}},ipa={miːfu̥kasuɴ},pos={動},rui={2},para={みふかさぬ},acc={R}]}%
	{%
		\MeaningsBlock{%
	\ryumeaning%
		{見通す。見透かす}{}{}%
	}%
	}%
	{\tailPrint[]}%
\entry
	{\headPrint[kana={みーぶりん\tl{˩˥}},ipa={miːburiɴ},pos={動},rui={3},para={みーぶるぬ},acc={R}]}%
	{%
		\MeaningsBlock{%
	\ryumeaning%
		{見惚れる。見とれる}{}{}%
	}%
	}%
	{\tailPrint[]}%
\entry
	{\headPrint[kana={みーむぬ\tl{˩˥}},ipa={miːmunu},pos={名},acc={R}]}%
	{%
		\MeaningsBlock{%
	\ryumeaning%
		{見物。見せ物}{}{}%
	}%
	}%
	{\tailPrint[]}%
\entry
	{\headPrint[kana={みーむぬ\tl{˩˥}},ipa={miːmunu},pos={名},acc={R}]}%
	{%
		\MeaningsBlock{%
	\ryumeaning%
		{雌}{}{}%
	}%
	}%
	{\tailPrint[]}%
\entry
	{\headPrint[kana={みーむん\tl{˥˩}},ipa={miːmuɴ},pos={名},acc={F}]}%
	{%
		\MeaningsBlock{%
	\ryumeaning%
		{新しい物。新品}{}{}%
	}%
	}%
	{\tailPrint[]}%
\entry
	{\headPrint[kana={みーやっ\tl{˩˥}さーん},ipa={miːjassaːɴ},pos={形},acc={R}]}%
	{%
		\MeaningsBlock{%
	\ryumeaning%
		{見やすい}{}{}%
	}%
	}%
	{\tailPrint[]}%
\entry
	{\headPrint[kana={みーん\tl{˩˥}ぱちま\tl{˥˩}はん},ipa={miːmpḁtʃimahaɴ},pos={形},acc={R+F}]}%
	{%
		\MeaningsBlock{%
	\ryumeaning%
		{目が眩しい}{}{}%
	}%
	}%
	{\tailPrint[]}%
\entry
	{\headPrint[kana={みあてぃ\tl{˥}},ipa={mi.ati},pos={名},acc={H}]}%
	{%
		\MeaningsBlock{%
	\ryumeaning%
		{目当て。目標}{}{}%
	}%
	}%
	{\tailPrint[]}%
\entry
	{\headPrint[kana={みぐとぅ\tl{˩˥}},ipa={migutu},pos={名},acc={R}]}%
	{%
		\MeaningsBlock{%
	\ryumeaning%
		{見事。立派}{}{}%
	}%
	}%
	{\tailPrint[]}%
\entry
	{\headPrint[kana={みぐらすん\tl{˥˩}},ipa={migurasuɴ},pos={動},rui={2},para={みぐらはぬ},acc={F}]}%
	{%
		\MeaningsBlock{%
	\ryumeaning%
		{めぐらす。回す}{}{}%
	}%
	}%
	{\tailPrint[]}%
\entry
	{\headPrint[kana={みぐりしゃ\tl{˩˥}はん},ipa={miguriʃahaɴ},pos={形},para={みぐりしゃへぬ},acc={R}]}%
	{%
		\MeaningsBlock{%
	\ryumeaning%
		{見苦しい。みっともない}{}{}%
	}%
	}%
	{\tailPrint[]}%
\entry
	{\headPrint[kana={みぐるん\tl{˥˩}},ipa={miguruɴ},pos={動},rui={1},para={みぐらぬ},acc={F}]}%
	{%
		\MeaningsBlock{%
	\ryumeaning%
		{めぐる。回る}{}{}%
	}%
	}%
	{\tailPrint[]}%
\entry
	{\headPrint[kana={みざし\tl{˩˥}},ipa={midzaʃi},pos={副},acc={R}]}%
	{%
		\MeaningsBlock{%
	\ryumeaning%
		{大変な。たいそう}{}{}%
	}%
	}%
	{\tailPrint[]}%
\entry
	{\headPrint[kana={みざし\tl{˩˥}},ipa={midzaʃi},pos={名},acc={R}]}%
	{%
		\MeaningsBlock{%
	\ryumeaning%
		{\ruby[g]{目差}{めざし}}{王府時代の下級役人}{}%
	}%
	}%
	{\tailPrint[]}%
\entry
	{\headPrint[kana={みざしくとぅ},ipa={midzaʃiku̥tu},pos={名}]}%
	{%
		\MeaningsBlock{%
	\ryumeaning%
		{大事。大変なこと}{}{}%
	}%
	}%
	{\tailPrint[]}%
\entry
	{\headPrint[kana={みざとぅ\tl{˩˥}},ipa={midzatu},pos={名},acc={R}]}%
	{%
		\MeaningsBlock{%
	\ryumeaning%
		{\ruby[g]{真栄}{まえ}\ruby[g]{里}{ざと}}{石垣島の集落名}{}%
	}%
	}%
	{\tailPrint[]}%
\entry
	{\headPrint[kana={みし\tl{˥˩}},ipa={miʃi},pos={名},acc={F}]}%
	{%
		\MeaningsBlock{%
	\ryumeaning%
		{\ruby[g]{神酒}{みき}}{}{}%
	}%
	}%
	{\tailPrint[]}%
\entry
	{\headPrint[kana={みしーみし},ipa={miʃiːmiʃi},pos={擬}]}%
	{%
		\MeaningsBlock{%
	\ryumeaning%
		{みすみす。まんまと}{}{}%
	}%
	}%
	{\tailPrint[]}%
\entry
	{\headPrint[kana={みじぃ\tl{˥˩}},ipa={midʒi},pos={名},acc={F}]}%
	{%
		\MeaningsBlock{%
	\ryumeaning%
		{水。淡水}{}{}%
	}%
	}%
	{\tailPrint[]}%
\entry
	{\headPrint[kana={みじぃがーさ\tl{˩˥}},ipa={midzugaːsa},pos={名},acc={R}]}%
	{%
		\MeaningsBlock{%
	\ryumeaning%
		{水疱瘡}{\emb{みじがーさ}とも発音する}{}%
	}%
	}%
	{\tailPrint[]}%
\entry
	{\headPrint[kana={みじぃがーみ\tl{˩˥}},ipa={midzugaːmi},pos={名},acc={R}]}%
	{%
		\MeaningsBlock{%
	\ryumeaning%
		{\ruby[g]{水}{みず}\ruby[g]{甕}{がめ}}{}{}%
	}%
	}%
	{\tailPrint[]}%
\entry
	{\headPrint[kana={みじぃがーり\tl{˩˥}},ipa={midzugaːri},pos={名},acc={R}]}%
	{%
		\MeaningsBlock{%
	\ryumeaning%
		{水当たり}{}{}%
	}%
	}%
	{\tailPrint[]}%
\entry
	{\headPrint[kana={みしきるん\tl{˩˥}},ipa={miʃi̥kiruɴ},pos={動},para={めすからぬ},acc={R}]}%
	{%
		\MeaningsBlock{%
	\ryumeaning%
		{見つける}{}{}%
	}%
	}%
	{\tailPrint[]}%
\entry
	{\headPrint[kana={みしくるみん\tl{˩˥}},ipa={miʃi̥kurumiɴ},pos={名},acc={R}]}%
	{%
		\MeaningsBlock{%
	\ryumeaning%
		{耳たぶ}{}{}%
	}%
	}%
	{\tailPrint[]}%
\entry
	{\headPrint[kana={みしこーみしこ},ipa={miʃi̥koːmiʃi̥ko},pos={擬}]}%
	{%
		\MeaningsBlock{%
	\ryumeaning%
		{慎重に。丁寧に}{}{}%
	}%
	}%
	{\tailPrint[]}%
\entry
	{\headPrint[kana={みしこ\tl{˥}はん},ipa={miʃi̥kohaɴ},pos={形},acc={H}]}%
	{%
		\MeaningsBlock{%
	\ryumeaning%
		{神聖だ。神々しい}{}{}%
	}%
	}%
	{\tailPrint[]}%
\entry
	{\headPrint[kana={みしゃー\tl{˥}\ws{}なるん\tl{˩˥}},ipa={miʃaː naruɴ},pos={句},acc={H R}]}%
	{%
		\MeaningsBlock{%
	\ryumeaning%
		{良くなる。治る}{}{}%
	}%
	}%
	{\tailPrint[]}%
\entry
	{\headPrint[kana={みしゃ\tl{˥}はん},ipa={miʃahaɴ},pos={形},para={みしゃへぬ},acc={H}]}%
	{%
		\MeaningsBlock{%
	\ryumeaning%
		{良い。よろしい}{}{}%
	}%
	}%
	{\tailPrint[]}%
\entry
	{\headPrint[kana={みしゅ\tl{˩˥}},ipa={miʃu},pos={名},acc={R}]}%
	{%
		\MeaningsBlock{%
	\ryumeaning%
		{\ruby[g]{味噌}{みそ}}{}{}%
	}%
	}%
	{\tailPrint[]}%
\entry
	{\headPrint[kana={みしゅ\tl{˩˥}どぅり\tl{˥˩}},ipa={miʃuduri},pos={名},acc={R+F}]}%
	{%
		\MeaningsBlock{%
	\ryumeaning%
		{雀}{}{}%
	}%
	}%
	{\tailPrint[]}%
\entry
	{\headPrint[kana={みしるん\tl{˩˥}},ipa={miʃiruɴ},pos={動},rui={3},para={みしゅぬ},acc={R}]}%
	{%
		\MeaningsBlock{%
	\ryumeaning%
		{見せる}{}{}%
	}%
	}%
	{\tailPrint[]}%
\entry
	{\headPrint[kana={みずあみ},ipa={midzu.ami},pos={名},acc={-}]}%
	{%
		\MeaningsBlock{%
	\ryumeaning%
		{水浴}{}{}%
	}%
	}%
	{\tailPrint[]}%
\entry
	{\headPrint[kana={みずぬ\ws{}こー},ipa={midzunu koː},pos={名}]}%
	{%
		\MeaningsBlock{%
	\ryumeaning%
		{水の香}{法要時の邪霊除け}{}%
	}%
	}%
	{\tailPrint[]}%
\entry
	{\headPrint[kana={みずら\tl{˩˥}はん},ipa={midzurahaɴ},pos={形},para={みじらへぬ},acc={R}]}%
	{%
		\MeaningsBlock{%
	\ryumeaning%
		{珍しい。不思議だ}{}{}%
	}%
	}%
	{\tailPrint[]}%
\entry
	{\headPrint[kana={みぞーり\tl{˥˩}},ipa={midzoːri},pos={名},acc={F}]}%
	{%
		\MeaningsBlock{%
	\ryumeaning%
		{溝。排水溝}{}{}%
	}%
	}%
	{\tailPrint[]}%
\entry
	{\headPrint[kana={みだつん\tl{˩˥}},ipa={midatsuɴ},pos={動},para={みだたぬ},acc={R}]}%
	{%
		\MeaningsBlock{%
	\ryumeaning%
		{目立つ。際立つ}{目立って見える}{}%
	}%
	}%
	{\tailPrint[]}%
\entry
	{\headPrint[kana={みたてぃるん\tl{˥}},ipa={mitatiruɴ},pos={動},acc={H}]}%
	{%
		\MeaningsBlock{%
	\ryumeaning%
		{見たてる}{見て選び定める}{}%
	}%
	}%
	{\tailPrint[]}%
\entry
	{\headPrint[kana={みちぃ\tl{˥˩}},ipa={mitʃi},pos={名},acc={F}]}%
	{%
		\MeaningsBlock{%
	\ryumeaning%
		{道。道路}{}{}%
	}%
	}%
	{\tailPrint[]}%
\entry
	{\headPrint[kana={みちぃくせー\tl{˥˩}},ipa={mitʃikuseː},pos={名},acc={F}]}%
	{%
		\MeaningsBlock{%
	\ryumeaning%
		{道普請。道路修理}{}{}%
	}%
	}%
	{\tailPrint[]}%
\entry
	{\headPrint[kana={みちさはん},ipa={mitʃi̥sahaɴ},pos={形},para={みちさへぬ}]}%
	{%
		\MeaningsBlock{%
	\ryumeaning%
		{\ruby[g]{太}{ふと}い}{}{}%
	}%
	}%
	{\tailPrint[]}%
\entry
	{\headPrint[kana={みちすねー\tl{˥˩}},ipa={mitʃisu̥neː},pos={名},acc={F}]}%
	{%
		\MeaningsBlock{%
	\ryumeaning%
		{（旧盆の）仮装行列}{}{}%
	}%
	}%
	{\tailPrint[]}%
\entry
	{\headPrint[kana={みちなりん},ipa={mitʃi̥nariɴ},pos={名},acc={-}]}%
	{%
		\MeaningsBlock{%
	\ryumeaning%
		{一昨年}{}{}%
	}%
	}%
	{\tailPrint[]}%
\entry
	{\headPrint[kana={みっくゎ\tl{˩˥}},ipa={mikkwa},pos={名},acc={R}]}%
	{%
		\MeaningsBlock{%
	\ryumeaning%
		{めくら。盲人}{}{}%
	}%
	}%
	{\tailPrint[]}%
\entry
	{\headPrint[kana={みっちん\tl{˩˥}},ipa={mittsiɴ},pos={名},acc={R}]}%
	{%
		\MeaningsBlock{%
	\ryumeaning%
		{目玉。眼球}{}{}%
	}%
	}%
	{\tailPrint[note={語中の二重子音は口蓋化していない}]}%
\entry
	{\headPrint[kana={みっと\tl{˩˥}はん},ipa={mittohaɴ},pos={形},acc={R}]}%
	{%
		\MeaningsBlock{%
	\ryumeaning%
		{見苦しい。みっともない}{}{}%
	}%
	}%
	{\tailPrint[]}%
\entry
	{\headPrint[kana={\josho{}みっふぁー\kako{}むん},ipa={miffaːmuɴ},pos={名},acc={H+L}]}%
	{%
		\MeaningsBlock{%
	\ryumeaning%
		{嫌われ者}{}{}%
	}%
	}%
	{\tailPrint[]}%
\entry
	{\headPrint[kana={みっふぁはん},ipa={miffahaɴ},pos={形},para={にふっあへぬ},acc={-}]}%
	{%
		\MeaningsBlock{%
	\ryumeaning%
		{憎らしい。嫌い}{}{}%
	}%
	}%
	{\tailPrint[]}%
\entry
	{\headPrint[kana={みっ\josho{}ふぁ\kako{}ん},ipa={miffaɴ},pos={副},acc={LHL}]}%
	{%
		\MeaningsBlock{%
	\ryumeaning%
		{どんどん。一目散に}{}{}%
	}%
	}%
	{\tailPrint[]}%
\entry
	{\headPrint[kana={みとぅどぅぎるん\tl{˩˥}},ipa={mitudugiɴ},pos={動},rui={3},para={みとぅどぅぐぬ},acc={R}]}%
	{%
		\MeaningsBlock{%
	\ryumeaning%
		{見届ける}{}{}%
	}%
	}%
	{\tailPrint[]}%
\entry
	{\headPrint[kana={みとぅみるん\tl{˩˥}},ipa={mitumiruɴ},pos={動},rui={1},para={みとぅむぬ},acc={R}]}%
	{%
		\MeaningsBlock{%
	\ryumeaning%
		{認める}{}{}%
	}%
	}%
	{\tailPrint[]}%
\entry
	{\headPrint[kana={みどぅむ\tl{˩˥}},ipa={midumu},pos={名},acc={R}]}%
	{%
		\MeaningsBlock{%
	\ryumeaning%
		{女。女性}{}{}%
	}%
	}%
	{\tailPrint[]}%
\entry
	{\headPrint[kana={みどぅんたま\tl{˩˥}},ipa={miduntama},pos={名},acc={R}]}%
	{%
		\MeaningsBlock{%
	\ryumeaning%
		{女の子。娘}{}{}%
	}%
	}%
	{\tailPrint[]}%
\entry
	{\headPrint[kana={みなー\tl{˥˩}},ipa={minaː},pos={名},acc={F}]}%
	{%
		\MeaningsBlock{%
	\ryumeaning%
		{庭}{}{}%
	}%
	}%
	{\tailPrint[]}%
\entry
	{\headPrint[kana={みなが\tl{˥˩}},ipa={minaga},pos={名},acc={F}]}%
	{%
		\MeaningsBlock{%
	\ryumeaning%
		{庭}{}{}%
	}%
	}%
	{\tailPrint[]}%
\entry
	{\headPrint[kana={みなすけぬ\ws{}ぱんだー},ipa={minasu̥kenu pandaː},pos={句}]}%
	{%
		\MeaningsBlock{%
	\ryumeaning%
		{豊年祭時の神司達}{}{}%
	}%
	}%
	{\tailPrint[]}%
\entry
	{\headPrint[kana={みなぶ\tl{˩˥}},ipa={minabu},pos={名},acc={R}]}%
	{%
		\MeaningsBlock{%
	\ryumeaning%
		{\ruby[g]{藁}{わら}\ruby[g]{蓆}{むしろ}}{穀物を干すのに使用}{}%
	}%
	}%
	{\tailPrint[]}%
\entry
	{\headPrint[kana={みなれー\tl{˩˥}},ipa={minareː},pos={名},acc={R}]}%
	{%
		\MeaningsBlock{%
	\ryumeaning%
		{見習い}{}{}%
	}%
	}%
	{\tailPrint[]}%
\entry
	{\headPrint[kana={みなん\tl{˩˥}},ipa={minaɴ},pos={名},acc={R}]}%
	{%
		\MeaningsBlock{%
	\ryumeaning%
		{貝}{貝類の総称}{}%
	}%
	}%
	{\tailPrint[]}%
\entry
	{\headPrint[kana={みなんぬ\tl{˩˥}\ws{}くー\tl{˥}},ipa={minannu kuː},pos={句},acc={R H}]}%
	{%
		\MeaningsBlock{%
	\ryumeaning%
		{貝殻}{}{}%
	}%
	}%
	{\tailPrint[]}%
\entry
	{\headPrint[kana={みぬがすん\tl{˩˥}},ipa={minugasuɴ},pos={動},rui={2},para={みぬがはぬ},acc={R}]}%
	{%
		\MeaningsBlock{%
	\ryumeaning%
		{見逃す。見落とす}{}{}%
	}%
	}%
	{\tailPrint[]}%
\entry
	{\headPrint[kana={みのーすん\tl{˩˥}},ipa={minoːsuɴ},pos={動},rui={2},para={みのーさぬ},acc={R}]}%
	{%
		\MeaningsBlock{%
	\ryumeaning%
		{見直す}{}{}%
	\ryumeaning%
		{見て許す}{}{}%
	}%
	}%
	{\tailPrint[]}%
\entry
	{\headPrint[kana={みまちごーん\tl{˩˥}},ipa={mimatʃigoːɴ},pos={動},rui={1},para={みーまちげーらぬ},acc={R}]}%
	{%
		\MeaningsBlock{%
	\ryumeaning%
		{見間違える}{}{}%
	\ryumeaning%
		{見損なう}{}{}%
	}%
	}%
	{\tailPrint[]}%
\entry
	{\headPrint[kana={みみじぃ\tl{˩˥}},ipa={mimidʒi},pos={名},acc={R}]}%
	{%
		\MeaningsBlock{%
	\ryumeaning%
		{ミミズ}{}{}%
	}%
	}%
	{\tailPrint[]}%
\entry
	{\headPrint[kana={みゃーぐ\tl{˥˩}},ipa={mjaːgu},pos={名},acc={F}]}%
	{%
		\MeaningsBlock{%
	\ryumeaning%
		{宮古}{宮古島}{}%
	}%
	}%
	{\tailPrint[]}%
\entry
	{\headPrint[kana={みゃぐ\tl{˩˥}},ipa={mjagu},pos={名},acc={R}]}%
	{%
		\MeaningsBlock{%
	\ryumeaning%
		{脈}{脈拍}{}%
	}%
	}%
	{\tailPrint[]}%
\entry
	{\headPrint[kana={みやくつぇー\tl{˥˩}},ipa={mijaku̥tseː},pos={名},acc={F}]}%
	{%
		\MeaningsBlock{%
	\ryumeaning%
		{（御嶽や神道の）除草清掃}{}{}%
	}%
	}%
	{\tailPrint[]}%
\entry
	{\headPrint[kana={みやらび\tl{˩˥}},ipa={mijarabi},pos={名},acc={R}]}%
	{%
		\MeaningsBlock{%
	\ryumeaning%
		{乙女。女わらべ。娘}{}{}%
	}%
	}%
	{\tailPrint[]}%
\entry
	{\headPrint[kana={みゃん},ipa={mjaɴ},pos={動（継）}]}%
	{%
		\MeaningsBlock{%
	\ryumeaning%
		{膿む}{}{}%
	}%
	}%
	{\tailPrint[]}%
\entry
	{\headPrint[kana={みゃん\tl{˥}},ipa={mjaɴ},pos={動（継）},para={まーぬ},acc={H}]}%
	{%
		\MeaningsBlock{%
	\ryumeaning%
		{熟する。熟れる}{}{}%
	}%
	}%
	{\tailPrint[]}%
\entry
	{\headPrint[kana={みょーでん\tl{˥˩}},ipa={mjoːdeɴ},pos={名},acc={F}]}%
	{%
		\MeaningsBlock{%
	\ryumeaning%
		{名代。代理}{}{}%
	}%
	}%
	{\tailPrint[]}%
\entry
	{\headPrint[kana={みり\tl{˩˥}\ws{}みらぬ\tl{˩˥}},ipa={miri miranu},pos={句},para={みりみん},acc={R R}]}%
	{%
		\MeaningsBlock{%
	\ryumeaning%
		{見覚えがない}{}{}%
	}%
	}%
	{\tailPrint[]}%
\entry
	{\headPrint[kana={みるく\tl{˥˩}},ipa={miruku},pos={名},acc={F}]}%
	{%
		\MeaningsBlock{%
	\ryumeaning%
		{弥勒神。弥勒菩薩}{}{}%
	}%
	}%
	{\tailPrint[]}%
\entry
	{\headPrint[kana={みるくゆー\tl{˥˩}},ipa={mirukujuː},pos={名},acc={F}]}%
	{%
		\MeaningsBlock{%
	\ryumeaning%
		{豊穣で平和の世}{「弥勒神の世」の義}{}%
	}%
	}%
	{\tailPrint[]}%
\entry
	{\headPrint[kana={みるん\tl{˩˥}},ipa={miruɴ},pos={動},rui={1},para={みらぬ},acc={R}]}%
	{%
		\MeaningsBlock{%
	\ryumeaning%
		{見る}{}{}%
	}%
	}%
	{\tailPrint[]}%
\entry
	{\headPrint[kana={みん\tl{˩˥}},ipa={miɴ},pos={名},acc={R}]}%
	{%
		\MeaningsBlock{%
	\ryumeaning%
		{目}{}{}%
	}%
	}%
	{\tailPrint[]}%
\entry
	{\headPrint[kana={みん\tl{˩˥}},ipa={miɴ},pos={名},acc={R}]}%
	{%
		\MeaningsBlock{%
	\ryumeaning%
		{耳}{}{}%
	}%
	}%
	{\tailPrint[]}%
\entry
	{\headPrint[kana={みん\tl{˩˥}},ipa={miɴ},pos={名},acc={R}]}%
	{%
		\MeaningsBlock{%
	\ryumeaning%
		{穴}{}{}%
	}%
	}%
	{\tailPrint[]}%
\entry
	{\headPrint[kana={みんかー\tl{˥}},ipa={miŋkaː},pos={名},acc={H}]}%
	{%
		\MeaningsBlock{%
	\ryumeaning%
		{\ruby[g]{聾}{つんぼ}}{}{}%
	}%
	}%
	{\tailPrint[]}%
\entry
	{\headPrint[kana={みんくぱり\tl{˩˥}},ipa={miŋku̥pari},pos={-},acc={R}]}%
	{%
		\MeaningsBlock{%
	\ryumeaning%
		{目がさえる}{}{}%
	}%
	}%
	{\tailPrint[]}%
\entry
	{\headPrint[kana={みんぐる\tl{˩˥}},ipa={miŋguru},pos={名},acc={R}]}%
	{%
		\MeaningsBlock{%
	\ryumeaning%
		{キクラゲ}{食用になる茸}{}%
	}%
	}%
	{\tailPrint[]}%
\entry
	{\headPrint[kana={みんぐるみん\tl{˩˥}},ipa={miŋgurumin},pos={名},acc={R}]}%
	{%
		\MeaningsBlock{%
	\ryumeaning%
		{キクラゲ}{食用になる茸}{}%
	}%
	}%
	{\tailPrint[]}%
\entry
	{\headPrint[kana={みん\tl{˩}さまるん\tl{˥}},ipa={minsḁmaruɴ},pos={句},acc={L+H}]}%
	{%
		\MeaningsBlock{%
	\ryumeaning%
		{目覚める}{}{}%
	}%
	}%
	{\tailPrint[]}%
\entry
	{\headPrint[kana={みんすさ\tl{˩˥}はーん},ipa={minsusahaːɴ},pos={形},acc={R}]}%
	{%
		\MeaningsBlock{%
	\ryumeaning%
		{目がかたい}{夜遅くまでよく起きる}{}%
	}%
	}%
	{\tailPrint[]}%
\entry
	{\headPrint[kana={みんぜー\tl{˩˥}},ipa={mindzeː},pos={名},acc={R}]}%
	{%
		\MeaningsBlock{%
	\ryumeaning%
		{ものもらい}{}{}%
	}%
	}%
	{\tailPrint[]}%
\entry
	{\headPrint[kana={みんたま\tl{˥}},ipa={mintama},pos={名},acc={H}]}%
	{%
		\MeaningsBlock{%
	\ryumeaning%
		{目玉。眼球}{}{}%
	}%
	}%
	{\tailPrint[]}%
\entry
	{\headPrint[kana={みんだれー\tl{˩˥}},ipa={mindareː},pos={名},acc={R}]}%
	{%
		\MeaningsBlock{%
	\ryumeaning%
		{洗面器}{}{}%
	}%
	}%
	{\tailPrint[note={びん-たらいの転訛}]}%
\entry
	{\headPrint[kana={みんとーら\tl{˩˥}},ipa={mintoːra},pos={-},acc={R}]}%
	{%
		\MeaningsBlock{%
	\ryumeaning%
		{聞こえない。難聴になる}{}{}%
	}%
	}%
	{\tailPrint[]}%
\entry
	{\headPrint[kana={みんとーりむん\tl{˩˥}},ipa={mintoːrimuɴ},pos={名},acc={R}]}%
	{%
		\MeaningsBlock{%
	\ryumeaning%
		{難聴になった人}{}{}%
	}%
	}%
	{\tailPrint[]}%
\entry
	{\headPrint[kana={みんどぅー\tl{˩˥}},ipa={minduː},pos={名},acc={R}]}%
	{%
		\MeaningsBlock{%
	\ryumeaning%
		{面倒}{}{}%
	}%
	}%
	{\tailPrint[]}%
\entry
	{\headPrint[kana={みんとぅーさーん},ipa={mintuːsaːɴ},pos={形}]}%
	{%
		\MeaningsBlock{%
	\ryumeaning%
		{耳が遠い}{}{}%
	}%
	}%
	{\tailPrint[]}%
\entry
	{\headPrint[kana={みんどな\tl{˩˥}},ipa={mindona},pos={-},acc={R}]}%
	{%
		\MeaningsBlock{%
	\ryumeaning%
		{面倒だ}{}{}%
	}%
	}%
	{\tailPrint[]}%
\entry
	{\headPrint[kana={みんぬ\tl{˩˥}\ws{}かー\tl{˥}},ipa={minnu kaː},pos={句},acc={R H}]}%
	{%
		\MeaningsBlock{%
	\ryumeaning%
		{\ruby[g]{瞼}{まぶた}}{}{}%
	}%
	}%
	{\tailPrint[note={直訳「目の皮」}]}%
\entry
	{\headPrint[kana={みんぴきるん},ipa={mimpikiruɴ},pos={動}]}%
	{%
		\MeaningsBlock{%
	\ryumeaning%
		{穴があく}{}{}%
	}%
	}%
	{\tailPrint[]}%
\entry
	{\headPrint[kana={みんみり\tl{˩˥}},ipa={mimmiri},pos={句},para={みんみらぬ},acc={R}]}%
	{%
		\MeaningsBlock{%
	\ryumeaning%
		{老眼になる}{}{}%
	}%
	}%
	{\tailPrint[]}%
\LetterHead{む}
\entry
	{\headPrint[kana={むーじ\tl{˥}},ipa={muːdʒi},pos={名},acc={H}]}%
	{%
		\MeaningsBlock{%
	\ryumeaning%
		{田芋}{}{}%
	}%
	}%
	{\tailPrint[]}%
\entry
	{\headPrint[kana={むーち\tl{˥˩}},ipa={muːtʃi},pos={名},acc={F}]}%
	{%
		\MeaningsBlock{%
	\ryumeaning%
		{六つ}{}{}%
	}%
	}%
	{\tailPrint[note={簡略化\emb{むー}}]}%
\entry
	{\headPrint[kana={むーる},ipa={muːru},pos={名}]}%
	{%
		\MeaningsBlock{%
	\ryumeaning%
		{皆}{}{}%
	}%
	}%
	{\tailPrint[]}%
\entry
	{\headPrint[kana={むい\tl{˩˥}かぷん\tl{˥}},ipa={muikḁpuɴ},pos={句},para={むいかぱぬ},acc={R+H}]}%
	{%
		\MeaningsBlock{%
	\ryumeaning%
		{生い茂る}{}{}%
	}%
	}%
	{\tailPrint[]}%
\entry
	{\headPrint[kana={むいきし\tl{˥}},ipa={muiki̥ʃi},pos={句},acc={H}]}%
	{%
		\MeaningsBlock{%
	\ryumeaning%
		{思い切る}{}{}%
	}%
	}%
	{\tailPrint[]}%
\entry
	{\headPrint[kana={むいくみ\tl{˥}},ipa={muikumi},pos={動},para={むいくまぬ},acc={H}]}%
	{%
		\MeaningsBlock{%
	\ryumeaning%
		{思い込む}{}{}%
	}%
	}%
	{\tailPrint[]}%
\entry
	{\headPrint[kana={むいるん\tl{˩˥}},ipa={muiruɴ},pos={動},rui={3},para={むーぬ},acc={R}]}%
	{%
		\MeaningsBlock{%
	\ryumeaning%
		{生える}{}{}%
	}%
	}%
	{\tailPrint[]}%
\entry
	{\headPrint[kana={むえー\tl{˥˩}},ipa={mueː},pos={名},acc={F}]}%
	{%
		\MeaningsBlock{%
	\ryumeaning%
		{模合。頼母子}{}{}%
	}%
	}%
	{\tailPrint[]}%
\entry
	{\headPrint[kana={むがし\tl{˩˥}},ipa={mugaʃi},pos={名},acc={R}]}%
	{%
		\MeaningsBlock{%
	\ryumeaning%
		{\ruby[g]{昔}{むかし}}{}{}%
	}%
	}%
	{\tailPrint[]}%
\entry
	{\headPrint[kana={むがじ\tl{˥˩}},ipa={mugadʒi},pos={名},acc={F}]}%
	{%
		\MeaningsBlock{%
	\ryumeaning%
		{百足}{}{}%
	}%
	}%
	{\tailPrint[]}%
\entry
	{\headPrint[kana={むがしぃがら\tl{˩˥}},ipa={mugasu̥gara},pos={句},acc={R}]}%
	{%
		\MeaningsBlock{%
	\ryumeaning%
		{昔から。古来}{}{}%
	}%
	}%
	{\tailPrint[]}%
\entry
	{\headPrint[kana={むがしぃなれー\tl{˩˥}},ipa={mugaʃinareː},pos={名},acc={R}]}%
	{%
		\MeaningsBlock{%
	\ryumeaning%
		{旧慣。旧習}{昔からのしきたり}{}%
	}%
	}%
	{\tailPrint[]}%
\entry
	{\headPrint[kana={むがしぃ\tl{˩˥}ぱなしぃ\tl{˥}},ipa={mugaʃipḁnaʃi},pos={名},acc={R+H}]}%
	{%
		\MeaningsBlock{%
	\ryumeaning%
		{昔話}{}{}%
	}%
	}%
	{\tailPrint[]}%
\entry
	{\headPrint[kana={むがしぃ\tl{˩˥}ぴぃとぅ\tl{˥˩}},ipa={mugaʃipi̥tu},pos={名},acc={R+F}]}%
	{%
		\MeaningsBlock{%
	\ryumeaning%
		{昔の人。古人}{}{}%
	}%
	}%
	{\tailPrint[]}%
\entry
	{\headPrint[kana={むがしぃ\tl{˩˥}ゆー\tl{˥˩}},ipa={mugaʃijuː},pos={名},acc={R+F}]}%
	{%
		\MeaningsBlock{%
	\ryumeaning%
		{昔の世。神の世。豊穣の世}{}{}%
	}%
	}%
	{\tailPrint[]}%
\entry
	{\headPrint[kana={むがし\tl{˩˥}くとぅば\tl{˥}},ipa={mugaʃiku̥tuba},pos={名},acc={R+H}]}%
	{%
		\MeaningsBlock{%
	\ryumeaning%
		{諺。格言}{}{}%
	}%
	}%
	{\tailPrint[]}%
\entry
	{\headPrint[kana={むぎゃん\tl{˥˩}},ipa={mugjaɴ},pos={動（継）},para={むがーぬ},acc={F}]}%
	{%
		\MeaningsBlock{%
	\ryumeaning%
		{向かう}{}{}%
	}%
	}%
	{\tailPrint[]}%
\entry
	{\headPrint[kana={むぐ\tl{˩˥}},ipa={mugu},pos={名},acc={R}]}%
	{%
		\MeaningsBlock{%
	\ryumeaning%
		{\ruby[g]{婿}{むこ}}{}{}%
	}%
	}%
	{\tailPrint[]}%
\entry
	{\headPrint[kana={むぐちょーでー\tl{˩˥}},ipa={mugutʃoːdeː},pos={名},acc={R}]}%
	{%
		\MeaningsBlock{%
	\ryumeaning%
		{婿兄弟}{}{}%
	}%
	}%
	{\tailPrint[]}%
\entry
	{\headPrint[kana={むぐぶざ\tl{˩˥}},ipa={mugubudza},pos={名},acc={R}]}%
	{%
		\MeaningsBlock{%
	\ryumeaning%
		{\ruby[g]{婿}{むこ}}{}{}%
	}%
	}%
	{\tailPrint[]}%
\entry
	{\headPrint[kana={むぐん\tl{˩˥}},ipa={muguɴ},pos={動},rui={1},para={むがぬ},acc={R}]}%
	{%
		\MeaningsBlock{%
	\ryumeaning%
		{（皮を）剥く}{}{}%
	}%
	}%
	{\tailPrint[]}%
\entry
	{\headPrint[kana={むげーかち\tl{˥˩}},ipa={mugeːkḁtʃi},pos={名},acc={F}]}%
	{%
		\MeaningsBlock{%
	\ryumeaning%
		{向かい風}{}{}%
	}%
	}%
	{\tailPrint[]}%
\entry
	{\headPrint[kana={むげーぴん\tl{˥˩}},ipa={mugeːpiɴ},pos={名},acc={F}]}%
	{%
		\MeaningsBlock{%
	\ryumeaning%
		{迎えの日。盆の初日}{}{}%
	}%
	}%
	{\tailPrint[]}%
\entry
	{\headPrint[kana={むげーるん\tl{˥˩}},ipa={mugeːruɴ},pos={動},rui={1},para={むげーらぬ},acc={F}]}%
	{%
		\MeaningsBlock{%
	\ryumeaning%
		{迎える}{}{}%
	}%
	}%
	{\tailPrint[]}%
\entry
	{\headPrint[kana={むごん\tl{˩˥}},ipa={mugoɴ},pos={名},acc={R}]}%
	{%
		\MeaningsBlock{%
	\ryumeaning%
		{やし蟹}{}{}%
	}%
	}%
	{\tailPrint[]}%
\entry
	{\headPrint[kana={むざ\tl{˥}},ipa={mudza},pos={名},acc={H}]}%
	{%
		\MeaningsBlock{%
	\ryumeaning%
		{\ruby[g]{猪}{いのしし}}{}{}%
	}%
	}%
	{\tailPrint[]}%
\entry
	{\headPrint[kana={むさん\tl{˥˩}},ipa={musaɴ},pos={名},acc={F}]}%
	{%
		\MeaningsBlock{%
	\ryumeaning%
		{大きなうねり。大波}{}{}%
	}%
	}%
	{\tailPrint[]}%
\entry
	{\headPrint[kana={むし},ipa={muʃi},pos={接尾}]}%
	{%
		\MeaningsBlock{%
	\ryumeaning%
		{～回}{回数を示す。\emb{ぴとむしぃ}「一回」など}{}%
	}%
	}%
	{\tailPrint[]}%
\entry
	{\headPrint[kana={むし\tl{˥˩}},ipa={muʃi},pos={名},acc={F}]}%
	{%
		\MeaningsBlock{%
	\ryumeaning%
		{虫。昆虫}{}{}%
	}%
	}%
	{\tailPrint[]}%
\entry
	{\headPrint[kana={むじかーるん\tl{˩˥}},ipa={mudʒikaːruɴ},pos={動},rui={1},para={むじかーらぬ},acc={R}]}%
	{%
		\MeaningsBlock{%
	\ryumeaning%
		{捩じれる。捩じ曲がる}{}{}%
	\ryumeaning%
		{拗ねる}{}{}%
	}%
	}%
	{\tailPrint[]}%
\entry
	{\headPrint[kana={むじかるん\tl{˩˥}},ipa={mudʒikaruɴ},pos={動},rui={1},para={むじからぬ},acc={R}]}%
	{%
		\MeaningsBlock{%
	\ryumeaning%
		{捩じれる。ねじ曲がる}{}{}%
	\ryumeaning%
		{拗ねる}{}{}%
	}%
	}%
	{\tailPrint[]}%
\entry
	{\headPrint[kana={むじまーすん\tl{˩˥}},ipa={mudʒimaːsuɴ},pos={動},rui={2b},para={むじまーさぬ},acc={R}]}%
	{%
		\MeaningsBlock{%
	\ryumeaning%
		{捩じる。ねじ曲げる。捻る}{}{}%
	}%
	}%
	{\tailPrint[]}%
\entry
	{\headPrint[kana={むしゃまー\tl{˥}},ipa={muʃamaː},pos={名},acc={H}]}%
	{%
		\MeaningsBlock{%
	\ryumeaning%
		{旧盆の\emb{ゆにげー}「世願い」行事}{}{}%
	}%
	}%
	{\tailPrint[]}%
\entry
	{\headPrint[kana={むすかっさ\tl{˩˥}はん},ipa={musu̥kassahaɴ},pos={形},para={むすかさへぬ},acc={R}]}%
	{%
		\MeaningsBlock{%
	\ryumeaning%
		{難しい。困難だ}{}{}%
	}%
	}%
	{\tailPrint[]}%
\entry
	{\headPrint[kana={むすなすん\tl{˩˥}},ipa={musu̥nasuɴ},pos={動},rui={2},para={むすまはぬ},acc={R}]}%
	{%
		\MeaningsBlock{%
	\ryumeaning%
		{（牛馬を）繋ぎ変える。（牛馬を）移す}{}{}%
	}%
	}%
	{\tailPrint[]}%
\entry
	{\headPrint[kana={むすぷん\tl{˥˩}},ipa={musu̥puɴ},pos={動},rui={1},para={むすぱぬ},acc={F}]}%
	{%
		\MeaningsBlock{%
	\ryumeaning%
		{結ぶ。縛る}{}{}%
	}%
	}%
	{\tailPrint[]}%
\entry
	{\headPrint[kana={むすめー\tl{˥˩}},ipa={musu̥meː},pos={名},acc={F}]}%
	{%
		\MeaningsBlock{%
	\ryumeaning%
		{もち米。うるち米}{}{}%
	}%
	}%
	{\tailPrint[]}%
\entry
	{\headPrint[kana={むたーりるん\tl{˥˩}},ipa={mutaːriruɴ},pos={動},acc={F}]}%
	{%
		\MeaningsBlock{%
	\ryumeaning%
		{もたれる。寄り掛かる}{}{}%
	\ryumeaning%
		{頼る}{}{}%
	}%
	}%
	{\tailPrint[]}%
\entry
	{\headPrint[kana={むたすくん\tl{˩˥}},ipa={mutasu̥kuɴ},pos={動},rui={1},para={むたすかぬ},acc={R}]}%
	{%
		\MeaningsBlock{%
	\ryumeaning%
		{取り扱う}{}{}%
	}%
	}%
	{\tailPrint[]}%
\entry
	{\headPrint[kana={むだすけー\tl{˥˩}},ipa={mudasu̥keː},pos={名},acc={F}]}%
	{%
		\MeaningsBlock{%
	\ryumeaning%
		{無駄遣い}{}{}%
	}%
	}%
	{\tailPrint[]}%
\entry
	{\headPrint[kana={むたすん\tl{˩˥}},ipa={mutasuɴ},pos={動},rui={2},para={むたはぬ},acc={R}]}%
	{%
		\MeaningsBlock{%
	\ryumeaning%
		{持たす。持参させる}{}{}%
	}%
	}%
	{\tailPrint[]}%
\entry
	{\headPrint[kana={むち\tl{˥}},ipa={mutʃi},pos={名},acc={H}]}%
	{%
		\MeaningsBlock{%
	\ryumeaning%
		{顔。顔面}{}{}%
	}%
	}%
	{\tailPrint[]}%
\entry
	{\headPrint[kana={むち\tl{˥˩}},ipa={mutʃi},pos={名},acc={F}]}%
	{%
		\MeaningsBlock{%
	\ryumeaning%
		{餅}{}{}%
	}%
	}%
	{\tailPrint[]}%
\entry
	{\headPrint[kana={むち\tl{˥˩}},ipa={mutʃi},pos={名},acc={F}]}%
	{%
		\MeaningsBlock{%
	\ryumeaning%
		{\ruby[g]{漆喰}{しっくい}}{}{}%
	}%
	}%
	{\tailPrint[]}%
\entry
	{\headPrint[kana={むち\tl{˩˥}\ws{}あんぎるん\tl{˥˩}},ipa={mutʃi aŋgiruɴ},pos={句},acc={R F}]}%
	{%
		\MeaningsBlock{%
	\ryumeaning%
		{（顔を）もたげる}{}{}%
	}%
	}%
	{\tailPrint[]}%
\entry
	{\headPrint[kana={むちあんぎるん\tl{˩˥}},ipa={mutʃi.aŋgiruɴ},pos={動},para={むちあんぐぬ},acc={R}]}%
	{%
		\MeaningsBlock{%
	\ryumeaning%
		{持ち上げる}{}{}%
	}%
	}%
	{\tailPrint[]}%
\entry
	{\headPrint[kana={むちあんぎるん\tl{˩˥}},ipa={mutʃi.aŋgiruɴ},pos={動},rui={3},para={むちあんぐぬ},acc={R}]}%
	{%
		\MeaningsBlock{%
	\ryumeaning%
		{（上へ）持ち上げる}{}{}%
	\ryumeaning%
		{引き立てる}{}{}%
	}%
	}%
	{\tailPrint[]}%
\entry
	{\headPrint[kana={むちぃむぬ\tl{˩˥}},ipa={mutʃimunu},pos={名},acc={R}]}%
	{%
		\MeaningsBlock{%
	\ryumeaning%
		{持ち物}{}{}%
	}%
	}%
	{\tailPrint[]}%
\entry
	{\headPrint[kana={むちのーすん\tl{˩˥}},ipa={mutʃinoːsuɴ},pos={動},rui={2b},para={むちのーさぬ},acc={R}]}%
	{%
		\MeaningsBlock{%
	\ryumeaning%
		{持ち直す}{いい方向へ戻す}{}%
	}%
	}%
	{\tailPrint[]}%
\entry
	{\headPrint[kana={むちのーるん\tl{˩˥}},ipa={mutʃinoːruɴ},pos={動},rui={1},para={むちのーらぬ},acc={R}]}%
	{%
		\MeaningsBlock{%
	\ryumeaning%
		{持ち直る。回復する。立ち直る}{}{}%
	}%
	}%
	{\tailPrint[]}%
\entry
	{\headPrint[kana={むち\ws{}んぐん},ipa={mutʃi ŋguɴ},pos={句},acc={-}]}%
	{%
		\MeaningsBlock{%
	\ryumeaning%
		{持っていく}{}{}%
	}%
	}%
	{\tailPrint[]}%
\entry
	{\headPrint[kana={むっさ\tl{˥}はん},ipa={mussahaɴ},pos={形},para={むっさへぬ},acc={H}]}%
	{%
		\MeaningsBlock{%
	\ryumeaning%
		{面白い}{}{}%
	}%
	}%
	{\tailPrint[]}%
\entry
	{\headPrint[kana={むっしるん\tl{˥˩}},ipa={muʃʃiruɴ},pos={動},rui={3},para={むっすぬ},acc={F}]}%
	{%
		\MeaningsBlock{%
	\ryumeaning%
		{切れる。千切れる}{}{}%
	}%
	}%
	{\tailPrint[]}%
\entry
	{\headPrint[kana={むっす\tl{˩˥}},ipa={mussu},pos={名},acc={R}]}%
	{%
		\MeaningsBlock{%
	\ryumeaning%
		{\ruby[g]{蓆}{むしろ}}{}{}%
	}%
	}%
	{\tailPrint[]}%
\entry
	{\headPrint[kana={むっすん\tl{˥˩}},ipa={mussuɴ},pos={動},rui={3},para={むっさぬ},acc={F}]}%
	{%
		\MeaningsBlock{%
	\ryumeaning%
		{むしる。摘み取る}{}{}%
	}%
	}%
	{\tailPrint[]}%
\entry
	{\headPrint[kana={むっつぁーすん\tl{˩˥}},ipa={muttsaːsuɴ},pos={動},rui={1},para={むっつぁーらぬ},acc={R}]}%
	{%
		\MeaningsBlock{%
	\ryumeaning%
		{くっ付ける。ひっ付ける}{}{}%
	}%
	}%
	{\tailPrint[]}%
\entry
	{\headPrint[kana={むっつぁーるん\tl{˩˥}},ipa={muttsaːruɴ},pos={動},rui={1},para={むっつぁーらぬ},acc={R}]}%
	{%
		\MeaningsBlock{%
	\ryumeaning%
		{くっ付く。ひっ付く}{}{}%
	}%
	}%
	{\tailPrint[]}%
\entry
	{\headPrint[kana={むっつぁ\tl{˩˥}はん},ipa={muttsahaɴ},pos={形},para={むっつぁへぬ},acc={R}]}%
	{%
		\MeaningsBlock{%
	\ryumeaning%
		{粘っこい。ねばねばする}{}{}%
	}%
	}%
	{\tailPrint[]}%
\entry
	{\headPrint[kana={むっつぁるん\tl{˥˩}},ipa={muttsaruɴ},pos={動},rui={1},para={むっつぁらぬ},acc={F}]}%
	{%
		\MeaningsBlock{%
	\ryumeaning%
		{粘付く}{}{}%
	}%
	}%
	{\tailPrint[]}%
\entry
	{\headPrint[kana={むっつぁるん\tl{˩˥}},ipa={muttsaruɴ},pos={動},rui={1},para={むっつぁらぬ},acc={R}]}%
	{%
		\MeaningsBlock{%
	\ryumeaning%
		{くっ付く。ひっ付く}{}{}%
	}%
	}%
	{\tailPrint[]}%
\entry
	{\headPrint[kana={むっとぅ\tl{˩˥}},ipa={muttu},pos={副},acc={R}]}%
	{%
		\MeaningsBlock{%
	\ryumeaning%
		{全く。さっぱり}{}{}%
	}%
	}%
	{\tailPrint[]}%
\entry
	{\headPrint[kana={むっ\josho{}とぅ\kako{}む},ipa={muttumu},pos={副},acc={LHL}]}%
	{%
		\MeaningsBlock{%
	\ryumeaning%
		{もっとも}{}{}%
	}%
	}%
	{\tailPrint[]}%
\entry
	{\headPrint[kana={むつまっ\tl{˩˥}さーん},ipa={mutsumassaːɴ},pos={形},acc={R}]}%
	{%
		\MeaningsBlock{%
	\ryumeaning%
		{睦まじい。仲がいい}{}{}%
	}%
	}%
	{\tailPrint[]}%
\entry
	{\headPrint[kana={むつまっさ\tl{˩˥}はん},ipa={mutsumassahaɴ},pos={形},para={むつまっさへぬ},acc={R}]}%
	{%
		\MeaningsBlock{%
	\ryumeaning%
		{睦まじい。仲がいい}{}{}%
	}%
	}%
	{\tailPrint[]}%
\entry
	{\headPrint[kana={むつん\tl{˩˥}},ipa={mutsuɴ},pos={動},rui={2},para={むつぁぬ},acc={R}]}%
	{%
		\MeaningsBlock{%
	\ryumeaning%
		{持つ。持参する}{}{}%
	}%
	}%
	{\tailPrint[]}%
\entry
	{\headPrint[kana={むとぅ\tl{˩˥}},ipa={mutu},pos={名},acc={R}]}%
	{%
		\MeaningsBlock{%
	\ryumeaning%
		{元。幹}{}{}%
	}%
	}%
	{\tailPrint[]}%
\entry
	{\headPrint[kana={むとぅがら\tl{˩˥}},ipa={mutugara},pos={句},acc={R}]}%
	{%
		\MeaningsBlock{%
	\ryumeaning%
		{元々。元来}{}{}%
	}%
	}%
	{\tailPrint[]}%
\entry
	{\headPrint[kana={むとぅぐい\tl{˩˥}},ipa={mutugui},pos={名},acc={R}]}%
	{%
		\MeaningsBlock{%
	\ryumeaning%
		{元肥}{}{}%
	}%
	}%
	{\tailPrint[]}%
\entry
	{\headPrint[kana={むどぅすん\tl{˩˥}},ipa={mudusuɴ},pos={動},rui={2b},para={むどぅさぬ},acc={R}]}%
	{%
		\MeaningsBlock{%
	\ryumeaning%
		{戻す}{}{}%
	\ryumeaning%
		{吐く}{}{}%
	}%
	}%
	{\tailPrint[]}%
\entry
	{\headPrint[kana={むとぅだぎ\tl{˥˩}},ipa={mutudagi},pos={名},acc={F}]}%
	{%
		\MeaningsBlock{%
	\ryumeaning%
		{\ruby[g]{於}{お}\ruby[g]{茂}{も}\ruby[g]{登}{と}岳}{石垣島の山の名}{}%
	}%
	}%
	{\tailPrint[]}%
\entry
	{\headPrint[kana={むとぅみるん\tl{˩˥}},ipa={mutumiruɴ},pos={動},rui={3},para={むとぅむぬ},acc={R}]}%
	{%
		\MeaningsBlock{%
	\ryumeaning%
		{求める}{}{}%
	}%
	}%
	{\tailPrint[]}%
\entry
	{\headPrint[kana={むどぅりみち\tl{˩˥}},ipa={mudurimitʃi},pos={名},acc={R}]}%
	{%
		\MeaningsBlock{%
	\ryumeaning%
		{戻り道。帰り道}{}{}%
	}%
	}%
	{\tailPrint[]}%
\entry
	{\headPrint[kana={むどぅるん\tl{˩˥}},ipa={muduruɴ},pos={動},rui={1},para={むどぅらぬ},acc={R}]}%
	{%
		\MeaningsBlock{%
	\ryumeaning%
		{戻る}{}{}%
	}%
	}%
	{\tailPrint[]}%
\entry
	{\headPrint[kana={むに\tl{˩˥}},ipa={muni},pos={名},acc={R}]}%
	{%
		\MeaningsBlock{%
	\ryumeaning%
		{言葉。言語}{}{}%
	}%
	}%
	{\tailPrint[]}%
\entry
	{\headPrint[kana={むに\tl{˩˥}あら\tl{˥˩}はーん},ipa={muniarahaːɴ},pos={形},acc={R F}]}%
	{%
		\MeaningsBlock{%
	\ryumeaning%
		{言葉が荒い。言葉が荒っぽい}{}{}%
	}%
	}%
	{\tailPrint[]}%
\entry
	{\headPrint[kana={むに\tl{˩˥}\ws{}ゆむん\tl{˩˥}},ipa={muni jumuɴ},pos={句},acc={R R}]}%
	{%
		\MeaningsBlock{%
	\ryumeaning%
		{愚痴を言う}{}{}%
	}%
	}%
	{\tailPrint[]}%
\entry
	{\headPrint[kana={むぬ\tl{˩˥}},ipa={munu},pos={名},acc={R}]}%
	{%
		\MeaningsBlock{%
	\ryumeaning%
		{物。食べ物。食事}{}{}%
	}%
	}%
	{\tailPrint[]}%
\entry
	{\headPrint[kana={むぬうぶい\tl{˩˥}},ipa={munu.ubui},pos={名},acc={R}]}%
	{%
		\MeaningsBlock{%
	\ryumeaning%
		{物覚え。記憶力}{}{}%
	}%
	}%
	{\tailPrint[]}%
\entry
	{\headPrint[kana={むぬかんげー\tl{˩˥}},ipa={munukangeː},pos={名},acc={R}]}%
	{%
		\MeaningsBlock{%
	\ryumeaning%
		{物考え。思案}{}{}%
	}%
	}%
	{\tailPrint[]}%
\entry
	{\headPrint[kana={むぬぐとぅ\tl{˩˥}},ipa={munugutu},pos={名},acc={R}]}%
	{%
		\MeaningsBlock{%
	\ryumeaning%
		{物事}{}{}%
	}%
	}%
	{\tailPrint[]}%
\entry
	{\headPrint[kana={むぬしらし\tl{˩˥}},ipa={munuʃiraʃi},pos={名},acc={R}]}%
	{%
		\MeaningsBlock{%
	\ryumeaning%
		{予兆。神仏の啓示}{}{}%
	}%
	}%
	{\tailPrint[]}%
\entry
	{\headPrint[kana={むぬしり\tl{˩˥}},ipa={munuʃiri},pos={名},acc={R}]}%
	{%
		\MeaningsBlock{%
	\ryumeaning%
		{物知り}{}{}%
	}%
	}%
	{\tailPrint[]}%
\entry
	{\headPrint[kana={むぬすくり\tl{˩˥}},ipa={munusu̥kuri},pos={名},acc={R}]}%
	{%
		\MeaningsBlock{%
	\ryumeaning%
		{農耕。栽培}{「物作り」の義}{}%
	}%
	}%
	{\tailPrint[]}%
\entry
	{\headPrint[kana={むぬ\tl{˩˥}すくん\tl{˥˩}},ipa={munusu̥kuɴ},pos={句},para={むぬすかぬ},acc={R+F}]}%
	{%
		\MeaningsBlock{%
	\ryumeaning%
		{尋ねる。物事を聞く}{}{}%
	}%
	}%
	{\tailPrint[]}%
\entry
	{\headPrint[kana={むぬなれー\tl{˩˥}},ipa={mununareː},pos={名},acc={R}]}%
	{%
		\MeaningsBlock{%
	\ryumeaning%
		{物習い。学習。勉強}{}{}%
	}%
	}%
	{\tailPrint[]}%
\entry
	{\headPrint[kana={むぬばっさ},ipa={munubassa},pos={句}]}%
	{%
		\MeaningsBlock{%
	\ryumeaning%
		{物忘れ}{}{}%
	}%
	}%
	{\tailPrint[]}%
\entry
	{\headPrint[kana={むぬぱなし\tl{˩˥}},ipa={munupanaʃi},pos={名},acc={R}]}%
	{%
		\MeaningsBlock{%
	\ryumeaning%
		{おしゃべり。世間話}{}{}%
	}%
	}%
	{\tailPrint[]}%
\entry
	{\headPrint[kana={むぬへーずく\tl{˩˥}},ipa={munuheːdzuku},pos={名},acc={R}]}%
	{%
		\MeaningsBlock{%
	\ryumeaning%
		{\ruby[g]{生}{なり}\ruby[g]{業}{わい}}{}{}%
	}%
	}%
	{\tailPrint[]}%
\entry
	{\headPrint[kana={むぬみずらはん},ipa={munumidzurahaɴ},pos={形},para={むぬみじらへぬ},acc={-}]}%
	{%
		\MeaningsBlock{%
	\ryumeaning%
		{珍しい。物珍しい}{}{}%
	}%
	}%
	{\tailPrint[]}%
\entry
	{\headPrint[kana={むぬむい\tl{˩˥}},ipa={munumui},pos={名},acc={R}]}%
	{%
		\MeaningsBlock{%
	\ryumeaning%
		{物思い。心配}{}{}%
	}%
	}%
	{\tailPrint[]}%
\entry
	{\headPrint[kana={むぬむち\tl{˩˥}},ipa={munumutʃi},pos={名},acc={R}]}%
	{%
		\MeaningsBlock{%
	\ryumeaning%
		{物持ち。資産家}{}{}%
	}%
	}%
	{\tailPrint[]}%
\entry
	{\headPrint[kana={むみん\tl{˥}},ipa={mumiɴ},pos={名},acc={H}]}%
	{%
		\MeaningsBlock{%
	\ryumeaning%
		{木綿}{}{}%
	}%
	}%
	{\tailPrint[]}%
\entry
	{\headPrint[kana={むむん\tl{˥˩}},ipa={mumuɴ},pos={動},rui={1},para={むまぬ},acc={F}]}%
	{%
		\MeaningsBlock{%
	\ryumeaning%
		{揉む。揉みほぐす}{}{}%
	}%
	}%
	{\tailPrint[]}%
\entry
	{\headPrint[kana={むやすん\tl{˩˥}},ipa={mujasuɴ},pos={動},rui={2},para={もーはぬ},acc={R}]}%
	{%
		\MeaningsBlock{%
	\ryumeaning%
		{（草を）生やす}{}{}%
	}%
	}%
	{\tailPrint[]}%
\entry
	{\headPrint[kana={むよー\tl{˥˩}},ipa={mujoː},pos={名},acc={F}]}%
	{%
		\MeaningsBlock{%
	\ryumeaning%
		{模様。様子}{}{}%
	}%
	}%
	{\tailPrint[]}%
\entry
	{\headPrint[kana={むら\tl{˥˩}},ipa={mura},pos={名},acc={F}]}%
	{%
		\MeaningsBlock{%
	\ryumeaning%
		{村。村落}{}{}%
	}%
	}%
	{\tailPrint[]}%
\entry
	{\headPrint[kana={むらぐとぅ\tl{˥˩}},ipa={muragutu},pos={名},acc={F}]}%
	{%
		\MeaningsBlock{%
	\ryumeaning%
		{村の行事}{}{}%
	}%
	}%
	{\tailPrint[]}%
\entry
	{\headPrint[kana={むらたたまり},ipa={muratḁtamari},pos={名}]}%
	{%
		\MeaningsBlock{%
	\ryumeaning%
		{村人の集会}{王府時代の村人集会}{}%
	}%
	}%
	{\tailPrint[]}%
\entry
	{\headPrint[kana={むらばぎん\tl{˩˥}},ipa={murabagiɴ},pos={名},acc={R}]}%
	{%
		\MeaningsBlock{%
	\ryumeaning%
		{分村}{王府時代に行われた強制的な分村}{}%
	}%
	}%
	{\tailPrint[]}%
\entry
	{\headPrint[kana={むり\tl{˥˩}},ipa={muri},pos={名},acc={F}]}%
	{%
		\MeaningsBlock{%
	\ryumeaning%
		{山。森。丘。石盛}{}{}%
	}%
	}%
	{\tailPrint[]}%
\entry
	{\headPrint[kana={むりあがるん\tl{˥˩}},ipa={muri.agaruɴ},pos={動},rui={1},para={むりあがらぬ},acc={F}]}%
	{%
		\MeaningsBlock{%
	\ryumeaning%
		{盛り上がる。膨れ上がる。持ち上がる}{}{}%
	}%
	}%
	{\tailPrint[]}%
\entry
	{\headPrint[kana={むりあぎるん\tl{˥˩}},ipa={muri.agiruɴ},pos={動},rui={3},para={むりあぐぬ},acc={F}]}%
	{%
		\MeaningsBlock{%
	\ryumeaning%
		{盛り上げる}{}{}%
	}%
	}%
	{\tailPrint[]}%
\entry
	{\headPrint[kana={むりあま\tl{˩˥}},ipa={muri.ama},pos={名},acc={R}]}%
	{%
		\MeaningsBlock{%
	\ryumeaning%
		{子守の娘}{(姉)}{}%
	}%
	}%
	{\tailPrint[]}%
\entry
	{\headPrint[kana={む\josho{}る\kako{}むるし},ipa={murumuruʃi},pos={副},acc={重複}]}%
	{%
		\MeaningsBlock{%
	\ryumeaning%
		{丸い。丸っこい}{}{}%
	}%
	}%
	{\tailPrint[]}%
\entry
	{\headPrint[kana={むるん},ipa={muruɴ},pos={動},rui={1},para={むらぬ},acc={-}]}%
	{%
		\MeaningsBlock{%
	\ryumeaning%
		{盛る。盛り上げる}{}{}%
	}%
	}%
	{\tailPrint[]}%
\entry
	{\headPrint[kana={むるん\tl{˥}},ipa={muruɴ},pos={名},acc={H}]}%
	{%
		\MeaningsBlock{%
	\ryumeaning%
		{もろみ}{}{}%
	}%
	}%
	{\tailPrint[]}%
\entry
	{\headPrint[kana={むん\tl{˥}},ipa={muɴ},pos={動},rui={1},para={もわぬ},acc={H}]}%
	{%
		\MeaningsBlock{%
	\ryumeaning%
		{思う。考える}{}{}%
	}%
	}%
	{\tailPrint[]}%
\entry
	{\headPrint[kana={むん\tl{˩˥}},ipa={muɴ},pos={名},acc={R}]}%
	{%
		\MeaningsBlock{%
	\ryumeaning%
		{麦}{麦類の総称}{}%
	}%
	}%
	{\tailPrint[]}%
\entry
	{\headPrint[kana={むんぐる\tl{˩˥}},ipa={muŋguru},pos={名},acc={R}]}%
	{%
		\MeaningsBlock{%
	\ryumeaning%
		{麦わら}{}{}%
	}%
	}%
	{\tailPrint[]}%
\entry
	{\headPrint[kana={むんじる\tl{˩˥}},ipa={mundʒiru},pos={動},acc={R}]}%
	{%
		\MeaningsBlock{%
	\ryumeaning%
		{思い出す}{}{}%
	}%
	}%
	{\tailPrint[]}%
\entry
	{\headPrint[kana={むんだに\tl{˩˥}},ipa={mundani},pos={名},acc={R}]}%
	{%
		\MeaningsBlock{%
	\ryumeaning%
		{餌。釣り餌}{}{}%
	}%
	}%
	{\tailPrint[]}%
\entry
	{\headPrint[kana={むんだらし\tl{˩˥}},ipa={mundaraʃi},pos={名},acc={R}]}%
	{%
		\MeaningsBlock{%
	\ryumeaning%
		{腿。太腿}{}{}%
	}%
	}%
	{\tailPrint[]}%
\entry
	{\headPrint[kana={むんどぅー\tl{˩˥}},ipa={munduː},pos={名},acc={R}]}%
	{%
		\MeaningsBlock{%
	\ryumeaning%
		{問答。争議。口論}{}{}%
	}%
	}%
	{\tailPrint[]}%
\entry
	{\headPrint[kana={むんどぅー\ws{}すん},ipa={munduː suɴ},pos={名}]}%
	{%
		\MeaningsBlock{%
	\ryumeaning%
		{争議する。口論する}{}{}%
	}%
	}%
	{\tailPrint[]}%
\entry
	{\headPrint[kana={むんぬ\tl{˩˥}\ws{}くー\tl{˥}},ipa={munnu kuː},pos={句},acc={R H}]}%
	{%
		\MeaningsBlock{%
	\ryumeaning%
		{小麦粉。メリケン粉}{}{}%
	}%
	}%
	{\tailPrint[]}%
\entry
	{\headPrint[kana={むんみ\tl{˥}},ipa={munmi},pos={名},acc={H}]}%
	{%
		\MeaningsBlock{%
	\ryumeaning%
		{\ruby[g]{匁}{もんめ}}{重さの単位}{}%
	}%
	}%
	{\tailPrint[]}%
\LetterHead{め}
\entry
	{\headPrint[kana={めー\tl{˥˩}},ipa={meː},pos={名},acc={F}]}%
	{%
		\MeaningsBlock{%
	\ryumeaning%
		{米}{作物の稲にも言う}{}%
	}%
	}%
	{\tailPrint[]}%
\entry
	{\headPrint[kana={めー\tl{˩˥}},ipa={meː},pos={名},acc={R}]}%
	{%
		\MeaningsBlock{%
	\ryumeaning%
		{前。前方。以前}{}{}%
	}%
	}%
	{\tailPrint[]}%
\entry
	{\headPrint[kana={めー\josho{}が\kako{}めーにち},ipa={meːgameːnitʃi},pos={副},acc={LHL}]}%
	{%
		\MeaningsBlock{%
	\ryumeaning%
		{毎日。何時も}{}{}%
	}%
	}%
	{\tailPrint[]}%
\entry
	{\headPrint[kana={めーかり\tl{˥˩}},ipa={meːkḁri},pos={名},acc={F}]}%
	{%
		\MeaningsBlock{%
	\ryumeaning%
		{稲刈り}{}{}%
	}%
	}%
	{\tailPrint[]}%
\entry
	{\headPrint[kana={めーがり\tl{˩˥}},ipa={meːgari},pos={名},acc={R}]}%
	{%
		\MeaningsBlock{%
	\ryumeaning%
		{前借}{}{}%
	}%
	}%
	{\tailPrint[]}%
\entry
	{\headPrint[kana={めー\tl{˩˥}しぴ\tl{˥˩}},ipa={meːʃipi},pos={名},acc={R+F}]}%
	{%
		\MeaningsBlock{%
	\ryumeaning%
		{前後。逆}{}{}%
	}%
	}%
	{\tailPrint[]}%
\entry
	{\headPrint[kana={めーすん\tl{˥˩}},ipa={meːsuɴ},pos={動},rui={2},para={めーはぬ},acc={F}]}%
	{%
		\MeaningsBlock{%
	\ryumeaning%
		{燃やす}{}{}%
	}%
	}%
	{\tailPrint[]}%
\entry
	{\headPrint[kana={めーだら\tl{˩˥}},ipa={meːdara},pos={名},acc={R}]}%
	{%
		\MeaningsBlock{%
	\ryumeaning%
		{米俵}{}{}%
	}%
	}%
	{\tailPrint[]}%
\entry
	{\headPrint[kana={めーどぅし\tl{˩˥}},ipa={meːduʃi},pos={名},acc={R}]}%
	{%
		\MeaningsBlock{%
	\ryumeaning%
		{毎年}{}{}%
	}%
	}%
	{\tailPrint[]}%
\entry
	{\headPrint[kana={めー\tl{˩˥}\ws{}なすん\tl{˩˥}},ipa={meː nasuɴ},pos={句},acc={R R}]}%
	{%
		\MeaningsBlock{%
	\ryumeaning%
		{前にする}{}{}%
	\ryumeaning%
		{先にする}{}{}%
	\ryumeaning%
		{前に置く}{}{}%
	}%
	}%
	{\tailPrint[]}%
\entry
	{\headPrint[kana={めー\tl{˩˥}\ws{}なるん\tl{˩˥}},ipa={meː naruɴ},pos={句},acc={R R}]}%
	{%
		\MeaningsBlock{%
	\ryumeaning%
		{前を行く。先になる}{}{}%
	}%
	}%
	{\tailPrint[]}%
\entry
	{\headPrint[kana={めーにち\tl{˩˥}},ipa={meːnitʃi},pos={名},acc={R}]}%
	{%
		\MeaningsBlock{%
	\ryumeaning%
		{毎日}{}{}%
	}%
	}%
	{\tailPrint[]}%
\entry
	{\headPrint[kana={めーぬ\tl{˥˩}\ws{}いー\tl{˥}},ipa={meːnu iː},pos={句},acc={F H}]}%
	{%
		\MeaningsBlock{%
	\ryumeaning%
		{米の飯。米飯}{}{}%
	}%
	}%
	{\tailPrint[]}%
\entry
	{\headPrint[kana={めーぬ\tl{˩˥}\ws{}しき\tl{˥}},ipa={meːnu ʃi̥ki},pos={句},acc={R H}]}%
	{%
		\MeaningsBlock{%
	\ryumeaning%
		{先月}{直訳「前の月」}{}%
	}%
	}%
	{\tailPrint[]}%
\entry
	{\headPrint[kana={めー\tl{˩˥}ぬすかるん\tl{˥˩}},ipa={meːnusu̥karuɴ},pos={句},acc={R+F}]}%
	{%
		\MeaningsBlock{%
	\ryumeaning%
		{間近に迫る。近づく}{}{}%
	}%
	}%
	{\tailPrint[]}%
\entry
	{\headPrint[kana={めーぱん\tl{˩˥}},ipa={meːpaɴ},pos={名},acc={R}]}%
	{%
		\MeaningsBlock{%
	\ryumeaning%
		{前足}{}{}%
	}%
	}%
	{\tailPrint[]}%
\entry
	{\headPrint[kana={めーぱん\tl{˩˥}},ipa={meːpaɴ},pos={名}]}%
	{%
		\MeaningsBlock{%
	\ryumeaning%
		{前歯}{}{}%
	}%
	}%
	{\tailPrint[]}%
\entry
	{\headPrint[kana={めーむち},ipa={meːmutʃi},pos={副}]}%
	{%
		\MeaningsBlock{%
	\ryumeaning%
		{前もって}{}{}%
	}%
	}%
	{\tailPrint[]}%
\entry
	{\headPrint[kana={めーむぬ\tl{˥˩}},ipa={meːmunu},pos={名},acc={F}]}%
	{%
		\MeaningsBlock{%
	\ryumeaning%
		{新しい物。新品}{}{}%
	}%
	}%
	{\tailPrint[]}%
\entry
	{\headPrint[kana={めーむるし\tl{˩˥}},ipa={meːmuruʃi},pos={名},acc={R}]}%
	{%
		\MeaningsBlock{%
	\ryumeaning%
		{古墳}{古墓}{}%
	}%
	}%
	{\tailPrint[]}%
\entry
	{\headPrint[kana={めーら\tl{˥˩}},ipa={meːra},pos={名},acc={F}]}%
	{%
		\MeaningsBlock{%
	\ryumeaning%
		{\ruby[g]{宮良}{みやら}}{石垣島の集落名}{}%
	}%
	}%
	{\tailPrint[]}%
\entry
	{\headPrint[kana={めーりぴてー\tl{˩˥}},ipa={meːripi̥teː},pos={名},acc={R}]}%
	{%
		\MeaningsBlock{%
	\ryumeaning%
		{肥沃な畑}{}{}%
	}%
	}%
	{\tailPrint[]}%
\entry
	{\headPrint[kana={めーるん\tl{˥˩}},ipa={meːruɴ},pos={動},rui={1},para={めーらぬ},acc={F}]}%
	{%
		\MeaningsBlock{%
	\ryumeaning%
		{（芽などが）太る。実る}{}{}%
	}%
	}%
	{\tailPrint[]}%
\entry
	{\headPrint[kana={めーるん\tl{˥˩}},ipa={meːruɴ},pos={動},rui={1},para={めーらぬ},acc={F}]}%
	{%
		\MeaningsBlock{%
	\ryumeaning%
		{燃える}{}{}%
	}%
	}%
	{\tailPrint[]}%
\entry
	{\headPrint[kana={めすかるん\tl{˩˥}},ipa={mesu̥karuɴ},pos={句},para={めすからぬ},acc={R}]}%
	{%
		\MeaningsBlock{%
	\ryumeaning%
		{見つかる}{}{}%
	}%
	}%
	{\tailPrint[]}%
\entry
	{\headPrint[kana={めっさ\tl{˩˥}はん},ipa={messahaɴ},pos={形},para={めっさへぬ},acc={R}]}%
	{%
		\MeaningsBlock{%
	\ryumeaning%
		{心地よい。気持ちいい}{}{}%
	}%
	}%
	{\tailPrint[]}%
\entry
	{\headPrint[kana={めった},ipa={metta},pos={名}]}%
	{%
		\MeaningsBlock{%
	\ryumeaning%
		{アンツク。入れ物の一種}{縄で編んだ物入れ}{}%
	}%
	}%
	{\tailPrint[]}%
\entry
	{\headPrint[kana={めっつぁ\tl{˩˥}},ipa={mettsa},pos={名},acc={R}]}%
	{%
		\MeaningsBlock{%
	\ryumeaning%
		{女のふんどし}{}{}%
	}%
	}%
	{\tailPrint[]}%
\entry
	{\headPrint[kana={めふな\tl{˥˩}},ipa={mefu̥na},pos={名},acc={F}]}%
	{%
		\MeaningsBlock{%
	\ryumeaning%
		{利口な}{}{}%
	}%
	}%
	{\tailPrint[]}%
\entry
	{\headPrint[kana={めんさ\tl{˩˥}},ipa={mensa},pos={名},acc={R}]}%
	{%
		\MeaningsBlock{%
	\ryumeaning%
		{\ruby[g]{蓑}{みの}}{}{}%
	}%
	}%
	{\tailPrint[]}%
\LetterHead{も}
\entry
	{\headPrint[kana={もー\tl{˥}},ipa={moː},pos={名},acc={H}]}%
	{%
		\MeaningsBlock{%
	\ryumeaning%
		{ここ。この辺}{}{}%
	}%
	}%
	{\tailPrint[]}%
\entry
	{\headPrint[kana={もーが\tl{˥}},ipa={moːga},pos={名},acc={H}]}%
	{%
		\MeaningsBlock{%
	\ryumeaning%
		{儲け。利益}{}{}%
	}%
	}%
	{\tailPrint[]}%
\entry
	{\headPrint[kana={もーぎ\tl{˩˥}},ipa={moːgi},pos={名},acc={R}]}%
	{%
		\MeaningsBlock{%
	\ryumeaning%
		{儲け。利益}{}{}%
	}%
	}%
	{\tailPrint[]}%
\entry
	{\headPrint[kana={もーや\tl{˩˥}},ipa={moːja},pos={名},acc={R}]}%
	{%
		\MeaningsBlock{%
	\ryumeaning%
		{即興踊り}{}{}%
	}%
	}%
	{\tailPrint[]}%
\entry
	{\headPrint[kana={もーら\tl{˥}},ipa={moːra},pos={句},acc={H}]}%
	{%
		\MeaningsBlock{%
	\ryumeaning%
		{こちらから}{}{}%
	}%
	}%
	{\tailPrint[]}%
\entry
	{\headPrint[kana={もーらすん\tl{˩˥}},ipa={moːrasuɴ},pos={動},rui={2},para={もーらはぬ},acc={R}]}%
	{%
		\MeaningsBlock{%
	\ryumeaning%
		{漏らす}{}{}%
	}%
	}%
	{\tailPrint[]}%
\entry
	{\headPrint[kana={もーるん\tl{˩˥}},ipa={moːruɴ},pos={動},rui={1},para={もーらぬ},acc={R}]}%
	{%
		\MeaningsBlock{%
	\ryumeaning%
		{漏る。漏れる}{}{}%
	}%
	}%
	{\tailPrint[]}%
\entry
	{\headPrint[kana={もが\tl{˥}},ipa={moga},pos={句},acc={H}]}%
	{%
		\MeaningsBlock{%
	\ryumeaning%
		{こちらへ}{}{}%
	}%
	}%
	{\tailPrint[]}%
\entry
	{\headPrint[kana={もぬ\tl{˥}\ws{}まーり\tl{˥˩}},ipa={monu maːri},pos={句},acc={H F}]}%
	{%
		\MeaningsBlock{%
	\ryumeaning%
		{この辺。この付近}{}{}%
	}%
	}%
	{\tailPrint[]}%
\entry
	{\headPrint[kana={もみ\tl{˥˩}},ipa={momi},pos={名},acc={F}]}%
	{%
		\MeaningsBlock{%
	\ryumeaning%
		{\ruby[g]{籾}{もみ}}{}{}%
	}%
	}%
	{\tailPrint[]}%
\LetterHead{や}
\entry
	{\headPrint[kana={や\tl{˥˩}},ipa={ja},pos={名},acc={F}]}%
	{%
		\MeaningsBlock{%
	\ryumeaning%
		{矢。\ruby[g]{楔}{くさび}}{}{}%
	}%
	}%
	{\tailPrint[]}%
\entry
	{\headPrint[kana={やー},ipa={jaː},pos={感}]}%
	{%
		\MeaningsBlock{%
	\ryumeaning%
		{やぁー。おーい}{呼びかけの言葉}{}%
	}%
	}%
	{\tailPrint[]}%
\entry
	{\headPrint[kana={やーち\tl{˥}},ipa={jaːtʃi},pos={名},acc={H}]}%
	{%
		\MeaningsBlock{%
	\ryumeaning%
		{八つ}{}{}%
	}%
	}%
	{\tailPrint[note={簡略化\emb{やー}}]}%
\entry
	{\headPrint[kana={やーなれー},ipa={jaːnareː},pos={名}]}%
	{%
		\MeaningsBlock{%
	\ryumeaning%
		{家の習慣。家風}{}{}%
	}%
	}%
	{\tailPrint[]}%
\entry
	{\headPrint[kana={やーにんじゅ\tl{˩˥}},ipa={jaːnindʒu},pos={名},acc={R}]}%
	{%
		\MeaningsBlock{%
	\ryumeaning%
		{家族}{}{}%
	}%
	}%
	{\tailPrint[]}%
\entry
	{\headPrint[kana={やーは\tl{˩˥}\ws{}すん\tl{˥˩}},ipa={jaːha suɴ},pos={句},acc={R F}]}%
	{%
		\MeaningsBlock{%
	\ryumeaning%
		{ひもじい。空腹だ}{}{}%
	}%
	}%
	{\tailPrint[]}%
\entry
	{\headPrint[kana={やー\tl{˩˥}はん},ipa={jaːhaɴ},pos={形},para={やーへぬ},acc={R}]}%
	{%
		\MeaningsBlock{%
	\ryumeaning%
		{ひもじい。空腹だ}{}{}%
	}%
	}%
	{\tailPrint[]}%
\entry
	{\headPrint[kana={やーばん\tl{˩˥}},ipa={jaːbaɴ},pos={名},acc={R}]}%
	{%
		\MeaningsBlock{%
	\ryumeaning%
		{家紋。家判}{}{}%
	}%
	}%
	{\tailPrint[]}%
\entry
	{\headPrint[kana={やーぶさ\tl{˩˥}},ipa={jaːbusa},pos={名},acc={R}]}%
	{%
		\MeaningsBlock{%
	\ryumeaning%
		{山補佐}{王府時代の百姓役目。杣山筆者の補佐役}{}%
	}%
	}%
	{\tailPrint[]}%
\entry
	{\headPrint[kana={やーむとぅ\tl{˩˥}},ipa={jaːmutu},pos={名},acc={R}]}%
	{%
		\MeaningsBlock{%
	\ryumeaning%
		{家元。本家。実家}{}{}%
	}%
	}%
	{\tailPrint[]}%
\entry
	{\headPrint[kana={やーむんどぅ},ipa={jaːmundu},pos={名}]}%
	{%
		\MeaningsBlock{%
	\ryumeaning%
		{家庭内争議}{}{}%
	}%
	}%
	{\tailPrint[]}%
\entry
	{\headPrint[kana={やーらぎるん\tl{˥˩}},ipa={jaːragiruɴ},pos={動},rui={3},para={やーらぐぬ},acc={F}]}%
	{%
		\MeaningsBlock{%
	\ryumeaning%
		{和らげる。柔らかくする}{}{}%
	}%
	}%
	{\tailPrint[]}%
\entry
	{\headPrint[kana={やーらぐん\tl{˥˩}},ipa={jaːraguɴ},pos={動},rui={1},para={やーらがぬ},acc={F}]}%
	{%
		\MeaningsBlock{%
	\ryumeaning%
		{柔らぐ。柔らかになる}{}{}%
	}%
	}%
	{\tailPrint[]}%
\entry
	{\headPrint[kana={やーりん\tl{˩˥}},ipa={jaːriɴ},pos={動},rui={3},para={やるぬ},acc={R}]}%
	{%
		\MeaningsBlock{%
	\ryumeaning%
		{破れる。裂ける}{}{}%
	}%
	}%
	{\tailPrint[]}%
\entry
	{\headPrint[kana={やいっと},ipa={jaitto},pos={擬}]}%
	{%
		\MeaningsBlock{%
	\ryumeaning%
		{がばと。すっくと}{力を入れて勢いよく行動を開始するさま}{}%
	}%
	}%
	{\tailPrint[]}%
\entry
	{\headPrint[kana={やいま\tl{˥˩}},ipa={jaima},pos={名},acc={F}]}%
	{%
		\MeaningsBlock{%
	\ryumeaning%
		{八重山}{地名}{}%
	}%
	}%
	{\tailPrint[]}%
\entry
	{\headPrint[kana={やいやいた},ipa={jaijaita},pos={擬}]}%
	{%
		\MeaningsBlock{%
	\ryumeaning%
		{せっせと働くさま}{}{}%
	}%
	}%
	{\tailPrint[]}%
\entry
	{\headPrint[kana={やがまっさ\tl{˩˥}はーん},ipa={jagamassahaːɴ},pos={形},acc={F}]}%
	{%
		\MeaningsBlock{%
	\ryumeaning%
		{やかましい}{}{}%
	}%
	}%
	{\tailPrint[]}%
\entry
	{\headPrint[kana={やがりるん\tl{˥˩}},ipa={jagariruɴ},pos={動},rui={3},para={やがるぬ},acc={F}]}%
	{%
		\MeaningsBlock{%
	\ryumeaning%
		{焼ける}{}{}%
	}%
	}%
	{\tailPrint[]}%
\entry
	{\headPrint[kana={やきー\tl{˩˥}},ipa={jakiː},pos={名},acc={R}]}%
	{%
		\MeaningsBlock{%
	\ryumeaning%
		{マラリア。熱病}{\emb{やえやま-やきー}と呼ばれた}{}%
	}%
	}%
	{\tailPrint[]}%
\entry
	{\headPrint[kana={やぎしっしるん},ipa={jagissiruɴ},pos={動}]}%
	{%
		\MeaningsBlock{%
	\ryumeaning%
		{焼き捨てる}{}{}%
	}%
	}%
	{\tailPrint[]}%
\entry
	{\headPrint[kana={やぎむぬ\tl{˩˥}},ipa={jagimunu},pos={名},acc={R}]}%
	{%
		\MeaningsBlock{%
	\ryumeaning%
		{焼き物。陶磁器}{}{}%
	}%
	}%
	{\tailPrint[]}%
\entry
	{\headPrint[kana={やく\tl{˩˥}},ipa={jaku},pos={名},acc={R}]}%
	{%
		\MeaningsBlock{%
	\ryumeaning%
		{役職}{}{}%
	}%
	}%
	{\tailPrint[]}%
\entry
	{\headPrint[kana={やく\tl{˩˥}},ipa={jaku},pos={名},acc={R}]}%
	{%
		\MeaningsBlock{%
	\ryumeaning%
		{厄。災い}{}{}%
	}%
	}%
	{\tailPrint[]}%
\entry
	{\headPrint[kana={やぐい\tl{˩˥}},ipa={jagui},pos={名},acc={R}]}%
	{%
		\MeaningsBlock{%
	\ryumeaning%
		{掛声}{}{}%
	}%
	}%
	{\tailPrint[]}%
\entry
	{\headPrint[kana={やく\tl{˩˥}たつん\tl{˥}},ipa={jakutḁtsuɴ},pos={動},para={やくたたぬ},acc={R+H}]}%
	{%
		\MeaningsBlock{%
	\ryumeaning%
		{役に立つ}{}{}%
	}%
	}%
	{\tailPrint[]}%
\entry
	{\headPrint[kana={やくどぅし\tl{˥˩}},ipa={jakuduʃi},pos={名},acc={F}]}%
	{%
		\MeaningsBlock{%
	\ryumeaning%
		{厄年}{}{}%
	}%
	}%
	{\tailPrint[]}%
\entry
	{\headPrint[kana={やくば\tl{˩˥}},ipa={jakuba},pos={名},acc={R}]}%
	{%
		\MeaningsBlock{%
	\ryumeaning%
		{役場。役所}{}{}%
	}%
	}%
	{\tailPrint[]}%
\entry
	{\headPrint[kana={やぐん\tl{˥˩}},ipa={jaguɴ},pos={動},rui={1},para={やがぬ},acc={F}]}%
	{%
		\MeaningsBlock{%
	\ryumeaning%
		{焼く。焼き捨てる}{}{}%
	}%
	}%
	{\tailPrint[]}%
\entry
	{\headPrint[kana={やし\tl{˩˥}},ipa={jaʃi},pos={名},acc={R}]}%
	{%
		\MeaningsBlock{%
	\ryumeaning%
		{\ruby[g]{椰子}{やし}}{植物名}{}%
	}%
	}%
	{\tailPrint[]}%
\entry
	{\headPrint[kana={やしぃねーふぁ},ipa={jaʃi̥neːfa},pos={名}]}%
	{%
		\MeaningsBlock{%
	\ryumeaning%
		{養子。養い子}{}{}%
	}%
	}%
	{\tailPrint[]}%
\entry
	{\headPrint[kana={やしき\tl{˩˥}},ipa={jaʃi̥ki},pos={名},acc={R}]}%
	{%
		\MeaningsBlock{%
	\ryumeaning%
		{屋敷}{}{}%
	}%
	}%
	{\tailPrint[]}%
\entry
	{\headPrint[kana={やしねー},ipa={jaʃi̥neː},pos={名}]}%
	{%
		\MeaningsBlock{%
	\ryumeaning%
		{養い}{}{}%
	}%
	}%
	{\tailPrint[]}%
\entry
	{\headPrint[kana={やしねーうや},ipa={jaʃi̥neː.uja},pos={名}]}%
	{%
		\MeaningsBlock{%
	\ryumeaning%
		{養い親。養父母}{}{}%
	}%
	}%
	{\tailPrint[]}%
\entry
	{\headPrint[kana={やじまるん\tl{˥˩}},ipa={jadʒimaruɴ},pos={動},rui={1},para={やじまらぬ},acc={F}]}%
	{%
		\MeaningsBlock{%
	\ryumeaning%
		{孕まない。不妊だ}{}{}%
	}%
	}%
	{\tailPrint[]}%
\entry
	{\headPrint[kana={やしみ\tl{˩˥}},ipa={jaʃi̥mi},pos={名},acc={R}]}%
	{%
		\MeaningsBlock{%
	\ryumeaning%
		{休み。休憩}{}{}%
	}%
	}%
	{\tailPrint[]}%
\entry
	{\headPrint[kana={やしむぬ\tl{˩˥}},ipa={jaʃi̥munu},pos={名},acc={R}]}%
	{%
		\MeaningsBlock{%
	\ryumeaning%
		{安物}{}{}%
	}%
	}%
	{\tailPrint[]}%
\entry
	{\headPrint[kana={やしゃ\tl{˩˥}},ipa={jaʃa},pos={名},acc={R}]}%
	{%
		\MeaningsBlock{%
	\ryumeaning%
		{次姉}{}{}%
	}%
	}%
	{\tailPrint[]}%
\entry
	{\headPrint[kana={やしり\tl{˩˥}},ipa={jaʃi̥ri},pos={名},acc={R}]}%
	{%
		\MeaningsBlock{%
	\ryumeaning%
		{\ruby[g]{鑢}{やすり}}{}{}%
	}%
	}%
	{\tailPrint[]}%
\entry
	{\headPrint[kana={やすぅむん\tl{˩˥}},ipa={jasu̥muɴ},pos={動},acc={R}]}%
	{%
		\MeaningsBlock{%
	\ryumeaning%
		{休む}{休憩の意ではなく、休暇を取るの意}{}%
	}%
	}%
	{\tailPrint[]}%
\entry
	{\headPrint[kana={やすなすん\tl{˥˩}},ipa={jasu̥nasuɴ},pos={動},rui={2},para={やすなはぬ},acc={F}]}%
	{%
		\MeaningsBlock{%
	\ryumeaning%
		{養う。養育する}{}{}%
	}%
	}%
	{\tailPrint[]}%
\entry
	{\headPrint[kana={やすんじるん\tl{˩˥}},ipa={jasundʒiruɴ},pos={動},acc={R}]}%
	{%
		\MeaningsBlock{%
	\ryumeaning%
		{安んじる。安心する}{}{}%
	}%
	}%
	{\tailPrint[]}%
\entry
	{\headPrint[kana={やせー\tl{˩˥}},ipa={jaseː},pos={名},acc={R}]}%
	{%
		\MeaningsBlock{%
	\ryumeaning%
		{野菜}{}{}%
	}%
	}%
	{\tailPrint[]}%
\entry
	{\headPrint[kana={やせーぴてー},ipa={jaseːpi̥teː},pos={名},acc={-}]}%
	{%
		\MeaningsBlock{%
	\ryumeaning%
		{野菜畑}{}{}%
	}%
	}%
	{\tailPrint[]}%
\entry
	{\headPrint[kana={やた\tl{˩˥}},ipa={jata},pos={名},acc={R}]}%
	{%
		\MeaningsBlock{%
	\ryumeaning%
		{\ruby[g]{脇}{わき}}{}{}%
	}%
	}%
	{\tailPrint[]}%
\entry
	{\headPrint[kana={やたかび},ipa={jatakabi},pos={名},acc={-}]}%
	{%
		\MeaningsBlock{%
	\ryumeaning%
		{浜木綿}{}{}%
	}%
	}%
	{\tailPrint[]}%
\entry
	{\headPrint[kana={や\josho{}た\kako{}ぷ},ipa={jatapu},pos={名},acc={特殊}]}%
	{%
		\MeaningsBlock{%
	\ryumeaning%
		{唐黍。高粱}{}{}%
	}%
	}%
	{\tailPrint[]}%
\entry
	{\headPrint[kana={やたふちぃ},ipa={jatafu̥tʃi},pos={名}]}%
	{%
		\MeaningsBlock{%
	\ryumeaning%
		{\ruby[g]{蓬}{よもぎ}}{}{}%
	}%
	}%
	{\tailPrint[]}%
\entry
	{\headPrint[kana={やたぶに\tl{˩˥}},ipa={jatabuni},pos={名},acc={R}]}%
	{%
		\MeaningsBlock{%
	\ryumeaning%
		{肋骨}{}{}%
	}%
	}%
	{\tailPrint[]}%
\entry
	{\headPrint[kana={やち\tl{˩˥}},ipa={jatʃi},pos={名},acc={R}]}%
	{%
		\MeaningsBlock{%
	\ryumeaning%
		{灸。お灸}{}{}%
	}%
	}%
	{\tailPrint[]}%
\entry
	{\headPrint[kana={やっけー\tl{˩˥}},ipa={jakkeː},pos={名},acc={R}]}%
	{%
		\MeaningsBlock{%
	\ryumeaning%
		{厄介。迷惑}{}{}%
	}%
	}%
	{\tailPrint[]}%
\entry
	{\headPrint[kana={やっさ},ipa={jassa},pos={助}]}%
	{%
		\MeaningsBlock{%
	\ryumeaning%
		{～だ。～だろう}{}{}%
	}%
	}%
	{\tailPrint[]}%
\entry
	{\headPrint[kana={やっさ\tl{˩˥}はん},ipa={jassahaɴ},pos={形},para={やっさへぬ},acc={R}]}%
	{%
		\MeaningsBlock{%
	\ryumeaning%
		{安い。安価だ}{}{}%
	\ryumeaning%
		{た易い}{}{}%
	}%
	}%
	{\tailPrint[]}%
\entry
	{\headPrint[kana={やっちるん\tl{˩˥}},ipa={jattʃiɴ},pos={動},rui={3},para={やっつぬ},acc={R}]}%
	{%
		\MeaningsBlock{%
	\ryumeaning%
		{痩せる。衰弱する}{}{}%
	}%
	}%
	{\tailPrint[]}%
\entry
	{\headPrint[kana={やっ\josho{}ち\kako{}ん},ipa={jattʃiɴ},pos={副},acc={LHL}]}%
	{%
		\MeaningsBlock{%
	\ryumeaning%
		{必ず}{}{}%
	}%
	}%
	{\tailPrint[]}%
\entry
	{\headPrint[kana={やっとぅ},ipa={jattu},pos={副}]}%
	{%
		\MeaningsBlock{%
	\ryumeaning%
		{やっと。ようやく}{}{}%
	}%
	}%
	{\tailPrint[]}%
\entry
	{\headPrint[kana={やっぱり\tl{˥˩}},ipa={jappari},pos={副},acc={F}]}%
	{%
		\MeaningsBlock{%
	\ryumeaning%
		{言い張る。主張する}{}{}%
	}%
	}%
	{\tailPrint[]}%
\entry
	{\headPrint[kana={やとーすん\tl{˩˥}},ipa={jatoːsuɴ},pos={動},rui={2},para={やとーはぬ},acc={R}]}%
	{%
		\MeaningsBlock{%
	\ryumeaning%
		{移植する}{}{}%
	}%
	}%
	{\tailPrint[]}%
\entry
	{\headPrint[kana={やどぅ\tl{˩˥}},ipa={jadu},pos={名},acc={R}]}%
	{%
		\MeaningsBlock{%
	\ryumeaning%
		{戸。雨戸}{}{}%
	}%
	}%
	{\tailPrint[]}%
\entry
	{\headPrint[kana={やどぅ\tl{˩˥}},ipa={jadu},pos={名},acc={R}]}%
	{%
		\MeaningsBlock{%
	\ryumeaning%
		{宿。宿屋}{}{}%
	}%
	}%
	{\tailPrint[]}%
\entry
	{\headPrint[kana={やどぅとぅるん},ipa={jaduturuɴ},pos={句},acc={R}]}%
	{%
		\MeaningsBlock{%
	\ryumeaning%
		{宿る。泊まる。宿泊する}{}{}%
	}%
	}%
	{\tailPrint[]}%
\entry
	{\headPrint[kana={やとすん\tl{˩˥}},ipa={jatosuɴ},pos={動},rui={2},para={やとはぬ},acc={R}]}%
	{%
		\MeaningsBlock{%
	\ryumeaning%
		{雇う。雇用する}{}{}%
	}%
	}%
	{\tailPrint[]}%
\entry
	{\headPrint[kana={やな\tl{˥˩}},ipa={jana},pos={形},acc={F}]}%
	{%
		\MeaningsBlock{%
	\ryumeaning%
		{悪い。劣悪だ}{}{}%
	}%
	}%
	{\tailPrint[]}%
\entry
	{\headPrint[kana={やなかー\tl{˥˩}},ipa={janakaː},pos={名},acc={F}]}%
	{%
		\MeaningsBlock{%
	\ryumeaning%
		{悪臭}{}{}%
	}%
	}%
	{\tailPrint[]}%
\entry
	{\headPrint[kana={やなかたち\tl{˥˩}},ipa={janakḁtatʃi},pos={名},acc={F}]}%
	{%
		\MeaningsBlock{%
	\ryumeaning%
		{悪い形。不格好}{形が悪いこと}{}%
	}%
	}%
	{\tailPrint[]}%
\entry
	{\headPrint[kana={やなくしー\tl{˥˩}},ipa={janaku̥ʃiː},pos={名},acc={F}]}%
	{%
		\MeaningsBlock{%
	\ryumeaning%
		{悪い癖。悪癖}{}{}%
	}%
	}%
	{\tailPrint[]}%
\entry
	{\headPrint[kana={やなくとぅ\tl{˩˥}},ipa={janaku̥tu},pos={名},acc={R}]}%
	{%
		\MeaningsBlock{%
	\ryumeaning%
		{悪事。凶事。不幸}{}{}%
	}%
	}%
	{\tailPrint[]}%
\entry
	{\headPrint[kana={やなごるん\tl{˥˩}},ipa={janagoruɴ},pos={動},rui={1},para={やなごらぬ},acc={F}]}%
	{%
		\MeaningsBlock{%
	\ryumeaning%
		{濁る。汚濁になる}{}{}%
	}%
	}%
	{\tailPrint[]}%
\entry
	{\headPrint[kana={やなたくみ\tl{˩˥}},ipa={janatḁkumi},pos={名},acc={R}]}%
	{%
		\MeaningsBlock{%
	\ryumeaning%
		{悪だくみ}{}{}%
	}%
	}%
	{\tailPrint[]}%
\entry
	{\headPrint[kana={やななれー\tl{˩˥}},ipa={jananareː},pos={名},acc={R}]}%
	{%
		\MeaningsBlock{%
	\ryumeaning%
		{悪い習慣。因習}{}{}%
	}%
	}%
	{\tailPrint[]}%
\entry
	{\headPrint[kana={やなふちぃ\tl{˩˥}},ipa={janafu̥tʃi},pos={名},acc={R}]}%
	{%
		\MeaningsBlock{%
	\ryumeaning%
		{悪口}{}{}%
	}%
	}%
	{\tailPrint[]}%
\entry
	{\headPrint[kana={やなむしぃ\tl{˥˩}},ipa={janamuʃi},pos={名},acc={F}]}%
	{%
		\MeaningsBlock{%
	\ryumeaning%
		{害虫}{}{}%
	}%
	}%
	{\tailPrint[]}%
\entry
	{\headPrint[kana={やなむん\tl{˥˩}},ipa={janamuɴ},pos={名},acc={F}]}%
	{%
		\MeaningsBlock{%
	\ryumeaning%
		{悪者。魔物}{}{}%
	}%
	}%
	{\tailPrint[]}%
\entry
	{\headPrint[kana={やなわしき\tl{˥˩}},ipa={janawaʃi̥ki},pos={名},acc={F}]}%
	{%
		\MeaningsBlock{%
	\ryumeaning%
		{天気が悪い。悪天候}{}{}%
	}%
	}%
	{\tailPrint[]}%
\entry
	{\headPrint[kana={やなわらび\tl{˩˥}},ipa={janawarabi},pos={名},acc={R}]}%
	{%
		\MeaningsBlock{%
	\ryumeaning%
		{悪戯っ子。悪童}{}{}%
	}%
	}%
	{\tailPrint[]}%
\entry
	{\headPrint[kana={やにっしゃ\tl{˩˥}はん},ipa={janiʃʃahaɴ},pos={形},para={やにっしゃへぬ},acc={R}]}%
	{%
		\MeaningsBlock{%
	\ryumeaning%
		{汚ない。汚れている}{}{}%
	}%
	}%
	{\tailPrint[]}%
\entry
	{\headPrint[kana={やば\tl{˩˥}はん},ipa={jabahaɴ},pos={形},para={やばへぬ},acc={R}]}%
	{%
		\MeaningsBlock{%
	\ryumeaning%
		{上手。優れている}{}{}%
	}%
	}%
	{\tailPrint[]}%
\entry
	{\headPrint[kana={やぴくん\tl{˥˩}},ipa={japikuɴ},pos={動},rui={3},acc={F}]}%
	{%
		\MeaningsBlock{%
	\ryumeaning%
		{火傷をする}{}{}%
	}%
	}%
	{\tailPrint[]}%
\entry
	{\headPrint[kana={やびむん\tl{˩˥}},ipa={jabimuɴ},pos={名},acc={R}]}%
	{%
		\MeaningsBlock{%
	\ryumeaning%
		{ろくでなし}{}{}%
	}%
	}%
	{\tailPrint[]}%
\entry
	{\headPrint[kana={やふ\tl{˩˥}},ipa={jafu},pos={名},acc={R}]}%
	{%
		\MeaningsBlock{%
	\ryumeaning%
		{災い。凶事}{}{}%
	}%
	}%
	{\tailPrint[]}%
\entry
	{\headPrint[kana={やふぁらぐん},ipa={jafaraguɴ},pos={動}]}%
	{%
		\MeaningsBlock{%
	\ryumeaning%
		{和らぐ。柔らかになる}{}{}%
	}%
	}%
	{\tailPrint[]}%
\entry
	{\headPrint[kana={やぶす\ws{}なるん},ipa={jabusu̥ naruɴ},pos={句},rui={1},para={やぶすならぬ},acc={-}]}%
	{%
		\MeaningsBlock{%
	\ryumeaning%
		{安心する。安堵する}{}{}%
	}%
	}%
	{\tailPrint[]}%
\entry
	{\headPrint[kana={やぶりん\tl{˩˥}},ipa={jaburiɴ},pos={動},rui={3},para={やぶるぬ},acc={R}]}%
	{%
		\MeaningsBlock{%
	\ryumeaning%
		{破れる。壊れる}{}{}%
	}%
	}%
	{\tailPrint[]}%
\entry
	{\headPrint[kana={やぶるん\tl{˩˥}},ipa={jaburuɴ},pos={動},rui={1},para={やぶらぬ},acc={R}]}%
	{%
		\MeaningsBlock{%
	\ryumeaning%
		{破る。壊す}{}{}%
	}%
	}%
	{\tailPrint[]}%
\entry
	{\headPrint[kana={やま\tl{˩˥}},ipa={jama},pos={名},acc={R}]}%
	{%
		\MeaningsBlock{%
	\ryumeaning%
		{山。森}{}{}%
	}%
	}%
	{\tailPrint[]}%
\entry
	{\headPrint[kana={や\josho{}まー\kako{}し},ipa={jamaːʃi},pos={副},acc={LHL}]}%
	{%
		\MeaningsBlock{%
	\ryumeaning%
		{静かに。ゆっくりと}{}{}%
	}%
	}%
	{\tailPrint[]}%
\entry
	{\headPrint[kana={やまかたなし\tl{˩˥}},ipa={jamakḁtanaʃi},pos={名},acc={R}]}%
	{%
		\MeaningsBlock{%
	\ryumeaning%
		{山刀}{}{}%
	}%
	}%
	{\tailPrint[]}%
\entry
	{\headPrint[kana={やまぐ\tl{˩˥}},ipa={jamagu},pos={名},acc={R}]}%
	{%
		\MeaningsBlock{%
	\ryumeaning%
		{乱暴者。暴れ者。腕白者。横着な者。悪戯っぽい者}{}{}%
	}%
	}%
	{\tailPrint[]}%
\entry
	{\headPrint[kana={やまぐ\tl{˩˥}\ws{}すん\tl{˥˩}},ipa={jamagu suɴ},pos={句},acc={R F}]}%
	{%
		\MeaningsBlock{%
	\ryumeaning%
		{乱暴を働く}{}{}%
	}%
	}%
	{\tailPrint[]}%
\entry
	{\headPrint[kana={やまじぃま\tl{˩˥}},ipa={jamadʒima},pos={名},acc={R}]}%
	{%
		\MeaningsBlock{%
	\ryumeaning%
		{山島。山国}{}{}%
	}%
	}%
	{\tailPrint[]}%
\entry
	{\headPrint[kana={やましぐとぅ\tl{˩˥}},ipa={jamaʃigutu},pos={名},acc={R}]}%
	{%
		\MeaningsBlock{%
	\ryumeaning%
		{山仕事}{}{}%
	}%
	}%
	{\tailPrint[]}%
\entry
	{\headPrint[kana={やますん\tl{˩˥}},ipa={jamasuɴ},pos={動},rui={2},para={やまはぬ},acc={R}]}%
	{%
		\MeaningsBlock{%
	\ryumeaning%
		{痛める。傷つける}{}{}%
	}%
	}%
	{\tailPrint[]}%
\entry
	{\headPrint[kana={やまとぅ\tl{˩˥}},ipa={jamatu},pos={名},acc={R}]}%
	{%
		\MeaningsBlock{%
	\ryumeaning%
		{大和。日本。日本本土}{}{}%
	}%
	}%
	{\tailPrint[]}%
\entry
	{\headPrint[kana={やまとぅ\tl{˩˥}ぴぃとぅ\tl{˥˩}},ipa={jamatupi̥tu},pos={名},acc={R+F}]}%
	{%
		\MeaningsBlock{%
	\ryumeaning%
		{日本人。内地人}{}{}%
	}%
	}%
	{\tailPrint[]}%
\entry
	{\headPrint[kana={やまどぅみ\tl{˩˥}},ipa={jamadumi},pos={名},acc={R}]}%
	{%
		\MeaningsBlock{%
	\ryumeaning%
		{山止め}{王府時代の禁忌}{}%
	}%
	}%
	{\tailPrint[]}%
\entry
	{\headPrint[kana={やまとぅむに\tl{˩˥}},ipa={jamatumuni},pos={名},acc={R}]}%
	{%
		\MeaningsBlock{%
	\ryumeaning%
		{日本語。標準語}{}{}%
	}%
	}%
	{\tailPrint[]}%
\entry
	{\headPrint[kana={やまにんじゅ\tl{˩˥}},ipa={jamanindʒu},pos={名},acc={R}]}%
	{%
		\MeaningsBlock{%
	\ryumeaning%
		{御嶽の祭祀に参加する村人}{}{}%
	}%
	}%
	{\tailPrint[]}%
\entry
	{\headPrint[kana={やまんぐ},ipa={jamangu},pos={名}]}%
	{%
		\MeaningsBlock{%
	\ryumeaning%
		{乱暴者。暴れ者。腕白者。横着な者。悪戯っぽい者}{}{}%
	}%
	}%
	{\tailPrint[]}%
\entry
	{\headPrint[kana={やまんぐ\tl{˩˥}さーん},ipa={jamaŋgusaːɴ},pos={形},acc={R}]}%
	{%
		\MeaningsBlock{%
	\ryumeaning%
		{暴れん坊だ}{}{}%
	\ryumeaning%
		{腕白だ。悪戯好きだ}{}{}%
	}%
	}%
	{\tailPrint[]}%
\entry
	{\headPrint[kana={やみ\tl{˥˩}},ipa={jami},pos={名},acc={F}]}%
	{%
		\MeaningsBlock{%
	\ryumeaning%
		{闇。真っ暗}{}{}%
	}%
	}%
	{\tailPrint[]}%
\entry
	{\headPrint[kana={やみ\tl{˩˥}},ipa={jami},pos={名},acc={R}]}%
	{%
		\MeaningsBlock{%
	\ryumeaning%
		{\ruby[g]{病}{やまい}。病気}{}{}%
	}%
	}%
	{\tailPrint[]}%
\entry
	{\headPrint[kana={やみるん\tl{˥˩}},ipa={jamiruɴ},pos={動},rui={3},para={やむぬ},acc={F}]}%
	{%
		\MeaningsBlock{%
	\ryumeaning%
		{止める}{}{}%
	\ryumeaning%
		{辞職する。辞任する}{}{}%
	}%
	}%
	{\tailPrint[]}%
\entry
	{\headPrint[kana={やむん\tl{˥˩}},ipa={jamuɴ},pos={動},rui={1},para={やまぬ},acc={F}]}%
	{%
		\MeaningsBlock{%
	\ryumeaning%
		{（風雨が）止む}{}{}%
	}%
	}%
	{\tailPrint[]}%
\entry
	{\headPrint[kana={やむん\tl{˩˥}},ipa={jamuɴ},pos={動},rui={1},para={やまぬ},acc={R}]}%
	{%
		\MeaningsBlock{%
	\ryumeaning%
		{痛む。病気に罹る}{}{}%
	}%
	}%
	{\tailPrint[]}%
\entry
	{\headPrint[kana={やらすん\tl{˥˩}},ipa={jarasuɴ},pos={動},rui={2},para={やらはぬ},acc={F}]}%
	{%
		\MeaningsBlock{%
	\ryumeaning%
		{よこす。遣わす}{}{}%
	}%
	}%
	{\tailPrint[]}%
\entry
	{\headPrint[kana={やら\tl{˩˥}はん},ipa={jarahaɴ},pos={形},para={やらへぬ},acc={R}]}%
	{%
		\MeaningsBlock{%
	\ryumeaning%
		{柔らかい}{}{}%
	}%
	}%
	{\tailPrint[]}%
\entry
	{\headPrint[kana={やらび\tl{˩˥}},ipa={jarabi},pos={名},acc={R}]}%
	{%
		\MeaningsBlock{%
	\ryumeaning%
		{童。子供}{}{}%
	}%
	}%
	{\tailPrint[]}%
\entry
	{\headPrint[kana={やらびなー\tl{˩˥}},ipa={jarabinaː},pos={名},acc={R}]}%
	{%
		\MeaningsBlock{%
	\ryumeaning%
		{童名。幼名}{}{}%
	}%
	}%
	{\tailPrint[]}%
\entry
	{\headPrint[kana={やらぶ\tl{˥}},ipa={jarabu},pos={名},acc={H}]}%
	{%
		\MeaningsBlock{%
	\ryumeaning%
		{テリハボク}{植物名}{}%
	}%
	}%
	{\tailPrint[]}%
\entry
	{\headPrint[kana={やらふぁー\tl{˩˥}\ws{}すん\tl{˥˩}},ipa={jarafaː suɴ},pos={句},acc={R F}]}%
	{%
		\MeaningsBlock{%
	\ryumeaning%
		{柔らかくする}{}{}%
	}%
	}%
	{\tailPrint[]}%
\entry
	{\headPrint[kana={やらふぁー\tl{˩˥}\ws{}なるん\tl{˩˥}},ipa={jarafaː naruɴ},pos={句},acc={R R}]}%
	{%
		\MeaningsBlock{%
	\ryumeaning%
		{柔らかくなる}{}{}%
	}%
	}%
	{\tailPrint[]}%
\entry
	{\headPrint[kana={やり\tl{˥˩}},ipa={jari},pos={名},acc={F}]}%
	{%
		\MeaningsBlock{%
	\ryumeaning%
		{槍}{}{}%
	}%
	}%
	{\tailPrint[]}%
\entry
	{\headPrint[kana={やりしちるん\tl{˩˥}},ipa={jariʃi̥tʃiruɴ},pos={動},rui={3},para={やりしつぬ},acc={R}]}%
	{%
		\MeaningsBlock{%
	\ryumeaning%
		{破って捨てる}{}{}%
	}%
	}%
	{\tailPrint[]}%
\entry
	{\headPrint[kana={やりすぅぬ\tl{˩˥}},ipa={jarisu̥nu},pos={名},acc={R}]}%
	{%
		\MeaningsBlock{%
	\ryumeaning%
		{ぼろの着物}{}{}%
	}%
	}%
	{\tailPrint[]}%
\entry
	{\headPrint[kana={やりふちぃ\tl{˩˥}},ipa={jarifu̥tʃi},pos={名},acc={R}]}%
	{%
		\MeaningsBlock{%
	\ryumeaning%
		{破れ目。裂け目}{}{}%
	}%
	}%
	{\tailPrint[]}%
\entry
	{\headPrint[kana={やりぽーるん\tl{˩˥}},ipa={jaripoːruɴ},pos={動},rui={1},para={やりぽーらぬ},acc={R}]}%
	{%
		\MeaningsBlock{%
	\ryumeaning%
		{破り散らかす}{}{}%
	}%
	}%
	{\tailPrint[]}%
\entry
	{\headPrint[kana={やりるん\tl{˩˥}},ipa={jariruɴ},pos={動},rui={3},para={やるぬ},acc={R}]}%
	{%
		\MeaningsBlock{%
	\ryumeaning%
		{破れる}{}{}%
	}%
	}%
	{\tailPrint[]}%
\entry
	{\headPrint[kana={やるはく\tl{˩˥}},ipa={jaruhḁku},pos={名},acc={R}]}%
	{%
		\MeaningsBlock{%
	\ryumeaning%
		{ぼろ。ぼろ切れ}{}{}%
	}%
	}%
	{\tailPrint[]}%
\entry
	{\headPrint[kana={やるん\tl{˥˩}},ipa={jaruɴ},pos={動},rui={1},para={あらぬ},acc={F}]}%
	{%
		\MeaningsBlock{%
	\ryumeaning%
		{～だ。～である}{}{}%
	}%
	}%
	{\tailPrint[]}%
\entry
	{\headPrint[kana={やん\tl{˩˥}},ipa={jaɴ},pos={名},acc={R}]}%
	{%
		\MeaningsBlock{%
	\ryumeaning%
		{\ruby[g]{病}{やまい}。病気}{}{}%
	}%
	}%
	{\tailPrint[]}%
\entry
	{\headPrint[kana={やんばるしん\tl{˩˥}},ipa={jambaruʃiɴ},pos={名},acc={R}]}%
	{%
		\MeaningsBlock{%
	\ryumeaning%
		{\ruby[g]{山}{やん}\ruby[g]{原}{ばる}\ruby[g]{船}{せん}}{沖縄在来の帆船}{}%
	}%
	}%
	{\tailPrint[]}%
\LetterHead{ゆ}
\entry
	{\headPrint[kana={ゆー\tl{˥˩}},ipa={juː},pos={名},acc={F}]}%
	{%
		\MeaningsBlock{%
	\ryumeaning%
		{魚。魚類}{}{}%
	}%
	}%
	{\tailPrint[]}%
\entry
	{\headPrint[kana={ゆー\tl{˥˩}},ipa={juː},pos={名},acc={F}]}%
	{%
		\MeaningsBlock{%
	\ryumeaning%
		{世。時代}{}{}%
	}%
	}%
	{\tailPrint[]}%
\entry
	{\headPrint[kana={ゆー\tl{˩˥}},ipa={juː},pos={名},acc={R}]}%
	{%
		\MeaningsBlock{%
	\ryumeaning%
		{\ruby[g]{粥}{かゆ}}{}{}%
	}%
	}%
	{\tailPrint[]}%
\entry
	{\headPrint[kana={ゆー\tl{˩˥}\ws{}あうん\tl{˥}},ipa={juː auɴ},pos={句},acc={R H}]}%
	{%
		\MeaningsBlock{%
	\ryumeaning%
		{よく合う。よく似合う}{}{}%
	}%
	}%
	{\tailPrint[]}%
\entry
	{\headPrint[kana={ゆーあみ\tl{˩˥}},ipa={juː.ami},pos={名},acc={R}]}%
	{%
		\MeaningsBlock{%
	\ryumeaning%
		{夜雨}{夜間に降る雨。豊年の印。ゆがふ雨として喜ばれた}{}%
	}%
	}%
	{\tailPrint[]}%
\entry
	{\headPrint[kana={ゆーがーり\tl{˥˩}},ipa={juːgaːri},pos={名},acc={F}]}%
	{%
		\MeaningsBlock{%
	\ryumeaning%
		{世変り}{}{}%
	}%
	}%
	{\tailPrint[]}%
\entry
	{\headPrint[kana={ゆーかどぅ\tl{˥˩}},ipa={juːkadu},pos={名},acc={F}]}%
	{%
		\MeaningsBlock{%
	\ryumeaning%
		{四ツ角。四隅}{}{}%
	}%
	}%
	{\tailPrint[]}%
\entry
	{\headPrint[kana={ゆー\tl{˩˥}さんちゃ\tl{˥}},ipa={juːsantʃa},pos={句},acc={R H}]}%
	{%
		\MeaningsBlock{%
	\ryumeaning%
		{もしかしたら}{}{}%
	}%
	}%
	{\tailPrint[]}%
\entry
	{\headPrint[kana={ゆーじ\tl{˩˥}},ipa={juːdʒi},pos={名},acc={R}]}%
	{%
		\MeaningsBlock{%
	\ryumeaning%
		{用事。行事}{}{}%
	}%
	}%
	{\tailPrint[]}%
\entry
	{\headPrint[kana={ゆーしき},ipa={juːʃi̥ki},pos={名}]}%
	{%
		\MeaningsBlock{%
	\ryumeaning%
		{魚突き}{銛で突く漁法}{}%
	}%
	}%
	{\tailPrint[]}%
\entry
	{\headPrint[kana={ゆーち\tl{˥˩}},ipa={juːtʃi},pos={名},acc={F}]}%
	{%
		\MeaningsBlock{%
	\ryumeaning%
		{四つ}{簡略化：ゆー}{}%
	}%
	}%
	{\tailPrint[]}%
\entry
	{\headPrint[kana={ゆーとぅり\tl{˥˩}},ipa={juːtu̥ri},pos={名},acc={F}]}%
	{%
		\MeaningsBlock{%
	\ryumeaning%
		{魚獲り。漁業}{}{}%
	}%
	}%
	{\tailPrint[]}%
\entry
	{\headPrint[kana={ゆーどぅり\tl{˩˥}},ipa={juːduri},pos={名},acc={R}]}%
	{%
		\MeaningsBlock{%
	\ryumeaning%
		{夕凪}{}{}%
	}%
	}%
	{\tailPrint[]}%
\entry
	{\headPrint[kana={ゆーび\tl{˩˥}},ipa={juːbi},pos={名},acc={R}]}%
	{%
		\MeaningsBlock{%
	\ryumeaning%
		{夕べ。昨夜}{}{}%
	}%
	}%
	{\tailPrint[]}%
\entry
	{\headPrint[kana={ゆーべー\tl{˩˥}},ipa={juːbeː},pos={名},acc={R}]}%
	{%
		\MeaningsBlock{%
	\ryumeaning%
		{\ruby[g]{妾}{めかけ}。情婦}{}{}%
	}%
	}%
	{\tailPrint[]}%
\entry
	{\headPrint[kana={ゆーべーかち\tl{˩˥}},ipa={juːbeːkḁtʃi},pos={名},acc={R}]}%
	{%
		\MeaningsBlock{%
	\ryumeaning%
		{戻り風}{風向が元にもどる風}{}%
	}%
	}%
	{\tailPrint[]}%
\entry
	{\headPrint[kana={ゆーべーぷち},ipa={juːbeːpu̥tʃi},pos={名}]}%
	{%
		\MeaningsBlock{%
	\ryumeaning%
		{流れ星}{直訳「夜這い星」}{}%
	}%
	}%
	{\tailPrint[]}%
\entry
	{\headPrint[kana={ゆーほすん},ipa={juːhosuɴ},pos={句},para={ゆーほはぬ}]}%
	{%
		\MeaningsBlock{%
	\ryumeaning%
		{魚釣り}{}{}%
	}%
	}%
	{\tailPrint[]}%
\entry
	{\headPrint[kana={ゆい\tl{˥˩}},ipa={jui},pos={動},para={ゆわぬ},acc={F}]}%
	{%
		\MeaningsBlock{%
	\ryumeaning%
		{叱る}{}{}%
	}%
	}%
	{\tailPrint[]}%
\entry
	{\headPrint[kana={ゆい\tl{˥˩}},ipa={jui},pos={名},acc={F}]}%
	{%
		\MeaningsBlock{%
	\ryumeaning%
		{\ruby[g]{結}{ゆい}。共同作業}{}{}%
	}%
	}%
	{\tailPrint[]}%
\entry
	{\headPrint[kana={ゆうん\tl{˥˩}},ipa={juːɴ},pos={動},para={ゆわぬ},acc={F}]}%
	{%
		\MeaningsBlock{%
	\ryumeaning%
		{結う。結ぶ}{}{}%
	}%
	}%
	{\tailPrint[]}%
\entry
	{\headPrint[kana={ゆがふ\tl{˥˩}},ipa={jugafu},pos={名},acc={F}]}%
	{%
		\MeaningsBlock{%
	\ryumeaning%
		{豊年}{直訳「世果報」}{}%
	}%
	}%
	{\tailPrint[]}%
\entry
	{\headPrint[kana={ゆがみるん\tl{˥˩}},ipa={jugamiruɴ},pos={動},acc={F}]}%
	{%
		\MeaningsBlock{%
	\ryumeaning%
		{歪める}{}{}%
	}%
	}%
	{\tailPrint[]}%
\entry
	{\headPrint[kana={ゆがむん\tl{˥˩}},ipa={jugamuɴ},pos={動},rui={1},para={ゆがまぬ},acc={F}]}%
	{%
		\MeaningsBlock{%
	\ryumeaning%
		{\ruby[g]{歪}{ゆが}む}{}{}%
	}%
	}%
	{\tailPrint[]}%
\entry
	{\headPrint[kana={ゆぐ\tl{˩˥}},ipa={jugu},pos={名},acc={R}]}%
	{%
		\MeaningsBlock{%
	\ryumeaning%
		{欲}{}{}%
	}%
	}%
	{\tailPrint[]}%
\entry
	{\headPrint[kana={ゆぐ\tl{˩˥}さん},ipa={jugusaɴ},pos={形},para={ゆぐさへぬ},acc={R}]}%
	{%
		\MeaningsBlock{%
	\ryumeaning%
		{欲張りだ。欲深い}{}{}%
	\ryumeaning%
		{けちだ。しみったれだ}{}{}%
	}%
	}%
	{\tailPrint[]}%
\entry
	{\headPrint[kana={ゆくしむに\tl{˩˥}},ipa={jukuʃi̥muni},pos={名},acc={R}]}%
	{%
		\MeaningsBlock{%
	\ryumeaning%
		{うそ。いつわり}{}{}%
	}%
	}%
	{\tailPrint[]}%
\entry
	{\headPrint[kana={ゆくすか\tl{˩˥}},ipa={jukusu̥ka},pos={動（接）},acc={R}]}%
	{%
		\MeaningsBlock{%
	\ryumeaning%
		{欲張る}{}{}%
	}%
	}%
	{\tailPrint[]}%
\entry
	{\headPrint[kana={ゆぐすん\tl{˥˩}},ipa={jugusuɴ},pos={動},rui={2b},para={ゆぐさぬ},acc={F}]}%
	{%
		\MeaningsBlock{%
	\ryumeaning%
		{汚す}{}{}%
	}%
	}%
	{\tailPrint[]}%
\entry
	{\headPrint[kana={ゆくたーるん\tl{˩˥}},ipa={juku̥taːruɴ},pos={動},rui={1},para={ゆくたーらぬ},acc={R}]}%
	{%
		\MeaningsBlock{%
	\ryumeaning%
		{横たわる。横になる}{}{}%
	\ryumeaning%
		{寝そべる}{}{}%
	}%
	}%
	{\tailPrint[]}%
\entry
	{\headPrint[kana={ゆぐぬ\ws{}ねーぬ},ipa={jugunu neːnu},pos={句}]}%
	{%
		\MeaningsBlock{%
	\ryumeaning%
		{無欲だ。欲がない}{食欲について淡白だ}{}%
	}%
	}%
	{\tailPrint[]}%
\entry
	{\headPrint[kana={ゆくばるん},ipa={jukubaruɴ},pos={動}]}%
	{%
		\MeaningsBlock{%
	\ryumeaning%
		{欲張る}{}{}%
	}%
	}%
	{\tailPrint[]}%
\entry
	{\headPrint[kana={ゆぐり\tl{˥˩}},ipa={juguri},pos={名},acc={F}]}%
	{%
		\MeaningsBlock{%
	\ryumeaning%
		{汚れ}{}{}%
	}%
	}%
	{\tailPrint[]}%
\entry
	{\headPrint[kana={ゆぐりるん\tl{˥˩}},ipa={juguriruɴ},pos={動},rui={3},para={ゆぐるぬ},acc={F}]}%
	{%
		\MeaningsBlock{%
	\ryumeaning%
		{汚れる。けがれる}{}{}%
	}%
	}%
	{\tailPrint[]}%
\entry
	{\headPrint[kana={ゆごー\tl{˥˩}},ipa={jugoː},pos={名},acc={F}]}%
	{%
		\MeaningsBlock{%
	\ryumeaning%
		{隅。片隅}{東の村では\emb{やごー}}{}%
	}%
	}%
	{\tailPrint[]}%
\entry
	{\headPrint[kana={ゆごー\tl{˩˥}\ws{}すん\tl{˥˩}},ipa={jugoː suɴ},pos={句},acc={R F}]}%
	{%
		\MeaningsBlock{%
	\ryumeaning%
		{お休みになる}{「休む」、「憩う」の敬語}{}%
	}%
	}%
	{\tailPrint[]}%
\entry
	{\headPrint[kana={ゆさし\tl{˥˩}},ipa={jusaʃi},pos={名},acc={F}]}%
	{%
		\MeaningsBlock{%
	\ryumeaning%
		{四つ星}{}{}%
	}%
	}%
	{\tailPrint[]}%
\entry
	{\headPrint[kana={ゆしき\tl{˥˩}},ipa={juʃi̥ki},pos={名},acc={F}]}%
	{%
		\MeaningsBlock{%
	\ryumeaning%
		{ススキ}{植物名}{}%
	}%
	}%
	{\tailPrint[]}%
\entry
	{\headPrint[kana={ゆしぐとぅ\tl{˥˩}},ipa={juʃi̥gutu},pos={名},acc={F}]}%
	{%
		\MeaningsBlock{%
	\ryumeaning%
		{忠言。訓話}{}{}%
	}%
	}%
	{\tailPrint[]}%
\entry
	{\headPrint[kana={ゆしるん\tl{˥˩}},ipa={juʃi̥ruɴ},pos={動},rui={3},para={ゆすぬ},acc={F}]}%
	{%
		\MeaningsBlock{%
	\ryumeaning%
		{寄せる。移動する}{}{}%
	}%
	}%
	{\tailPrint[]}%
\entry
	{\headPrint[kana={ゆす\tl{˩˥}},ipa={jusu},pos={名},acc={R}]}%
	{%
		\MeaningsBlock{%
	\ryumeaning%
		{露。朝露}{}{}%
	}%
	}%
	{\tailPrint[]}%
\entry
	{\headPrint[kana={ゆすぱり\tl{˩˥}},ipa={jusu̥pari},pos={名},acc={R}]}%
	{%
		\MeaningsBlock{%
	\ryumeaning%
		{寝小便}{}{}%
	}%
	}%
	{\tailPrint[]}%
\entry
	{\headPrint[kana={ゆずるん\tl{˥˩}},ipa={judzuruɴ},pos={動},rui={1},para={ゆずらぬ},acc={F}]}%
	{%
		\MeaningsBlock{%
	\ryumeaning%
		{譲る}{}{}%
	}%
	}%
	{\tailPrint[]}%
\entry
	{\headPrint[kana={ゆた\tl{˥˩}},ipa={juta},pos={名},acc={F}]}%
	{%
		\MeaningsBlock{%
	\ryumeaning%
		{易者。占い師}{}{}%
	}%
	}%
	{\tailPrint[]}%
\entry
	{\headPrint[kana={ゆだ\tl{˥˩}},ipa={juda},pos={名},acc={F}]}%
	{%
		\MeaningsBlock{%
	\ryumeaning%
		{ユウナ}{植物名}{}%
	}%
	}%
	{\tailPrint[]}%
\entry
	{\headPrint[kana={ゆだぱ\tl{˥˩}},ipa={judapa},pos={名},acc={F}]}%
	{%
		\MeaningsBlock{%
	\ryumeaning%
		{枝葉}{}{}%
	}%
	}%
	{\tailPrint[]}%
\entry
	{\headPrint[kana={ゆだばり\tl{˥˩}},ipa={judabari},pos={名},acc={F}]}%
	{%
		\MeaningsBlock{%
	\ryumeaning%
		{黙っていること。ぐずぐずすること}{}{}%
	}%
	}%
	{\tailPrint[]}%
\entry
	{\headPrint[kana={ゆだり\tl{˥˩}},ipa={judari},pos={名},acc={F}]}%
	{%
		\MeaningsBlock{%
	\ryumeaning%
		{涎}{}{}%
	}%
	}%
	{\tailPrint[]}%
\entry
	{\headPrint[kana={ゆどぅあみ\tl{˩˥}},ipa={judu.ami},pos={名},acc={R}]}%
	{%
		\MeaningsBlock{%
	\ryumeaning%
		{梅雨。五月雨}{}{}%
	}%
	}%
	{\tailPrint[]}%
\entry
	{\headPrint[kana={ゆどぅし},ipa={juduʃi},pos={副}]}%
	{%
		\MeaningsBlock{%
	\ryumeaning%
		{夜通し。一晩中}{}{}%
	}%
	}%
	{\tailPrint[]}%
\entry
	{\headPrint[kana={ゆどぅむん\tl{˩˥}},ipa={judumuɴ},pos={動},rui={1},para={ゆどぅまぬ},acc={R}]}%
	{%
		\MeaningsBlock{%
	\ryumeaning%
		{淀む。停滞する}{}{}%
	}%
	}%
	{\tailPrint[]}%
\entry
	{\headPrint[kana={ゆなが\tl{˩˥}},ipa={junaga},pos={名},acc={R}]}%
	{%
		\MeaningsBlock{%
	\ryumeaning%
		{夜中。真夜中}{}{}%
	}%
	}%
	{\tailPrint[]}%
\entry
	{\headPrint[kana={ゆにげー\tl{˥˩}},ipa={junigeː},pos={名},acc={F}]}%
	{%
		\MeaningsBlock{%
	\ryumeaning%
		{世願い}{豊作祈願など神行事}{}%
	}%
	}%
	{\tailPrint[]}%
\entry
	{\headPrint[kana={ゆぬ},ipa={junu},pos={接頭}]}%
	{%
		\MeaningsBlock{%
	\ryumeaning%
		{同じ}{}{}%
	}%
	}%
	{\tailPrint[]}%
\entry
	{\headPrint[kana={ゆぬぐ\tl{˩˥}},ipa={junugu},pos={名},acc={R}]}%
	{%
		\MeaningsBlock{%
	\ryumeaning%
		{（小麦の）はったい粉}{}{}%
	}%
	}%
	{\tailPrint[]}%
\entry
	{\headPrint[kana={ゆぬしゅく\tl{˩˥}},ipa={junuʃu̥ku},pos={名},acc={R}]}%
	{%
		\MeaningsBlock{%
	\ryumeaning%
		{同じ程度}{}{}%
	}%
	}%
	{\tailPrint[]}%
\entry
	{\headPrint[kana={ゆぬとぅし\tl{˩˥}},ipa={junutuʃi},pos={名},acc={R}]}%
	{%
		\MeaningsBlock{%
	\ryumeaning%
		{同じ年。同年}{}{}%
	}%
	}%
	{\tailPrint[]}%
\entry
	{\headPrint[kana={ゆぬむぬ},ipa={junumunu},pos={名}]}%
	{%
		\MeaningsBlock{%
	\ryumeaning%
		{同じ。同じもの。同じこと}{}{}%
	}%
	}%
	{\tailPrint[]}%
\entry
	{\headPrint[kana={ゆぬめー\tl{˥˩}},ipa={junumeː},pos={名},acc={F}]}%
	{%
		\MeaningsBlock{%
	\ryumeaning%
		{魚の目}{足などに出るいぼ}{}%
	}%
	}%
	{\tailPrint[]}%
\entry
	{\headPrint[kana={ゆねん\tl{˩˥}},ipa={juneɴ},pos={名},acc={R}]}%
	{%
		\MeaningsBlock{%
	\ryumeaning%
		{夕方。夕暮れ}{}{}%
	}%
	}%
	{\tailPrint[]}%
\entry
	{\headPrint[kana={ゆのー\tl{˥˩}},ipa={junoː},pos={名},acc={F}]}%
	{%
		\MeaningsBlock{%
	\ryumeaning%
		{与那国}{与那国島。八重山諸島の島の名。最も西の方に位置する}{}%
	}%
	}%
	{\tailPrint[]}%
\entry
	{\headPrint[kana={ゆぴくん\tl{˩˥}},ipa={jupi̥kuɴ},pos={動},rui={1},para={ゆぴかぬ},acc={R}]}%
	{%
		\MeaningsBlock{%
	\ryumeaning%
		{茹でる。湯がく}{}{}%
	}%
	}%
	{\tailPrint[]}%
\entry
	{\headPrint[kana={ゆびだすん\tl{˥˩}},ipa={jubidasuɴ},pos={動},rui={2},para={ゆびんだはぬ},acc={F}]}%
	{%
		\MeaningsBlock{%
	\ryumeaning%
		{吸い出す}{}{}%
	}%
	}%
	{\tailPrint[]}%
\entry
	{\headPrint[kana={ゆふる\tl{˩˥}},ipa={jufu̥ru},pos={名},acc={R}]}%
	{%
		\MeaningsBlock{%
	\ryumeaning%
		{風呂}{}{}%
	}%
	}%
	{\tailPrint[]}%
\entry
	{\headPrint[kana={ゆぶん\tl{˥˩}},ipa={jubuɴ},pos={動},rui={1},para={ゆばぬ},acc={F}]}%
	{%
		\MeaningsBlock{%
	\ryumeaning%
		{呼ぶ}{}{}%
	}%
	}%
	{\tailPrint[]}%
\entry
	{\headPrint[kana={ゆぶん\tl{˥˩}},ipa={jubuɴ},pos={動},rui={1},para={ゆばぬ},acc={F}]}%
	{%
		\MeaningsBlock{%
	\ryumeaning%
		{（水気を）吸う。啜る}{}{}%
	}%
	}%
	{\tailPrint[]}%
\entry
	{\headPrint[kana={ゆみ\tl{˥˩}},ipa={jumi},pos={名},acc={F}]}%
	{%
		\MeaningsBlock{%
	\ryumeaning%
		{嫁}{}{}%
	}%
	}%
	{\tailPrint[]}%
\entry
	{\headPrint[kana={ゆむぬ\tl{˩˥}},ipa={jumunu},pos={名},acc={R}]}%
	{%
		\MeaningsBlock{%
	\ryumeaning%
		{夕飯。夕食}{}{}%
	}%
	}%
	{\tailPrint[]}%
\entry
	{\headPrint[kana={ゆむん\tl{˩˥}},ipa={jumuɴ},pos={動},rui={1},para={ゆまぬ},acc={R}]}%
	{%
		\MeaningsBlock{%
	\ryumeaning%
		{読む}{}{}%
	\ryumeaning%
		{数える}{}{}%
	}%
	}%
	{\tailPrint[]}%
\entry
	{\headPrint[kana={ゆらし\tl{˥˩}},ipa={juraʃi},pos={名},acc={F}]}%
	{%
		\MeaningsBlock{%
	\ryumeaning%
		{（選別用の）箕}{}{}%
	}%
	}%
	{\tailPrint[]}%
\entry
	{\headPrint[kana={ゆらすん\tl{˥˩}},ipa={jurasuɴ},pos={動},rui={2},para={ゆらはぬ},acc={F}]}%
	{%
		\MeaningsBlock{%
	\ryumeaning%
		{揺する。振るう}{}{}%
	}%
	}%
	{\tailPrint[]}%
\entry
	{\headPrint[kana={ゆりしき\tl{˥˩}},ipa={juriʃi̥ki},pos={名},acc={F}]}%
	{%
		\MeaningsBlock{%
	\ryumeaning%
		{閏月}{}{}%
	}%
	}%
	{\tailPrint[]}%
\entry
	{\headPrint[kana={ゆりどぅし\tl{˥˩}},ipa={juriduʃi},pos={名},acc={F}]}%
	{%
		\MeaningsBlock{%
	\ryumeaning%
		{閏年}{}{}%
	}%
	}%
	{\tailPrint[]}%
\entry
	{\headPrint[kana={ゆりはかるん\tl{˥˩}},ipa={jurihḁkaruɴ},pos={動},rui={1},para={ゆりはからぬ},acc={F}]}%
	{%
		\MeaningsBlock{%
	\ryumeaning%
		{よりかかる。もたれかかる}{}{}%
	}%
	}%
	{\tailPrint[]}%
\entry
	{\headPrint[kana={ゆる\tl{˩˥}},ipa={juru},pos={名},acc={R}]}%
	{%
		\MeaningsBlock{%
	\ryumeaning%
		{夜。夜間}{}{}%
	}%
	}%
	{\tailPrint[]}%
\entry
	{\headPrint[kana={ゆるあみ\tl{˩˥}},ipa={juru.ami},pos={名},acc={R}]}%
	{%
		\MeaningsBlock{%
	\ryumeaning%
		{夜雨}{夜間に降る雨。豊年の印。ゆがふ雨として喜ばれた}{}%
	}%
	}%
	{\tailPrint[]}%
\entry
	{\headPrint[kana={ゆるさーるん},ipa={jurusaːruɴ},pos={動}]}%
	{%
		\MeaningsBlock{%
	\ryumeaning%
		{許される}{}{}%
	}%
	}%
	{\tailPrint[]}%
\entry
	{\headPrint[kana={ゆるさ\tl{˩˥}はん},ipa={jurusahaɴ},pos={形},para={ゆるさへぬ},acc={R}]}%
	{%
		\MeaningsBlock{%
	\ryumeaning%
		{緩い}{}{}%
	}%
	}%
	{\tailPrint[]}%
\entry
	{\headPrint[kana={ゆるし\tl{˩˥}},ipa={juruʃi},pos={名},acc={R}]}%
	{%
		\MeaningsBlock{%
	\ryumeaning%
		{許し。許可}{}{}%
	}%
	}%
	{\tailPrint[]}%
\entry
	{\headPrint[kana={ゆるすん\tl{˩˥}},ipa={jurusuɴ},pos={動},rui={2},para={ゆるさぬ},acc={R}]}%
	{%
		\MeaningsBlock{%
	\ryumeaning%
		{許す。許可する}{}{}%
	}%
	}%
	{\tailPrint[]}%
\entry
	{\headPrint[kana={ゆるみるん\tl{˩˥}},ipa={jurumiruɴ},pos={動},rui={3},para={ゆるむぬ},acc={R}]}%
	{%
		\MeaningsBlock{%
	\ryumeaning%
		{緩める。弱める}{}{}%
	}%
	}%
	{\tailPrint[]}%
\entry
	{\headPrint[kana={ゆるみん\tl{˩˥}},ipa={jurumiɴ},pos={動},rui={3},para={ゆるむぬ},acc={R}]}%
	{%
		\MeaningsBlock{%
	\ryumeaning%
		{緩める。弱める}{}{}%
	}%
	}%
	{\tailPrint[]}%
\entry
	{\headPrint[kana={ゆるん\tl{˥˩}},ipa={juruɴ},pos={動},rui={1},para={ゆらぬ},acc={F}]}%
	{%
		\MeaningsBlock{%
	\ryumeaning%
		{寄る}{寄って来る}{}%
	\ryumeaning%
		{寄り合う。集まる}{}{}%
	}%
	}%
	{\tailPrint[]}%
\entry
	{\headPrint[kana={ゆん\tl{˥˩}},ipa={juɴ},pos={名},acc={F}]}%
	{%
		\MeaningsBlock{%
	\ryumeaning%
		{弓}{}{}%
	}%
	}%
	{\tailPrint[]}%
\entry
	{\headPrint[kana={ゆんぐとぅ\tl{˩˥}},ipa={juŋgutu},pos={名},acc={R}]}%
	{%
		\MeaningsBlock{%
	\ryumeaning%
		{古謡の種類}{}{}%
	}%
	}%
	{\tailPrint[]}%
\entry
	{\headPrint[kana={ゆんさん\tl{˩˥}},ipa={junsaɴ},pos={名},acc={R}]}%
	{%
		\MeaningsBlock{%
	\ryumeaning%
		{巡査。駐在}{}{}%
	}%
	}%
	{\tailPrint[]}%
\entry
	{\headPrint[kana={ゆんた\tl{˩˥}},ipa={junta},pos={名},acc={R}]}%
	{%
		\MeaningsBlock{%
	\ryumeaning%
		{古謡の種類}{}{}%
	}%
	}%
	{\tailPrint[]}%
\entry
	{\headPrint[kana={ゆんたく\tl{˩˥}},ipa={juntaku},pos={名},acc={R}]}%
	{%
		\MeaningsBlock{%
	\ryumeaning%
		{おしゃべり。雑談}{}{}%
	}%
	}%
	{\tailPrint[]}%
\entry
	{\headPrint[kana={ゆんたま},ipa={juntama},pos={名}]}%
	{%
		\MeaningsBlock{%
	\ryumeaning%
		{小魚}{}{}%
	}%
	}%
	{\tailPrint[]}%
\entry
	{\headPrint[kana={ゆんちゅ\tl{˥˩}},ipa={juntʃu},pos={名},acc={F}]}%
	{%
		\MeaningsBlock{%
	\ryumeaning%
		{\ruby[g]{与人}{よんちゅ}}{王府時代の役人の位階名}{}%
	}%
	}%
	{\tailPrint[]}%
\LetterHead{よ}
\entry
	{\headPrint[kana={よー\tl{˥˩}},ipa={joː},pos={副},acc={F}]}%
	{%
		\MeaningsBlock{%
	\ryumeaning%
		{きっと。よくよく}{\emb{よーよー}は強調}{}%
	}%
	}%
	{\tailPrint[]}%
\entry
	{\headPrint[kana={よーがりやん},ipa={joːgarijaɴ}]}%
	{%
		\MeaningsBlock{%
	\ryumeaning%
		{痩せ衰える}{}{}%
	}%
	}%
	{\tailPrint[]}%
\entry
	{\headPrint[kana={よーがりるん\tl{˥˩}},ipa={joːgariruɴ},pos={動},rui={3},para={よーがるぬ},acc={F}]}%
	{%
		\MeaningsBlock{%
	\ryumeaning%
		{痩せる}{}{}%
	\ryumeaning%
		{くびれる}{}{}%
	}%
	}%
	{\tailPrint[]}%
\entry
	{\headPrint[kana={よーがりん\tl{˥˩}},ipa={joːgariɴ},pos={動},rui={3},para={よーがるぬ},acc={F}]}%
	{%
		\MeaningsBlock{%
	\ryumeaning%
		{痩せる。衰弱する}{}{}%
	}%
	}%
	{\tailPrint[]}%
\entry
	{\headPrint[kana={よーに\tl{˩˥}},ipa={joːni},pos={副},acc={R}]}%
	{%
		\MeaningsBlock{%
	\ryumeaning%
		{容易に}{}{}%
	}%
	}%
	{\tailPrint[]}%
\entry
	{\headPrint[kana={よーば\tl{˩˥}},ipa={joːba},pos={名},acc={R}]}%
	{%
		\MeaningsBlock{%
	\ryumeaning%
		{弱虫。虚弱}{}{}%
	}%
	}%
	{\tailPrint[]}%
\entry
	{\headPrint[kana={よー\tl{˩˥}はん},ipa={joːhaɴ},pos={形},para={よーへぬ},acc={R}]}%
	{%
		\MeaningsBlock{%
	\ryumeaning%
		{弱い}{}{}%
	}%
	}%
	{\tailPrint[]}%
\entry
	{\headPrint[kana={よーみるん\tl{˩˥}},ipa={joːmiruɴ},pos={動},rui={3},para={よーむぬ},acc={R}]}%
	{%
		\MeaningsBlock{%
	\ryumeaning%
		{弱める}{}{}%
	}%
	}%
	{\tailPrint[]}%
\entry
	{\headPrint[kana={よーよー},ipa={joːjoː},pos={副}]}%
	{%
		\MeaningsBlock{%
	\ryumeaning%
		{（強調）よくよく}{\emb{よー}の強調形}{}%
	}%
	}%
	{\tailPrint[]}%
\entry
	{\headPrint[kana={よーらさー\tl{˩˥}},ipa={joːrasaː},pos={名},acc={R}]}%
	{%
		\MeaningsBlock{%
	\ryumeaning%
		{ゴイサギ}{鳥の名}{}%
	}%
	}%
	{\tailPrint[]}%
\entry
	{\headPrint[kana={よーるん\tl{˩˥}},ipa={joːruɴ},pos={動},rui={1},para={よーらぬ},acc={R}]}%
	{%
		\MeaningsBlock{%
	\ryumeaning%
		{弱る。弱くなる。衰弱する}{}{}%
	}%
	}%
	{\tailPrint[]}%
\entry
	{\headPrint[kana={よい\tl{˩˥}},ipa={joi},pos={名},acc={R}]}%
	{%
		\MeaningsBlock{%
	\ryumeaning%
		{祝い。祝典}{}{}%
	}%
	}%
	{\tailPrint[]}%
\entry
	{\headPrint[kana={よいぬ\ws{}むぬ},ipa={joinu munu},pos={句}]}%
	{%
		\MeaningsBlock{%
	\ryumeaning%
		{祝いの引き出物}{}{}%
	}%
	}%
	{\tailPrint[]}%
\entry
	{\headPrint[kana={よがーしみるん},ipa={jogaːʃi̥miruɴ},pos={動}]}%
	{%
		\MeaningsBlock{%
	\ryumeaning%
		{休ませる}{}{}%
	}%
	}%
	{\tailPrint[]}%
\entry
	{\headPrint[kana={よがすん\tl{˩˥}},ipa={jogasuɴ},pos={動},rui={2},para={よがはぬ},acc={R}]}%
	{%
		\MeaningsBlock{%
	\ryumeaning%
		{休む。休憩する。休息する}{「寝る」の尊敬語}{}%
	}%
	}%
	{\tailPrint[]}%
\entry
	{\headPrint[kana={よみしゃ\tl{˩˥}はん},ipa={jomiʃahaɴ},pos={形},para={よみしゃへぬ},acc={R}]}%
	{%
		\MeaningsBlock{%
	\ryumeaning%
		{神聖だ。聖なる}{御嶽や拝所などの場所など}{}%
	}%
	}%
	{\tailPrint[note={宮古の jagumi- 「畏れ多い」と同源}]}%
\entry
	{\headPrint[kana={よんなよんな},ipa={jonnajonna},pos={擬}]}%
	{%
		\MeaningsBlock{%
	\ryumeaning%
		{ゆっくり。ゆらりと}{「ゆっくり」の強調}{}%
	}%
	}%
	{\tailPrint[]}%
\LetterHead{ら}
\entry
	{\headPrint[kana={ら},ipa={ra},pos={接尾}]}%
	{%
		\MeaningsBlock{%
	\ryumeaning%
		{～匹}{魚などを数える助数詞。\emb{ぴとーら}「一匹」}{}%
	}%
	}%
	{\tailPrint[]}%
\entry
	{\headPrint[kana={らく\tl{˩˥}},ipa={raku},pos={名},acc={R}]}%
	{%
		\MeaningsBlock{%
	\ryumeaning%
		{楽}{}{}%
	}%
	}%
	{\tailPrint[]}%
\entry
	{\headPrint[kana={らく\tl{˩˥} \ws{}すん\tl{˥˩}},ipa={raku suɴ},pos={句},para={らくさぬ},acc={R F}]}%
	{%
		\MeaningsBlock{%
	\ryumeaning%
		{楽をする}{}{}%
	}%
	}%
	{\tailPrint[]}%
\LetterHead{り}
\entry
	{\headPrint[kana={り},ipa={ri},pos={接尾}]}%
	{%
		\MeaningsBlock{%
	\ryumeaning%
		{～ろ}{命令接辞。\emb{みり}「見ろ」、\emb{ぱり}「走れ」など}{}%
	}%
	}%
	{\tailPrint[]}%
\entry
	{\headPrint[kana={りーぎ\tl{˥˩}},ipa={riːgi},pos={名},acc={F}]}%
	{%
		\MeaningsBlock{%
	\ryumeaning%
		{礼儀}{「礼儀」の転訛}{}%
	}%
	}%
	{\tailPrint[]}%
\entry
	{\headPrint[kana={りくち\tl{˩˥}},ipa={riku̥tʃi},pos={名},acc={R}]}%
	{%
		\MeaningsBlock{%
	\ryumeaning%
		{理屈}{}{}%
	}%
	}%
	{\tailPrint[note={理屈の転訛}]}%
\entry
	{\headPrint[kana={りくちぃむち\tl{˩˥}},ipa={rikutʃimutʃi},pos={名},acc={R}]}%
	{%
		\MeaningsBlock{%
	\ryumeaning%
		{悪知恵の者}{}{}%
	}%
	}%
	{\tailPrint[]}%
\LetterHead{る}
\entry
	{\headPrint[kana={るがい\tl{˥}},ipa={rugai},pos={名},acc={H}]}%
	{%
		\MeaningsBlock{%
	\ryumeaning%
		{\ruby[g]{竜}{りゅう}\ruby[g]{舌}{ぜつ}\ruby[g]{蘭}{らん}}{植物名}{}%
	}%
	}%
	{\tailPrint[]}%
\entry
	{\headPrint[kana={るくず\tl{˩˥}},ipa={rukudzu},pos={名},acc={R}]}%
	{%
		\MeaningsBlock{%
	\ryumeaning%
		{六十}{}{}%
	}%
	}%
	{\tailPrint[]}%
\LetterHead{わ}
\entry
	{\headPrint[kana={わー\tl{˥}},ipa={waː},pos={名},acc={H}]}%
	{%
		\MeaningsBlock{%
	\ryumeaning%
		{御嶽。拝所}{}{}%
	}%
	}%
	{\tailPrint[]}%
\entry
	{\headPrint[kana={わーわー},ipa={waːwaː},pos={擬}]}%
	{%
		\MeaningsBlock{%
	\ryumeaning%
		{わあわあ}{大声で泣くさま}{}%
	}%
	}%
	{\tailPrint[]}%
\entry
	{\headPrint[kana={わざ\tl{˩˥}},ipa={wadza},pos={名},acc={R}]}%
	{%
		\MeaningsBlock{%
	\ryumeaning%
		{技。技術}{}{}%
	\ryumeaning%
		{業。行為}{}{}%
	}%
	}%
	{\tailPrint[]}%
\entry
	{\headPrint[kana={わしき\tl{˥˩}},ipa={waʃi̥ki},pos={名},acc={F}]}%
	{%
		\MeaningsBlock{%
	\ryumeaning%
		{天気。天候}{}{}%
	}%
	}%
	{\tailPrint[]}%
\entry
	{\headPrint[kana={わしく\tl{˩˥}},ipa={waʃi̥ku},pos={名},acc={R}]}%
	{%
		\MeaningsBlock{%
	\ryumeaning%
		{悪さ。悪戯。妨害}{}{}%
	}%
	}%
	{\tailPrint[]}%
\entry
	{\headPrint[kana={わっか\tl{˥˩}\ws{}んぐん\tl{˥˩}},ipa={wakka ŋguɴ},pos={句},para={わっかんがぬ},acc={F F}]}%
	{%
		\MeaningsBlock{%
	\ryumeaning%
		{追いかける}{}{}%
	}%
	}%
	{\tailPrint[]}%
\entry
	{\headPrint[kana={わっか\tl{˥˩}\ws{}んだすん\tl{˥}},ipa={wakka ndasuɴ},pos={句},acc={F H}]}%
	{%
		\MeaningsBlock{%
	\ryumeaning%
		{追い出す}{}{}%
	}%
	}%
	{\tailPrint[]}%
\entry
	{\headPrint[kana={わっきらいるん\tl{˥˩}},ipa={wakkirairuɴ},pos={動},rui={3},acc={F}]}%
	{%
		\MeaningsBlock{%
	\ryumeaning%
		{追われる}{}{}%
	}%
	}%
	{\tailPrint[]}%
\entry
	{\headPrint[kana={わっきるん\tl{˥˩}},ipa={wakkiruɴ},pos={動},rui={3},para={わっくぬ},acc={F}]}%
	{%
		\MeaningsBlock{%
	\ryumeaning%
		{追う。追い払う}{}{}%
	}%
	}%
	{\tailPrint[]}%
\entry
	{\headPrint[kana={わっきん\tl{˥˩}},ipa={wakkiɴ},pos={動},rui={3},para={わっくぬ},acc={F}]}%
	{%
		\MeaningsBlock{%
	\ryumeaning%
		{追う。追い払う}{}{}%
	}%
	}%
	{\tailPrint[]}%
\entry
	{\headPrint[kana={わっさ\tl{˩˥}はん},ipa={wassahaɴ},pos={形},para={わっさへぬ},acc={R}]}%
	{%
		\MeaningsBlock{%
	\ryumeaning%
		{悪い}{}{}%
	}%
	}%
	{\tailPrint[]}%
\entry
	{\headPrint[kana={わりあてぃん\tl{˥}},ipa={wari.atiɴ},pos={動},rui={3},para={わりあとぅぬ},acc={H}]}%
	{%
		\MeaningsBlock{%
	\ryumeaning%
		{割り当てる。振り当てる}{}{}%
	}%
	}%
	{\tailPrint[]}%
\LetterHead{ん}
\entry
	{\headPrint[kana={んーむん\tl{˥}},ipa={ɴɴmuɴ},pos={動},rui={1},para={うーまぬ},acc={H}]}%
	{%
		\MeaningsBlock{%
	\ryumeaning%
		{績む}{}{}%
	}%
	}%
	{\tailPrint[]}%
\entry
	{\headPrint[kana={んーんー},ipa={ɴɴ.ɴɴ},pos={感}]}%
	{%
		\MeaningsBlock{%
	\ryumeaning%
		{ううんううん}{うなる声の形容}{}%
	}%
	}%
	{\tailPrint[]}%
\entry
	{\headPrint[kana={んぎちげーるん\tl{˥˩}},ipa={ŋgitʃigeːruɴ},pos={動},rui={1},para={んぎちげーらぬ},acc={F}]}%
	{%
		\MeaningsBlock{%
	\ryumeaning%
		{行き違う。すれ違う}{}{}%
	}%
	}%
	{\tailPrint[]}%
\entry
	{\headPrint[kana={んくむん\tl{˥˩}},ipa={ŋkumuɴ},pos={動},rui={1},para={んくまぬ},acc={F}]}%
	{%
		\MeaningsBlock{%
	\ryumeaning%
		{力む。いきむ}{}{}%
	}%
	}%
	{\tailPrint[]}%
\entry
	{\headPrint[kana={んぐん\tl{˥˩}},ipa={ŋguɴ},pos={動},rui={1},para={んがぬ},acc={F}]}%
	{%
		\MeaningsBlock{%
	\ryumeaning%
		{行く。去る}{}{}%
	\ryumeaning%
		{時が経つ}{}{}%
	}%
	}%
	{\tailPrint[]}%
\entry
	{\headPrint[kana={んげーるん\tl{˥˩}},ipa={ŋgeːruɴ},pos={動},rui={1},para={んげーらぬ},acc={F}]}%
	{%
		\MeaningsBlock{%
	\ryumeaning%
		{召し上がる}{「食べる」の尊敬語}{}%
	}%
	}%
	{\tailPrint[]}%
\entry
	{\headPrint[kana={ん\josho{}ごー\kako{}び},ipa={ŋgoːbi},pos={副},acc={LHL}]}%
	{%
		\MeaningsBlock{%
	\ryumeaning%
		{たくさん。一杯}{}{}%
	}%
	}%
	{\tailPrint[]}%
\entry
	{\headPrint[kana={んさ\tl{˩˥}はん},ipa={nsahaɴ},pos={形},para={んさへぬ},acc={R}]}%
	{%
		\MeaningsBlock{%
	\ryumeaning%
		{重い。重たい}{}{}%
	}%
	}%
	{\tailPrint[]}%
\entry
	{\headPrint[kana={んじたつん\tl{˥}},ipa={ndʒitḁtsuɴ},pos={動},para={んじたたぬ},acc={H}]}%
	{%
		\MeaningsBlock{%
	\ryumeaning%
		{出で立つ。出発する}{}{}%
	}%
	}%
	{\tailPrint[]}%
\entry
	{\headPrint[kana={んじふち\tl{˥}},ipa={ndʒifu̥tʃi},pos={名},acc={H}]}%
	{%
		\MeaningsBlock{%
	\ryumeaning%
		{出口}{}{}%
	}%
	}%
	{\tailPrint[]}%
\entry
	{\headPrint[kana={んじふに\tl{˥}},ipa={ndʒifu̥ni},pos={名},acc={H}]}%
	{%
		\MeaningsBlock{%
	\ryumeaning%
		{出船}{}{}%
	}%
	}%
	{\tailPrint[]}%
\entry
	{\headPrint[kana={んじるん\tl{˥}},ipa={ndʒiruɴ},pos={動},rui={3}]}%
	{%
		\MeaningsBlock{%
	\ryumeaning%
		{出る。出席する}{}{}%
	}%
	}%
	{\tailPrint[]}%
\entry
	{\headPrint[kana={んじん\tl{˥}},ipa={ndʒiɴ},pos={動},rui={3},para={んどぅぬ},acc={H}]}%
	{%
		\MeaningsBlock{%
	\ryumeaning%
		{出る。出席する}{}{}%
	}%
	}%
	{\tailPrint[]}%
\entry
	{\headPrint[kana={んた\tl{˩˥}},ipa={nta},pos={名},acc={R}]}%
	{%
		\MeaningsBlock{%
	\ryumeaning%
		{土。粘土}{}{}%
	}%
	}%
	{\tailPrint[]}%
\entry
	{\headPrint[kana={んだ},ipa={nda},pos={接尾}]}%
	{%
		\MeaningsBlock{%
	\ryumeaning%
		{～\ruby[g]{達}{たち}}{複数の人を示す接尾語。\emb{いやんだ}「父親たち」など}{}%
	}%
	}%
	{\tailPrint[]}%
\entry
	{\headPrint[kana={んだすん\tl{˥}},ipa={ndasuɴ},pos={動},rui={2},para={んだはぬ},acc={H}]}%
	{%
		\MeaningsBlock{%
	\ryumeaning%
		{出す。差し出す}{}{}%
	}%
	}%
	{\tailPrint[]}%
\entry
	{\headPrint[kana={んたちるん\tl{˥˩}},ipa={ntatʃiruɴ},pos={動},acc={F}]}%
	{%
		\MeaningsBlock{%
	\ryumeaning%
		{（全部）\ruby[g]{零}{こぼ}す}{}{}%
	}%
	}%
	{\tailPrint[]}%
\entry
	{\headPrint[kana={んちん\tl{˩˥}},ipa={ntʃiɴ},pos={動},rui={3},para={んつぬ},acc={R}]}%
	{%
		\MeaningsBlock{%
	\ryumeaning%
		{満ちる。一杯になる}{}{}%
	}%
	}%
	{\tailPrint[]}%
\entry
	{\headPrint[kana={んつぁすん\tl{˩˥}},ipa={ntsasuɴ},pos={動},rui={2},para={んつぁはぬ},acc={R}]}%
	{%
		\MeaningsBlock{%
	\ryumeaning%
		{満たす}{}{}%
	}%
	}%
	{\tailPrint[]}%
\entry
	{\headPrint[kana={んな\tl{˥˩}},ipa={nna},pos={名},acc={F}]}%
	{%
		\MeaningsBlock{%
	\ryumeaning%
		{空っぽ}{}{}%
	}%
	}%
	{\tailPrint[]}%
\entry
	{\headPrint[kana={んなあわり\tl{˥˩}},ipa={nna.awari},pos={名},acc={F}]}%
	{%
		\MeaningsBlock{%
	\ryumeaning%
		{無駄骨。徒労}{}{}%
	}%
	}%
	{\tailPrint[]}%
\entry
	{\headPrint[kana={んなすー\tl{˥˩}},ipa={nnassu},pos={名},acc={F}]}%
	{%
		\MeaningsBlock{%
	\ryumeaning%
		{中味のない汁物}{}{}%
	}%
	}%
	{\tailPrint[]}%
\entry
	{\headPrint[kana={んなふち\tl{˥˩}},ipa={nnafu̥tʃi},pos={名},acc={F}]}%
	{%
		\MeaningsBlock{%
	\ryumeaning%
		{無駄口}{}{}%
	}%
	}%
	{\tailPrint[]}%
\entry
	{\headPrint[kana={んなん},ipa={nnaɴ},pos={名}]}%
	{%
		\MeaningsBlock{%
	\ryumeaning%
		{遺言}{}{}%
	}%
	}%
	{\tailPrint[]}%
\entry
	{\headPrint[kana={んば},ipa={mba},pos={感}]}%
	{%
		\MeaningsBlock{%
	\ryumeaning%
		{\ruby[g]{嫌}{いや}}{\emb{あーい}より否定の程度が強い}{}%
	}%
	}%
	{\tailPrint[]}%
\entry
	{\headPrint[kana={んぶすん\tl{˥}},ipa={mbusuɴ},pos={動},rui={2b},para={んぶさぬ}]}%
	{%
		\MeaningsBlock{%
	\ryumeaning%
		{煮しめる。味付けをして煮る}{}{}%
	}%
	}%
	{\tailPrint[]}%
\entry
	{\headPrint[kana={んま\tl{˥}},ipa={mma},pos={名},acc={H}]}%
	{%
		\MeaningsBlock{%
	\ryumeaning%
		{\ruby[g]{午}{うま}}{十二支の午}{}%
	}%
	}%
	{\tailPrint[]}%
\entry
	{\headPrint[kana={んまん\tl{˥}},ipa={mmaɴ},pos={名},acc={H}]}%
	{%
		\MeaningsBlock{%
	\ryumeaning%
		{馬}{}{}%
	}%
	}%
	{\tailPrint[]}%
 % new settings    
  \end{multicols*}
  \clearpage
\endgroup

% =========================
% ---- REVERSE LOOKUP ----
% =========================
%\pagestyle{empty}
%\cleardoublepage 

% Add title of dictionary part (if necessary)
\addcontentsline{toc}{part}{日本語引き}\pagestyle{empty}\thispagestyle{empty}
\part*{日本語引き}

\begingroup
  % Typesetting parameters for dictionary part
  \ltjsetparameter{
    kanjiskip  = 0pt plus 0.02em minus 0.08em,
    xkanjiskip = 0.10em plus 0.03em minus 0.08em
  }
  \pagestyle{dictionary}
  \phantomsection
  \addcontentsline{toc}{chapter}{日本語・琉球}
  \SetDictBlockWidth
  
  % Suppress paragraph indent for the entries
  \setlength{\parindent}{0pt}
  
  % Set the size of paragraph skip
  \setlength{\parskip}{0pt}
  
  \begin{multicols*}{2}
   \rentry{ああ}{嗚呼}{あー}
\rentry{ああとうとし}{ああ尊し}{うーとーとぅ}
\rentry{あい}{藍}{あい\tl{˥˩}}
\rentry{あいご}{アイゴ}{えぬふぁ\tl{˥˩}}
\rentry{あいさつ}{挨拶}{あいさち\tl{˥˩}、すさりぶち\tl{˥˩}}
\rentry{あいず}{合図}{あいず\tl{˥}}
\rentry{あいぞめ}{藍染}{あいずみ}
\rentry{あいだ}{間}{ふたびし\tl{˥˩}}
\rentry{あいつら}{}{うしたま\tl{˥}}
\rentry{あいて}{相手}{あいてぃ\tl{˥}}
\rentry{あいぼう}{相棒}{ぐー\tl{˩˥}}
\rentry{あいま}{}{あいま\tl{˥}}
\rentry{あいま}{合間}{ふたびし\tl{˥˩}}
\rentry{あいらしい}{愛らしい}{かな\tl{˥}さーん}
\rentry{あえもの}{和え物}{まんず}
\rentry{あえる}{会える}{あいるん\tl{˥}}
\rentry{あえる}{和える}{あいるん\tl{˥}}
\rentry{あえる}{逢える}{あいるん\tl{˥}}
\rentry{あおあおと}{青々と}{お\josho{}ー\kako{}し}
\rentry{あおがんぴ}{青雁皮}{ふがじ\tl{˥}}
\rentry{あおぐ}{扇ぐ}{おんぐん\tl{˥}}
\rentry{あおくさい}{青臭い}{おーふつぁ\tl{˥˩}はーん}
\rentry{あおくなる}{青くなる}{おーむん\tl{˥˩}}
\rentry{あおだいしょう}{青大将}{おーじゃなー\tl{˥}}
\rentry{あおばと}{青鳩}{おーぱーと\tl{˥}}
\rentry{あおばな}{青洟}{ぱなだり\tl{˥˩}}
\rentry{あおばむ}{青ばむ}{おー\kako{}し\ws{}なるん\tl{˩˥}、おーむん\tl{˥˩}}
\rentry{あおむ}{青む}{おーむん\tl{˥˩}}
\rentry{あおむく}{仰向く}{あばんぐん\tl{˥˩}}
\rentry{あおむけになる}{仰向けになる}{あばんぐん\tl{˥˩}}
\rentry{あか}{垢}{ふたり\tl{˥}}
\rentry{あか}{淦}{あか\tl{˥}}
\rentry{あかい}{赤い}{あが\tl{˥˩}はん}
\rentry{あかいねんど}{赤い粘土}{あが\tl{˥˩}んた\tl{˩˥}}
\rentry{あかいろ}{赤色}{あかーいる}
\rentry{あかく}{赤く}{あ\josho{}かー\kako{}し}
\rentry{あかくする}{赤くする}{あがますん\tl{˥˩}}
\rentry{あかくなる}{赤くなる}{あ\josho{}かー\kako{}し\ws{}なるん\tl{˩˥}、あかー\ws{}すん、あがむん\tl{˥˩}}
\rentry{あかご}{赤子}{あがたま\tl{˥}、あがんたま\tl{˥}}
\rentry{あかす}{明かす}{〔夜を\ntilde{}〕あがらすん\tl{˥˩}}
\rentry{あかだに}{赤ダニ}{あがだん\tl{˥˩}}
\rentry{あかつき}{暁}{あかしきん\tl{˥}}
\rentry{あかつち}{赤土}{あが\tl{˥˩}んた\tl{˩˥}}
\rentry{あかてつ}{アカテツ}{とぅのー\tl{˥}}
\rentry{あがめる}{崇める}{あがみるん\tl{˥}}
\rentry{あからむ}{赤らむ}{あ\josho{}かー\kako{}し\ws{}なるん\tl{˩˥}、あかー\ws{}すん、あがむん\tl{˥˩}}
\rentry{あかり}{灯り}{あがり\tl{˥˩}、とぅり\tl{˥}}
\rentry{あがる}{上がる}{あがるん\tl{˥˩}、ぬぶるん\tl{˥˩}}
\rentry{あかるい}{明るい}{あがりゃん}
\rentry{あかるくする}{明るくする}{あがらすん\tl{˥˩}}
\rentry{あかるくなる}{明るくなる}{あがるん\tl{˥˩}}
\rentry{あかんぼう}{赤ん坊}{あがたま\tl{˥}、あがんたま\tl{˥}}
\rentry{あき}{秋}{すさんち\tl{˥}}
\rentry{あきや}{空家}{あぎひ\tl{˥˩}}
\rentry{あきやしき}{空き屋敷}{あぎやしき\tl{˥˩}}
\rentry{あきらめる}{諦める}{うむいきすん\tl{˥}}
\rentry{あきる}{飽きる}{あきりるん\tl{˥˩}、あきりん\tl{˥˩}}
\rentry{あきれはてる}{呆れ果てる}{あきるん\tl{˥˩}}
\rentry{あきれる}{呆れる}{あきるん\tl{˥˩}}
\rentry{あく}{開く}{あぐん\tl{˥˩}}
\rentry{あくじ}{悪事}{やなくとぅ\tl{˩˥}}
\rentry{あくしゅう}{悪臭}{ふつぁりかー、やなかー\tl{˥˩}}
\rentry{あくてんこう}{悪天候}{やなわしき\tl{˥˩}}
\rentry{あくどう}{悪童}{やなわらび\tl{˩˥}}
\rentry{あくへき}{悪癖}{やなくしー\tl{˥˩}}
\rentry{あくりょうばらいのひとつ}{悪霊払いの一つ}{いたすきばら\tl{˥}}
\rentry{あけがた}{明け方}{あぎさり}
\rentry{あけのみょうじょう}{明けの明星}{あかしきぶし}
\rentry{あける}{明ける}{あがるん\tl{˥˩}}
\rentry{あける}{空ける}{いたちるん}
\rentry{あける}{開ける}{〔戸を\ntilde{}〕あぎるん\tl{˥˩}、あぎん\tl{˥˩}}
\rentry{あげる}{}{ひるん\tl{˥}}
\rentry{あげる}{上げる}{あんぎるん\tl{˥˩}}
\rentry{あげる}{揚げる}{〔油で\ntilde{}〕あぎるん\tl{˥˩}}
\rentry{あご}{顎}{はこち\tl{˥˩}}
\rentry{あこう}{アコウ}{あごん\tl{˥˩}}
\rentry{あさ}{朝}{すとぅむち\tl{˥}}
\rentry{あさい}{浅い}{あさ\tl{˥˩}はん}
\rentry{あざける}{嘲る}{あなどるん\tl{˥}、ばくるん\tl{˩˥}}
\rentry{あさせ}{浅瀬}{〔磯に続く\ntilde{}〕びしょー\tl{˥˩}}
\rentry{あさって}{明後日}{あすぅとぅ\tl{˥}、あつぁあすとぅ\tl{˥}}
\rentry{あさって}{近いうち}{あつぁあすとぅ\tl{˥}}
\rentry{あさつゆ}{朝露}{ゆす\tl{˩˥}}
\rentry{あざなう}{}{あんじるん\tl{˥}}
\rentry{あさなぎ}{朝凪}{あさどぅり\tl{˥}}
\rentry{あざみ}{薊}{あざん\tl{˥˩}}
\rentry{あざみさんご}{アザミサンゴ}{すぷるいし\tl{˥}}
\rentry{あざむく}{欺く}{あぎまーすん\tl{˥˩}}
\rentry{あさる}{漁る}{あさるん\tl{˥}、ふつるん\tl{˥}}
\rentry{あし}{足}{ぱん\tl{˥}}
\rentry{あじ}{味}{あじ\tl{˥˩}}
\rentry{あじ}{按司}{あざまぐ\tl{˥}}
\rentry{あしあと}{足跡}{ぴょん\tl{˥˩}}
\rentry{あしがおそい}{足が遅い}{ぱんぬ\ws{}にふつぁー}
\rentry{あしくび}{足首}{ぱんぬ\ws{}ふき}
\rentry{あした}{明日}{あっつぁー\tl{˥}}
\rentry{あしたのあさ}{明日の朝}{あつぁすとぅむち\tl{˥}}
\rentry{あしたのよる}{明日の夜}{あつぁゆ\tl{˥}}
\rentry{あじつけをしてにる}{味付けをして煮る}{んぶすん\tl{˥}}
\rentry{あしのこう}{足の甲}{ぱんぬ\ws{}こー、ぱんぬ\ws{}ぴぃさ、ぴぃさ\tl{˥˩}}
\rentry{あじわう}{味わう}{あじ\ws{}すん}
\rentry{あす}{}{あっつぁー\tl{˥}}
\rentry{あずかる}{預かる}{あすかるん\tl{˥}}
\rentry{あずき}{小豆}{あがまーみ\tl{˥˩}}
\rentry{あずける}{預ける}{あしきるん\tl{˥}}
\rentry{あせ}{汗}{あし\tl{˥}}
\rentry{あぜ}{畦}{あぶし\tl{˥}}
\rentry{あせくさい}{汗臭い}{あし\tl{˥}ふさ\tl{˥}はーん}
\rentry{あせばむ}{汗ばむ}{あし\ws{}ふきん}
\rentry{あせみず}{汗水}{あしみじ\tl{˥}}
\rentry{あせも}{汗疹}{あしゃぼー\tl{˥}}
\rentry{あせをかく}{汗をかく}{あし\ws{}ふきん}
\rentry{あそこ}{}{はー}
\rentry{あそこに}{}{はな\tl{˥}}
\rentry{あそび}{遊び}{あしぴ\tl{˥˩}}
\rentry{あたえる}{与える}{とらしみん、ひるん\tl{˥}}
\rentry{あたたかい}{暖かい}{のーさ\tl{˩˥}はん}
\rentry{あたたかく}{温かく}{あ\josho{}てー\kako{}し}
\rentry{あたたかくなる}{暖かくなる}{のーさ\tl{˩˥}\ws{}なるん\tl{˩˥}}
\rentry{あたたまる}{暖まる}{のーさ\tl{˩˥}\ws{}なるん\tl{˩˥}}
\rentry{あたたまる}{温まる}{あつぁは\tl{˥}\ws{}なるん\tl{˩˥}、ぬるみゃん\tl{˩˥}}
\rentry{あたためる}{温める}{あつぁーすん\tl{˥}、〔食べ物を\ntilde{}〕あっつぁすん\tl{˥}}
\rentry{あたま}{頭}{あまっすくる\tl{˥}、かしら\tl{˥}}
\rentry{あたまがいたい}{頭が痛い}{あまっすくる\tl{˥}やみ\tl{˩˥}}
\rentry{あたらしい}{新しい}{みー\tl{˥˩}}
\rentry{あたらしいもの}{新しい物}{みーむん\tl{˥˩}、めーむぬ\tl{˥˩}}
\rentry{あたり}{当たり}{あたりゃん}
\rentry{あたる}{当たる}{あたるん\tl{˥˩}、〔風に\ntilde{}〕ふかりるん\tl{˥}}
\rentry{あだん}{阿旦}{あんだに\tl{˥}}
\rentry{あだんのきこん}{アダンの気根}{あんだしぃ\tl{˥}}
\rentry{あちら}{}{はー}
\rentry{あちらこちら}{}{あま\tl{˥}くま\tl{˥˩}}
\rentry{あちらに}{}{はな\tl{˥}}
\rentry{あつい}{暑い}{あっつぁ\tl{˥}はん}
\rentry{あつい}{熱い}{あっつぁ\tl{˥}はん}
\rentry{あつがりだ}{暑がりやだ}{あつぁむさ\tl{˥}はん}
\rentry{あつくなる}{熱くなる}{あつぁは\tl{˥}\ws{}なるん\tl{˩˥}}
\rentry{あつまらせる}{集まらせる}{するいるん\tl{˥}}
\rentry{あつまる}{集まる}{あすまるん\tl{˥}、ゆるん\tl{˥˩}}
\rentry{あつめる}{集める}{あしみるん\tl{˥}}
\rentry{あつらえる}{誂える}{あちらいん\tl{˥}}
\rentry{あてにする}{当てにする}{あちすん\tl{˥˩}}
\rentry{あてはめる}{当てはめる}{ぱみるん\tl{˥}}
\rentry{あてもの}{当て物}{〔頭上運搬の\ntilde{}〕ししき\tl{˥}}
\rentry{あてる}{当てる}{あちるん\tl{˥˩}}
\rentry{あと}{後}{あとぅ\tl{˥}、しぴ\tl{˥˩}、しぴゃた\tl{˥˩}}
\rentry{あと}{跡}{あとぅ\tl{˥}}
\rentry{あとあと}{後々}{あとぅあとぅ\tl{˩˥}}
\rentry{あとさき}{後先}{あとさき\tl{˥˩}}
\rentry{あとつぎ}{後継ぎ}{あとぅ\tl{˥}ちぎ\tl{˥˩}}
\rentry{あな}{穴}{みん\tl{˩˥}}
\rentry{あながあく}{穴があく}{みんぴきるん}
\rentry{あなどる}{侮る}{あなどるん\tl{˥}、ばくるん\tl{˩˥}}
\rentry{あなをあける}{穴をあける}{ぴくん\tl{˥}}
\rentry{あに}{兄}{しゃま\tl{˥}}
\rentry{あね}{姉}{あま\tl{˥}}
\rentry{あの}{}{うぬ\tl{˥˩}}
\rentry{あのよ}{あの世}{ぐそー\tl{˩˥}}
\rentry{あばらぼね}{肋骨}{やたぶに\tl{˩˥}}
\rentry{あばれもの}{暴れ者}{やまぐ\tl{˩˥}、やまんぐ}
\rentry{あばれる}{暴れる}{ありるん\tl{˥˩}}
\rentry{あばれんぼうだ}{暴れん坊だ}{やまんぐ\tl{˩˥}さーん}
\rentry{あびせる}{浴びせる}{あますん}
\rentry{あひる}{家鴨}{あぴら\tl{˥}}
\rentry{あびる}{浴びる}{あみん\tl{˥˩}}
\rentry{あぶってかんそうする}{焙って乾燥する}{〔火で\ntilde{}〕かんがーすん\tl{˥}}
\rentry{あぶない}{危ない}{ぴこ\tl{˥˩}はん}
\rentry{あぶら}{油}{あば\tl{˥}}
\rentry{あぶらぜみ}{アブラゼミ}{しゃんしゃん\tl{˥}}
\rentry{あぶらっこい}{脂っこい}{あばずー\tl{˥}さーん}
\rentry{あぶらみそ}{油味噌}{あんだみしゅ\tl{˥}}
\rentry{あぶらむし}{油虫}{かむし\tl{˥}}
\rentry{あぶる}{炙る}{あぶるん\tl{˥}}
\rentry{あふれる}{溢れる}{あふらん}
\rentry{あほらしい}{}{ほーらきし\tl{˥}}
\rentry{あまい}{甘い}{あずま\tl{˥˩}はん}
\rentry{あまぐも}{雨雲}{あまふもん\tl{˥}}
\rentry{あまごいきがん}{雨乞い祈願}{あみにげー\tl{˥}}
\rentry{あまごいきがんのかめんしん}{雨乞い祈願の仮面神}{ふつぁまら\tl{˥}}
\rentry{あます}{余す}{あまらすん\tl{˥}}
\rentry{あまったれ}{}{ふんでー\tl{˥˩}}
\rentry{あまど}{雨戸}{やどぅ\tl{˩˥}}
\rentry{あまどい}{雨樋}{はきじゃー\tl{˥}}
\rentry{あまのがわ}{天の川}{じんぬ\tl{˥˩}\ws{}ふかー\tl{˥}}
\rentry{あまもり}{雨漏り}{あまむり\tl{˥}}
\rentry{あまやかす}{甘やかす}{ぷんでー\ws{}すむん}
\rentry{あまらす}{余らす}{あまらすん\tl{˥}}
\rentry{あまり}{余り}{あまり\tl{˥}、おーばー\tl{˥˩}}
\rentry{あまりにひどい}{あまりに酷い}{どぅぐ}
\rentry{あまる}{余る}{あまるん\tl{˥}、のがるん\tl{˩˥}}
\rentry{あみ}{網}{あん\tl{˥}}
\rentry{あみでとる}{網で捕る}{〔魚を\ntilde{}〕はからすん\tl{˥}}
\rentry{あむ}{編む}{あむん\tl{˥}、ふむん\tl{˥}}
\rentry{あめ}{雨}{あみ\tl{˥}}
\rentry{あやかる}{肖る}{あやかーるん\tl{˥}、あやかるん\tl{˥}}
\rentry{あやしい}{怪しい}{あやっ\tl{˥˩}さーん}
\rentry{あやまり}{誤り}{まちげー\tl{˩˥}}
\rentry{あやまる}{誤る}{ばっぺー\ws{}すん}
\rentry{あら}{粗}{あら\tl{˥˩}}
\rentry{あらい}{粗い}{あら\tl{˥˩}はん}
\rentry{あらい}{荒い}{〔波、動作、粉末などが\ntilde{}〕あら\tl{˥˩}はん}
\rentry{あらいことば}{荒い言葉}{あら\tl{˥˩}むに\tl{˩˥}}
\rentry{あらいそ}{荒磯}{すび\tl{˥˩}}
\rentry{あらう}{洗う}{あらすん\tl{˥˩}}
\rentry{あらうみ}{荒海}{あらなん\tl{˥˩}、なんぶらん}
\rentry{あらぐすくじま}{新城島}{ぱなり\tl{˥}、ぱなりぬ\ws{}すま\tl{˥}}
\rentry{あらす}{荒らす}{あらすん\tl{˥˩}}
\rentry{あらそう}{争う}{えんだりるん\tl{˥}、えんだりん\tl{˥}}
\rentry{あらそう}{戦争する}{いふつぁー\tl{˥}\ws{}すん\tl{˥˩}}
\rentry{あらだてる}{荒立てる}{あらだてぃるん\tl{˥˩}}
\rentry{あらためる}{改める}{あらたみるん\tl{˥}}
\rentry{あらなみ}{荒波}{あらなん\tl{˥˩}}
\rentry{あらわれる}{現れる}{あらわりるん\tl{˥}}
\rentry{あり}{蟻}{あーり\tl{˥˩}}
\rentry{ありがたい}{}{ふくら\tl{˥˩}はん、ふこーら\tl{˥˩}はーん}
\rentry{ありがとう}{}{にーはい、ふくら\tl{˥˩}はん}
\rentry{ありがとうございます}{}{にーはいゆー}
\rentry{ある}{有る}{あん\tl{˥}}
\rentry{あるく}{歩く}{あるぐん\tl{˥}}
\rentry{あれ}{}{うり\tl{˥˩}}
\rentry{あれる}{荒れる}{ありるん\tl{˥˩}、〔波が\ntilde{}〕ぶりるん\tl{˩˥}}
\rentry{あわ}{泡}{あすぺー\tl{˥}}
\rentry{あわ}{粟}{あん\tl{˥}}
\rentry{あわせる}{会わせる}{あーしみるん\tl{˥}}
\rentry{あわせる}{合わせる}{あーすん\tl{˥}}
\rentry{あわてさせる}{慌てさせる}{さわがすん\tl{˥˩}}
\rentry{あわてふためく}{慌てふためく}{あばちかんち\tl{˥}\ws{}すん\tl{˥˩}}
\rentry{あわてる}{慌てる}{あばちるん\tl{˥˩}}
\rentry{あわれだ}{哀れだ}{すむやむーん}
\rentry{あん}{餡}{あん\tl{˥}、ばた\tl{˩˥}}
\rentry{あんかだ}{安価だ}{やっさ\tl{˩˥}はん}
\rentry{あんこ}{餡子}{あん\tl{˥}、ばた\tl{˩˥}}
\rentry{あんしょう}{暗礁}{じり\tl{˩˥}}
\rentry{あんしん}{安心}{\josho{}うなー\kako{}ぐ、とぅくっとぅ\tl{˥}}
\rentry{あんしんした}{安心する}{やぶす\ws{}なるん}
\rentry{あんしんする}{安心する}{\josho{}うなー\kako{}ぐ\ws{}なるん\tl{˩˥}、やすんじるん\tl{˩˥}}
\rentry{あんつく}{アンツク}{めった}
\rentry{あんど}{安堵}{\josho{}うなー\kako{}ぐ、とぅくっとぅ\tl{˥}}
\rentry{あんどした}{安堵する}{やぶす\ws{}なるん}
\rentry{あんどする}{安堵する}{\josho{}うなー\kako{}ぐ\ws{}なるん\tl{˩˥}}
\rentry{あんないする}{案内する}{しけーすん\tl{˥˩}}
\rentry{あんまりだ}{}{どぅ\josho{}ぐ\kako{}どぅぐ}
\rentry{あんもち}{餡餅}{あんむち\tl{˥}}
\rentry{い}{亥}{びー\tl{˩˥}}
\rentry{い}{胃}{ぶーばた\tl{˥}}
\rentry{いいあわせる}{言い合わせる}{そんだん\tl{˥}\ws{}すん\tl{˥˩}}
\rentry{いいえ}{}{あーい、あらぬ}
\rentry{いいなおす}{言いなおす}{いーのーすん\tl{˩˥}}
\rentry{いいはる}{言い張る}{やっぱり\tl{˥˩}}
\rentry{いいまかす}{言い負かす}{いーまかすん\tl{˥˩}}
\rentry{いいわけ}{言い訳}{いーばぎ\tl{˥˩}}
\rentry{いう}{言う}{えぬん\tl{˥˩}}
\rentry{いえ}{家}{ひー\tl{˩˥}}
\rentry{いえづくり}{家造り}{ひすくり\tl{˩˥}}
\rentry{いえのしゅうかん}{家の習慣}{やーなれー}
\rentry{いえのらくせいいわい}{家の落成祝い}{ひすくりよい\tl{˩˥}}
\rentry{いえもと}{家元}{やーむとぅ\tl{˩˥}}
\rentry{いかす}{生かす}{いがすん\tl{˥}}
\rentry{いかなる}{如何なる}{ねーしゃるん}
\rentry{いかり}{怒り}{くんぞー\tl{˥}}
\rentry{いき}{息}{いし\tl{˥}}
\rentry{いきおいづく}{勢いづく}{ばざーるん\tl{˩˥}}
\rentry{いきおいよくあふれる}{勢いよく溢れる}{ふきんじるん\tl{˥}}
\rentry{いきおいよくうつすさま}{勢いよく移すさま}{〔液体や小粒の物を\ntilde{}〕いたちり\tl{˥˩}}
\rentry{いきかえる}{生き返る}{ぬちぃむやん}
\rentry{いきがはずむ}{息がはずむ}{あふくん\tl{˥˩}}
\rentry{いきぎれになる}{息切れになる}{あふくん\tl{˥˩}}
\rentry{いきちがう}{行き違う}{んぎちげーるん\tl{˥˩}}
\rentry{いきている}{生きている}{いぎだる}
\rentry{いきむ}{}{んくむん\tl{˥˩}}
\rentry{いきもの}{生き物}{いすむし\tl{˥}}
\rentry{いきりょう}{生霊}{いきろー\tl{˥}}
\rentry{いきる}{生きる}{いぎるん\tl{˥}、いぎん\tl{˥}}
\rentry{いく}{行く}{ぱるん\tl{˥}、んぐん\tl{˥˩}}
\rentry{いく～}{幾～}{うー}
\rentry{いくたり}{}{うたーり\tl{˥˩}}
\rentry{いくつ}{幾つ}{うーちぃ\tl{˥˩}、うーび\tl{˥˩}}
\rentry{いくら}{幾ら}{うーび\tl{˥˩}}
\rentry{いける}{生ける}{〔花などを\ntilde{}〕いぎるん\tl{˥}}
\rentry{いざってすすむ}{いざって進む}{しきっつぁーるん}
\rentry{いさり}{漁り}{いざり\tl{˥˩}}
\rentry{いざる}{}{しきつぁーすん\tl{˥˩}}
\rentry{いし}{石}{いし\tl{˥˩}}
\rentry{いじ}{意地}{いじぃ\tl{˥˩}}
\rentry{いしうす}{石臼}{いしうし\tl{˥˩}}
\rentry{いしがき}{石垣}{いしまーし\tl{˥˩}、ふく\tl{˥}}
\rentry{いしがきし}{石垣市}{しか\tl{˥}}
\rentry{いしがきじま}{石垣島}{いさすぅま\tl{˥}}
\rentry{いしき}{意識}{そー\tl{˥}}
\rentry{いじめる}{虐める}{いじみるん\tl{˥˩}}
\rentry{いしもり}{石盛}{むり\tl{˥˩}}
\rentry{いしゃ}{医者}{いしゃん\tl{˥}}
\rentry{いしょう}{衣装}{いしょー\tl{˥˩}}
\rentry{いしょくする}{移植する}{やとーすん\tl{˩˥}}
\rentry{いじる}{弄る}{だーぶん\tl{˩˥}}
\rentry{いじをはる}{意地を張る}{こっぱるん\tl{˥}}
\rentry{いずみ}{泉}{ばぎみじ\tl{˥˩}、ぱりみじぃ\tl{˥}}
\rentry{いぜん}{以前}{めー\tl{˩˥}}
\rentry{いそ}{礒}{すび\tl{˥˩}}
\rentry{いそいで}{急いで}{ぺーしゃな\tl{˥}}
\rentry{いそがしい}{忙しい}{ぱんたっさ\tl{˥}はん}
\rentry{いそがせる}{急がせる}{あーらすん\tl{˥˩}、いすがすん\tl{˥}}
\rentry{いそぐ}{急ぐ}{あばちるん\tl{˥˩}、いすぐん\tl{˥}}
\rentry{いそべ}{礒辺}{\josho{}ぱま\kako{}ふちぃ}
\rentry{いた}{板}{いた\tl{˥}}
\rentry{いたい}{痛い}{あがー}
\rentry{いたくてかゆいかんじがする}{痛くて痒い感じがする}{〔稲、麦などの芒が皮膚をつきさすような\ntilde{}〕ぱすぅこー\tl{˥}はん}
\rentry{いたずら}{悪戯}{がんまり\tl{˩˥}、わしく\tl{˩˥}}
\rentry{いたずらずきだ}{悪戯好きだ}{やまんぐ\tl{˩˥}さーん}
\rentry{いたずらっこ}{悪戯っ子}{やなわらび\tl{˩˥}}
\rentry{いたずらっぽい}{悪戯っぽい}{がま\tl{˩˥}さーん、がんまり\ws{}しゃーん}
\rentry{いたずらっぽいもの}{悪戯っぽい者}{やまぐ\tl{˩˥}、やまんぐ}
\rentry{いたずらに}{}{あだり\tl{˥˩}}
\rentry{いただき}{頂}{ちじ\tl{˥}}
\rentry{いただきもの}{頂き物}{\josho{}たぼらり\kako{}むぬ}
\rentry{いただく}{頂く}{たぼらいるん\tl{˥}、たぼるん\tl{˥}}
\rentry{いたっ}{}{あがー}
\rentry{いたみなげくさま}{痛み嘆くさま}{あがよー}
\rentry{いたむ}{傷む}{〔肉類が\ntilde{}〕さがるん\tl{˥}}
\rentry{いたむ}{痛む}{やむん\tl{˩˥}}
\rentry{いためる}{痛める}{やますん\tl{˩˥}}
\rentry{いたゆか}{板床}{ふんた\tl{˥˩}}
\rentry{いちごうます}{一合枡}{なもり\tl{˩˥}}
\rentry{いちぞく}{一族}{いちむん\tl{˥˩}}
\rentry{いちだいじ}{一大事}{うーぐとぅ}
\rentry{いちど}{一度}{ぴぃとぅむし\tl{˥}}
\rentry{いちにち}{一日}{ぴてん\tl{˥}}
\rentry{いちにちじゅう}{一日中}{ぴてんぴじゅ}
\rentry{いちばん}{一番}{いすぱん\tl{˥}}
\rentry{いちばんきょうげん}{一番狂言}{いすぱんこんぎ\tl{˥}}
\rentry{いちもくさんに}{一目散に}{びょすた、みっ\josho{}ふぁ\kako{}ん}
\rentry{いちもん}{一門}{いちむん\tl{˥˩}}
\rentry{いつ}{何時}{いちぃ\tl{˥˩}}
\rentry{いっかい}{一回}{ぴぃとぅむし\tl{˥}}
\rentry{いっしゅうき}{一周忌}{なんが\tl{˩˥}}
\rentry{いっしょ}{一緒}{まーじぃ\tl{˥˩}}
\rentry{いっしょうけんめい}{一生懸命}{いざんだ\tl{˩˥}}
\rentry{いっしょうます}{一升枡}{しゃー\tl{˥}}
\rentry{いっしょにいく}{一緒に行く}{そーり\tl{˥˩}\ws{}んぐん\tl{˥˩}}
\rentry{いっしょにくる}{一緒に来る}{そーり\tl{˥˩}\ws{}くん\tl{˥}}
\rentry{いっしょにする}{一緒にする}{まんざーすん}
\rentry{いつつ}{五つ}{いっし\tl{˥}}
\rentry{いってきかせる}{言って聞かせる}{えにすかすん\tl{˥˩}}
\rentry{いっとき}{}{あ\josho{}たー\kako{}すま}
\rentry{いっぱい}{一杯}{まんしん\tl{˥˩}、まんどん\tl{˩˥}、ん\josho{}ごー\kako{}び}
\rentry{いっぱいになる}{一杯になる}{んちん\tl{˩˥}}
\rentry{いっぴん}{一品}{ぴぃとぅすな\tl{˥}}
\rentry{いっぽう}{一方}{ぴぃとぅかた\tl{˥}}
\rentry{いつも}{何時も}{い\josho{}ちぃ\kako{}ん、しゃー\tl{˥}、めー\josho{}が\kako{}めーにち}
\rentry{いつわり}{}{ゆくしむに\tl{˩˥}}
\rentry{いでたつ}{出で立つ}{んじたつん\tl{˥}}
\rentry{いと}{糸}{いとぅ\tl{˥}}
\rentry{いど}{井戸}{けー\tl{˥˩}}
\rentry{いどうする}{移動する}{うがすん\tl{˥}、ゆしるん\tl{˥˩}}
\rentry{いとこ}{従兄弟}{いちふ\tl{˥}}
\rentry{いどころ}{居所}{ぶーどぅり\tl{˩˥}}
\rentry{いどむ}{挑む}{しかきるん\tl{˥˩}}
\rentry{いなくなる}{居なくなる}{のーなさん}
\rentry{いなづま}{稲妻}{ぷちり\tl{˥}}
\rentry{いなびかり}{稲光}{ぷちり\tl{˥}}
\rentry{いなむら}{稲むら}{しら\tl{˥˩}}
\rentry{いぬ}{戌}{いぬ、いん\tl{˥}}
\rentry{いぬ}{犬}{いぬ、いん\tl{˥}}
\rentry{いぬまき}{イヌマキ}{きゃんぎ\tl{˥˩}}
\rentry{いね}{稲}{いに\tl{˥}}
\rentry{いねかり}{稲刈り}{めーかり\tl{˥˩}}
\rentry{いねむり}{居眠り}{にんぶり\tl{˥˩}}
\rentry{いのしし}{猪}{むざ\tl{˥}}
\rentry{いのち}{命}{ぬちぃ\tl{˩˥}}
\rentry{いのちづよい}{命強い}{ぬちぃずー\tl{˩˥}さーん}
\rentry{いのちのおんじん}{命の恩人}{ぬちぃぬうや\tl{˩˥}}
\rentry{いのり}{祈り}{うやっすり\tl{˥}}
\rentry{いのる}{祈る}{にげー\tl{˩˥}しきるん\tl{˥}、にげるん\tl{˩˥}、にんじるん\tl{˥˩}}
\rentry{いはい}{位牌}{いべー\tl{˥˩}}
\rentry{いばる}{威張る}{いばりすくん\tl{˥˩}、いばるん\tl{˥˩}、がーりすきん\tl{˩˥}、たかぶるん\tl{˥}、ほーらし}
\rentry{いびき}{鼾}{ぱな\tl{˥˩}ふき\tl{˩˥}}
\rentry{いびる}{}{はき\tl{˥}まんじるん\tl{˩˥}}
\rentry{いぶくろ}{胃袋}{ぶーばた\tl{˥}}
\rentry{いま}{今}{なま\tl{˥}、まな\tl{˥}}
\rentry{いまごろ}{今頃}{まなばら\tl{˥˩}}
\rentry{いましめる}{戒める}{いましみるん\tl{˥}、いましみん\tl{˥}、すかり\tl{˥˩}すくん\tl{˥}}
\rentry{いますぐ}{今すぐ}{ふたぎな\tl{˥˩}、まなまな}
\rentry{いまのよ}{今の世}{まなぬ\tl{˥}\ws{}ゆー\tl{˥˩}}
\rentry{いむ}{忌}{いみ\tl{˥}}
\rentry{いもうと}{妹}{うとぅーとぅ\tl{˥}}
\rentry{いや}{嫌}{あーい、んば}
\rentry{いやだ}{嫌だ}{ふつくれー}
\rentry{いらっしゃる}{}{おーるん\tl{˥}}
\rentry{いりおもて}{西表}{いりむち\tl{˥˩}}
\rentry{いりおもてじま}{西表島}{いりむち\tl{˥˩}}
\rentry{いりぐち}{入口}{ぺりふちぃ\tl{˥}}
\rentry{いりひ}{入日}{いりしな\tl{˥˩}}
\rentry{いりびたる}{入り浸る}{いりびたるん\tl{˥˩}}
\rentry{いりむこ}{入り婿}{いり\tl{˥˩}むぐ\tl{˩˥}}
\rentry{いる}{射る}{いるん\tl{˥}}
\rentry{いる}{居る}{ぶん\tl{˥˩}}
\rentry{いる}{煎る}{いりぐん\tl{˥}}
\rentry{いる}{要る}{いるん\tl{˥˩}}
\rentry{いるいいれ}{衣類入れ}{なーばぐ\tl{˩˥}}
\rentry{いれかえする}{入れ換えする}{いりかい\tl{˥˩}\ws{}すん\tl{˥˩}}
\rentry{いれかえる}{入れ換える}{いりかい\tl{˥˩}\ws{}すん\tl{˥˩}、いりかいるん\tl{˥˩}}
\rentry{いれかわる}{入れ代わる}{いりかーるん\tl{˥˩}}
\rentry{いれもの}{入れ物}{いりむん\tl{˥˩}}
\rentry{いれもののいっしゅ}{入れ物の一種}{めった}
\rentry{いれる}{入れる}{いりん\tl{˥˩}}
\rentry{いろ}{色}{いる\tl{˥}}
\rentry{いろあせる}{色褪せる}{ぱぎるん\tl{˥}}
\rentry{いろがあせる}{色があせる}{すさむん\tl{˥}}
\rentry{いろがしろい}{色が白い}{いるすそー\tl{˥}はん}
\rentry{いろきちがい}{色気違い}{しきべー}
\rentry{いろり}{囲炉裏}{すか\tl{˥}}
\rentry{いわい}{祝い}{よい\tl{˩˥}}
\rentry{いわいのひきでもの}{祝いの引き出物}{よいぬ\ws{}むぬ}
\rentry{いわがんぺき}{岩岸壁}{すぱに\tl{˥˩}}
\rentry{いわれ}{謂れ}{いわり\tl{˥˩}}
\rentry{いん}{印}{いん\tl{˥}、ぱん\tl{˥˩}}
\rentry{いんかん}{印鑑}{いん\tl{˥}、ぱん\tl{˥˩}}
\rentry{いんけい}{陰茎}{〔幼児語\ntilde{}〕こっち、まら\tl{˩˥}}
\rentry{いんしゅう}{因習}{やななれー\tl{˩˥}}
\rentry{いんしょくしつくす}{飲食し尽くす}{〔飲食物を残らず\ntilde{}〕たいらぎるん\tl{˥}}
\rentry{う}{卯}{うー\tl{˥}}
\rentry{ういきょう}{ウイキョウ}{に\josho{}じ\kako{}きょー}
\rentry{ううんううん}{}{んーんー}
\rentry{うえ}{上}{うい\tl{˥˩}}
\rentry{うえからかける}{上からかける}{かふちるん\tl{˥}}
\rentry{うえへ}{上へ}{ういが\tl{˥˩}}
\rentry{うえる}{植える}{いびるん\tl{˥˩}}
\rentry{うえる}{飢える}{かちりるん\tl{˥}、かちりん\tl{˥}}
\rentry{うおのめ}{魚の目}{ゆぬめー\tl{˥˩}}
\rentry{うかせる}{浮かせる}{おーぎるん\tl{˥˩}}
\rentry{うかぶ}{浮かぶ}{おーがーるん\tl{˥˩}、おーがるん\tl{˥˩}}
\rentry{うかべる}{浮かべる}{おーぎるん\tl{˥˩}}
\rentry{うかれる}{浮かれる}{おーがーるん\tl{˥˩}}
\rentry{うけおう}{請け負う}{うきるん\tl{˥}}
\rentry{うけとる}{受け取る}{うきとぅるん\tl{˥}}
\rentry{うけもつ}{受け持つ}{うきむつん\tl{˥}}
\rentry{うける}{受ける}{うきるん\tl{˥}}
\rentry{うける}{浮ける}{おーぎるん\tl{˥˩}}
\rentry{うごかす}{動かす}{うがすん\tl{˥}}
\rentry{うごく}{動く}{うぐん\tl{˥}}
\rentry{うし}{丑}{うし\tl{˥˩}}
\rentry{うし}{牛}{うし\tl{˥˩}}
\rentry{うじ}{蛆}{うじ\tl{˥}}
\rentry{うじこ}{氏子}{〔御嶽の\ntilde{}〕\josho{}ぱなぬ\kako{}\ws{}ふぁー}
\rentry{うしなう}{失う}{うすなすん\tl{˥˩}、ねーな\tl{˩˥}さん\tl{˥˩}}
\rentry{うしろ}{後ろ}{しなた\tl{˥˩}、しぴ\tl{˥˩}、しぴゃた\tl{˥˩}}
\rentry{うす}{臼}{うし\tl{˥}}
\rentry{うすあじだ}{薄味だ}{あふぁ\tl{˥}さん}
\rentry{うすぐらい}{薄暗い}{あやっふぁ\tl{˥}さーん}
\rentry{うすぐらくなる}{薄暗くなる}{ふぁどぅまるん\tl{˥˩}}
\rentry{うずたかく}{うず高く}{た\josho{}か\kako{}たかし}
\rentry{うずまる}{埋まる}{うずまるん\tl{˥˩}}
\rentry{うすめ}{薄い}{ぴぃさ\tl{˥˩}はん}
\rentry{うすめだ}{薄めだ}{ぴぃさ\tl{˥˩}はん}
\rentry{うすめる}{薄める}{〔水などで\ntilde{}〕ばーすん\tl{˥˩}、ばいるん\tl{˥˩}}
\rentry{うずめる}{埋める}{うずみるん\tl{˥˩}}
\rentry{うずら}{ウズラ}{うずら\tl{˥}}
\rentry{うそ}{}{ゆくしむに\tl{˩˥}}
\rentry{うそ}{嘘}{ぷりむに\tl{˥˩}}
\rentry{うそだ}{嘘だ}{がーん}
\rentry{うそなき}{うそ泣き}{なぎまーべ\tl{˥˩}}
\rentry{うた}{歌}{うた\tl{˥˩}}
\rentry{うたがう}{疑う}{うたげーるん\tl{˥˩}、うたごーん\tl{˥˩}}
\rentry{うたがわしい}{疑わしい}{あやっ\tl{˥˩}さーん}
\rentry{うたき}{御嶽}{わー\tl{˥}}
\rentry{うたきのさいしにさんかするむらびと}{御嶽の祭祀に参加する村人}{やまにんじゅ\tl{˩˥}}
\rentry{うたさんしん}{歌三線}{うたさんしん\tl{˥˩}}
\rentry{うち}{内}{うちな\tl{˥˩}}
\rentry{うちがわ}{内側}{うちな\tl{˥˩}}
\rentry{うちくだく}{打ち砕く}{くだくん}
\rentry{うちこむあめ}{打ち込む雨}{〔室内に\ntilde{}〕うちあみ\tl{˥}}
\rentry{うちつける}{打ち付ける}{うっつぁーすん\tl{˥}}
\rentry{うちとける}{打ち解ける}{すむぴらすん\tl{˥}}
\rentry{うちまかす}{打ち負かす}{うちまかすん}
\rentry{うちみ}{打ち身}{うったち}
\rentry{うちわ}{団扇}{おん\tl{˥}、おんぎ\tl{˥}}
\rentry{うちわすれる}{うち忘れる}{うちばっさん}
\rentry{うつ}{打つ}{うつん\tl{˥}、くらすん\tl{˥˩}、たたぐん\tl{˥}}
\rentry{うつ}{撃つ}{いるん\tl{˥}}
\rentry{うっかり}{}{うかっと}
\rentry{うっかりする}{}{うかっと\ws{}すん\tl{˥˩}}
\rentry{うつくしい}{美しい}{〔容姿が\ntilde{}〕あばりしゃ\tl{˥}はん、けしゃ\tl{˥}はん}
\rentry{うつす}{写す}{うつすん\tl{˥}}
\rentry{うつす}{移す}{〔牛馬を\ntilde{}〕むすなすん\tl{˩˥}}
\rentry{うったえる}{訴える}{うったいん\tl{˥}}
\rentry{うつむく}{俯く}{うすふくん\tl{˥˩}}
\rentry{うつる}{}{〔病気が\ntilde{}〕うつるん\tl{˥}}
\rentry{うつる}{映る}{うつるん\tl{˥}}
\rentry{うつる}{移る}{うつるん\tl{˥}}
\rentry{うで}{腕}{うじ\tl{˥}}
\rentry{うてん}{雨天}{あまおしき\tl{˥}}
\rentry{うなぎ}{ウナギ}{うなん\tl{˥}}
\rentry{うなりがみ}{姉妹神}{ぶなりかん\tl{˥˩}}
\rentry{うなる}{唸る}{たけーるん\tl{˥}}
\rentry{うに}{海栗}{うん\tl{˥˩}、かち\tl{˥}}
\rentry{うぬぼれる}{自惚れる}{ほーらし}
\rentry{うね}{畝}{ぱか\tl{˥}}
\rentry{うばいとる}{奪い取る}{ばがいとぅるん\tl{˩˥}、ばがへ\ws{}とぅるん\tl{˥}}
\rentry{うばう}{奪う}{とぅるん\tl{˥}、ばがすん\tl{˩˥}}
\rentry{うほう}{右方}{ねーり\tl{˥˩}}
\rentry{うま}{午}{んま\tl{˥}}
\rentry{うま}{馬}{んまん\tl{˥}}
\rentry{うまい}{上手い}{ぞーじ\tl{˩˥}}
\rentry{うまい}{旨い}{まー\tl{˥}はん}
\rentry{うまくない}{旨くない}{にしゃ\tl{˥˩}はん}
\rentry{うまくまとまる}{}{とぅとぅのーん\tl{˥}}
\rentry{うまれ}{生まれ}{まり\tl{˥˩}}
\rentry{うまれじま}{生まれ島}{まりじま\tl{˥˩}}
\rentry{うまれどし}{生まれ年}{まりどぅしぃ\tl{˥˩}}
\rentry{うまれる}{生まれる}{まりるん\tl{˥˩}}
\rentry{うみ}{海}{いなー\tl{˥}、いなが\tl{˥}}
\rentry{うみどり}{海鳥}{\josho{}いん\kako{}どぅり}
\rentry{うみべ}{海辺}{\josho{}いなん\kako{}ぱた}
\rentry{うみへのおりぐち}{海への降り口}{うだち\tl{˥}}
\rentry{うむ}{産む}{なすん\tl{˩˥}}
\rentry{うむ}{績む}{んーむん\tl{˥}}
\rentry{うむ}{膿む}{みゃん}
\rentry{うめあわせる}{埋め合わせる}{たらーすん\tl{˥˩}}
\rentry{うもれる}{埋もれる}{うずまるん\tl{˥˩}}
\rentry{うやまう}{敬う}{あがみるん\tl{˥}、うやまいるん\tl{˥˩}、うやめーるん\tl{˥˩}}
\rentry{うようよ}{}{うじゃうじゃ}
\rentry{うらがえす}{裏返す}{うらがいすん\tl{˥}、うらげーすん\tl{˥}}
\rentry{うらさびしい}{うら寂しい}{すとぅる\tl{˥}さーん}
\rentry{うらないし}{占い師}{さんぎんそー\tl{˥}、ゆた\tl{˥˩}}
\rentry{うらむ}{恨む}{にたむん\tl{˩˥}}
\rentry{うらやましい}{羨ましい}{うらみしゃはん、べーちゃー}
\rentry{うる}{売る}{かしみるん\tl{˥˩}}
\rentry{うるうづき}{閏月}{ゆりしき\tl{˥˩}}
\rentry{うるうどし}{閏年}{ゆりどぅし\tl{˥˩}}
\rentry{うるおい}{潤い}{うりー\tl{˥}}
\rentry{うるちまい}{うるち米}{むすめー\tl{˥˩}}
\rentry{うれしい}{嬉しい}{さにしゃ\tl{˥˩}はん}
\rentry{うれしがる}{嬉しがる}{さにしゃー\ws{}すん}
\rentry{うれる}{熟れる}{みゃん\tl{˥}}
\rentry{うろこ}{鱗}{いりぎ\tl{˥}}
\rentry{うろたえる}{}{ざまどぅるん\tl{˩˥}、どまんぐるん\tl{˩˥}}
\rentry{うわきおんな}{浮気女}{さんごなー\tl{˥}}
\rentry{うわさ}{噂}{さた\tl{˥}}
\rentry{うわやく}{上役}{ういぬ\tl{˥˩}\ws{}ぴぃとぅ\tl{˥˩}}
\rentry{うん}{運}{うん\tl{˥˩}、ふん\tl{˥}}
\rentry{うんがない}{運がない}{うんぬ\tl{˥˩}\ws{}ねーぬ\tl{˩˥}}
\rentry{え}{絵}{えー\tl{˥˩}}
\rentry{えいようをつける}{栄養をつける}{くんき\ws{}しきるん}
\rentry{えきしゃ}{易者}{さんぎんそー\tl{˥}、ゆた\tl{˥˩}}
\rentry{えきびょうよけのぎょうじ}{疫病除けの行事}{しまふさら\tl{˥}}
\rentry{えさ}{餌}{むんだに\tl{˩˥}}
\rentry{えだは}{枝葉}{ゆだぱ\tl{˥˩}}
\rentry{えび}{海老}{いび\tl{˥˩}}
\rentry{えびずる}{}{かなぶ\tl{˥˩}}
\rentry{えもの}{獲物}{てぃがら\tl{˥}}
\rentry{えものがおおい}{獲物が多い}{だーさ\tl{˩˥}はん、まだ\tl{˥˩}はん}
\rentry{えらいひと}{偉い人}{たかぴぃとぅ\tl{˥}、ぶーぴぃとぅ}
\rentry{えらぶ}{選ぶ}{いらぶん\tl{˥}}
\rentry{えんきする}{延期する}{ぬばすん\tl{˩˥}}
\rentry{えんさい}{エンサイ}{うんつぇー}
\rentry{えんぽう}{遠方}{とぅけー\tl{˥˩}}
\rentry{えんぽうだ}{遠方だ}{とぅさ\tl{˥˩}はん}
\rentry{お}{尾}{ぶすぷ\tl{˩˥}}
\rentry{おーい}{}{やー}
\rentry{おあがりになる}{お上がりになる}{おーし\ws{}おるん}
\rentry{おい}{甥}{ぶい\tl{˥˩}、ぶいふぁー\tl{˥˩}}
\rentry{おい}{老い}{うい\tl{˥}}
\rentry{おいかける}{追いかける}{わっか\tl{˥˩}\ws{}んぐん\tl{˥˩}}
\rentry{おいこす}{追い越す}{ういぬぐん\tl{˥˩}}
\rentry{おいこむ}{追い込む}{ういくむん\tl{˥}}
\rentry{おいしい}{美味しい}{まー\tl{˥}はん}
\rentry{おいしげる}{生い茂る}{かぷん\tl{˥}、むい\tl{˩˥}かぷん\tl{˥}}
\rentry{おいだす}{追い出す}{ういんだすん\tl{˥}、わっか\tl{˥˩}\ws{}んだすん\tl{˥}}
\rentry{おいたてる}{追い立てる}{ういまーすん\tl{˥˩}、だふつるん\tl{˩˥}}
\rentry{おいちらす}{追い散らす}{ういわっきるん\tl{˥˩}、だふつり\tl{˩˥}わっきん\tl{˥˩}}
\rentry{おいつく}{追いつく}{ういすくん\tl{˥˩}}
\rentry{おいておく}{置いておく}{すくん\tl{˥}}
\rentry{おいぬく}{追い抜く}{ういぬぐん\tl{˥˩}}
\rentry{おいはらう}{追い払う}{わっきるん\tl{˥˩}、わっきん\tl{˥˩}}
\rentry{おいぼれる}{老いぼれる}{ういぶり\tl{˥}}
\rentry{おいまわす}{追い回す}{ういまーすん\tl{˥˩}}
\rentry{おいめ}{負い目}{うか\tl{˥}}
\rentry{おいる}{老いる}{ういるん\tl{˥}}
\rentry{おいること}{老いること}{うい\tl{˥}}
\rentry{おう}{追う}{わっきるん\tl{˥˩}、わっきん\tl{˥˩}}
\rentry{おうぎ}{扇}{おん\tl{˥}、おんぎ\tl{˥}}
\rentry{おうし}{雄牛}{〔大きな\ntilde{}〕ぐつぇー\tl{˩˥}、びーうしぃ\tl{˥˩}}
\rentry{おうちゃくなもの}{横着な者}{なまふつぁりむぬ\tl{˩˥}、やまぐ\tl{˩˥}、やまんぐ}
\rentry{おうべいじん}{欧米人}{うらんだー\tl{˥}}
\rentry{おえる}{終える}{しー\ws{}うすくまん、すますん\tl{˥}、とぅじみるん\tl{˥˩}}
\rentry{おお}{大}{ぶー}
\rentry{おおあわてする}{大慌てする}{あばちかんち\tl{˥}\ws{}すん\tl{˥˩}}
\rentry{おおい}{多い}{ぶさ\tl{˥}はん}
\rentry{おおう}{覆う}{うすん\tl{˥˩}、かふちるん\tl{˥}}
\rentry{おおかぜ}{大風}{ぶーかち\tl{˥}}
\rentry{おおがたのえい}{大型エイ}{かまんた\tl{˥}}
\rentry{おおがたのかつお}{大型のカツオ}{だいばん\tl{˩˥}}
\rentry{おおかわ}{大川}{ふーがー\tl{˥}}
\rentry{おおきい}{大きい}{ぶー、ぶさ\tl{˥}はん、まーぎ}
\rentry{おおきくなる}{大きくなる}{ぶさは\tl{˥}\ws{}なるん\tl{˩˥}}
\rentry{おおきなうねり}{大きなうねり}{むさん\tl{˥˩}}
\rentry{おおきなかなてこ}{大きな金梃子}{しんがら\tl{˥}}
\rentry{おおきなは}{大きな葉}{かなぱ\tl{˥˩}}
\rentry{おおきなひと}{大きな人}{ぶーぴぃとぅ}
\rentry{おおくする}{多くする}{ぶーさ\tl{˥}\ws{}なすん\tl{˩˥}}
\rentry{おおくなる}{多くなる}{ぶーさ\tl{˥}\ws{}なるん\tl{˩˥}}
\rentry{おおごえ}{大声}{まーぎ\tl{˩˥}くい\tl{˥}}
\rentry{おおざけのみ}{大酒飲み}{さきじょーぐ\tl{˥˩}}
\rentry{おおざっぱに}{大雑把に}{\josho{}ざっ\kako{}と}
\rentry{おおざら}{大皿}{ぶーさら}
\rentry{おおしお}{大潮}{ぶーすー\tl{˥}}
\rentry{おおたにわたり}{オオタニワタリ}{ふつがら\tl{˥}}
\rentry{おおなみ}{大波}{ぶーなん\tl{˥}、むさん\tl{˥˩}}
\rentry{おおばかもの}{大馬鹿者}{まーぷりむん\tl{˥˩}}
\rentry{おおはま}{大浜}{ぼーま\tl{˥}}
\rentry{おおみそかのよる}{大晦日の夜}{とぅしぃぬ\tl{˥}\ws{}ゆー\tl{˩˥}、とぅしぃぬ\tl{˥}\ws{}ゆる\tl{˩˥}}
\rentry{おおむかし}{大昔}{かーま\tl{˥˩}むがし\tl{˩˥}}
\rentry{おおもうけ}{大儲け}{あらもーぎ\tl{˥˩}}
\rentry{おおわらい}{大笑い}{ぷりばーり\tl{˥˩}}
\rentry{おか}{丘}{むり\tl{˥˩}}
\rentry{おかしい}{}{いふなー\tl{˥˩}}
\rentry{おかしい}{可笑しい}{ばーしゃ\tl{˩˥}はん}
\rentry{おかず}{}{かちむん\tl{˥}}
\rentry{おかね}{お金}{じん\tl{˩˥}}
\rentry{おかねいれ}{お金入れ}{じんふくる\tl{˩˥}}
\rentry{おかねをためる}{金を貯める}{じん\tl{˩˥}たみるん\tl{˥˩}}
\rentry{おがみ}{拝み}{うがみ\tl{˥}}
\rentry{おがむ}{拝む}{うがむん\tl{˥}、こーみ\tl{˥}}
\rentry{おがむこと}{拝むこと}{うがみ\tl{˥}}
\rentry{おかよい}{陸酔い}{じふねー\tl{˩˥}}
\rentry{おがわら}{雄瓦}{びーがーら\tl{˥˩}}
\rentry{おき}{沖}{とー\tl{˥˩}}
\rentry{おぎなう}{補う}{たらすん\tl{˥˩}}
\rentry{おきなわ}{沖縄}{うすな\tl{˥}}
\rentry{おきなわほんとう}{沖縄本島}{うすな\tl{˥}}
\rentry{おきゃく}{お客}{しん\tl{˥}}
\rentry{おきゅう}{お灸}{やち\tl{˩˥}}
\rentry{おきる}{起きる}{うぎるん、うぎん\tl{˥}}
\rentry{おく}{奥}{うく\tl{˥˩}、すく\tl{˥˩}}
\rentry{おく}{置く}{うすくん\tl{˩˥}、すくん\tl{˥}}
\rentry{おくのほう}{奥の方}{うく\tl{˥˩}}
\rentry{おくびょうだ}{臆病だ}{いじぃぬ\tl{˥˩}\ws{}ねーぬ\tl{˩˥}、いじぃ\ws{}ねーぬ、うくびょー\tl{˥}さーん}
\rentry{おくらせる}{遅らせる}{うくらすん\tl{˥˩}}
\rentry{おくる}{送る}{うぐるん\tl{˥˩}}
\rentry{おくれる}{遅れる}{うくりん\tl{˥˩}、にふつぁ\tl{˥˩}はん}
\rentry{おけ}{桶}{たんぐ\tl{˥}}
\rentry{おこう}{お香}{こー\tl{˥}}
\rentry{おこげ}{}{なびしき\tl{˩˥}}
\rentry{おこす}{熾す}{〔火を\ntilde{}〕てしきるん\tl{˥˩}}
\rentry{おこす}{起こす}{うがすん\tl{˥}}
\rentry{おこたる}{怠る}{うくたるん\tl{˥˩}}
\rentry{おこりんぼ}{怒りんぼ}{くんぞーむぬ}
\rentry{おこる}{怒る}{うくるん\tl{˥}}
\rentry{おこる}{起こる}{うくるん\tl{˥}}
\rentry{おこること}{怒ること}{くんぞー\tl{˥}}
\rentry{おこわ}{}{かしき\tl{˥}}
\rentry{おさえこむ}{押さえ込む}{うししきん\tl{˥˩}}
\rentry{おさえる}{押さえる}{うすん\tl{˥˩}}
\rentry{おさけ}{お酒}{ぐしん\tl{˩˥}}
\rentry{おさないこ}{幼い子}{いしゃがーたま\tl{˥}}
\rentry{おさまる}{治まる}{うさまるん\tl{˥}}
\rentry{おさめる}{治める}{うさみるん\tl{˥}}
\rentry{おさめる}{納める}{うさみるん\tl{˥}}
\rentry{おじ}{伯父}{ぶざま\tl{˩˥}}
\rentry{おじ}{叔父}{ぶざま\tl{˩˥}}
\rentry{おしあげる}{押し上げる}{うしあんぎるん\tl{˥˩}}
\rentry{おしい}{惜しい}{あがやー、あたら\tl{˥˩}さん}
\rentry{おじいさん}{}{うしとぅぶや、ぶや\tl{˥}}
\rentry{おしいれ}{押入れ}{うちくる\tl{˥˩}}
\rentry{おしえる}{教える}{ならすん\tl{˩˥}}
\rentry{おじおばのまご}{叔父叔母の孫}{\josho{}いちふ\kako{}ぶい}
\rentry{おしかえす}{押し返す}{うしかいすん\tl{˥˩}}
\rentry{おじき}{お辞儀}{こーみ\tl{˥}}
\rentry{おじぎする}{お辞儀する}{こむん\tl{˥}}
\rentry{おじけづく}{怖気づく}{たんきるん\tl{˥}}
\rentry{おしこむ}{押し込む}{うしくむん\tl{˥˩}}
\rentry{おしこめる}{押し込める}{うしくみるん\tl{˥˩}}
\rentry{おしころがす}{押し転がす}{うしくるばすん\tl{˥˩}}
\rentry{おしころばす}{押し転ばす}{うしくるばすん\tl{˥˩}}
\rentry{おしたおす}{押し倒す}{うしとーすん}
\rentry{おしだまる}{押し黙る}{とぅなばるん\tl{˥˩}}
\rentry{おしつける}{押し付ける}{うししきるん\tl{˥˩}}
\rentry{おしつぶす}{押しつぶす}{びらがすん\tl{˥˩}}
\rentry{おしてひらたくする}{押して平たくする}{ぴさんたらすん\tl{˥˩}、びらがすん\tl{˥˩}}
\rentry{おしはかる}{推し量る}{ばがるん\tl{˩˥}}
\rentry{おしもどす}{押し戻す}{うしむどぅすん\tl{˥˩}}
\rentry{おじや}{}{ずーし\tl{˩˥}}
\rentry{おしゃべり}{}{むぬぱなし\tl{˩˥}、ゆんたく\tl{˩˥}}
\rentry{おしゃれだ}{}{はいから\tl{˥}さーん}
\rentry{おしよせる}{押し寄せる}{うしゆしるん}
\rentry{おす}{押す}{うすん\tl{˥˩}}
\rentry{おす}{雄}{びーむぬ\tl{˥˩}}
\rentry{おぜん}{お膳}{じん\tl{˥˩}}
\rentry{おそい}{遅い}{にふつぁ\tl{˥˩}はん}
\rentry{おそなえする}{お供えする}{しきるん\tl{˥˩}}
\rentry{おそろしい}{恐ろしい}{うすまさん、うとぅるさ\tl{˥}はん、ごー\tl{˥}はん、ぬぐりしゃ\tl{˩˥}はん}
\rentry{おだく}{汚濁になる}{やなごるん\tl{˥˩}}
\rentry{おたふくかぜ}{}{とーしんべー\tl{˥}}
\rentry{おたま}{お玉}{すーぺー\tl{˥}}
\rentry{おたまじゃくし}{オタマジャクシ}{たらぐ\tl{˥}}
\rentry{おだやかになる}{穏やかになる}{なだらぐん\tl{˥˩}}
\rentry{おちつきがない}{落ち着かない}{すむあーり}
\rentry{おちつく}{落ち付く}{うちすくん\tl{˥}}
\rentry{おちぶれる}{落ちぶれる}{さぼーりん\tl{˥˩}}
\rentry{おちゃうけ}{お茶うけ}{さうき\tl{˥}}
\rentry{おちる}{落ちる}{うちん\tl{˥}}
\rentry{おっかけていく}{追っかけていく}{ういんぐん\tl{˥˩}}
\rentry{おっかぶせる}{}{すかふちるん\tl{˥}}
\rentry{おっこちる}{落っこちる}{すかうちるん\tl{˥}}
\rentry{おっと}{夫}{ぶとぅ\tl{˥˩}}
\rentry{おっぱらう}{追っ払う}{ういぱろーん\tl{˥˩}}
\rentry{おつゆ}{}{すー\tl{˥}}
\rentry{おてんば}{お転婆}{ぱたらしやみどぅむ、ぴゃんが\tl{˥}}
\rentry{おてんばむすめ}{お転婆娘}{ぴゃんが\tl{˥}}
\rentry{おと}{音}{うとー\tl{˥˩}}
\rentry{おとうと}{弟}{うとぅーとぅ\tl{˥}}
\rentry{おどける}{}{ぴらすん\tl{˥}}
\rentry{おとこ}{男}{びどぅむ\tl{˥˩}}
\rentry{おとこのきょうだい}{男の兄弟}{びぎり\tl{˥˩}}
\rentry{おとこのこのあいしょう}{男の児の愛称}{こー\josho{}に\kako{}ー}
\rentry{おとこのさいこん}{男の再婚}{またそーり\tl{˥˩}}
\rentry{おとさた}{音沙汰}{うとぅさた\tl{˥˩}}
\rentry{おとす}{落とす}{うたすん\tl{˥}}
\rentry{おどす}{脅す}{うどーし}
\rentry{おととい}{一昨日}{ぶとぅち\tl{˩˥}}
\rentry{おととし}{一昨年}{みちなりん}
\rentry{おとなしい}{大人しい}{うとぅなさ\tl{˥}はん}
\rentry{おとめ}{乙女}{みやらび\tl{˩˥}}
\rentry{おどり}{踊り}{ぶどぅり\tl{˥˩}、ぶんどぅり}
\rentry{おとる}{劣る}{うとぅるん\tl{˥}、ぴな\tl{˥}はん}
\rentry{おどる}{踊る}{ぶどぅるん\tl{˥˩}、ぶんどぅるん\tl{˥˩}}
\rentry{おとろえる}{衰える}{うとぅるいるん\tl{˥}、〔植物の勢いが\ntilde{}〕だーりるん\tl{˩˥}}
\rentry{おどろかす}{驚かす}{うどぅがすん\tl{˥}}
\rentry{おどろき}{驚き}{うどぅるぎ\tl{˥}}
\rentry{おどろく}{驚く}{うどぅるぐん\tl{˥}}
\rentry{おなかをこわしやすい}{お腹をこわしやすい}{ばた\tl{˩˥}よー\tl{˩˥}さーん}
\rentry{おなじ}{同じ}{ゆぬ、ゆぬむぬ}
\rentry{おなじこと}{同じこと}{ゆぬむぬ}
\rentry{おなじていど}{同じ程度}{ゆぬしゅく\tl{˩˥}}
\rentry{おなじとし}{同じ年}{ゆぬとぅし\tl{˩˥}}
\rentry{おなじもの}{同じもの}{ゆぬむぬ}
\rentry{おに}{鬼}{うん\tl{˥}}
\rentry{おにおこぜ}{オニオコゼ}{いしゃば\tl{˥˩}}
\rentry{おの}{斧}{ぶぬ\tl{˩˥}}
\rentry{おば}{伯母}{ぶば\tl{˥}}
\rentry{おば}{叔母}{ぶば\tl{˥}}
\rentry{おばあさん}{}{ぱー\tl{˥}}
\rentry{おばさんたち}{おばさん達}{ぱんだー\tl{˥}}
\rentry{おはつ}{お初}{ぱち\tl{˥˩}}
\rentry{おび}{帯}{すくび\tl{˥˩}}
\rentry{おぼえ}{覚え}{うぶい\tl{˥}}
\rentry{おぼえる}{覚える}{うぶいるん\tl{˥}、うぶいん}
\rentry{おぼれる}{溺れる}{うぶりるん\tl{˥˩}}
\rentry{おぼん}{お盆}{そーりん\tl{˥}}
\rentry{おぼんのさいしゅうび}{お盆の最終日}{うぐるぴん\tl{˥˩}}
\rentry{おまえ}{お前}{だー\tl{˥˩}}
\rentry{おまえは}{お前は}{だんざー\tl{˥}}
\rentry{おまえら}{お前ら}{だいま\tl{˥˩}}
\rentry{おみき}{お神酒}{ぐしん\tl{˩˥}}
\rentry{おみやげ}{お土産}{すぅとぅ\tl{˥}}
\rentry{おむかえする}{お迎えする}{しけーすん\tl{˥˩}}
\rentry{おもい}{重い}{んさ\tl{˩˥}はん}
\rentry{おもいあたる}{思い当たる}{うむいあたるん}
\rentry{おもいきる}{思い切る}{うむいきすん\tl{˥}、むいきし\tl{˥}}
\rentry{おもいこがれる}{思い焦がれる}{うむいくがりん\tl{˥}}
\rentry{おもいこむ}{思い込む}{うむいくむん\tl{˥}、むいくみ\tl{˥}}
\rentry{おもいだす}{思い出す}{うむいんだすん\tl{˥}、むんじる\tl{˩˥}}
\rentry{おもいたつ}{思い立つ}{うむいたつん\tl{˥}}
\rentry{おもいつく}{思い付く}{うむい\tl{˥}すくん\tl{˥˩}}
\rentry{おもいつめる}{思い詰める}{うむいつみるん\tl{˥}、うむいつみん\tl{˥}}
\rentry{おもいなおす}{思い直す}{うむい\tl{˥}のーすん\tl{˩˥}}
\rentry{おもいのこす}{思い残す}{うむい\tl{˥}ぬぐすん\tl{˩˥}}
\rentry{おもいやり}{思いやり}{なさき\tl{˩˥}}
\rentry{おもう}{思う}{かんげーるん\tl{˥}、むん\tl{˥}}
\rentry{おもしろい}{面白い}{むっさ\tl{˥}はん}
\rentry{おもたい}{重たい}{んさ\tl{˩˥}はん}
\rentry{おもちゃ}{玩具}{だーびむぬ\tl{˩˥}}
\rentry{おもとだけ}{於茂登岳}{むとぅだぎ\tl{˥˩}}
\rentry{おもや}{母屋}{ぶー\tl{˥}ひー\tl{˩˥}}
\rentry{おや}{}{あい}
\rentry{おや}{親}{うや\tl{˥}}
\rentry{おやかた}{親方}{うやかた\tl{˥}}
\rentry{おやこ}{親子}{\josho{}うや\kako{}ふぁ}
\rentry{おやすみになる}{お休みになる}{ゆごー\tl{˩˥}\ws{}すん\tl{˥˩}}
\rentry{おやゆび}{親指}{ぶーび\tl{˥}}
\rentry{およぐ}{泳ぐ}{うい\tl{˥}}
\rentry{およぶ}{及ぶ}{ういぶん\tl{˥˩}、うゆぶん\tl{˥˩}}
\rentry{おりいど}{降り井戸}{うりげー\tl{˥}}
\rentry{おりこうさん}{お利口さん}{ぼーりぼーり}
\rentry{おりたたむ}{折り畳む}{たたむん}
\rentry{おりもの}{織物}{ぬぬ\tl{˥˩}}
\rentry{おる}{折る}{ぶるん\tl{˩˥}}
\rentry{おる}{織る}{うるん\tl{˥˩}}
\rentry{おれ}{俺}{ばぬ\tl{˩˥}}
\rentry{おれやすい}{折れやすい}{さぱ\tl{˥˩}はん}
\rentry{おれる}{折れる}{ぶりるん\tl{˩˥}}
\rentry{おろす}{下ろす}{うらすん\tl{˥}}
\rentry{おろす}{降ろす}{うらすん\tl{˥}}
\rentry{おわせる}{負わせる}{〔全責任を\ntilde{}〕ぱかすん\tl{˥˩}}
\rentry{おわり}{終わり}{おび\tl{˥}}
\rentry{おわる}{終わる}{うわるん\tl{˥˩}、おわるん\tl{˥˩}、しまいるん\tl{˥˩}、すむん\tl{˥}}
\rentry{おわれる}{追われる}{わっきらいるん\tl{˥˩}}
\rentry{おん}{恩}{おん\tl{˥}、ぶん\tl{˩˥}、ぶんぎ\tl{˩˥}}
\rentry{おんぎ}{恩義}{おん\tl{˥}、ぶん\tl{˩˥}、ぶんぎ\tl{˩˥}}
\rentry{おんぎょく}{音曲}{うたさんしん\tl{˥˩}}
\rentry{おんしん}{音信}{うとぅさた\tl{˥˩}}
\rentry{おんな}{女}{みどぅむ\tl{˩˥}}
\rentry{おんなのこ}{女の子}{みどぅんたま\tl{˩˥}}
\rentry{おんなのさいこん}{女の再婚}{またむち\tl{˥˩}}
\rentry{おんなのふんどし}{女のふんどし}{めっつぁ\tl{˩˥}}
\rentry{おんなわらべ}{女わらべ}{みやらび\tl{˩˥}}
\rentry{おんぶする}{}{かにるん\tl{˥}}
\rentry{おんわだ}{穏和だ}{うとぅなさ\tl{˥}はん}
\rentry{おんをかえす}{恩を返す}{ぶん\tl{˩˥}かいすん\tl{˥}}
\rentry{か}{蚊}{がんざん\tl{˩˥}}
\rentry{が}{我}{がー\tl{˩˥}}
\rentry{が}{蛾}{ぱぴる\tl{˥}}
\rentry{かーかー}{}{がーがー}
\rentry{かい}{櫂}{いやぐ\tl{˥}}
\rentry{かい}{貝}{みなん\tl{˩˥}}
\rentry{がい}{害}{がい\tl{˥˩}}
\rentry{～かい}{～回}{むし}
\rentry{かいが}{絵画}{えー\tl{˥˩}}
\rentry{かいがら}{貝殻}{みなんぬ\tl{˩˥}\ws{}くー\tl{˥}}
\rentry{かいがんふきん}{海岸付近}{\josho{}いなん\kako{}ぱた}
\rentry{かいこんする}{開墾する}{ありしあぎるん\tl{˥˩}}
\rentry{かいこんはたち}{開墾畑地}{ありしぴてー\tl{˥˩}}
\rentry{がいしゅつ}{外出}{ふかまーり\tl{˥}}
\rentry{かいすい}{海水}{すー\tl{˥}}
\rentry{かいすいよく}{海水浴}{おんだ\tl{˥}}
\rentry{がいする}{害する}{がい\tl{˥˩}\ws{}すん\tl{˥˩}}
\rentry{かいぜんする}{改善する}{あらたみるん\tl{˥}}
\rentry{かいたいする}{解体する}{ばずらすん\tl{˥˩}}
\rentry{かいだん}{階段}{\josho{}だん\kako{}だん}
\rentry{かいちゅう}{回虫}{ばたぬ\tl{˩˥}\ws{}むしぃ\tl{˥˩}}
\rentry{がいちゅう}{害虫}{やなむしぃ\tl{˥˩}}
\rentry{かいてんさせる}{回転させる}{まーすん\tl{˥˩}}
\rentry{かいてんする}{回転する}{まーるん\tl{˥˩}}
\rentry{かいふくする}{回復する}{むちのーるん\tl{˩˥}}
\rentry{かう}{買う}{こーん\tl{˥˩}}
\rentry{かう}{飼う}{すかのーすん\tl{˥}}
\rentry{かえりみち}{帰り道}{むどぅりみち\tl{˩˥}}
\rentry{かえる}{代える}{かれーるん\tl{˥˩}}
\rentry{かえる}{変える}{かれーるん\tl{˥˩}}
\rentry{かえる}{帰る}{けーるん\tl{˥}}
\rentry{かえる}{替える}{かれーるん\tl{˥˩}}
\rentry{かえる}{蛙}{おった\tl{˥}}
\rentry{かお}{顔}{むち\tl{˥}}
\rentry{かおく}{家屋}{ひー\tl{˩˥}}
\rentry{かおり}{香り}{かー\tl{˥˩}}
\rentry{かおる}{薫る}{かぱ\tl{˥˩}はん}
\rentry{かがせる}{嗅がせる}{かぱすん\tl{˥˩}}
\rentry{かかと}{踵}{あどぅ\tl{˥}}
\rentry{かがみ}{鏡}{かんがん\tl{˥}}
\rentry{かがむ}{屈む}{すくまるん\tl{˥}}
\rentry{かかる}{}{〔禁忌に\ntilde{}〕はかるん\tl{˥}}
\rentry{かかる}{掛かる}{はかるん\tl{˥}}
\rentry{かかる}{罹る}{〔病気に\ntilde{}〕はかるん\tl{˥}}
\rentry{かき}{夏季}{なちぃ\tl{˥˩}}
\rentry{かき}{書く}{はくん\tl{˥}}
\rentry{かきあつめてとる}{掻き集めて取る}{はかじ\tl{˥}\ws{}とぅるん\tl{˥}}
\rentry{かきあつめる}{掻き集める}{はきあしみるん\tl{˥}}
\rentry{かきいれる}{書き入れる}{はき\tl{˥}いりるん\tl{˥˩}}
\rentry{かきおとす}{書き落とす}{はきもらすん\tl{˥}}
\rentry{かきそえる}{書き添える}{はきたすん\tl{˥}}
\rentry{かきたす}{書き足す}{はきたすん\tl{˥}}
\rentry{かきだす}{掻き出す}{はきんだすん\tl{˥}}
\rentry{かきちらかす}{掻き散らかす}{はきぽーるん\tl{˥}}
\rentry{かきとどめる}{書き留める}{はきとぅみるん\tl{˥}}
\rentry{かぎのて}{鉤の手}{かーぎ\tl{˥}}
\rentry{かきまぜる}{掻き混ぜる}{きんつぁーらすん\tl{˥}、はき\tl{˥}まんじるん\tl{˩˥}}
\rentry{かきまわす}{掻き回す}{はきまーすん\tl{˥}}
\rentry{かきみだす}{掻き乱す}{きんだるん\tl{˥˩}、はき\tl{˥}まんじるん\tl{˩˥}}
\rentry{かきもらす}{書き漏らす}{はきもらすん\tl{˥}}
\rentry{かく}{描く}{はくん\tl{˥}}
\rentry{かく}{搔く}{はくん\tl{˥}}
\rentry{かぐ}{嗅ぐ}{〔匂いを\ntilde{}〕かぷん\tl{˥˩}}
\rentry{かくげん}{格言}{むがし\tl{˩˥}くとぅば\tl{˥}}
\rentry{かくじ}{各自}{なーどぅーどぅ}
\rentry{かくしご}{隠し子}{ぐんぼー\tl{˩˥}}
\rentry{かくじで}{各自で}{どぅしけな}
\rentry{がくしゅう}{学習}{むぬなれー\tl{˩˥}}
\rentry{かくす}{隠す}{かつぁみるん\tl{˥}、はこすん\tl{˥}}
\rentry{かくだいする}{拡大する}{ぴそーぎるん}
\rentry{がくどう}{学童}{がくたま}
\rentry{かくにんする}{確認する}{たしかみるん\tl{˥}}
\rentry{がくもん}{学問}{がくむん\tl{˥}}
\rentry{かくれる}{隠れる}{くまるん\tl{˥}、はこりん\tl{˥}}
\rentry{かけ}{賭け}{かーき\tl{˥}}
\rentry{かげ}{影}{かぎ\tl{˥}}
\rentry{かげ}{陰}{かぎ\tl{˥}、けー}
\rentry{かけがいする}{掛け買いする}{さがらすん\tl{˥}}
\rentry{かけごえ}{掛声}{やぐい\tl{˩˥}}
\rentry{かげぜん}{陰膳}{かぎじん\tl{˥}}
\rentry{かけまわる}{駆け回る}{はきまーるん\tl{˥}}
\rentry{かけら}{欠片}{はき\tl{˥˩}}
\rentry{かける}{掛ける}{はきるん\tl{˥}}
\rentry{かける}{欠ける}{はきるん\tl{˥˩}、はきん\tl{˥˩}}
\rentry{かける}{駆ける}{かしきるん\tl{˥}、はきしぃきるん\tl{˥}}
\rentry{かご}{籠}{かぐ\tl{˥}}
\rentry{かこつける}{託ける}{なしきるん\tl{˥}}
\rentry{かこませる}{囲ませる}{はこますん\tl{˥˩}}
\rentry{かこむ}{囲む}{はこますん\tl{˥˩}、はこむん\tl{˥˩}}
\rentry{かさ}{傘}{さな\tl{˥}}
\rentry{かさ}{暈}{ぴながん\tl{˥}}
\rentry{かさ}{笠}{かつぁ\tl{˥}}
\rentry{かざかみへ}{風上へ}{おーらが\tl{˥˩}}
\rentry{かざぐるま}{風車}{かざまやー\tl{˥}}
\rentry{かざしも}{風下}{すとーま\tl{˥˩}}
\rentry{かさなる}{重なる}{かさばるん\tl{˥˩}、かつぁなるん\tl{˥˩}}
\rentry{かさねる}{重ねる}{かさにるん\tl{˥˩}、かさびるん\tl{˥˩}、かつぁびるん\tl{˥˩}}
\rentry{かさむ}{嵩む}{かつぁむん\tl{˥˩}}
\rentry{かざむき}{風向}{たむき\tl{˥}}
\rentry{かざり}{飾り}{かつぁり\tl{˥˩}}
\rentry{かし}{菓子}{こーし\tl{˥}}
\rentry{かじ}{火事}{ぴー\tl{˥}}
\rentry{かじ}{舵}{かち\tl{˥˩}}
\rentry{かじ}{鍛冶}{かつぇー\tl{˥˩}}
\rentry{がし}{餓死}{がーし\tl{˥˩}}
\rentry{かしこい}{賢い}{そーいり\tl{˥˩}}
\rentry{がしのとし}{餓死の年}{がしぃどぅしぃ}
\rentry{かじぼう}{舵棒}{てーに\tl{˥}}
\rentry{かじや}{鍛冶屋}{かつぇー\tl{˥˩}}
\rentry{かじやまつり}{鍛冶屋祭}{かつぇーぶな\tl{˥˩}}
\rentry{かじゅ}{果樹}{なーりむぬ\tl{˩˥}}
\rentry{がじゅまる}{ガジュマル}{がざまに\tl{˩˥}}
\rentry{かす}{貸す}{からすん\tl{˥˩}}
\rentry{かず}{数}{かじ\tl{˥}}
\rentry{かずら}{蔓}{かちぃら\tl{˥˩}}
\rentry{かぜ}{風}{かち\tl{˥˩}}
\rentry{かせい}{加勢}{かしー\tl{˥}}
\rentry{かぜがつよい}{風が強い}{かちょ\tl{˥˩}はん}
\rentry{かぜがはやる}{風邪が流行る}{ぱなしきぬ\ws{}はやるん}
\rentry{かぜまわり}{風廻り}{かちまーり\tl{˥˩}}
\rentry{かぜよけ}{風除け}{かたが\tl{˥}}
\rentry{かせん}{河川}{かーら\tl{˥˩}}
\rentry{かそうぎょうれつ}{仮装行列}{〔旧盆の\ntilde{}〕みちすねー\tl{˥˩}}
\rentry{かぞえる}{数える}{かずいるん\tl{˥}、ゆむん\tl{˩˥}}
\rentry{かぞく}{家族}{やーにんじゅ\tl{˩˥}}
\rentry{かた}{型}{かた\tl{˥˩}}
\rentry{かた}{肩}{かた\tl{˥}}
\rentry{かたあし}{片足}{かたぱん\tl{˥}}
\rentry{かたい}{固い}{こー\tl{˥}はん}
\rentry{かたい}{硬い}{こー\tl{˥}はん}
\rentry{かたいっぽう}{片一方}{かたぐ\tl{˥}}
\rentry{かたおや}{片親}{かたうや\tl{˥}}
\rentry{がたがた}{}{がたがた}
\rentry{かたがわ}{片側}{ぴぃとぅかた\tl{˥}}
\rentry{かたく}{固く}{こーし}
\rentry{かたくなる}{固くなる}{こっぱるん\tl{˥}}
\rentry{かたくなる}{硬くなる}{かたまるん\tl{˥˩}、こーるん\tl{˥}、こすぱるん\tl{˥}}
\rentry{かたぐるま}{肩車}{かんご\tl{˥}}
\rentry{かたすみ}{片隅}{ゆごー\tl{˥˩}}
\rentry{かたち}{形}{かたち\tl{˥˩}}
\rentry{かたづけ}{片付け}{かたすか\tl{˥}}
\rentry{かたづけてきれいにする}{片付けてきれいにする}{あざーぎるん\tl{˥}}
\rentry{かたづける}{片付ける}{まだぎるん\tl{˩˥}、まどぅぎるん\tl{˩˥}}
\rentry{かたつむり}{カタツムリ}{すたみん\tl{˥}}
\rentry{かたな}{刀}{かたな\tl{˥}}
\rentry{かたぶり}{片降り}{かたぶり\tl{˥}}
\rentry{かたほう}{片方}{かたぐ\tl{˥}}
\rentry{かたまる}{固まる}{かたまるん\tl{˥˩}、こーるん\tl{˥}、こすぱるん\tl{˥}}
\rentry{かたむく}{傾く}{かたふくん\tl{˥}}
\rentry{かたむける}{傾ける}{かたふかすん、かたふきるん\tl{˥}}
\rentry{かためる}{固める}{かたみるん\tl{˥˩}}
\rentry{かたらう}{語らう}{かたるん\tl{˥˩}}
\rentry{かたる}{語る}{かたるん\tl{˥˩}、ぱなすん\tl{˥}}
\rentry{かたわら}{傍ら}{ぱた\tl{˥˩}}
\rentry{かち}{勝ち}{かち\tl{˥}}
\rentry{かつ}{勝つ}{かつん\tl{˥}}
\rentry{かつえる}{餓える}{かちりるん\tl{˥}、かちりん\tl{˥}}
\rentry{かつお}{鰹}{かちゅ\tl{˥˩}}
\rentry{かつおぎょせん}{鰹漁船}{かちゅぶに\tl{˥˩}}
\rentry{かつおこうじょう}{鰹工場}{なや\tl{˥˩}}
\rentry{かつおのえさのこざかな}{鰹の餌の小魚}{じゃく\tl{˩˥}}
\rentry{かつおのけずりかす}{鰹の削りかす}{ぴじがら\tl{˥}}
\rentry{かつおぶし}{鰹節}{かちゅぶし\tl{˥˩}}
\rentry{かつぐ}{担ぐ}{〔肩に\ntilde{}〕かたみるん\tl{˥}、かにるん\tl{˥}}
\rentry{かっこう}{恰好}{かたち\tl{˥˩}}
\rentry{がっこう}{学校}{がく}
\rentry{かっこうわるい}{恰好悪い}{うきしゃ\tl{˥}はん}
\rentry{かつこと}{勝つこと}{かち\tl{˥}}
\rentry{がっちり}{}{だ\josho{}ん\kako{}た}
\rentry{かってほうだいだ}{勝手放題だ}{ふんだやー}
\rentry{かっぱらう}{}{ばぎとぅるん\tl{˥}}
\rentry{かてい}{家庭}{きに\tl{˥˩}}
\rentry{かていそうぎ}{家庭内争議}{やーむんどぅ}
\rentry{かど}{角}{かどぅ\tl{˥}}
\rentry{かなう}{叶う}{ういぶん\tl{˥˩}、かのーん\tl{˥}}
\rentry{かなえる}{叶える}{かねーるん\tl{˥}}
\rentry{かなしい}{悲しい}{しから\tl{˥}はん、なぐりしゃ\tl{˩˥}はん}
\rentry{かなしがる}{悲しがる}{しからさ\tl{˥}\ws{}すん\tl{˥˩}}
\rentry{かなしみふさぐ}{悲しみふさぐ}{しから\tl{˥}はーん}
\rentry{かなづち}{金槌}{かにち\tl{˥˩}}
\rentry{かならず}{必ず}{やっ\josho{}ち\kako{}ん}
\rentry{かに}{蟹}{かん\tl{˥˩}}
\rentry{かね}{鉦}{かに\tl{˥˩}}
\rentry{かね}{鐘}{かに\tl{˥˩}}
\rentry{かねじゃく}{曲尺}{ばんぞーがに\tl{˩˥}}
\rentry{かねもうけ}{金儲け}{じんもーき\tl{˩˥}}
\rentry{かねもち}{金持ち}{うやぎ\tl{˥}、じんむちゃー\tl{˩˥}}
\rentry{かねもちのいえ}{金持ちの家}{\josho{}うやぎ\kako{}ひー}
\rentry{がばと}{}{やいっと}
\rentry{かび}{黴}{かぺー\tl{˥}}
\rentry{かびん}{花瓶}{ぱないぎ\tl{˥}}
\rentry{かふう}{家風}{やーなれー}
\rentry{かふう}{火風}{ぴーかち\tl{˥}}
\rentry{かぶせる}{被せる}{かふちるん\tl{˥}}
\rentry{かぶる}{被る}{かぷん\tl{˥}}
\rentry{かぶれ}{}{まぎ\tl{˥˩}}
\rentry{かぶわけする}{株分けする}{にんばぎるん\tl{˩˥}}
\rentry{かほうへ}{下方へ}{したいが\tl{˥˩}}
\rentry{かま}{釜}{かま\tl{˥˩}}
\rentry{かま}{鎌}{がっきゃ\tl{˩˥}}
\rentry{かまえる}{構える}{かまいるん\tl{˥}}
\rentry{かまぼこ}{}{かまぐ\tl{˥˩}}
\rentry{かまわない}{構わない}{だすたすくん\tl{˩˥}}
\rentry{がまんつよい}{我慢強い}{がーずさはん}
\rentry{がまんできない}{我慢できない}{ふしがらぬ}
\rentry{かみ}{神}{かん\tl{˥}}
\rentry{かみ}{髪}{あまじ\tl{˥}}
\rentry{かみうらない}{神占い}{うがじぃ\tl{˥}}
\rentry{かみぎょうじ}{神行事}{かぬめー\tl{˥}、きさり\tl{˥˩}}
\rentry{かみぎょうじのぎょろうがかり}{神行事の漁労係}{いしょーぶさ\tl{˥˩}}
\rentry{かみきる}{噛み切る}{じんぱるん\tl{˥˩}}
\rentry{かみくだく}{噛み砕く}{かんぱるん\tl{˥}}
\rentry{かみさま}{神様}{かん\tl{˥}}
\rentry{かみそり}{剃刀}{かんすり\tl{˥}}
\rentry{かみなり}{雷}{かんなり\tl{˥}}
\rentry{かみのみち}{神の道}{かんみちぃ\tl{˥}}
\rentry{かみのよ}{神の世}{かんぬ\tl{˥}\ws{}ゆー\tl{˥˩}、むがしぃ\tl{˩˥}ゆー\tl{˥˩}}
\rentry{かみのれいげんがたかい}{神の霊験が高い}{かんだが\tl{˥}さーん}
\rentry{かむ}{咬む}{かむん\tl{˥}}
\rentry{かむ}{噛む}{かむん\tl{˥}、かんだるん\tl{˥}、〔歯で\ntilde{}〕ふしきるん\tl{˥}}
\rentry{かめ}{亀}{かみ\tl{˥}}
\rentry{かめ}{甕}{かみ\tl{˥}}
\rentry{かもされる}{醸される}{〔酒が\ntilde{}〕ばぐん\tl{˥˩}}
\rentry{かもん}{家紋}{やーばん\tl{˩˥}}
\rentry{かや}{茅}{がや\tl{˩˥}}
\rentry{かや}{蚊帳}{かつぁ\tl{˥˩}}
\rentry{がやがや}{}{がーがー}
\rentry{かやぶきのかおく}{萱ぶきの家屋}{がやふきひー\tl{˩˥}}
\rentry{かゆ}{粥}{ゆー\tl{˩˥}}
\rentry{かゆい}{痒い}{びょー\tl{˩˥}はん}
\rentry{かよう}{通う}{かよーん\tl{˥˩}}
\rentry{からい}{辛い}{から\tl{˥}はん}
\rentry{からかう}{}{ぱくるん\tl{˥}、ばっくるん\tl{˩˥}}
\rentry{からから}{}{がらがら}
\rentry{からげる}{絡げる}{からぎるん\tl{˥}}
\rentry{からす}{カラス}{がらし\tl{˩˥}}
\rentry{からだ}{体}{ぐてー\tl{˩˥}、どぅー\tl{˩˥}}
\rentry{からだがけんこうだ}{体が健康だ}{どぅー\tl{˩˥}がんじゅー\tl{˩˥}さーん}
\rentry{からだがじょうぶだ}{体が丈夫だ}{どぅー\tl{˩˥}ずさ\tl{˩˥}はん}
\rentry{からだがつよくなる}{体が強くなる}{がんずさ\tl{˩˥}\ws{}なるん\tl{˩˥}}
\rentry{からっぽ}{空っぽ}{んな\tl{˥˩}}
\rentry{からにする}{空にする}{いたちるん}
\rentry{からになる}{空になる}{から\tl{˥˩}\ws{}なるん\tl{˩˥}}
\rentry{からまる}{絡まる}{あんざるん\tl{˥}、からまぐん\tl{˥}}
\rentry{からみつく}{絡み付く}{からまぐん\tl{˥}、からむん\tl{˥}}
\rentry{からめる}{絡める}{からまぐん\tl{˥}}
\rentry{かりる}{借りる}{かりるん\tl{˥˩}、かるん\tl{˥˩}}
\rentry{かる}{刈る}{かるん\tl{˥˩}}
\rentry{かるい}{軽い}{かろ\tl{˥˩}はん}
\rentry{かるくおもう}{軽く思う}{かるんじるん\tl{˥˩}}
\rentry{かれい}{嘉例}{かりー\tl{˥}}
\rentry{かれら}{彼ら}{うしたま\tl{˥}}
\rentry{かれる}{嗄れる}{〔声が\ntilde{}〕かりるん\tl{˥˩}}
\rentry{かれる}{枯れる}{かりるん\tl{˥˩}}
\rentry{かろんじる}{軽んじる}{かるんじるん\tl{˥˩}}
\rentry{かわ}{川}{かーら\tl{˥˩}}
\rentry{かわいい}{可愛い}{かな\tl{˥}さーん、しんだら\tl{˥}さーん}
\rentry{かわいそうだ}{可哀そうだ}{すむ\tl{˥}いた\tl{˥}はん、すむ\tl{˥}ぐり\tl{˥}さん、どぅぐりしゃ\tl{˩˥}はーん}
\rentry{かわいそうなおもいがする}{可哀そうそうな思いがする}{すむ\tl{˥}ぐりしゃ\tl{˥}はーん}
\rentry{かわいらしい}{可愛らしい}{あたら\tl{˥˩}さん、しんだら\tl{˥}さーん}
\rentry{かわかす}{乾かす}{かりがすん\tl{˥˩}、ぷつん\tl{˥}}
\rentry{かわく}{乾く}{かりぐん\tl{˥˩}}
\rentry{かわっている}{変わっている}{いふなー\ws{}やっさー\tl{˥˩}}
\rentry{かわはぎ}{カワハギ}{ふくらび\tl{˥}}
\rentry{かわや}{厠}{ふーる\tl{˥}、ふるや\tl{˥}}
\rentry{かわら}{瓦}{かーら\tl{˥}}
\rentry{かわらぶきのいえ}{瓦葺きの家}{かーら\tl{˥}ひー\tl{˩˥}}
\rentry{かわり}{代わり}{かーり\tl{˥˩}}
\rentry{かわり}{変わり}{かわり}
\rentry{かわる}{代わる}{かわるん\tl{˥˩}}
\rentry{かわる}{変わる}{かわるん\tl{˥˩}、ちごーん\tl{˥˩}}
\rentry{がん}{龕}{がんだるごー\tl{˥˩}}
\rentry{かんおけ}{棺桶}{かん\tl{˥˩}}
\rentry{かんがい}{旱害}{ぴぱり\tl{˥˩}、ぺーりまぎ\tl{˥˩}}
\rentry{かんがえ}{考え}{かんげー\tl{˥}}
\rentry{かんがえる}{考える}{かんげーるん\tl{˥}、むん\tl{˥}}
\rentry{かんき}{寒気}{かん\tl{˥}、ぴらぐ\tl{˥}}
\rentry{がんきゅう}{眼球}{みっちん\tl{˩˥}、みんたま\tl{˥}}
\rentry{がんこもの}{頑固者}{がんくむぬ}
\rentry{かんしょ}{甘蔗}{あますな\tl{˥}}
\rentry{かんしょ}{甘藷}{あがん\tl{˩˥}}
\rentry{がんしょう}{岩礒}{すぱに\tl{˥˩}}
\rentry{がんじょうだ}{頑丈だ}{がんずさ\tl{˩˥}はん}
\rentry{かんしんがある}{関心がある}{きぬ\tl{˥˩}\ws{}あん\tl{˥}}
\rentry{かんせんさせる}{感染させる}{うつすん\tl{˥}}
\rentry{かんせんする}{感染する}{うつるん\tl{˥}}
\rentry{がんそ}{元祖}{がんす\tl{˩˥}}
\rentry{かんぞう}{肝臓}{きむ\tl{˥}}
\rentry{かんたんに}{簡単に}{\josho{}ざっ\kako{}と}
\rentry{かんちょうじ}{干潮時}{すっち\tl{˥}}
\rentry{かんちょうになる}{干潮になる}{すーぴすん\tl{˥}}
\rentry{かんづかさ}{神司}{すかさ\tl{˥}}
\rentry{かんな}{鉋}{かな\tl{˥}}
\rentry{かんぱ}{寒波}{ぴらぐ\tl{˥}}
\rentry{かんばしい}{}{ぱからっ\tl{˥}さ}
\rentry{かんばしい}{芳しい}{かぱ\tl{˥˩}はん}
\rentry{かんばしくない}{芳しくない}{ぱからっさー\tl{˥}\ws{}ねーぬ\tl{˩˥}}
\rentry{かんばつ}{旱魃}{ぺーり\tl{˥˩}}
\rentry{かんばつだ}{旱魃だ}{ぺーり\ws{}しゃーん}
\rentry{がんばる}{頑張る}{ぎばるん\tl{˥˩}}
\rentry{かんまんである}{緩慢である}{〔動作が\ntilde{}〕どぅなー\tl{˩˥}さーん}
\rentry{がんめん}{顔面}{むち\tl{˥}}
\rentry{がんらい}{元来}{むとぅがら\tl{˩˥}}
\rentry{き}{木}{きー\tl{˥}}
\rentry{き}{気}{きー\tl{˥˩}}
\rentry{きーる}{キール}{かーら\tl{˥}、まつら\tl{˥}}
\rentry{きいろ}{黄色}{きーる\tl{˥˩}}
\rentry{きいろになる}{黄色になる}{きんき\ws{}すん\tl{˥˩}}
\rentry{きおく}{記憶}{うぶい\tl{˥}}
\rentry{きおくりょく}{記憶力}{むぬうぶい\tl{˩˥}}
\rentry{きおけ}{木桶}{うき\tl{˥}}
\rentry{きおちする}{気落ちする}{すむだーりるん\tl{˥˩}}
\rentry{きがある}{気がある}{きぬ\tl{˥˩}\ws{}あん\tl{˥}}
\rentry{きがえる}{着替える}{すぅぬ\tl{˥}\ws{}かれーるん\tl{˥˩}}
\rentry{きがきかなくなる}{気が利かなくなる}{じんぶん\tl{˩˥}\ws{}とーらん\tl{˥˩}}
\rentry{きがくるう}{気が狂う}{ぷりるん\tl{˥˩}}
\rentry{きかざる}{着飾る}{しだすん\tl{˥˩}}
\rentry{きかせる}{聞かせる}{すかすん\tl{˥˩}}
\rentry{きがたかまる}{気が高まる}{ばざるん\tl{˥˩}}
\rentry{きがちいさい}{気が小さい}{すむ\tl{˥}いしゃが\tl{˥}はーん}
\rentry{きかれる}{聞かれる}{すかいるん\tl{˥˩}}
\rentry{きがん}{祈願}{うがみ\tl{˥}、うがん\tl{˥}、うやっすり\tl{˥}、にげー\tl{˩˥}}
\rentry{きがんする}{祈願する}{にげるん\tl{˩˥}}
\rentry{きがんのことば}{祈願の言葉}{にげーふちぃ\tl{˩˥}}
\rentry{ききぐるしい}{聞き苦しい}{しきぐり\tl{˥˩}さーん}
\rentry{ききん}{飢饉}{がーし\tl{˥˩}}
\rentry{きく}{聞く}{すくん\tl{˥˩}}
\rentry{きぐらいがたかい}{気位が高い}{すむ\tl{˥}だが\tl{˥}はん}
\rentry{きくらげ}{キクラゲ}{みんぐる\tl{˩˥}、みんぐるみん\tl{˩˥}}
\rentry{きける}{聞ける}{すかいるん\tl{˥˩}}
\rentry{きけんだ}{危険だ}{ぬぐりしゃ\tl{˩˥}はん、ぴこ\tl{˥˩}はん}
\rentry{きこえない}{聞こえない}{すかーるぬ、みんとーら\tl{˩˥}}
\rentry{きこえる}{聞こえる}{すかいるん\tl{˥˩}}
\rentry{きこつ}{気骨}{きがい\tl{˥˩}}
\rentry{きさまは}{}{だんざー\tl{˥}}
\rentry{きざむ}{刻む}{うらすん\tl{˥}、きつぁむん\tl{˥˩}}
\rentry{ぎじゅつ}{技術}{わざ\tl{˩˥}}
\rentry{きずつける}{傷つける}{あやみるん\tl{˥}、いたますん\tl{˥}、いたみるん\tl{˥}、やますん\tl{˩˥}}
\rentry{きせいちゅう}{寄生虫}{ばたぬ\tl{˩˥}\ws{}むしぃ\tl{˥˩}}
\rentry{きせつ}{季節}{しち\tl{˥}}
\rentry{きぜつする}{気絶する}{ふきーるん\tl{˥˩}}
\rentry{きせる}{煙管}{きしり\tl{˥}}
\rentry{きそく}{規則}{さだみ\tl{˥}}
\rentry{きた}{北}{にし\tl{˥˩}}
\rentry{きたいする}{期待する}{あちすん\tl{˥˩}}
\rentry{きたかぜ}{北風}{にしかち\tl{˥˩}}
\rentry{きたがわ}{北側}{にした\tl{˥˩}}
\rentry{きだて}{気立て}{きがい\tl{˥˩}}
\rentry{きたない}{汚い}{あだら\tl{˥˩}、すぷた\tl{˥}はん}
\rentry{きたない}{汚ない}{やにっしゃ\tl{˩˥}はん}
\rentry{きたむら}{北村}{にしむら\tl{˥˩}}
\rentry{きちがい}{気違い}{しんけー\tl{˥}、ぷりむぬ\tl{˥˩}}
\rentry{きちじつ}{吉日}{ぴーる}
\rentry{きちゅう}{忌中}{いみうち}
\rentry{きちゅう}{忌中だ}{いみはかるん\tl{˥}}
\rentry{きちれい}{吉例}{かりー\tl{˥}}
\rentry{きちんとせいりされる}{きちんと整理される}{かたずぐん\tl{˥}}
\rentry{きちんとなる}{}{とぅとぅのーん\tl{˥}}
\rentry{きっかかる}{引っかかる}{ぴきはかるん}
\rentry{きづく}{気付く}{さとぅるん\tl{˥}}
\rentry{きづち}{木槌}{あいち\tl{˥}}
\rentry{きっと}{}{よー\tl{˥˩}}
\rentry{きぬ}{絹}{いちゅ\tl{˥}}
\rentry{きぬのきもの}{絹の着物}{いちゅすぬ}
\rentry{きね}{杵}{いなしき\tl{˥}}
\rentry{きねんさい}{祈年祭}{あみじわー\tl{˥}}
\rentry{きのう}{昨日}{しぬ\tl{˥˩}、すぅぬ}
\rentry{きのこ}{茸}{なば\tl{˩˥}}
\rentry{きのどくだ}{気の毒だ}{すむ\tl{˥}いた\tl{˥}はん、すむ\tl{˥}ぐり\tl{˥}さん、どぅぐりしゃ\tl{˩˥}はーん}
\rentry{きのみき}{木の幹}{きーぬ\ws{}むとぅ}
\rentry{きばむ}{黄ばむ}{きんき\ws{}すん\tl{˥˩}}
\rentry{きぶんがわるい}{気分が悪い}{おま\tl{˥}はん}
\rentry{きぼう}{希望}{ぬずみ\tl{˥˩}}
\rentry{きまずい}{気まずい}{どぅぐりしゃ\tl{˩˥}はーん}
\rentry{きまる}{決まる}{きまるん\tl{˥}}
\rentry{きみ}{君}{だー\tl{˥˩}}
\rentry{きみたち}{君たち}{だいま\tl{˥˩}}
\rentry{きめる}{決める}{きみるん\tl{˥}、さだみるん\tl{˥}}
\rentry{きも}{肝}{きむ\tl{˥}、すむ\tl{˥}}
\rentry{きもいり}{肝入り}{すむ\tl{˥}いり\tl{˥˩}}
\rentry{きもち}{気持}{きー\tl{˥˩}}
\rentry{きもちいい}{気持ちいい}{めっさ\tl{˩˥}はん}
\rentry{きもの}{着物}{すぅぬ\tl{˥}}
\rentry{きものをきる}{着物を着る}{すぅぬ\tl{˥}\ws{}すっすん\tl{˥˩}}
\rentry{ぎゃく}{逆}{さかさー、めー\tl{˩˥}しぴ\tl{˥˩}}
\rentry{ぎゃくたいする}{虐待する}{いじみるん\tl{˥˩}}
\rentry{きゅう}{灸}{やち\tl{˩˥}}
\rentry{きゅうかん}{旧慣}{むがしぃなれー\tl{˩˥}}
\rentry{きゅうくつだ}{窮屈だ}{くちさ\tl{˥˩}はん}
\rentry{きゅうけい}{休憩}{やしみ\tl{˩˥}}
\rentry{きゅうけいする}{休憩する}{どぅー\tl{˩˥}よがすん\tl{˩˥}、よがすん\tl{˩˥}}
\rentry{きゅうこんする}{求婚する}{くいるん\tl{˥}}
\rentry{きゅうし}{急死}{あたしに}
\rentry{きゅうしゅう}{旧習}{むがしぃなれー\tl{˩˥}}
\rentry{きゅうじゅう}{九十}{ぐじゅー\tl{˩˥}}
\rentry{きゅうじゅうななさいのちょうじゅいわい}{九十七歳の長寿祝い}{かじまやー\tl{˥}}
\rentry{きゅうしょ}{急所}{すぅぷ\tl{˥˩}}
\rentry{きゅうす}{急須}{ぞっか\tl{˩˥}}
\rentry{きゅうそくする}{休息する}{よがすん\tl{˩˥}}
\rentry{きゅうに}{急に}{あた\tl{˥˩}}
\rentry{きゅうにちいさくなるさま}{急に小さくなる様}{だった}
\rentry{きゅうになぐさま}{急に凪ぐ様}{だった}
\rentry{きゅうぼん}{旧盆}{そーりん\tl{˥}}
\rentry{きゅうぼんのゆにげーぎょうじ}{旧盆の〈ゆにげー〉「世願い」行事}{むしゃまー\tl{˥}}
\rentry{きゅうよう}{休養}{ぷにやしみ\tl{˥}}
\rentry{きゅうり}{胡瓜}{うりん\tl{˥}}
\rentry{きゅうれきのさんがつみっか}{旧暦の三月三日}{さにち\tl{˥˩}}
\rentry{きゅうれきのたなばた}{旧暦の七夕}{なんがそーりん\tl{˩˥}}
\rentry{きょう}{今日}{きゅー\tl{˥}}
\rentry{きょうかい}{境界}{さけー\tl{˥˩}}
\rentry{きょうくんする}{教訓する}{えにすかすん\tl{˥˩}}
\rentry{きょうげん}{狂言}{こんぎ\tl{˥}}
\rentry{きょうさくのとし}{凶作の年}{がしぃどぅしぃ}
\rentry{きょうじ}{凶事}{やなくとぅ\tl{˩˥}、やふ\tl{˩˥}}
\rentry{ぎょうじ}{行事}{ゆーじ\tl{˩˥}}
\rentry{きょうしゅくだ}{恐縮だ}{どぅぐりしゃ\tl{˩˥}はーん}
\rentry{きようだ}{器用だ}{〔手が\ntilde{}〕くま\tl{˥}はん、しー\tl{˥}ぐま\tl{˩˥}はん}
\rentry{きょうだいしまい}{兄弟姉妹}{うちざ\tl{˥}}
\rentry{きょうどうさぎょう}{共同作業}{ぼー\tl{˩˥}、ゆい\tl{˥˩}}
\rentry{ぎょうむ}{業務}{しぐとぅ\tl{˥˩}、すかま\tl{˥}}
\rentry{きょうり}{郷里}{すぅま\tl{˥}}
\rentry{きょうりょくだ}{強力だ}{すさ\tl{˥}はん}
\rentry{ぎょうれつ}{行列}{ぎょーれつ\tl{˥˩}}
\rentry{ぎょうれつをつくる}{行列をつくる}{するいるん\tl{˥}}
\rentry{きょか}{許可}{ゆるし\tl{˩˥}}
\rentry{きょかする}{許可する}{ゆるすん\tl{˩˥}}
\rentry{ぎょぎょう}{漁業}{ゆーとぅり\tl{˥˩}}
\rentry{ぎょぐ}{漁具}{いしょーだぐ\tl{˥˩}}
\rentry{きょげん}{虚言}{ぷりむに\tl{˥˩}}
\rentry{きょじゃく}{虚弱}{びーら、よーば\tl{˩˥}}
\rentry{きょせいする}{去勢する}{たに\tl{˥}\ws{}とぅるん\tl{˥}}
\rentry{きょねん}{去年}{くつん\tl{˥}}
\rentry{きよめる}{清める}{きゆみるん\tl{˥}}
\rentry{ぎょるい}{魚類}{ゆー\tl{˥˩}}
\rentry{ぎょろう}{漁労}{いしょー\tl{˥˩}}
\rentry{きらい}{嫌い}{みっふぁはん}
\rentry{きらわれもの}{嫌われ者}{\josho{}みっふぁー\kako{}むん}
\rentry{きり}{錐}{いり\tl{˥}}
\rentry{ぎりかたい}{義理堅い}{ぎりかた\tl{˩˥}はん}
\rentry{きりたおす}{切り倒す}{きりとーすん\tl{˥}、けーり\tl{˥}\ws{}とーすん\tl{˥}、けーるん\tl{˥}、なぎとーすん\tl{˩˥}}
\rentry{きりとる}{切り取る}{しっしとぅるん\tl{˥}}
\rentry{きりひらく}{切り開く}{さくるん\tl{˥}}
\rentry{きる}{切る}{しっすん\tl{˥}}
\rentry{きる}{着る}{しっすん\tl{˥˩}}
\rentry{きれい}{}{けしゃ\tl{˥}はん}
\rentry{きれいずきだ}{きれい好きだ}{あざぎしゃ\tl{˥}はん}
\rentry{きれいに}{}{けーし}
\rentry{きれる}{切れる}{むっしるん\tl{˥˩}}
\rentry{きわだつ}{際立つ}{みだつん\tl{˩˥}}
\rentry{きをつける}{気をつける}{きー\tl{˥˩}\ws{}しきるん\tl{˥}、くくるいるん}
\rentry{きん}{斤}{きん\tl{˥˩}}
\rentry{ぎん}{銀}{なんざ\tl{˩˥}}
\rentry{ぎんが}{銀河}{じんぬ\tl{˥˩}\ws{}ふかー\tl{˥}}
\rentry{きんきょり}{近距離}{すか\tl{˥}はん}
\rentry{きんすう}{斤数}{きん\tl{˥˩}}
\rentry{きんぞく}{金属}{かに\tl{˥˩}}
\rentry{ぎんみ}{吟味}{ぎんみ\tl{˩˥}}
\rentry{きんむする}{勤務する}{すとぅみるん\tl{˥}}
\rentry{くいあます}{食い余す}{へー\tl{˥˩}\ws{}あますん\tl{˥}}
\rentry{くいいじがきたない}{食い意地が汚い}{へだま\tl{˥˩}さーん}
\rentry{くいきる}{食い切る}{ふしきるん\tl{˥}}
\rentry{くいしんぼう}{食いしん坊}{へだま\tl{˥˩}}
\rentry{くいつく}{食いつく}{ふしきるん\tl{˥}}
\rentry{くいのこす}{食い残す}{へー\tl{˥˩}\ws{}あますん\tl{˥}}
\rentry{くいはずす}{食い外す}{へーぱんちるん\tl{˥˩}}
\rentry{ぐうぐう}{}{ごーごー}
\rentry{くうふくだ}{空腹だ}{やーは\tl{˩˥}\ws{}すん\tl{˥˩}、やー\tl{˩˥}はん}
\rentry{くかく}{区画}{ぱか\tl{˥}、〔田の\ntilde{}〕まーしぃ\tl{˥˩}}
\rentry{くき}{茎}{ふき\tl{˥}}
\rentry{くぎ}{釘}{ふん\tl{˥˩}}
\rentry{くぐりぬける}{潜り抜ける}{ふきるん\tl{˥˩}}
\rentry{くくる}{括る}{ふくるん\tl{˥˩}}
\rentry{くぐる}{潜る}{ふきるん\tl{˥˩}}
\rentry{くげん}{苦言}{すーむに\tl{˥}}
\rentry{くさ}{草}{ふつぁ\tl{˥}、ふつぁんだに\tl{˥}}
\rentry{くさい}{臭い}{ふつぁ\tl{˥}はん}
\rentry{くさび}{楔}{ふつぁび、や\tl{˥˩}}
\rentry{くさむら}{草むら}{ふつぁらし\tl{˥}}
\rentry{くさらせる}{腐らせる}{ふつぁらすん\tl{˥}}
\rentry{くさりかかる}{腐りかかる}{にーまらん\tl{˥˩}、にんまるん\tl{˥˩}}
\rentry{ぐさりと}{}{ざっふぁた}
\rentry{くさる}{腐る}{〔食物が\ntilde{}〕にーまらん\tl{˥˩}、ふつぁりん\tl{˥}}
\rentry{くし}{櫛}{ふち\tl{˥}}
\rentry{くじ}{公事}{うやだり\tl{˥}}
\rentry{くしゃみをする}{}{ぱな\tl{˥˩}ぴぃすん\tl{˥}}
\rentry{くじら}{鯨}{ぐじぃら\tl{˩˥}}
\rentry{くじる}{抉る}{ぴんくるん\tl{˥}}
\rentry{ぐずぐずすること}{}{ゆだばり\tl{˥˩}}
\rentry{くすぐったい}{擽ったい}{こさ\tl{˥}はん}
\rentry{くすぐる}{擽る}{ぐちゅるん\tl{˥˩}}
\rentry{くずす}{崩す}{こっふぁすん\tl{˥˩}}
\rentry{くすり}{薬}{ふちり\tl{˥}}
\rentry{くずれる}{崩れる}{こーりるん\tl{˥}、こっふぃるん\tl{˥˩}}
\rentry{くそ}{糞}{ふつ\tl{˥}}
\rentry{くそくさい}{糞臭い}{ふつ\tl{˥}ふつぁ\tl{˥}はーん}
\rentry{くそをたれる}{糞を垂れる}{ふつ\tl{˥}まるん\tl{˩˥}}
\rentry{くだく}{砕く}{くだくん、しきだるん\tl{˥}、ふくなすん\tl{˥}、ふなぐん\tl{˥}}
\rentry{くださる}{下さる}{たぼらいるん\tl{˥}、たぼるん\tl{˥}}
\rentry{くだす}{下す}{〔お腹を\ntilde{}〕くなすん\tl{˥˩}}
\rentry{ぐだっと}{}{だーり}
\rentry{くたばる}{}{うがんちゅむん\tl{˥}}
\rentry{くたびれ}{草臥れ}{くたんでー\tl{˥}}
\rentry{くたびれる}{草臥れる}{くたんでぃるん\tl{˥}}
\rentry{くだもの}{果物}{なーり\tl{˩˥}、なーりむぬ\tl{˩˥}}
\rentry{くだりざか}{下り坂}{さんがり\tl{˥}}
\rentry{くち}{口}{ふち\tl{˥˩}}
\rentry{くちがおもい}{口が重い}{ふちぃくぱはん}
\rentry{くちがかるい}{口が軽い}{ふちぃかろ\tl{˥˩}はん}
\rentry{くちぎたない}{口汚い}{ふちぃやにっしゃーん}
\rentry{くちたっしゃだ}{口達者だ}{ふちぃ\tl{˥˩}やば\tl{˩˥}はん}
\rentry{くちにかきこむ}{口にかき込む}{はき\tl{˥}いりるん\tl{˥˩}}
\rentry{くちのおおきなみずがめ}{口の大きな水甕}{ばんどー\tl{˥˩}}
\rentry{くちびる}{唇}{すぅぱ\tl{˥}、ふちぬ\tl{˥˩}\ws{}すぱ\tl{˥}}
\rentry{くちべただ}{口下手だ}{ふちぃくぱはん}
\rentry{くちる}{朽ちる}{ふとぅちるん\tl{˥}}
\rentry{ぐちをいう}{愚痴を言う}{むに\tl{˩˥}\ws{}ゆむん\tl{˩˥}}
\rentry{くつ}{靴}{ふちぃ\tl{˥}}
\rentry{ぐったり}{}{だーり}
\rentry{ぐったりする}{}{だーりるん\tl{˩˥}}
\rentry{くっつく}{くっ付く}{だっくゎーるん\tl{˥}、むっつぁーるん\tl{˩˥}、むっつぁるん\tl{˩˥}}
\rentry{くっつける}{くっ付ける}{すぱしきるん\tl{˥˩}、むっつぁーすん\tl{˩˥}}
\rentry{くば}{クバ}{くぱ\tl{˥˩}}
\rentry{くばがさ}{クバ笠}{くばかつぁ\tl{˥˩}}
\rentry{くばる}{配る}{くぱるん\tl{˥}}
\rentry{くび}{首}{ぬぶしん\tl{˩˥}、まんしん\tl{˩˥}}
\rentry{くびれる}{}{よーがりるん\tl{˥˩}}
\rentry{くふう}{工夫}{しかき\tl{˥˩}}
\rentry{くべる}{}{〔薪などを\ntilde{}〕てしきるん\tl{˥˩}}
\rentry{くぼち}{窪地}{とー\tl{˥˩}}
\rentry{くぼませる}{窪ませる}{とまらすん}
\rentry{くぼむ}{窪む}{とぅまるん\tl{˥˩}}
\rentry{くまい}{供米}{ぱなぐみ\tl{˥}}
\rentry{くみする}{組する}{かたすん\tl{˥˩}}
\rentry{くむ}{汲む}{〔水を\ntilde{}〕ふむん\tl{˥˩}}
\rentry{くめん}{工面}{さんだん}
\rentry{くめんする}{工面する}{さんだん\ws{}すん}
\rentry{くも}{クモ}{すぅふく\tl{˥}}
\rentry{くも}{雲}{ふもん\tl{˥}}
\rentry{くものいと}{クモの糸}{すぅふく\tl{˥}}
\rentry{くもる}{曇る}{ふゎどぅまり\tl{˥˩}\ws{}なるん\tl{˩˥}、ふゎむん\tl{˥˩}}
\rentry{くら}{倉}{ふっふぁ\tl{˥}}
\rentry{くら}{鞍}{ふっふぁ\tl{˥}}
\rentry{くらい}{暗い}{ふわ\tl{˥˩}はん}
\rentry{くらくなる}{暗くなる}{ふゎどぅまり\tl{˥˩}\ws{}なるん\tl{˩˥}、ふゎむん\tl{˥˩}}
\rentry{くらげ}{クラゲ}{いら\tl{˥}}
\rentry{くらべる}{比べる}{くらびるん\tl{˥˩}}
\rentry{くりぶね}{くり舟}{いたふに\tl{˥}、さばに\tl{˥}}
\rentry{くる}{来る}{くん\tl{˥}}
\rentry{くるしい}{苦しい}{くちさ\tl{˥˩}はん}
\rentry{くるとし}{来る年}{えん\tl{˩˥}、げん}
\rentry{くるまぼう}{車棒}{くるまぼー\tl{˥}}
\rentry{くるりぼう}{くるり棒}{くるまぼー\tl{˥}}
\rentry{くれる}{暮れる}{〔日が\ntilde{}〕ふゎむん\tl{˥˩}}
\rentry{くろい}{黒い}{ふっかばり\tl{˥}}
\rentry{くろう}{苦労}{あわり\tl{˥}}
\rentry{くろき}{黒木}{きな\tl{˥˩}}
\rentry{くろく}{黒く}{\josho{}ふー\kako{}し}
\rentry{くろくなる}{黒くなる}{\josho{}ふー\kako{}し\ws{}なるん\tl{˩˥}}
\rentry{くろざとう}{黒砂糖}{ふーさた}
\rentry{くろしま}{黒島}{ふすぅま\tl{˥}}
\rentry{くろずむ}{黒ずむ}{\josho{}ふー\kako{}し、\josho{}ふー\kako{}し\ws{}なるん\tl{˩˥}、ふっかばり\tl{˥}}
\rentry{くろつぐ}{クロツグ}{まーに\tl{˩˥}}
\rentry{くろつぐのせんい}{クロツグの繊維}{ふがら\tl{˥}}
\rentry{くろむ}{黒む}{\josho{}ふー\kako{}し\ws{}なるん\tl{˩˥}}
\rentry{くろよな}{クロヨナ}{ぶがま}
\rentry{くわ}{鍬}{ぺー\tl{˥˩}}
\rentry{くわえる}{}{〔口に\ntilde{}〕ふしきるん\tl{˥}}
\rentry{くわえる}{加える}{くわいるん\tl{˥}、まんざーすん}
\rentry{くわずいも}{クワズイモ}{びりかなぱ\tl{˩˥}}
\rentry{くわのき}{桑の木}{こんぎ\tl{˥}}
\rentry{くんかい}{訓戒}{ういし\tl{˥}}
\rentry{くんかいする}{訓戒する}{ういし\ws{}うがますん}
\rentry{くんわ}{訓話}{ゆしぐとぅ\tl{˥˩}}
\rentry{け}{毛}{きー\tl{˥˩}}
\rentry{けいかする}{経過する}{たつん\tl{˥}}
\rentry{けいこ}{稽古}{きっく\tl{˥}}
\rentry{けいさつ}{警察}{きっしい\tl{˥˩}}
\rentry{けいさん}{計算}{さんみん\tl{˥}}
\rentry{けいじ}{慶事}{\josho{}いー\kako{}くとぅ}
\rentry{けいしゃする}{傾斜する}{かたふくん\tl{˥}}
\rentry{けいそくする}{計測する}{ぱかるん\tl{˥}}
\rentry{けいそつだ}{軽率だ}{うかっと\ws{}すん}
\rentry{げいのう}{芸能}{ぎなむぬ\tl{˩˥}}
\rentry{けいべつする}{軽蔑する}{うせーるん\tl{˥˩}}
\rentry{けいらん}{鶏卵}{ごかけー\tl{˥}}
\rentry{けいれんする}{痙攣する}{ぴきしぃきん\tl{˥˩}}
\rentry{けがれる}{}{ゆぐりるん\tl{˥˩}}
\rentry{げし}{夏至}{かーち\tl{˥}}
\rentry{けす}{消す}{けすん\tl{˥˩}}
\rentry{けずる}{削る}{きちぃるん\tl{˥˩}、ぴぐん\tl{˥}}
\rentry{けた}{桁}{きた\tl{˥}}
\rentry{げた}{下駄}{あした\tl{˥}}
\rentry{けちだ}{}{ゆぐ\tl{˩˥}さん}
\rentry{けつえき}{血液}{じぃー\tl{˥˩}}
\rentry{けっかん}{血管}{じぃぬ\ws{}みち\tl{˥˩}}
\rentry{けつがん}{結願}{きちがん\tl{˥˩}}
\rentry{けっこんいわい}{結婚祝い}{\josho{}あいなー\kako{}よい}
\rentry{けっこんしき}{結婚式}{\josho{}あいなー\kako{}よい、にーびち\tl{˩˥}}
\rentry{けっこんする}{結婚する}{ぶとぅむち}
\rentry{けっしんする}{決心する}{うむいきすん\tl{˥}}
\rentry{けっそんする}{欠損する}{そん\tl{˥}\ws{}すん\tl{˥˩}}
\rentry{けっていする}{決定する}{きみるん\tl{˥}}
\rentry{けっとう}{血統}{ぴき\tl{˥˩}、まり\tl{˥˩}}
\rentry{げっとう}{月桃}{さーに\tl{˥}}
\rentry{けとばす}{蹴飛ばす}{きりとぅぱすん\tl{˥}}
\rentry{けなす}{貶す}{うせーん\tl{˥˩}}
\rentry{げなん}{下男}{ばたさ\tl{˩˥}}
\rentry{けぶる}{煙る}{きぷさ\tl{˥˩}はーん}
\rentry{けむたい}{煙たい}{きぷさ\tl{˥˩}はーん}
\rentry{けむりをたてる}{煙を立てる}{ふもーらすん\tl{˥}}
\rentry{げりする}{下痢する}{くなすん\tl{˥˩}、さぎるん\tl{˥}、とろふかすん\tl{˥}}
\rentry{けりたおす}{蹴り倒す}{きりとーすん\tl{˥}}
\rentry{ける}{蹴る}{きるん\tl{˥}}
\rentry{けれども}{}{えすが}
\rentry{げん}{弦}{すーる\tl{˥˩}}
\rentry{けんかする}{喧嘩する}{えんだりるん\tl{˥}、えんだりん\tl{˥}}
\rentry{げんきがなくなる}{元気がなくなる}{だりるん\tl{˩˥}}
\rentry{げんきだ}{元気だ}{がんずさ\tl{˩˥}はん}
\rentry{げんきづける}{元気付ける}{きーいりん\tl{˥˩}}
\rentry{げんきをなくす}{元気をなくす}{がんどーりん\tl{˩˥}}
\rentry{げんご}{言語}{くとぅば\tl{˥}、むに\tl{˩˥}}
\rentry{けんこうきがん}{健康祈願}{\josho{}どぅばだ\kako{}にげー}
\rentry{けんこうだ}{健康だ}{どぅー\tl{˩˥}ずさ\tl{˩˥}はん}
\rentry{げんこつ}{拳骨}{こーさー\tl{˥}}
\rentry{けんさ}{検査}{きんさ\tl{˥}}
\rentry{げんざい}{現在}{なま\tl{˥}、まな\tl{˥}}
\rentry{けんじつに}{堅実に}{こーし}
\rentry{げんしょうする}{減少する}{ぴなるん\tl{˥˩}}
\rentry{けんじょうひん}{献上品}{うさぎ\tl{˥˩}むぬ\tl{˩˥}、おーすむぬ\tl{˥˩}}
\rentry{げんだい}{現代}{まなぬ\tl{˥}\ws{}ゆー\tl{˥˩}}
\rentry{けんとう}{検討}{ぎんみ\tl{˩˥}}
\rentry{げんのう}{玄翁}{かにち\tl{˥˩}}
\rentry{けんまする}{研磨する}{とぅぐん\tl{˥}}
\rentry{こ}{子}{うたま\tl{˥}}
\rentry{～こ}{～個}{ちー}
\rentry{こい}{濃い}{かた\tl{˥˩}はん}
\rentry{ごいさぎ}{ゴイサギ}{よーらさー\tl{˩˥}}
\rentry{こいし}{小石}{いしふく\tl{˥˩}、いしん\tl{˥˩}たま\tl{˥}}
\rentry{こう}{請う}{くいるん\tl{˥}}
\rentry{こうい}{行為}{しわざ\tl{˥˩}、わざ\tl{˩˥}}
\rentry{こういか}{甲烏賊}{くしみ\tl{˥}}
\rentry{こううん}{幸運}{うん\tl{˥˩}}
\rentry{こうか}{硬貨}{かにじん\tl{˥˩}}
\rentry{こうかいあんぜん}{航海安全}{かりゆし\tl{˥}}
\rentry{こうかな}{高価な}{でーだが\tl{˥˩}}
\rentry{こうかんする}{交換する}{かれーるん\tl{˥˩}}
\rentry{こうぎ}{抗議}{だんぱん\tl{˥˩}}
\rentry{こうぎする}{抗議する}{かきあうん\tl{˥}}
\rentry{ごうけい}{合計}{すーだが\tl{˥}}
\rentry{こうけいしゃ}{後継者}{あとぅ\tl{˥}ちぎ\tl{˥˩}}
\rentry{こうごうしい}{神々しい}{みしこ\tl{˥}はん}
\rentry{こうさい}{交際}{ぴぃとぅぴれー\tl{˥}、ぴらい\tl{˥}、ぴらす}
\rentry{こうさく}{耕作}{さくほー\tl{˥}}
\rentry{こうささせる}{交差させる}{あんざらすん\tl{˥}、あんじるん\tl{˥}}
\rentry{こうさんする}{降参する}{まきるん\tl{˥˩}}
\rentry{こうじつにする}{口実にする}{なしきるん\tl{˥}}
\rentry{こうじょう}{口上}{すさりぶち\tl{˥˩}}
\rentry{ごうじょうだ}{強情だ}{がーずさはん}
\rentry{こうしょうぶる}{高尚ぶる}{いばりすくん\tl{˥˩}}
\rentry{こうしょくのじょし}{好色の女子}{しきべー}
\rentry{こうせい}{後生}{ぐそー\tl{˩˥}}
\rentry{こうぞう}{構造}{すくり\tl{˥}}
\rentry{こうたいする}{後退する}{さがるん\tl{˥}}
\rentry{こうち}{高地}{うがり\tl{˥˩}、た\josho{}か\kako{}た}
\rentry{こうちょくする}{硬直する}{こっぱるん\tl{˥}}
\rentry{こうつうきかん}{交通機関}{ぬりむぬ\tl{˥˩}}
\rentry{こうつごう}{好都合}{\josho{}いー\kako{}ば}
\rentry{こうとうぶ}{後頭部}{うっす\tl{˥}}
\rentry{こうびする}{交尾する}{ずーぶん\tl{˥}}
\rentry{こうへい}{公平}{まーたき}
\rentry{こうほう}{後方}{しなた\tl{˥˩}、しぴゃた\tl{˥˩}}
\rentry{こうむ}{公務}{うやだり\tl{˥}}
\rentry{こうもり}{コウモリ}{かぷる\tl{˥}}
\rentry{こうもん}{肛門}{しぴぬ\tl{˥˩}みん\tl{˩˥}}
\rentry{ごうりきしゃ}{強力者}{ぶし\tl{˥˩}}
\rentry{こうりゃん}{高粱}{や\josho{}た\kako{}ぷ}
\rentry{こうろへいあん}{航路平安}{かりゆし\tl{˥}}
\rentry{こうろん}{口論}{むんどぅー\tl{˩˥}}
\rentry{こうろんする}{口論する}{むんどぅー\ws{}すん}
\rentry{こえ}{声}{くい\tl{˥}}
\rentry{こえだめ}{肥溜}{こいすぷ\tl{˥˩}}
\rentry{こえび}{小エビ}{せー\tl{˥}}
\rentry{こえる}{肥える}{ぱなたるん\tl{˥}、ぱんたるん\tl{˥}}
\rentry{こえる}{越える}{くいるん\tl{˥˩}、こすん\tl{˥˩}}
\rentry{こがす}{焦がす}{くがすん\tl{˥˩}}
\rentry{こがたな}{小刀}{しぐ\tl{˥}}
\rentry{こがね}{黄金}{くがに\tl{˥˩}}
\rentry{こがま}{小鎌}{〔粟刈りの\ntilde{}〕いら\tl{˥}}
\rentry{こがれる}{焦がれる}{くがりるん\tl{˥˩}}
\rentry{こかん}{股間}{またびし\tl{˩˥}}
\rentry{ごきぶり}{ゴキブリ}{かむし\tl{˥}}
\rentry{こきょう}{故郷}{まりじま\tl{˥˩}}
\rentry{こく}{}{〔屁を\ntilde{}〕ぴぃすん\tl{˥}}
\rentry{こく}{扱く}{〔稲を\ntilde{}〕しゅるん\tl{˥}}
\rentry{こぐ}{漕ぐ}{くん\tl{˥}}
\rentry{こくおう}{国王}{うしゅがなし\tl{˥˩}}
\rentry{こくする}{濃くする}{〔酒などを\ntilde{}〕しぃみるん\tl{˥}}
\rentry{こくとう}{黒糖}{さた\tl{˥}、ふーさた}
\rentry{こくのあるうまみ}{コクのある旨味}{すまじしゃ\tl{˥}はん}
\rentry{こくもつ}{穀物}{\josho{}ぷーぬ\kako{}\ws{}むぬ}
\rentry{ごくんごくん}{}{ごんた}
\rentry{こけ}{苔}{ぬり\tl{˩˥}}
\rentry{こげる}{焦げる}{くがりるん\tl{˥˩}、なびふつぁりん\tl{˩˥}}
\rentry{ここ}{}{なー\tl{˥}、もー\tl{˥}}
\rentry{ごご}{午後}{ぴぃすあとぅ}
\rentry{こごえる}{凍える}{〔寒さで\ntilde{}〕くぱるん\tl{˥}、ぴんぐるん\tl{˥}}
\rentry{ここちよい}{心地よい}{めっさ\tl{˩˥}はん}
\rentry{ここに}{此処に}{なー\tl{˥}}
\rentry{ここのつ}{九つ}{くくぬち\tl{˥}}
\rentry{こころ}{心}{くくる\tl{˥}、すむ\tl{˥}}
\rentry{こころえる}{心得る}{くくるいるん}
\rentry{こころがいたむ}{心が痛む}{すむやむーん}
\rentry{こころがうつくしい}{心が美しい}{すむ\tl{˥}けしゃ\tl{˥}はーん}
\rentry{こころがける}{心掛ける}{くくるいるん、くくるがきるん}
\rentry{こころがやさしい}{心が優しい}{すむ\tl{˥}けしゃ\tl{˥}はーん}
\rentry{こころづよい}{心強い}{すむ\tl{˥}ずさ\tl{˥}はーん、すむ\tl{˥}すさ\tl{˥}はん}
\rentry{こころね}{心根}{きむぐくる\tl{˥}}
\rentry{こさえる}{}{くつぁすん\tl{˥˩}}
\rentry{こざかな}{小魚}{ゆんたま}
\rentry{こさせる}{越させる}{こすん\tl{˥˩}}
\rentry{こざっぱりする}{}{あざぎしゃ\tl{˥}はん}
\rentry{こし}{腰}{くち\tl{˥˩}}
\rentry{ごしゃくまい}{五勺米}{ぐそーみ\tl{˩˥}}
\rentry{ごじゅう}{五十}{ぐんず\tl{˩˥}}
\rentry{こしらえる}{拵える}{〔魚などを\ntilde{}〕くつぁーすん\tl{˥˩}、くつぁすん\tl{˥˩}、しこーるん\tl{˥˩}}
\rentry{こじん}{古人}{むがしぃ\tl{˩˥}ぴぃとぅ\tl{˥˩}}
\rentry{こす}{漉す}{こすん\tl{˥˩}}
\rentry{こす}{濾す}{〔篩で\ntilde{}〕ふかすん\tl{˥˩}}
\rentry{こす}{越す}{こすん\tl{˥˩}}
\rentry{こする}{擦る}{しっすするん\tl{˥˩}、すするん\tl{˥˩}、とーぴくん\tl{˥}}
\rentry{こそだてする}{子育てする}{うたま\tl{˥}\ws{}すかなすん\tl{˥}}
\rentry{こたえる}{答える}{くたいるん\tl{˥}}
\rentry{こだかいだいち}{小高い台地}{うがり\tl{˥˩}}
\rentry{こたち}{子たち}{\josho{}うたま\kako{}んじ}
\rentry{ごちそう}{ご馳走}{こっきー\tl{˥}、まさむぬ\tl{˥}}
\rentry{こちょこちょ}{}{ぐちゅぐちゅ}
\rentry{こちらから}{}{もーら\tl{˥}}
\rentry{こちらへ}{}{もが\tl{˥}}
\rentry{こづかい}{小使い}{ばたさ\tl{˩˥}}
\rentry{こっかく}{骨格}{ぷに\tl{˥}}
\rentry{こつぶだ}{小粒だ}{ぐま\tl{˩˥}はん}
\rentry{こつん}{}{ごーんた}
\rentry{こときれる}{こと切れる}{ぬちぃ\ws{}しさん}
\rentry{ことし}{今年}{くとぅしぃ\tl{˥˩}}
\rentry{ことづける}{言付ける}{くとぅしきるん\tl{˥}}
\rentry{ことば}{言葉}{くとぅば\tl{˥}、むに\tl{˩˥}}
\rentry{ことばがあらい}{言葉が荒い}{むに\tl{˩˥}あら\tl{˥˩}はーん}
\rentry{ことばがあらっぽい}{言葉が荒っぽい}{むに\tl{˩˥}あら\tl{˥˩}はーん}
\rentry{こども}{子供}{うたま\tl{˥}、やらび\tl{˩˥}}
\rentry{こどもたち}{子供たち}{\josho{}うたま\kako{}んじ}
\rentry{こどもをうむ}{子供を産む}{うたま\tl{˥}\ws{}なすん\tl{˩˥}}
\rentry{ことわざ}{諺}{むがし\tl{˩˥}くとぅば\tl{˥}}
\rentry{ことわり}{}{どーり\tl{˩˥}}
\rentry{ことわる}{断る}{くとぅばるん\tl{˥}}
\rentry{こな}{粉}{く\tl{˥}、ふく\tl{˥}}
\rentry{こなごなにする}{粉々にする}{ふくなすん\tl{˥}}
\rentry{こなごなになる}{粉々になる}{ふくなるん\tl{˩˥}}
\rentry{こなごなにわれる}{粉々に割れる}{ふく\ws{}なるん}
\rentry{こなす}{}{くなすん\tl{˥˩}}
\rentry{この}{}{くぬ\tl{˥˩}}
\rentry{このあいだ}{この間}{くんた\tl{˥˩}ばり\tl{˩˥}}
\rentry{このふきん}{この付近}{もぬ\tl{˥}\ws{}まーり\tl{˥˩}}
\rentry{このへん}{この辺}{もー\tl{˥}、もぬ\tl{˥}\ws{}まーり\tl{˥˩}}
\rentry{このむ}{好む}{すくん\tl{˥}}
\rentry{ごはん}{ご飯}{いー\tl{˥}}
\rentry{こぶし}{拳}{しふく\tl{˥}}
\rentry{こふん}{古墳}{めーむるし\tl{˩˥}}
\rentry{ごぼう}{牛蒡}{ぐぼん\tl{˩˥}}
\rentry{こぼす}{零す}{〔液体を残りなく\ntilde{}〕いたちるん\tl{˥˩}、くぱすん\tl{˥}、〔全部\ntilde{}〕んたちるん\tl{˥˩}}
\rentry{こぼれる}{零れる}{くぷりるん\tl{˥}、くぷりん\tl{˥}}
\rentry{こぼれるようす}{零れる様子}{〔粒など小さな物が\ntilde{}〕ざらざら}
\rentry{こま}{独楽}{こーろ\tl{˥}}
\rentry{こまかい}{細かい}{くま\tl{˥}はん}
\rentry{こまったこと}{困ったこと}{ざーふぇー\tl{˩˥}}
\rentry{ごみ}{}{あるふた\tl{˥˩}}
\rentry{こむぎ}{小麦}{まーむん}
\rentry{こむぎこ}{小麦粉}{むんぬ\tl{˩˥}\ws{}くー\tl{˥}}
\rentry{こめ}{米}{めー\tl{˥˩}}
\rentry{こめこのむしがし}{米粉の蒸し菓子}{ありすむち\tl{˥˩}}
\rentry{こめだわら}{米俵}{めーだら\tl{˩˥}}
\rentry{こめぬか}{米糠}{ぬが\tl{˩˥}}
\rentry{こめのめし}{米の飯}{めーぬ\tl{˥˩}\ws{}いー\tl{˥}}
\rentry{ごめんください}{ご免下さい}{すされー}
\rentry{こもりのむすめ}{子守の娘}{むりあま\tl{˩˥}}
\rentry{こもる}{籠る}{くまるん\tl{˥}}
\rentry{こやし}{肥し}{こい\tl{˥}}
\rentry{こやすがい}{子安貝}{しぴ\tl{˥˩}}
\rentry{こゆび}{小指}{びびんたま\tl{˥}}
\rentry{こようする}{雇用する}{やとすん\tl{˩˥}}
\rentry{こようのいっしゅ}{古謡の一種}{あよー\tl{˥}}
\rentry{こようのしゅるい}{古謡の種類}{ゆんぐとぅ\tl{˩˥}、ゆんた\tl{˩˥}}
\rentry{こようのろうどうか}{古謡の労働歌}{じらば\tl{˩˥}}
\rentry{ごようふ}{御用布}{ぐいふ\tl{˩˥}}
\rentry{こらい}{古来}{むがしぃがら\tl{˩˥}}
\rentry{こられる}{来られる}{おーるん\tl{˥}}
\rentry{これの}{}{くぬ\tl{˥˩}}
\rentry{ころ}{}{じぶん\tl{˥˩}}
\rentry{ころがす}{転がす}{くるばすん\tl{˥˩}}
\rentry{ころがる}{転がる}{くるぶん\tl{˥˩}、まるぶちん\tl{˥˩}}
\rentry{ころす}{殺す}{くらすん\tl{˥˩}、すなすん\tl{˥˩}}
\rentry{ころばす}{転ばす}{くるばすん\tl{˥˩}}
\rentry{ころぶ}{転ぶ}{くるぶん\tl{˥˩}、まるぶちん\tl{˥˩}、まるぶん\tl{˥˩}}
\rentry{こわい}{怖い}{うとぅるさ\tl{˥}はん、ごー\tl{˥}はん}
\rentry{こわす}{壊す}{こーすん、こっふぁすん\tl{˥˩}、やぶるん\tl{˩˥}}
\rentry{こわばる}{強張る}{こっぱるん\tl{˥}}
\rentry{こわめし}{こわ飯}{かしき\tl{˥}}
\rentry{こわれる}{壊れる}{こーりるん\tl{˥}、やぶりん\tl{˩˥}}
\rentry{こんがらがる}{}{あんざるん\tl{˥}}
\rentry{こんき}{根気}{くんき\tl{˥˩}}
\rentry{こんごうする}{混合する}{まんざるん\tl{˩˥}、まんじるん\tl{˩˥}}
\rentry{こんじょう}{根性}{がー\tl{˩˥}}
\rentry{こんちゅう}{昆虫}{むし\tl{˥˩}}
\rentry{こんなんだ}{困難だ}{むすかっさ\tl{˩˥}はん}
\rentry{こんばん}{今晩}{にが\tl{˩˥}}
\rentry{さ}{差}{さ\tl{˥}}
\rentry{ざ}{座}{ざー\tl{˩˥}}
\rentry{～さ}{}{さ}
\rentry{さあー}{}{せー}
\rentry{ざぁざぁ}{ザァザァ}{ぞーりぞーり}
\rentry{～さい}{～歳}{ちー}
\rentry{ざいさん}{財産}{ざいさん\tl{˥˩}}
\rentry{さいしゅする}{採種する}{たに\tl{˥}\ws{}とぅるん\tl{˥}}
\rentry{さいそくする}{催促する}{いみるん\tl{˥}}
\rentry{ざいにん}{罪人}{とぅがにん}
\rentry{さいばい}{栽培}{むぬすくり\tl{˩˥}}
\rentry{さいふ}{財布}{じんふくる\tl{˩˥}}
\rentry{さいほう}{裁縫}{ぬーん\tl{˩˥}、ぬいむぬ\tl{˩˥}}
\rentry{ざいもく}{材木}{ざいぎ\tl{˩˥}}
\rentry{さいれい}{祭礼}{ぶなが\tl{˥}}
\rentry{さえずる}{囀る}{〔鳥が\ntilde{}〕なーぐん\tl{˥˩}、ふきん\tl{˥}}
\rentry{さお}{竿}{そー\tl{˥}}
\rentry{さおばかり}{竿秤}{ぴき\tl{˥˩}}
\rentry{さかい}{境}{さけー\tl{˥˩}}
\rentry{さかえる}{栄える}{さかいるん\tl{˥˩}}
\rentry{さかさ}{逆さ}{さかさー}
\rentry{さがしまわる}{探し回る}{ふつるん\tl{˥}}
\rentry{さがしもとめる}{捜し求める}{たしかみるん\tl{˥}}
\rentry{さがす}{探す}{あさるん\tl{˥}、さばくん\tl{˥}、たじぃにるん\tl{˥}、とぅみるん\tl{˥}、とぅみん\tl{˥}}
\rentry{さかずき}{盃}{さかすけー\tl{˥˩}}
\rentry{さかな}{肴}{うせー\tl{˥}}
\rentry{さかな}{魚}{ゆー\tl{˥˩}}
\rentry{さかなつき}{魚突き}{ゆーしき}
\rentry{さかなつり}{魚釣り}{ゆーほすん}
\rentry{さかなとり}{魚獲り}{ゆーとぅり\tl{˥˩}}
\rentry{さかなのえら}{えら}{〔魚の\ntilde{}〕ぴむ\tl{˥}}
\rentry{さかなのはらにく}{魚の腹肉}{はらごー\tl{˥}}
\rentry{さかなのひもの}{魚の干物}{てー\tl{˥}}
\rentry{さかり}{盛り}{ばんじぃん\tl{˥˩}}
\rentry{さがる}{下がる}{さがるん\tl{˥}}
\rentry{さがん}{砂岩}{あーいし\tl{˥}}
\rentry{さき}{先}{さき\tl{˥˩}、すら\tl{˥}、ぱな\tl{˥˩}、ぱなた\tl{˥}}
\rentry{さき}{崎}{さき\tl{˥˩}、ぱな\tl{˥˩}}
\rentry{さきにする}{先にする}{めー\tl{˩˥}\ws{}なすん\tl{˩˥}}
\rentry{さきになる}{先になる}{めー\tl{˩˥}\ws{}なるん\tl{˩˥}}
\rentry{さきまわり}{先回り}{さきまーり\tl{˥˩}}
\rentry{さきやま}{崎山}{さきやま\tl{˥}}
\rentry{さぎょうぎ}{作業着}{ぴてーすぅぬ\tl{˥}}
\rentry{さく}{咲く}{さくん\tl{˥˩}}
\rentry{さく}{裂く}{さくん\tl{˥}}
\rentry{さくせいする}{作成する}{すくるん\tl{˥}}
\rentry{さくねん}{昨年}{くつん\tl{˥}}
\rentry{さくや}{昨夜}{ゆーび\tl{˩˥}}
\rentry{さくりと}{}{ぷちんた}
\rentry{さぐる}{探る}{さぐるん\tl{˥˩}}
\rentry{さけ}{酒}{さき\tl{˥˩}}
\rentry{さけくさい}{酒臭い}{さきふつぁ\tl{˥˩}はーん}
\rentry{さけぐせ}{酒癖}{さきぐし\tl{˥˩}}
\rentry{さけじょうご}{酒上戸}{さきじょーぐ\tl{˥˩}}
\rentry{さけなどがはっこうする}{酒などが発酵する}{ばーぐん\tl{˩˥}}
\rentry{さけによう}{酒に酔う}{びたこら\tl{˩˥}}
\rentry{さけぶ}{叫ぶ}{たけーるん\tl{˥}}
\rentry{さけめ}{裂け目}{やりふちぃ\tl{˩˥}}
\rentry{さけやすい}{裂けやすい}{さぱ\tl{˥˩}はん}
\rentry{さける}{裂ける}{〔布などが古くなって\ntilde{}〕ざーりるん\tl{˩˥}、さきるん\tl{˥}、ざりるん\tl{˩˥}、やーりん\tl{˩˥}}
\rentry{ささえる}{支える}{しけーるん\tl{˥˩}}
\rentry{ささえるちから}{支える力}{てー\tl{˥}}
\rentry{さしあげる}{差し上げる}{うさぎるん\tl{˥˩}、おーすん\tl{˥˩}}
\rentry{ざしき}{座敷}{ざしき\tl{˩˥}}
\rentry{さしず}{指図}{ぎし\tl{˩˥}}
\rentry{さしずする}{指図する}{えにすかすん\tl{˥˩}}
\rentry{さしだす}{差し出す}{んだすん\tl{˥}}
\rentry{さしば}{鸇}{たか\tl{˥˩}}
\rentry{さしみ}{刺身}{なましぃ\tl{˩˥}}
\rentry{さす}{刺す}{さすん\tl{˥}、ぬぐん\tl{˩˥}}
\rentry{さす}{差す}{さすん\tl{˥}、しっすん\tl{˥˩}}
\rentry{さす}{挿す}{さすん\tl{˥}}
\rentry{さずける}{授ける}{たぼらいるん\tl{˥}}
\rentry{さそり}{蠍}{しんがま\tl{˥}}
\rentry{さだまる}{定まる}{きまるん\tl{˥}}
\rentry{さだめ}{定め}{さだみ\tl{˥}}
\rentry{さだめる}{定める}{さだみるん\tl{˥}}
\rentry{さっさと}{}{さーった}
\rentry{さっさといきおいよくあるくようす}{さっさと勢いよく歩く様子}{さーっさーった}
\rentry{ざっそう}{雑草}{ふつぁ\tl{˥}、ふつぁんだに\tl{˥}}
\rentry{ざっそうをねこそぎにとる}{雑草を根こそぎに取る}{〔ヘラで\ntilde{}〕そーるん\tl{˥˩}}
\rentry{ざつだん}{雑談}{ゆんたく\tl{˩˥}}
\rentry{さっとかんつうするさま}{さっと貫通するさま}{ぼふぁった、ぼふった}
\rentry{さっぱり}{}{むっとぅ\tl{˩˥}}
\rentry{さつまいも}{さつま芋}{あがん\tl{˩˥}}
\rentry{さとう}{砂糖}{さた\tl{˥}}
\rentry{さとうきび}{サトウキビ}{あますな\tl{˥}}
\rentry{さとうてんぷら}{砂糖てんぷら}{さたぱんべー\tl{˥}}
\rentry{さとのし}{里之子}{さとぬし\tl{˥}}
\rentry{さとる}{悟る}{さとぅるん\tl{˥}}
\rentry{さばく}{}{さばくん\tl{˥}}
\rentry{さび}{錆}{さび\tl{˥}}
\rentry{さびしがる}{寂しがる}{しからさ\tl{˥}\ws{}すん\tl{˥˩}、しから\tl{˥}はーん}
\rentry{さびつく}{錆付く}{さびすくん}
\rentry{さびれる}{寂れる}{さぼーりん\tl{˥˩}}
\rentry{ざぶざぶ}{}{ざっふぁざっふぁ}
\rentry{ざぶっと}{}{ざっふぁった、ざぼんた}
\rentry{ざぶんと}{}{だぼんた}
\rentry{さます}{冷ます}{さますん\tl{˥}、ぴらすん\tl{˥}}
\rentry{さます}{覚ます}{さますん\tl{˥}}
\rentry{さまたげる}{妨げる}{さまたぎるん\tl{˥˩}}
\rentry{さみしい}{淋しい}{しから\tl{˥}はん}
\rentry{さみだれ}{五月雨}{ゆどぅあみ\tl{˩˥}}
\rentry{さむい}{寒い}{ぴしゃ\tl{˥}はん}
\rentry{さむくてふるえる}{寒くて震える}{ふひつぁるん}
\rentry{さむけがする}{寒気がする}{ふぃーつぁるん}
\rentry{さむさ}{寒さ}{かん\tl{˥}}
\rentry{さめる}{冷める}{さみるん\tl{˥}}
\rentry{さめる}{覚める}{さみるん\tl{˥}}
\rentry{さめる}{醒める}{さまるん\tl{˥}}
\rentry{さゆ}{白湯}{さゆ\tl{˥}}
\rentry{さら}{皿}{さらー\tl{˥˩}}
\rentry{さらいねん}{再来年}{まー\ws{}みちん}
\rentry{ざらざらする}{}{ざらざーら\ws{}すん}
\rentry{さらす}{晒す}{さらすん\tl{˥˩}、〔水に\ntilde{}〕さらすん\tl{˥˩}}
\rentry{さる}{去る}{んぐん\tl{˥˩}}
\rentry{さる}{猿}{さーる\tl{˥}}
\rentry{さる}{申}{さり\tl{˥}}
\rentry{ざる}{}{ばーき\tl{˩˥}}
\rentry{さるかけみかん}{サルカケミカン}{じろーし\tl{˥˩}}
\rentry{さわがしくする}{騒がしくする}{あーらすん\tl{˥˩}}
\rentry{さわがす}{騒がす}{あーらすん\tl{˥˩}、さわがすん\tl{˥˩}}
\rentry{さわぎ}{騒ぎ}{そーどー\tl{˥˩}}
\rentry{さわぐ}{騒ぐ}{あーるん\tl{˥˩}、なーるん\tl{˥˩}}
\rentry{さわら}{サワラ}{さーら\tl{˥}}
\rentry{さわる}{触る}{さーるん\tl{˥˩}}
\rentry{さん}{桟}{さん\tl{˥}}
\rentry{さんきらい}{山帰来}{さんきら\tl{˥}}
\rentry{さんご}{珊瑚}{うる\tl{˥˩}}
\rentry{さんごのいし}{珊瑚の石}{うる\tl{˥˩}}
\rentry{さんじゅう}{三十}{さんず\tl{˥}}
\rentry{さんじゅうさんねんき}{三十三年忌}{あぎごっこー\tl{˥˩}}
\rentry{さんしん}{三線}{さんしん\tl{˥˩}}
\rentry{ざんねん}{残念}{あがやー、いなむん}
\rentry{ざんぶ}{残部}{ぬぐり\tl{˩˥}、のごり}
\rentry{さんぽう}{三方}{\josho{}さん\kako{}ぼー}
\rentry{さんらんさせる}{散乱させる}{しきぽーるん\tl{˥˩}}
\rentry{じ}{字}{じぃー\tl{˥˩}}
\rentry{しあげる}{仕上げる}{しあぎるん\tl{˥˩}}
\rentry{しあん}{思案}{かんげー\tl{˥}、むぬかんげー\tl{˩˥}}
\rentry{しいくする}{飼育する}{すかなすん\tl{˥}}
\rentry{しいる}{強いる}{うししきるん\tl{˥˩}}
\rentry{じうたい}{地謡}{じうてー\tl{˩˥}}
\rentry{しお}{塩}{まーす\tl{˩˥}}
\rentry{しお}{潮}{すー\tl{˥}}
\rentry{しおがひくこと}{潮が引く}{すーぴすん\tl{˥}}
\rentry{しおがみちる}{潮が満ちる}{すー\tl{˥}んつん\tl{˩˥}}
\rentry{しおからい}{塩辛い}{さこら\tl{˥}はん}
\rentry{しおどき}{潮時}{すっち\tl{˥}}
\rentry{しおに}{塩煮}{まーすにー\tl{˩˥}}
\rentry{しおひがり}{潮干狩り}{あさらごー\tl{˥}}
\rentry{しおみず}{潮水}{ぶす\tl{˥}}
\rentry{しおみぶそく}{塩味不足}{あふぁ\tl{˥}さん}
\rentry{しおれる}{萎れる}{〔植物が\ntilde{}〕だーりん\tl{˩˥}、だりるん\tl{˩˥}}
\rentry{しかけ}{仕掛け}{しかき\tl{˥˩}}
\rentry{しかける}{仕掛ける}{しかきるん\tl{˥˩}}
\rentry{しかし}{}{えすが}
\rentry{しかた}{仕方}{しかた\tl{˥˩}}
\rentry{しかむ}{}{しかむん\tl{˥}}
\rentry{しかりつける}{叱りつける}{しゃみしきるん、しゃみしきん\tl{˥˩}、すかり\tl{˥˩}すくん\tl{˥}}
\rentry{しかりつける}{𠮟りつける}{しゃーみるん\tl{˥˩}}
\rentry{しかる}{叱る}{うくるん\tl{˥}、ゆい\tl{˥˩}}
\rentry{しきさい}{色彩}{いる\tl{˥}}
\rentry{しく}{敷く}{すくん\tl{˥˩}}
\rentry{しげきしておこらせる}{刺激して怒らせる}{〔人を\ntilde{}〕おこらすん\tl{˥}}
\rentry{しげる}{繁る}{さかいるん\tl{˥˩}}
\rentry{しげる}{茂る}{かぷん\tl{˥}}
\rentry{しごきおとす}{しごき落とす}{しゅるん\tl{˥}}
\rentry{しごく}{}{しゅるん\tl{˥}}
\rentry{しごと}{仕事}{しかま\tl{˥}、しぐとぅ\tl{˥˩}、すかま\tl{˥}}
\rentry{しごとする}{仕事する}{ぱたらぐん\tl{˥˩}}
\rentry{しごとをかんりょうする}{仕事を完了する}{しー\tl{˥}うすくま\tl{˥˩}}
\rentry{しこむ}{仕込む}{しくむん\tl{˥˩}}
\rentry{しさい}{司祭}{にんぶち\tl{˩˥}}
\rentry{しさんか}{資産家}{むぬむち\tl{˩˥}}
\rentry{じさんさせる}{持参させる}{むたすん\tl{˩˥}}
\rentry{じさんする}{持参する}{むつん\tl{˩˥}}
\rentry{しし}{獅子}{しーし\tl{˥}}
\rentry{じし}{次姉}{やしゃ\tl{˩˥}}
\rentry{ししょうじ}{指小辞}{たま}
\rentry{じしょくする}{辞職する}{やみるん\tl{˥˩}}
\rentry{じしん}{地震}{にん\tl{˩˥}}
\rentry{しずかに}{静かに}{や\josho{}まー\kako{}し}
\rentry{しずかになる}{静かになる}{とぅりるん\tl{˥˩}}
\rentry{しずませる}{沈ませる}{すますん\tl{˥}}
\rentry{しずむ}{沈む}{すむん\tl{˥}}
\rentry{しせいじ}{私生児}{ぐんぼー\tl{˩˥}}
\rentry{した}{舌}{した\tl{˥}}
\rentry{したあご}{下あご}{すたはこち\tl{˥˩}}
\rentry{じだい}{時代}{ゆー\tl{˥˩}}
\rentry{したくする}{支度する}{したく\ws{}すん}
\rentry{したくない}{}{しぼへぬ}
\rentry{したたらせる}{滴らせる}{すたらすん\tl{˥}}
\rentry{したたる}{滴る}{あまだりん\tl{˥}、すたるん\tl{˥}}
\rentry{したてる}{仕立てる}{したてぃるん\tl{˥˩}}
\rentry{したにむけておく}{下に向けて置く}{うすふかすん\tl{˥˩}}
\rentry{したへ}{下へ}{したいが\tl{˥˩}}
\rentry{じっか}{実家}{やーむとぅ\tl{˩˥}}
\rentry{しっかり}{}{\josho{}うん\kako{}とぅ}
\rentry{しつかり}{しっかり}{\josho{}しかい\kako{}とぅ、だ\josho{}ん\kako{}た}
\rentry{しっかりと}{}{だーんた}
\rentry{しっくい}{漆喰}{むち\tl{˥˩}}
\rentry{しつける}{躾ける}{すかり\tl{˥˩}すくん\tl{˥}}
\rentry{じつげんする}{実現する}{かのーん\tl{˥}}
\rentry{しっこうする}{膝行する}{しきっつぁーるん}
\rentry{しっしんする}{失神する}{ふきーるん\tl{˥˩}}
\rentry{しったことか}{知ったことか}{たすたがー\tl{˥˩}}
\rentry{しっている}{知っている}{ししゃん\tl{˥˩}}
\rentry{しっとする}{嫉妬する}{たい\tl{˥}\ws{}すん\tl{˥˩}}
\rentry{しっぽ}{尻尾}{ぶすぷ\tl{˩˥}}
\rentry{しつれい}{失礼}{ぶりー\tl{˩˥}}
\rentry{してやられる}{}{しらいん\tl{˥˩}}
\rentry{しとげる}{し遂げる}{しー\ws{}うすくまん、すますん\tl{˥}}
\rentry{しとしと}{}{すとぅりすとぅり}
\rentry{しな}{品}{しな\tl{˥}}
\rentry{しなおす}{し直す}{しー\tl{˥˩}のーすん\tl{˩˥}}
\rentry{しなのき}{シナノキ}{きなし\tl{˥}}
\rentry{しなびる}{萎びる}{しぃぴりるん\tl{˥˩}、しぴりん\tl{˥˩}}
\rentry{しにはてる}{死に果てる}{しにぶたん}
\rentry{じにんする}{辞任する}{やみるん\tl{˥˩}}
\rentry{しぬ}{死ぬ}{すぅぬん\tl{˥˩}}
\rentry{しばしばいく}{しばしば行く}{しゃー\tl{˥}\ws{}んぐん\tl{˥˩}}
\rentry{しはらう}{支払う}{ぱらすん\tl{˥}}
\rentry{しばらく}{暫く}{あ\josho{}たー\kako{}すま}
\rentry{しばりつける}{縛りつける}{さましきるん}
\rentry{しばる}{縛る}{さまるん\tl{˥}、ふくるん\tl{˥˩}、むすぷん\tl{˥˩}}
\rentry{しびれる}{痺れる}{ぴくら\tl{˥˩}やむん\tl{˩˥}、ぺしくらやむん\tl{˩˥}}
\rentry{しぶい}{渋い}{すぽ\tl{˥}はん}
\rentry{しぶみがある}{渋みがある}{すぽ\tl{˥}はん}
\rentry{じぶん}{時分}{じぶん\tl{˥˩}}
\rentry{じぶん}{自分}{どぅー\tl{˩˥}、なーどぅーどぅ、ばー\tl{˩˥}}
\rentry{じぶんかって}{自分勝手}{どぅー\tl{˩˥}かってぃ\tl{˥˩}}
\rentry{じぶんかってしている}{自分勝手している}{どぅーかってぃ\ws{}しゃーん}
\rentry{じぶんかってだ}{自分勝手だ}{どぅーかってぃ\ws{}しゃーん}
\rentry{じぶんで}{自分で}{\josho{}どぅー\kako{}し、どぅしけな}
\rentry{しぼうする}{死亡する}{ぬちぃ\ws{}しさん}
\rentry{しぼむ}{萎む}{〔花などが\ntilde{}〕だーりるん\tl{˩˥}}
\rentry{しぼる}{搾る}{すぷるん\tl{˥}}
\rentry{しま}{島}{すぅま\tl{˥}}
\rentry{しまい}{姉妹}{ぶなり\tl{˥˩}}
\rentry{しまいこむ}{仕舞い込む}{かつぁみるん\tl{˥}}
\rentry{しまう}{仕舞う}{うすくみん\tl{˥˩}、うわるん\tl{˥˩}}
\rentry{しまことば}{島言葉}{すま\tl{˥}むに\tl{˩˥}}
\rentry{しみこむ}{しみ込む}{〔味などが\ntilde{}〕ふくむん}
\rentry{しみったれだ}{}{ゆぐ\tl{˩˥}さん}
\rentry{じめじめ}{}{びちゃびちゃ}
\rentry{じめじめする}{}{〔畑が\ntilde{}〕ふこ\tl{˥˩}はん}
\rentry{しめつける}{締めつける}{しみしきるん\tl{˥}、しみるん\tl{˥}}
\rentry{しめっぽくなる}{湿っぽくなる}{しみるん\tl{˥}}
\rentry{しめらせる}{湿らせる}{しみらすん\tl{˥}}
\rentry{しめる}{湿る}{しみるん\tl{˥}、すたるん\tl{˥}}
\rentry{しめる}{締める}{しみるん\tl{˥}}
\rentry{しめる}{閉める}{ふーん\tl{˥˩}}
\rentry{しゃくし}{杓子}{すーぺー\tl{˥}}
\rentry{しゃくにさわる}{}{くさむくん\tl{˥}}
\rentry{しゃくようする}{借用する}{かるん\tl{˥˩}}
\rentry{しゃげきする}{射撃する}{うつん\tl{˥}}
\rentry{しゃこがい}{シャコガイ}{あしけー\tl{˥}}
\rentry{じゃまする}{邪魔する}{さまたぎるん\tl{˥˩}}
\rentry{じゃまになる}{邪魔になる}{ふさがるん\tl{˥˩}}
\rentry{しゃもじ}{杓文字}{しゃしゃびら\tl{˥}}
\rentry{じゃり}{砂利}{いしん\tl{˥˩}たま\tl{˥}}
\rentry{しゅいうふやく}{首里大屋子}{しなばく\tl{˥˩}}
\rentry{しゅうい}{周囲}{まーましぃ\tl{˥˩}、まーり\tl{˩˥}}
\rentry{じゅういち}{十一}{とぅーぴとぅち}
\rentry{じゅういっこ}{十一個}{とぅーぴとぅち}
\rentry{しゅうかくする}{収穫する}{ぶるん\tl{˩˥}}
\rentry{しゅうかん}{習慣}{なれー\tl{˩˥}}
\rentry{しゅうき}{秋季}{すさんち\tl{˥}}
\rentry{じゅうぎょういん}{従業員}{しんか\tl{˥}}
\rentry{しゅうごうさせる}{集合させる}{するいるん\tl{˥}}
\rentry{じゅうごや}{十五夜}{じぐや\tl{˩˥}}
\rentry{しゅうじつ}{終日}{ぴてんぴじゅ}
\rentry{じゅうしょ}{住所}{ぶーどぅり\tl{˩˥}}
\rentry{しゅうのうする}{収納する}{うすくみん\tl{˥˩}}
\rentry{じゅうばこ}{重箱}{じばぐ\tl{˩˥}}
\rentry{じゅうぶんにじゅくする}{十分に熟する}{けった\ws{}みゃーん}
\rentry{しゅうりょうする}{終了する}{おわるん\tl{˥˩}}
\rentry{じゅうろくにちさい}{十六日祭}{じりく\tl{˩˥}にち\tl{˥˩}}
\rentry{じゅくしきる}{熟しきる}{けった\ws{}みゃーん}
\rentry{じゅくすいする}{熟睡する}{たーりん\tl{˥˩}}
\rentry{じゅくする}{熟する}{みゃん\tl{˥}}
\rentry{しゅくてん}{祝典}{よい\tl{˩˥}}
\rentry{しゅくはくする}{宿泊する}{とぅまるん\tl{˥˩}、やどぅとぅるん}
\rentry{じゅごん}{ジュゴン}{ざん\tl{˩˥}、ざんぬ\tl{˩˥}ゆ\tl{˥˩}}
\rentry{しゅし}{種子}{たに\tl{˥}}
\rentry{じゅし}{呪詞}{ぱん\tl{˥}}
\rentry{しゅちょうする}{主張する}{やっぱり\tl{˥˩}}
\rentry{しゅっこうする}{出港する}{ふにんじるん}
\rentry{しゅっさん}{出産する}{うたま\tl{˥}\ws{}なすん\tl{˩˥}}
\rentry{しゅっせきする}{出席する}{んじるん\tl{˥}、んじん\tl{˥}}
\rentry{しゅっぱつする}{出発する}{んじたつん\tl{˥}}
\rentry{じゅもく}{樹木}{きー\tl{˥}}
\rentry{しゅるい}{種類}{しな\tl{˥}}
\rentry{しゅろ}{棕櫚}{しゅる\tl{˥˩}}
\rentry{じゅんかいする}{巡回する}{まーり\tl{˥˩}\ws{}みるん\tl{˩˥}}
\rentry{しゅんき}{春季}{うりじん\tl{˥}}
\rentry{じゅんさ}{巡査}{ゆんさん\tl{˩˥}}
\rentry{しゅんじに}{瞬時に}{たでーま}
\rentry{じゅんばん}{順番}{ばん\tl{˩˥}、まーり\tl{˥˩}}
\rentry{じゅんびする}{準備する}{しくむん\tl{˥˩}、しこーるん\tl{˥˩}、するいるん\tl{˥}}
\rentry{じょいん}{女陰}{ぴー\tl{˥}}
\rentry{しょうがつ}{正月}{さごち}
\rentry{じょうがぶかい}{情が深い}{じょー\tl{˥˩}ふか\tl{˥˩}はーん}
\rentry{しょうかふりょう}{消化不良}{ばたふくら}
\rentry{しょうき}{正気}{そー\tl{˥}}
\rentry{じょうき}{蒸気}{きぷ\tl{˥˩}}
\rentry{じょうぎ}{定規}{じょーぎ\tl{˩˥}}
\rentry{じょうきげんになる}{上機嫌になる}{さーふーふー\ws{}すん\tl{˥˩}}
\rentry{しょうきょする}{消去する}{けすん\tl{˥˩}}
\rentry{しょうきをうしなう}{正気を失う}{どまんぐるん\tl{˩˥}}
\rentry{じょうげぎゃく}{上下逆}{すたうい}
\rentry{しょうご}{正午}{ぴぃすまり\tl{˥˩}}
\rentry{じょうご}{上戸}{じょーぐ\tl{˩˥}}
\rentry{じょうし}{上司}{ういぬ\tl{˥˩}\ws{}ぴぃとぅ\tl{˥˩}}
\rentry{しょうじきだ}{正直だ}{まとーば\tl{˥˩}}
\rentry{じょうず}{上手}{ぞーじ\tl{˩˥}、やば\tl{˩˥}はん}
\rentry{しようする}{使用する}{すこーん\tl{˥˩}}
\rentry{しょうそく}{消息}{うとぅさた\tl{˥˩}、さた\tl{˥}}
\rentry{しょうたいきゃく}{招待客}{しん\tl{˥}}
\rentry{じょうだん}{冗談}{ふばむに\tl{˥}}
\rentry{しょうたんだ}{小胆だ}{いじぃぬ\tl{˥˩}\ws{}ねーぬ\tl{˩˥}、いじぃ\ws{}ねーぬ、すむ\tl{˥}いしゃが\tl{˥}はーん}
\rentry{しょうち}{礁池}{いのー\tl{˥˩}}
\rentry{しょうちょう}{小腸}{ばちばた\tl{˩˥}}
\rentry{しょうてん}{商店}{まちやー\tl{˥˩}}
\rentry{じょうとう}{上等}{ぞーとぅ\tl{˩˥}}
\rentry{しょうにゅうせき}{鍾乳石}{いんぬ\tl{˥˩}\ws{}まら\tl{˩˥}}
\rentry{しょうぶ}{勝負}{しーぶ\tl{˥}}
\rentry{じょうふ}{情婦}{ゆーべー\tl{˩˥}}
\rentry{じょうぶだ}{丈夫だ}{がんずさ\tl{˩˥}はん}
\rentry{じょうぶになる}{丈夫になる}{がんずさ\tl{˩˥}\ws{}なるん\tl{˩˥}}
\rentry{じょうほう}{上方}{うい\tl{˥˩}}
\rentry{じょうほうへ}{上方へ}{ういが\tl{˥˩}}
\rentry{しょうまん}{小満}{すーまん\tl{˥}}
\rentry{しょうめん}{正面}{たんか\tl{˥˩}}
\rentry{しょうゆ}{醤油}{したち\tl{˥}}
\rentry{しょか}{初夏}{ばがなち\tl{˩˥}}
\rentry{しょき}{暑気}{ぷーき}
\rentry{しょき}{書記}{ぴっしゃ\tl{˥}}
\rentry{しょくじ}{食事}{むぬ\tl{˩˥}}
\rentry{しょくす}{食酢}{ぺーる\tl{˩˥}}
\rentry{しょくぶん}{職分}{しゅくぶん\tl{˩˥}}
\rentry{しょくゆ}{食油}{あば\tl{˥}}
\rentry{しょくよくがないこと}{食欲がないこと}{ふち\tl{˥˩}やぶら\tl{˩˥}}
\rentry{しょくりょう}{食糧}{ぱめー\tl{˥˩}、ほーむぬ\tl{˥˩}}
\rentry{しょげかえる}{悄げかえる}{だりるん\tl{˩˥}}
\rentry{しょげる}{悄げる}{すむだーりるん\tl{˥˩}}
\rentry{じょせい}{女性}{みどぅむ\tl{˩˥}}
\rentry{じょそうする}{除草する}{〔へらで\ntilde{}〕そるん\tl{˥˩}}
\rentry{じょそうせいそう}{除草清掃}{〔御嶽や神道の\ntilde{}〕みやくつぇー\tl{˥˩}}
\rentry{しょっき}{食器}{かまさ\tl{˥}}
\rentry{しょもうする}{所望する}{ぬずむん}
\rentry{しょもつ}{書物}{すむち\tl{˥˩}}
\rentry{しらが}{白髪}{しっせー\tl{˥}}
\rentry{しらさぎ}{白鷺}{そん\tl{˥}さみゃ\tl{˩˥}}
\rentry{しらせる}{知らせる}{いーしきるん\tl{˥˩}、すかすん\tl{˥˩}}
\rentry{しらない}{知らない}{ばがらぬ\tl{˩˥}}
\rentry{しらべる}{調べる}{すさびるん\tl{˥}}
\rentry{しらほ}{白保}{すさぶ\tl{˥}}
\rentry{しらみ}{虱}{さん\tl{˥}}
\rentry{しらむ}{白む}{すさむん\tl{˥}}
\rentry{しり}{尻}{しぴ\tl{˥˩}}
\rentry{しりがおもい}{尻が重い}{しぴんさ\tl{˥˩}はーん}
\rentry{しりがかるい}{尻が軽い}{しぴかろ\tl{˥˩}はん}
\rentry{しりがるおんな}{尻軽女}{さんごなー\tl{˥}}
\rentry{しりぞく}{退く}{ぬぎるん\tl{˩˥}}
\rentry{しりのあな}{尻の穴}{しぴぬ\tl{˥˩}みん\tl{˩˥}}
\rentry{しる}{汁}{すー\tl{˥}}
\rentry{しる}{知る}{しっすん\tl{˥˩}}
\rentry{しるし}{印}{しるし\tl{˥˩}}
\rentry{しろあり}{白蟻}{すさーり\tl{˥}}
\rentry{しろい}{白い}{\josho{}すそー\kako{}し}
\rentry{しろいきのこ}{白いキノコ}{しっすみん\tl{˥}}
\rentry{しろくなる}{白くなる}{すさむん\tl{˥}}
\rentry{しろっぽい}{白っぽい}{\josho{}すそー\kako{}し}
\rentry{しろみをおびる}{白みを帯びる}{すさむん\tl{˥}}
\rentry{しわ}{皺}{ぴぃな\tl{˥}}
\rentry{しわざ}{仕業}{しわざ\tl{˥˩}}
\rentry{じわれ}{地割れ}{ぴぱり\tl{˥˩}}
\rentry{しん～}{真～}{まー}
\rentry{しんこうする}{信仰する}{しんじるん\tl{˥}}
\rentry{しんじつ}{真実}{ふんとー\tl{˥}}
\rentry{しんじょう}{心情}{きむぐくる\tl{˥}、すむぐくる\tl{˥}}
\rentry{しんじる}{信じる}{しんじるん\tl{˥}}
\rentry{しんせいだ}{神聖だ}{みしこ\tl{˥}はん、よみしゃ\tl{˩˥}はん}
\rentry{しんせき}{親戚}{うちざ\tl{˥}まり\tl{˥˩}、うやぐ\tl{˥}}
\rentry{しんぜんのくもつのな}{神前の供物の名}{くぱん\tl{˥}}
\rentry{しんぞう}{心臓}{まーみ\tl{˩˥}}
\rentry{しんぞく}{親族}{うやぐ\tl{˥}}
\rentry{しんちょう}{身長}{たき\tl{˥}}
\rentry{しんちょうに}{慎重に}{みしこーみしこ}
\rentry{しんどうする}{振動する}{ふっふん\tl{˥˩}}
\rentry{しんねん}{新年}{あらとぅし\tl{˥}}
\rentry{しんぱい}{心配}{しわー\tl{˥˩}、むぬむい\tl{˩˥}}
\rentry{しんぱいする}{心配する}{すむ\tl{˥}やむん\tl{˩˥}}
\rentry{しんぴん}{新品}{みーむん\tl{˥˩}、めーむぬ\tl{˥˩}}
\rentry{しんぶつのけいじ}{神仏の啓示}{むぬしらし\tl{˩˥}}
\rentry{しんゆう}{親友}{いすぱんどぅし\tl{˥}}
\rentry{しんようする}{信用する}{しんじるん\tl{˥}}
\rentry{しんりょく}{新緑}{ばがぱ\tl{˩˥}}
\rentry{しんるい}{親類}{うちざ\tl{˥}まり\tl{˥˩}}
\rentry{す}{巣}{しー\tl{˥}}
\rentry{す}{酢}{ぺーる\tl{˩˥}}
\rentry{すいじ}{炊事}{ぞーしき\tl{˥˩}}
\rentry{すいじごや}{炊事小屋}{とーら\tl{˥˩}}
\rentry{すいじする}{炊事する}{たくん\tl{˥˩}、まがすん\tl{˩˥}}
\rentry{すいじば}{炊事場}{ふちめー\tl{˥}}
\rentry{すいじゃくする}{衰弱する}{やっちるん\tl{˩˥}、よーがりん\tl{˥˩}、よーるん\tl{˩˥}}
\rentry{すいせい}{彗星}{ぽーぎぶしぃ\tl{˥}}
\rentry{すいだす}{吸い出す}{ゆびだすん\tl{˥˩}}
\rentry{すいちゅうめがね}{水中眼鏡}{がんきょー\tl{˥}、みーかんがん\tl{˩˥}}
\rentry{すいでん}{水田}{たなー\tl{˥}、たなが\tl{˥}}
\rentry{すいとる}{吸い取る}{すぷるん\tl{˥˩}}
\rentry{すいひおけ}{水肥おけ}{こいたんぐ\tl{˥˩}}
\rentry{すいふ}{水夫}{はこー\tl{˥˩}、ふなはこー}
\rentry{すいもの}{吸い物}{しむぬ\tl{˥˩}}
\rentry{すう}{吸う}{すぷるん\tl{˥˩}、〔水気を\ntilde{}〕ゆぶん\tl{˥˩}}
\rentry{すえる}{据える}{びしるん\tl{˥˩}}
\rentry{すえる}{饐える}{しーるん\tl{˥}}
\rentry{ずが}{図画}{えー\tl{˥˩}}
\rentry{すがた}{姿}{かたち\tl{˥˩}}
\rentry{すきぐし}{梳き櫛}{かとーし\tl{˥˩}}
\rentry{すぎさる}{過ぎ去る}{しぎるん\tl{˥}、すぎゃん\tl{˥}}
\rentry{すきだ}{好きだ}{すくん\tl{˥}}
\rentry{すきとおってみえる}{透き通って見える}{みーふかさりん\tl{˩˥}}
\rentry{すきま}{隙間}{あいま\tl{˥}}
\rentry{すく}{梳く}{〔櫛で髪を\ntilde{}〕きちぃるん\tl{˥˩}}
\rentry{すく}{鋤く}{すかすん\tl{˥˩}}
\rentry{すくう}{掬う}{すこーすん\tl{˥˩}}
\rentry{すくう}{救う}{たしきるん\tl{˥}}
\rentry{すくない}{少ない}{いきら\tl{˥}さん}
\rentry{すぐに}{}{たでーま、ぴく}
\rentry{すくむ}{}{すくむん}
\rentry{すぐれた}{}{かねー}
\rentry{すぐれている}{優れている}{やば\tl{˩˥}はん}
\rentry{すぐれる}{優れる}{すぐりん\tl{˥˩}、まさるん\tl{˥˩}}
\rentry{すこし}{少し}{べー\tl{˩˥}、べ\josho{}ー\kako{}な、\josho{}べー\kako{}び}
\rentry{すこしまえ}{少し前}{くんた\tl{˥˩}ばり\tl{˩˥}}
\rentry{すじ}{}{かち\tl{˥}}
\rentry{ずじょうにのせてはこぶ}{頭上に乗せて運ぶ}{かみん\tl{˥}}
\rentry{すす}{煤}{ししー\tl{˥}}
\rentry{すずかぜ}{涼風}{ぴりかち\tl{˥}}
\rentry{すすき}{ススキ}{ゆしき\tl{˥˩}}
\rentry{すずしい}{涼しい}{ぴりしゃ\tl{˥}はん}
\rentry{すずむ}{涼む}{ぴりしゃー\ws{}すん}
\rentry{すずめ}{雀}{みしゅ\tl{˩˥}どぅり\tl{˥˩}}
\rentry{すずめばち}{スズメバチ}{あがぱーち\tl{˥˩}}
\rentry{すずり}{硯}{しじり\tl{˥}}
\rentry{すする}{}{すぷるん\tl{˥˩}}
\rentry{すする}{啜る}{ゆぶん\tl{˥˩}}
\rentry{すそ}{裾}{すっす\tl{˥}}
\rentry{ずたずたにさく}{ずたずたに裂く}{さきつぁーるん\tl{˥˩}}
\rentry{ずつう}{頭痛}{あまっすくる\tl{˥}やみ\tl{˩˥}}
\rentry{ずつうがする}{頭痛がする}{おま\tl{˥}はん}
\rentry{すっかりわすれきる}{すっかり忘れきる}{けった\ws{}ばっさーん}
\rentry{すっくと}{}{やいっと}
\rentry{すっぱい}{酸っぱい}{しさ\tl{˥}はん}
\rentry{すてておく}{捨てておく}{だすたすくん\tl{˩˥}}
\rentry{すでに}{既に}{きさ\tl{˥}}
\rentry{すてる}{捨てる}{ししん\tl{˥˩}、なんが\tl{˩˥}\ws{}っしん\tl{˩}}
\rentry{すな}{砂}{いしょん\tl{˥˩}}
\rentry{すなはま}{砂浜}{ぱま\tl{˥}}
\rentry{すねる}{}{ひんすん\tl{˥˩}}
\rentry{すねる}{拗ねる}{むじかーるん\tl{˩˥}、むじかるん\tl{˩˥}}
\rentry{すばしこい}{}{からばっさ\tl{˥}はん}
\rentry{すばやく}{素早く}{さーった}
\rentry{すべらせる}{滑らせる}{ししらすん\tl{˥}}
\rentry{すべりやすい}{滑りやすい}{なふこ\tl{˩˥}はん}
\rentry{すべる}{滑る}{ししりん\tl{˥}、しっしりるん\tl{˥}}
\rentry{すます}{済ます}{すますん\tl{˥}}
\rentry{すます}{澄ます}{すますん\tl{˥}}
\rentry{すませる}{済ませる}{すますん\tl{˥}}
\rentry{すみ}{墨}{しん\tl{˥}}
\rentry{すみ}{炭}{たん\tl{˥˩}}
\rentry{すみ}{隅}{かどぅ\tl{˥}}
\rentry{すむ}{済む}{しまいるん\tl{˥˩}、すむん\tl{˥}}
\rentry{すむ}{澄む}{すむん\tl{˥}、ぴりがるん\tl{˥}}
\rentry{すもう}{相撲}{さぱぁ\tl{˥}}
\rentry{ずりおちる}{ずり落ちる}{しっしりるん\tl{˥}}
\rentry{する}{}{すん\tl{˥˩}}
\rentry{する}{擦る}{しっすするん\tl{˥˩}}
\rentry{するどくとがった}{鋭く尖った}{ぴんとまり}
\rentry{すれちがう}{すれ違う}{んぎちげーるん\tl{˥˩}}
\rentry{すわる}{座る}{びるん\tl{˥˩}}
\rentry{すんぽう}{寸法}{すんぽー\tl{˥}}
\rentry{せいいっぱい}{精一杯}{\josho{}うん\kako{}とぅ、しーっぺー\tl{˥}、\josho{}しかい\kako{}とぅ}
\rentry{せいけつだ}{清潔だ}{あざぎしゃ\tl{˥}はん}
\rentry{せいこうだ}{精巧だ}{〔技が\ntilde{}〕くま\tl{˥}はん}
\rentry{せいざ}{正座}{ぺしきびり\tl{˥˩}}
\rentry{せいしん}{精神}{くくる\tl{˥}、すむぐくる\tl{˥}}
\rentry{せいち}{生地}{まり\tl{˥˩}}
\rentry{せいちょうする}{成長する}{ぶーさ\tl{˥}\ws{}なるん\tl{˩˥}、ぶさは\tl{˥}\ws{}なるん\tl{˩˥}、ぶすとぅ\ws{}なるん}
\rentry{せいてい}{正丁}{ふだにん\tl{˥}}
\rentry{せいと}{生徒}{がくたま}
\rentry{せいとうする}{製糖する}{さたたくん}
\rentry{せいなる}{聖なる}{よみしゃ\tl{˩˥}はん}
\rentry{せいねん}{生年}{まりどぅしぃ\tl{˥˩}}
\rentry{せいねん}{青年}{ばがむぬ\tl{˩˥}}
\rentry{せいねんいわい}{生年祝い}{まりぶなー\tl{˥˩}}
\rentry{せいほう}{西方}{いらた、いり\tl{˥˩}}
\rentry{せいめい}{生命}{ぬちぃ\tl{˩˥}}
\rentry{せいようじん}{西洋人}{うらんだー\tl{˥}}
\rentry{せいり}{整理}{かたすか\tl{˥}}
\rentry{せいりする}{整理する}{まだぎるん\tl{˩˥}}
\rentry{せいれつする}{整列する}{ならぶん\tl{˥˩}}
\rentry{せいろ}{蒸籠}{せいろー\tl{˥}}
\rentry{せおわす}{背負わす}{〔牛馬に荷物を\ntilde{}〕おーすん\tl{˥}}
\rentry{せがむ}{}{きぬん\tl{˥˩}}
\rentry{せき}{咳}{さこー\tl{˥}}
\rentry{せき}{席}{ざー\tl{˩˥}}
\rentry{せきうん}{積雲}{ぬりふもん\tl{˩˥}}
\rentry{せきたてる}{急き立てる}{あーらすん\tl{˥˩}、あばたしみるん}
\rentry{せきたん}{石炭}{しきたん\tl{˥}}
\rentry{せきにんをおう}{責任を負う}{しきにんむつん、ぱくん\tl{˥˩}}
\rentry{せきにんをもつ}{責任を持つ}{しきにんむつん}
\rentry{せきゆ}{石油}{しきゆ\tl{˥}}
\rentry{せきらんうん}{積乱雲}{ぬりふもん\tl{˩˥}}
\rentry{せけん}{世間}{しきん\tl{˥}}
\rentry{せけんばなし}{世間話}{むぬぱなし\tl{˩˥}}
\rentry{せたい}{世帯}{きに\tl{˥˩}}
\rentry{せつ}{節}{しち\tl{˥}}
\rentry{せっかい}{石灰}{うるんぺー\tl{˥˩}}
\rentry{せっきょうする}{説教する}{いましみるん\tl{˥}、いましみん\tl{˥}}
\rentry{せっけん}{石鹸}{さっぷん\tl{˥}}
\rentry{せっせとはたらくさま}{せっせと働くさま}{やいやいた}
\rentry{せったいする}{接待する}{とぅりむつん\tl{˥}}
\rentry{せっちゃくする}{接着する}{すぱしきるん\tl{˥˩}}
\rentry{せっとくする}{説得する}{とぅしきるん\tl{˥}}
\rentry{せっぱん}{折半}{ふたばぎ\tl{˥˩}}
\rentry{せつまつり}{節祭}{ししん\tl{˥}}
\rentry{せつまつりでこぐくりぶね}{節祭で漕ぐくり舟}{しっしんふに\tl{˥}}
\rentry{ぜつめつする}{絶滅する}{たに\ws{}しさん}
\rentry{せなか}{背中}{いしなが\tl{˥}}
\rentry{ぜに}{銭}{じん\tl{˩˥}}
\rentry{せのたかいひと}{背の高い人}{たかぴぃとぅ\tl{˥}}
\rentry{せのび}{背伸び}{ぬーび\tl{˩˥}}
\rentry{せばめる}{狭める}{しぱみるん\tl{˥˩}、すぱみるん\tl{˥}}
\rentry{ぜひとも}{是非とも}{たんでー\tl{˥}、\josho{}どー\kako{}でぃん}
\rentry{せぼね}{背骨}{なーぶに\tl{˩˥}}
\rentry{せまくする}{狭くする}{しぱみるん\tl{˥˩}}
\rentry{せまる}{}{〔間近に、目前に\ntilde{}〕ぬすかるん\tl{˥˩}}
\rentry{せみ}{セミ}{しゃんしゃん\tl{˥}}
\rentry{せめつける}{責めつける}{しみしきるん\tl{˥}}
\rentry{せめる}{攻める}{しみるん\tl{˥}}
\rentry{せめる}{責める}{しみるん\tl{˥}}
\rentry{せん}{千}{しん\tl{˥˩}}
\rentry{ぜん}{膳}{じん\tl{˥˩}}
\rentry{せんい}{繊維}{かち\tl{˥}}
\rentry{せんいん}{船員}{はこー\tl{˥˩}、ふなはこー}
\rentry{せんげつ}{先月}{めーぬ\tl{˩˥}\ws{}しき\tl{˥}}
\rentry{ぜんご}{前後}{あとさき\tl{˥˩}、めー\tl{˩˥}しぴ\tl{˥˩}}
\rentry{せんこう}{線香}{こー\tl{˥}}
\rentry{せんこつ}{洗骨}{しんくつぇー\tl{˥}}
\rentry{せんじやく}{煎じ薬}{\josho{}しんじ\kako{}ふちり}
\rentry{せんしゅざいがうえにでたふね}{船首材が上に出た船}{かんぽーしん\tl{˥}}
\rentry{せんじる}{煎じる}{しんじるん\tl{˥}}
\rentry{せんせい}{先生}{しんしん\tl{˥}}
\rentry{せんぞ}{先祖}{がんす\tl{˩˥}}
\rentry{せんそう}{戦争}{いふつぁ\tl{˥}}
\rentry{せんぞしん}{先祖神}{うやぴぃとぅ\tl{˥}}
\rentry{せんぞのれい}{先祖の霊}{うやぴぃとぅ\tl{˥}}
\rentry{せんたい}{浅堆}{すねー\tl{˥}}
\rentry{せんたくする}{洗濯する}{あらすん\tl{˥˩}、すぅぬ\tl{˥}\ws{}あらすん\tl{˥˩}}
\rentry{せんたくする}{選択する}{いらぶん\tl{˥}}
\rentry{せんたん}{先端}{さき\tl{˥˩}、すら\tl{˥}、ぱなた\tl{˥}}
\rentry{せんどう}{船頭}{しんどー\tl{˥}}
\rentry{せんぱい}{先輩}{しゃまかた\tl{˥}}
\rentry{ぜんぶ}{全部}{\josho{}がす\kako{}た、けーら、ぼふた}
\rentry{ぜんぽう}{前方}{めー\tl{˩˥}}
\rentry{せんめんき}{洗面器}{みんだれー\tl{˩˥}}
\rentry{ぜんら}{全裸}{まるばい\tl{˥˩}}
\rentry{せんりょう}{染料}{あい\tl{˥˩}}
\rentry{ぜんりょくではしる}{全力で走る}{かしきるん\tl{˥}}
\rentry{そあくな}{粗悪な}{そーべー\tl{˥}}
\rentry{そう}{}{えー}
\rentry{そうおうする}{相応する}{うゆぶん\tl{˥˩}}
\rentry{そうか}{}{えーなー}
\rentry{そうがく}{総額}{すーだが\tl{˥}}
\rentry{そうぎ}{争議}{むんどぅー\tl{˩˥}}
\rentry{そうぎする}{争議する}{むんどぅー\ws{}すん}
\rentry{そうさせる}{}{\josho{}えー\kako{}\ws{}しみるん\tl{˥˩}}
\rentry{そうじ}{掃除}{そーじ\tl{˥}}
\rentry{そうしか}{}{\josho{}えー\kako{}ばぎる}
\rentry{そうしき}{葬式}{そーしき\tl{˥˩}}
\rentry{そうしきをする}{葬式をする}{うぐるん\tl{˥˩}}
\rentry{そうじする}{掃除する}{ぽーぐん\tl{˥}}
\rentry{そうして}{}{\josho{}え\kako{}した}
\rentry{そうしょく}{装飾}{かつぁり\tl{˥˩}}
\rentry{ぞうすい}{雑炊}{ずーし\tl{˩˥}}
\rentry{そうする}{}{えー\tl{˥}\ws{}すん\tl{˥˩}}
\rentry{そうせいする}{早世する}{はやじに\tl{˥}\ws{}すん\tl{˥˩}}
\rentry{そうそふ}{曽祖父}{ういぶや\tl{˥}}
\rentry{そうそぼ}{曽祖母}{ういぱー\tl{˥}}
\rentry{そうだ}{}{えーるやろ}
\rentry{そうだから}{}{\josho{}えす\kako{}がら、えちー\tl{˥˩}、えちる}
\rentry{そうだよ}{}{えーさ}
\rentry{そうだん}{相談}{そんだん\tl{˥}}
\rentry{そうだんする}{相談する}{そんだん\tl{˥}\ws{}すん\tl{˥˩}}
\rentry{そうではない}{}{あらぬ\tl{˥}}
\rentry{そうでも}{}{\josho{}えー\kako{}\ws{}やばん}
\rentry{そうどう}{騒動}{そーどー\tl{˥˩}}
\rentry{そうとうする}{相当する}{ひきあうん\tl{˥˩}}
\rentry{そうなら}{}{\josho{}えー\kako{}\ws{}やちゃら}
\rentry{そうめん}{素麵}{そーみん\tl{˥}}
\rentry{ぞうり}{草履}{さぱん\tl{˥˩}}
\rentry{そうりょ}{僧侶}{ぼーず\tl{˩˥}}
\rentry{そぐ}{}{しゅるん\tl{˥}}
\rentry{そくざに}{即座に}{たでーま}
\rentry{そこ}{底}{すく\tl{˥˩}}
\rentry{そこなう}{損なう}{あやみるん\tl{˥}、いたますん\tl{˥}、いたみるん\tl{˥}}
\rentry{そぜい}{租税}{ぞーぬ\tl{˩˥}}
\rentry{そせき}{礎石}{びしじ\tl{˥˩}}
\rentry{そそぐ}{注ぐ}{しっすん\tl{˥˩}}
\rentry{そだつ}{育つ}{すだつん\tl{˥}、ぶすとぅ\ws{}なるん}
\rentry{そだてる}{育てる}{すだちるん\tl{˥}、すなちるん\tl{˥}}
\rentry{そっきょうおどり}{即興踊り}{もーや\tl{˩˥}}
\rentry{そつぎょうする}{卒業する}{ぬがーるん\tl{˩˥}}
\rentry{そっちょくに}{率直に}{まんがたんが}
\rentry{そっとうする}{卒倒する}{ふきん\tl{˥˩}}
\rentry{そてつ}{蘇鉄}{すとぅち\tl{˥}}
\rentry{そと}{外}{ふか\tl{˥}}
\rentry{そとまご}{外孫}{ふかまー\tl{˥}}
\rentry{そとまわり}{外回り}{ふかまーり\tl{˥}}
\rentry{そね}{曽根}{すねー\tl{˥}}
\rentry{その}{}{うぬ\tl{˥˩}}
\rentry{そのていど}{その程度}{うだぎ、だぎ}
\rentry{そのとおりだ}{その通りだ}{えーるやろ}
\rentry{そのはず}{その筈}{えぬぱち}
\rentry{そのまま}{}{うぬ\ws{}まーま}
\rentry{そば}{側}{ぱた\tl{˥˩}}
\rentry{そふ}{祖父}{ぶや\tl{˥}}
\rentry{そぼ}{祖母}{ぱー\tl{˥}}
\rentry{そまつにする}{粗末にする}{すまっち\ws{}すん}
\rentry{そまる}{染まる}{すまるん\tl{˥˩}}
\rentry{そむく}{背く}{すむくん\tl{˥}}
\rentry{そめる}{染める}{しみん\tl{˥˩}}
\rentry{そら}{空}{じん\tl{˥˩}}
\rentry{そる}{剃る}{するん\tl{˥}}
\rentry{それ}{}{うり\tl{˥˩}}
\rentry{それだけ}{}{うだぎ、おび\tl{˥}}
\rentry{それなら}{}{\josho{}えー\kako{}\ws{}やちゃら}
\rentry{それほど}{}{うしゅく、うんしゅく\tl{˩˥}}
\rentry{それまで}{}{おび\tl{˥}}
\rentry{そろえる}{揃える}{するいるん\tl{˥}、するいん\tl{˥}}
\rentry{そん}{損}{すん\tl{˥}}
\rentry{そんけいする}{尊敬する}{うやまいるん\tl{˥˩}、うやめーるん\tl{˥˩}}
\rentry{ぞんざいにあつかう}{ぞんざいに扱う}{すまっち\ws{}すん}
\rentry{そんする}{損する}{そん\tl{˥}\ws{}すん\tl{˥˩}}
\rentry{そんな}{}{\josho{}えー\kako{}ぬ}
\rentry{そんなこと}{}{\josho{}えー\kako{}ぬ}
\rentry{そんなもの}{}{えぬ\ws{}むぬ}
\rentry{そんらく}{村落}{むら\tl{˥˩}}
\rentry{～だ}{}{やっさ、やるん\tl{˥˩}}
\rentry{だい}{大}{ぶー}
\rentry{たいおん}{体温}{にちぃ\tl{˩˥}}
\rentry{たいかい}{大海}{とー\tl{˥˩}}
\rentry{たいがい}{大概}{たげー\tl{˥}}
\rentry{たいかく}{体格}{ぐてー\tl{˩˥}}
\rentry{だいきん}{代金}{でん\tl{˥˩}}
\rentry{だいく}{大工}{せーぐ\tl{˥}、でーぐ\tl{˩˥}}
\rentry{たいこ}{太古}{かーま\tl{˥˩}むがし\tl{˩˥}}
\rentry{たいこ}{太鼓}{てーく\tl{˥}}
\rentry{たいこうする}{対抗する}{たい\tl{˥}\ws{}すん\tl{˥˩}}
\rentry{だいこん}{大根}{でーぐに\tl{˩˥}}
\rentry{だいじ}{大事}{うーぐとぅ、でーじ\tl{˩˥}、みざしくとぅ}
\rentry{だいじにしまう}{大事にしまう}{\josho{}あためー\kako{}すん、かちみるん\tl{˥}}
\rentry{たいせつに}{大切に}{あたらは\tl{˥˩}}
\rentry{たいそう}{}{\josho{}けっ\kako{}た、みざし\tl{˩˥}}
\rentry{だいたい}{太腿}{むんだらし\tl{˩˥}}
\rentry{たいだなのうふ}{怠惰な農夫}{ぴらすか\tl{˥}}
\rentry{だいち}{台地}{た\josho{}か\kako{}た}
\rentry{だいちじょうのへいち}{台地上の平地}{ぴぃせー\tl{˥˩}}
\rentry{たいちょうがわるい}{体調が悪い}{まーさかさー\ws{}ねーぬ}
\rentry{たいとう}{対等}{まーたき}
\rentry{だいどころ}{台所}{ふちめー\tl{˥}}
\rentry{だいのう}{大脳}{のー\tl{˥˩}}
\rentry{だいぶ}{}{たげー\tl{˥}}
\rentry{たいふう}{台風}{かちふき\tl{˥˩}、ぶーかち\tl{˥}}
\rentry{たいふうのふきかえし}{台風の吹き返し}{ふきけーし\tl{˥˩}}
\rentry{たいへん}{大変}{でーじ\tl{˩˥}}
\rentry{たいへんだ}{大変だ}{さってぃむ}
\rentry{たいへんな}{大変な}{みざし\tl{˩˥}}
\rentry{たいへんなこと}{大変なこと}{みざしくとぅ}
\rentry{たいまつ}{松明}{てー\tl{˥}}
\rentry{たいも}{田芋}{たーむじぃ\tl{˥}、むーじ\tl{˥}}
\rentry{たいもう}{体毛}{きー\tl{˥˩}}
\rentry{たいよう}{太陽}{しな\tl{˥}}
\rentry{たいようのかさ}{太陽の暈}{ぴながん\tl{˥}}
\rentry{たいらにならす}{平らにならす}{なだらぎるん\tl{˥˩}}
\rentry{だいり}{代理}{みょーでん\tl{˥˩}}
\rentry{たいりょく}{体力}{くんき\tl{˥˩}}
\rentry{たいんのじょし}{多淫の女子}{しきべー}
\rentry{たうえをする}{田植えをする}{たな\tl{˥}いびるん\tl{˥˩}}
\rentry{だえき}{唾液}{しん\tl{˥}}
\rentry{たえる}{絶える}{しっしるん\tl{˥}}
\rentry{たおす}{倒す}{とーすん\tl{˥}}
\rentry{たおる}{タオル}{しっし\tl{˥˩}}
\rentry{たおれる}{倒れる}{とーりるん\tl{˥}}
\rentry{だが}{}{えすが}
\rentry{たかい}{高い}{たか\tl{˥}はん}
\rentry{たかくとまる}{高く止まる}{たかぶるん\tl{˥}}
\rentry{たかだかと}{高々と}{た\josho{}か\kako{}たかし}
\rentry{たかなきする}{高鳴きする}{ふきん\tl{˥}}
\rentry{たかぶる}{高ぶる}{たかぶるん\tl{˥}}
\rentry{たがやす}{耕す}{けーすん\tl{˥}}
\rentry{たから}{宝}{たから\tl{˥}}
\rentry{だから}{}{\josho{}えす\kako{}がら}
\rentry{たからがい}{タカラガイ}{しぴ\tl{˥˩}}
\rentry{たきぎ}{薪}{たむぬ\tl{˥}、たむん}
\rentry{たきつける}{焚き付ける}{てしきるん\tl{˥˩}}
\rentry{たぎらせる}{}{たぎらすん\tl{˥˩}}
\rentry{たぎる}{}{たぎるん\tl{˥˩}}
\rentry{たく}{炊く}{たくん\tl{˥˩}、〔米飯を\ntilde{}〕ぴぃさすん\tl{˥}}
\rentry{だく}{抱く}{だぐん\tl{˥˩}}
\rentry{たぐさ}{田草}{とーさ\tl{˥}}
\rentry{たくさん}{}{ん\josho{}ごー\kako{}び}
\rentry{たくさん}{沢山}{まんどん\tl{˩˥}}
\rentry{～たくない}{}{ぼへぬ}
\rentry{たくらむ}{企む}{たくらむん、はかるん\tl{˥}}
\rentry{たぐる}{手繰る}{たぐるん\tl{˥}}
\rentry{たくわえる}{蓄える}{たみるん\tl{˥˩}}
\rentry{たくわえる}{貯える}{たこーすん\tl{˥˩}}
\rentry{たけ}{丈}{たき\tl{˥}、なぎ\tl{˩˥}}
\rentry{たけ}{岳}{だき\tl{˥˩}}
\rentry{たけ}{竹}{たき\tl{˥˩}}
\rentry{～だけ}{}{がーし}
\rentry{たけざる}{竹ざる}{ばーき\tl{˩˥}}
\rentry{たけとみ}{竹富}{たきどぅん\tl{˥}}
\rentry{たけぶえ}{竹笛}{ぴん\tl{˥˩}}
\rentry{たけゆか}{竹床}{たきふんた}
\rentry{たこ}{凧}{\josho{}たー\kako{}こ}
\rentry{たこ}{蛸}{たく\tl{˥}}
\rentry{だし}{出汁}{だーし\tl{˩˥}}
\rentry{たしかめる}{確かめる}{たしかみるん\tl{˥}、ただすん\tl{˥}}
\rentry{だしもの}{出しもの}{ぎなむぬ\tl{˩˥}}
\rentry{だす}{出す}{まるん\tl{˩˥}、んだすん\tl{˥}}
\rentry{たすかる}{助かる}{ぬちぃむやん}
\rentry{たすき}{}{あじまぎ\tl{˥}}
\rentry{たすける}{助ける}{たしきるん\tl{˥}}
\rentry{たずねる}{尋ねる}{たじぃにるん\tl{˥}、とーん\tl{˥˩}、むぬ\tl{˩˥}すくん\tl{˥˩}}
\rentry{～だそうだ}{}{えーちゅー、ちゅー}
\rentry{たそがれどき}{たそがれ時}{あやっふゎみ\tl{˥}}
\rentry{たたかい}{戦い}{いふつぁ\tl{˥}}
\rentry{たたかわせる}{戦わせる}{あーすん\tl{˥}}
\rentry{たたききる}{叩き切る}{〔固いものを\ntilde{}〕けーるん\tl{˥}、〔固い物、刃物で\ntilde{}〕だきすん\tl{˩˥}}
\rentry{たたきつける}{叩きつける}{だすきん\tl{˩˥}、なんぎるん\tl{˩˥}}
\rentry{たたきのめす}{叩きのめす}{どみんがすん\tl{˩˥}}
\rentry{たたきわる}{叩き割る}{たたぎばるん\tl{˥}}
\rentry{たたく}{叩く}{しきだるん\tl{˥}、たたぐん\tl{˥}、ふなぐん\tl{˥}}
\rentry{ただしい}{正しい}{あたりゃん、あたるん\tl{˥˩}}
\rentry{ただす}{質す}{ただすん\tl{˥}}
\rentry{ただちに}{直ちに}{ふたぎな\tl{˥˩}}
\rentry{ただのもの}{只のもの}{いたんだ\tl{˥}}
\rentry{たたみ}{畳}{たためー\tl{˥˩}}
\rentry{たたむ}{畳む}{たたむん}
\rentry{ただよわす}{漂わす}{おーぎるん\tl{˥˩}、なんぶりるん}
\rentry{たたる}{祟る}{たたーるん\tl{˥}}
\rentry{ただれる}{爛れる}{ただりん}
\rentry{～たち}{～達}{んだ}
\rentry{たちなおる}{立ち直る}{たちのーるん\tl{˥}、むちのーるん\tl{˩˥}}
\rentry{たつ}{建つ}{たつん\tl{˥}}
\rentry{たつ}{発つ}{ぱるん\tl{˥}}
\rentry{たつ}{立つ}{だっち、たつん\tl{˥}}
\rentry{たつ}{経つ}{たつん\tl{˥}}
\rentry{たつ}{辰}{たち\tl{˥˩}}
\rentry{だっこくする}{脱穀する}{しゅるん\tl{˥}}
\rentry{だっした}{脱した}{ぬがりゃん\tl{˩˥}}
\rentry{たつまき}{竜巻}{いのー\tl{˥˩}}
\rentry{たてあみのいっしゅ}{建て網の一種}{ふつぁあん}
\rentry{たてる}{建てる}{たちるん\tl{˥}}
\rentry{たてる}{立てる}{たちるん\tl{˥}}
\rentry{たとえる}{例える}{たとぅいん\tl{˥}}
\rentry{たどる}{辿る}{たどぅるん\tl{˥}}
\rentry{たな}{棚}{たな\tl{˥˩}}
\rentry{たにん}{他人}{ぴぃとぅ\tl{˥˩}}
\rentry{たね}{種}{たに\tl{˥}}
\rentry{たねがなくなる}{種がなくなる}{たに\ws{}しさん}
\rentry{たねとりさい}{種取祭}{たにどぅり\tl{˥˩}}
\rentry{たねとりさいのやまもりめし}{種取祭の山盛飯}{いばち\tl{˥˩}}
\rentry{たのあぜ}{田の畔}{あぶし\tl{˥}}
\rentry{たのしむ}{楽しむ}{あまいるん\tl{˥˩}、うむっさーすん、さにしゃー\ws{}すん}
\rentry{たのみ}{頼み}{たぬみ\tl{˥}}
\rentry{たのむ}{頼む}{たぬむん\tl{˥}}
\rentry{たのもし}{頼母子}{むえー\tl{˥˩}}
\rentry{たのもしくおもう}{頼もしく思う}{すむ\tl{˥}ずさ\tl{˥}はーん}
\rentry{たば}{束}{たぱ\tl{˥}}
\rentry{たばこ}{煙草}{たかぶ\tl{˥}}
\rentry{たばこいれ}{煙草入れ}{ぷつぉー\tl{˥}}
\rentry{たばねる}{束ねる}{さまるん\tl{˥}、まりぐん\tl{˥˩}}
\rentry{たび}{旅}{たぴ\tl{˥˩}}
\rentry{たびたび}{度々}{す\josho{}たー\kako{}すたー}
\rentry{たべごろ}{食べ頃}{へーじぶん\tl{˥˩}}
\rentry{たべつくす}{食べ尽くす}{へぶたん}
\rentry{たべもの}{食べ物}{ほーむぬ\tl{˥˩}、むぬ\tl{˩˥}}
\rentry{たべものをいれるたけかご}{食べ物を入れる竹籠}{ふたじろー\tl{˥˩}}
\rentry{たべる}{食べる}{ほーん\tl{˥˩}}
\rentry{たぼうだ}{多忙だ}{ぱんたっさ\tl{˥}はん}
\rentry{たま}{弾}{たまん\tl{˥}}
\rentry{たま}{玉}{たまん\tl{˥}}
\rentry{たまご}{卵}{けー\tl{˥}、〔蟹の\ntilde{}〕ぱりゃん\tl{˥}}
\rentry{たましい}{魂}{たましぃ\tl{˥}、たまち\tl{˥}、まーぶり\tl{˩˥}}
\rentry{だます}{騙す}{あぎまーすん\tl{˥˩}}
\rentry{たまたま}{}{た\josho{}まー\kako{}たま}
\rentry{だまっていること}{黙っていること}{ゆだばり\tl{˥˩}}
\rentry{たまに}{}{た\josho{}まー\kako{}たま}
\rentry{たまる}{溜まる}{たまるん\tl{˥˩}}
\rentry{たまる}{貯まる}{たまるん\tl{˥˩}}
\rentry{だまる}{黙る}{とぅなばるん\tl{˥˩}}
\rentry{たまわる}{賜る}{たぼーらりん\tl{˥}、たぼらいるん\tl{˥}}
\rentry{たむし}{田虫}{うるばたぎ\tl{˥}}
\rentry{ため}{為}{たみ\tl{˥}}
\rentry{ためいけ}{溜め池}{あなぶ\tl{˥}}
\rentry{ためす}{試す}{たみすん\tl{˥}}
\rentry{ためる}{溜める}{たみるん\tl{˥˩}}
\rentry{ためる}{貯める}{たみるん\tl{˥˩}}
\rentry{たもあみ}{たも網}{たぶ\tl{˥}}
\rentry{たもつ}{保つ}{たむつん\tl{˥}}
\rentry{たやす}{絶やす}{しっさん\tl{˥}}
\rentry{たやすい}{た易い}{やっさ\tl{˩˥}はん}
\rentry{たより}{頼り}{たぬみ\tl{˥}}
\rentry{たよりにする}{頼りにする}{たなぎるん}
\rentry{たよる}{頼る}{たなぐん、たゆるん\tl{˥}、むたーりるん\tl{˥˩}}
\rentry{たらい}{盥}{たれー\tl{˥}}
\rentry{だらけ}{}{ぶった}
\rentry{だらしないすがたをする}{だらしない姿をする}{〔帯が下に下がって\ntilde{}〕ふきつぁーるん\tl{˥˩}}
\rentry{たらす}{垂らす}{すさーらすん\tl{˥}}
\rentry{たらす}{足らす}{たらすん\tl{˥˩}}
\rentry{だらだら}{}{だらだら}
\rentry{たりない}{足りない}{たらーぬ\tl{˥˩}}
\rentry{たりる}{足りる}{たりるん\tl{˥˩}}
\rentry{だるい}{}{だる\tl{˩˥}さーん}
\rentry{たるき}{垂木}{きちぃ\tl{˥}}
\rentry{だれ}{誰}{たー\tl{˥˩}}
\rentry{だれが}{誰が}{たるざー\tl{˥}}
\rentry{たれさがる}{垂れ下がる}{すさーるん}
\rentry{たれる}{垂れる}{すさーるん、〔雫が\ntilde{}〕すたるん\tl{˥}、たりるん\tl{˥}}
\rentry{だれる}{}{だーりるん\tl{˩˥}、だーりん\tl{˩˥}}
\rentry{～だろう}{}{やっさ}
\rentry{たわむ}{}{まんかるん\tl{˥˩}}
\rentry{たわら}{俵}{たーら\tl{˥}、たらぐ\tl{˥}}
\rentry{たん}{反}{たん\tl{˥˩}}
\rentry{だん}{壇}{だん\tl{˩˥}}
\rentry{だん}{段}{だん\tl{˩˥}}
\rentry{たんきもの}{短気者}{たんきむぬ}
\rentry{たんさくする}{探索する}{さぐるん\tl{˥˩}、さばくん\tl{˥}}
\rentry{たんじょういわい}{誕生祝い}{まりぶなー\tl{˥˩}}
\rentry{たんじょうする}{誕生する}{まりるん\tl{˥˩}}
\rentry{たんじょうび}{誕生日}{まりぴん\tl{˥˩}}
\rentry{たんすい}{淡水}{あまみじ\tl{˥˩}、みじぃ\tl{˥˩}}
\rentry{だんせい}{男性}{びどぅむ\tl{˥˩}}
\rentry{だんちく}{ダンチク}{だちご\tl{˥˩}}
\rentry{だんどく}{ダンドク}{まるぶさ\tl{˥˩}}
\rentry{だんぱつ}{断髪}{だんぱち\tl{˥}}
\rentry{だんぱん}{談判}{だんぱん\tl{˥˩}}
\rentry{だんぱんする}{談判する}{かきあうん\tl{˥}}
\rentry{たんぼ}{田んぼ}{たなー\tl{˥}、たなが\tl{˥}}
\rentry{たんめい}{短命}{ぬちぃまろ\tl{˩˥}はん}
\rentry{ち}{地}{じぃー\tl{˩˥}}
\rentry{ち}{血}{じぃー\tl{˥˩}}
\rentry{ちいさい}{小さい}{いしゃが\tl{˥}はん、ぐま\tl{˩˥}はん}
\rentry{ちいさくなる}{小さくなる}{しぃぴりるん\tl{˥˩}、すくまるん\tl{˥}}
\rentry{ちえ}{知恵}{じんぶん\tl{˩˥}}
\rentry{ちかい}{近い}{すか\tl{˥}はん}
\rentry{ちがう}{違う}{あらぬ\tl{˥}、ちごーん\tl{˥˩}}
\rentry{ちかづく}{近づく}{すかは\tl{˥}\ws{}なるん\tl{˩˥}、ぬすかるん\tl{˥˩}、めー\tl{˩˥}ぬすかるん\tl{˥˩}}
\rentry{ちかみち}{近道}{すかみち\tl{˥}}
\rentry{ちがや}{チガヤ}{がや\tl{˩˥}}
\rentry{ちぎる}{千切る}{すむん\tl{˥}}
\rentry{ちぎれる}{千切れる}{むっしるん\tl{˥˩}}
\rentry{ちくせいのたべものかご}{竹製の食べ物籠}{びらぐ\tl{˩˥}}
\rentry{ちくどぅん}{筑登之}{ちくどぅん\tl{˥}}
\rentry{ちすじ}{血筋}{ぴき\tl{˥˩}}
\rentry{ちち}{乳}{じぃ\tl{˥}}
\rentry{ちち}{父}{いや\tl{˥}}
\rentry{ちちおや}{父親}{いや\tl{˥}}
\rentry{ちちかたのしんせき}{父方の親戚}{たにかたうやぐ\tl{˥}}
\rentry{ちぢこまる}{縮こまる}{すくまるん\tl{˥}}
\rentry{ちぶさ}{乳房}{じぃ\tl{˥}、じっち}
\rentry{ちほう}{痴呆}{ういぷら\tl{˥}}
\rentry{ちゃ}{茶}{さ\tl{˥}}
\rentry{ちゃがし}{茶菓子}{さうき\tl{˥}}
\rentry{ちゃづつ}{茶筒}{じんぎり\tl{˩˥}}
\rentry{ちゃとう}{茶湯}{さとー\tl{˥}}
\rentry{ちゃわん}{茶碗}{さばん\tl{˥}}
\rentry{ちゅういする}{注意する}{きー\tl{˥˩}\ws{}しきるん\tl{˥}}
\rentry{ちゅうげん}{忠言}{すーむに\tl{˥}、ゆしぐとぅ\tl{˥˩}}
\rentry{ちゅうごく}{中国}{とー\tl{˥}}
\rentry{ちゅうざい}{駐在}{ゆんさん\tl{˩˥}}
\rentry{ちゅうさいする}{仲裁する}{あがすん\tl{˥˩}}
\rentry{ちゅうざら}{中皿}{なかざら\tl{˩˥}}
\rentry{ちゅうしょく}{昼食}{ぴぃすまりむぬ\tl{˩˥}}
\rentry{ちゅうどくさせる}{中毒させる}{びらーすん\tl{˩˥}}
\rentry{ちゅうどくする}{中毒する}{びーるん\tl{˩˥}}
\rentry{ちょう}{蝶}{ぱぴる\tl{˥}}
\rentry{ちょうけい}{長兄}{ぶしゃ\tl{˥}}
\rentry{ちょうげん}{調弦}{ちんだみ}
\rentry{ちょうこう}{兆候}{しるし\tl{˥˩}}
\rentry{ちょうし}{長姉}{ぼーま\tl{˥}}
\rentry{ちょうじゅ}{長寿}{ちょーみー\tl{˥˩}、ながいき\tl{˥˩}}
\rentry{ちょうじょう}{頂上}{ちじ\tl{˥}}
\rentry{ちょうしょく}{朝食}{あさむぬ\tl{˥}}
\rentry{ちょうだいする}{頂戴する}{たぼーらりん\tl{˥}}
\rentry{ちょうていする}{調停する}{さばくん\tl{˥}}
\rentry{ちょうなん}{長男}{さこーし\tl{˥}}
\rentry{ちょうめい}{長命}{ちょーみー\tl{˥˩}}
\rentry{ちょうりゅう}{潮流}{すーぴき\tl{˥˩}}
\rentry{ちょきんする}{貯金する}{じん\tl{˩˥}たみるん\tl{˥˩}}
\rentry{ちょぞうする}{貯蔵する}{たこーすん\tl{˥˩}}
\rentry{ちょま}{苧麻}{ぶー\tl{˩˥}}
\rentry{ちらかす}{散らかす}{しきぽーるん\tl{˥˩}}
\rentry{ちらかる}{散らかる}{ぽーりゃん\tl{˥˩}}
\rentry{ちらす}{散らす}{ぽーるん\tl{˥˩}}
\rentry{ちらばる}{散らばる}{ぽーりゃん\tl{˥˩}}
\rentry{ちり}{塵}{あるふた\tl{˥˩}}
\rentry{ちりょうする}{治療する}{いしゃんが\tl{˥}\ws{}はかるん\tl{˥}}
\rentry{ちんでんする}{沈澱する}{ぴりがるん\tl{˥}}
\rentry{ついていく}{付いて行く}{しき\tl{˥}\ws{}んぐん\tl{˥˩}}
\rentry{ついばむ}{}{しきほーん\tl{˥˩}}
\rentry{ついひ}{追肥}{またぐい\tl{˥˩}}
\rentry{ついやす}{費やす}{すこーん\tl{˥˩}}
\rentry{つうじがよくなる}{通じが良くなる}{さぎるん\tl{˥}}
\rentry{つうようする}{通用する}{とーるん\tl{˥}}
\rentry{つえ}{杖}{ぐさん\tl{˩˥}}
\rentry{つか}{柄}{すか\tl{˥}}
\rentry{つかいでがある}{使い出がある}{すけーとー\tl{˥˩}\ws{}あん\tl{˥}}
\rentry{つかう}{使う}{すこーん\tl{˥˩}}
\rentry{つかさ}{司}{すかさ\tl{˥}}
\rentry{つかまえる}{捕える}{かみしきるん\tl{˥˩}、かみるん\tl{˥}、かみん\tl{˥}}
\rentry{つかみだす}{掴み出す}{すかみんだすん\tl{˥}}
\rentry{つかむ}{掴む}{かみるん\tl{˥}、かみん\tl{˥}、すかむん\tl{˥}、〔箸などを\ntilde{}〕ぱつぁむん\tl{˥}}
\rentry{つかる}{漬かる}{すかるん\tl{˥˩}}
\rentry{つかれ}{疲れ}{ぼーがり\tl{˩˥}}
\rentry{つかれはてる}{疲れ果てる}{がんどーりん\tl{˩˥}}
\rentry{つかれる}{疲れる}{くたんでぃるん\tl{˥}、だーりん\tl{˩˥}、ぶがりるん\tl{˩˥}、ぼたりるん\tl{˩˥}、ぼたりん\tl{˩˥}}
\rentry{つかれをとる}{疲れをとる}{ぼーり\tl{˩˥}のがすん\tl{˩˥}}
\rentry{つかわす}{遣わす}{やらすん\tl{˥˩}}
\rentry{つき}{月}{〔天体の\ntilde{}〕しけん\tl{˥}、〔月日の\ntilde{}〕ちき\tl{˥}}
\rentry{つきあい}{付き合い}{ぴぃとぅぴれー\tl{˥}、ぴらい\tl{˥}、ぴらす}
\rentry{つきあげる}{突き上げる}{しきあんぎるん\tl{˥˩}}
\rentry{つきあたる}{突き当たる}{しきあたるん\tl{˥˩}}
\rentry{つきくずす}{突き崩す}{どみがすん\tl{˩˥}}
\rentry{つきさす}{突き刺す}{ぬぐん\tl{˩˥}}
\rentry{つきださせる}{突き出させる}{しきんだすん}
\rentry{つきだす}{突き出す}{しきんだすん}
\rentry{つきでる}{突き出る}{しき\tl{˥˩}んじるん\tl{˥}、とぅんちるん\tl{˥˩}}
\rentry{つきとおす}{突き通す}{とふかすん\tl{˥˩}}
\rentry{つきとばす}{突き飛ばす}{うしとーすん、しきとぅぱすん\tl{˥˩}}
\rentry{つきよ}{月夜}{しけんゆー}
\rentry{つく}{搗く}{すくん\tl{˥}、〔臼・杵で\ntilde{}〕すさぎるん\tl{˥}}
\rentry{つく}{着く}{すくん\tl{˥}}
\rentry{つく}{突く}{しくん\tl{˥˩}、すくん\tl{˥˩}}
\rentry{つぐ}{注ぐ}{ぱらすん\tl{˥}}
\rentry{つぐ}{継ぐ}{つぐん\tl{˥˩}}
\rentry{つくす}{尽くす}{ぎばるん\tl{˥˩}}
\rentry{つぐなう}{償う}{ぱくん\tl{˥˩}}
\rentry{つくり}{造り}{すくり\tl{˥}}
\rentry{つくる}{作る}{しこーるん\tl{˥˩}、すくるん\tl{˥}}
\rentry{つける}{漬ける}{しきるん\tl{˥˩}}
\rentry{つける}{点ける}{〔灯りや火を\ntilde{}〕しきるん\tl{˥˩}、〔薪に火を\ntilde{}〕てしきるん\tl{˥˩}}
\rentry{つげる}{告げる}{いーしきるん\tl{˥˩}}
\rentry{つたない}{拙い}{すぅな\tl{˥}はん}
\rentry{つち}{土}{んた\tl{˩˥}}
\rentry{つづかせる}{続かせる}{すながらすん\tl{˥˩}}
\rentry{つづく}{続く}{ちじくん\tl{˥˩}}
\rentry{つづけてみる}{続けて見る}{みーとーすん\tl{˩˥}}
\rentry{つづける}{続ける}{ちじきるん\tl{˥˩}}
\rentry{つつしむ}{慎む}{ちちすむん\tl{˥˩}}
\rentry{つったつ}{突っ立つ}{だっち}
\rentry{つつむ}{包む}{すむん\tl{˥}}
\rentry{つとめる}{勤める}{すとぅみるん\tl{˥}}
\rentry{つな}{綱}{すな\tl{˥}}
\rentry{つなぎかえる}{繋ぎ変える}{〔牛馬を\ntilde{}〕むすなすん\tl{˩˥}}
\rentry{つなぎとめる}{繋ぎとめる}{ばつぁみるん\tl{˩˥}}
\rentry{つなぐ}{繋ぐ}{〔牛馬などの動物を\ntilde{}〕ばつぁみるん\tl{˩˥}}
\rentry{つなみ}{津波}{ぶーなん\tl{˥}}
\rentry{つねに}{常に}{い\josho{}ちぃ\kako{}ん、しゃー\tl{˥}}
\rentry{つねる}{}{すむん\tl{˥}}
\rentry{つねる}{抓る}{じんぱるん}
\rentry{つの}{角}{しのー\tl{˥}}
\rentry{つば}{唾}{しん\tl{˥}}
\rentry{つばさ}{翼}{ぱに\tl{˥˩}}
\rentry{つばめ}{燕}{またさ\tl{˩˥}}
\rentry{つぶて}{礫}{いしふく\tl{˥˩}}
\rentry{つぶれる}{潰れる}{ぴさんたらすん\tl{˥˩}}
\rentry{つぼ}{坪}{ちぶ\tl{˥˩}}
\rentry{つぼ}{壺}{すぅぷ\tl{˥˩}}
\rentry{つま}{妻}{とぅん\tl{˥˩}}
\rentry{つまさきだち}{つま先立ち}{とぅんたち\tl{˥˩}}
\rentry{つまだつ}{爪立つ}{とぅんたちびるん\tl{˥˩}}
\rentry{つまにする}{妻にする}{そーるん\tl{˥˩}}
\rentry{つまみおる}{摘み折る}{ぶるん\tl{˩˥}}
\rentry{つまる}{詰まる}{すぅまるん\tl{˥}}
\rentry{つまをめとる}{娶る}{〔妻を\ntilde{}〕とぅみん\tl{˥}}
\rentry{つみ}{罪}{とぅが\tl{˥}、ばつぁ\tl{˩˥}}
\rentry{つみかさねる}{積み重ねる}{すむん\tl{˥˩}}
\rentry{つみきる}{摘み切る}{すむん\tl{˥}}
\rentry{つみとる}{摘み取る}{むっすん\tl{˥˩}}
\rentry{つむ}{摘む}{すむん\tl{˥}、〔木の実などを\ntilde{}〕ぶるん\tl{˥˩}}
\rentry{つむ}{積む}{すむん\tl{˥˩}、ぬしん\tl{˥˩}}
\rentry{つむじかぜ}{つむじ風}{まーじゃかち\tl{˥˩}}
\rentry{つめ}{爪}{しみ\tl{˥˩}}
\rentry{つめこむ}{詰め込む}{うしくむん\tl{˥˩}}
\rentry{つめたい}{冷たい}{ぴじゅる\tl{˥}さーん、ぴり\tl{˥}}
\rentry{つめたいかぜ}{冷たい風}{ぴりかち\tl{˥}}
\rentry{つもり}{}{しゃーみ\tl{˥}}
\rentry{つゆ}{梅雨}{ゆどぅあみ\tl{˩˥}}
\rentry{つゆ}{露}{ゆす\tl{˩˥}}
\rentry{つよい}{強い}{すさ\tl{˥}はん}
\rentry{つよくたたく}{強く叩く}{まるばすん\tl{˥˩}}
\rentry{つよくたばねる}{強く束ねる}{ふんまるぐん\tl{˥˩}}
\rentry{つよくなげいれる}{強く投げ入れる}{〔液体の中に物を\ntilde{}〕だぼんがすん\tl{˩˥}}
\rentry{つよくふる}{強く振る}{すくふぁすん\tl{˥˩}}
\rentry{つよめる}{強める}{しぃみるん\tl{˥}}
\rentry{つらぬく}{貫く}{ぬぐん\tl{˩˥}}
\rentry{つらねる}{連ねる}{すながらすん\tl{˥˩}、すながるん\tl{˥˩}}
\rentry{つりあう}{つり合う}{あたるん\tl{˥˩}}
\rentry{つりいと}{釣り糸}{なん\tl{˩˥}}
\rentry{つりえさ}{釣り餌}{むんだに\tl{˩˥}}
\rentry{つりざお}{釣り竿}{じぼー\tl{˥˩}}
\rentry{つりざおのくふう}{釣竿の工夫}{あんざらだーぐ\tl{˥}}
\rentry{つりさげる}{吊下げる}{さんぎるん\tl{˥}}
\rentry{つりばり}{釣針}{じばり\tl{˥˩}}
\rentry{つる}{釣る}{ほすん\tl{˥}}
\rentry{つるす}{吊るす}{さんぎるん\tl{˥}}
\rentry{つるべ}{釣瓶}{ぶら\tl{˥˩}}
\rentry{つれ}{連れ}{ぐー\tl{˩˥}}
\rentry{つれていく}{連れていく}{そり\tl{˥˩}\ws{}んぐん\tl{˥˩}、そんぐん\tl{˥˩}}
\rentry{つれる}{連れる}{そーるん\tl{˥˩}}
\rentry{つんぼ}{聾}{みんかー\tl{˥}}
\rentry{て}{手}{しー\tl{˥}}
\rentry{てあし}{手足}{しーぱん\tl{˥}}
\rentry{てあらい}{手荒い}{しー\tl{˥}あらはーん}
\rentry{～である}{}{やるん\tl{˥˩}}
\rentry{ていたいする}{停滞する}{ゆどぅむん\tl{˩˥}}
\rentry{ていねいに}{丁寧に}{みしこーみしこ}
\rentry{ていへいち}{低平地}{とー\tl{˥˩}}
\rentry{てがあらっぽい}{手が粗っぽい}{しー\tl{˥}あらはーん}
\rentry{てがける}{手がける}{ししきるん}
\rentry{でかした}{}{したい}
\rentry{てがすぐでる}{手がすぐ出る}{しー\tl{˥}ぺしゃ\tl{˥}はん}
\rentry{てがすばしこい}{手がすばしこい}{しー\tl{˥}ぺしゃ\tl{˥}はん}
\rentry{てがみ}{手紙}{しがみ\tl{˥}}
\rentry{てがら}{手柄}{てぃがら\tl{˥}}
\rentry{てきたいする}{敵対する}{たい\tl{˥}\ws{}すん\tl{˥˩}}
\rentry{てきぱきと}{}{さっと}
\rentry{できる}{出来る}{なーるん\tl{˩˥}}
\rentry{でぐち}{出口}{んじふち\tl{˥}}
\rentry{てくび}{手首}{\josho{}しぬ\kako{}\ws{}ふき}
\rentry{てじな}{手品}{ほっか\tl{˥}}
\rentry{てだすけ}{手助け}{しがねー\tl{˥}}
\rentry{てつだい}{手伝い}{かしー\tl{˥}、しがねー\tl{˥}}
\rentry{てっぽう}{鉄砲}{しゅぷ\tl{˥˩}}
\rentry{てぬぐい}{手拭い}{しっし\tl{˥˩}}
\rentry{てのしごとがおそい}{手の仕事が遅い}{しーにふつはん}
\rentry{てのゆび}{手の指}{しんび\tl{˥}}
\rentry{てばやい}{手早い}{しー\tl{˥}ぺしゃ\tl{˥}はん}
\rentry{てばやく}{手早く}{さっと}
\rentry{でぶ}{}{ばたぶたー\tl{˩˥}}
\rentry{でぶね}{出船}{んじふに\tl{˥}}
\rentry{てま}{手間}{しぃま\tl{˥}}
\rentry{てまちん}{手間賃}{しぃま\tl{˥}}
\rentry{てまどる}{手間取る}{しまはかるん\tl{˥}}
\rentry{てらす}{照らす}{あがらすん\tl{˥˩}}
\rentry{てりはぼく}{テリハボク}{やらぶ\tl{˥}}
\rentry{てる}{照る}{〔太陽が\ntilde{}〕しるん\tl{˥}}
\rentry{でる}{出る}{んじるん\tl{˥}、んじん\tl{˥}}
\rentry{てをあわせる}{手を合わせる}{しー\tl{˥}うさぎるん\tl{˥˩}}
\rentry{てをつける}{手をつける}{ししきるん}
\rentry{てん}{天}{じん\tl{˥˩}}
\rentry{てんき}{天気}{おしき\tl{˥˩}、わしき\tl{˥˩}}
\rentry{てんきがわるい}{天気が悪い}{やなわしき\tl{˥˩}}
\rentry{てんこう}{天候}{おしき\tl{˥˩}、わしき\tl{˥˩}}
\rentry{でんごんする}{伝言する}{くとぅしきるん\tl{˥}}
\rentry{てんじょう}{天井}{しんじょー\tl{˥}}
\rentry{でんせんびょう}{伝染病}{はやりやん\tl{˥}}
\rentry{てんてこまい}{てんてこ舞い}{ぱったらげー}
\rentry{てんとうする}{転倒する}{まるぶちん\tl{˥˩}}
\rentry{てんぷら}{天ぷら}{しんぷら\tl{˥}、ぱんべー\tl{˥}}
\rentry{てんません}{伝馬船}{てぃんま\tl{˥˩}}
\rentry{と}{戸}{やどぅ\tl{˩˥}}
\rentry{とあみ}{投網}{うちあん\tl{˥}}
\rentry{といきく}{問い聞く}{とぅいすくん\tl{˥˩}}
\rentry{といし}{砥石}{とぅし\tl{˥}}
\rentry{といせめる}{問い責める}{とぅしみるん\tl{˥}}
\rentry{どいつが}{}{たるざー\tl{˥}}
\rentry{といつめる}{問い詰める}{とぅいしみるん\tl{˥˩}、とぅしみるん\tl{˥}}
\rentry{とう}{唐}{とー\tl{˥}}
\rentry{とう}{問う}{たじぃにるん\tl{˥}、とーん\tl{˥˩}}
\rentry{とう}{籐}{くち\tl{˥}}
\rentry{どう}{}{ねー\tl{˩˥}、\josho{}ねー\kako{}や}
\rentry{どうか}{}{たんでー\tl{˥}、\josho{}ちゃー\kako{}が}
\rentry{どうかね}{}{\josho{}ねー\kako{}や}
\rentry{とうがらし}{トウガラシ}{ぐす\tl{˥˩}}
\rentry{とうがん}{冬瓜}{すぷりん\tl{˥}}
\rentry{とうき}{冬季}{ふっひ\tl{˥˩}}
\rentry{どうきせい}{同期生}{まーじぃ\tl{˥˩}}
\rentry{とうきび}{唐黍}{や\josho{}た\kako{}ぷ}
\rentry{どうぐ}{道具}{だーぐ\tl{˩˥}}
\rentry{どうくつ}{洞窟}{いん\tl{˥˩}}
\rentry{どうけつ}{洞穴}{いん\tl{˥˩}、〔海中の\ntilde{}〕がまん\tl{˥˩}}
\rentry{とうじ}{冬至}{とぅんじー\tl{˥}}
\rentry{とうじき}{陶磁器}{やぎむぬ\tl{˩˥}}
\rentry{どうして}{}{ねーき\tl{˩˥}、ねーきる}
\rentry{どうしても}{}{ぬーしん、\josho{}のー\kako{}しん}
\rentry{どうぞ}{}{\josho{}どー\kako{}でぃん}
\rentry{どうたい}{胴体}{どぅー\tl{˩˥}}
\rentry{どうだい}{}{\josho{}ちゃー\kako{}が}
\rentry{とうちする}{統治する}{うさみるん\tl{˥}}
\rentry{とうちゃくする}{到着する}{すくん\tl{˥}}
\rentry{どうとう}{同等}{まーたき}
\rentry{どうどうとあるくさま}{堂々と歩くさま}{だっふぁだっふぁ}
\rentry{どうねん}{同年}{ゆぬとぅし\tl{˩˥}}
\rentry{とうはつ}{頭髪}{あまじ\tl{˥}}
\rentry{とうひょう}{投票}{ふだいり\tl{˥˩}}
\rentry{とうぶ}{頭部}{あまっすくる\tl{˥}}
\rentry{どうぶつ}{動物}{いすむし\tl{˥}}
\rentry{とうほう}{東方}{あーり\tl{˥˩}、あらた、ありかた\tl{˥˩}}
\rentry{とうぼうする}{逃亡する}{ぴんぎるん\tl{˥}}
\rentry{どうめい}{童名}{やらびなー\tl{˩˥}}
\rentry{どうり}{道理}{どーり\tl{˩˥}}
\rentry{とうりょう}{頭領}{かしら\tl{˥}}
\rentry{どうろ}{道路}{みちぃ\tl{˥˩}}
\rentry{どうろしゅうり}{道路修理}{みちぃくせー\tl{˥˩}}
\rentry{とお}{十}{とー\tl{˥˩}}
\rentry{とおい}{遠い}{とぅさ\tl{˥˩}はん}
\rentry{とおいところ}{遠い所}{とぅけー\tl{˥˩}}
\rentry{とおす}{通す}{とーすん\tl{˥}}
\rentry{とおのく}{遠のく}{ぬぎるん\tl{˩˥}}
\rentry{とおりにくい}{通りにくい}{〔棘木やつるが多くて\ntilde{}〕あざらはーん}
\rentry{とおる}{通る}{とーるん\tl{˥}}
\rentry{とおれない}{通れない}{〔荒れて\ntilde{}〕あざらはーん}
\rentry{とが}{咎}{ばつぁ\tl{˩˥}}
\rentry{とかげ}{トカゲ}{ばずら\tl{˩˥}}
\rentry{とかす}{溶かす}{たらすん\tl{˥˩}}
\rentry{とがめる}{咎める}{とぅがみるん\tl{˥}}
\rentry{とがらせる}{尖らせる}{ぴんとまらすん\tl{˥}}
\rentry{とがる}{尖る}{ぴんとまるん\tl{˥}}
\rentry{とがをかぶる}{咎を被る}{ばつぁかぷん}
\rentry{ときがたつ}{時が経つ}{んぐん\tl{˥˩}}
\rentry{とききかせる}{説き聞かせる}{えにすかすん\tl{˥˩}}
\rentry{どきどき}{}{どんどん、ふとぅふとぅ}
\rentry{とぎれる}{途切れる}{しっしるん\tl{˥}}
\rentry{とく}{得}{とぅく\tl{˥}}
\rentry{とく}{徳}{とぅく\tl{˥}}
\rentry{とぐ}{研ぐ}{あーらすん、とぅぐん\tl{˥}}
\rentry{どく}{毒}{どぅく\tl{˩˥}}
\rentry{どくしん}{独身}{たんがーむぬ\tl{˥}}
\rentry{とくそくする}{督促する}{いみるん\tl{˥}}
\rentry{とげ}{棘}{じぃー\tl{˩˥}}
\rentry{とける}{溶ける}{たりるん\tl{˥}、とぅきるん\tl{˥}}
\rentry{とける}{融ける}{とぅきるん\tl{˥}}
\rentry{どこ}{何処}{ざー\tl{˩˥}}
\rentry{どこの}{}{ざぬ}
\rentry{どこの}{何処の}{ざーぬ\tl{˩˥}}
\rentry{とこのま}{床の間}{ざとく\tl{˩˥}、とぅく\tl{˥˩}}
\rentry{とこのまのかみだな}{床の間の神棚}{ぶざすけー\tl{˥}}
\rentry{とこや}{床屋}{だんぱちーやー}
\rentry{ところ}{所}{とぅくる\tl{˥}}
\rentry{とさつする}{屠殺する}{ばずらすん\tl{˥˩}}
\rentry{とし}{年}{とぅちぃ\tl{˥}}
\rentry{とし}{歳}{とぅちぃ\tl{˥}}
\rentry{としうえ}{年上}{しじゃー\tl{˥}、しゃま\tl{˥}}
\rentry{とししたのきょうだい}{年下の兄弟}{うとぅーとぅ\tl{˥}}
\rentry{としより}{年寄り}{うしとぅ\tl{˥}}
\rentry{としよりがたばこをすうどうぐ}{年寄りが煙草を吸う道具}{たかぼん}
\rentry{とじる}{閉じる}{ふーん\tl{˥˩}、ふさぐん\tl{˥˩}、〔目を\ntilde{}〕ふつぁるん\tl{˥˩}}
\rentry{としをとる}{年を取る}{ういるん\tl{˥}、とぅちぃとぅるん\tl{˥}}
\rentry{としをとる}{歳を取る}{うしとぅ\tl{˥}\ws{}なるん\tl{˩˥}}
\rentry{どすん}{}{だーんた}
\rentry{どだい}{土台}{びしじ\tl{˥˩}}
\rentry{とち}{土地}{じぃー\tl{˩˥}}
\rentry{どっかと}{どかっと}{どっふぁった}
\rentry{どっかり}{}{どっふぁった}
\rentry{とっくに}{}{きさ\tl{˥}}
\rentry{とつぜん}{突然}{あた\tl{˥˩}}
\rentry{とってすてる}{取って捨てる}{とぅり\tl{˥}\ws{}っしるん\tl{˥˩}、\josho{}とぅり\kako{}\ws{}っしん\tl{˥˩}}
\rentry{とつべんだ}{訥弁だ}{ふちぃくぱはん}
\rentry{とてもいたい}{とても痛い}{あがよー}
\rentry{とどく}{届く}{とぅとぅぐん\tl{˥}}
\rentry{とどける}{届ける}{とぅとぅぎるん\tl{˥}}
\rentry{ととのう}{整う}{とぅとぅのーん\tl{˥}}
\rentry{とどろかせる}{轟かせる}{とぅゆますん\tl{˥}}
\rentry{どなた}{}{たー\tl{˥˩}}
\rentry{どなる}{怒鳴る}{たけーるん\tl{˥}、どぅげーるん\tl{˩˥}}
\rentry{どの}{}{ざーぬ\tl{˩˥}、ざぬ}
\rentry{とのしろ}{登野城}{とぅぬすく\tl{˥˩}}
\rentry{どのような}{}{ぬ\josho{}ー\kako{}しゃる}
\rentry{どのように}{}{ねーる\tl{˩˥}}
\rentry{とび}{鳶}{ぺふつぁ\tl{˥}}
\rentry{とびうお}{飛魚}{とぅぴゆ\tl{˥˩}}
\rentry{とびおりる}{跳び下りる}{ぶんちうちん\tl{˩˥}}
\rentry{とびこえる}{飛び越える}{とぅぴくいるん\tl{˥˩}}
\rentry{とびこむ}{飛び込む}{とぅぷちん\tl{˥˩}}
\rentry{とびでる}{飛び出る}{とぅんじるん\tl{˥˩}}
\rentry{とぶ}{跳ぶ}{ふんつん\tl{˥˩}}
\rentry{とぶ}{飛ぶ}{とぅぷん\tl{˥˩}}
\rentry{とぼしい}{貧しい}{ぴんそー\tl{˥}}
\rentry{どぼんと}{}{だぼんた}
\rentry{とまどう}{}{すまどぅるん\tl{˥˩}}
\rentry{とまらせる}{泊まらせる}{とぅまらすん\tl{˥˩}}
\rentry{とまる}{泊まる}{とぅまるん\tl{˥˩}、やどぅとぅるん}
\rentry{とめる}{止める}{とぅみるん\tl{˥}、やみるん\tl{˥˩}}
\rentry{とも}{艫}{とぅむ\tl{˥}}
\rentry{ともだち}{友達}{どぅし\tl{˥˩}}
\rentry{どもりがちだ}{吃りがちだ}{ふちぃくぱはん}
\rentry{どもる}{吃る}{したくぱるん\tl{˥}}
\rentry{どやしつける}{}{しゃみしきるん}
\rentry{どやす}{}{うどぅがすん\tl{˥}}
\rentry{とよます}{}{とぅゆますん\tl{˥}}
\rentry{とよむ}{}{とぅゆまりん\tl{˥}}
\rentry{とら}{虎}{とぅら\tl{˥˩}}
\rentry{どら}{銅鑼}{どらん\tl{˥˩}}
\rentry{とらえる}{捕える}{すかむん\tl{˥}}
\rentry{とらせる}{取らせる}{とらしみん}
\rentry{とり}{酉}{とぅり\tl{˥˩}}
\rentry{とり}{鳥}{とぅり\tl{˥˩}}
\rentry{どりーね}{ドリーネ}{あぶ\tl{˥˩}}
\rentry{とりあつかう}{取り扱う}{むたすくん\tl{˩˥}}
\rentry{とりかえす}{取り返す}{とぅりかいすん\tl{˥}、とぅりけーすん\tl{˥}}
\rentry{とりかえる}{取り替える}{とぅりかいるん\tl{˥}}
\rentry{とりくむ}{取り組む}{とぅりくむん\tl{˥}}
\rentry{とりこむ}{取り込む}{とぅりくむん}
\rentry{とりさる}{取り去る}{\josho{}とぅり\kako{}\ws{}っしん\tl{˥˩}}
\rentry{とりしまる}{取り締まる}{とぅりしまるん\tl{˥}}
\rentry{とりだす}{取り出す}{とぅりんだすん\tl{˥}}
\rentry{とりたてる}{取り立てる}{とぅりたちん\tl{˥}}
\rentry{とりちがえる}{取り違える}{まちがいん\tl{˥}}
\rentry{とりにがす}{取り逃がす}{とぅりぴんがすん\tl{˥}}
\rentry{とりのす}{鳥の巣}{しー\tl{˥}}
\rentry{とりのぞく}{取り除く}{とぅり\tl{˥}\ws{}っしるん\tl{˥˩}}
\rentry{とりはからう}{取り計らう}{とぅりはかろーん\tl{˥}}
\rentry{とりはずす}{取り外す}{とぅり\tl{˥}ぱんつん\tl{˥˩}}
\rentry{とりまぜる}{取り混ぜる}{ぶくむん\tl{˥}}
\rentry{とりもどす}{取り戻す}{とぅりかいすん\tl{˥}、とぅりむどぅすん\tl{˥}}
\rentry{どりょくする}{努力する}{ぎばるん\tl{˥˩}}
\rentry{とりよせる}{取り寄せる}{とぅりゆしるん}
\rentry{とる}{取る}{とぅるん\tl{˥}}
\rentry{どれ}{}{ざー\tl{˩˥}}
\rentry{どれほど}{}{いかすく}
\rentry{どろ}{泥}{どぅーる\tl{˩˥}}
\rentry{とろう}{徒労}{んなあわり\tl{˥˩}}
\rentry{とろかす}{}{たらすん\tl{˥˩}}
\rentry{どろだらけになる}{泥だらけになる}{どぅるびちゃー\tl{˩˥}\ws{}なるん\tl{˩˥}}
\rentry{どろぼう}{泥棒}{ぬすとぅり\tl{˩˥}}
\rentry{どろみち}{泥道}{どぅるみち\tl{˩˥}}
\rentry{どんてん}{曇天}{あまおしき\tl{˥}}
\rentry{どんと}{}{だーんた}
\rentry{どんどん}{}{みっ\josho{}ふぁ\kako{}ん}
\rentry{どんな}{}{ぬ\josho{}ー\kako{}しゃる、ねー\tl{˩˥}、ねーしゃるん}
\rentry{どんなふうに}{どんな風に}{ねーる\tl{˩˥}}
\rentry{どんぶり}{丼}{どんぶり\tl{˩˥}}
\rentry{とんぼ}{蜻蛉}{あいじ\tl{˥}}
\rentry{な}{名}{なー\tl{˥˩}、なん\tl{˥˩}}
\rentry{な}{菜}{なん\tl{˩˥}}
\rentry{なーらぼん}{ナーラボン}{なーらぼん\tl{˥}}
\rentry{ない}{無い}{ねーぬ\tl{˩˥}}
\rentry{ないがしろにする}{}{かるんじるん\tl{˥˩}}
\rentry{ないしむら}{長石村}{なすむら\tl{˥˩}}
\rentry{ないぞう}{内臓}{ばた\tl{˩˥}}
\rentry{ないちじん}{内地人}{やまとぅ\tl{˩˥}ぴぃとぅ\tl{˥˩}}
\rentry{ないふ}{ナイフ}{しぐ\tl{˥}}
\rentry{なう}{綯う}{なーすん\tl{˩˥}、なすん\tl{˩˥}}
\rentry{なえる}{萎える}{だーりるん\tl{˩˥}}
\rentry{なおす}{治す}{のーすん\tl{˩˥}}
\rentry{なおす}{直す}{のーすん\tl{˩˥}}
\rentry{なおる}{治る}{のーるん\tl{˩˥}、みしゃー\tl{˥}\ws{}なるん\tl{˩˥}}
\rentry{なおる}{直る}{のーるん\tl{˩˥}}
\rentry{なか}{中}{なが\tl{˥˩}}
\rentry{ながあめ}{長雨}{ながあみ\tl{˥˩}}
\rentry{ながい}{永い}{なー\tl{˩˥}はん}
\rentry{ながい}{長い}{なー\tl{˩˥}はん}
\rentry{ながい}{長居}{ながびり\tl{˥˩}}
\rentry{ながいあいだ}{長い間}{なげーくとぅ}
\rentry{ながいき}{長生き}{ながいき\tl{˥˩}}
\rentry{ながいきする}{長生きする}{ながいぎ\tl{˥˩}\ws{}すん\tl{˥˩}}
\rentry{ながいこと}{長いこと}{なげーくとぅ}
\rentry{なかがいい}{仲がいい}{むつまっ\tl{˩˥}さーん、むつまっさ\tl{˩˥}はん}
\rentry{ながさ}{長さ}{なぎ\tl{˩˥}}
\rentry{なかす}{泣かす}{なーがすん\tl{˥˩}}
\rentry{ながす}{流す}{なんぶらすん\tl{˩˥}、〔血、汗などを\ntilde{}〕ぱらすん\tl{˥}}
\rentry{なかど}{中戸}{ながやどぅ\tl{˩˥}}
\rentry{なかなおりする}{仲直りする}{すむぴらすん\tl{˥}}
\rentry{ながながとした}{長々とした}{な\josho{}ー\kako{}なー\ws{}しゃん\tl{˥˩}}
\rentry{なかなかみることのない}{中々見ることのない}{みーどぅー\tl{˩˥}さーん}
\rentry{なかに}{中に}{なが\tl{˥˩}}
\rentry{ながびく}{長引く}{ながびくん\tl{˥˩}}
\rentry{なかま}{仲間}{ぐー\tl{˩˥}}
\rentry{なかみのないしるもの}{中味のない汁物}{んなすー\tl{˥˩}}
\rentry{ながめる}{眺める}{ながみん\tl{˥}}
\rentry{ながもち}{長持ち}{なーばぐ\tl{˩˥}}
\rentry{ながもちする}{長持ちする}{たむつん\tl{˥}}
\rentry{ながれだす}{流れ出す}{なんぶりん\tl{˩˥}}
\rentry{ながれぼし}{流れ星}{ゆーべーぷち}
\rentry{ながれる}{流れる}{なーりん\tl{˩˥}}
\rentry{ながわずらい}{長患い}{ながやみ\tl{˥˩}}
\rentry{なきかかる}{泣きかかる}{なぎすかり\tl{˥˩}}
\rentry{なきくたびれる}{泣きくたびれる}{なぎぼたりるん\tl{˩˥}}
\rentry{なきこがれる}{泣き焦がれる}{なきくがりん}
\rentry{なぎさ}{}{すーぬ\tl{˥}\ws{}ふち\tl{˥˩}、\josho{}ぱま\kako{}ふちぃ}
\rentry{なきしおれる}{泣きしおれる}{なぎぼたりるん\tl{˩˥}}
\rentry{なぎたおす}{薙ぎ倒す}{けーり\tl{˥}\ws{}とーすん\tl{˥}、なぎとーすん\tl{˩˥}}
\rentry{なきつかれる}{泣き疲れる}{なぎぼたりるん\tl{˩˥}}
\rentry{なきつく}{泣きつく}{なぎすかり\tl{˥˩}}
\rentry{なきまね}{泣き真似}{なぎまーべ\tl{˥˩}}
\rentry{なきむしだ}{泣き虫だ}{なちぶさー\tl{˥˩}}
\rentry{なく}{泣く}{なーぐん\tl{˥˩}}
\rentry{なく}{鳴く}{なーぐん\tl{˥˩}}
\rentry{なぐ}{凪ぐ}{とぅりるん\tl{˥˩}、なだらぐん\tl{˥˩}}
\rentry{なぐさめる}{慰める}{なぐさみん\tl{˥˩}}
\rentry{なくす}{}{うすなすん\tl{˥˩}、ねーな\tl{˩˥}さん\tl{˥˩}}
\rentry{なくなった}{}{ねーな\tl{˩˥}さん\tl{˥˩}}
\rentry{なくなる}{}{ねー\tl{˩˥}\ws{}なるん\tl{˩˥}}
\rentry{なくなる}{亡くなる}{まらすん\tl{˥˩}}
\rentry{なぐる}{殴る}{くらすん\tl{˥˩}、すなすん\tl{˥˩}、たたぐん\tl{˥}、ばみがすん\tl{˥˩}}
\rentry{なげあみ}{投げ網}{うちあん\tl{˥}}
\rentry{なげいれる}{投げ入れる}{だくむん\tl{˩˥}}
\rentry{なげきかなしむさま}{嘆き悲しむさま}{あがよー}
\rentry{なげこむ}{投げ込む}{うちくむん\tl{˥}}
\rentry{なげすてる}{投げ捨てる}{すか\ws{}っしりん、なんが\tl{˩˥}\ws{}ししるん\tl{˥˩}、なんが\tl{˩˥}\ws{}っしん\tl{˩}、ぼんぎるん\tl{˩˥}}
\rentry{なげてたたきつける}{投げて叩きつける}{たたしきるん\tl{˥}}
\rentry{なげる}{投げる}{なんぎるん\tl{˩˥}}
\rentry{なこうど}{仲人}{なかだち\tl{˩˥}}
\rentry{なごりおしい}{名残惜しい}{なぐりしゃ\tl{˩˥}はん}
\rentry{なさけ}{情け}{なさき\tl{˩˥}}
\rentry{なじむ}{}{なすぅくん\tl{˥˩}}
\rentry{なすび}{茄子}{なしぴ\tl{˩˥}}
\rentry{なぜ}{何故}{ねーき\tl{˩˥}、ねーきる}
\rentry{なぞなぞ}{謎々}{ぬーとぅ}
\rentry{なだかい}{名高い}{うとぅだが\tl{˥˩}はん、なー\tl{˥˩}たか\tl{˥}はん}
\rentry{なだめる}{宥める}{なだみるん\tl{˥˩}}
\rentry{なだらかだ}{}{なだらぎゃーん}
\rentry{なだらかにする}{}{なだらぎるん\tl{˥˩}}
\rentry{なつ}{夏}{なちぃ\tl{˥˩}}
\rentry{なつかしい}{懐かしい}{なすかっさ\tl{˩˥}はん}
\rentry{なつぎ}{夏着}{なちすぅぬ\tl{˥˩}}
\rentry{なつく}{懐く}{なすぅくん\tl{˥˩}}
\rentry{なづけおや}{名付け親}{なーしきうや\tl{˩˥}}
\rentry{なつまけ}{夏負け}{なちぃまき\tl{˥˩}}
\rentry{なでる}{撫でる}{しっつぁーすん\tl{˥}}
\rentry{ななじゅう}{七十}{ななず\tl{˩˥}}
\rentry{ななつ}{七つ}{ななち\tl{˩˥}}
\rentry{なに}{何}{ぬー\tl{˩˥}}
\rentry{なにが}{何が}{ぬ\josho{}ー\kako{}や}
\rentry{なにごと}{何事}{ぬーぐとぅ}
\rentry{なにも}{何も}{ぬ\josho{}ー\kako{}ん}
\rentry{なのか}{七日}{なんが\tl{˩˥}}
\rentry{なのる}{名乗る}{なん\tl{˥˩}\ws{}えぬん\tl{˥˩}}
\rentry{なべ}{鍋}{なび\tl{˩˥}}
\rentry{なべをこがす}{鍋を焦がす}{なびふつぁらすん\tl{˩˥}}
\rentry{なま}{生}{なま\tl{˥}}
\rentry{なまえ}{名前}{なー\tl{˥˩}、なん\tl{˥˩}}
\rentry{なまき}{生木}{なまき\tl{˩˥}}
\rentry{なまぐさい}{生臭い}{おーふつぁ\tl{˥˩}はーん、なまふささーん、なまふつぁ\tl{˩˥}はん}
\rentry{なまけ}{怠け}{ふゆー\tl{˥}}
\rentry{なまけもの}{怠け者}{なまふつぁりむぬ\tl{˩˥}、ふゆーむぬ\tl{˥}}
\rentry{なまこ}{海鼠}{しきり\tl{˥˩}}
\rentry{なまごめ}{生米}{なまぐみ\tl{˩˥}}
\rentry{なまもの}{生物}{なまむぬ\tl{˩˥}}
\rentry{なまる}{鈍る}{なまるん\tl{˩˥}}
\rentry{なみ}{波}{なん\tl{˥˩}}
\rentry{なみうちぎわ}{波打ち際}{すーぬ\tl{˥}\ws{}ふち\tl{˥˩}}
\rentry{なみだ}{涙}{なだ\tl{˩˥}、なんだ\tl{˩˥}}
\rentry{なみだがでる}{涙が出る}{なんだ\tl{˩˥}\ws{}んじるん\tl{˥˩}}
\rentry{なみだぐむ}{涙ぐむ}{なんだ\tl{˩˥}\ws{}んじるん\tl{˥˩}}
\rentry{なめらかだ}{}{なふこ\tl{˩˥}はん}
\rentry{なめる}{舐める}{しぃぴるん\tl{˥˩}、すぷるん\tl{˥˩}、ぬっふぇーるん\tl{˥˩}}
\rentry{なやむ}{悩む}{すむ\tl{˥}やむん\tl{˩˥}}
\rentry{ならい}{}{なれー\tl{˩˥}}
\rentry{ならう}{習う}{ならすん\tl{˩˥}}
\rentry{ならす}{均す}{なだぎるん\tl{˥˩}、なだらがすん\tl{˥˩}}
\rentry{ならす}{鳴らす}{ならすん\tl{˥˩}}
\rentry{ならぶ}{並ぶ}{ならぶん\tl{˥˩}}
\rentry{ならべる}{並べる}{すながるん\tl{˥˩}、ならびるん\tl{˥˩}}
\rentry{なりわい}{生業}{むぬへーずく\tl{˩˥}}
\rentry{なる}{鳴る}{なーるん\tl{˥˩}}
\rentry{なるほど}{}{えーなー}
\rentry{なれる}{慣れる}{なーれるん\tl{˩˥}、ならすん\tl{˩˥}}
\rentry{なわ}{縄}{すな\tl{˥}}
\rentry{なわしろ}{苗代}{なっす\tl{˩˥}}
\rentry{なわしろいちご}{ナワシロイチゴ}{てーし\tl{˥}}
\rentry{なわしろだ}{苗代田}{なっすだー}
\rentry{なをつける}{名を付ける}{なーしきん\tl{˩˥}}
\rentry{なん～}{何～}{うー}
\rentry{なんかい}{何回}{うーむし}
\rentry{なんぎ}{難儀}{あわり\tl{˥}、なんぎ\tl{˩˥}}
\rentry{なんきんぶくろ}{南京袋}{\josho{}かしがー\kako{}ふくる}
\rentry{なんさい}{何歳}{うーちぃ\tl{˥˩}}
\rentry{なんせいほう}{南西方}{\josho{}さん\kako{}ば}
\rentry{なんたることか}{}{あきさみよー、はっさみよー}
\rentry{なんちょうになったひと}{難聴になった人}{みんとーりむん\tl{˩˥}}
\rentry{なんちょうになる}{難聴になる}{みんとーら\tl{˩˥}}
\rentry{なんでも}{}{ぬ\josho{}ー\ws{}やばん\tl{˩}}
\rentry{なんと}{}{ぬ\josho{}ー\kako{}た、ぬーたる}
\rentry{なんといってるか}{何と言ってるか}{ぬ\josho{}ー\kako{}た、ぬーたる}
\rentry{なんにん}{何人}{うたーり\tl{˥˩}}
\rentry{なんぽう}{南方}{ぺー\tl{˥}、\josho{}ぺーぬ\kako{}\ws{}かた、\josho{}ぺったぬ\kako{}\ws{}かた}
\rentry{にあう}{似合う}{うつるん\tl{˥}、ぱいるん\tl{˥}}
\rentry{にいさん}{兄さん}{びらま\tl{˩˥}}
\rentry{にえたぎらせる}{煮えたぎらせる}{たぎらすん\tl{˥˩}}
\rentry{にえたぎる}{煮えたぎる}{たぎるん\tl{˥˩}}
\rentry{にえる}{煮える}{ねーるん\tl{˥˩}}
\rentry{におい}{匂い}{かー\tl{˥˩}}
\rentry{においがする}{匂いがする}{かーすん\tl{˥˩}}
\rentry{においがつよい}{匂いが強い}{かーぬ\ws{}すさはぬ}
\rentry{におう}{匂う}{かーすん\tl{˥˩}}
\rentry{にがい}{苦い}{ざー\tl{˩˥}はん}
\rentry{にかえす}{煮返す}{ふつぁーすん\tl{˥}}
\rentry{にがおもい}{荷が重い}{にんぬ\ws{}んさはん}
\rentry{にがす}{逃がす}{ぴんがすん\tl{˥}}
\rentry{にがたけ}{苦竹}{ざーだぎ\tl{˩˥}}
\rentry{にぎやかにさわぐ}{賑やかに騒ぐ}{はなやかすん\tl{˥}}
\rentry{にぎわす}{賑わす}{はなやかすん\tl{˥}}
\rentry{にく}{肉}{にく\tl{˩˥}}
\rentry{にくい}{憎い}{にた\tl{˩˥}はん}
\rentry{にくしん}{肉親}{うや\tl{˥}}
\rentry{にくまれぐち}{憎まれ口}{ばたふちりむぬ\tl{˩˥}}
\rentry{にくらしい}{憎らしい}{みっふぁはん}
\rentry{にげまわる}{逃げ回る}{ぴんぎまーるん\tl{˥}}
\rentry{にげる}{逃げる}{ぴんぎるん\tl{˥}}
\rentry{にこにことわらう}{にこにこと笑う}{さにさに\ws{}しゃん}
\rentry{にごる}{濁る}{やなごるん\tl{˥˩}}
\rentry{にし}{西}{いり\tl{˥˩}}
\rentry{にじ}{虹}{のーじん\tl{˥˩}}
\rentry{にしむら}{西村}{いりむら}
\rentry{にしめる}{煮しめる}{んぶすん\tl{˥}}
\rentry{にじゅう}{二十}{にんず}
\rentry{にしゅうき}{二週忌}{ふたなんが\tl{˥˩}}
\rentry{にじりよる}{にじり寄る}{しきつぁーすん\tl{˥˩}}
\rentry{にせる}{似せる}{にーらすん}
\rentry{にたつ}{煮立つ}{ふつん\tl{˥}}
\rentry{にちぼつ}{日没になる}{しな\tl{˥}いるん\tl{˥˩}}
\rentry{にっすう}{日数}{ぴんちぃ\tl{˥˩}}
\rentry{にっちゅう}{日中}{ぴぃす\tl{˥˩}、まーびる\tl{˥˩}}
\rentry{にている}{似ている}{にししゃ\tl{˥˩}はん}
\rentry{にとうぶん}{二等分}{ふたばぎ\tl{˥˩}}
\rentry{にばい}{二倍}{まーおーび}
\rentry{にぶる}{鈍る}{〔切れ味が\ntilde{}〕なまるん\tl{˩˥}}
\rentry{にほん}{日本}{やまとぅ\tl{˩˥}}
\rentry{にほんご}{日本語}{やまとぅむに\tl{˩˥}}
\rentry{にほんじん}{日本人}{やまとぅ\tl{˩˥}ぴぃとぅ\tl{˥˩}}
\rentry{にほんほんど}{日本本土}{やまとぅ\tl{˩˥}}
\rentry{にもつ}{荷物}{にむちぃ\tl{˩˥}}
\rentry{にゅうどうぐも}{入道雲}{たかふもん\tl{˥}}
\rentry{にら}{韮}{びら\tl{˩˥}}
\rentry{にらむ}{睨む}{みーしきるん}
\rentry{にる}{似る}{にししゃ\tl{˥˩}はん}
\rentry{にる}{煮る}{ねっすん\tl{˥˩}、ばがすん\tl{˥˩}}
\rentry{にわ}{庭}{みなー\tl{˥˩}、みなが\tl{˥˩}}
\rentry{にわかあめ}{俄雨}{あたあみ}
\rentry{にわとり}{鶏}{ごっか\tl{˩˥}}
\rentry{にわとりのたまご}{ニワトリの卵}{ごかけー\tl{˥}}
\rentry{にんげん}{人間}{にんぎん\tl{˩˥}、ぴぃとぅ\tl{˥˩}}
\rentry{にんじょうがない}{人情がない}{すむ\tl{˥}\ws{}ねーぬ\tl{˩˥}}
\rentry{にんじん}{ニンジン}{あがでーぐに\tl{˥˩}}
\rentry{にんしんする}{妊娠する}{ぱるむん\tl{˥}}
\rentry{にんずう}{人数}{にんずー\tl{˩˥}}
\rentry{にんたい}{忍耐}{がー\tl{˩˥}}
\rentry{にんたいりょくがつよい}{忍耐力が強い}{がーずさはん}
\rentry{にんにく}{大蒜}{ぴる\tl{˥˩}}
\rentry{にんむ}{任務}{しゅくぶん\tl{˩˥}}
\rentry{ぬいもの}{縫い物}{ぬいむぬ\tl{˩˥}}
\rentry{ぬう}{縫う}{ぬーん\tl{˩˥}}
\rentry{ぬか}{糠}{ぬが\tl{˩˥}}
\rentry{ぬきだす}{抜き出す}{ぴき\tl{˥˩}ぬぐん\tl{˩˥}}
\rentry{ぬきや}{貫き家}{ぬきひー}
\rentry{ぬきんでる}{抜きんでる}{ぬぎんだん}
\rentry{ぬく}{抜く}{ぬぐん\tl{˩˥}}
\rentry{ぬぐ}{脱ぐ}{ぬぐん\tl{˩˥}、〔着物を\ntilde{}〕ぱちるん\tl{˥}}
\rentry{ぬくく}{温く}{あ\josho{}てー\kako{}し}
\rentry{ぬける}{抜ける}{ぬぎるん\tl{˩˥}}
\rentry{ぬすっと}{盗人}{ぬすとぅり\tl{˩˥}}
\rentry{ぬすむ}{盗む}{ぬすむん\tl{˩˥}}
\rentry{ぬの}{布}{ぬぬ\tl{˥˩}}
\rentry{ぬのさらし}{布晒し}{ぬのさらし\tl{˥˩}}
\rentry{ぬま}{沼}{あなぶ\tl{˥}}
\rentry{ぬらす}{濡らす}{ぞっふゎらすん\tl{˥˩}、ぞっふゎるん\tl{˥˩}}
\rentry{ぬりたくる}{塗りたくる}{ぬりだっくゎーすん}
\rentry{ぬりつける}{塗りつける}{だっくゎすん}
\rentry{ぬる}{塗る}{ぬるん\tl{˥˩}}
\rentry{ぬるくする}{温くする}{ぬるみるん\tl{˩˥}}
\rentry{ぬれる}{濡れる}{ぞっふゎりん\tl{˥˩}}
\rentry{ね}{子}{にー\tl{˥}}
\rentry{ね}{根}{にん\tl{˩˥}}
\rentry{ねいる}{寝入る}{たーりん\tl{˥˩}、ぬっふたりるん}
\rentry{ねがいでる}{願い出る}{にげー\tl{˩˥}\ws{}んじるん\tl{˥}}
\rentry{ねがう}{願い}{にげー\tl{˩˥}}
\rentry{ねかしつける}{寝かしつける}{ぬふちるん}
\rentry{ねかせる}{寝かせる}{ぬふちるん}
\rentry{ねぎ}{葱}{しびる\tl{˥˩}}
\rentry{ねこ}{猫}{まゆ\tl{˥˩}}
\rentry{ねこそぎとる}{根こそぎ取る}{こーすん\tl{˥˩}}
\rentry{ねじまがる}{ねじ曲がる}{むじかるん\tl{˩˥}}
\rentry{ねじまがる}{捩じ曲がる}{むじかーるん\tl{˩˥}}
\rentry{ねじまげる}{ねじ曲げる}{むじまーすん\tl{˩˥}}
\rentry{ねしょうべん}{寝小便}{ゆすぱり\tl{˩˥}}
\rentry{ねじる}{捩じる}{むじまーすん\tl{˩˥}}
\rentry{ねじれる}{捩じれる}{むじかーるん\tl{˩˥}、むじかるん\tl{˩˥}}
\rentry{ねそべる}{寝そべる}{ゆくたーるん\tl{˩˥}}
\rentry{ねたましい}{妬ましい}{にた\tl{˩˥}はん}
\rentry{ねたむ}{妬む}{にたむん\tl{˩˥}}
\rentry{ねだやしになる}{根絶やしになる}{〔作物が\ntilde{}〕たに\ws{}しさん}
\rentry{ねだる}{}{いみるん\tl{˥}、きぬん\tl{˥˩}}
\rentry{ねだん}{値段}{でん\tl{˥˩}}
\rentry{ねつ}{熱}{にちぃ\tl{˩˥}}
\rentry{ねつがでる}{熱が出る}{にち\tl{˩˥}\ws{}んじるん\tl{˥}}
\rentry{ねづく}{根付く}{にんうりん、にん\tl{˩˥}すくん\tl{˥}}
\rentry{ねっしん}{熱心}{すむ\tl{˥}いり\tl{˥˩}}
\rentry{ねっちゅう}{熱中}{まーぷり\tl{˥˩}}
\rentry{ねっとう}{熱湯}{あちぃゆ\tl{˥}}
\rentry{ねつびょう}{熱病}{ぷーき\tl{˥˩}、やきー\tl{˩˥}}
\rentry{ねどこ}{寝床}{とぅく\tl{˥˩}}
\rentry{ねばつく}{粘付く}{むっつぁるん\tl{˥˩}}
\rentry{ねばっこい}{粘っこい}{むっつぁ\tl{˩˥}はん}
\rentry{ねばねばする}{}{むっつぁ\tl{˩˥}はん}
\rentry{ねばりけのないさま}{粘り気のない様}{さぱさぱ}
\rentry{ねぼう}{寝坊}{にーぶやー\tl{˥˩}}
\rentry{ねぼすけ}{寝坊助}{にーぶやー\tl{˥˩}}
\rentry{ねむたい}{眠たい}{ぬふた\tl{˩˥}はーん}
\rentry{ねむる}{眠る}{ぬっふん\tl{˥˩}}
\rentry{ねもと}{根元}{にんむとぅ\tl{˩˥}}
\rentry{ねる}{寝る}{ぬっふん\tl{˥˩}}
\rentry{ねわけする}{根分けする}{にんばぎるん\tl{˩˥}}
\rentry{ねん}{年}{にん\tl{˥˩}}
\rentry{ねんぐ}{年貢}{ぞーぬ\tl{˩˥}、にんぐ\tl{˩˥}}
\rentry{ねんずる}{念ずる}{にんじるん\tl{˥˩}}
\rentry{ねんちょう}{年長}{しじゃー\tl{˥}}
\rentry{ねんど}{粘土}{んた\tl{˩˥}}
\rentry{ねんねん}{年々}{にんにん}
\rentry{ねんぶつしゃ}{念仏者}{にんぶち\tl{˩˥}}
\rentry{ねんりょうゆ}{燃料油}{あば\tl{˥}}
\rentry{ねんれい}{年齢}{とぅちぃ\tl{˥}}
\rentry{の}{野}{ぬー\tl{˩˥}}
\rentry{のいちご}{野いちご}{てーし\tl{˥}}
\rentry{のう}{脳}{のー\tl{˥˩}}
\rentry{のうこう}{農耕}{むぬすくり\tl{˩˥}}
\rentry{のうこうだ}{濃厚だ}{かた\tl{˥˩}はん}
\rentry{のうさぎょう}{農作業}{ぴてー\tl{˥}すくるん\tl{˥}}
\rentry{のうさく}{農作}{さくほー\tl{˥}}
\rentry{のうにゅうする}{納入する}{うさみるん\tl{˥}}
\rentry{のうべんだ}{能弁だ}{ふちぃ\tl{˥˩}やば\tl{˩˥}はん}
\rentry{のうりつがあがる}{能率が上がる}{だっつぁぐん\tl{˥˩}}
\rentry{のからむし}{ノカラムシ}{ばごー\tl{˩˥}}
\rentry{のきさき}{軒先}{あまぱんぎ\tl{˥}}
\rentry{のきば}{軒端}{あまぱんぎ\tl{˥}}
\rentry{のこぎり}{鋸}{ぬぎり\tl{˩˥}}
\rentry{のこす}{残す}{のがすん\tl{˥˩}}
\rentry{のこり}{残り}{ぬぐり\tl{˩˥}、のごり}
\rentry{のこる}{残る}{のがるん\tl{˩˥}}
\rentry{のせる}{乗せる}{ぬしん\tl{˥˩}}
\rentry{のぞみ}{望み}{ぬずみ\tl{˥˩}}
\rentry{のぞむ}{望む}{ぬずむん}
\rentry{のど}{喉}{ぬどぅ\tl{˩˥}}
\rentry{のどがかわく}{喉が渇く}{ぬどぅ\tl{˩˥}\ws{}かりるん\tl{˥˩}}
\rentry{のどもと}{喉元}{ぬぶしん\tl{˩˥}}
\rentry{のばす}{延ばす}{ぬばすん\tl{˩˥}、ぬびるん\tl{˩˥}、ぴそーぎるん}
\rentry{のはら}{野原}{ぬー\tl{˩˥}}
\rentry{のびる}{伸びる}{ぬぶん\tl{˥˩}}
\rentry{のぶどう}{野ブドウ}{かなぶ\tl{˥˩}}
\rentry{のべる}{伸べる}{ぬびるん\tl{˩˥}}
\rentry{のべる}{延べる}{ぴそーぎるん}
\rentry{～のほうへ}{～の方へ}{が}
\rentry{のぼせる}{逆上せる}{ぬぶしるん\tl{˥˩}}
\rentry{のぼりざか}{上り坂}{さか\tl{˥˩}}
\rentry{のぼる}{昇る}{あがるん\tl{˥˩}}
\rentry{のぼる}{登る}{ぬぶるん\tl{˥˩}}
\rentry{のみ}{蚤}{ぬん\tl{˩˥}}
\rentry{のみ}{鑿}{ぬん\tl{˩˥}}
\rentry{のみおろす}{飲み下ろす}{ぬみ\ws{}うらへ}
\rentry{のみこむ}{飲み込む}{ぬみくむん\tl{˩˥}}
\rentry{のみつくす}{飲み尽くす}{ぬぶたん\tl{˩˥}}
\rentry{のみのこす}{飲み残す}{ぬみのがすん\tl{˩˥}}
\rentry{のむ}{飲む}{ぬむん\tl{˩˥}}
\rentry{のらぎ}{野良着}{ぴてーすぅぬ\tl{˥}}
\rentry{のらんそう}{卵巣}{〔鰹など魚の\ntilde{}〕はらみ\tl{˥}}
\rentry{のり}{糊}{ぬり\tl{˩˥}}
\rentry{のりおくれる}{乗り遅れる}{ぬりうくりるん}
\rentry{のりと}{祈詞}{ぱん\tl{˥}}
\rentry{のりはずす}{乗りはずす}{ぬりうくりるん}
\rentry{のりもの}{乗り物}{ぬりむぬ\tl{˥˩}}
\rentry{のる}{乗る}{ぬぶるん\tl{˥˩}、ぬるん\tl{˥˩}}
\rentry{のる}{載る}{ぬるん\tl{˥˩}}
\rentry{のろい}{呪い}{いきろー\tl{˥}}
\rentry{のろい}{鈍い}{どぅなはん}
\rentry{のろのろしている}{}{どぅなー\tl{˩˥}さーん}
\rentry{は}{歯}{ぱん\tl{˥}、ふちぃぬ\tl{˥˩}\ws{}ぱん\tl{˥}}
\rentry{は}{葉}{ぱー\tl{˥˩}}
\rentry{ばあさん}{おばあさん}{うしとぅぱー\tl{˥}}
\rentry{ぱあっと}{}{ぱーった}
\rentry{はあはあ}{}{はーはー}
\rentry{はい}{}{おー}
\rentry{ばい}{倍}{まーおーび}
\rentry{はいからだ}{ハイカラだ}{はいから\tl{˥}さーん}
\rentry{はいきび}{ハイキビ}{にんどーしぃ\tl{˩˥}}
\rentry{ばいしゃくにん}{媒酌人}{なかだち\tl{˩˥}}
\rentry{はいしょ}{拝所}{うがんじゅ、わー\tl{˥}}
\rentry{はいすいこう}{排水溝}{みぞーり\tl{˥˩}}
\rentry{はいせつしちらかす}{排泄し散らかす}{〔大便を\ntilde{}〕ふつ\tl{˥}\ws{}まりぽーるん\tl{˩˥}}
\rentry{はいせつしちらす}{排泄し散らす}{〔大小便を\ntilde{}〕まりぽーるん\tl{˥˩}}
\rentry{はいぞう}{肺臓}{ふく\tl{˥}}
\rentry{はいつくばう}{這いつくばう}{ぺーつぁーるん\tl{˥}、ぺっつぁるん\tl{˥}}
\rentry{はいべんする}{排便する}{ふつ\tl{˥}まるん\tl{˩˥}}
\rentry{はいりこむ}{入り込む}{ぺーりくむん\tl{˥}}
\rentry{はいる}{入る}{ぺーるん\tl{˥}、ぺるん\tl{˥}}
\rentry{はう}{這う}{ぺーつぁーるん\tl{˥}、ぺっかるん\tl{˥}}
\rentry{はえ}{蠅}{ぺー\tl{˥˩}}
\rentry{はえきび}{ハエキビ}{なざぎ\tl{˩˥}}
\rentry{はえる}{映える}{ぱいるん\tl{˥}}
\rentry{はえる}{生える}{むいるん\tl{˩˥}}
\rentry{はか}{墓}{ぱか\tl{˥˩}}
\rentry{ばか}{馬鹿}{ふらー\tl{˥˩}}
\rentry{はがす}{剥がす}{ばがすん\tl{˩˥}}
\rentry{ばかす}{化かす}{ざまどぅらすん\tl{˩˥}}
\rentry{はかせる}{吐かせる}{ぱかすん\tl{˥}}
\rentry{はかどる}{}{だっつぁぐん\tl{˥˩}}
\rentry{ばかにする}{馬鹿にする}{うせーん\tl{˥˩}、ばっくるん\tl{˩˥}}
\rentry{はがね}{鋼}{はーがに\tl{˥}、はがに}
\rentry{はかま}{袴}{ぱかま\tl{˥˩}}
\rentry{はがま}{羽釜}{ぱんがま\tl{˥˩}}
\rentry{ばかもの}{馬鹿者}{ふらー\tl{˥˩}、ぷりむぬ\tl{˥˩}}
\rentry{～ばかり}{}{がーし}
\rentry{はかる}{計る}{ぱかるん\tl{˥}}
\rentry{はかる}{謀る}{はかるん\tl{˥}}
\rentry{はがれる}{剥がれる}{ぱがりるん\tl{˥}、ぱなりるん\tl{˥}}
\rentry{ばかわらい}{馬鹿笑い}{ぷりばーり\tl{˥˩}}
\rentry{はきけがする}{吐き気がする}{〔胸がむかついて\ntilde{}〕いがばりゃん}
\rentry{はきだす}{吐き出す}{ぱきんだすん\tl{˥}}
\rentry{はぎとる}{剥ぎ取る}{ぱぎとぅるん\tl{˥}、ぱぐん\tl{˥}}
\rentry{はきもの}{履物}{ふむむぬ\tl{˥˩}}
\rentry{はく}{}{〔履物を\ntilde{}〕ふむん\tl{˥˩}}
\rentry{はく}{佩く}{ぱくん\tl{˥˩}}
\rentry{はく}{吐く}{ぱくん\tl{˥˩}、むどぅすん\tl{˩˥}}
\rentry{はく}{掃く}{はくん\tl{˥}、ぽーぐん\tl{˥}}
\rentry{はぐ}{剥ぐ}{ぱぐん\tl{˥}}
\rentry{はくじょう}{薄情}{すむ\tl{˥}\ws{}ねーぬ\tl{˩˥}}
\rentry{はくめい}{薄命}{ぬちぃまろ\tl{˩˥}はん}
\rentry{はげあたま}{禿頭}{ぱぎすぷる\tl{˥}}
\rentry{ばけつ}{バケツ}{ばぎし\tl{˩˥}}
\rentry{はげる}{禿げる}{ぱぎん\tl{˥}}
\rentry{はこ}{箱}{ぱく\tl{˥˩}}
\rentry{はごたえがある}{歯ごたえがある}{こー\tl{˥}はん、こすぱり}
\rentry{はさき}{刃先}{ぱん\tl{˥}}
\rentry{はさみ}{鋏}{ぱつぁん\tl{˥}}
\rentry{はさむ}{挟む}{ぱつぁむん\tl{˥}}
\rentry{はし}{橋}{ぱち\tl{˥˩}}
\rentry{はし}{端}{ぱた\tl{˥˩}、ぱんた\tl{˥˩}}
\rentry{はし}{箸}{まーし\tl{˥}}
\rentry{はし}{隅}{ゆごー\tl{˥˩}}
\rentry{はじく}{弾く}{ぱんくん\tl{˥}}
\rentry{はじける}{弾ける}{ぱしゅくりん\tl{˥}}
\rentry{はしご}{梯子}{ぱち\tl{˥˩}}
\rentry{はしっこ}{端っこ}{ぱんた\tl{˥˩}}
\rentry{はじまる}{始まる}{はじまるん\tl{˥˩}}
\rentry{はじめる}{始める}{はじみるん\tl{˥˩}}
\rentry{はしゃぐ}{}{ばざるん\tl{˥˩}}
\rentry{ばしょ}{場所}{とぅくる\tl{˥}、ばしゅ\tl{˩˥}}
\rentry{ばしょう}{芭蕉}{ばーさ\tl{˩˥}}
\rentry{ばしょうふのきもの}{芭蕉布の着物}{ばさすぅぬ\tl{˩˥}}
\rentry{はしら}{柱}{ぱら\tl{˥}}
\rentry{はしらせる}{走らせる}{ぱらすん\tl{˥}}
\rentry{はしる}{走る}{はきしぃきるん\tl{˥}、ぱるん\tl{˥}}
\rentry{はずかしい}{恥ずかしい}{〔人に対して\ntilde{}〕ぱすぅか\tl{˥}さーん、ぱすかはん}
\rentry{はずす}{外す}{ぱんつん}
\rentry{はすのはぎり}{ハスノハギリ}{さこだち\tl{˥˩}}
\rentry{はずれる}{外れる}{ぬぎるん\tl{˩˥}、ぱんちん\tl{˥˩}}
\rentry{はぜのき}{ハゼノキ}{ぱんじ\tl{˥˩}}
\rentry{はた}{ハタ}{にばり\tl{˩˥}}
\rentry{はた}{旗}{ぱた\tl{˥˩}}
\rentry{はたおりき}{機織り機}{ぱたむん\tl{˥˩}}
\rentry{はだか}{裸}{あばだり\tl{˥}}
\rentry{はたけ}{畑}{ぴてー\tl{˥}、ぴてーぎ\tl{˥}}
\rentry{はたけごや}{畑小屋}{ぴてーひ\tl{˥}}
\rentry{はたけしごと}{畑仕事}{ぴてー\tl{˥}すくるん\tl{˥}}
\rentry{はだし}{裸足}{からぱん\tl{˥˩}}
\rentry{はたす}{果たす}{うむいきすん\tl{˥}}
\rentry{はたち}{二十歳}{ぱたち\tl{˥˩}}
\rentry{はたらく}{働く}{ぱたらぐん\tl{˥˩}}
\rentry{はち}{蜂}{ぱーち\tl{˥˩}}
\rentry{ばち}{}{〔大鼓の\ntilde{}〕ばち}
\rentry{はちじゅう}{八十}{ぱちず\tl{˥}}
\rentry{はちまき}{鉢巻}{さち\tl{˥˩}}
\rentry{ばつ}{罰}{とぅが\tl{˥}、ばち\tl{˥˩}、ばつぁ\tl{˩˥}}
\rentry{はつじょうする}{発情する}{ずーぶるん\tl{˥}}
\rentry{ばっする}{罰する}{とぅがみるん\tl{˥}}
\rentry{ばった}{バッタ}{かたん\tl{˥˩}}
\rentry{はったいこ}{はったい粉}{〔小麦の\ntilde{}〕ゆぬぐ\tl{˩˥}}
\rentry{はってすすむ}{這って進む}{ぴゃっかるん\tl{˥}}
\rentry{はってんする}{発展する}{ひらきるん\tl{˥}}
\rentry{はつどうきせん}{発動機船}{きかんしん\tl{˥}}
\rentry{はつねつする}{発熱する}{にち\tl{˩˥}\ws{}んじるん\tl{˥}}
\rentry{はつほまつり}{初穂祭}{すくまん\tl{˥}}
\rentry{はつもの}{初物}{ぱち\tl{˥˩}}
\rentry{ばつをかぶる}{罰を被る}{ばつぁかぷん}
\rentry{はて}{}{あい}
\rentry{はてるま}{波照間}{ぱちるま、はてろー\tl{˥˩}、ぱてろーま\tl{˥˩}}
\rentry{はてるまじま}{波照間島}{はてろー\tl{˥˩}、ぱてろーま\tl{˥˩}、〔島民間で\ntilde{}〕べーすま}
\rentry{はと}{鳩}{ぱとん\tl{˥}}
\rentry{はとま}{鳩間}{ぱとーま\tl{˥}}
\rentry{はとまじま}{鳩間島}{ぱとーま\tl{˥}}
\rentry{はな}{花}{ぱな\tl{˥}}
\rentry{はな}{鼻}{ぱな\tl{˥˩}}
\rentry{はないけ}{花生け}{ぱないぎ\tl{˥}}
\rentry{はなし}{話}{ぱなし\tl{˥}}
\rentry{はなじる}{鼻汁}{ぱなだり\tl{˥˩}}
\rentry{はなす}{話す}{かたるん\tl{˥˩}、ぱなすん\tl{˥}}
\rentry{はなす}{離す}{ばがすん\tl{˩˥}、ぱなすん\tl{˥}}
\rentry{はなつ}{放つ}{〔糞を\ntilde{}〕まるん\tl{˩˥}}
\rentry{はなづな}{鼻綱}{〔牛の\ntilde{}〕ぱなじな\tl{˥}}
\rentry{ばなな}{バナナ}{ばーさ\tl{˩˥}}
\rentry{ばななのみ}{バナナの実}{ばーさぬ\tl{˩˥}\ws{}なーり\tl{˩˥}}
\rentry{はなよめ}{花嫁}{あいなー\tl{˥}}
\rentry{はなれ}{離れ}{ぱなり\tl{˥}}
\rentry{はなれのしま}{離れの島}{ぱなりぬ\ws{}すま\tl{˥}}
\rentry{はなれる}{離れる}{ぬぎるん\tl{˩˥}、ぱなりるん\tl{˥}、ぱなりん\tl{˥}、ひだみるん\tl{˥}、ひだみん\tl{˥}}
\rentry{はね}{羽}{ぱに\tl{˥˩}}
\rentry{はねかえす}{跳ね返す}{〔物に突き当たって\ntilde{}〕はにかいすん\tl{˥}}
\rentry{はねかえる}{跳ね返る}{はにかいるん\tl{˥}}
\rentry{はは}{母}{あぶゎ\tl{˥}}
\rentry{ぱぱいあ}{パパイア}{まんじょー\tl{˥}}
\rentry{ははおや}{母親}{あぶゎ\tl{˥}}
\rentry{ははかたしんせき}{母方親戚}{じーかたうやぐ、ちーかたうやぐ\tl{˥}}
\rentry{はへん}{破片}{はき\tl{˥˩}}
\rentry{はま}{浜}{ぱま\tl{˥}}
\rentry{はまかに}{浜蟹}{ぱまかん\tl{˥}}
\rentry{はましたん}{浜紫檀}{すぱんこっち}
\rentry{はまじゃり}{浜砂利}{な\josho{}り\kako{}さ}
\rentry{はまゆう}{浜木綿}{やたかび}
\rentry{はまる}{}{ぱまるん\tl{˥}}
\rentry{はめこむ}{はめ込む}{ぱみるん\tl{˥}}
\rentry{はめる}{}{ぱみるん\tl{˥}}
\rentry{ばめん}{場面}{ばしゅ\tl{˩˥}}
\rentry{はやうまれ}{早生まれ}{はやまり\tl{˥}}
\rentry{はやく}{早く}{がんがん、ぴく、ぴくだり\tl{˥˩}、ぺーしゃな\tl{˥}}
\rentry{はやくする}{早くする}{はやみるん\tl{˥}}
\rentry{はやくなる}{早くなる}{はやまるん\tl{˥}}
\rentry{はやじにする}{早死する}{はやじに\tl{˥}\ws{}すん\tl{˥˩}}
\rentry{はやしのなか}{林の中}{きぬ\tl{˥}\ws{}みー\tl{˥˩}}
\rentry{はやす}{囃す}{ぱーすん\tl{˩˥}、ぱやすん\tl{˩˥}}
\rentry{はやす}{生やす}{〔草を\ntilde{}〕むやすん\tl{˩˥}}
\rentry{はやまる}{早まる}{はやまるん\tl{˥}、ぺーまるん\tl{˥}}
\rentry{はやめに}{早めに}{ぴくだり\tl{˥˩}}
\rentry{はやめる}{早める}{はやみるん\tl{˥}}
\rentry{はやらせる}{流行らせる}{ひるみるん\tl{˥}、ひるみん\tl{˥}}
\rentry{はやり}{流行り}{はやり\tl{˥}}
\rentry{はやる}{流行る}{はやーるん\tl{˥}}
\rentry{はら}{腹}{ばた\tl{˩˥}}
\rentry{はらいっぱいになる}{腹一杯になる}{ばた\tl{˩˥}んち\tl{˩˥}}
\rentry{はらう}{払う}{ぱらすん\tl{˥}}
\rentry{はらう}{祓う}{〔不吉を\ntilde{}〕ぱちるん\tl{˥}}
\rentry{はらがしぼるようにいたむ}{腹がしぼるように痛む}{すぶりやみ}
\rentry{はらがふくれる}{腹が膨れる}{〔消化不良で\ntilde{}〕ばたふくりるん\tl{˩˥}}
\rentry{はらがもたれる}{腹がもたれる}{ばたふくら}
\rentry{ばらす}{}{ぽっつぁらすん}
\rentry{ばらばらにする}{}{ぽっつぁらすん}
\rentry{ばらばらになる}{}{ぽっつぁーるん\tl{˥}}
\rentry{ばらばらにわれる}{ばらばらに割れる}{さんざらごーなるん}
\rentry{はらまない}{孕まない}{やじまるん\tl{˥˩}}
\rentry{はらむ}{孕む}{ぱるむん\tl{˥}}
\rentry{はらわた}{腹わた}{ばた\tl{˩˥}}
\rentry{はり}{梁}{かむい\tl{˥}}
\rentry{はり}{針}{ぱり\tl{˥}}
\rentry{はりせんぼん}{ハリセンボン}{いしかぼ\tl{˥}}
\rentry{はる}{張る}{ぱるん\tl{˥˩}}
\rentry{はる}{春}{うりじん\tl{˥}}
\rentry{はるののげし}{ハルノノゲシ}{ふくな\tl{˥˩}}
\rentry{はれがひく}{腫れがひく}{しぴりん\tl{˥˩}}
\rentry{はれつする}{破裂する}{ぱしゅくりん\tl{˥}}
\rentry{はれもの}{腫れ物}{にーぶた\tl{˩˥}}
\rentry{はれる}{晴れる}{〔天気が\ntilde{}〕ぱりん\tl{˥}}
\rentry{はれる}{腫れる}{ふくりるん\tl{˥˩}}
\rentry{はろう}{波浪}{なん\tl{˥˩}}
\rentry{はわす}{這わす}{〔釣り糸などを\ntilde{}〕ぺーすん\tl{˥}}
\rentry{ばん}{番}{ばん\tl{˩˥}}
\rentry{はんえい}{繁栄}{はんじょー\tl{˥}}
\rentry{はんえいする}{繁栄する}{さかいるん\tl{˥˩}}
\rentry{ばんく}{バンク}{すねー\tl{˥}}
\rentry{はんこうする}{反抗する}{たい\tl{˥}\ws{}すん\tl{˥˩}}
\rentry{ばんざくろ}{バンザクロ}{〔野生の\ntilde{}〕とっきん、ばんしゅる\tl{˩˥}}
\rentry{はんじょう}{繁昌}{はんじょー\tl{˥}}
\rentry{はんせん}{船舶}{ふに\tl{˥}}
\rentry{はんにえ}{半煮え}{なまにー\tl{˩˥}}
\rentry{はんにち}{半日}{ぱんにち\tl{˥˩}}
\rentry{はんばいする}{販売する}{かしみるん\tl{˥˩}}
\rentry{はんぶん}{半分}{ぱんぶん\tl{˥}}
\rentry{はんぶんわけ}{半分分け}{ぱんぶんばぎ\tl{˥}}
\rentry{はんもする}{繁茂する}{かぷん\tl{˥}}
\rentry{ひ}{日}{しな\tl{˥}、ぴん\tl{˥˩}}
\rentry{ひ}{火}{ぴー\tl{˥}}
\rentry{びーちろっく}{ビーチロック}{なんだら\tl{˩˥}}
\rentry{ぴーなつ}{ピーナツ}{じぃーまーみ\tl{˩˥}}
\rentry{ひえ}{稗}{しん\tl{˥}、ぴん\tl{˥˩}}
\rentry{ひえらび}{日選び}{ぴーる}
\rentry{ひえる}{冷える}{ぴーるん\tl{˥}、ぴしゃ\tl{˥}はん、ぴり\tl{˥}、ぴんぐるん\tl{˥}}
\rentry{ひかくする}{比較する}{くらびるん\tl{˥˩}}
\rentry{ひかげ}{日蔭}{けー}
\rentry{ひがし}{東}{あーり\tl{˥˩}}
\rentry{ひがしがわ}{東側}{あらた、ありかた\tl{˥˩}}
\rentry{ひがしずむ}{日が沈む}{しな\tl{˥}いるん\tl{˥˩}}
\rentry{ひがしむら}{東村}{ありむら}
\rentry{ひがてる}{日が照る}{しな\tl{˥}しっすん\tl{˥}}
\rentry{ひがのぼる}{日が昇る}{しな\tl{˥}あがるん\tl{˥˩}}
\rentry{ひからせる}{光らせる}{ぴからすん\tl{˥}}
\rentry{ひかり}{光る}{ぴかるん\tl{˥}}
\rentry{ひがん}{彼岸}{ぴんがん\tl{˥}}
\rentry{びき}{引き}{ぴき\tl{˥˩}}
\rentry{～ひき}{～匹}{がら、ら}
\rentry{ひきあう}{引き合う}{ひきあうん\tl{˥˩}}
\rentry{ひきあげる}{引き揚げる}{ひきあんぎるん\tl{˥˩}}
\rentry{ひきうける}{引き受ける}{うきとぅるん\tl{˥}、ひき\tl{˥˩}うきるん\tl{˥}}
\rentry{ひきうす}{引き臼}{ぴきうし\tl{˥˩}}
\rentry{ひきおろす}{引き下ろす}{ぴきうらすん、ぴきうるすん}
\rentry{ひきかえす}{引き返す}{ぴきけーすん\tl{˥˩}}
\rentry{ひきさく}{引き裂く}{ぴきさくん\tl{˥˩}}
\rentry{ひきしお}{引き潮}{ぴきすー\tl{˥˩}}
\rentry{ひきずる}{引きずる}{さふくん\tl{˥˩}、そんぐん\tl{˥˩}}
\rentry{ひきたおす}{引き倒す}{ぴきとーすん}
\rentry{ひきだす}{引き出す}{ぴきだーすん、ぴきんだすん}
\rentry{ひきたてる}{引き立てる}{むちあんぎるん\tl{˩˥}}
\rentry{ひきつける}{引き付ける}{ぴきしぃきん\tl{˥˩}}
\rentry{ひきつれる}{引き連れる}{ぴきそーるん\tl{˥˩}}
\rentry{ひきど}{引き戸}{ぴきやどぅ\tl{˩˥}}
\rentry{ひきとる}{引き取る}{ぴきとぅるん\tl{˥˩}}
\rentry{ひきなわりょう}{曳き縄漁}{ぴきなん\tl{˩˥}}
\rentry{ひきぬく}{引き抜く}{ぴき\tl{˥˩}ぬぐん\tl{˩˥}}
\rentry{ひきのばす}{引き伸ばす}{ぴき\tl{˥˩}ぬばすん\tl{˩˥}}
\rentry{ひきはなす}{引き離す}{ばがすん\tl{˩˥}}
\rentry{ひきやぶる}{引き破る}{ぴき\tl{˥˩}やぶるん\tl{˩˥}}
\rentry{ひきよせる}{引き寄せる}{ぴきゆしるん\tl{˥˩}}
\rentry{ひく}{引く}{〔潮が\ntilde{}〕ぴぃすん\tl{˥}、ぴくん\tl{˥}}
\rentry{ひく}{弾く}{ぴくん\tl{˥}}
\rentry{ひく}{挽く}{〔木を\ntilde{}〕ばぐん\tl{˥}、〔臼を\ntilde{}〕ぴくん\tl{˥˩}}
\rentry{ひくい}{低い}{まろ\tl{˩˥}はん}
\rentry{ひげ}{鬚}{ぴに\tl{˥˩}}
\rentry{ひご}{庇護}{かたが\tl{˥}}
\rentry{ひこうする}{飛行する}{とぅぷん\tl{˥˩}}
\rentry{ひざ}{膝}{すぷしん\tl{˥˩}}
\rentry{ひさし}{庇}{あまだり\tl{˥}}
\rentry{ひさしくあわない}{久しく会わない}{みーどぅー\tl{˩˥}さーん}
\rentry{ひざまずき}{}{ぺしきびり\tl{˥˩}}
\rentry{ひしゃく}{柄杓}{にぶ\tl{˩˥}}
\rentry{ひしゃげる}{}{ぴさんたらすん\tl{˥˩}}
\rentry{びじんだ}{美人だ}{あばりしゃ\tl{˥}はん}
\rentry{ひせ}{干瀬}{ぴー\tl{˥}}
\rentry{ひせのわれめ}{干瀬の割れ目}{\josho{}ぴーぬ\kako{}\ws{}ばり}
\rentry{ひたい}{額}{ふてー\tl{˥˩}}
\rentry{ひだり}{左}{ぴなり\tl{˥}}
\rentry{ひだりかた}{左方}{ぴなり\tl{˥}}
\rentry{ひたる}{浸る}{すかるん\tl{˥˩}}
\rentry{ひっかく}{引っ掻く}{はかじるん\tl{˥}、〔爪で\ntilde{}〕はかずるん\tl{˥}}
\rentry{ひっかける}{引っ掛かる}{はかるん\tl{˥}}
\rentry{ひっかぶる}{}{ぱくん\tl{˥˩}}
\rentry{ひっくりかえす}{引っくり返す}{ぐるっけーすん\tl{˩˥}}
\rentry{ひっくりかえる}{引っくり返る}{ぐるっけーらん\tl{˩˥}}
\rentry{ひづけ}{日付}{ぴんちぃ\tl{˥˩}}
\rentry{ひっこす}{引っ越す}{うつるん\tl{˥}}
\rentry{びっこになる}{}{ないぐん\tl{˩˥}}
\rentry{びっこをひく}{}{ないぐん\tl{˩˥}}
\rentry{ひっさげる}{引っ提げる}{ぴきさんぎるん\tl{˥˩}}
\rentry{ひつじ}{未}{ぴちぃ\tl{˥˩}}
\rentry{ひっしゃ}{筆者}{ぴっしゃ\tl{˥}}
\rentry{びっしょり}{}{けった\ws{}ぞっふぁら}
\rentry{ひっつかむ}{}{かみしきるん\tl{˥˩}}
\rentry{ひっつく}{ひっ付く}{だっくゎーるん\tl{˥}、むっつぁーるん\tl{˩˥}、むっつぁるん\tl{˩˥}}
\rentry{ひっつける}{ひっ付ける}{むっつぁーすん\tl{˩˥}}
\rentry{ひっつける}{引っつける}{だっくゎすん}
\rentry{ひっぱたく}{}{〔鞭で\ntilde{}〕だすますん\tl{˩˥}}
\rentry{ひっぱる}{引っ張る}{ぴきしぃきん\tl{˥˩}}
\rentry{ひつようだ}{必要だ}{いるん\tl{˥˩}}
\rentry{ひでり}{日照り}{ぺーり\tl{˥˩}}
\rentry{ひと}{人}{にんぎん\tl{˩˥}、ぴぃとぅ\tl{˥˩}}
\rentry{ひどい}{酷い}{どぅぐ}
\rentry{ひといき}{一息}{ぴぃとぅいし\tl{˥}}
\rentry{ひとうね}{一畝}{ぴぃとぅきぱ\tl{˥}}
\rentry{ひとえぐさ}{ヒトエグサ}{あーさ\tl{˥}}
\rentry{ひとえまぶた}{一重瞼}{かーみん\tl{˥}}
\rentry{ひとおじする}{人怖じする}{ぴぃとぅごはすん}
\rentry{ひときれ}{一切れ}{ぴぃとぅっし\tl{˥}}
\rentry{ひとさしゆび}{人差し指}{ぴぃささしー\tl{˥}}
\rentry{ひとしごと}{一仕事}{ぴぃとぅしぐとぅ\tl{˥}}
\rentry{ひとだのみ}{人頼み}{ぴぃとぅたなぐん}
\rentry{ひとだま}{人魂}{ぴんだま\tl{˥}}
\rentry{ひとつ}{一つ}{ぴとぅち\tl{˥}}
\rentry{ひとなみに}{人並に}{ぴぃとぅだぎ\tl{˥˩}}
\rentry{ひとにたよる}{人に頼る}{ぴぃとぅたなぐん}
\rentry{ひとばん}{一晩}{ぴぃとぅゆー\tl{˥}}
\rentry{ひとばんじゅう}{一晩中}{ゆどぅし}
\rentry{ひとみしりをする}{人見知りをする}{ぴぃとぅごはすん、ぴとぅみしり\ws{}すん}
\rentry{ひとよ}{一夜}{ぴぃとぅゆー\tl{˥}}
\rentry{ひとり}{一人}{ぴぃとぅり\tl{˥}}
\rentry{ひとり}{独り}{ぴぃとぅり\tl{˥}}
\rentry{ひとりごと}{独り言}{ぴぃとぅりふち\tl{˥}}
\rentry{ひとりっこ}{一人っ子}{ぴぃとぅりたま\tl{˥}}
\rentry{ひとりもの}{独り者}{たんがーむぬ\tl{˥}}
\rentry{ひとをつかう}{人を使う}{ぴぃとぅすこーん}
\rentry{ひなた}{日向}{しなぬ\tl{˥}\ws{}みー\tl{˥˩}}
\rentry{ひなわ}{火縄}{ぴーなん\tl{˥}}
\rentry{ひなんする}{非難する}{しみるん\tl{˥}}
\rentry{ひにち}{日にち}{ぴん\tl{˥˩}}
\rentry{ひねる}{捻る}{ぴにるん\tl{˥}、むじまーすん\tl{˩˥}}
\rentry{ひのかみ}{火の神}{ぴなかん\tl{˥}}
\rentry{ひのたま}{火の球}{ぴんだま\tl{˥}}
\rentry{ひのでになる}{日の出になる}{しな\tl{˥}あがるん\tl{˥˩}}
\rentry{ひはつもどき}{ヒハツモドキ}{ぴぱちぃ\tl{˥}}
\rentry{ひばんもり}{火番盛}{\josho{}ひばん\kako{}むり}
\rentry{ひびる}{}{しかむん\tl{˥}}
\rentry{ひま}{暇}{ぴま\tl{˥˩}}
\rentry{ひまご}{曽孫}{またまー\tl{˥˩}}
\rentry{ひまんしゃ}{肥満者}{ばたぶたー\tl{˩˥}}
\rentry{ひまんになる}{肥満になる}{ぱなたるん\tl{˥}、ぱんたるん\tl{˥}}
\rentry{ひも}{紐}{ぶー\tl{˩˥}}
\rentry{ひもじい}{}{やーは\tl{˩˥}\ws{}すん\tl{˥˩}、やー\tl{˩˥}はん}
\rentry{ひやかす}{}{ぱくるん\tl{˥}}
\rentry{ひゃく}{百}{ぴゃーぐ\tl{˥}}
\rentry{ひゃくしょう}{百姓}{ぶざ\tl{˥}}
\rentry{ひやす}{冷やす}{ぴらすん\tl{˥}}
\rentry{ひやとい}{日雇い}{ひよー\tl{˥}}
\rentry{ひやみず}{冷や水}{ぴーみじぃ\tl{˥}}
\rentry{ひやめし}{冷や飯}{ぴりふつぁりむぬ}
\rentry{ひゅうひゅう}{}{びゅーびゅー}
\rentry{びゅうびゅう}{}{ばーばー}
\rentry{びょうき}{病気}{ぱなしき\tl{˥˩}、やみ\tl{˩˥}、やん\tl{˩˥}}
\rentry{びょうきにかかる}{病気に罹る}{やむん\tl{˩˥}}
\rentry{びょうじゃくだ}{病弱だ}{びな\tl{˩˥}さーん}
\rentry{びょうじゃくなこども}{病弱な子供}{びーらー}
\rentry{ひょうじゅんご}{標準語}{やまとぅむに\tl{˩˥}}
\rentry{ひょうたん}{ヒョウタン}{かなばりん\tl{˥}}
\rentry{ひょうはくする}{漂白する}{さらすん\tl{˥˩}}
\rentry{ひょうばん}{評判}{さた\tl{˥}}
\rentry{ひょうりゅうする}{漂流する}{なーりん\tl{˩˥}}
\rentry{ひよくなはたけ}{肥沃な畑}{めーりぴてー\tl{˩˥}}
\rentry{ひより}{日和}{ぴんちぃ\tl{˥˩}}
\rentry{ひよわ}{ひ弱}{びーら}
\rentry{ひらえ}{平得}{ぴぃせー\tl{˥˩}}
\rentry{ひらきど}{開き戸}{まーしとぅ\tl{˥˩}}
\rentry{ひらける}{開ける}{ひらきるん\tl{˥}}
\rentry{ひらたい}{平たい}{ぴぃさんたり\tl{˥˩}}
\rentry{ひらたくなる}{平たくなる}{ぴさんたらすん\tl{˥˩}}
\rentry{ひらてでうつ}{平手で打つ}{ぱちみがすん\tl{˥}}
\rentry{ひらべったい}{平べったぃ}{ぴぃさんたり\tl{˥˩}}
\rentry{ひらみれもん}{ヒラミレモン}{ざがふなぶ\tl{˩˥}}
\rentry{ひりょう}{肥料}{こい\tl{˥}}
\rentry{ひる}{}{〔屁を\ntilde{}〕ぴぃすん\tl{˥}}
\rentry{ひる}{昼}{ぴる\tl{˥˩}}
\rentry{ひるご}{昼後}{ぴぃすあとぅ}
\rentry{ひるね}{昼寝}{ぴぃすまぬっふぃ\tl{˥˩}}
\rentry{ひるのうちに}{昼の内に}{ぴるぬ\ws{}うち}
\rentry{ひるま}{昼間}{ぴぃす\tl{˥˩}、ぴる\tl{˥˩}、ぴるぬ\ws{}うち}
\rentry{ひるむ}{}{たんきるん\tl{˥}}
\rentry{ひるめし}{昼飯}{ぴぃすまりむぬ\tl{˩˥}}
\rentry{ひるやすみ}{昼休み}{ぴぃすまり\tl{˥˩}やしみ\tl{˩˥}}
\rentry{ひれ}{鰭}{ぱに\tl{˥˩}}
\rentry{ひろい}{広い}{ぴぃそ\tl{˥}はん}
\rentry{ひろいもの}{拾い物}{\josho{}ぴぃしゃー\kako{}むぬ}
\rentry{ひろう}{拾う}{とぅみるん\tl{˥}、ぴぃすん\tl{˥˩}}
\rentry{ひろう}{疲労}{くたんでー\tl{˥}、ぼーがり\tl{˩˥}}
\rentry{びろう}{ビロウ}{くぱ\tl{˥˩}}
\rentry{ひろうぎみだ}{疲労ぎみだ}{だる\tl{˩˥}さーん}
\rentry{ひろうする}{疲労する}{ぶがりるん\tl{˩˥}、ぼたりるん\tl{˩˥}}
\rentry{ひろがる}{広がる}{ひるまるん\tl{˥}、ぴるまるん\tl{˥}}
\rentry{ひろげる}{広げる}{ぴそーぎるん、ひるみるん\tl{˥}、ひるみん\tl{˥}}
\rentry{ひろまる}{広まる}{ひるまるん\tl{˥}、ぴるまるん\tl{˥}}
\rentry{ひろめる}{広める}{ひるみるん\tl{˥}、ひるみん\tl{˥}}
\rentry{ひんじゃくだ}{貧弱だ}{ぴな\tl{˥}はん}
\rentry{びんしょうだ}{敏捷だ}{からばっさ\tl{˥}はん}
\rentry{びんぼう}{貧乏}{ぴんそー\tl{˥}}
\rentry{ふあん}{不安}{しわー\tl{˥˩}}
\rentry{ふいごまつり}{ふいご祭}{かつぇーぶな\tl{˥˩}}
\rentry{ふうしん}{風疹}{とーしんべー\tl{˥}}
\rentry{ふうすい}{風水}{ふんしー\tl{˥}}
\rentry{ふうすいがく}{風水学}{ふんしー\tl{˥}}
\rentry{ふうたい}{風袋}{ふーたい\tl{˥˩}}
\rentry{ふうとう}{封筒}{じょーぶくろ\tl{˥˩}}
\rentry{ふうふ}{夫婦}{とぅんぶとぅ\tl{˥˩}}
\rentry{ふうんだ}{不運だ}{うんぬ\tl{˥˩}\ws{}ねーぬ\tl{˩˥}}
\rentry{ふえ}{笛}{ぴん\tl{˥˩}}
\rentry{ふえる}{増える}{ぶーさ\tl{˥}\ws{}なるん\tl{˩˥}}
\rentry{ふか}{鱶}{さぱ\tl{˥˩}}
\rentry{ぶか}{部下}{しんか\tl{˥}}
\rentry{ふかい}{深い}{ふか\tl{˥˩}はん}
\rentry{ふかくおもう}{深く思う}{うむいくがりん\tl{˥}}
\rentry{ぶかっこう}{不格好}{やなかたち\tl{˥˩}}
\rentry{ふかむら}{冨嘉村}{ふかむら\tl{˥˩}}
\rentry{ふかれる}{吹かれる}{〔風に\ntilde{}〕ふかりるん\tl{˥}}
\rentry{ふきあれる}{吹き荒れる}{ふきありるん\tl{˥}}
\rentry{ふきでもの}{吹き出物}{にーぶた\tl{˩˥}}
\rentry{ふきとる}{拭き取る}{しっするん\tl{˥˩}}
\rentry{ふきまくる}{吹きまくる}{ふきありるん\tl{˥}}
\rentry{ぶきようだ}{不器用だ}{くぱ\tl{˥}はん、しー\tl{˥}くぱ\tl{˥}はーん}
\rentry{ふく}{拭く}{しっするん\tl{˥˩}、すするん\tl{˥˩}}
\rentry{ふく}{葺く}{〔屋根を\ntilde{}〕ふくん\tl{˥˩}}
\rentry{ふくぎ}{福木}{ふこん\tl{˥}}
\rentry{ふくつう}{腹痛}{ばた\tl{˩˥}やみ\tl{˩˥}}
\rentry{ぶくぶく}{ブクブク}{ぶくぶく}
\rentry{ふくませる}{含ませる}{ふくますん\tl{˥}}
\rentry{ふくむ}{含む}{ふくむん\tl{˥}}
\rentry{ふくめん}{覆面}{こーがき\tl{˥}}
\rentry{ふくや}{福屋}{とーら\tl{˥˩}}
\rentry{ふくらます}{膨らます}{ふくらすん\tl{˥˩}}
\rentry{ふくらませる}{膨らませる}{ふくらますん\tl{˥˩}}
\rentry{ふくらむ}{膨らむ}{ふくりん\tl{˥˩}}
\rentry{ふくれあがる}{膨れ上がる}{むりあがるん\tl{˥˩}}
\rentry{ふくれる}{膨れる}{ふくりるん\tl{˥˩}、〔水を含んで\ntilde{}〕ふくりん\tl{˥˩}}
\rentry{ふくろ}{袋}{ふくる\tl{˥}}
\rentry{ふくろう}{梟}{すくく\tl{˥}}
\rentry{ふけつだ}{不潔だ}{すぷた\tl{˥}はん、すぷ\tl{˥}つぁーん}
\rentry{ふこう}{不幸}{やなくとぅ\tl{˩˥}}
\rentry{ふごうかく}{不合格}{ふごー\tl{˥˩}}
\rentry{ふさい}{負債}{うか\tl{˥}}
\rentry{ふさがる}{塞がる}{ふさがるん\tl{˥˩}、ふつぁるん\tl{˥˩}}
\rentry{ふさぐ}{塞ぐ}{ふさぐん\tl{˥˩}}
\rentry{ふし}{節}{ふし\tl{˥}}
\rentry{ぶし}{武士}{ぶし\tl{˥˩}}
\rentry{ふしうた}{節歌}{ふしうた\tl{˥}}
\rentry{ふしぎだ}{不思議だ}{ぴるまさん、ぴるま\tl{˥}はん、みずら\tl{˩˥}はん}
\rentry{ふしゅう}{腐臭}{ふつぁりかー}
\rentry{ぶしょう}{不精}{ふゆー\tl{˥}}
\rentry{ぶしょうだ}{不精だ}{しぴんさ\tl{˥˩}はーん}
\rentry{ふしょうちだ}{不承知だ}{ふつくれー}
\rentry{ぶしょうもの}{不精者}{ふゆーむぬ\tl{˥}、ふゆな\ws{}むぬ}
\rentry{ふせぐ}{防ぐ}{ふしぐん\tl{˥}}
\rentry{ふせる}{伏せる}{うすふかすん\tl{˥˩}}
\rentry{ふそく}{不足}{ふすく\tl{˥˩}}
\rentry{ふそくする}{不足する}{たらーぬ\tl{˥˩}}
\rentry{ふた}{蓋}{ふた\tl{˥˩}}
\rentry{ふだ}{札}{ふだー\tl{˥˩}}
\rentry{ぶた}{豚}{うわー\tl{˥}}
\rentry{ふだいれ}{札入れ}{ふだいり\tl{˥˩}}
\rentry{ぶたごや}{豚小屋}{うわん\tl{˥}ひー\tl{˩˥}}
\rentry{ふたたび}{再び}{また\tl{˥˩}}
\rentry{ふたたびふっとうさせる}{再び沸騰させる}{ふつぁーすん\tl{˥}}
\rentry{ふたつ}{二つ}{ふたーち\tl{˥˩}、ふたち\tl{˥˩}}
\rentry{ふたつとも}{二つとも}{ふたーちぃ\tl{˥˩}ばぎ\tl{˩˥}}
\rentry{ふだにん}{札人}{ふだにん\tl{˥}}
\rentry{ぶたのえさいれ}{豚の餌入れ}{とーに\tl{˥}}
\rentry{ふたまた}{二又}{ふたまた\tl{˩˥}}
\rentry{ふたまた}{二股}{ふたまた\tl{˩˥}}
\rentry{ふたをする}{蓋をする}{ふーん\tl{˥˩}}
\rentry{ふだんぎ}{普段着}{ひーすぅぬ\tl{˩˥}}
\rentry{ふち}{縁}{ふちぃ\tl{˥˩}}
\rentry{ぶちあたる}{ぶち当たる}{だっさが}
\rentry{ぶちあてる}{ぶち当てる}{だっさが}
\rentry{ぶちこむ}{ぶち込む}{だくむん\tl{˩˥}、ぶくむん\tl{˥}}
\rentry{ぶちこわす}{ぶち壊す}{すくっふゎすん\tl{˥˩}}
\rentry{ぶっかける}{}{〔水を\ntilde{}〕いがふちるん\tl{˥˩}}
\rentry{ふっくらとしたさま}{ふっくらとした様}{ふくるふくる}
\rentry{ぶつぜん}{仏前}{ぷとぅぎんぬ\tl{˥}\ws{}めー\tl{˩˥}}
\rentry{ぶつだん}{仏壇}{とぅく\tl{˥˩}}
\rentry{ぶっつける}{}{うったち}
\rentry{ふっておとす}{振って落とす}{ふい\tl{˥˩}\ws{}っしん\tl{˩}}
\rentry{ふっとうさせる}{沸騰させる}{たぎらすん\tl{˥˩}}
\rentry{ふっとうする}{沸騰する}{たぎるん\tl{˥˩}、ふつん\tl{˥}}
\rentry{ふっとうするさま}{沸騰する様}{ぶとぅぶとぅ}
\rentry{ぷつりと}{}{ぷちんた}
\rentry{ぷつんと}{}{ぷちんた}
\rentry{ふで}{筆}{ふち\tl{˥˩}}
\rentry{ふとい}{太い}{みちさはん}
\rentry{ふところ}{懐}{ふちくる}
\rentry{ふとる}{太る}{ぱなたるん\tl{˥}、ぱんたるん\tl{˥}、〔芽などが\ntilde{}〕めーるん\tl{˥˩}}
\rentry{ふとん}{布団}{うず\tl{˥˩}}
\rentry{ふなうら}{船浦}{ふのーら\tl{˥˩}}
\rentry{ふなだいく}{船大工}{ふなでーぐ}
\rentry{ふなたび}{船旅}{ふなたぴ\tl{˥}}
\rentry{ふなだま}{船霊}{ふなだま\tl{˥}}
\rentry{ふなだまり}{船溜り}{ふなだまり\tl{˥}}
\rentry{ふなよい}{船酔い}{ふねー\tl{˥}}
\rentry{ふにんだ}{不妊だ}{やじまるん\tl{˥˩}}
\rentry{ふね}{船}{ふに\tl{˥}}
\rentry{ふねのていはくばしょ}{船の停泊場所}{ふなだまり\tl{˥}}
\rentry{ふねのよこいた}{船の横板}{ふなばに\tl{˥}}
\rentry{ふねのりゅうこつ}{船の竜骨}{かーら\tl{˥}、まつら\tl{˥}}
\rentry{ふはいする}{腐敗する}{ふつぁりん\tl{˥}}
\rentry{ふびじん}{不美人}{うきしゃ\tl{˥}はん}
\rentry{ふみつける}{踏みつける}{ふんしきるん\tl{˥˩}、ふんつぁーすん\tl{˥˩}}
\rentry{ふみつぶす}{踏み潰す}{うすぴらがすん\tl{˥˩}、ふんびらがすん\tl{˥˩}}
\rentry{ふみにじる}{踏みにじる}{ふんつぁーすん\tl{˥˩}}
\rentry{ふみはずす}{踏み外す}{とぅぷちん\tl{˥˩}、ふんぱんちるん\tl{˥˩}、ふんぱんつぁすん\tl{˥˩}}
\rentry{ふむ}{踏む}{〔足で\ntilde{}〕ふむん}
\rentry{ふやす}{増やす}{ぶーさ\tl{˥}\ws{}なすん\tl{˩˥}}
\rentry{ふゆ}{冬}{ふっひ\tl{˥˩}}
\rentry{ふゆう}{富裕}{うやぎ\tl{˥}}
\rentry{ふゆうしゃ}{富裕者}{\josho{}うやぎ\kako{}ひー}
\rentry{ふゆうする}{浮遊する}{おーがるん\tl{˥˩}}
\rentry{ぶよう}{舞踊}{ぶどぅり\tl{˥˩}、ぶんどぅり}
\rentry{ぶらさがる}{ぶら下がる}{さがるん\tl{˥}、さんがるん\tl{˥}}
\rentry{ぶらさげる}{ぶら下げる}{すさーらすん\tl{˥}、ぴき\tl{˥˩}さんぎん\tl{˩˥}}
\rentry{ふらふらっと}{}{ふらふら}
\rentry{ふりあてる}{振り当てる}{ばりあてぃん\tl{˥˩}、わりあてぃん\tl{˥}}
\rentry{ふりかえる}{振り返る}{とぅんけーるん}
\rentry{ふりかける}{振りかける}{ふりはきん\tl{˥˩}}
\rentry{ふりすてる}{振り捨てる}{ふっひ\ws{}ししるん}
\rentry{ふりはらう}{振り払う}{ふっひ\ws{}ししるん}
\rentry{ふりまわす}{振り回す}{ふりまーすん\tl{˥˩}}
\rentry{ふりむく}{振り向く}{とぅんけーるん}
\rentry{ふる}{振る}{ふっふん\tl{˥˩}}
\rentry{ふる}{降る}{〔雨が\ntilde{}〕ふっふん\tl{˥}}
\rentry{ふるい}{篩}{しのー\tl{˥}}
\rentry{ふるう}{振るう}{ゆらすん\tl{˥˩}}
\rentry{ふるぎ}{古着}{ふーすぅぬ\tl{˥}}
\rentry{ふるくする}{古くする}{ふー\tl{˥}\ws{}なすん\tl{˩˥}}
\rentry{ふるくなる}{古くなる}{ふー\tl{˥}\ws{}なるん\tl{˩˥}}
\rentry{ふるさと}{古里}{まりじま\tl{˥˩}}
\rentry{ふるびる}{古びる}{ふー\tl{˥}\ws{}なるん\tl{˩˥}}
\rentry{ぶるぶる}{}{ふとぅふとぅ}
\rentry{ふるむ}{古む}{ふー\tl{˥}\ws{}なるん\tl{˩˥}}
\rentry{ぶれい}{無礼}{ぶりー\tl{˩˥}}
\rentry{ふれる}{触れる}{さーるん\tl{˥˩}}
\rentry{ふろ}{風呂}{ゆふる\tl{˩˥}}
\rentry{ふろしき}{風呂敷}{うしぴ\tl{˥}}
\rentry{ふんがいする}{憤慨する}{くさむくん\tl{˥}}
\rentry{ぶんけ}{分家}{きにばがら\tl{˩˥}}
\rentry{ぶんけする}{分家する}{きにばがりるん\tl{˩˥}}
\rentry{ふんじばる}{ふん縛る}{ふんさまるん\tl{˥˩}}
\rentry{ぶんそん}{分村}{むらばぎん\tl{˩˥}}
\rentry{ふんどう}{分銅}{ぴきぬ\ws{}ふり}
\rentry{ふんどし}{褌}{さねー\tl{˥˩}}
\rentry{ぶんなぐる}{ぶん殴る}{ばみがすん\tl{˥˩}}
\rentry{ぶんぱいする}{分配する}{くぱるん\tl{˥}}
\rentry{ぶんべつ}{分別}{じんぶん\tl{˩˥}}
\rentry{ぶんべつする}{分別する}{ばぎん\tl{˩˥}}
\rentry{ふんまつ}{粉末}{く\tl{˥}、ふく\tl{˥}}
\rentry{ふんまつにする}{粉末にする}{ふくなすん\tl{˥}}
\rentry{ふんわりする}{}{ふくら\tl{˥˩}さーん、ふくるさ\tl{˥}はん}
\rentry{へ}{屁}{ぴん\tl{˥}}
\rentry{べいじゅ}{米寿}{とーかち\tl{˥}}
\rentry{へいたんだ}{平坦だ}{なだらが\tl{˥˩}はん}
\rentry{へいたんなしま}{平坦な島}{ぴぃさじま\tl{˥˩}}
\rentry{へいたんにする}{平坦にする}{なだらがすん\tl{˥˩}}
\rentry{へいたんになる}{平坦になる}{なだらぐん\tl{˥˩}}
\rentry{べいはん}{米飯}{めーぬ\tl{˥˩}\ws{}いー\tl{˥}}
\rentry{へいみん}{平民}{ぶざ\tl{˥}}
\rentry{へぇー}{}{べー}
\rentry{へくさい}{屁臭い}{ぴんふつぁ\tl{˥}はーん}
\rentry{へこむ}{}{とぅまるん\tl{˥˩}}
\rentry{へしこむ}{へし込む}{うしくむん\tl{˥˩}}
\rentry{へそ}{臍}{ぷつ\tl{˥˩}}
\rentry{へただ}{下手だ}{くぱ\tl{˥}はん}
\rentry{へだてる}{隔てる}{ひだみるん\tl{˥}、ひだみん\tl{˥}}
\rentry{へたなのうふ}{下手な農夫}{ぴらすか\tl{˥}}
\rentry{べたぼれ}{べた惚れ}{まーぷり\tl{˥˩}}
\rentry{へちま}{糸瓜}{なーびら\tl{˩˥}}
\rentry{べつ}{別}{びちぃ\tl{˩˥}}
\rentry{べったりすわる}{べったり座る}{しきだーすん\tl{˥˩}}
\rentry{べつに}{別に}{ふかな}
\rentry{べつの}{別の}{びちぃぬ\tl{˩˥}}
\rentry{べつべつ}{別々}{びちぃびちぃ}
\rentry{へび}{蛇}{ぱく\tl{˥˩}}
\rentry{へら}{箆}{ぴら\tl{˥}}
\rentry{べら}{ベラ}{ふつぁび\tl{˥}}
\rentry{へる}{減る}{ぴなるん\tl{˥˩}}
\rentry{へんかする}{変化する}{かわるん\tl{˥˩}}
\rentry{べんきょう}{勉強}{むぬなれー\tl{˩˥}}
\rentry{へんじ}{返事}{ひんとー\tl{˥˩}}
\rentry{べんじょ}{便所}{ふーる\tl{˥}、ふるや\tl{˥}}
\rentry{べんしょうさせる}{弁償させる}{はかすん\tl{˥˩}}
\rentry{べんしょうする}{弁償する}{ぱくん\tl{˥˩}}
\rentry{へんずつう}{片頭痛}{かたすぷる\tl{˥}やみ\tl{˩˥}}
\rentry{へんだ}{変だ}{いふなー\tl{˥˩}、いふなー\ws{}やっさー\tl{˥˩}}
\rentry{へんたんにする}{平坦にする}{なだぎるん\tl{˥˩}}
\rentry{へんとう}{返答}{ひんとー\tl{˥˩}}
\rentry{べんとう}{弁当}{びんとー\tl{˥˩}}
\rentry{へんなぐあいだ}{変な具合だ}{〔胸のあたりが何となく\ntilde{}〕いがばりゃん}
\rentry{ほ}{帆}{ぽー\tl{˥˩}}
\rentry{ほ}{穂}{ぷー\tl{˥}}
\rentry{ぼう}{棒}{ぼー\tl{˩˥}}
\rentry{～ほう}{～方}{た}
\rentry{ぼうがい}{妨害}{わしく\tl{˩˥}}
\rentry{ぼうかんぎ}{防寒着}{ふくた\tl{˥}}
\rentry{ほうき}{箒}{ぽーち\tl{˥}}
\rentry{ほうきぼし}{ほうき星}{ぽーぎぶしぃ\tl{˥}}
\rentry{ほうげん}{方言}{すま\tl{˥}むに\tl{˩˥}}
\rentry{ぼうげん}{暴言}{あら\tl{˥˩}むに\tl{˩˥}}
\rentry{ほうこう}{奉公}{ふーく\tl{˥}}
\rentry{ほうじ}{法事}{こっこー\tl{˥}、しょっこー\tl{˥}}
\rentry{ぼうしゅ}{芒種}{ぼーす\tl{˩˥}}
\rentry{ほうしゅう}{報酬}{しぃま\tl{˥}}
\rentry{ぼうじゅつ}{棒術}{ぼー\tl{˩˥}}
\rentry{ほうじょうでへいわのよ}{豊穣で平和の世}{みるくゆー\tl{˥˩}}
\rentry{ほうじょうのよ}{豊穣の世}{かんぬ\tl{˥}\ws{}ゆー\tl{˥˩}、むがしぃ\tl{˩˥}ゆー\tl{˥˩}}
\rentry{ぼうず}{坊主}{ぼーず\tl{˩˥}}
\rentry{ほうちする}{放置する}{だすたすくん\tl{˩˥}}
\rentry{ほうちょう}{包丁}{かたな\tl{˥}、ぽっつぁ\tl{˥}}
\rentry{ほうっておく}{放っておく}{だすたすくん\tl{˩˥}}
\rentry{ほうねん}{豊年}{のーり\tl{˩˥}ゆ\tl{˥˩}、ゆがふ\tl{˥˩}}
\rentry{ほうねんさい}{豊年祭}{ぷーりん\tl{˥}}
\rentry{ほうねんさいじのかみつかさたち}{豊年祭時の神司達}{みなすけぬ\ws{}ぱんだー}
\rentry{ぼうふう}{暴風}{かちふき\tl{˥˩}}
\rentry{ほうほう}{方法}{しかた\tl{˥˩}}
\rentry{ぼうぼう}{}{ばーばー}
\rentry{ほうよう}{法要}{こっこー\tl{˥}、しょっこー\tl{˥}}
\rentry{ほうりこむ}{放り込む}{うちくむん\tl{˥}、だくむん\tl{˩˥}、たたっくむん\tl{˥}、ぶくむん\tl{˥}}
\rentry{ほうりなげる}{放り投げる}{ぼんぎるん\tl{˩˥}}
\rentry{ほおかぶり}{頬かぶり}{こーがき\tl{˥}}
\rentry{ほか}{他}{ふか\tl{˥}}
\rentry{ほかに}{外に}{ふかな}
\rentry{ほかの}{他の}{びちぃぬ\tl{˩˥}}
\rentry{ほかんする}{保管する}{\josho{}あためー\kako{}すん、かちみるん\tl{˥}}
\rentry{ぽきっと}{}{ぷちんた}
\rentry{ぽきりと}{}{ぷちんた}
\rentry{ぼくじゅう}{墨汁}{しん\tl{˥}}
\rentry{ほくとしちせい}{北斗七星}{にしななち\tl{˩˥}}
\rentry{ほくとせい}{北斗星}{にしななち\tl{˩˥}}
\rentry{ほこり}{埃}{ふくじぃ\tl{˥}}
\rentry{ほし}{星}{ぷち\tl{˥˩}}
\rentry{ほしい}{欲しい}{ふつぁはん}
\rentry{ほしがる}{欲しがる}{ふつぁはん}
\rentry{ほす}{干す}{ぷつん\tl{˥}}
\rentry{ほせん}{帆船}{ぽーしん\tl{˥˩}}
\rentry{ほそい}{細い}{ばちさはん}
\rentry{ぽたぽた}{}{すとぅりすとぅり}
\rentry{ほたる}{蛍}{しっすぴかり\tl{˥}}
\rentry{ほっきょくせい}{北極星}{にんばぷち\tl{˩˥}}
\rentry{ほっする}{欲する}{ぬずむん}
\rentry{ほったてごや}{掘っ立て小屋}{\josho{}あなばり\kako{}ひー}
\rentry{ほっとする}{}{〔心配事がなくなり\ntilde{}〕うなーぐ\ws{}しゃーん、\josho{}うなー\kako{}ぐ\ws{}なるん\tl{˩˥}}
\rentry{ほづな}{帆綱}{しーなん\tl{˥}}
\rentry{ほっぽう}{北方}{にし\tl{˥˩}、にした\tl{˥˩}}
\rentry{ほどく}{解く}{ぷとぅぐん\tl{˥}}
\rentry{ほとけ}{仏}{ぷとぅぎん\tl{˥}}
\rentry{ほとけのまえ}{仏の前}{ぷとぅぎんぬ\tl{˥}\ws{}めー\tl{˩˥}}
\rentry{ほどける}{解ける}{ぱんちるん\tl{˥˩}、ぷとぅぎるん}
\rentry{ほね}{骨}{ぷに\tl{˥}}
\rentry{ほねおしみ}{骨惜しみ}{どぅー\tl{˩˥}たんき\tl{˥}}
\rentry{ほねやすめ}{骨休め}{ぷにやしみ\tl{˥}}
\rentry{ほほえむ}{微笑む}{ばーりゃん\tl{˥˩}}
\rentry{ほめあげる}{褒めあげる}{ふみあぎるん\tl{˥}}
\rentry{ほめたたえる}{褒め讃える}{ふみあぎるん\tl{˥}}
\rentry{ほめる}{褒める}{ふみるん\tl{˥}}
\rentry{ほら}{法螺}{〔楽器の\ntilde{}〕ぶら}
\rentry{ほらがい}{法螺貝}{さぶら\tl{˥}、ぶら\tl{˩˥}}
\rentry{ほりとる}{掘り取る}{こーすん\tl{˥˩}}
\rentry{ほる}{彫る}{ぷるん\tl{˥}}
\rentry{ほる}{掘る}{ぷるん\tl{˥}}
\rentry{ほれる}{惚れる}{ぷりるん\tl{˥˩}}
\rentry{ぼろ}{}{やるはく\tl{˩˥}}
\rentry{ぼろきれ}{ぼろ切れ}{やるはく\tl{˩˥}}
\rentry{ぼろのきもの}{ぼろの着物}{ふくた\tl{˥}、やりすぅぬ\tl{˩˥}}
\rentry{ほろぶ}{滅ぶ}{とーりるん\tl{˥}、ふるぶん\tl{˥˩}}
\rentry{ほろぼす}{滅ぼす}{ふるばすん\tl{˥˩}}
\rentry{ほろよいになる}{ほろ酔いになる}{さーふーふー\ws{}すん\tl{˥˩}}
\rentry{ほん}{本}{すむち\tl{˥˩}}
\rentry{ほん～}{本～}{まー}
\rentry{ほんき}{本気}{まーすむ\tl{˥˩}}
\rentry{ほんけ}{本家}{やーむとぅ\tl{˩˥}}
\rentry{ほんけんちく}{本建築}{ぬきひー}
\rentry{ほんじつ}{本日}{きゅー\tl{˥}}
\rentry{ほんしん}{本心}{まーすむ\tl{˥˩}}
\rentry{ほんとう}{本当}{ふんとー\tl{˥}}
\rentry{ほんとうだ}{本当だ}{ましん\tl{˥˩}}
\rentry{ぼんのしょにち}{盆の初日}{むげーぴん\tl{˥˩}}
\rentry{ほんむすび}{本結び}{まーむしぃぴ\tl{˥˩}}
\rentry{ほんもの}{本物}{まーむぬ\tl{˥˩}}
\rentry{ぼんやりする}{}{うかっと\ws{}すん\tl{˥˩}、ぷ\josho{}り\kako{}ぷり\ws{}すん\tl{˥˩}}
\rentry{ま}{間}{まー\tl{˥˩}}
\rentry{まーらんせん}{馬艦船}{まー\josho{}らん\kako{}ぶに}
\rentry{まいとし}{毎年}{にんにん、めーどぅし\tl{˩˥}}
\rentry{まいにち}{毎日}{めー\josho{}が\kako{}めーにち、めーにち\tl{˩˥}}
\rentry{まいる}{参る}{うがんちゅむん\tl{˥}}
\rentry{まう}{舞う}{ぶんどぅるん\tl{˥˩}}
\rentry{まうしろ}{真後ろ}{まとむ\tl{˥˩}}
\rentry{まえ}{前}{めー\tl{˩˥}}
\rentry{まえあし}{前足}{めーぱん\tl{˩˥}}
\rentry{まえがり}{前借}{めーがり\tl{˩˥}}
\rentry{まえさと}{真栄里}{みざとぅ\tl{˩˥}}
\rentry{まえにおく}{前に置く}{めー\tl{˩˥}\ws{}なすん\tl{˩˥}}
\rentry{まえにする}{前にする}{めー\tl{˩˥}\ws{}なすん\tl{˩˥}}
\rentry{まえば}{前歯}{めーぱん\tl{˩˥}}
\rentry{まえむら}{前村}{なー\tl{˩˥}むら\tl{˥˩}}
\rentry{まえもって}{前もって}{めーむち}
\rentry{まえをいく}{前を行く}{めー\tl{˩˥}\ws{}なるん\tl{˩˥}}
\rentry{まかす}{負かす}{まかすん\tl{˥˩}}
\rentry{まかせる}{任せる}{まかすん\tl{˥˩}}
\rentry{まがたま}{勾玉}{がーら\tl{˩˥}だま\tl{˥˩}}
\rentry{まがりくねって}{曲がりくねって}{まんかりかんかり}
\rentry{まがる}{曲がる}{まんかるん\tl{˥˩}}
\rentry{まきおどり}{巻き踊り}{ますぶどぅり\tl{˥˩}}
\rentry{まきちらかす}{蒔き散らかす}{ぽーりすくん}
\rentry{まきちらす}{撒き散らす}{まぎぽーるん\tl{˩˥}}
\rentry{まきつける}{巻き付ける}{からまぐん\tl{˥}}
\rentry{まぎり}{間切}{まぎり\tl{˥˩}}
\rentry{まく}{巻く}{まーぐん\tl{˥˩}}
\rentry{まく}{幕}{まく\tl{˩˥}}
\rentry{まく}{播く}{まーぐん\tl{˥˩}}
\rentry{まく}{蒔く}{まーぐん\tl{˥˩}}
\rentry{まくら}{枕}{まっふぁ\tl{˩˥}}
\rentry{まくる}{捲る}{からぎるん\tl{˥}}
\rentry{まくれる}{捲れる}{ぱんきん\tl{˥}}
\rentry{まぐろ}{鮪}{しび\tl{˥˩}}
\rentry{まけ}{}{まぎ\tl{˥˩}}
\rentry{まける}{負ける}{まきるん\tl{˥˩}}
\rentry{まげる}{曲げる}{まんきるん\tl{˥˩}}
\rentry{まご}{孫}{まー\tl{˥}}
\rentry{まごころ}{真心}{まぐくる\tl{˥˩}}
\rentry{まごつく}{}{ざまどぅるん\tl{˩˥}}
\rentry{まこと}{誠}{まくとぅ\tl{˩˥}}
\rentry{まさか}{}{がーん}
\rentry{まさる}{勝る}{すぐりん\tl{˥˩}、まさるん\tl{˥˩}}
\rentry{まざる}{混ざる}{まんざるん\tl{˩˥}}
\rentry{まし}{}{〔～より\ntilde{}〕まーし\tl{˩˥}}
\rentry{ましょうめん}{真正面}{まそーみ\tl{˥˩}}
\rentry{ます}{増す}{ぶーさ\tl{˥}\ws{}なすん\tl{˩˥}}
\rentry{ます}{枡}{しゃー\tl{˥}}
\rentry{まずい}{不味い}{にしゃ\tl{˥˩}はん}
\rentry{まずいこと}{}{ざーふぇー\tl{˩˥}}
\rentry{まぜあわせる}{混ぜ合わせる}{あいるん\tl{˥}、まんざーすん}
\rentry{まぜる}{混ぜる}{あいるん\tl{˥}、まんじるん\tl{˩˥}}
\rentry{また}{又}{また\tl{˥˩}}
\rentry{また}{叉}{また\tl{˥˩}}
\rentry{また}{股}{また\tl{˥˩}}
\rentry{またいとこ}{又従兄弟}{またいちふ\tl{˥˩}}
\rentry{またがる}{跨る}{またんがるん\tl{˩˥}}
\rentry{またした}{股下}{またびし\tl{˩˥}}
\rentry{まだらかになる}{なだらかになる}{なだらぐん\tl{˥˩}}
\rentry{まちがい}{間違い}{ばっぺー\tl{˩˥}、まちげー\tl{˩˥}}
\rentry{まちがう}{間違う}{ばっぺー\ws{}すん}
\rentry{まちがえる}{間違える}{まちがいん\tl{˥}、まちげーるん\tl{˥}}
\rentry{まぢかにせまる}{間近に迫る}{めー\tl{˩˥}ぬすかるん\tl{˥˩}}
\rentry{まちこがれる}{待ち焦がれる}{まちくがりん\tl{˩˥}}
\rentry{まちどおしくおもう}{待ち遠しく思う}{まちぼさーん}
\rentry{まちわびる}{待ちわびる}{まちぼさーん}
\rentry{まつ}{待つ}{まつん\tl{˩˥}}
\rentry{まつ}{松}{まーち\tl{˩˥}}
\rentry{まっくら}{真っ暗}{やみ\tl{˥˩}}
\rentry{まつげ}{睫毛}{まちゅ\tl{˩˥}}
\rentry{まっさかり}{真っ盛り}{ばんじぃん\tl{˥˩}}
\rentry{まっすぐ}{真っ直ぐ}{まんが\tl{˩˥}}
\rentry{まっすぐにする}{真っ直ぐにする}{たみらすん、たみるん\tl{˥}}
\rentry{まっすぐになる}{真っ直ぐになる}{〔曲がっている物が\ntilde{}〕たまるん\tl{˥}}
\rentry{まったく}{全く}{ぬ\josho{}ー\kako{}ん、むっとぅ\tl{˩˥}}
\rentry{まったくりこうでない}{全く利口でない}{まーぱからさー\ws{}ねーぬ}
\rentry{まっち}{燐寸}{しきだぎ\tl{˥}}
\rentry{まつのねのしん}{松の根の芯}{あがしぃ\tl{˥˩}}
\rentry{まつり}{祭}{ぶなが\tl{˥}、まちり\tl{˥˩}}
\rentry{～まで}{}{ばぎ}
\rentry{まとまる}{}{まとぅまるん\tl{˥˩}}
\rentry{まとめる}{}{まとぅみん\tl{˥˩}}
\rentry{まとも}{}{まとむ\tl{˥˩}}
\rentry{まどわす}{惑わす}{ざまどぅらすん\tl{˩˥}}
\rentry{まないた}{まな板}{まんつぁ\tl{˩˥}}
\rentry{まなぶ}{学ぶ}{ならすん\tl{˩˥}}
\rentry{まにあう}{間に合う}{まにあうん}
\rentry{まにあわせる}{間に合わせる}{〔時間に\ntilde{}〕まにあーすん\tl{˥}}
\rentry{まぬがれさせる}{免れさせる}{ぬがーらすん}
\rentry{まぬがれる}{免れる}{ぬがーるん\tl{˩˥}}
\rentry{まぬけ}{間抜け}{ほーらきし\tl{˥}}
\rentry{まね}{真似}{まーべー\tl{˥˩}}
\rentry{まねする}{真似する}{まーべー\tl{˥˩}\ws{}すん\tl{˥˩}}
\rentry{まねる}{真似る}{まーべー\tl{˥˩}\ws{}すん\tl{˥˩}}
\rentry{まひる}{真昼}{ぴぃすまり\tl{˥˩}、まーびる\tl{˥˩}}
\rentry{まぶしい}{眩しい}{ぱちま\tl{˥˩}はん、ぱちまはん}
\rentry{まぶた}{瞼}{みんぬ\tl{˩˥}\ws{}かー\tl{˥}}
\rentry{まみず}{真水}{あまみじ\tl{˥˩}}
\rentry{まむかい}{真向かい}{たんか\tl{˥˩}}
\rentry{まむすび}{真結び}{まーむしぃぴ\tl{˥˩}}
\rentry{まめ}{豆}{まーみ\tl{˩˥}}
\rentry{まもの}{魔物}{まざむん\tl{˥˩}、やなむん\tl{˥˩}}
\rentry{まもる}{守る}{まむるん\tl{˥˩}}
\rentry{まゆ}{眉}{まちゅ\tl{˩˥}}
\rentry{まゆ}{繭}{まゆ\tl{˩˥}}
\rentry{まよう}{迷う}{ざまどぅるん\tl{˩˥}、すまどぅるん\tl{˥˩}}
\rentry{まよけのむすび}{魔除けの結び}{さん\tl{˥}}
\rentry{まよなか}{真夜中}{ゆなが\tl{˩˥}}
\rentry{まよわす}{迷わす}{ざまどぅらすん\tl{˩˥}}
\rentry{まらりあ}{マラリア}{ぷーき\tl{˥˩}、やきー\tl{˩˥}}
\rentry{まるい}{丸い}{む\josho{}る\kako{}むるし}
\rentry{まるっこい}{丸っこい}{む\josho{}る\kako{}むるし}
\rentry{まるでおろかだ}{まるで愚かだ}{まーぱからさー\ws{}ねーぬ}
\rentry{まるはだか}{丸裸}{まるばい\tl{˥˩}}
\rentry{まろぶ}{}{まるぶん\tl{˥˩}}
\rentry{まわす}{回す}{まーすん\tl{˥˩}、みぐらすん\tl{˥˩}}
\rentry{まわり}{}{まーましぃ\tl{˥˩}}
\rentry{まわり}{周り}{まーり\tl{˩˥}}
\rentry{まわる}{回る}{まーるん\tl{˥˩}、みぐるん\tl{˥˩}}
\rentry{まん}{万}{まん\tl{˥˩}}
\rentry{まんいち}{万一}{まんいち\tl{˥˩}}
\rentry{まんいっさいのたんじょうび}{満一歳の誕生日}{たんかー\tl{˥}}
\rentry{まんげつ}{満月}{ぶーしけん}
\rentry{まんさい}{満載}{まんしん\tl{˥˩}}
\rentry{まんた}{マンタ}{かまんた\tl{˥}}
\rentry{まんちょうになる}{満潮になる}{すー\tl{˥}んつん\tl{˩˥}}
\rentry{まんなか}{真ん中}{ま\josho{}んなー\kako{}が}
\rentry{まんぱいになる}{満杯になる}{まんしん\ws{}なるん}
\rentry{まんぷくになる}{満腹になる}{ばた\tl{˩˥}んち\tl{˩˥}}
\rentry{まんまと}{}{みしーみし}
\rentry{み}{実}{なーり\tl{˩˥}}
\rentry{み}{巳}{みー\tl{˥}}
\rentry{み}{箕}{そーぎ\tl{˥}、〔選別用の\ntilde{}〕ゆらし\tl{˥˩}}
\rentry{みあやまり}{見誤り}{みーばっぺー\tl{˩˥}}
\rentry{みおとす}{見落とす}{みーうとすん、みぬがすん\tl{˩˥}}
\rentry{みおぼえがない}{見覚えがない}{みり\tl{˩˥}\ws{}みらぬ\tl{˩˥}}
\rentry{みがく}{磨く}{とーぴくん\tl{˥}}
\rentry{みかたする}{味方する}{かたすん\tl{˥˩}}
\rentry{みがなる}{実がなる}{なーるん\tl{˩˥}}
\rentry{みがまえる}{身構える}{かまいるん\tl{˥}}
\rentry{みがるだ}{身軽だ}{からばっさ\tl{˥}はん}
\rentry{みかん}{蜜柑}{ふなぶ\tl{˥}}
\rentry{みき}{幹}{むとぅ\tl{˩˥}}
\rentry{みき}{神酒}{みし\tl{˥˩}}
\rentry{みぎ}{右}{ねーり\tl{˥˩}}
\rentry{みぎがわ}{右側}{ねーり\tl{˥˩}}
\rentry{みぐるしい}{見苦しい}{うかさ\tl{˥}はーん、みーぐりしゃーん、みぐりしゃ\tl{˩˥}はん、みっと\tl{˩˥}はん}
\rentry{みごと}{見事}{みぐとぅ\tl{˩˥}}
\rentry{みさき}{岬}{さき\tl{˥˩}、ぱな\tl{˥˩}}
\rentry{みじかい}{短い}{まろ\tl{˩˥}はん}
\rentry{みじめなおもいがする}{みじめな思いがする}{すむ\tl{˥}ぐりしゃ\tl{˥}はーん}
\rentry{みず}{水}{みじぃ\tl{˥˩}}
\rentry{みずあたり}{水当たり}{みじぃがーり\tl{˩˥}}
\rentry{みずあび}{水浴}{みずあみ}
\rentry{みずおけ}{水桶}{うき\tl{˥}、たんぐ\tl{˥}}
\rentry{みすかされる}{見透かされる}{みーふかさりん\tl{˩˥}}
\rentry{みすかす}{見透かす}{みーふかすん\tl{˩˥}}
\rentry{みずがめ}{水甕}{みじぃがーみ\tl{˩˥}}
\rentry{みずがわく}{湧く}{〔水が\ntilde{}〕ばぐん\tl{˥˩}}
\rentry{みずがんぴ}{ミズガンピ}{すぱんこっち}
\rentry{みすてる}{見捨てる}{みー\tl{˩˥}\ws{}しちん\tl{˥˩}}
\rentry{みずのこう}{水の香}{みずぬ\ws{}こー}
\rentry{みずぼうそう}{水疱瘡}{みじぃがーさ\tl{˩˥}}
\rentry{みすみす}{}{みしーみし}
\rentry{みずをきる}{水を切る}{すたらすん\tl{˥}}
\rentry{みせ}{店}{まちやー\tl{˥˩}}
\rentry{みせもの}{見せ物}{みーむぬ\tl{˩˥}}
\rentry{みせる}{見せる}{みしるん\tl{˩˥}}
\rentry{みそ}{味噌}{みしゅ\tl{˩˥}}
\rentry{みぞ}{溝}{みぞーり\tl{˥˩}}
\rentry{みそこなう}{見損なう}{みまちごーん\tl{˩˥}}
\rentry{みだしなみ}{身だしなみ}{すがい\tl{˥}}
\rentry{みたす}{満たす}{んつぁすん\tl{˩˥}}
\rentry{みたてる}{見たてる}{みたてぃるん\tl{˥}}
\rentry{みだらだ}{淫らだ}{さんごなー\tl{˥}\ws{}すん\tl{˥˩}}
\rentry{みち}{道}{みちぃ\tl{˥˩}}
\rentry{みちぶしん}{道普請}{みちぃくせー\tl{˥˩}}
\rentry{みちる}{満ちる}{んちん\tl{˩˥}}
\rentry{みつかる}{見つかる}{めすかるん\tl{˩˥}}
\rentry{みつける}{見つける}{とぅみるん\tl{˥}、とぅみん\tl{˥}、みしきるん\tl{˩˥}}
\rentry{みっしゅうしている}{密集している}{かた\tl{˥˩}はん}
\rentry{みっつ}{三つ}{みーち\tl{˥˩}}
\rentry{みっともない}{}{みぐりしゃ\tl{˩˥}はん、みっと\tl{˩˥}はん}
\rentry{みてゆるす}{見て許す}{みのーすん\tl{˩˥}}
\rentry{みとおす}{見通す}{みーとーすん\tl{˩˥}、みーふかすん\tl{˩˥}}
\rentry{みとどける}{見届ける}{みーとぅどぅきるん\tl{˩˥}、みとぅどぅぎるん\tl{˩˥}}
\rentry{みとめられる}{認められる}{とーるん\tl{˥}}
\rentry{みとめる}{認める}{みとぅみるん\tl{˩˥}}
\rentry{みどりいろに}{緑色に}{お\josho{}ー\kako{}し}
\rentry{みとれる}{見とれる}{みーぶりん\tl{˩˥}}
\rentry{みなおす}{見直す}{みのーすん\tl{˩˥}}
\rentry{みなさま}{皆さま}{けーら\tl{˥}ねーら\tl{˩˥}}
\rentry{みなみ}{南}{ぺー\tl{˥}}
\rentry{みなみかぜ}{南風}{ぺーかち\tl{˥}}
\rentry{みなみなさま}{皆皆様}{けーら\tl{˥}ねーら\tl{˩˥}}
\rentry{みなみむら}{南村}{ぺーむら\tl{˥˩}}
\rentry{みならい}{見習い}{みなれー\tl{˩˥}}
\rentry{みならう}{見習う}{みーならすん\tl{˩˥}}
\rentry{みなり}{身なり}{すがい\tl{˥}}
\rentry{みなれる}{見慣れる}{みーなりゃん}
\rentry{みにくい}{見にくい}{みーぐりしゃーん}
\rentry{みにくい}{醜い}{うかさ\tl{˥}はーん}
\rentry{みにこたえる}{身に応える}{あたるん\tl{˥˩}}
\rentry{みの}{蓑}{めんさ\tl{˩˥}}
\rentry{みのがす}{見逃す}{みーぬがすん\tl{˩˥}、みぬがすん\tl{˩˥}}
\rentry{みのり}{実り}{のーり\tl{˩˥}}
\rentry{みのる}{実る}{のーるん、めーるん\tl{˥˩}}
\rentry{みはりする}{見張りする}{ばん\tl{˩˥}}
\rentry{みぶるいする}{身震いする}{ふぃーつぁるん}
\rentry{みほれる}{見惚れる}{みーぶりん\tl{˩˥}}
\rentry{みまちがい}{見間違い}{みーばっぺー\tl{˩˥}}
\rentry{みまちがえる}{見間違える}{みまちごーん\tl{˩˥}}
\rentry{みまわる}{見回る}{まーり\tl{˥˩}\ws{}みるん\tl{˩˥}}
\rentry{みみ}{耳}{みん\tl{˩˥}}
\rentry{みみがとおい}{耳が遠い}{みんとぅーさーん}
\rentry{みみず}{ミミズ}{みみじぃ\tl{˩˥}}
\rentry{みみずく}{ミミズク}{ましかく\tl{˥˩}}
\rentry{みみたぶ}{耳たぶ}{みしくるみん\tl{˩˥}}
\rentry{みもの}{見物}{みーむぬ\tl{˩˥}}
\rentry{みゃく}{脈}{みゃぐ\tl{˩˥}}
\rentry{みやこ}{宮古}{みゃーぐ\tl{˥˩}}
\rentry{みやすい}{見やすい}{みーやっ\tl{˩˥}さーん}
\rentry{みやら}{宮良}{めーら\tl{˥˩}}
\rentry{みょうだい}{名代}{みょーでん\tl{˥˩}}
\rentry{みょうばん}{明晩}{あつぁゆ\tl{˥}}
\rentry{みる}{見る}{みるん\tl{˩˥}}
\rentry{みるにしのびない}{見るに忍びない}{すむやむーん}
\rentry{みろくがみ}{弥勒神}{みるく\tl{˥˩}}
\rentry{みろくぼさつ}{弥勒菩薩}{みるく\tl{˥˩}}
\rentry{みをつける}{目を付ける}{みーしきるん}
\rentry{みんな}{皆}{\josho{}がす\kako{}た、けーら、むーる}
\rentry{みんよう}{民謡}{うた\tl{˥˩}}
\rentry{むいに}{無為に}{あだり\tl{˥˩}}
\rentry{むかいかぜ}{向かい風}{むげーかち\tl{˥˩}}
\rentry{むかう}{向かう}{むぎゃん\tl{˥˩}}
\rentry{むかえのひ}{迎えの日}{むげーぴん\tl{˥˩}}
\rentry{むかえる}{迎える}{むげーるん\tl{˥˩}}
\rentry{むかし}{昔}{むがし\tl{˩˥}}
\rentry{むかしから}{昔から}{むがしぃがら\tl{˩˥}}
\rentry{むかしのひと}{昔の人}{むがしぃ\tl{˩˥}ぴぃとぅ\tl{˥˩}}
\rentry{むかしのふなまちごや}{昔の船待ち小屋}{いなさ\tl{˥}ひー\tl{˩˥}}
\rentry{むかしのよ}{昔の世}{むがしぃ\tl{˩˥}ゆー\tl{˥˩}}
\rentry{むかしばなし}{昔話}{むがしぃ\tl{˩˥}ぱなしぃ\tl{˥}}
\rentry{むかで}{百足}{むがじ\tl{˥˩}}
\rentry{むぎ}{麦}{むん\tl{˩˥}}
\rentry{むぎわら}{麦わら}{むんぐる\tl{˩˥}}
\rentry{むく}{剥く}{〔皮を\ntilde{}〕むぐん\tl{˩˥}}
\rentry{むこ}{婿}{むぐ\tl{˩˥}、むぐぶざ\tl{˩˥}}
\rentry{むこきょうだい}{婿兄弟}{むぐちょーでー\tl{˩˥}}
\rentry{むし}{虫}{むし\tl{˥˩}}
\rentry{むしば}{虫歯}{ふとぅちぱん\tl{˥}}
\rentry{むしる}{}{むっすん\tl{˥˩}}
\rentry{むしろ}{蓆}{むっす\tl{˩˥}}
\rentry{むす}{蒸す}{あーらすん\tl{˥˩}}
\rentry{むずかしい}{難しい}{むすかっさ\tl{˩˥}はん}
\rentry{むすぶ}{結ぶ}{むすぷん\tl{˥˩}、ゆうん\tl{˥˩}}
\rentry{むすめ}{娘}{みどぅんたま\tl{˩˥}、みやらび\tl{˩˥}}
\rentry{むだい}{無代}{いたんだ\tl{˥}}
\rentry{むだぐち}{無駄口}{んなふち\tl{˥˩}}
\rentry{むだづかい}{無駄遣い}{むだすけー\tl{˥˩}}
\rentry{むだぼね}{無駄骨}{んなあわり\tl{˥˩}}
\rentry{むち}{鞭}{すむとぅ\tl{˥}}
\rentry{むちうつ}{鞭打つ}{すんぐるん\tl{˥}}
\rentry{むっつ}{六つ}{むーち\tl{˥˩}}
\rentry{むつまじい}{睦まじい}{むつまっ\tl{˩˥}さーん、むつまっさ\tl{˩˥}はん}
\rentry{むなさわぎ}{胸騒ぎ}{すむあーり}
\rentry{むね}{胸}{にちぃ\tl{˥˩}}
\rentry{むねのいたむおもいがする}{胸の痛む思いがする}{すむやむーん}
\rentry{むねやけ}{胸焼け}{いがばり\tl{˥˩}}
\rentry{むねん}{無念}{いなむん}
\rentry{むよくだ}{無欲だ}{ゆぐぬ\ws{}ねーぬ}
\rentry{むら}{村}{むら\tl{˥˩}}
\rentry{むらじむしょ}{村事務所}{おーしゃ\tl{˥}}
\rentry{むらのぎょうじ}{村の行事}{むらぐとぅ\tl{˥˩}}
\rentry{むらばんしょ}{村番所}{おーしゃ\tl{˥}}
\rentry{むらびとのしゅうかい}{村人の集会}{むらたたまり}
\rentry{め}{目}{みん\tl{˩˥}}
\rentry{め}{芽}{ふき\tl{˥}、べー\tl{˩˥}}
\rentry{めあて}{目当て}{みあてぃ\tl{˥}}
\rentry{めい}{姪}{ぶい\tl{˥˩}、ぶいふぁー\tl{˥˩}}
\rentry{めいにち}{命日}{みーにち\tl{˥˩}}
\rentry{めいめいする}{命名する}{なーしきん\tl{˩˥}}
\rentry{めいれい}{命令}{ぎし\tl{˩˥}}
\rentry{めいわく}{迷惑}{やっけー\tl{˩˥}}
\rentry{めうえ}{目上}{うい\tl{˥˩}}
\rentry{めうし}{牝牛}{ういなん\tl{˥}}
\rentry{めおと}{夫婦}{とぅんぶとぅ\tl{˥˩}}
\rentry{めがかたい}{目がかたい}{みんすさ\tl{˩˥}はーん}
\rentry{めかくしのいしがき}{目隠しの石垣}{なーふく\tl{˩˥}}
\rentry{めかけ}{妾}{ゆーべー\tl{˩˥}}
\rentry{めがさえる}{目がさえる}{みんくぱり\tl{˩˥}}
\rentry{めがね}{眼鏡}{がんきょー\tl{˥}、みーかんがん\tl{˩˥}}
\rentry{めがまぶしい}{目が眩しい}{みーん\tl{˩˥}ぱちま\tl{˥˩}はん}
\rentry{めくら}{}{みっくゎ\tl{˩˥}}
\rentry{めぐらす}{}{みぐらすん\tl{˥˩}}
\rentry{めぐる}{}{みぐるん\tl{˥˩}}
\rentry{めざし}{目差}{みざし\tl{˩˥}}
\rentry{めざめる}{目覚める}{うぎるん、うぎん\tl{˥}、さまるん\tl{˥}、みん\tl{˩}さまるん\tl{˥}}
\rentry{めし}{飯}{いー\tl{˥}}
\rentry{めしあがる}{召し上がる}{おーし\ws{}おるん、んげーるん\tl{˥˩}}
\rentry{めしぢゃわん}{飯茶碗}{いーまーり\tl{˥}}
\rentry{めじろ}{メジロ}{ふなどぅり\tl{˥˩}}
\rentry{めす}{雌}{みーむぬ\tl{˩˥}}
\rentry{めずらしい}{珍しい}{ぴるまさん、ぴるま\tl{˥}はん、みずら\tl{˩˥}はん、むぬみずらはん}
\rentry{めだつ}{目立つ}{みだつん\tl{˩˥}}
\rentry{めだま}{目玉}{みっちん\tl{˩˥}、みんたま\tl{˥}}
\rentry{めつぼうする}{滅亡する}{とーりるん\tl{˥}}
\rentry{めとる}{娶る}{あいなー\tl{˥}\ws{}すん\tl{˥˩}、そーるん\tl{˥˩}、〔妻を\ntilde{}〕そるん\tl{˥˩}、とぅんそるん\tl{˥˩}}
\rentry{めりけんこ}{メリケン粉}{むんぬ\tl{˩˥}\ws{}くー\tl{˥}}
\rentry{めるい}{温い}{ぬるさ\tl{˩˥}はん}
\rentry{めん}{綿}{ばた\tl{˩˥}}
\rentry{めんどう}{面倒}{みんどぅー\tl{˩˥}}
\rentry{めんどうだ}{面倒だ}{みんどな\tl{˩˥}}
\rentry{も}{喪}{いみ\tl{˥}}
\rentry{もあい}{模合}{むえー\tl{˥˩}}
\rentry{もう}{}{とー}
\rentry{もういちど}{もう一度}{まー\ws{}ぴぃとぅむし}
\rentry{もういやだ}{}{まきしょー}
\rentry{もうけ}{儲け}{もーが\tl{˥}、もーぎ\tl{˩˥}}
\rentry{もうしあげる}{申し上げる}{すさりん\tl{˥˩}}
\rentry{もうしいれる}{申し入れる}{にげー\tl{˩˥}\ws{}んじるん\tl{˥}}
\rentry{もうしょ}{猛暑}{ぷーき}
\rentry{もうじん}{盲人}{みっくゎ\tl{˩˥}}
\rentry{もうすこし}{もう少し}{ま\josho{}べ\kako{}ー}
\rentry{もうたいへん}{もう大変}{さってぃむ}
\rentry{もうちょっと}{}{ま\josho{}べ\kako{}ー}
\rentry{もうひつ}{毛筆}{ふち\tl{˥˩}}
\rentry{もうふ}{毛布}{きっとー\tl{˥}}
\rentry{もうろく}{耄碌}{ういぷら\tl{˥}}
\rentry{もうろくしている}{耄碌している}{そー\ws{}ねーぬ}
\rentry{もえだす}{燃えだす}{〔火が\ntilde{}〕てすかるん\tl{˥˩}}
\rentry{もえる}{燃える}{めーるん\tl{˥˩}}
\rentry{もぎとる}{もぎ取る}{ぶるん\tl{˩˥}}
\rentry{もくざい}{木材}{ざいぎ\tl{˩˥}}
\rentry{もくたん}{木炭}{たん\tl{˥˩}}
\rentry{もくひょう}{目標}{みあてぃ\tl{˥}}
\rentry{もぐらせる}{潜らせる}{すますん\tl{˥}}
\rentry{もぐる}{潜る}{すむん\tl{˥}}
\rentry{もじ}{文字}{じぃー\tl{˥˩}}
\rentry{もしかしたら}{}{ゆー\tl{˩˥}さんちゃ\tl{˥}}
\rentry{もしも}{}{まんいち\tl{˥˩}}
\rentry{もたげる}{}{〔顔を\ntilde{}〕むち\tl{˩˥}\ws{}あんぎるん\tl{˥˩}}
\rentry{もたす}{持たす}{むたすん\tl{˩˥}}
\rentry{もたれかかる}{}{ゆりはかるん\tl{˥˩}}
\rentry{もたれる}{}{むたーりるん\tl{˥˩}}
\rentry{もち}{餅}{むち\tl{˥˩}}
\rentry{もちあがる}{持ち上がる}{むりあがるん\tl{˥˩}}
\rentry{もちあげる}{持ち上げる}{むちあんぎるん\tl{˩˥}、〔上へ\ntilde{}〕むちあんぎるん\tl{˩˥}}
\rentry{もちごめ}{もち米}{むすめー\tl{˥˩}}
\rentry{もちなおす}{持ち直す}{むちのーすん\tl{˩˥}}
\rentry{もちなおる}{持ち直る}{むちのーるん\tl{˩˥}}
\rentry{もちぶん}{持分}{たまち\tl{˥}}
\rentry{もちもの}{持ち物}{むちぃむぬ\tl{˩˥}}
\rentry{もつ}{持つ}{むつん\tl{˩˥}}
\rentry{もっこ}{畚}{おんだ\tl{˥}}
\rentry{もっていく}{持っていく}{むち\ws{}んぐん}
\rentry{もっと}{}{まーび}
\rentry{もっとも}{}{むっ\josho{}とぅ\kako{}む}
\rentry{もつれる}{}{あんざーりるん\tl{˥}、あんざるん\tl{˥}}
\rentry{もてあそぶ}{弄ぶ}{だーぶん\tl{˩˥}}
\rentry{もてなす}{持てなす}{とぅりむつん\tl{˥}}
\rentry{もと}{元}{むとぅ\tl{˩˥}}
\rentry{もとごえ}{元肥}{むとぅぐい\tl{˩˥}}
\rentry{もどす}{戻す}{むどぅすん\tl{˩˥}}
\rentry{もとめる}{求める}{くいるん\tl{˥}、とぅみるん\tl{˥}、むとぅみるん\tl{˩˥}}
\rentry{もともと}{元々}{むとぅがら\tl{˩˥}}
\rentry{もどりかぜ}{戻り風}{ゆーべーかち\tl{˩˥}}
\rentry{もどりみち}{戻り道}{むどぅりみち\tl{˩˥}}
\rentry{もどる}{戻る}{むどぅるん\tl{˩˥}}
\rentry{もにふくする}{喪に服する}{いみはかるん\tl{˥}}
\rentry{もの}{物}{むぬ\tl{˩˥}}
\rentry{ものおと}{物音}{うとー\tl{˥˩}}
\rentry{ものおぼえ}{物覚え}{むぬうぶい\tl{˩˥}}
\rentry{ものおもい}{物思い}{むぬむい\tl{˩˥}}
\rentry{ものかんがえ}{物考え}{むぬかんげー\tl{˩˥}}
\rentry{ものごと}{物事}{むぬぐとぅ\tl{˩˥}}
\rentry{ものごとをきく}{物事を聞く}{むぬ\tl{˩˥}すくん\tl{˥˩}}
\rentry{ものさし}{物差し}{じょーぎ\tl{˩˥}}
\rentry{ものしり}{物知り}{むぬしり\tl{˩˥}}
\rentry{ものすごい}{}{うすまさん}
\rentry{ものならい}{物習い}{むぬなれー\tl{˩˥}}
\rentry{ものめずらしい}{物珍しい}{むぬみずらはん}
\rentry{ものもち}{物持ち}{むぬむち\tl{˩˥}}
\rentry{ものもらい}{}{いっふぇー\tl{˥}、みんぜー\tl{˩˥}}
\rentry{ものわすれ}{物忘れ}{むぬばっさ}
\rentry{ものわらいになる}{物笑いになる}{ばーらいるん\tl{˥˩}}
\rentry{ものをおぎなう}{物を補う}{たらーすん\tl{˥˩}}
\rentry{もみ}{籾}{もみ\tl{˥˩}}
\rentry{もみがら}{籾殻}{すくぶ\tl{˥}}
\rentry{もみつぶ}{籾粒}{〔白米中の\ntilde{}〕あら\tl{˥˩}}
\rentry{もみほぐす}{揉みほぐす}{むむん\tl{˥˩}}
\rentry{もむ}{揉む}{むむん\tl{˥˩}}
\rentry{もめん}{木綿}{ばた\tl{˩˥}、むみん\tl{˥}}
\rentry{もも}{腿}{むんだらし\tl{˩˥}}
\rentry{もやし}{モヤシ}{まーみなん\tl{˩˥}}
\rentry{もやす}{燃やす}{めーすん\tl{˥˩}}
\rentry{もよう}{模様}{むよー\tl{˥˩}}
\rentry{もらいもの}{貰い物}{\josho{}たぼらり\kako{}むぬ}
\rentry{もらす}{漏らす}{もーらすん\tl{˩˥}}
\rentry{もり}{森}{むり\tl{˥˩}、やま\tl{˩˥}}
\rentry{もり}{銛}{うい\tl{˥}}
\rentry{もりあがる}{盛り上がる}{むりあがるん\tl{˥˩}}
\rentry{もりあげる}{盛り上げる}{むりあぎるん\tl{˥˩}、むるん}
\rentry{もる}{漏る}{もーるん\tl{˩˥}}
\rentry{もる}{盛る}{むるん}
\rentry{もれる}{漏れる}{もーるん\tl{˩˥}}
\rentry{もろみ}{}{むるん\tl{˥}}
\rentry{もん}{門}{ぞー\tl{˩˥}}
\rentry{もんどう}{問答}{むんどぅー\tl{˩˥}}
\rentry{もんぱのき}{モンパノキ}{すさざぎ\tl{˥}}
\rentry{もんめ}{匁}{むんみ\tl{˥}}
\rentry{や}{矢}{や\tl{˥˩}}
\rentry{やぁー}{}{やー}
\rentry{やいば}{刃}{ぱん\tl{˥}}
\rentry{やえやま}{八重山}{やいま\tl{˥˩}}
\rentry{やえやまあおき}{ヤエヤマアオキ}{ぶ\josho{}が\kako{}き}
\rentry{やかましい}{}{やがまっさ\tl{˩˥}はーん}
\rentry{やかましい}{喧しい}{かつぁま\tl{˥˩}はーん}
\rentry{やかん}{夜間}{ゆる\tl{˩˥}}
\rentry{やぎ}{山羊}{ぴみざ\tl{˥}}
\rentry{やきすてる}{焼き捨てる}{やぎしっしるん、やぐん\tl{˥˩}}
\rentry{やきもの}{焼き物}{やぎむぬ\tl{˩˥}}
\rentry{やぎょう}{夜業}{ぶなび}
\rentry{やく}{厄}{やく\tl{˩˥}}
\rentry{やく}{焼く}{あぶるん\tl{˥}、やぐん\tl{˥˩}}
\rentry{やくしょ}{役所}{やくば\tl{˩˥}}
\rentry{やくしょく}{役職}{やく\tl{˩˥}}
\rentry{やくそく}{約束}{かーき\tl{˥}}
\rentry{やくどし}{厄年}{やくどぅし\tl{˥˩}}
\rentry{やくにたつ}{役に立つ}{やく\tl{˩˥}たつん\tl{˥}}
\rentry{やくば}{役場}{やくば\tl{˩˥}}
\rentry{やくひん}{薬品}{ふちり\tl{˥}}
\rentry{やけどをする}{火傷をする}{やぴくん\tl{˥˩}}
\rentry{やける}{焼ける}{やがりるん\tl{˥˩}}
\rentry{やさい}{野菜}{やせー\tl{˩˥}}
\rentry{やさいばたけ}{野菜畑}{やせーぴてー}
\rentry{やし}{椰子}{やし\tl{˩˥}}
\rentry{やしがに}{やし蟹}{むごん\tl{˩˥}}
\rentry{やしき}{屋敷}{はこー\tl{˥˩}、やしき\tl{˩˥}}
\rentry{やしない}{養い}{やしねー}
\rentry{やしないおや}{養い親}{すかなーうや\tl{˥}、やしねーうや}
\rentry{やしないご}{養い子}{やしぃねーふぁ}
\rentry{やしなう}{養う}{すかのーすん\tl{˥}、やすなすん\tl{˥˩}}
\rentry{やすい}{安い}{やっさ\tl{˩˥}はん}
\rentry{やすっぽい}{安っぽい}{そーべー\tl{˥}}
\rentry{やすませる}{休ませる}{よがーしみるん}
\rentry{やすみ}{休み}{やしみ\tl{˩˥}}
\rentry{やすむ}{休む}{どぅー\tl{˩˥}よがすん\tl{˩˥}、やすぅむん\tl{˩˥}、よがすん\tl{˩˥}}
\rentry{やすもの}{安物}{やしむぬ\tl{˩˥}}
\rentry{やすり}{鑢}{やしり\tl{˩˥}}
\rentry{やすんじる}{安んじる}{やすんじるん\tl{˩˥}}
\rentry{やせおとろえる}{痩せ衰える}{がんどらん\tl{˩˥}、よーがりやん}
\rentry{やせち}{痩せ地}{ぱぎじ\tl{˥}}
\rentry{やせる}{痩せる}{やっちるん\tl{˩˥}、よーがりるん\tl{˥˩}、よーがりん\tl{˥˩}}
\rentry{やっかい}{厄介}{やっけー\tl{˩˥}}
\rentry{やったー}{}{したい}
\rentry{やっつ}{八つ}{やーち\tl{˥}}
\rentry{やっつける}{}{しゃーみるん\tl{˥˩}}
\rentry{やっと}{}{やっとぅ}
\rentry{やど}{宿}{やどぅ\tl{˩˥}}
\rentry{やとう}{雇う}{やとすん\tl{˩˥}}
\rentry{やどかり}{ヤドカリ}{あまん\tl{˥}}
\rentry{やどや}{宿屋}{やどぅ\tl{˩˥}}
\rentry{やどる}{宿る}{とぅまるん\tl{˥˩}、やどぅとぅるん}
\rentry{やね}{屋根}{ぴぃさ\tl{˥˩}}
\rentry{やはん}{家判}{やーばん\tl{˩˥}}
\rentry{やぶってすてる}{破って捨てる}{やりしちるん\tl{˩˥}}
\rentry{やぶのなか}{藪の中}{きぬ\tl{˥}\ws{}みー\tl{˥˩}}
\rentry{やぶりちらかす}{破り散らかす}{やりぽーるん\tl{˩˥}}
\rentry{やぶる}{破る}{やぶるん\tl{˩˥}}
\rentry{やぶれる}{破れる}{ざりるん\tl{˩˥}、やーりん\tl{˩˥}、やぶりん\tl{˩˥}、やりるん\tl{˩˥}}
\rentry{やま}{山}{だき\tl{˥˩}、むり\tl{˥˩}、やま\tl{˩˥}}
\rentry{やまい}{病}{ぱなしき\tl{˥˩}、やみ\tl{˩˥}、やん\tl{˩˥}}
\rentry{やまいも}{山芋}{かりょん\tl{˥}}
\rentry{やまがたな}{山刀}{やまかたなし\tl{˩˥}}
\rentry{やまぐに}{山国}{やまじぃま\tl{˩˥}}
\rentry{やましごと}{山仕事}{やましぐとぅ\tl{˩˥}}
\rentry{やまじま}{山島}{やまじぃま\tl{˩˥}}
\rentry{やまと}{大和}{やまとぅ\tl{˩˥}}
\rentry{やまどめ}{山止め}{やまどぅみ\tl{˩˥}}
\rentry{やまのないへいたんなしま}{山のない平坦な島}{ぬぐじま\tl{˥˩}}
\rentry{やまほさ}{山補佐}{やーぶさ\tl{˩˥}}
\rentry{やみ}{闇}{やみ\tl{˥˩}}
\rentry{やむ}{止む}{〔風雨が\ntilde{}〕やむん\tl{˥˩}}
\rentry{やもり}{守宮}{ふたしみ\tl{˥}}
\rentry{やり}{槍}{やり\tl{˥˩}}
\rentry{やりこめる}{やり込める}{しゃーみるん\tl{˥˩}}
\rentry{やりたくない}{}{しぼへぬ}
\rentry{やる}{}{すん\tl{˥˩}}
\rentry{やるもんか}{}{まきしょー}
\rentry{やれめ}{破れ目}{やりふちぃ\tl{˩˥}}
\rentry{やわらかい}{柔らかい}{やら\tl{˩˥}はん}
\rentry{やわらかくする}{柔らかくする}{くなすん\tl{˥˩}、やーらぎるん\tl{˥˩}、やらふぁー\tl{˩˥}\ws{}すん\tl{˥˩}}
\rentry{やわらかくなる}{柔らかくなる}{やらふぁー\tl{˩˥}\ws{}なるん\tl{˩˥}}
\rentry{やわらかになる}{柔らかになる}{やーらぐん\tl{˥˩}、やふぁらぐん}
\rentry{やわらぐ}{和らぐ}{やふぁらぐん}
\rentry{やわらぐ}{柔らぐ}{やーらぐん\tl{˥˩}}
\rentry{やわらげる}{和らげる}{やーらぎるん\tl{˥˩}}
\rentry{やんばるせん}{山原船}{やんばるしん\tl{˩˥}}
\rentry{ゆ}{湯}{あちぃゆ\tl{˥}}
\rentry{ゆい}{結}{ゆい\tl{˥˩}}
\rentry{ゆいごん}{遺言}{んなん}
\rentry{ゆう}{結う}{ゆうん\tl{˥˩}}
\rentry{ゆうがおのみ}{夕顔の実}{かなばりん\tl{˥}}
\rentry{ゆうかく}{遊廓}{さかなやー\tl{˥˩}}
\rentry{ゆうがた}{夕方}{ゆねん\tl{˩˥}}
\rentry{ゆうき}{勇気}{いじぃ\tl{˥˩}}
\rentry{ゆうぐれ}{夕暮れ}{ゆねん\tl{˩˥}}
\rentry{ゆうしょく}{夕食}{ゆむぬ\tl{˩˥}}
\rentry{ゆうじん}{友人}{どぅし\tl{˥˩}}
\rentry{ゆうせいだ}{優勢だ}{すさ\tl{˥}はん}
\rentry{ゆうだち}{夕立ち}{ぺーしゅら\tl{˥}}
\rentry{ゆうな}{ユウナ}{ゆだ\tl{˥˩}}
\rentry{ゆうなぎ}{夕凪}{ゆーどぅり\tl{˩˥}}
\rentry{ゆうのうな}{有能な}{かねー}
\rentry{ゆうはん}{夕飯}{ゆむぬ\tl{˩˥}}
\rentry{ゆうべ}{夕べ}{ゆーび\tl{˩˥}}
\rentry{ゆうめいだ}{有名だ}{うとぅだが\tl{˥˩}はん、なー\tl{˥˩}たか\tl{˥}はん}
\rentry{ゆうれい}{幽霊}{まざむん\tl{˥˩}}
\rentry{ゆか}{床}{ふんた\tl{˥˩}}
\rentry{ゆがく}{湯がく}{ゆぴくん\tl{˩˥}}
\rentry{ゆかした}{床下}{いためー\tl{˥}}
\rentry{ゆがむ}{歪む}{ゆがむん\tl{˥˩}}
\rentry{ゆがめる}{歪める}{ゆがみるん\tl{˥˩}}
\rentry{ゆきあう}{行き会う}{いげーるん\tl{˥˩}}
\rentry{ゆくすえ}{行く末}{あとぅあとぅ\tl{˩˥}}
\rentry{ゆげ}{湯気}{きぷ\tl{˥˩}}
\rentry{ゆする}{揺する}{くいつぁーすん\tl{˥}、ゆらすん\tl{˥˩}}
\rentry{ゆずる}{譲る}{ゆずるん\tl{˥˩}}
\rentry{ゆだんする}{油断する}{うくたるん\tl{˥˩}}
\rentry{ゆっくり}{}{よんなよんな}
\rentry{ゆっくりと}{}{や\josho{}まー\kako{}し}
\rentry{ゆでる}{茹でる}{ゆぴくん\tl{˩˥}}
\rentry{ゆのみ}{湯呑}{さばん\tl{˥}}
\rentry{ゆびぶえ}{指笛}{しーふき\tl{˥}}
\rentry{ゆびわ}{指輪}{うるんがに\tl{˩˥}}
\rentry{ゆみ}{弓}{ゆん\tl{˥˩}}
\rentry{ゆめ}{夢}{いみ\tl{˥}}
\rentry{ゆらい}{由来}{いわり\tl{˥˩}}
\rentry{ゆらす}{揺らす}{くいつぁーすん\tl{˥}、すくふぁすん\tl{˥˩}}
\rentry{ゆらりと}{}{よんなよんな}
\rentry{ゆるい}{緩い}{ゆるさ\tl{˩˥}はん}
\rentry{ゆるされる}{許される}{ぬがーるん\tl{˩˥}、ゆるさーるん}
\rentry{ゆるし}{許し}{ゆるし\tl{˩˥}}
\rentry{ゆるす}{許す}{ぬがーらすん、ゆるすん\tl{˩˥}}
\rentry{ゆるむ}{緩む}{ふきっつぁーるん\tl{˥˩}}
\rentry{ゆるめる}{緩める}{ゆるみるん\tl{˩˥}、ゆるみん\tl{˩˥}}
\rentry{ゆれうごくさま}{揺れ動くさま}{がたがた}
\rentry{ゆれる}{揺れる}{ふっふん\tl{˥˩}}
\rentry{よ}{世}{ゆー\tl{˥˩}}
\rentry{よあけまえ}{夜明け前}{あぎさり}
\rentry{よい}{良い}{まーし\tl{˩˥}、みしゃ\tl{˥}はん}
\rentry{よいこと}{良いこと}{\josho{}いー\kako{}くとぅ}
\rentry{よいつぶれる}{酔いつぶれる}{びーたこりるん\tl{˩˥}}
\rentry{よいのみょうじょう}{宵の明星}{\josho{}しかま\kako{}ぷち}
\rentry{よいもの}{良い物}{しょーむぬ\tl{˥˩}}
\rentry{よういくする}{養育する}{すかなすん\tl{˥}、やすなすん\tl{˥˩}}
\rentry{よういする}{用意する}{しこーるん\tl{˥˩}}
\rentry{よういに}{容易に}{よーに\tl{˩˥}}
\rentry{ようかい}{妖怪}{まざむん\tl{˥˩}}
\rentry{ようがさ}{洋傘}{かぷるさな\tl{˥}、さな\tl{˥}、だんがさ\tl{˥}}
\rentry{ようき}{容器}{いりむん\tl{˥˩}}
\rentry{ようきめい}{容器名}{がいじぃ\tl{˩˥}}
\rentry{ようし}{養子}{やしぃねーふぁ}
\rentry{ようじ}{幼児}{いしゃがーたま\tl{˥}}
\rentry{ようじ}{用事}{ゆーじ\tl{˩˥}}
\rentry{ようす}{様子}{むよー\tl{˥˩}}
\rentry{ようせき}{陽石}{びっちゅる\tl{˩˥}}
\rentry{ようち}{幼稚}{すぅな\tl{˥}はん}
\rentry{ようはい}{遥拝}{たんかにげー\tl{˩˥}}
\rentry{ようふぼ}{養父母}{すかなーうや\tl{˥}、やしねーうや}
\rentry{ようみょう}{幼名}{やらびなー\tl{˩˥}}
\rentry{ようやく}{}{やっとぅ}
\rentry{よがわり}{世変り}{ゆーがーり\tl{˥˩}}
\rentry{よく}{好く}{すくん\tl{˥}}
\rentry{よく}{欲}{ゆぐ\tl{˩˥}}
\rentry{よくあう}{よく合う}{ゆー\tl{˩˥}\ws{}あうん\tl{˥}}
\rentry{よくいく}{よく行く}{しゃー\tl{˥}\ws{}んぐん\tl{˥˩}}
\rentry{よくいたずらをする}{よく悪戯をする}{がんまり\ws{}しゃーん}
\rentry{よくがない}{欲がない}{ゆぐぬ\ws{}ねーぬ}
\rentry{よくかむ}{よく噛む}{かんだるん\tl{˥}}
\rentry{よくちょう}{翌朝}{あつぁすとぅむち\tl{˥}}
\rentry{よくない}{良くない}{ぱからっさー\tl{˥}\ws{}ねーぬ\tl{˩˥}}
\rentry{よくなる}{良くなる}{みしゃー\tl{˥}\ws{}なるん\tl{˩˥}}
\rentry{よくにあう}{よく似合う}{ゆー\tl{˩˥}\ws{}あうん\tl{˥}}
\rentry{よくねむる}{よく眠る}{ぬっふたりるん}
\rentry{よくはたらく}{よく働く}{ぎばるん\tl{˥˩}}
\rentry{よくばりだ}{欲張りだ}{ゆぐ\tl{˩˥}さん}
\rentry{よくばる}{欲張る}{ゆくすか\tl{˩˥}、ゆくばるん}
\rentry{よくぶかい}{欲深い}{ゆぐ\tl{˩˥}さん}
\rentry{よくよく}{}{よー\tl{˥˩}、〔強調\ntilde{}〕よーよー}
\rentry{よこうえんしゅう}{予行演習}{しらす\tl{˥˩}}
\rentry{よこす}{}{やらすん\tl{˥˩}}
\rentry{よごす}{汚す}{ゆぐすん\tl{˥˩}}
\rentry{よこたえる}{横たえる}{くるばすん\tl{˥˩}}
\rentry{よこたわる}{横たわる}{ゆくたーるん\tl{˩˥}}
\rentry{よこになる}{横になる}{ゆくたーるん\tl{˩˥}}
\rentry{よこべり}{横舷}{ふなばに\tl{˥}}
\rentry{よごれ}{汚れ}{ふたり\tl{˥}、ゆぐり\tl{˥˩}}
\rentry{よごれている}{汚れている}{やにっしゃ\tl{˩˥}はん}
\rentry{よごれる}{汚れる}{ゆぐりるん\tl{˥˩}}
\rentry{よさめ}{夜雨}{ゆーあみ\tl{˩˥}、ゆるあみ\tl{˩˥}}
\rentry{よすみ}{四隅}{ゆーかどぅ\tl{˥˩}}
\rentry{よせあつめる}{寄せ集める}{ぴきゆしるん\tl{˥˩}}
\rentry{よせる}{寄せる}{ゆしるん\tl{˥˩}}
\rentry{よだれ}{涎}{ゆだり\tl{˥˩}}
\rentry{よちょう}{予兆}{むぬしらし\tl{˩˥}}
\rentry{よつかど}{四ツ角}{ゆーかどぅ\tl{˥˩}}
\rentry{よっつ}{四つ}{ゆーち\tl{˥˩}}
\rentry{よっぱらい}{酔っ払い}{びーちゃー\tl{˩˥}}
\rentry{よつぼし}{四つ星}{ゆさし\tl{˥˩}}
\rentry{よてい}{予定}{しゃーみ\tl{˥}}
\rentry{よどおし}{夜通し}{ゆどぅし}
\rentry{よどむ}{淀む}{ゆどぅむん\tl{˩˥}}
\rentry{よなあべごや}{よなべ小屋}{ぶなびやー\tl{˥}}
\rentry{よなか}{夜中}{ゆなが\tl{˩˥}}
\rentry{よなぐに}{与那国}{ゆのー\tl{˥˩}}
\rentry{よなべ}{}{ぶなび}
\rentry{よになりひびく}{世に鳴り響く}{とぅゆまりん\tl{˥}}
\rentry{よねがい}{世願い}{ゆにげー\tl{˥˩}}
\rentry{よのなか}{世の中}{しきん\tl{˥}}
\rentry{よぶ}{呼ぶ}{ゆぶん\tl{˥˩}}
\rentry{よぶん}{余分}{あまり\tl{˥}、おーばー\tl{˥˩}}
\rentry{よみがえらせる}{蘇らせる}{いがすん}
\rentry{よむ}{読む}{ゆむん\tl{˩˥}}
\rentry{よめ}{嫁}{とぅん\tl{˥˩}、ゆみ\tl{˥˩}}
\rentry{よめにいく}{嫁に行く}{ぶとぅむち}
\rentry{よめをさがす}{探す}{〔嫁を\ntilde{}〕とぅみん\tl{˥}}
\rentry{よもぎ}{蓬}{やたふちぃ}
\rentry{よりあう}{寄り合う}{ゆるん\tl{˥˩}}
\rentry{よりかかる}{}{ゆりはかるん\tl{˥˩}}
\rentry{よりかかる}{寄り掛かる}{むたーりるん\tl{˥˩}}
\rentry{よる}{夜}{ゆる\tl{˩˥}}
\rentry{よる}{寄る}{ゆるん\tl{˥˩}}
\rentry{よろこぶ}{喜ぶ}{あまいるん\tl{˥˩}、さにしゃー\ws{}すん}
\rentry{よろしい}{}{みしゃ\tl{˥}はん}
\rentry{よわい}{弱い}{よー\tl{˩˥}はん}
\rentry{よわくなる}{弱くなる}{よーるん\tl{˩˥}}
\rentry{よわむし}{弱虫}{よーば\tl{˩˥}}
\rentry{よわめる}{弱める}{ゆるみるん\tl{˩˥}、ゆるみん\tl{˩˥}、よーみるん\tl{˩˥}}
\rentry{よわる}{弱る}{よーるん\tl{˩˥}}
\rentry{よんじゅう}{四十}{しんず\tl{˥}}
\rentry{よんちゅ}{与人}{ゆんちゅ\tl{˥˩}}
\rentry{らいおん}{ライオン}{しーし\tl{˥}}
\rentry{らいねん}{来年}{えん\tl{˩˥}、げん}
\rentry{らいねんのなつ}{来年の夏}{くなちぃ\tl{˥}}
\rentry{らいめい}{雷鳴}{かんなり\tl{˥}}
\rentry{らく}{楽}{らく\tl{˩˥}}
\rentry{らくをする}{楽をする}{らく\tl{˩˥} \ws{}すん\tl{˥˩}}
\rentry{らたい}{裸体}{あばだり\tl{˥}}
\rentry{らっかせい}{落花生}{じぃーまーみ\tl{˩˥}}
\rentry{らっきょう}{ラッキョウ}{だっきょん\tl{˩˥}}
\rentry{らんがさ}{蘭傘}{かぷるさな\tl{˥}}
\rentry{らんそう}{卵巣}{〔魚の\ntilde{}〕ぱりゃん\tl{˥}}
\rentry{らんぷ}{ランプ}{とぅり\tl{˥}}
\rentry{らんぼうもの}{乱暴者}{やまぐ\tl{˩˥}、やまんぐ}
\rentry{らんぼうをはたらく}{乱暴を働く}{やまぐ\tl{˩˥}\ws{}すん\tl{˥˩}}
\rentry{りーふ}{リーフ}{じり\tl{˩˥}}
\rentry{りえき}{利益}{もーが\tl{˥}、もーぎ\tl{˩˥}}
\rentry{りかいする}{理解する}{ばがるん\tl{˩˥}}
\rentry{りきむ}{力む}{んくむん\tl{˥˩}}
\rentry{りくつ}{理屈}{りくち\tl{˩˥}}
\rentry{りこう}{利口}{そーいり\tl{˥˩}}
\rentry{りこうな}{利口な}{めふな\tl{˥˩}}
\rentry{りこんする}{離婚する}{ばがりん\tl{˩˥}}
\rentry{りちぎだ}{律義だ}{ぎりかた\tl{˩˥}はん}
\rentry{りっぱ}{立派}{ぞーとぅ\tl{˩˥}、ぱからっ\tl{˥}さ、みぐとぅ\tl{˩˥}}
\rentry{りっぱだ}{立派だ}{だーさ\tl{˩˥}はん、まだ\tl{˥˩}はん}
\rentry{りっぱなもの}{立派な物}{しょーむぬ\tl{˥˩}}
\rentry{りっぱに}{立派に}{けーし、こーし}
\rentry{りとく}{利得}{とぅく\tl{˥}}
\rentry{りはーさる}{リハーサル}{しらす\tl{˥˩}}
\rentry{りはつ}{理髪}{だんぱち\tl{˥}}
\rentry{りはつてん}{理髪店}{だんぱちーやー}
\rentry{りゅうきゅうこくおう}{琉球国王}{うしゅがなし\tl{˥˩}}
\rentry{りゅうきゅうこくたん}{琉球黒檀}{きな\tl{˥˩}}
\rentry{りゅうこう}{流行}{はやり\tl{˥}}
\rentry{りゅうこうする}{流行する}{はやーるん\tl{˥}}
\rentry{りゅうしゅつする}{流失する}{なんぶりん\tl{˩˥}}
\rentry{りゅうぜつらん}{竜舌蘭}{るがい\tl{˥}}
\rentry{りょうがえする}{両替する}{かれーるん\tl{˥˩}}
\rentry{りょうしん}{両親}{ふたうや\tl{˥˩}}
\rentry{りょうて}{両手}{ふたし\tl{˥˩}}
\rentry{りょうてい}{料亭}{さかなやー\tl{˥˩}}
\rentry{りょうてで}{両手で}{ふたし\tl{˥˩}}
\rentry{りょうほう}{両方}{ふたかた\tl{˥˩}}
\rentry{りょうほうとも}{両方とも}{ふたーちぃ\tl{˥˩}ばぎ\tl{˩˥}}
\rentry{りょうめ}{両目}{ふたみん\tl{˥˩}}
\rentry{りょこう}{旅行}{たぴ\tl{˥˩}}
\rentry{りんげつ}{臨月}{まりしき\tl{˥˩}}
\rentry{るす}{留守}{ひな\tl{˩˥}\ws{}のーぬ\tl{˩˥}}
\rentry{るすばん}{留守番}{ひぬばん\tl{˩˥}}
\rentry{れいかんのうりょくがたかい}{霊感能力が高い}{さーだが\tl{˥˩}はん}
\rentry{れいぎ}{礼儀}{りーぎ\tl{˥˩}}
\rentry{れいげんあらたかだ}{霊験あらたかだ}{さーだが\tl{˥˩}はん}
\rentry{れいこん}{霊魂}{たましぃ\tl{˥}、たまち\tl{˥}、まーぶり\tl{˩˥}}
\rentry{れいはいする}{礼拝する}{しー\tl{˥}うさぎるん\tl{˥˩}}
\rentry{れつあくだ}{劣悪だ}{やな\tl{˥˩}}
\rentry{れつをととのえる}{列を整える}{するいるん\tl{˥}}
\rentry{れんしゅう}{練習}{きっく\tl{˥}}
\rentry{ろ}{炉}{すか\tl{˥}}
\rentry{～ろ}{}{り}
\rentry{ろうえき}{労役}{すかま\tl{˥}}
\rentry{ろうがんになる}{老眼になる}{みんみり\tl{˩˥}}
\rentry{ろうじん}{老人}{うしとぅ\tl{˥}}
\rentry{ろうじんになる}{老人になる}{うしとぅ\tl{˥}\ws{}なるん\tl{˩˥}}
\rentry{ろうば}{老婆}{うしとぅぱー\tl{˥}}
\rentry{ろうや}{老爺}{うしとぅぶや}
\rentry{ろくじゅう}{六十}{るくず\tl{˩˥}}
\rentry{ろくでなし}{}{やびむん\tl{˩˥}}
\rentry{わぁー}{}{べー}
\rentry{わあわあ}{}{わーわー}
\rentry{わが}{我が}{べー\tl{˩˥}}
\rentry{わかい}{若い}{ばが\tl{˩˥}はん}
\rentry{わかいさせる}{和解させる}{ばぎるん\tl{˩˥}}
\rentry{わかがえる}{若返る}{ばがけーる\tl{˩˥}}
\rentry{わがしま}{我が島}{べーすま}
\rentry{わかしらが}{若白髪}{ばがっせー\tl{˩˥}}
\rentry{わかす}{沸かす}{〔湯を\ntilde{}〕ふかすん\tl{˥}}
\rentry{わかつき}{若月}{ばがすけん\tl{˩˥}}
\rentry{わかなつ}{若夏}{ばがなち\tl{˩˥}}
\rentry{わかば}{若葉}{ばがぱ\tl{˩˥}}
\rentry{わがまま}{我がまま}{ふんでー\tl{˥˩}}
\rentry{わがままだ}{我がままだ}{ふんだやー}
\rentry{わかもの}{若者}{ばがむぬ\tl{˩˥}}
\rentry{わかる}{分かる}{ばがるん\tl{˩˥}}
\rentry{わかれる}{別れる}{ばがりん\tl{˩˥}、ぱなりるん\tl{˥}}
\rentry{わき}{脇}{ぐちみん\tl{˩˥}、やた\tl{˩˥}}
\rentry{わきでる}{湧き出る}{ばーぐん\tl{˩˥}}
\rentry{わきのした}{脇の下}{ぐちみん\tl{˩˥}}
\rentry{わきみず}{湧水}{ばぎみじ\tl{˥˩}、ぱりみじぃ\tl{˥}}
\rentry{わきめ}{脇芽}{またべー\tl{˩˥}}
\rentry{わけまえ}{分け前}{くまた\tl{˥}、ばぎだま\tl{˩˥}}
\rentry{わける}{分ける}{ばぎるん\tl{˩˥}、ばぎん\tl{˩˥}}
\rentry{わざ}{技}{わざ\tl{˩˥}}
\rentry{わざ}{業}{わざ\tl{˩˥}}
\rentry{わざわい}{災い}{やく\tl{˩˥}、やふ\tl{˩˥}}
\rentry{わし}{鷲}{ばしぃ\tl{˥˩}}
\rentry{わずか}{}{いきら\tl{˥}さん、べー\tl{˩˥}、べ\josho{}ー\kako{}な、\josho{}べー\kako{}び}
\rentry{わずらう}{患う}{はかるん\tl{˥}}
\rentry{わすれっぽい}{忘れっぽい}{ばっし\tl{˥˩}めさ\tl{˩˥}はん}
\rentry{わすれもの}{忘れ物}{ばっしむぬ}
\rentry{わすれる}{忘れる}{ばっしるん\tl{˥˩}、ばっしん\tl{˥˩}}
\rentry{わたし}{私}{ばー\tl{˩˥}、ばぬ\tl{˩˥}}
\rentry{わたしたち}{私たち}{ばいま\tl{˩˥}}
\rentry{わたしの}{私の}{ばー\tl{˩˥}}
\rentry{わたしのいえ}{私の家}{ば\tl{˥}ひー\tl{˩˥}}
\rentry{わたしのこ}{私の子}{ばー\ws{}うたま}
\rentry{わたす}{渡す}{ばたすん\tl{˥˩}}
\rentry{わたる}{渡る}{ばたるん\tl{˥˩}}
\rentry{わってしまう}{割ってしまう}{ばり\ws{}っしん}
\rentry{わら}{藁}{ばら\tl{˩˥}}
\rentry{わらう}{笑う}{ばーるん\tl{˥˩}}
\rentry{わらざん}{藁算}{ばらざん\tl{˩˥}}
\rentry{わらじ}{草鞋}{ふちぃ\tl{˥}}
\rentry{わらせいのわらじ}{藁製の草鞋}{ばらふちぃ}
\rentry{わらたば}{藁束}{ばらふた}
\rentry{わらなわ}{藁縄}{ばらじぃな\tl{˩˥}}
\rentry{わらべ}{童}{やらび\tl{˩˥}}
\rentry{わらむしろ}{藁蓆}{みなぶ\tl{˩˥}}
\rentry{わらわれる}{笑われる}{ばーらいるん\tl{˥˩}}
\rentry{わりあて}{割当て}{たまち\tl{˥}}
\rentry{わりあてる}{割り当てる}{ばりあてぃん\tl{˥˩}、わりあてぃん\tl{˥}}
\rentry{わりくだく}{割り砕く}{ばりくだくん\tl{˩˥}}
\rentry{わる}{割る}{さくん\tl{˥}}
\rentry{わるい}{悪い}{やな\tl{˥˩}、わっさ\tl{˩˥}はん}
\rentry{わるいかたち}{悪い形}{やなかたち\tl{˥˩}}
\rentry{わるいくせ}{悪い癖}{やなくしー\tl{˥˩}}
\rentry{わるいこと}{悪いこと}{あく\tl{˥}}
\rentry{わるいしゅうかん}{悪い習慣}{やななれー\tl{˩˥}}
\rentry{わるぐち}{悪口}{あく\tl{˥}、やなふちぃ\tl{˩˥}}
\rentry{わるぐちをいうこと}{悪口を言うこと}{ばたふちり\tl{˩˥}}
\rentry{わるさ}{悪さ}{がんまり\tl{˩˥}、わしく\tl{˩˥}}
\rentry{わるじえのもの}{悪知恵の者}{りくちぃむち\tl{˩˥}}
\rentry{わるだくみ}{悪だくみ}{やなたくみ\tl{˩˥}}
\rentry{わるもの}{悪者}{やなむん\tl{˥˩}}
\rentry{われ}{我}{ばぬ\tl{˩˥}、べー\tl{˩˥}}
\rentry{われめ}{割れ目}{ばーり\tl{˥˩}}
\rentry{わん}{湾}{うら\tl{˥}}
\rentry{わん}{碗}{まーり\tl{˩˥}}
\rentry{わんぱくだ}{腕白だ}{がま\tl{˩˥}さーん、やまんぐ\tl{˩˥}さーん}
\rentry{わんぱくもの}{腕白者}{やまぐ\tl{˩˥}、やまんぐ}

  \end{multicols*}
  \clearpage
\endgroup

% =================================
% ---------- BACK MATTER ---------
% =================================

\pagestyle{empty} % switch to 'plain' for having page numbers on blank pages as well

% !TEX root = ../main.tex
\chapter*{あとがき}\label{atogaki}
\thispagestyle{plain}
\addcontentsline{toc}{chapter}{あとがき}

〇〇〇〇〇〇〇〇〇〇〇〇〇〇〇〇〇〇〇〇〇〇〇〇〇〇〇〇〇〇〇〇〇〇〇〇〇〇〇〇〇〇〇〇〇〇〇〇〇〇〇〇〇〇〇〇〇〇〇〇〇〇〇〇〇〇〇〇〇〇〇〇〇

〇〇〇〇〇〇〇〇〇〇〇〇〇〇〇〇〇〇〇〇〇〇〇〇〇〇〇〇〇〇〇〇〇〇〇〇〇〇〇〇〇〇〇〇〇〇〇〇〇〇〇〇〇〇〇〇〇〇〇〇〇〇〇〇〇〇〇〇〇〇〇〇〇〇〇〇〇〇〇〇〇〇〇〇〇〇〇〇〇〇〇〇〇〇〇〇〇〇〇〇〇〇〇〇〇〇〇〇〇〇〇〇〇

〇〇〇〇〇〇〇〇〇〇〇〇〇〇〇〇〇〇〇〇〇〇〇〇〇〇〇〇〇〇〇〇〇〇〇〇〇〇〇〇〇〇〇〇〇〇〇〇〇〇〇〇〇〇〇〇〇〇〇〇〇〇〇〇〇〇〇〇〇〇〇〇〇〇〇〇〇〇〇〇〇〇〇〇〇〇〇〇〇〇〇〇〇〇〇〇〇〇〇〇〇〇〇〇〇〇〇〇〇〇〇〇〇

\begin{flushright}
令和〇年〇月〇日\\
辞書太郎
\end{flushright}

% !TEX root = ../main.tex
\chapter*{謝辞}\thispagestyle{plain}
\addcontentsline{toc}{chapter}{謝辞}

この辞典は数多くの方々の助けなきには生まれなかったに相違ありません。〇〇〇〇〇〇〇〇〇〇〇〇〇〇〇〇〇〇〇〇〇〇〇〇〇〇〇〇〇〇〇〇〇〇〇〇〇〇〇〇〇〇〇〇〇〇〇〇〇〇〇〇〇〇。

この辞典はJSPS科研費〇〇、〇〇〇〇〇〇〇〇〇〇〇〇〇〇〇〇〇の助成を受けたものです。

\clearpage
\ifodd\value{page}\else\null\newpage\fi % for 奥付 on the last odd page
% !TEX root = ../main.tex
\thispagestyle{empty}
\ws{}

% 著者情報
\begin{table}[t!]%
	\centering%
	\begin{tabular}{lcll}%
	    \multicolumn{4}{l}{ {\textbf{辞書太郎}（じしょたろう） }}\\
	    \multicolumn{4}{l}{ {\small 1910年10月09日　沖縄県〇〇〇〇}}\\
	    \vspace{2pt}
		%\hline
		\small 経歴 &&\\
        \small 〇〇－〇〇年　〇〇〇〇 && \\
        \small 〇〇－〇〇年　〇〇〇〇 && \\
        \small 〇〇－〇〇年　〇〇〇〇 && \\
        \small 〇〇－〇〇年　〇〇〇〇 && \\
        \small 〇〇－〇〇年　〇〇〇〇 && \\
        \small 〇〇－〇〇年　〇〇〇〇 && \\
	\end{tabular}%
\end{table}%

\;\vspace{80pt} % space for vertical adjustment

\begin{table}[h!]%
	\centering%
	\begin{tabular}{lcll}%
		\multicolumn{4}{l}{ {\bfseries{\large{〇〇〇方言辞典}}} }	\\
		\hline
		\multicolumn{4}{l}{2020年3月31日発行}\\
		%発行日 && 2020年3月30日 & \\
		著　者 && 〇〇〇〇、〇〇〇〇 \\
		発行者 && 〇〇〇〇 &	\\
		発行所 &&  琉和辞典研究所 &	\\
		&& 〒190-0014\hspace{2pt}東京都立川市緑町10 - 2\\
		&& TEL\hspace{2pt}0555-555-555　FAX\hspace{2pt}055-555-5555\\
		印刷所 && 〇〇〇〇 & \\
		ISBN && ISBN 〇〇〇〇\\
		&& C3581\\
		\hline
		\multicolumn{4}{l}{\textcopyright \hspace{2pt}2020, 〇〇〇〇 \& 〇〇〇〇\hspace{54pt}Printed in Japan}
	\end{tabular}%
\end{table}%

\end{document}

